% !TEX root = relative_reconstruction_models.tex
\documentclass[a4paper,11pt]{ltjsarticle}
\usepackage[margin=25mm,headheight=16pt]{geometry}
\usepackage{amsmath,amssymb,amsthm,mathtools,bm,mathrsfs}
\usepackage{booktabs,longtable,array,enumitem}
\usepackage{xcolor}
\usepackage[unicode,colorlinks=true,linkcolor=blue!40!black,citecolor=blue!40!black,urlcolor=blue!45!black]{hyperref}
\usepackage{fancyhdr}
\usepackage{microtype}
\setlength{\parindent}{1em}
\setlength{\parskip}{0.25em}
\setlist{itemsep=0.2em,topsep=0.4em}
\allowdisplaybreaks[2]
\pagestyle{fancy}
\fancyhf{}
\fancyhead[L]{\small 相対復元・有限対角・明示模型}
\fancyhead[R]{\small 2026-11-13 / v1}
\fancyfoot[C]{\thepage}
\newcommand{\CC}{\mathbb C}
\newcommand{\ZZ}{\mathbb Z}
\newcommand{\PP}{\mathbb P}
\newcommand{\OO}{\mathcal O}
\newcommand{\cK}{\mathcal K}
\newcommand{\cP}{\mathcal P}
\newcommand{\cD}{\mathcal D}
\newcommand{\cC}{\mathcal C}
\newcommand{\cE}{\mathcal E}
\newcommand{\cF}{\mathcal F}
\newcommand{\cM}{\mathcal M}
\newcommand{\fm}{\mathfrak m}
\newcommand{\Perf}{\operatorname{Perf}}
\newcommand{\Ku}{\operatorname{Ku}}
\newcommand{\MF}{\operatorname{MF}}
\newcommand{\HH}{\operatorname{HH}}
\newcommand{\QGr}{\operatorname{QGr}}
\newcommand{\qgr}{\operatorname{qgr}}
\newcommand{\Dsg}{D_{\mathrm{sg}}}
\newcommand{\RHom}{\operatorname{RHom}}
\newcommand{\REnd}{\operatorname{REnd}}
\newcommand{\Hom}{\operatorname{Hom}}
\newcommand{\Ext}{\operatorname{Ext}}
\newcommand{\Ind}{\operatorname{Ind}}
\newcommand{\Def}{\operatorname{Def}}
\newcommand{\thick}{\operatorname{thick}}
\newcommand{\Spec}{\operatorname{Spec}}
\newcommand{\Spf}{\operatorname{Spf}}
\newcommand{\Id}{\operatorname{Id}}
\newcommand{\sym}{\operatorname{sym}}
\newcommand{\Cone}{\operatorname{Cone}}
\newcommand{\rank}{\operatorname{rank}}
\newcommand{\Pic}{\operatorname{Pic}}
\newcommand{\otL}{\mathbin{\otimes^{\mathbf L}}}
\newcommand{\DocDate}{2026年11月9日}
\newcommand{\DocVersion}{論文分割版 v6}
\theoremstyle{plain}
\newtheorem{theorem}{定理}[section]
\newtheorem{proposition}[theorem]{命題}
\newtheorem{lemma}[theorem]{補題}
\newtheorem{corollary}[theorem]{系}
\theoremstyle{definition}
\newtheorem{definition}[theorem]{定義}
\newtheorem{problem}[theorem]{問題}
\newtheorem{example}[theorem]{例}
\theoremstyle{remark}
\newtheorem{remark}[theorem]{注意}
\renewcommand{\proofname}{証明}
\hypersetup{pdfauthor={},pdfsubject={3次超曲面のPoisson構造・非可換量子化・rotation付き残余圏}}

\setlength{\emergencystretch}{3em}
\numberwithin{equation}{section}
\newcommand{\Sym}{\operatorname{Sym}}

\newcommand{\gr}{\operatorname{gr}}
\newcommand{\dd}{^{\vee}}
\newcommand{\LL}{\mathbf L}
\newcommand{\RR}{\mathbf R}
\theoremstyle{definition}
\newtheorem{hypothesis}[theorem]{仮定}
\newcommand{\C}{\mathbb C}
\newcommand{\Gr}{\operatorname{Gr}}

\hypersetup{pdftitle={非可換三次対の相対復元と有限対角・明示模型}}
\title{非可換三次対の相対復元と有限対角・明示模型}
\author{}
\date{2026年11月13日・研究草稿 v1}
\begin{document}
\maketitle
\begin{abstract}
非可換三次対の一般点復元とは独立に、Artin形式変形とétale局所実現、
有限パラメータの楕円ブロック、多ブロックの生成性、曲率付き有限対角とbar積、
片側Morita核の代数化、混合量子化のjetを扱う。
形式的存在、指定有限jetの代数化、実際の代数的圏の族の存在を区別し、
各比較の入力と残る障害を明記する。
\end{abstract}
\section*{各部の関係と既知入力}
本稿は五つの関連した問題を扱う。第I部は一般の局所Artin基底と、
実在するsmooth properな代数的族を区別する。
第II部は分離した楕円ブロックの有限パラメータにおける明示模型を構成し、
二つ以上の無限位数因子についての生成性を未解決のまま区別する。
第III部は一変数の全有限切断におけるMF同値の存在を構成し、
その存在定理とArtin相対windowを合成して可換残余pairを得る。
比較同値を仮定するだけの条件付き移送で、この全次数存在の主張を置き換えない。
第IV部は対角・bar・sector・核の乗法性と底変更を扱う。
第V部は一般点への適用と別構成を整理する。
これらは一般点主定理と同じ結果を別の新成果として数えるものではない。

一般次元のPoisson分類・量子化の正本はnc-cubicの論文A、
一般点可換復元とその必要証明の正本は同論文Bである。
本稿の各命題の適用条件、全coherent圏と線束生成圏の違い、
相対smoothnessの仮定をそれぞれの定理文で保持する。
基礎となるformality、MF比較、Hochschild計算、対象moduliは引用した既知入力である。
本稿に証明を記載することと、それら全ての独立監査、新規性の確定を区別する。
先行研究上の初出と全証明の独立監査は未了である。
\tableofcontents

\part{Artin形式変形とétale局所実現}
\section{Artin形式復元の定式化と結果の概要}\label{sec:intro}

滑らかなcubicのKuznetsov成分を
\[
 D^b(X)=\langle\Ku(X),\OO_X,\OO_X(1),\ldots,\OO_X(n-2)\rangle
\]
の規約で取る。rotationは
\[
 \rho_X=\mathbf L_{\OO_X}\circ(-\otimes\OO_X(1))
\]
である。すべての圏はdg enhancementと冪等完備化を保って扱う。
可換cubicについては、\((\Ku(X),\rho_X)\) が \(X\) の射影同型類を決定する
\cite[Theorem 1.7]{LZ}。このTorelli定理は可換モデルの一意性を与えるが、
任意の非可換な圏の組がその像に属することを述べない。
幾何側のenhanced同値は、包含・射影と合成してFourier--Mukai核で表されるため、
同定理のFourier--Mukai条件を満たす。
ここでは、主定理\ref{gdr:thm:generic}とは
基底と仮定の異なる二種類の結果を先に整理する。
第一にArtin局所基底上の標識付き形式復元、
第二にsmooth properな代数的dg圏のétale局所実現である。
有限係数の具体模型と楕円ブロックの主定理は
第\ref{sec:explicit-block-results}節以降で扱う。
形式的な存在だけから一般の有限パラメータでの存在が
自動的に従うわけではない。

\subsection{Artin基底での標識付き形式復元}

\begin{theorem}[次数付き変形の形式的可換モデル]\label{thm:main}
 \(X=V(F)\subset\PP^{n+1}_{\CC}\) を滑らかなcubic、\(n\ge3\) とする。
 \(R\) を剰余体が \(\CC\) の有限次元可換局所Artin代数とし、
 \(B_R=\bigoplus_{d\ge0}(B_R)_d\) を中心的な \(R\)-代数で、
 \((B_R)_0=R\)、各次数部分が有限自由、かつ指定された次数付き同型
 \[
 B_R\otimes_R\CC\simeq\CC[x_1,\ldots,x_{n+2}]/(F)
 \]
 を持つものとする。次を定義する。
 \[
 \cP_R=D(\QGr B_R)^c,\qquad
 \cK_R=\langle\OO_R,\ldots,\OO_R(n-2)\rangle^\perp\subset\cP_R,\qquad
 \rho_R=\mathbf L_{\OO_R}\circ(1).
 \]
 このとき \((\cK_R,\rho_R)\) はsmooth properな \(R\)-linear dg圏と自己同値の組である。
 あるmarkedな滑らかなcubic変形
 \(X_R\subset\PP^{n+1}_R\) と、markingを保つ同値
 \[
 (\cK_R,\rho_R)\simeq(\Ku(X_R),\rho_{X_R})
 \]
 が存在する。\(X_R\) のmarked射影同型類は一意である。
 これらはArtin底の変更と両立する。
\end{theorem}

ここで \(\QGr\) は全次数付き加群をtorsion加群で割ったGrothendieck圏であり、
上付き \(c\) はcompact対象を表す。Artin底上では
\(\cP_R\) を \(D^b(\qgr B_R)\) 全体に置き換えない。
例えば可換な定数族でも、底 \(R\) が非正則なら両者は異なる。
後述する体上の有限パラメータでは、残余圏を通常のgraded matrix factorization圏、
すなわちOrlov同値で \(D^b(\qgr B)\) に埋め込まれる成分として扱う。
特異なファイバーで、この成分と \(D(\QGr B)^c\) の直交成分を混同しない。

\begin{corollary}[形式族]\label{cor:formal}
 \(B_\hbar\) が次数ごとに \(\CC[[\hbar]]\)-flatな、定理の斉次座標環の形式的変形なら、
 \[
 F_\hbar\in\CC[[\hbar]][x_1,\ldots,x_{n+2}]_3,\qquad F_\hbar\bmod\hbar=F
 \]
 が存在し、各 \(\CC[\hbar]/(\hbar^{r+1})\) 上で残余圏の組を
 \(V(F_\hbar)\) の組と同一視できる。これらの同値は互換的に選べる。
 形式的cubicのmarked同型類は一意である。
\end{corollary}

これは形式的な主張である。係数の収束や、指定された代数的パラメータ空間への
降下、任意の非零パラメータでの可換モデルの存在を含まない。
また \(F_\hbar=F\) を意味しない。

\subsection{smooth properな代数的族のétale局所復元}

定理\ref{thm:etale}は、有限型複素底 \(S\) 上のsmooth properな代数的dg圏の族
\(\cK_S\) と、可逆perfect bimoduleで表された自己同値 \(\rho_S\) を入力とする。
ある点 \(s\) で \((\cK_s,\rho_s)\simeq(\Ku(X),\rho_X)\) なら、
\(s\) のétale近傍 \(T\) 上に滑らかなcubic \(Y/T\) と
\[
 (\cK_T,\rho_T)\simeq(\Ku(Y/T),\rho_{Y/T})
\]
が存在する。従って、可換cubicで実現される複素点はZariski開部分をなす。
これはgraded flatな代数の族から、相対smooth properなdg模型を作る課題を
省略するものではない。

また \(n\ne4\) では \(\HH^2(\Ku(X))=J(F)_3\)、\(\HH^1=0\) であり、
rotationを忘れた形式的復元も成立する（定理\ref{thm:unpolarized-formal}）。
同じ条件下では、代数的dg圏の族についてもrotationを最初から指定せずに
étale局所モデルを作れる（系\ref{cor:etale-unpolarized}）。
これらは局所的な実現の結果であり、rotationなしの大域的な一意性とは分けて扱う。

\subsection{形式的残余圏の比較に用いる既知結果}

Orlov \cite{Orlov} とCassidy--Conner--Kirkman--Moore \cite{CCKM} は、
中心正則元による商、特異圏、matrix factorization、残余成分の比較を与える。
体上のwindowのdg模型にはCăldăraru--Tu \cite[Theorem 5.8]{CT} を用いる。
本稿の形式的な構成では、これらの体上の同値をArtin底上へ無条件に延長せず、
De Deken--Lowen \cite{DDL} のabelian deformationとcompact generationを用いる。
smoothnessとpropernessの持ち上げはnilpotent descent
\cite{Mathew}による。

残余圏のHochschild cohomologyはBallard--Favero--Katzarkov \cite{BFK}、
幾何学的Kodaira--Spencer写像はHuybrechts--Rennemo \cite{HR} と、
三次元の場合にはJacovskis--Lin--Liu--Zhang \cite{JLLZ} を用いる。
四次元の余分な二方向に対するrotationの作用は
Huybrechts \cite[Proposition 3.12]{Huy} から確認する。
actual dg変形をHochschild形式変形と同一視する際のcurvature問題
\cite{KL}は、Lurie \cite{Lurie} とBlanc--Katzarkov--Pandit
\cite{BKP} の比較写像の単射性で処理する。

\paragraph{形式的結果の射程と証明の所在.}
以下ではArtin形式基底における残余dg圏と標識付き同値を構成する。
定理\ref{thm:main}の証明は、相対残余圏の構成（第\ref{sec:relative}節）と
rotation付き形式変形の比較（第\ref{sec:formal}節）からなる。
これらの組合せの先行研究上の初出は未確認であり、
基礎となる既知定理の新規性は主張しない。
有限係数の可換模型については、第\ref{sec:explicit-block-results}節で
別の仮定の下に述べる。

\section{次数付き変形から相対残余圏を構成する}\label{sec:relative}

以下では、Artin底上で用いる残余dg圏を、
通常の有界coherent導来圏と区別して構成する。
特に、compact対象の底変更、相対例外列とrotationの順に議論する。

\subsection{nilpotentな底変更によるcompactnessの検出}

\begin{lemma}\label{lem:descent}
 \(R\) を上の局所Artin代数、\(k=R/\fm\) とする。
 \(A\) を \(R\)-linear dg代数、\(M\) を \(A\)-加群とし、
 \(A_0=A\otL_R k\)、\(M_0=M\otL_R k\) と書く。
 \(M_0\) が \(A_0\)-compactなら \(M\) は \(A\)-compactである。
 また導来底変更 \(-\otL_R k\) は保守的である。
\end{lemma}
\begin{proof}
 \(\fm\) の冪による有限濾過から
 \(R\in\thick_{D(R)}(k)\) である。
 \(A\)-加群の任意の族 \((N_\alpha)\) に対し、
 \[
 \bigoplus_\alpha\RHom_A(M,P\otL_RN_\alpha)
 \longrightarrow
 \RHom_A\!\left(M,P\otL_R\bigoplus_\alpha N_\alpha\right)
 \]
 が同値となる \(P\in D(R)\) の全体はthickである。
 \(P=k\) では導来随伴と \(M_0\) のcompactnessによって同値となるので、
 \(P=R\) でも同値となる。保守性も、
 \(P\otL_RM=0\) となる \(P\) のthick部分圏を使って同様に従う。
 これは \cite[Propositions 3.28, 3.35; Remark 3.29]{Mathew}
 の特別な場合である。
\end{proof}

\begin{remark}
 この補題は \(k\) が \(R\)-perfectであるとは仮定していない。
 また普通の商 \(A/\fm A\) だけを調べる命題ではなく、
 必ず導来ファイバー \(A\otL_Rk\) を使う。
\end{remark}

\subsection{abelian deformationと明示的な生成子}

\begin{proposition}\label{prop:qgr}
 定理\ref{thm:main}の \(B_R\) について、
 \(\QGr B_R\) は \(\operatorname{Qcoh}(X)\) のflat abelian deformationである。
 \[
 E_R=\bigoplus_{i=0}^{n+1}\OO_R(i)
 \]
 は \(D(\QGr B_R)\) のcompact generatorであり、
 \[
 \Lambda_R=\REnd(E_R),\qquad
 \Lambda_R\otL_Rk\simeq\REnd_X(E_0)
 \tag{\(\ast\)}\label{eq:derived-basechange}
 \]
 がdg代数として成り立つ。
\end{proposition}
\begin{proof}
 まず \(B_R\) は左右noetherianである。
 \(\fm^aB_R/\fm^{a+1}B_R\) は、
 flatnessにより \(B_0\) の有限個の直和である。
 有限濾過と \(B_0\) のnoetherian性を使えばよい。
 次数付き生成元を持ち上げる同じ議論から、
 \(B_R\) は有限生成で、付随する \(\ZZ\)-代数のtailsも通常のtorsion商で記述できる。

 \(X\) の射影正規性により \(B_0\) の \(\ZZ\)-代数は
 \((\OO_X(i))_{i\in\ZZ}\) のdirected Hom-categoryである。
 この列は \(\ZZ\)-generatingであり、\(i\le j\) に対して
 \[
 \Ext_X^{1,2}(\OO_X(i),\OO_X(j))
 =H^{1,2}(X,\OO_X(j-i))=0.
 \]
 hypersurfaceの完全列と射影空間のcohomologyからこの消滅が従う。
 よって \cite[Theorem 4.5]{DDL} を適用できる。
 \(\OO_R(i)\) は明示的なflat liftである。
 実際 \(B_R(i)\) との有限 \(R\)-加群のテンソル積は完全で、
 \(\operatorname{Gr}B_R\to\QGr B_R\) は完全である。
 これらの対象はabelian圏も生成するので、必要なflat resolutionを取れる。

 閉埋め込み \(j:X\hookrightarrow\PP^{n+1}\) を考える。
 \(E_0=j^*(\bigoplus_{i=0}^{n+1}\OO_{\PP^{n+1}}(i))\) はcompactである。
 \(\RHom_X(E_0,M)=0\) なら、随伴とBeilinsonの生成定理
 \cite{Beilinson}から \(Rj_*M=0\)、従って \(M=0\) となる。
 したがって \(E_0\) はcompact generatorである。

 \(E_R\) は有界flat liftなので、
 \cite[Propositions 5.8, 5.9]{DDL} によりcompact generatorとなる。
 同論文Proposition 5.5と導来随伴は
 \[
 k\otL_R\RHom(E_R,E_R)
 \simeq\RHom_X(k\otL_RE_R,k\otL_RE_R)
 \]
 を与える。同論文Theorem 5.13の証明でも同じ底変更が使われている。
 dg enhancement上の導来還元関手からこの比較を作るため、
 比較は合成と単位を保つ。必要ならhomotopically \(R\)-flatなdg模型を選ぶ。
 これで \eqref{eq:derived-basechange} は複体だけでなくdg代数として成立する。
\end{proof}

\begin{corollary}\label{cor:smooth-proper}
 \(\cP_R=D(\QGr B_R)^c\) はsmooth properであり、
 \(\cP_R\otL_Rk\simeq\Perf(X)\) である。
\end{corollary}
\begin{proof}
 導来Morita理論 \cite{Keller} によって
 \(\cP_R\simeq\Perf(\Lambda_R)\) である。
 \(X\) のpropernessから \(\Lambda_0\) は \(k\)-perfectなので、
 補題\ref{lem:descent}で \(\Lambda_R\) は \(R\)-perfectとなる。
 次に
 \[
 A=\Lambda_R^{\mathrm{op}}\otL_R\Lambda_R,\qquad M=\Lambda_R
 \]
 とする。この導来ファイバーは \(\Lambda_0\) のenveloping algebraと対角加群である。
 \(X\) のsmoothnessにより対角加群はその上でperfectであり、
 同じ補題から \(\Lambda_R\) の対角もperfectとなる。
\end{proof}

\subsection{例外列、mutation、rotation}

\begin{proposition}\label{prop:relative-pair}
 \(\OO_R,\ldots,\OO_R(n-2)\) は \(\cP_R\) 内のrelative exceptional collectionである。
 その右直交 \(\cK_R\) はsmooth properで、導来特殊ファイバーは \(\Ku(X)\) である。
 \(\rho_R=\mathbf L_{\OO_R}\circ(1)\) は \(\cK_R\) の自己同値となる。
 これらの構成はArtin底変更と両立する。
\end{proposition}
\begin{proof}
 Hom複体の導来底変更と補題\ref{lem:descent}から
 \[
 \RHom(\OO_R(i),\OO_R(i))\simeq R,\qquad
 \RHom(\OO_R(j),\OO_R(i))=0\quad(j>i,\ 0\le i,j\le n-2)
 \]
 である。自己Homについては単位写像のconeを還元する。
 propernessにより評価・余評価はperfect対象を保つため、各例外成分はadmissibleである。
 評価のconeによるmutationを順に取って
 \[
 p_R=\mathbf L_{\OO_R}\cdots\mathbf L_{\OO_R(n-2)}
 \]
 と置くと、\(G_R=p_R(E_R)\) は \(\cK_R\) を生成する。
 mutationの導来底変更は \cite[Lemma 5.14]{DDL} による。
 \(G_R\) の導来特殊ファイバーは \(\Ku(X)\) の生成子なので、
 その自己準同型dg代数へ補題\ref{lem:descent}を適用すればsmooth proper性も従う。

 \(T=(1)\) は次数付き射影圏の自己同値で、\(T\OO_R(i)=\OO_R(i+1)\) である。
 \(M\in\cK_R\)、\(1\le i\le n-2\) に対して
 \[
 \RHom(\OO_R(i),TM)=\RHom(\OO_R(i-1),M)=0.
 \]
 \(\OO_R\) によるmutationはこの消滅を保ち、\(i=0\) の直交性も与える。
 よって \(\rho_R\) は \(\cK_R\) を保つ。
 その特殊ファイバーは通常の \(\rho_X\) である。
 Homの比較写像は還元後に同値なので、nilpotentな保守性から
 \(\rho_R\) は完全忠実である。

 \(\rho_R(G_R)\) の特殊ファイバーは生成子である。
 これに直交するInd対象を導来還元すると零になり、
 同じ保守性で元の対象も零になる。従って \(\rho_R(G_R)\) もcompact generatorである。
 冪等完備な圏からの完全忠実完全関手の像はthickなので、\(\rho_R\) は本質的全射でもある。

 底変更については、次数付き導来テンソル関手による生成子の自己Homの自然な比較を使う。
 剰余体上でこの比較は恒等写像となるため、変更後のArtin底上で同値となる。
 評価写像とそのconeも底変更と両立するので、残余圏とrotationまで両立する。
\end{proof}

\section{rotation付き残余圏の形式的可換復元}\label{sec:formal}

前節の相対残余圏に対し、変形関手とHochschild接写像を比較し、
Artin帰納法によって標識付き模型を構成する。
rotationを忘れた場合の結論は次元によって異なるため、
同じ変形問題として扱わない。

\subsection{正確な変形問題}

\begin{definition}
 \((\Ku(X),\rho_X)\) のmarked pair deformationとは、
 smooth properな \(R\)-linear dg圏 \(\cK_R\)、invertible bimoduleで表される
 自己同値 \(\rho_R\)、および特殊ファイバーでのenhanced pairの指定された同定からなる。
 同値にはこのmarkingとの両立を要求する。
 \(\rho_R^3\simeq[2]\) の指定されたcoherenceは定義に含めない。
\end{definition}

\begin{theorem}[pairの形式的復元]\label{thm:pair}
 \(X\) を滑らかなcubic \(n\)-fold、\(n\ge3\) とする。
 marked cubic deformationからmarked pair deformationへの写像は、
 各古典的局所Artin \(\CC\)-代数上で同型類の全単射を与える。
 互換的な形式族についても同型類の一意性と互換的な実現が成り立つ。
\end{theorem}

この主張はfull higher deformation groupoidの同値ではない。
pairには自然変換のscalar automorphismがあり、実現する圏同値も一意とは限らない。

\subsection{actual変形とHochschild形式変形}

\(\cC=\Ind(\Ku(X))\) とする。
actual compactly generated deformationの空間と、その形式moduli problemの間には
\[
 \theta_{\cC}:\operatorname{CatDef}_{\cC}^{c}
 \longrightarrow\operatorname{CatDef}_{\cC}^{\wedge}
\]
がある。右辺の接複体は \(CC^\bullet(\Ku(X))[2]\) である。
actual Morita変形が任意のHochschild classを実現するとは一般にはいえない
\cite{KL}。本稿ではその全射性を使わない。

\begin{lemma}\label{lem:theta}
 上の \(\theta_{\cC}\) は \((-1)\)-truncatedである。
 従って、右辺で同値となった二つのactual変形の同値はactual側へ戻せる。
 この性質はrotationによる \(\ZZ\)-homotopy fixed pointを取った後にも成り立つ。
\end{lemma}
\begin{proof}
 \(\Ku(X)\) はproperでcompact generatorを持つためtameである。
 \cite[Proposition 5.3.21, Theorem 5.1.9]{Lurie} を適用する。
 \cite[Definition 3.3, Section 3.3]{BKP} にも、
 \(\pi_0\) で単射、\(\pi_i\), \(i\ge1\), で同型という形でこの帰結が明記されている。
 比較写像は圏同値に自然であり、markingを移すrotationの作用と両立する。
 空間のmonomorphismはhomotopy limitで保たれる。
\end{proof}

一つのenhanced autoequivalenceを持ち上げる問題は、
markingに作用する \(\ZZ\) のhomotopy fixed pointである。
\(\ZZ\) は一つの可逆射で自由生成されるので、追加の有限群関係を課していない。
\[
 P=\left(\operatorname{CatDef}_{\cC}^{\wedge}\right)^{h\ZZ}
\]
と置くと、これはformal moduli problemであり、
\[
 T_P=R\Gamma(\ZZ,CC^\bullet(\Ku(X))[2]).
\]
\(\ZZ\) の自明加群の長さ一の自由分解から、次の完全列を得る。
\[
 0\longrightarrow
 \operatorname{coker}(1-\rho:\HH^1\to\HH^1)
 \longrightarrow H^0(T_P)
 \longrightarrow
 \ker(1-\rho:\HH^2\to\HH^2)
 \longrightarrow0.                                      \tag{1}\label{eq:tangent}
\]

\subsection{Hochschild計算とKodaira--Spencer写像}

\(N=n+2\) および
\[
 J(F)=\CC[x_1,\ldots,x_N]/(\partial_1F,\ldots,\partial_NF)
\]
を用いる。

\begin{lemma}\label{lem:HH}
 \(n\ge3\) のとき
 \[
 \HH^{<0}(\Ku(X))=0,\quad \HH^0(\Ku(X))=\CC,\quad
 \HH^1(\Ku(X))=0,
 \]
 \[
 \HH^2(\Ku(X))=
 \begin{cases}
 J(F)_3& n\ne4,\\
 J(F)_3\oplus\CC^2& n=4.
 \end{cases}
 \]
 幾何学的Kodaira--Spencer写像は同型
 \[
 H^1(X,T_X)\xrightarrow{\sim}
 \HH^2(\Ku(X))^{\rho_X}=J(F)_3
 \]
 を与える。その次元は \(\binom{n+2}{3}\) である。
\end{lemma}
\begin{proof}
 \cite[Theorem 5.39]{BFK} のKoszul cohomologyによる公式を
 三次のequivariant matrix factorizationへ適用する。
 恒等sectorは偶数次数 \(2r\) に \(J(F)_{3r}\) を与える。
 非自明な \(g\in\mu_3\) では固定点集合は原点だけであり、
 conormalのrankは \(N\)、determinant characterは \(-N\) である。
 cohomological degree \(q\) で残るweightは
 \[
 N+3(q-N)/2=(3q-N)/2.
 \]
 従って二つの非自明sectorは、\(3\mid N\) のときに限り、
 \(q=N/3\) に各一次元を与える。これで列記した群が得られる。

 BFKの恒等sectorのKoszul模型では、\(G\) の乗法が \([G]\) を表す。
 \(F+\varepsilon G\) の実際の変形類との一致は次の比較による。
 \(n\ge4\) では、periodicityから作るkernel morphismと
 Atiyah--Kodaira--Spencer morphismの一致を
 \cite[Lemma 2.19]{HR} とその前後の構成で確認できる。
 ここで三次に必要な次数条件 \(3\le(n+2)/2\) は満たされる。
 \(n=3\) では \cite[Theorem 4.4(2)]{JLLZ} が必要な接写像の同型を与える。
 normal sequenceとEuler sequenceは \(H^1(T_X)=J(F)_3\) を与える。

 幾何学的変形はrotationを持つので、この部分空間は不変である。
 \(n\ne4\) ではそれが \(\HH^2\) 全体である。
 \(n=4\) では \cite[Proposition 3.12]{Huy} により、
 rotationはMukai格子の \(A_2^\perp\) を固定し、\(A_2\) の三つの根を巡回置換する。
 invariantな2-Calabi--Yau volumeによる
 \(\HH^2\simeq\HH_0\) の下で、余分な二方向の固有値は
 \(\zeta,\zeta^2\) となり、不変部分はない。
 最後に偏微分は二次の正則列であるから
 \(H_{J(F)}(t)=(1+t)^N\) であり、次数三の次元は \(\binom N3\) である。
\end{proof}

\subsection{持ち上げの差と自己同型}

\begin{lemma}[Artin帰納法]\label{lem:induction}
 \(D\) をsmoothかつinfinitesimally homogeneousなmarked変形関手で、
 古典的Artin底上ではset-valuedとし、
 \(P\) をformal moduli problemとする。
 自然な写像 \(f:D\to P\) が接空間で同型を誘導すると仮定する。
 さらに、各small extension \(R'\to R\) と \(x'\in D(R')\) に対し、
 \(f(x'\vert_R)\) の任意の自己同型類が \(f(x')\) 上へ持ち上がると仮定する。
 このとき \(D(R)\to\pi_0P(R)\) は各古典的Artin底上で全単射である。
\end{lemma}
\begin{proof}
 Artin長について帰納法を行う。核 \(I\) のsmall extension上に
 \(y'\in P(R')\) があり、その還元を \(f(x)\) と同一視できたとする。
 \(D\) のsmoothnessにより \(x\) を \(x'\) に持ち上げる。
 固定した還元の同定を持つ \(f(x')\) と \(y'\) は、
 \(P(R')\to P(R)\) の同じhomotopy fibreに属する。
 このfibreが空でなければ、その連結成分は
 \(H^0(T_P)\otimes I\) のtorsorである。
 接写像の全射性によって \(x'\) を修正すれば差を消せる。

 単射性では、\(f(x')\simeq f(z')\) の還元を考える。
 帰納法により幾何学的還元は一致するが、
 圏同値の還元はその像の自己同型となり得る。
 仮定によってこの自己同型を上へ持ち上げ、その逆を合成する。
 すると還元を恒等に固定したfibreで二つのliftを比較できる。
 差の接ベクトルの像は零であり、接写像の単射性で元の差も零となる。
\end{proof}

\begin{lemma}\label{lem:automorphisms}
 cubicから得たpairのmarked deformationにおいて、
 small extensionに沿う上の自己同型持ち上げ条件が成り立つ。
\end{lemma}
\begin{proof}
 smooth properな圏ではHochschild複体はperfectで、その導来底変更が成り立つ。
 補題\ref{lem:HH}と局所Artin環上のminimal perfect complexを使うと
 \[
 \HH^0(\cK_R/R)=R,\qquad \HH^1(\cK_R/R)=0
 \]
 である。実際特殊ファイバーの負次数と次数一が零なので、
 minimal complexにはそれらの次数の項がなく、次数零は単位で生成される。

 恒等関手に還元されるinvertible bimoduleは相対対角と同型である。
 perfect bimoduleのliftの差はsmall extensionの核 \(I\) に対して
 \(\Ext^1(\Delta,\Delta)\otimes I=\HH^1(\Ku(X))\otimes I=0\)
 で測られ、対角自身がliftを与えるためである。
 この同型は各段階で指定した特殊ファイバーのmarkingを保つように選ぶ。
 従ってpairの自己同型の基礎となる圏同値は恒等と同型であり、
 残るのは \(\rho_R\Rightarrow\rho_R\) のscalar unitである。
 markingを固定すればこのscalarは \(1+\fm_R\) に入り、必ず \(R'\) へ持ち上がる。
 scalarな二次自己同型も持ち上がる。
 formal moduli problem側の自己同型をactual bimoduleで扱えることは、
 補題\ref{lem:theta}の \(\pi_1\) の同型による。
\end{proof}

\begin{proof}[定理\ref{thm:pair}の証明]
 cubicのmarked変形関手はsmoothである。方程式の係数を持ち上げればよく、
 \(H^1(\OO_X)=H^2(\OO_X)=0\) によって \(\OO_X(1)\) は一意に延長する。
 infinitesimalな射影自己同型もない。実際、二次の偏微分の正則列には
 一次係数の非自明なsyzygyがなく、Euler scalarを除けば \(H^0(T_X)=0\) である。
 これはmarkedな幾何学的変形をset-valuedとして扱うことを正当化する。

 式\eqref{eq:tangent}と補題\ref{lem:HH}により接写像は同型である。
 補題\ref{lem:automorphisms}と補題\ref{lem:induction}を適用し、
 最後に補題\ref{lem:theta}でactual pairの同値へ戻す。
 帰納法では各段階で指定した還元の同値を保てるので、
 \(\CC[\hbar]/(\hbar^{r+1})\) に沿った実現も互換的に選べる。
\end{proof}

\begin{proof}[定理\ref{thm:main}と系\ref{cor:formal}の証明]
 命題\ref{prop:relative-pair}の組へ定理\ref{thm:pair}を適用する。
 互換的な幾何学的liftの方程式の係数の逆極限が \(F_\hbar\) を与える。
 この最後の操作は形式族を作るものであり、非零複素数で評価する操作ではない。
\end{proof}

\begin{corollary}\label{cor:period}
 上で構成したArtin pairについて、指定を伴わない同型
 \(\rho_R^3\simeq[2]\) が存在する。
\end{corollary}
\begin{proof}
 \(\rho_R^3[-2]\) は恒等に還元されるinvertible bimoduleである。
 補題\ref{lem:automorphisms}で使った \(\HH^1\)-rigidityを適用する。
 または可換モデルからこの同型を移してもよい。
 特定のmatrix factorization表示が与えるperiodicity同型とのcoherenceの一致は、
 この系の主張に含めない。
\end{proof}

\begin{remark}
 sevenfoldでは \(\HH^3(\Ku(X))\simeq\CC^2\) があるが、
 上の帰納法は \(\HH^3=0\) を要求しない。
 比較している二つのliftは既に存在し、その差を接空間で修正する。
 higher operationsを無視したformalityの仮定も使っていない。
\end{remark}

\subsection{四次元以外ではrotationを忘れた形式復元}\label{sec:unpolarized-formal}

四次元以外では、形式近傍に限れば、復元の入力からrotationを省くことができる。
ここでいうmarked dg変形は、特殊ファイバーを \(\Ku(X)\) と同定した
smooth properな \(R\)-linear dg圏であり、同値にもこのmarkingとの両立を要求する。

\begin{theorem}[rotationを忘れた形式復元]\label{thm:unpolarized-formal}
 \(X\subset\PP^{n+1}_{\CC}\) を滑らかなcubicとし、\(n\ge3\)、\(n\ne4\) とする。
 任意の古典的局所Artin \(\CC\)-代数 \(R\) について、
 marked cubic変形から \(\Ku(X)\) のmarked smooth proper dg変形への写像は、
 同型類の全単射を与える。従って、任意のそのような \(\cK_R\) は
 一意なmarked cubic変形 \(X_R\) の \(\Ku(X_R)\) と同値である。

 また、固定した \(\cK_R\) 上で、\(\rho_X\) を持ち上げるinvertible bimodule
 \(\rho_R\) が存在し、そのmarked同型類は一意である。
 従ってrotationを忘れる写像
 \[
  \{\text{marked smooth proper pair変形}\}/\simeq
  \ \longrightarrow\
  \{\text{marked smooth proper dg変形}\}/\simeq
 \]
 も全単射である。これらの主張はArtin底変更と互換的な形式族についても成り立つ。
\end{theorem}
\begin{proof}
 \(Q=\operatorname{CatDef}_{\Ind\Ku(X)}^{\wedge}\) と置く。
 補題\ref{lem:HH}の幾何学的Kodaira--Spencer写像は、\(n\ne4\) では
 \[
  H^1(X,T_X)\xrightarrow{\sim}
  \HH^2(\Ku(X))=H^0(T_Q)
 \]
 である。cubicのmarked変形関手はsmoothかつset-valuedである。
 補題\ref{lem:induction}を適用するため、幾何学的な変形の像の自己同型を調べる。

 \(\Ku(X_R)\) の恒等に還元される自己同値を表すinvertible bimoduleは、
 \(\HH^1(\Ku(X))=0\) により相対対角bimoduleと同型である。
 これは補題\ref{lem:automorphisms}と同じsmall extensionの帰納法で従う。
 特殊ファイバー上の同型を指定しても、そこでのscalarは \(R^\times\) へ持ち上がるので、
 同型をmarkingと両立するように選べる。
 従ってmarkedな自己同値の同型類は恒等だけであり、
 必要な自己同型持ち上げ条件は満たされる。
 \(Q\) 側の自己同型についても、補題\ref{lem:theta}の比較写像が
 \(\pi_1\) 上の同型を誘導するため同じことが成り立つ。

 補題\ref{lem:induction}から、cubic変形と \(\pi_0Q(R)\) の全単射が得られる。
 幾何学的変形の像はactual smooth proper変形であり、
 補題\ref{lem:theta}によりactual変形間の同値を検出できる。
 よって最初の主張が従う。この議論はHochschild形式変形の一般的なactual実現性を
 仮定せず、この場合には幾何学的なliftで全てを実現している。

 得られた同値 \(\Ku(X_R)\simeq\cK_R\) を通じて
 \(\rho_{X_R}\) を移せば、\(\rho_X\) のliftが存在する。
 二つのmarkedなlift \(\rho_R,\rho'_R\) があれば、
 \(\rho'_R\rho_R^{-1}\) は恒等に還元されるinvertible bimoduleである。
 再び \(\HH^1=0\) による対角のrigidityを用いると、
 \(\rho'_R\rho_R^{-1}\simeq\mathrm{id}\) となり、markingを保つ同型も選べる。
 これがliftの同型類の一意性とforgetful mapの単射性を与える。
 各small extensionで還元上の同値を保って構成できるため、
 形式族での互換性も従う。
\end{proof}

\begin{remark}\label{rem:unpolarized-scope}
 この定理の一意性は同型類についてのものである。
 rotationとのintertwinerにはscalar自由度があり、
 rotationを忘れる写像がfull higher deformation groupoidの同値だとは主張しない。
 実際、\(\HH^0=\CC\)、\(\HH^1=0\) とrotationのscalar上の自明作用から、
 pairの接複体には \(H^{-1}(T_P)\simeq\CC\) がある一方、
 \(H^{-1}(T_Q)=0\) である。

 \(n=4\) では幾何学的な20方向に加えて \(\HH^2\) に二方向があり、
 rotationの固定部分を取って初めて接写像が同型となる。
 この次元でrotationを省く結論を上の証明から導くことはできない。
 ここでは、Hochschildの全22方向が各Artin底上でactual変形として
 実現されるとは主張していない。

 \(n=7\) では \(\HH^3\simeq\CC^2\) が残るが、
 幾何学的なliftと既に存在するliftを比較する帰納法により結論は成り立つ。
 障害を収め得るcohomology群が非零であることと、
 この変形問題で実際に障害が現れることは同じではない。
 一般の有限パラメータの圏からrotationを省けるという大域的主張も含めない。
\end{remark}


\section{飽和な代数的族に対するétale局所復元}\label{sec:algebraization}

前節の形式的復元を有限型底上の結論へ移すためには、
圏の族そのものの代数性と相対smooth proper性が必要になる。
本節ではこれらを仮定として明示し、
可換模型と圏同値をétale近傍上で構成する。
次数別のflatnessだけからこれらの仮定が従うわけではない。

\paragraph{既知結果と本節の構成.}
étale局所復元では、Toën--Vaquié \cite{TV} のperfect objectのmoduliと
smoothnessの無限小判定 \cite{Stacks} を用いる。
新たに示すのは、同値bimoduleとrotation intertwinerを同時にパラメータ化したスタックが、
形式的復元定理によって底上smoothになることである。
Pirozhkov \cite{Pirozhkov} およびLin--Zhang \cite{LZ} のTorelliは
二つの可換モデルの一意性に関する結果であり、この局所的な存在の証明とは役割が異なる。
Huybrechts \cite[Section 9]{HuyUpdate} の四次元における一般的な非可換変形の問題と、
ここで扱うrotation付き変形問題とは、比較する構造の水準が異なる。

\begin{definition}\label{def:algebraic-family}
 \(S\) を有限型複素スキームとする。
 本節で\emph{smooth properな代数的dg圏の族}とは、
 étale局所的に \(\Perf(A)\) と書け、\(A\) が底環上のdg代数で、
 \(A\) 自身が底上perfect、対角 \(A\) が \(A^{\mathrm{op}}\otL A\) 上perfectである族をいう。
 圏は小さく、pretriangulatedかつ冪等完備とし、すべての底変更を導来的に取る。
 自己同値 \(\rho\) の代数的族は、可逆なperfect bimoduleによって与える。
\end{definition}

\begin{theorem}[étale局所的な可換復元]\label{thm:etale}
 \(S\) を有限型複素スキーム、\(s\in S(\CC)\) とし、
 \((\cK_S,\rho_S)\) を定義\ref{def:algebraic-family}の組とする。
 滑らかなcubic \(X^n\subset\PP^{n+1}\)、\(n\ge3\)、およびenhanced pair同値
 \[
 \phi_s:(\cK_s,\rho_s)\xrightarrow{\sim}(\Ku(X),\rho_X)
 \]
 が与えられたとする。このとき、あるpointed étale近傍
 \((T,t)\to(S,s)\)、滑らかな相対cubic \(Y\subset\PP^{n+1}_T\)、および
 \[
 \Phi_T:(\cK_T,\rho_T)\xrightarrow{\sim}(\Ku(Y/T),\rho_{Y/T})
 \]
 が存在し、\(t\) 上で \(Y_t=X\)、\(\Phi_t\simeq\phi_s\) とできる。
 この同値はperfect bimoduleとrotationを結ぶ同型によって与えられる。

 従って、可換cubicのpairで実現できる複素点は \(S\) のZariski開部分をなす。
 その各複素点で、実現する滑らかなcubicの射影同型類は一意である。
\end{theorem}

この定理は、底 \(S\) 全体上のcubicの一意な降下や、同値の一意性を述べない。
また、\(S\) が連結であることだけでは、上の開部分が \(S\) 全体であるとはいえない。
\(n\ne4\) ではrotationを最初から与えなくても、同じ方法でétale局所的な
可換モデルが存在する（系\ref{cor:etale-unpolarized}）。

\subsection{指定された還元同値を保つ持ち上げ}

\begin{lemma}\label{lem:relative-realization}
 \(R'\twoheadrightarrow R\) を剰余体 \(\CC\) の局所Artin代数のsmall extensionとする。
 \((\cK_{R'},\rho_{R'})\) をsmooth proper pairとし、滑らかなcubic \(Y_R\) と同値
 \[
 \phi_R:(\cK_R,\rho_R)\simeq(\Ku(Y_R),\rho_{Y_R})
 \]
 を固定する。このとき \(Y_R\) を延長する \(Y_{R'}\) と、
 還元後に \(\phi_R\) と同型になるpair同値 \(\phi_{R'}\) が存在する。
 \(Y_R\) の射影空間への埋込みも指定してよい。
\end{lemma}
\begin{proof}
 特殊ファイバーでの \(\phi_R\) をmarkingとして固定する。
 まずcubicの係数を持ち上げ、幾何側に一つのlift \(Y^0_{R'}\) を取る。
 \(P\) を第\ref{sec:formal}節のpairのformal moduli problemとし、
 \(I=\ker(R'\to R)\) と置く。
 与えられた圏とこの幾何学的liftを、\(\phi_R\) によって
 \(P(R')\to P(R)\) の同じhomotopy fibre内で比較する。
 存在するliftの差は \(H^0(T_P)\otimes I\) のtorsorで測られる。
 幾何学的Kodaira--Spencer写像
 \(H^1(T_{Y_0})\simeq H^0(T_P)\) により \(Y^0_{R'}\) を修正し、その差を消す。
 補題\ref{lem:theta}の完全忠実性で、指定した還元同値を保つactual同値へ戻る。
 途中で還元上の自己同型が現れる表示を用いた場合には、
 補題\ref{lem:automorphisms}でそれを持ち上げて補正できる。

 抽象的な偏極cubicのliftに対しても、\(\OO(1)\) の延長と
 \(H^0(\OO(1))\) の自由な基底の持ち上げにより、指定された埋込みを延長できる。
 これは \(\operatorname{PGL}_{n+2}\) のsmoothnessを用いることと同じである。
\end{proof}

\subsection{同値とrotationの同型を表す代数的スタック}

\(U\subset\PP(\CC[x_1,\ldots,x_{n+2}]_3)\) を滑らかなcubicの方程式空間とし、
 \(\mathcal X\to U\) を普遍cubicとする。
 相対的な例外列とmutationにより
 \((\Ku(\mathcal X/U),\rho_U)\) はsmooth properな代数的pairの族となる。
 相対生成子として \(\bigoplus_{i=0}^{n+1}\OO_{\mathcal X}(i)\) の残余射影を使える。
 この生成性は、射影空間の相対Beilinson生成子の引戻しと、
 閉埋込みの押出しが保守的であることから従う。

\begin{proposition}\label{prop:equivalence-stack}
 \(S\)-スキーム \(T\) に対し、
 \[
 (f:T\to U,\ P,\ \alpha:P\circ\rho_T\xrightarrow{\sim}\rho_f\circ P)
 \]
 を対象とするスタック \(\mathcal E\) を考える。
 ここで \(P:\cK_T\simeq\Ku(\mathcal X_f/T)\) はperfect bimoduleで表される同値、
 射は \(\alpha\) と両立するbimoduleの同型である。
 このスタックは \(S\times U\) 上、従って \(S\) 上で局所有限表示のArtin 1-stackである。
\end{proposition}
\begin{proof}
 すべての構成はétale降下と両立するため、底をaffineにして生成子を選んでよい。
 \(\cD=\cK_S^{\mathrm{op}}\otL\Ku(\mathcal X/U)\) は \(S\times U\) 上smooth properである。
 Toën--VaquiéのTheorem 3.6とCorollary 3.15 \cite{TV}により、
 そのperfect bimoduleのmoduliは局所幾何的で局所有限表示である。
 同論文は一般の可換底環上で定式化されているので、定数族だけに限定されない。

 \(P\) が同値を与える条件は開である。
 実際、右随伴を表すperfect bimoduleを \(P^R\) とすると、
 単位と余単位
 \[
 \Delta_{\cK}\longrightarrow P^R P,
 \qquad PP^R\longrightarrow\Delta_{\Ku}
 \]
 のconeがともに零となることが必要十分である。
 生成子でdg代数を選んだ模型で、これらのconeを評価する。
 smooth proper性によりその底上の複体はperfectであり、消滅はそこで検出される。
 従って消滅座は開である。
 これは \cite[Lemma 3.25とCorollary 3.24の証明]{TV} の二つの圏に対する同じ議論である。

 同値 \(P\) の幾何ファイバーで
 \[
 \Ext^i(P,P)=\HH^i(\Ku(X_f)),\qquad i\in\ZZ.
 \]
 補題\ref{lem:HH}から負次数は消える。
 perfectな自己Hom複体のTor振幅の開性により、これは任意の底変更後にも負次数が
 消える開部分として扱える。\cite[Corollary 3.21]{TV}により、
 同値を表す部分は通常のArtin 1-stackとなる。

 最後に \(P\rho_T\) と \(\rho_fP\) の間の同型を付加する。
 射のmoduliの局所有限表示性は \cite[Corollary 3.19]{TV} による。
 その中でconeが零となる部分を取れば同型のmoduliを得る。
 同型が存在する点では両側は可逆bimoduleであり、
 負次数のHomは再び \(\HH^{<0}\) の消滅によって零である。
 従ってこの同型の選択は代数的空間で表され、
 \(\mathcal E\) はArtin 1-stackのままである。
 この構成には \(\rho^3\simeq[2]\) の選択や、そのcoherenceは課していない。
\end{proof}

\subsection{smoothnessとétale局所section}

\begin{proof}[定理\ref{thm:etale}の証明]
 \(\phi_s\) と \(X\) の方程式を選んで \(e\in\mathcal E(\CC)\) を得る。
 \(V\to\mathcal E\) をsmooth atlasとし、\(e\) の上に複素点 \(v\) を選ぶ。
 \(v\) における \(V\to S\) のsmall extensionの持ち上げ問題を考える。
 補題\ref{lem:relative-realization}により、対応する \(\mathcal E\) の対象を
 指定された還元を保って持ち上げられる。
 さらに \(V\to\mathcal E\) のsmoothnessにより、それを \(V\) へ持ち上げる。
 底が局所Artin環で剰余体が \(\CC\) なので、atlasへのétale局所持ち上げはsectionを持ち、
 この持ち上げを同じ底上で取れる。

 \(V\to S\) は局所有限型であるから、Noetherian底上の無限小判定
 \cite[Tag 02HX]{Stacks}により \(v\) でsmoothである。
 \(v\) の近傍を縮め、smooth morphismのétale局所section
 \cite[Tag 055Uとその証明中のLemma 37.38.5]{Stacks}を使うと、
 \(\kappa(v)=\kappa(s)=\CC\) であることから、指定した \(v\) を通る
 \((T,t)\to(S,s)\) と \(T\to V\to\mathcal E\) が得られる。
 これが求める \(Y\)、\(P\)、\(\alpha\) を同時に与える。

 \(T\to S\) の像はZariski開であり、その各複素点の上には複素点を取れる。
 従って可換モデルが存在する複素点の集合は開部分の複素点集合である。
 各点でのモデルの一意性は \cite[Theorem 1.7]{LZ} による。
\end{proof}

\begin{corollary}\label{cor:generic-realization}
 定理\ref{thm:etale}の仮定で \(S\) が既約なら、\(s\) を含む空でないZariski開部分
 \(S^\circ\) の全複素ファイバーに可換cubicモデルが存在する。
 幾何学的generic fibreにも、有限型étale底変更を通じてモデルが存在する。
\end{corollary}
\begin{proof}
 上のétale近傍の像を \(S^\circ\) とする。既約性からこれはdenseである。
 \(S\) を必要なら被約化してgeneric pointを考えると、
 \(T\) の支配的な成分のgeneric residue fieldは \(\CC(S)\) の有限分離拡大である。
 その上のcubicと同値を代数閉包へ底変更すればよい。
\end{proof}

\begin{remark}\label{rem:algebraization-boundary}
 本節では形式的なpower seriesの収束を証明していない。
 指定された代数的圏の族に対し、代数的な同値のmoduliのsmoothnessから
 別途、実際の代数的近傍と同値を作っている。
 その完備化は形式的復元定理のモデルとmarked同型類で一致する。
 しかし、任意に選ばれた形式的座標表示や形式的bimoduleそのものが
 元の底上へ一意に降下する、という結論ではない。
\end{remark}


\begin{corollary}[rotationを指定しない局所的存在]\label{cor:etale-unpolarized}
 定理\ref{thm:etale}で \(n\ne4\) とする。
 \(\rho_S\) は仮定せず、\(\cK_s\simeq\Ku(X)\) というenhanced同値だけを与えても、
 pointed étale近傍 \(T\) 上にcubic \(Y/T\) と、指定された特殊ファイバーの同値を
 延長する \(\cK_T\simeq\Ku(Y/T)\) が存在する。
\end{corollary}
\begin{proof}
 命題\ref{prop:equivalence-stack}から \(\alpha\) の指定を除いた同値のスタックを用いる。
 定理\ref{thm:unpolarized-formal}の証明を、還元同値を指定したhomotopy fibre内で行えば、
 補題\ref{lem:relative-realization}のrotationなしの版が得られる。
 後は定理\ref{thm:etale}と同じsmooth atlasとétale局所sectionの議論を適用する。
 得られたモデルからrotationを移せるが、その \(T\) 全体上での一意性は主張しない。
\end{proof}


\section{同値双加群の障害消去とétale局所実現}\label{er:sec:etale}
本節では、可換模型の存在に使う局所実現原理を証明する。
入力は実在するsmooth properな代数的dg族であり、proper性だけの族ではない。
圏の全変形をHochschild複体から実現するという主張は用いず、
実在するperfect双加群のmoduliと、その完全な障害理論を用いる。
ここでのpairは自己同値一つを指定したもので、$\rho^3\simeq[2]$という
選択された同型やそのcoherenceを追加データとして含めない。

\subsection{perfect加群の小拡大障害}
小拡大は剰余体$k=\CC$を持つ局所Artin代数の全射
$A'\twoheadrightarrow A$で、核$I$が$\mathfrak m_{A'}I=0$を満たすものとする。
$D'$は$A'$上smooth properなdg代数とし、$D=D'\otL_{A'}A$とする。

\begin{lemma}[加群の完全な障害理論と特徴作用]\label{er:lem:obstruction}
$P\in\Perf(D)$の持上げには自然な障害
\[
 o(P;D')\in\operatorname{Ext}^2_{D_k}(P_k,P_k)\otimes_k I
 \tag{O}\label{er:eq:obstruction}
\]
があり、零であることと$\Perf(D')$への持上げの存在は同値である。
障害はMorita関手と対象の同型に関して自然である。
$D'=C'^{\rm op}\otL_{A'}B'$で$P$が$C$から$B$への同値を表すとき、
$B'$を固定し、$C'$のliftをHochschild類$\xi\in\mathrm{HH}^2(C_k)\otimes I$だけ変えると、
障害の変化は、符号を一貫して選べば
\[
 o_{\rm new}-o_{\rm old}=-\chi_P(\xi),\qquad
 \chi_P:\mathrm{HH}^2(C_k)\xrightarrow{\sim}
       \operatorname{Ext}^2_{C_k^{\rm op}\otimes B_k}(P_k,P_k).
 \tag{V}\label{er:eq:variation}
\]
\end{lemma}
\begin{proof}
Toën--Vaquié\cite[Theorem 3.6, Corollaries 3.15, 3.17]{TV}を
可換底$A'$と$D'$に適用する。smoothnessによりpseudo-perfect加群と
perfect加群は一致し、その実際の対象moduliの相対接複体は$\RHom_D(P,P)[1]$である。
導来対象moduliの平方零拡大に関する障害理論は、その接複体の次数一に完全な障害を持つ。
従って障害は$\operatorname{Ext}^2_D(P,P\otL_A I)$に入り、
導来随伴と$\mathfrak m_A I=0$によって\eqref{er:eq:obstruction}に一致する。
これは単なる接空間の計算を代数化の代用にしているのではない。
実在する対象の導来moduliが無限小凝集的であることを使っており、
障害のnull-homotopyはそのmoduliの$A'$-点、すなわち実際のperfect加群を与える。

変化式も説明する。cofibrantな代数模型と加群のbar模型を選び、
構造を次数一のbar余微分$b$で表す。任意のlift$\widetilde b$の曲率
$\widetilde b^{\,2}$は$I$に値を取り、$[b,\widetilde b^{\,2}]=0$である。
持上げを$I$値の$h$だけ変えると、$I^2=0$により曲率は$[b,h]$だけ変わる。
この曲率類が上のmoduliの障害である。
代数側を$\xi$だけ変えた際の曲率の変化は、$\xi$を加群の構造写像に挿入したcochainであり、
それが特徴作用$\chi_P(\xi)$である。符号はsourceの反対代数とtorsorの規約に合わせて
\eqref{er:eq:variation}と取る。必要なのは同じ規約を両式で使うことである。
この記述はLowen\cite[Theorem 4.8]{ERLowen}の特徴作用による障害の記述を、
bar模型でdg加群と任意の小拡大に書いたものである。
$P_k$の逆Morita双加群をテンソルすれば、その自己Homと対角の自己Homが同定されるので、
$\chi_P$は同型である。曲率の構成、またはmoduliの障害の自然性から、最後の自然性も従う。
\end{proof}

\begin{lemma}[還元後の同値とintertwiner]\label{er:lem:lift}
$C',B'$は$A'$上smooth properとする。
$P'\in\Perf(C'^{\rm op}\otL B')$の還元がMorita同値なら$P'$自身も同値である。
さらに$Q'$もperfectで、$\operatorname{Ext}^1(P_k,Q_k)=0$なら、
任意の射$P'_A\to Q'_A$は$A'$上へ持ち上がる。元の射が同型ならliftも同型となる。
\end{lemma}
\begin{proof}
perfect双加群は両側に随伴を持ち、smooth proper性により随伴双加群もperfectである。
unitとcounitのconeは底上perfectで、還元後に零であるから導来Nakayamaで零となる。
射の持上げには三角$Q'\otL I\to Q'\to Q'_A$を$\RHom(P',-)$で評価する。
射の障害は$\operatorname{Ext}^1(P_k,Q_k)\otimes I$に入り、仮定で零である。
持ち上げた射のconeに同じNakayamaを適用すれば同型性も従う。
\end{proof}

\subsection{rotationを保つ障害消去}
$H$は複素数体上smoothな有限型スキーム、$h\in H(\CC)$とする。
$\mathcal A_H$はsmooth proper dg族、$r_H$は可逆perfect双加群で表された自己同値とする。
$\mathcal A=\mathcal A_h$、$r=r_h$と置く。局所的なdg代数模型を取って議論する。

\begin{theorem}[自己同値付き局所実現の判定]\label{er:thm:criterion}
次を仮定する。
\[
 \mathrm{HH}^{<0}(\mathcal A)=0,\qquad
 \mathrm{HH}^{1}(\mathcal A)=0,
 \qquad
 T_hH\longrightarrow\mathrm{HH}^2(\mathcal A)^r
 \text{ が全射}.
 \tag{KS}\label{er:eq:KS}
\]
最後の写像は$\mathcal A_H$のKodaira--Spencer写像である。
$S$は有限型複素スキーム、$(\mathcal B_S,s_S)$はsmooth proper dg族と
可逆perfect双加群のpairとする。$s\in S(\CC)$に指定した同値
$(\mathcal A,r)\simeq(\mathcal B_s,s_s)$があるなら、
pointed étale近傍$V\to(S,s)$、$f:V\to H$、$f(v)=h$と
指定同値を延長するpairの同値
\[
 (f^*\mathcal A_H,f^*r_H)\simeq(\mathcal B_V,s_V)
\]
が存在する。
\end{theorem}
\begin{proof}
まず小拡大$A'\twoheadrightarrow A$上の持上げを証明する。
$\operatorname{Spec}A'\to S$を固定し、$A$上の点$h_A$、同値双加群$P_A$、
intertwiner $\alpha_A:s_AP_A\xrightarrow{\sim}P_Ar_A$が与えられているとする。
$H$はsmoothなので$h_A$の任意のlift$h'_0$を取れる。
両側のliftされた自己同値から、双加群圏の自己同値
\[
 \mathcal F'(P)=s'\circ P\circ(r'_0)^{-1}
\]
を得る。$\alpha_A$は$\mathcal F_A(P_A)\simeq P_A$を与える。
補題\ref{er:lem:obstruction}の自然性より、$P_A$の持上げ障害$o$は
$\mathcal F_k$の作用で不変である。
Morita同定$\operatorname{Ext}^2(P_k,P_k)=\mathrm{HH}^2(\mathcal A)$の下で、
この作用は$r$による共役作用（規約によりその逆）である。
従って
\[
 o\in\mathrm{HH}^2(\mathcal A)^r\otimes I.
 \tag{I}\label{er:eq:invariant}
\]
この段階では$\alpha_A$を$A'$上へ持ち上げたとは仮定していない。

\eqref{er:eq:KS}と\eqref{er:eq:variation}により、$h'_0$を
$T_hH\otimes I$の元だけ修正して$o$を消せる。
修正後のsourceと固定したtargetの間にperfect双加群$P'$が存在し、
補題\ref{er:lem:lift}によって同値となる。
$\alpha_A$を持ち上げる障害は
\[
 \operatorname{Ext}^1(s_kP_k,P_kr_k)\otimes I
       \simeq \mathrm{HH}^1(\mathcal A)\otimes I=0
\]
なので、intertwinerも同型として持ち上がる。
自己同値一つのintertwiner以外に群関係は課していないので、ここで必要なデータは全て揃う。
任意のnilpotent拡大は小拡大に分解できる。

次に代数化する。$H\times S$上で、perfectな
$\mathcal A_H^{\rm op}\otimes\mathcal B_S$-加群$P$と
同型$s_SP\simeq Pr_H$のmoduliを考え、$P$が同値である部分を$\mathcal E$とする。
Toën--Vaquié\cite[Theorem 3.6, Corollaries 3.15, 3.19]{TV}により、
perfect対象とその射の導来moduliは局所幾何的かつ局所有限表示である。
unit、counitおよびintertwinerのconeは底上perfectだから、その消滅は開条件である。
従って$\mathcal E$も局所幾何的かつ局所有限表示である。
同値点での自己Extの負次数は$\mathrm{HH}^{<0}(\mathcal A)=0$であり、
その近傍でこの消滅を保てる。従って古典的切断はその近傍でArtin 1-stackとなる。
ここでは局所dg代数模型を用いてTVの定理を$H\times S$上の可換係数環に適用している。

上の小拡大持上げは、指定点$e=(h,s,P_k,\alpha_k)$で
$\mathcal E\to S$が形式的にsmoothであることを示す。
局所有限表示のArtin stackの無限小判定により、この射は$e$でsmoothである。
$e$を覆うsmooth atlas $U\to\mathcal E$を取り、点の近傍へ縮めると$U\to S$もsmoothである。
滑らかな射のpointed étale局所section\cite[Tags 02HX, 055U]{StacksER}を使い、
$V\to S$と$V\to U\to\mathcal E$を得る。
この普遍対象が所要の族、双加群およびintertwinerを与える。
\end{proof}

\begin{remark}[使わない仮定と先行研究]\label{er:rem:etale-scope}
この証明は$\mathrm{HH}^3=0$、Hochschild cochainのformality、
抽象的なHochschild変形の全てがactual Morita変形になること、または$\rho^3=[2]$の指定を使わない。
Lin--Zhang\cite[Theorem 2.2, Proposition 2.4]{ERLinZhang}の同値核の障害消去と同じ原理を用いるが、
ここではtargetを非幾何的なsmooth proper dg族として扱い、
自己同値が障害を不変部分へ制限することと、intertwinerの持上げを別々に証明した。
同論文の幾何的projectorの仮定を満たさないtargetへその定理を直接適用していない。
この一般定式化の新規性は未確認である。
\end{remark}

\subsection{滑らかなcubicへの適用}
\begin{lemma}[低次Hochschild群と幾何的接写像]\label{er:lem:cubic}
$X=V(F)\subset\mathbb P^{N-1}$を滑らかなcubic、$N=n+2\ge5$とし、
$J(F)=\CC[z_1,\ldots,z_N]/(\partial_1F,\ldots,\partial_NF)$とする。
\[
 \mathrm{HH}^{<0}(\Ku(X))=0,\quad
 \mathrm{HH}^0(\Ku(X))=\CC,\quad \mathrm{HH}^1(\Ku(X))=0,
\]
\[
 \mathrm{HH}^2(\Ku(X))=
 \begin{cases}J(F)_3&n\ne4,\\J(F)_3\oplus\CC^2&n=4,\end{cases}
 \qquad
 H^1(X,T_X)\xrightarrow{\sim}\mathrm{HH}^2(\Ku(X))^{\rho_X}=J(F)_3.
 \tag{H2}\label{er:eq:HH}
\]
$\dim J(F)_3=\binom N3$である。
\end{lemma}
\begin{proof}
BFK\cite[Theorem 5.39]{BFKER}のequivariant matrix factorizationの公式を用いる。
$G=\mathbb G_m$、各座標の重みは1、potentialのcharacterは3である。
恒等sectorのKoszul cohomologyは、偏微分が正則列であるためJacobian商だけで、
cohomological degree $2r$に$J(F)_{3r}$を与える。
非恒等$g\in\mu_3$の固定空間は原点、conormalの次元は$N$、determinantの重みは$-N$である。
次数$q$の重み零条件は$(3q-N)/2=0$であり、従って$3\mid N$かつ$q=N/3$のときだけ
各sectorから一次元が現れる。これで負次数、次数0、1、2の群が得られる。

幾何的Kodaira--Spencer写像の同定は、単なる次元一致ではない。
$n\ge4$では$3\le(n+2)/2$なので、Huybrechts--Rennemo
\cite[Lemma 2.19, \S2.5, \S3.3]{HRER}のAtiyah--Kodaira--Spencerと
三次周期性の比較を適用でき、$F+\epsilon G$の像は恒等sectorの$[G]\in J(F)_3$である。
$n=3$ではJacovskis--Lin--Liu--Zhang\cite[Theorem 4.4(2)]{JLLZER}の
幾何的写像の同型を用いる。normal sequenceとEuler sequenceから$H^1(T_X)=J(F)_3$である。
幾何的族にはrotationがあるのでこの像は不変である。
$n=4$で余分な二方向が不変でないことは、Huybrechts\cite[Proposition 3.12]{HuyER}による。
rotationはsymplecticで、Mukai格子の$A_2$上で三つの根を巡回する。
不変なCalabi--Yau類による$\mathrm{HH}^2\simeq\mathrm{HH}_0$の同定の下で、
その二方向の固有値は$\zeta,\zeta^2$である。よって不変部分は正確に$J(F)_3$となる。
最後に偏微分が二次正則列であることから$H_{J(F)}(t)=(1+t)^N$を得る。
\end{proof}

\begin{theorem}[cubicのrotation付きétale局所実現]\label{er:thm:ER}
$S$は有限型複素スキーム、$\mathcal K_S$は実在するsmooth properな代数的dg族で、
$\rho_S$は可逆perfect双加群により表される自己同値とする。
$s\in S(\CC)$で、$n\ge3$の滑らかなcubic $X$に対する指定同値
$(\mathcal K_s,\rho_s)\simeq(\Ku(X),\rho_X)$があるなら、
pointed étale近傍$V\to(S,s)$、滑らかなcubic族$Y/V$、および指定同値を延長する同値
\[
 (\mathcal K_V,\rho_V)\simeq(\Ku(Y/V),\rho_{Y/V})
 \tag{ER}\label{er:eq:ER}
\]
が存在する。
\end{theorem}
\begin{proof}
滑らかな三次形式の射影パラメータ空間の$[F]$を含むアフィン開を$H$とする。
普遍cubic $\mathcal X/H$のrelative perfect圏はsmooth properで、
$\OO,\ldots,\OO(n-2)$のrelative exceptional列の射影により
smooth properな相対Kuznetsov成分を得る。
射影とrotationは有限個の評価三角で構成され、底変更に可換である。
普遍cubic上の有限個のtwistから生成子を選べるので、局所的なdg代数模型も存在する。
$T_{[F]}H\to H^1(T_X)=J(F)_3$は全射である。
補題\ref{er:lem:cubic}と定理\ref{er:thm:criterion}を適用する。
\end{proof}

\begin{corollary}[rotationを最初に指定しない存在]\label{er:cor:unpolarized}
$n\ne4$なら、前定理で$\rho_S$を入力として指定せずとも、
$\mathcal K_s\simeq\Ku(X)$からétale局所的な可換cubic模型が得られる。
これはrotationなしの大域的な一意性の主張ではない。
\end{corollary}
\begin{proof}
補題\ref{er:lem:cubic}では幾何的接写像が$\mathrm{HH}^2$全体に全射である。
定理\ref{er:thm:criterion}の証明から不変性とintertwinerの段階を省けばよい。
\end{proof}

\section{Hochschildの奇数次障害とrotation不変部分}\label{ho:sec:parity}
可換復元における対角smoothnessを、形式変形の障害と混同しないために、
滑らかな三次超曲面のKuznetsov成分の低次Hochschild群を全次元で整理する。
以下は既知のequivariant matrix factorizationの計算からの帰結であり、
残余dg圏の相対smoothnessに関する新たな一般定理ではない。

\begin{proposition}[奇数次Hochschild群とrotation固有値]\label{ho:prop:sectors}
$X=V(F)\subset\mathbb P^{N-1}_\CC$を滑らかな三次超曲面、$N=n+2\ge5$とし、
$J(F)=\CC[z_1,\dots,z_N]/(\partial F)$とする。
$\operatorname{Ku}(X)$の標準rotationを$\rho$とする。
$\zeta$を原始三乗根と書けば、次の複素ベクトル空間の分解が成り立つ：
\[
 \operatorname{HH}^{q}(\operatorname{Ku}(X))\cong
 \begin{cases}
 J(F)_{3q/2},&q\text{ 偶数},\\ 0,&q\text{ 奇数}
 \end{cases}
 \quad\oplus\quad
 \begin{cases}
 \CC_{\zeta}\oplus\CC_{\zeta^2},&3\mid N,\ q=N/3,\\
 0,&\text{それ以外}.
 \end{cases}
 \tag{HS}\label{ho:eq:sectors}
\]
ここで$J(F)_d=0$（$d<0$）とし、右側の二つの一次元空間は
$\rho$の共役作用の非自明な固有空間である。
$J(F)$から来る項では$\rho$は恒等的に作用する。
特に
\[
 \operatorname{HH}^3(\operatorname{Ku}(X))
 \cong\begin{cases}\CC^2&n=7,\\0&n\ne7,\end{cases}
 \qquad
 \boxed{\operatorname{HH}^3(\operatorname{Ku}(X))^\rho=0\quad(n\ge3).}
 \tag{V3}\label{ho:eq:HH3}
\]
さらに$\operatorname{HH}^2(\operatorname{Ku}(X))^\rho\cong J(F)_3$で、
$\dim_\CC J(F)_3=\binom N3$である。
\end{proposition}
\begin{proof}
Orlovのgraded matrix factorization表示と
Ballard--Favero--Katzarkov\cite[Theorem 5.39, pp.~82--85]{BFKER}の
固定sector公式を用いる。
$G=\mathbb G_m$は全生成変数の重みが$1$となる作用で、
potentialのcharacterは三次である。$K_\chi=\mu_3$とする。
$g=1$において$V_g=\mathbb A^N$、$W_g=0$である。
偏微分$(\partial F)$は原点にのみ共通零点を持つ長さ$N$の正則列であり、
Koszul cohomologyは零次の$J(F)$に集中する。
重み零の条件とcohomological gradingから、恒等sectorは偶数次数$q=2r$に
$J(F)_{3r}$を与え、奇数次数には寄与しない。

$g\in\mu_3\setminus\{1\}$なら固定空間は原点で、$\dim W_g=N$、
$J(F_g)=\CC$である。
Theorem 5.39のtwistとdeterminant characterを代入すると、
$u=(q-N)/2$に対する重み零条件は$3u+N=0$となる。
従って$q=N/3$でのみ各非自明sectorが一次元を与える。
各sectorは標準的な文字twist $(1)$ による共役作用で
$g$の文字（規約によって逆文字）を受ける。
これは同論文のkernelの固定graph分解と
\cite[Theorem 6.5]{BFKER}のcharacter作用の計算と同じものである。
Orlov同値の下で次数shift $(1)$ は$\rho$に対応するため、
二つのsectorの固有値は順序を除いて$\zeta,\zeta^2$となる。
$J(F)$から来る恒等sectorではこの作用は自明である。
これで\eqref{ho:eq:sectors}を得る。

$q=3$は奇数なので、非零になり得るのは$N/3=3$、すなわち$N=9$
（$n=7$）だけである。しかもその二方向には$\rho$-不変ベクトルがない。
$q=2$の不変部分は常に$J(F)_3$である。
滑らかさから偏微分は二次正則列なので、
$H_{J(F)}(t)=(1+t)^N$となり$\dim J(F)_3=\binom N3$を得る。
\end{proof}

\begin{corollary}[非楕円cubic fivefoldの変形障害の所在]\label{ho:cor:fivefold}
滑らかな三次五次元超曲面$X^5$について
（特に節点型・尖点型の量子化で現れる滑らかな可換特殊ファイバーにも）、
\[
 (\dim\operatorname{HH}^0,\dim\operatorname{HH}^1,
 \dim\operatorname{HH}^2,\dim\operatorname{HH}^3,
 \dim\operatorname{HH}^4)=(1,0,35,0,7).
 \tag{F5}\label{ho:eq:fivefold}
\]
従って滑らかな可換特殊ファイバーのdg圏自身の通常のMorita変形に対する
\emph{一次の}Kuranishi障害空間$\operatorname{HH}^3$は零である。
ただし、量子化から構成した$\Gamma_R$の
$\Gamma_R^e$上の対角perfect性を、これだけからは結論できない。
\end{corollary}
\begin{proof}
$N=7$を命題\ref{ho:prop:sectors}に代入する。
$J(F)_0,J(F)_3,J(F)_6$の次元はそれぞれ$1,35,7$で、
それ以外の表示次数$q=1,3$は零となる。
Morita変形の一次障害が$\operatorname{HH}^3$に入ることは
標準的なHochschild変形理論の適用である。
\end{proof}

\begin{remark}[残余対角との論理的境界]\label{ho:rem:barrier}
$\operatorname{HH}^3(\Ku(X))^\rho=0$は重要な形式的情報であるが、
この不変部分の消滅だけで任意のrotation付きdg圏変形の厳密な
$\mathbb Z/3$-equivariant obstruction theoryを構成したとはしない。
また、Artin切断上の変形が可換的に説明できたとしても、
有限型係数環上の\emph{実在する}残余対角のperfect性、
すなわち$\Gamma_R\in\Perf(\Gamma_R^{\mathrm{op}}\otL_R\Gamma_R)$は別の問題である。
具体的な量子化族へ定理\ref{er:thm:ER}を適用する際には、
この残余対角のperfect性を別途確認しなければならない。
本節はBFKの既知の公式の具体的な帰結と、その本研究における障害の切分けを与える。
\end{remark}

% 2026-11-10 / v22
\section{全てのArtin切断のsmoothnessから完備基底のsmoothnessは従わない}
\label{as:sec:obstruction}

本節では、有限Artin基底での生成子の持上げと、
完備局所基底での生成性の間にある障害を明示する。
出発点となる$\mathbb P^1$上の生成対象の変形例は、
SashaのMathOverflow回答\cite{AMSasha}にある既知の例である。
以下の四次元strict dg模型、および全Artin切断の相対smoothnessと
generic fiberの非smoothnessを併記した整理の先行研究上の初出は未確認である。

この例は、非可換cubicから構成した本稿の$\Gamma_R$が非smoothであることを示すものではない。
示すのは、底上proper、滑らかな特殊ファイバー、
全ての有限Artin切断の相対smoothnessという抽象的条件だけでは、
完備基底上の相対smoothnessを結論できないという事実である。

\subsection{生成対象がgeneric fiberで生成しなくなる族}

$R=\CC[[h]]$、$K=\CC((h))$、$X=\mathbb P^1_R$とし、
$H^0(X,\mathcal O_X(1))$の基底を$s,t$とする。
ここで$s,t$は射影直線の一次sectionであり、
MFの周期変数$t$とは無関係である。
混同を避けるため、本節でのみこの記号を用いる。
二項perfect複体
\begin{equation}\label{as:eq:P}
 P_h=\bigl[\mathcal O_X(-1)\xrightarrow{\,hs\,}\mathcal O_X\bigr]
\end{equation}
をcohomological degrees $-1,0$に置く。

特殊ファイバーでは微分が零なので
\[
 P_0\simeq\mathcal O_{\mathbb P^1}\oplus
                  \mathcal O_{\mathbb P^1}(-1)[1].
\]
これは$\operatorname{Perf}(\mathbb P^1)$の古典的生成対象である。
一方、generic fiberでは$h$は可逆なので、
$s$の零点$p$について
\[
 P_h\otimes_R K\simeq\mathcal O_p.
\]
$\mathcal O_p$から有限coneと直和因子によって作る対象のsupportは$p$に含まれるため、
$\mathcal O_p$は$\operatorname{Perf}(\mathbb P^1_K)$を生成しない。
従って$P_h$も$\operatorname{Perf}(\mathbb P^1_R)$を生成しない。
実際、生成していれば、その有限生成操作を$K$へ底変更することで
generic fiberでも生成してしまう。

\begin{lemma}[各有限切断では生成する]\label{as:lem:artin-generation}
$R_N=R/(h^{N+1})$とし、$P_N=P_h\otimes_RR_N$とする。
このとき$P_N$は$\operatorname{Perf}(\mathbb P^1_{R_N})$を生成する。
\end{lemma}
\begin{proof}
射影直線の相対Beilinson分解により、perfect圏の底変更と
perfect対象間のHomの導来底変更は$R_N\to\CC$と両立する。
従って補題\ref{am:lem:generator}のInd圏における議論を、そのまま
$\operatorname{Perf}(\mathbb P^1_{R_N})$に適用できる。
具体的には、$P_N$と右直交するInd対象$M$について、
$P_N$のcompact性と底変更から
\[
 \operatorname{RHom}(P_0,M\otimes_{R_N}^{\mathbf L}\CC)=0
\]
となる。$P_0$は生成対象なので$M\otimes_{R_N}^{\mathbf L}\CC=0$である。
$R_N$の有限$h$-filtrationにより
$R_N\in\operatorname{thick}_{D(R_N)}(\CC)$だから、$M=0$となる。
これで$P_N$のcompact生成性、従ってperfect圏の古典的生成性を得る。

ここで用いた射影直線の有限模型は、対角の完全列
\[
 0\longrightarrow\mathcal O(-1,-1)
 \longrightarrow\mathcal O
 \longrightarrow\mathcal O_\Delta\longrightarrow0
\]
から得られる。従ってこの底変更の主張は、非正則なArtin基底でも成立する。
\end{proof}

\subsection{四次元strict dg代数}

\begin{proposition}[底上properな明示的dg模型]\label{as:prop:dg-model}
$\Gamma_R=\operatorname{RHom}_X(P_h,P_h)$は、次のstrict unital dg代数で表される：
\begin{equation}\label{as:eq:dg-algebra}
 \Gamma_R^0=Re_-\oplus Re_+,\qquad
 \Gamma_R^1=Rs\oplus Rt,\qquad
 \Gamma_R^q=0\quad(q\ne0,1),
\end{equation}
\[
 d(e_-)=hs,\qquad d(e_+)=-hs,\qquad d(s)=d(t)=0.
\]
$e_-,e_+$は互いに直交する冪等元で、単位は$e_-+e_+$である。
$s,t$は$e_-$から$e_+$へのdegree一の射であり、
\[
 e_+s=s=se_-,\qquad e_+t=t=te_-,
\]
その他の冪等元と射の積、および二本の射同士の積は零である。
特に$\Gamma_R$は$R$上perfectである。
この模型の形成は各$R_N$と$K$への導来底変更と両立する。
\end{proposition}
\begin{proof}
$P_h$の二つの項の間のsheaf Homは
$\mathcal O_X(-1),\mathcal O_X,\mathcal O_X(1)$の直和である。
これらの一次以上のcohomologyは零なので、
大域Homの二項複体が$\operatorname{RHom}$を実際に計算する。
degree $-1$の大域Homは$H^0(\mathcal O_X(-1))=0$であり、
degree零は二つの項のscalar endomorphism、degree一は
$H^0(\mathcal O_X(1))=Rs\oplus Rt$である。
複体の微分を$D=hs$とすると、endomorphismの微分は
$f\mapsto Df-(-1)^{|f|}fD$なので、表示した微分を得る。
積はendomorphismの合成そのものである。

この計算は有限自由$R$-加群とその積だけからなり、
同じline bundle cohomologyの消滅は各底変更後も成立する。
従って通常のtensorがこの模型の導来底変更を計算し、
底変更後の$P_h$のderived endomorphismを与える。
\end{proof}

\begin{lemma}[generic fiberの外積代数模型]\label{as:lem:generic}
$E=K[\epsilon]/(\epsilon^2)$に$|\epsilon|=1$、零微分を与える。
このとき
\[
 E\longrightarrow\Gamma_R\otimes_R K,\qquad
 1\longmapsto e_-+e_+,\qquad \epsilon\longmapsto t
\]
はdg代数のquasi-isomorphismである。
\end{lemma}
\begin{proof}
単位と積と微分を保つことは式\eqref{as:eq:dg-algebra}から直ちに分かる。
$h$が可逆なので
\[
 H^0(\Gamma_R\otimes_R K)=K(e_-+e_+),\qquad
 H^1(\Gamma_R\otimes_R K)=(Ks\oplus Kt)/Ks=K\cdot[t]
\]
であり、他のcohomologyは零である。
従って上の写像は全てのcohomologyで同型となる。
\end{proof}

\subsection{無限derived tensorによる非smoothness}

\begin{lemma}[degree一の外積代数はsmoothでない]
\label{as:lem:exterior-not-smooth}
前補題の$E$は$K$上properだが、homologically smoothではない。
\end{lemma}
\begin{proof}
augmentation $E\to K$によって$K$を左右の$E$-加群とする。
左$E$-加群のsemifree resolution
\[
 Q=\bigoplus_{n\ge0}Eu_n,\qquad |u_n|=0,\qquad
 d(u_0)=0,\quad d(u_n)=\epsilon u_{n-1}\quad(n\ge1)
\]
を考える。$Q\to K$は$u_0\mapsto1$、$u_n\mapsto0$（$n\ge1$）で与える。
$n$に関する増大filtrationはsemifree filtrationであり、
$d^2=0$は$\epsilon^2=0$による。
degree零のcycleは$Ku_0$だけであり、degree一の全ての元
$\epsilon u_n$は$d(u_{n+1})$である。
従って$Q\to K$はquasi-isomorphismである。

右からaugmentation加群$K$をtensorすると
\[
 K\otimes_E^{\mathbf L}K
 \simeq K\otimes_E Q
 =\bigoplus_{n\ge0}Ku_n
\]
となり、右辺の微分は零で、全ての$u_n$はdegree零にある。
従って$H^0(K\otimes_E^{\mathbf L}K)$は無限次元である。
もし$K$が$E$上perfectなら、このderived tensorは
有限個の$K$のshiftから有限coneと直和因子で作られ、
有限次元の有界cohomologyを持つはずなので矛盾する。

さらに、もし$E$の対角が$E^{\rm op}\otimes_KE$上perfectなら、
その有限cell表示の一方の$E$-作用へ$-\otimes_E^{\mathbf L}K$を施すと、
$K$の有限$E$-cell表示を得る。
これは$K$が$K$上perfectであることを用いる。
従ってsmoothnessは$K\in\operatorname{Perf}(E)$を含意するが、
直前でこれを否定した。
よって$E$はhomologically smoothでない。
一方、$E$自体は二次元の有界$K$-複体なのでproperである。
\end{proof}

\begin{theorem}[全Artin切断ではsmoothだが完備基底ではsmoothでない族]
\label{as:thm:obstruction}
命題\ref{as:prop:dg-model}の$\Gamma_R$は、次の全ての性質を持つ。
\begin{enumerate}
 \item $\Gamma_R$は$R$上perfectなstrict dg代数である。
 \item $\Gamma_R\otimes_R^{\mathbf L}\CC$は$\CC$上smoothである。
 \item 各$N\ge0$に対し
 $\Gamma_R\otimes_R^{\mathbf L}R_N$は$R_N$上smoothである。
 \item $\Gamma_R\otimes_R K$は$K$上smoothでなく、
 $\Gamma_R$も$R$上smoothでない。
\end{enumerate}
\end{theorem}
\begin{proof}
第一項は有限自由模型から従う。
補題\ref{as:lem:artin-generation}により$P_N$は
$\operatorname{Perf}(\mathbb P^1_{R_N})$の生成対象であり、
命題\ref{as:prop:dg-model}の底変更から
\[
 \operatorname{Perf}(\Gamma_R\otimes_R^{\mathbf L}R_N)
 \simeq\operatorname{Perf}(\mathbb P^1_{R_N})
\]
となる。
相対射影直線の対角は前述の二項線束分解を持つので、
このperfect圏は$R_N$上smoothである。
smoothnessはdg Morita同値で不変だから、第二・第三項を得る。

補題\ref{as:lem:generic}と
\ref{as:lem:exterior-not-smooth}からgeneric fiberはsmoothでない。
$\Gamma_R$が$R$上smoothなら、そのperfect対角の表示を
flatな底変更$R\to K$へ移してgeneric fiberもsmoothとなるため、矛盾する。
従って第四項も成立する。
\end{proof}

\begin{remark}[この反例の適用範囲]
$R_N\in\operatorname{thick}_{D(R_N)}(\CC)$はArtin基底での生成子の持上げを支える。
しかし$R=\CC[[h]]$は$D(R)$で$\CC$の有限coneと直和因子から作れない。
実際、そのような対象を$K$へtensorすると零になるが、$R\otimes_RK=K\ne0$である。
上の族では、この違いが生成性とsmoothnessの実際の破綻として現れる。

従って、全Artin比較の両立性だけから、完備基底上の生成性や相対smoothnessを補ってはならない。
可換側がsmooth properな射影族であることだけでも、変形した対象の生成性は保証されない。
中心三次非可換商の残余圏へこの原理を適用するには、
当該windowと生成対象についての有限生成証明、
または対角の具体的なperfect表示を別途用いる。
本例は、その特別な主張を否定してはいない。

なお本例の$R$上のcohomologyは
$H^0(\Gamma_R)=R$、$H^1(\Gamma_R)=Rt\oplus(R/hR)s$であり、
cohomologyの底上flatnessは満たさない。
この例が否定する仮定の範囲を、証明していない別の仮定へ広げてはならない。
例\ref{cg:ex:artin-singularity}はMFと絶対特異圏の区別を示すものであり、
本節のsmoothnessの持上げ障害とは別の反例である。
\end{remark}

\section{相対smoothnessと三次周期性の判定}
\label{sec:relative-saturation-reduction}
本節では滑らかな可換cubicを含むPoisson量子化族の相対saturationを、
全超曲面圏の対角の有限証明書へ還元する。
相対Beilinson仮定(RB1)--(RB2)、相対例外列、相対Orlov比較は
元論文の量子化定理からは自動的に従わない。
gluingに関する以下の命題は既知の一般原理の相対版であり、
新規な分類定理とは主張しない\cite[Theorem 3.25]{OrlovGlue}。

\subsection{例外列を除く前後の相対smoothness}
\begin{proposition}[相対smoothnessの同値]
\label{prop:saturation-reduction}
$R$ を可換Noether環、$\mathcal C$ を小さいidempotent completeな
$R$-linear dg圏とし、
$\RHom_{\mathcal C}(X,Y)\in\Perf(R)$ が任意の対象$X,Y$に対し成り立つとする。
有限個の相対例外対象 $E_1,\ldots,E_r$ と
dg半直交分解
\[
\mathcal C=\langle\mathcal K,E_1,\ldots,E_r\rangle
\]
が与えられ、さらに$\mathcal K$がcompact generatorを持つならば、
\[
\mathcal C\text{ は }R\text{-smooth}
\quad\Longleftrightarrow\quad
\mathcal K\text{ は }R\text{-smooth}.
\]
ゆえに相対propernessが既に分かっていれば、
両者の相対saturationは同値な条件である。
\end{proposition}
\begin{proof}
例外対象を一つずつ追加する帰納法により一個の場合に帰着する。
$\mathcal K$ のcompact generatorを$G$、例外対象を$E$とする。
この半直交分解のdg enhancementは$\mathcal K$と$\Perf(R)$を
\[
M=\RHom_{\mathcal C}(G,E)
\]
により三角dg gluingしたものとして表される（右左の規約によるoppositeを適宜取る）。
$\Gamma=\REnd_{\mathcal K}(G)$ と置けば、propernessより
$M$は$R$-perfectである。
$\mathcal K$がsmoothならば
$\Gamma\in\Perf(\Gamma^{\mathrm{op}}\otL_R\Gamma)$。
恒等式 $M\simeq M\otL_\Gamma\Gamma$ に対角の有限cell表示を適用すると、
その像は $M\otL_R\Gamma$ の有限回のshift、cone、直和因子から構成される。
$M$が$R$-perfectだから$M$は$\Gamma$-perfect、従って
gluing bimoduleがperfectとなる。三角dg圏の対角の標準triangleにより
$\mathcal C$はsmoothとなる。
逆に$\mathcal C$がsmoothなら、同じgluingの対角triangleを
両成分へ制限して$\mathcal K$の対角がperfectである。
以上は有限cell・導来テンソル・直和因子のみを使うので
$R$上でも成立する。体上の定理は\cite[Theorem 3.25]{OrlovGlue}。
\end{proof}

\begin{corollary}[必要な対角証明書は一つでよい]
\label{cor:only-ambient-diagonal}
量子化 $B_R=A_R/(\Omega_R)$ が研究ノートの相対Beilinson仮定
(RB1)--(RB2)と相対例外列を満たし、
$E=\bigoplus_{i=0}^{N-1}\mathcal O_{B_R}(i)$ がcompact generatorであるとする。
$\Lambda_R=\REnd(E)$は$R$-perfectであり、
その残余圏$\mathcal K_R$は、次の条件と同値に$R$-smoothである：
\[
\Lambda_R\in
\Perf(\Lambda_R^{\mathrm{op}}\otL_R\Lambda_R).
\tag{S}
\]
残余圏のEnd代数$\Gamma_R$と$\Lambda_R$の二つの対角を
別々に持ち上げる必要はない。
\end{corollary}
\begin{proof}
研究ノートの相対proper模型によって$\Perf(\Lambda_R)$はproperで、
$\mathcal K_R$と有限相対例外列による半直交分解を持つ。
命題\ref{prop:saturation-reduction}を適用する。
\end{proof}

\begin{remark}[閉ファイバーのsmoothnessだけでは足りない]
$R=\CC[\hbar]_{(\hbar)}$とし、
$\mathcal C_0\simeq\Perf(X)$が滑らかな可換cubicの圏でも、
その対角の有限cell分解を$R$へ持ち上げることは別の問題である。
$P\in\Perf(\Lambda_R^{\mathrm{op}}\otL_R\Lambda_R)$と
$P\to\Lambda_R$を構成し、$\hbar=0$で擬同型であることを証明できれば、
coneの$R$-perfect性とsupportの閉性により
零を含むZariski開近傍で(S)が成り立つ。
この有限証明書はまだ一般量子化族について構成されていない。
\end{remark}

\subsection{graded matrix factorizationの三次周期性}
\begin{lemma}[中心三次元の形式的周期性]
\label{lem:mf-three-period}
$A$を可換環$R$上の整数次数付き結合代数、
$\Omega\in Z(A)_3$を中心元とする。
有限射影次数付き加群のgraded matrix factorizationの
内部次数シフトを$(1)$とすると、自然なdg自己同値の同型
\[
(1)^3\simeq[2]
\]
があり、任意の基底変更と両立する。
\end{lemma}
\begin{proof}
因子分解を
\[
P_0\xrightarrow{d_0}P_1\xrightarrow{d_1}P_0(3),
\quad
d_1d_0=\Omega,\quad d_0(3)d_1=\Omega
\]
と書く。懸垂$[1]$は二項を入れ替え、片側を$(3)$移し、
二つの微分に負号を付ける。二回の懸垂では負号が消え、
両項および微分が一斉に$(3)$移る。
従って$[2]\simeq(3)=(1)^3$である。
次数移動と符号のみを用いるので基底変更とも可換である。
\end{proof}

\begin{corollary}[相対Orlov比較の下でのrotation]
\label{cor:relative-rotation-period}
$\mathcal K_R$と$\MF^{\mathrm{gr}}(A_R,\Omega_R)$の間に、
基底変更と両立する相対Orlov同値があり、
$(1)$が標準rotation$\rho_R$に対応すると仮定する。
すると$\rho_R$は自己同値で、相対的に
\[
\rho_R^3\simeq[2]
\]
が成立する。この段階では相対smoothnessを別途仮定しない。
\end{corollary}
\begin{proof}
補題\ref{lem:mf-three-period}を相対同値に沿って移す。
\end{proof}

\subsection{相対的な可換復元に必要な追加条件}
相対形式復元を代数的な族の全ファイバーへ拡張する際には、次の四点を区別する：
(RB1)--(RB2)を指定量子化族について証明すること、
全超曲面圏$\Perf(\Lambda_R)$の対角の有限証明書を可換点から持ち上げること、
相対Orlov比較でrotationを同定すること、
そして得られたsmooth proper pairが可換cubicのpairの代数的変形像に入ること。
前の三点が解けても最後の可換モデルの存在は別問題であり、
Zhang twistの障害を無視することはできない。


\part{有限パラメータの楕円ブロックと生成性}
\section{楕円ブロックの明示模型と有限係数の主定理}\label{sec:explicit-block-results}

\[
 S_h=\CC\langle x,y,z\rangle/
 ([x,y]-hz^2,\ [y,z]-hx^2,\ [z,x]-hy^2)
\]
と置き、\(T=\sym(xyz)\) を六つの順列の平均、
\(G_h=x^3+y^3+z^3-2h^2T\) とする。
中心変数 \(u_1,\ldots,u_m\) と斉次三次式 \(P(u)\) に対して
\[
 B_{h,P}=S_h[u_1,\ldots,u_m]/(G_h+P(u))
\]
を考える。\(P=\sum_i u_i^3\) の場合を \(B_{h,m}\) と略記する。

\begin{theorem}[一ブロックと任意の中心三次式の可換モデル]\label{thm:global}
 \(m\ge2\)、\(n=m+1\)、\(h\in\CC\)、\(h^6\ne-27\) とする。
 \(P\in\CC[u_1,\ldots,u_m]_3\) のアフィン臨界点が原点だけであると仮定する。
 \(B_{h,P}\) の残余dg圏とrotationは、滑らかなcubic
 \[
 Y_{h,P}=V\!\left(x^3+y^3+z^3+h^2xyz+P(u)\right)
 \subset\PP^{m+2}
 \]
 の \((\Ku(Y_{h,P}),\rho_{Y_{h,P}})\) と同値である。
 この \(Y_{h,P}\) は、同じrotation付き残余圏を持つ滑らかな可換cubicとして
 射影同型を除いて一意である。
\end{theorem}

この定理はsixfoldでは \(m=5\)、sevenfoldでは \(m=6\) を含む。
中心変数どうしの任意の三次混合項を許すが、一つのSklyaninブロックと中心部分が
分離されているという仮定を保持する。\(P=\sum u_i^3\) はその特殊例である。
複数の非可換ブロックについても、対応する可換cubicの残余圏からの完全忠実埋込みを得る
（命題\ref{prop:multiblock-embedding}）。次の部分的な有限位数条件の下では、その像が全体を生成する。

\begin{theorem}[一つの任意位数ブロックを許す多ブロック可換モデル]\label{thm:multiblock-model}
 \(r\ge2\)、\(m\ge0\)、\(n=3r+m-2\) とする。
 各 \(j\) について \(h_j^6\ne-27\) とし、楕円曲線
 \[
 E_{h_j}=V(x_j^3+y_j^3+z_j^3+h_j^2x_jy_jz_j),\qquad
 o_j=[1:-1:0],\quad p_j=[1:-1:-h_j]
 \]
 上の点 \(p_j\) は、高々一つを除いて有限位数であると仮定する。
 すなわち因子を並べ替えた後、\(j\ge2\) にのみ有限位数を仮定し、\(p_1\) の位数には条件を課さない。
 \(m>0\) の場合、\(P\in\CC[u_1,\ldots,u_m]_3\) は原点だけに臨界点を持つとし、
 \(m=0\) の場合は中心部分を省く。
 異なるSklyaninブロックの元を可換とした代数
 \[
 B_{\mathbf h,P}=
 (S_{h_1}\otimes\cdots\otimes S_{h_r})[u_1,\ldots,u_m]
 \big/\left(\sum_jG_{h_j}+P\right)
 \]
 のrotation付き残余dg圏は、滑らかなcubic
 \[
 Y_{\mathbf h,P}=V\!\left(
 \sum_{j=1}^r(x_j^3+y_j^3+z_j^3+h_j^2x_jy_jz_j)+P(u)
 \right)\subset\PP^{3r+m-1}
 \]
 の \((\Ku(Y_{\mathbf h,P}),\rho_{Y_{\mathbf h,P}})\) と同値である。
 この可換モデルは射影同型を除いて一意である。
\end{theorem}

ここで「有限パラメータ」は形式変数でなく実際の複素数であることを意味し、
楕円曲線上の点が「有限位数」であることとは異なる。
定理は一つの無限位数ブロックを許す。
具体的に \(h_1=2\) は無限位数（補題\ref{lem:order-infinite}）、
\(h_j^6=5\ (j\ge2)\) は位数五（補題\ref{lem:order-five}）となる。
従って二ブロックのsixfold、三ブロックのsevenfoldにも、真に無限位数の因子を含む例が得られる。
二つ以上の無限位数ブロックを任意に許す主張は、本定理には含まれない。
証明は第\ref{sec:finite-center}節に置く。

\begin{proposition}[例外パラメータのhomological singularity]\label{prop:bad-intro}
 \(m\ge0\) とする。Fermat中心部分の族 \(B_{h,m}\) で \(h^6=-27\) とすると、\(S_h[u]\) はAS-regularで
 \(G_h+\sum u_i^3\) は中心正則元であるが、そのgraded matrix factorization圏は
 properではなく、さらに
 \[
 \operatorname{hdim}(\qgr B_{h,m})=\infty
 \]
 である。ここでhomological dimensionは、非零のExtが存在する次数の上限である。
 従って、滑らかな射影cubicのKuznetsov成分とは同値にならない。
\end{proposition}

これはAS-regularな周囲、三次の中心正則元、期待されるHilbert series、
および \(\rho^3\simeq[2]\) だけでは、大域的な可換モデルの存在を保証しないことを示す。
本稿の可換点からのPoisson量子化族に限定した、良い有限パラメータ全体での存在問題は別に残る。
一般のAS-regular中心三次商については、smooth properかつCalabi--Yauな例にも
定理\ref{thm:zhang-cy-fourfold-counterexample}の反例がある。

\paragraph{証明の所在.}
一ブロックの定理\ref{thm:global}では、楕円曲線と次数シフトの比較
（第\ref{sec:elliptic}節）に続いて、中心変数と通常次数への降下
（第\ref{sec:TS}節）を用いる。
多ブロックの定理\ref{thm:multiblock-model}では、
外部積の完全忠実性と生成性の区別（第\ref{sec:multiblock}節）から、
一方有限中心の生成定理（第\ref{sec:finite-center}節）へ進む。
例外係数の命題\ref{prop:bad-intro}は、第\ref{sec:scope}節で証明する。

\subsection{有限係数模型と外部積生成に関する先行研究}

明示的族の楕円ブロックはEtingof--Ginzburg \cite{EG} と
Sklyaninの点曲線の公式 \cite{Walton} に依拠する。
本稿では次数シフトの共役、中心変数を加える際の生成性、
有限群による通常次数への降下を証明する。
可換な係数環についてのThom--Sebastiani定理を、
非可換な係数環に対して仮定を確認せず引用することはしない。

有限中心代数の外部積生成については、Elagin--Lunts--Schn\"urer
\cite[Theorem 4.18]{ELS} が、perfect field上のnoetherian J-2 schemesに載る
coherentな代数に対し、ungradedな有界導来圏の生成定理を証明している。
本稿では、\(\ZZ^r\) の任意の粗化次数を保ち、さらに一方の係数代数には
有限中心性を仮定しない形を証明する（定理\ref{thm:one-sided-boxes}）。
\cite[Theorem 4.18]{ELS} は双方の因子をcoherentな構造代数とするため、
任意のnoetherian係数代数を許すここでの主張は、その定理の直接の適用ではない。
新たな入力はBeilinsonの対角分解 \cite{Beilinson} を非可換係数付きで用いる
多項式拡大の補題である。中心のNoether正規化と分離冪等元によるretractを組み合わせ、
斉次中心局所化の核を次元について帰納する。
対象の細分次数へのliftや、係数代数の有限global dimensionは仮定しない。
有限位数THCRの有限中心性 \cite[Corollary 2.3]{SmithTate} をこの片側のみに適用し、
一つの任意位数ブロックを含む多ブロックのMF比較を得る。
He--Ma--Ye \cite{HMY} のtensor productに対する一般化Kn\"orrer周期性は、
中心元の次数が二である非可換quadricとClifford変形を扱う。
本稿の三次中心元への直接の適用ではない。
周囲のnoetherian性については、Chan \cite{ChanMulti} の多重Serre同値と
面・帯による議論を任意個の楕円因子へ展開し、有限位数の仮定なしで証明する。
Keeler \cite{KeelerMulti} のTheorem 2.7は任意個の可換な可逆bimoduleの
NC-amplenessを判定するが、同論文のTheorem 3.3とCorollary 3.5のnoetherian性は
二因子についての主張である。本稿では任意個の積に必要な帰納法を記す。
基礎となる対角分解と有限中心のungraded生成定理は既知である。
一方有限中心の次数付き定式化および本適用の先行研究上の新規性は確認継続中である。

\section{楕円一ブロックの明示的可換復元}\label{er:sec:intro}

全体を通じて基礎体は$\CC$である。
三次中心正則元$g$を持つ次数付きAS-regular代数$S$から
$S[u_1,\ldots,u_m]/(g+P(u))$を作るとき、その残余圏を
可換三次超曲面から復元できるかを考える。
中心変数$u_i$は$S$の全ての元と可換とし、$P$は次数三、
\begin{equation}\label{er:eq:isolated}
 V(\partial_{u_1}P,\ldots,\partial_{u_m}P)=\{0\}\subset\mathbb A^m,
 \qquad m\ge2
\end{equation}
を仮定する。これは$P$のアフィン臨界点が原点のみであるという条件である。

使用する量子化は、$t^3\ne-1$に対する
\[
 f_t(x,y,z)=\frac{x^3+y^3+z^3}{3}+txyz
\]
のJacobian Poisson括弧の量子化である。
$h$を次数零の変形パラメータとし、$S_{h,t}$を
\begin{equation}\label{er:eq:relations}
\begin{aligned}
 [y,z]&=h\bigl(x^2+\tfrac t2(yz+zy)\bigr),\\
 [z,x]&=h\bigl(y^2+\tfrac t2(zx+xz)\bigr),\\
 [x,y]&=h\bigl(z^2+\tfrac t2(xy+yx)\bigr)
\end{aligned}
\end{equation}
で定める。$x,y,z,u_i$の内部次数は一である。
次の多項式と有理関数を用いる：
\begin{equation}\label{er:eq:coefficients}
\begin{aligned}
 b(h,t)&=1-\frac{h^2t^2}{4},&
 a(h,t)&=t+\frac{h^2}{3}\left(1+\frac{t^3}{4}\right),\\
 \vartheta(h,t)&=\frac{a(h,t)}{b(h,t)},&
 \Delta(h,t)&=a(h,t)^3+b(h,t)^3,\\
 U(h,t)&=1+\frac{h^2t^2}{2}+\frac{h^4t}{6}-\frac{h^4t^4}{48},\\
 W(h,t)&=\frac{(4-h^2t^2)\bigl(12t+h^2(t^3-8)\bigr)}{48}.
\end{aligned}
\end{equation}
$\mathscr U_t=\{h:U(h,t)b(h,t)\Delta(h,t)\ne0\}$は$0$を含む開集合である。
$\sym(xyz)$は六つの順序付き積の平均を表し、
\begin{equation}\label{er:eq:center}
 \Omega_{h,t}=\frac{U(h,t)}3(x^3+y^3+z^3)+W(h,t)\sym(xyz)
\end{equation}
と置く。中心性・正則性は第\ref{er:sec:algebra}節で確認する。

\begin{theorem}[一ブロックの明示的可換復元]\label{er:thm:main}
$t^3\ne-1$、$h\in\mathscr U_t$および\eqref{er:eq:isolated}を仮定する。
\[
 B_{h,t,P}=S_{h,t}[u_1,\ldots,u_m]/(\Omega_{h,t}+P(u)),
 \qquad
 Y_{h,t,P}=V\bigl(f_{\vartheta(h,t)}(x,y,z)+P(u)\bigr)
 \subset\PP^{m+2}
\]
とおく。$Y_{h,t,P}$は滑らかな三次超曲面であり、dg enhancementを保つ同値
\begin{equation}\label{er:eq:main-pair}
 (\mathcal K_{B_{h,t,P}},\rho_B)
 \simeq (\Ku(Y_{h,t,P}),\rho_Y)
\end{equation}
が存在する。ここで$\mathcal K_B$はOrlovの残余成分、
$\rho_B=\mathbf L_{\OO_B}\circ(1)$である。
従って$\mathcal K_B$はsmooth properである。
この可換模型は射影同型を除いて一意である。
楕円点曲線上の平行移動の有限位数性は仮定しない。
\end{theorem}

\begin{remark}[主張の範囲]\label{er:rem:scope}
定理は各複素パラメータにおける同値である。
全ての同値を底上の一つの代数的dg同値に選んだとは主張していない。
また$u_i$と$x,y,z$の本質的な混合項や、二つ以上の独立な無限位数楕円ブロックは対象外である。
smooth proper性はこの特定の残余成分について同値から得る結論であり、
任意のAS-regular中心三次商からの一般的帰結ではない。
\end{remark}

\subsection{先行研究と論理的な依存関係}
三次元Sklyanin代数の中心三次元と楕円曲線による商の表示は古典的である
\cite{SmithTate,Walton,DeLaetLeBruyn}。
Orlovの定理はGorenstein parameter零の特異圏を射影圏と同一視し、
正のparameterでは特異圏を残余成分として実現する\cite[Theorem 2.5]{Orlov}。
非可換行列因子分解との比較には\cite[Theorem 5.8]{CCKM}を用いる。
Etingof--Ginzburgの\cite[Theorem 3.6.2]{EG}にも楕円曲線とgraded factorizationの比較が現れる。
本稿では同結果のgenericという仮定を全パラメータへ読み替えず、
各$h\in\mathscr U_t$で中心商とOrlovの適用条件を確認する。

本稿の比較の要点は、楕円曲線の次数シフトに現れる平行移動を直線束のテンソルで共役し、
その三次周期性を保存したまま中心変数を追加することである。
有限群の双対同変化には\cite[Theorem 4.2, Section 8]{Elagin}を用いる。
可換模型の一意性は\cite[Theorem 1.7]{LZ}によるが、これは存在を与える定理ではない。
本稿の結果は既存の一楕円ブロック比較をHesse量子化へ適用し、
その変形方向と一定の可換模型を持つ代数的曲線を明示したものである。
これらの組合せについて先行研究上の初出を確定したとは主張しない。

\subsection{圏と次数の規約}
右加群を用いる。$\MF^\Gamma(A,w)$は有限射影$\Gamma$次数付き行列因子分解のdg圏の
pretriangulated hullを取り冪等完備化したものとする。
ここで行列因子分解自体は$D^2=w$を満たすcurved objectであり、通常の複体ではない。
一方、二対象間のHomには$d(\varphi)=D_Q\varphi-(-1)^{|\varphi|}\varphi D_P$が入り、
$w$の中心性から$d^2=0$となる。
$\Gamma=\ZZ$、$\deg w=3$なら、内部シフトは自然な周期性
\begin{equation}\label{er:eq:period}
 (1)^3=(3)\simeq[2]
\end{equation}
を備える。以下の同値では、この周期性も記録する。

\section{中心三次元と楕円商}\label{er:sec:algebra}
\begin{proposition}\label{er:prop:center}
$h\in\mathscr U_t$では$S_{h,t}$はnoetherianな三次元Koszul AS-regular domainで、
Hilbert級数は$(1-z)^{-3}$である。
\eqref{er:eq:center}は中心正則元であり、$\Omega_{0,t}=f_t$。
さらに滑らかな三次曲線
\[
 E_{h,t}=V\bigl(x^3+y^3+z^3+3\vartheta(h,t)xyz\bigr)
\]
と$L=\OO_{E_{h,t}}(1)$について
\begin{equation}\label{er:eq:thcr}
 S_{h,t}/(\Omega_{h,t})\simeq B(E_{h,t},L,t_{p_{h,t}}),
 \qquad
 p_{h,t}=[1-ht/2:-1-ht/2:-h]
\end{equation}
が成り立つ。原点は$o=[1:-1:0]$である。$h=0$では$t_{p_{0,t}}=\Id$と解釈する。
\end{proposition}
\begin{proof}
$q=ht/2$、$\alpha=1-q$、$\beta=-1-q$、$\gamma=-h$とすると、関係式は
$\alpha yz+\beta zy+\gamma x^2=0$とその巡回置換である。
$h\ne0$では$\alpha\beta\gamma=h b(h,t)\ne0$かつ
\[
 (\alpha^3+\beta^3+\gamma^3)^3-(3\alpha\beta\gamma)^3
 =-27h^3\Delta(h,t)\ne0.
\]
従って楕円型Sklyanin代数の既知の定理を適用できる\cite{SmithTate,Walton}。
その点曲線は
$\alpha\beta\gamma(x^3+y^3+z^3)-
(\alpha^3+\beta^3+\gamma^3)xyz=0$である。
$\alpha^3+\beta^3+\gamma^3=-3h a(h,t)$を代入すると$E_{h,t}$を得る。
$1+\vartheta^3=\Delta/b^3\ne0$なので滑らかである。

中心元の規格化を確認するため、$S_3=x^3+y^3+z^3$、
$C=xyz+yzx+zxy$、$D=yxz+xzy+zyx$と書く。
\cite[Section 2, p.~3]{DeLaetLeBruyn}の中心三次元の表示は
\[
 G=\alpha(\beta^3-\gamma^3)C+
   \beta(\gamma^3-\alpha^3)D+
   \gamma(\alpha^3-\beta^3)S_3.
\]
関係式を$x,y,z$倍して足すと$C-D=hS_3+q(C+D)$である。
この式を用いて直接整理すると
\begin{equation}\label{er:eq:central-identity}
 G=-9h\left(\frac U3 S_3+\frac W6(C+D)\right)=-9h\Omega_{h,t}.
\end{equation}
従って$h\ne0$では$\Omega_{h,t}$は標準中心元の非零スカラー倍であり、
中心正則で、その商が\eqref{er:eq:thcr}となる。
$h=0$では代数は多項式環、$U(0,t)=1$、$W(0,t)=t$であるから
$\Omega_{0,t}=f_t$である。この場合は滑らかな平面三次曲線の通常の斉次座標環であり、
同じ商の記述が成立する。
\end{proof}

\begin{remark}
$h=0$における$S_{0,t}$自体のpoint schemeは$\PP^2$であって、$E_{0,t}$ではない。
$E_{0,t}$は中心三次元で割った商の射影曲線である。
この二つを混同すると零パラメータを含む族の記述を誤る。
また中心元の公式の証明では\cite{DeLaetLeBruyn}のSection 2の表示を使う。
そのarXiv v1の冒頭のTheorem 1に印刷された表示とSection 2の表示には係数の差があるため、
\eqref{er:eq:central-identity}を別途多項式として検算する。
\end{remark}

\section{楕円曲線上のrotationの共役}\label{er:sec:rotation}
この節では$E$を滑らかな平面三次曲線、$L=\OO_E(1)$、$\sigma=t_p$とし、
$S$をnoetherianな三次元二次AS-regular代数、$g\in Z(S)_3$を正則元で
$S/(g)=B(E,L,\sigma)$となるものとする。

\begin{lemma}\label{er:lem:block}
$\MF^\ZZ(S,g)$は$\Perf(E)$とdg Morita同値である。
この同値の下で内部シフトは
\[
 \rho_p=T_{\OO_E}\circ\sigma_*\circ(-\otimes L)
       =t_{-p}^*\rho_0,
 \qquad \rho_0=T_{\OO_E}\circ(-\otimes L)
\]
に対応する。$T$はspherical twistである。
\end{lemma}
\begin{proof}
三次元AS-regular代数のparameterは三であり、三次中心正則元で割ると
Gorenstein parameterは零になる。
\cite[Theorem 5.8]{CCKM}のcokernel関手および
\cite[Theorem 2.5(iii)]{Orlov}のwindowを合成すると
\[
 \MF^\ZZ(S,g)\simeq \Dsg^{\mathrm{gr}}(S/(g))
 \simeq D^b(\qgr(S/(g)))\simeq\Perf(E)
\]
を得る。ここではcokernelとwindowのdg模型を使う。
同じ有限射影分解・dg quotientによってHom複体の比較が得られるので、
三角圏の抽象的同値だけを用いているわけではない。
THCRの右加群の規約では、tailsの次数シフトは
$\sigma_*(-\otimes L)$である。
windowを一つ移動させる際の差は$\OO_E$の評価射のconeであり、
従って残余側の次数シフトは$T_{\OO_E}\sigma_*(-\otimes L)$になる。
最後に、平行移動は$\OO_E$を保つのでそのpullbackは$T_{\OO_E}$と可換である。
\end{proof}

\begin{lemma}[次数零直線束による共役]\label{er:lem:conjugacy}
$P_q=\OO_E(q-o)$、$A_q=(-\otimes P_q)$と置くと
\begin{equation}\label{er:eq:conjugacy}
 A_q\rho_0A_q^{-1}\simeq t_q^*\rho_0.
\end{equation}
特に$\rho_p=A_{-p}\rho_0A_{-p}^{-1}$であり、$A_p$は$(\Perf(E),\rho_p)$から
$(\Perf(E),\rho_0)$への同値を与える。$p$の位数に条件はない。
\end{lemma}
\begin{proof}
テンソル関手との共役により$A_qT_{\OO_E}A_q^{-1}=T_{P_q}$である。
sourceを第一因子、targetを第二因子とするFourier--Mukai規約で、
$T_{P_q}$のkernelは
\[
 (P_q^{-1}\boxtimes P_q)\otimes I_\Delta[1]
\]
である。$E\times E$の二つの方向のfiberへの制限を計算すると、see-sawの原理により
\[
 (P_q^{-1}\boxtimes P_q)\otimes\OO(-\Delta)
 \simeq\OO(-\Gamma_q),\qquad
 \Gamma_q=\{(x,y):y=x-q\}
\]
となる。実際、第一座標$x$を固定すると両辺の次数$-1$の直線束の類は
$-x+q$であり、第二座標$y$を固定すると類は$-y-q$である。
よって$T_{P_q}=t_q^*T_{\OO_E}$。
$A_q$は$(-\otimes L)$と可換なので\eqref{er:eq:conjugacy}を得る。
\end{proof}

\begin{lemma}[周期性を含むブロック比較]\label{er:lem:periodic-block}
$E=V(f_E)$ならば、内部シフトと\eqref{er:eq:period}を保つdg Morita同値
\[
 \MF^\ZZ(S,g)\simeq\MF^\ZZ(\CC[x,y,z],f_E)
\]
が存在する。
\end{lemma}
\begin{proof}
両側を補題\ref{er:lem:block}により$\Perf(E)$と同一視し、
補題\ref{er:lem:conjugacy}の$A_p$を用いる。
これにより内部シフトをintertwineする自然同型を得る。
二つの三次周期性をこの同型に沿って比較した差は、同値関手の可逆な自然自己変換である。
$\HH^0(E)=H^0(E,\OO_E)=\CC$だから、この差は非零スカラーである。
その三乗根を用いてintertwinerを再規格化すれば、三次周期性も一致する。
負次数のHochschild cohomologyは零であり、既にfactorization側で与えられた
内部シフトのcoherenceをこの同値で移す。
従って単に自己同値の同型類を合わせるだけでなく、次節の有限群作用に必要な
周期性のデータも保存できる。
\end{proof}

\section{中心変数の追加と有限次数差の降下}\label{er:sec:suspension}
この節は、ブロックの比較を全非可換cubicへ移すための生成性を示す。
$T=\CC[u_1,\ldots,u_m]$、$A=S\otimes T$、$w=g+P$とする。
次数群を
\[
 \Gamma=(\ZZ e_0\oplus\ZZ e_1)/(3e_0-3e_1)
\]
とし、$S_d$の次数を$de_0$、$T_d$の次数を$de_1$とする。

\begin{lemma}[中心変数に対する外部積生成]\label{er:lem:TS}
$S$が非負次数付きnoetherian AS-regular代数、$g\in Z(S)_3$が正則で、
$P$が\eqref{er:eq:isolated}を満たせば、外部積はdg Morita同値
\begin{equation}\label{er:eq:TS}
 \Perf\bigl(\MF^\ZZ(S,g)\otimes_\CC\MF^\ZZ(T,P)\bigr)
 \simeq\MF^\Gamma(A,w)
\end{equation}
を与える。第一・第二因子の内部シフトはそれぞれ$(e_0)$、$(e_1)$に対応する。
\end{lemma}
\begin{proof}
外部積の微分は$D_Q\otimes1+(-1)^{\deg}1\otimes D_K$であり、その二乗は$g+P$である。
有限自由factorization間のHomを行列成分ごとに計算すると、右辺のHom複体は左辺の
二つのHom複体のtensor totalizationに一致する。
非負の内部次数と固定された有限個の生成元のshiftにより、各Hom複体は下に有界であり、
各total次数では和は有限である。これが完全忠実性を与える。

本質的全射性には$u$方向の微分を使う。
任意のfactorization$F$の微分を$D$とすると、$S$と$u_i$が独立なので
$\partial_{u_i}$は$S[u]$の導分であり、
\[
 D(\partial_{u_i}D)+(\partial_{u_i}D)D
   =(\partial_{u_i}P)\Id_F.
\]
従ってJacobianイデアル$J(P)$は$F$上で零ホモトピックに作用する。
\eqref{er:eq:isolated}より、十分大きい$d$に対し全ての$u_i^d$は$J(P)$に属する。
$u_i^d$による乗法のconeを順に取ったKoszul対象は$F$のshiftを直和因子に持つ。
特異圏でその対象を有界複体として表すと、そのcohomologyは
$u_1^d,\ldots,u_m^d$で殺される。
$(u_1,\ldots,u_m)$の冪による有限filtrationの各商は
$C=S/(g)=A/(w,u)$上の有限生成次数付き加群である。
従って右辺は、$C$-加群のrestrictionをstabilizeした対象で生成される。

ここでrestrictionが特異圏に降りることを確認する。
$u_1,\ldots,u_m$は$A/(w)$上の中心正則列である。
実際、途中で一部の$u$を零にしても$g+$残りの三次式は
$S$次数によるfiltrationで正則な最高記号$g$を持つ。
一つの$u_i$についても、先に$u_i=0$とした商でこの元が正則であることから
$u_i$が非零因子となる。
従って$C$は$A/(w)$上有限射影次元を持ち、restrictionはperfectをperfectへ送る。

$C$-加群$M$の十分高いsyzygyを取れば、その$S$上の射影次元は高々一となる。
実際、自由$C$-加群の$S$上の射影次元は一なので、syzygyを一つ取るごとに
次数二以上の$\Ext_S$は一つdimension shiftし、$S$の有限大域次元により最後には消える。
その二項$S$-射影分解と$g$の零ホモトピーは
$Q\in\MF^\ZZ(S,g)$を与える。
$M$を$A$上で分解するには、この二項分解と$u_1,\ldots,u_m$のKoszul分解をtensorすればよい。
$w$の零ホモトピーには、前者の$g$の零ホモトピーと
\[
 P=\sum_i u_i\,\frac{\partial_{u_i}P}{3}
\]
から得るKoszulホモトピーを足す。
この有界分解を周期化すると、restrictionされた$M$のstabilizationは
\[
 Q\boxtimes K_P,\qquad
 K_P=\left(\bigwedge^\bullet\CC^m\otimes T,
       \sum_i u_i\iota_i+\frac{\partial_{u_i}P}{3}\,\varepsilon_i\right)
\]
に一致する。$\varepsilon_i$は外積、$\iota_i$は縮約で、$K_P$の微分の二乗は$P$である。
二項分解の完全性、Koszul分解の完全性、およびこの式は、周期化のcohomologyを
直接同定し、余分な有限射影尾部は特異圏で零となる。

最後に、$C$は$\ZZ e_0$に次数を持つ一方、
$\Gamma/\ZZ e_0\simeq\ZZ/3$である。
任意の$\Gamma$次数付き$C$-加群をこの三剰余類に分ければ、
上の$\ZZ$次数付き議論と有限個のshiftですべての対象が得られる。
従って粗化次数でのみ現れる加群を落とさず、外部積が全生成する。
\end{proof}

\begin{lemma}[通常次数への降下]\label{er:lem:descent}
補題\ref{er:lem:periodic-block}のように第一因子の内部シフトと三次周期性を保つdg同値は、
\eqref{er:eq:TS}を通じ、通常の整数次数について
\[
 \MF^\ZZ(S[u],g+P)\simeq
 \MF^\ZZ(\CC[x,y,z,u],f_E+P)
\]
を誘導する。この同値は通常次数の$(1)$を保つ。
\end{lemma}
\begin{proof}
$\Gamma\to\ZZ$を$e_0,e_1\mapsto1$で定めると、核は$\ZZ/3$である。
その双対群$H\simeq\mu_3$は、整数次数の代数において$u$だけを三乗根倍する。
$P$の次数が三なのでcurvatureは不変である。
$\Gamma$次数付きfactorizationは、整数次数付きfactorizationの$H$-同変対象と同一視される。

その同変圏には双対群$H^\vee\simeq\ZZ/3$が文字をtensorすることで作用する。
外部積の表示では生成元は$\rho_T\rho_S^{-1}$であり、その三乗の自明化は
二つの$(3)\simeq[2]$の比較である。
第一因子の同値がこのデータを保つので、$H^\vee$-同変化にも同値を誘導する。
有限群の位数が可逆であることと冪等完備性を用い、
双対同変化で元の圏を回復する定理\cite[Theorem 4.2]{Elagin}を適用する。
同論文Section 8のdg模型、あるいは有限群のaveraging idempotentで、この比較はdg同値として行える。
$e_0$のシフトは全過程で保存されるため、通常次数の$(1)$も保存される。
\end{proof}

\subsection{主定理の証明}
\begin{proof}[定理\ref{er:thm:main}の証明]
命題\ref{er:prop:center}により$(S,g)=(S_{h,t},\Omega_{h,t})$は
補題\ref{er:lem:block}の全ての仮定を満たす。
$E=V(f_{\vartheta})$なので、周期性を保つブロック比較と中心変数の追加・降下から
\[
 \MF^\ZZ(S_{h,t}[u],\Omega_{h,t}+P)
 \simeq\MF^\ZZ(\CC[x,y,z,u],f_{\vartheta}+P)
\]
を得る。$S_{h,t}[u]$は次元$m+3$のnoetherian AS-regular代数であり、
$\Omega_{h,t}+P$は中心正則である。
その商のGorenstein parameterは$m$であり、
Orlov windowはこのMF圏を
$\langle\OO_B,\ldots,\OO_B(m-1)\rangle^\perp$へ実現する。
degree shiftはその残余成分の$\mathbf L_{\OO_B}(1)$に移る。
可換側も同じwindowを用いるので\eqref{er:eq:main-pair}を得る。

可換方程式の全偏微分が零なら、まず\eqref{er:eq:isolated}により$u=0$である。
残る$f_{\vartheta}$は$1+\vartheta^3\ne0$だからアフィン臨界点は原点のみである。
従って$Y$は射影的に滑らかである。
$\Ku(Y)$は滑らかな射影多様体のadmissible成分なのでsmooth properである。
最後に、別の滑らかな可換cubicのpairと同値ならば、二つの可換pairのenhanced同値を得る。
これはadmissible成分のFourier--Mukai同値として表され、
\cite[Theorem 1.7]{LZ}から二つのcubicは同型になる。
次元$m+1\ge3$ではその同型は超平面束を保ち、射影同型である。
\end{proof}

\begin{corollary}[任意の滑らかな三変数三次式]\label{er:cor:ternary}
滑らかな三変数三次式$g$を
$g(x)=\lambda f_t(Mx)$、$\lambda\ne0$、$M\in\mathrm{GL}_3(\CC)$とHesse形にする。
$g$のJacobian括弧を量子化し、$h'=\lambda\det(M)h$、
$\Omega_h^g(x)=\lambda\Omega_{h',t}(Mx)$で中心元を正規化する。
このとき$h'\in\mathscr U_t$では、
$S_h^g[u]/(\Omega_h^g+P)$の可換pair模型は
\[
 V\bigl(\lambda f_{\vartheta(h',t)}(Mx)+P(u)\bigr)
\]
である。さらに$Mx$を中心変数の一次形式だけ平行移動した場合も、同じ座標置換で結論を移せる。
\end{corollary}
\begin{proof}
滑らかな平面三次曲線のHesse標準形は\cite[Lemma 1]{ArtebaniDolgachev}による。
$y=Mx$の下で定数三重ベクトルは$\det(M)$倍となり、対称化は線形変換と可換なので
$S_h^g\simeq S_{h',t}$である。
商の方程式を$\lambda$で割り、定理\ref{er:thm:main}を$P/\lambda$に適用する。
最後に方程式を$\lambda$倍して元の座標に戻す。
中心線形シフトは次数付き代数の同型であり、内部シフトも保つ。
\end{proof}

\section{復元されたcubicの二次から始まる変形}\label{er:sec:second-order}
次の主張は、定理\ref{er:thm:main}で明示した可換cubicの族自体に関するものである。
非可換族の全体との相対同値を仮定しない。

\begin{theorem}[markedな二次変形の非自明性]\label{er:thm:second-order}
$t$と$P$を固定する。$Y_{h,t,P}$の方程式は
\begin{equation}\label{er:eq:second}
 f_t+P+\frac{h^2(1+t^3)}{3(1-h^2t^2/4)}xyz
\end{equation}
である。従って$h$に関して一次のmarked変形は自明だが、
$\CC[h]/(h^3)$上のmarked変形は自明でない。
$q=h^2$をパラメータとする族のKodaira--Spencer類は
\[
 \frac{1+t^3}{3}[xyz]\in
 J(f_t+P)_3,
 \qquad J(F)=\CC[x,y,z,u]/(\partial F),
\]
であり、これは非零である。
\end{theorem}
\begin{proof}
\eqref{er:eq:coefficients}から直接
\[
 \vartheta(h,t)-t=
 \frac{h^2(1+t^3)}{3(1-h^2t^2/4)}
\]
を得る。滑らかな射影超曲面のembeddedな一次変形を射影座標変換と方程式のスカラー倍で割った空間は
Jacobian環の三次部分である。Euler等式により$F$自身もJacobianイデアルに含まれる。
次元三以上では、この記述は$H^1(T_Y)$による抽象変形の接空間とも一致する。

$f_t$の三つの偏微分は$x^2+tyz,y^2+tzx,z^2+txy$であり、正則列をなす。
従ってそのJacobian環はHilbert級数$(1+z)^3$を持ち、三次部分の次元は一である。
同環で
\[
 x^2y=-t y^2z=t^2xz^2=-t^3x^2y
\]
なので、$1+t^3\ne0$から$x^2y=0$となり、同様に全ての指数型$(2,1,0)$の単項式が零となる。
また$x^3=y^3=z^3=-txyz$である。
従って三次部分は$xyz$で生成され、次元一であるから$[xyz]\ne0$。
$t=0$でも同じ議論が成立し、$xyz$は三次の唯一の平方自由単項式である。
変数が分離しているので$J(f_t+P)=J(f_t)\otimes J(P)$であり、
$[xyz]\otimes1$は非零のままである。

一次の変形項が零なので、$h^2$での座標変換が消せるのはJacobianイデアルの三次部分だけである。
$h$で始まる座標変換が一次を保つ場合、その一次項は$F$の無限小射影自己同型である。
滑らかな三次超曲面では、偏微分の一次係数syzygyが存在しないこととEuler等式により
それはスカラーEuler方向に限られ、射影的には零である。
よって非零の上記Jacobian類は$\CC[h]/(h^3)$上で消せない。
\end{proof}

\begin{remark}
この結論はmarked変形、または局所moduli stackについてのものとして述べている。
特別なcubicの有限自己同型による商を取ったcoarse moduliの局所座標で、
同じ二次の次数がそのまま最初の非零次数になるとは主張しない。
また、ファイバーごとの\eqref{er:eq:main-pair}のみから、非可換族の相対Kodaira--Spencer類を
上の式と同一視してはいない。
\end{remark}

\section{同じ可換模型を持つ非可換な代数的曲線}\label{er:sec:constant-model}
固定した$t$では復元されたcubicが動く一方、$t$も動かすと正反対の現象が現れる。

\begin{theorem}[ファイバーごとに同じpairを持つ量子化族]\label{er:thm:isocategorical}
$\tau^3\ne-1$と\eqref{er:eq:isolated}を満たす$P$を固定し、
$Y_0=V(f_\tau+P)$とする。
$(h,t)=(0,\tau)$を通る平面曲線
\begin{equation}\label{er:eq:curve}
 \mathscr C_\tau:
 h^2t^3+3\tau h^2t^2+12t+4h^2-12\tau=0
\end{equation}
のその点を含む適切な滑らかな開部分上に、次数別平坦な代数的族
$B_{h,t,P}$がある。次を満たす。
\begin{enumerate}
 \item $h$への射は$(0,\tau)$でétaleであり、全ての複素ファイバーで
 $(\mathcal K_{B_{h,t,P}},\rho_B)\simeq(\Ku(Y_0),\rho_{Y_0})$。
 \item $B_{h,t,P}$の$(0,\tau)$における次数付き結合代数の一次変形は非自明である。
 特にこの族は可換代数の定数族ではない。
 \item この曲線の関数体上で、Sklyaninの平行移動点$p_{h,t}$は無限位数である。
\end{enumerate}
ここで主張しているのはファイバーごとの同値であり、
一つの相対同値による自明化ではない。
\end{theorem}
\begin{proof}
$\vartheta(h,t)=\tau$の分母を払うと\eqref{er:eq:curve}になる。
その$t$による偏微分は$(0,\tau)$で12である。
従ってこの点を含む唯一の既約成分を取り、適切な開近傍に制限すると、
曲線は滑らかで$h$への射はétaleとなる。
さらに$t^3+1$、$U$、$b$、$\Delta$が非零となるよう縮める。
いずれも$(0,\tau)$で非零である。
ここでは$\vartheta=\tau$なので、定理\ref{er:thm:main}の可換模型は常に$Y_0$である。

代数の表示はこの曲線上の正則関数を係数とする有限表示である。
各幾何ファイバーの$S$のHilbert級数は$(1-z)^{-3}$で一定である。
滑らかな曲線上の有限生成加群のtorsionはその台でfiber次元を増加させるので、
各次数成分は局所自由になる。
中心元による各次数の乗法は全fiberで単射であり、その余核も局所自由である。
従って商$B$は次数別平坦である。

局所環の完備化では$t=\tau+O(h^2)$である。
従って$\CC[h]/(h^2)$上で
\[
 [x,y]=h(z^2+\tau xy).
\]
右辺の係数は$B_0=\CC[x,y,z,u]/(f_\tau+P)$で零でない。
実際、唯一の定義方程式は三次であり、非零の二次式を消さない。
可換な定数変形へ同型な代数は可換なので、この一次変形は自明ではない。

最後に、点曲線は一定の$E_\tau$であり、$p_{h,t}$はその上の点である。
原点$o$の局所パラメータ$z/x$で
\[
 (z/x)(p_{h,t})=-\frac h{1-ht/2}=-h+O(h^2)
\]
となる。標数零の楕円曲線の形式群では、任意の正整数$r$について
$[r](v)=rv+O(v^2)$である。
従って$[r]p$の局所パラメータは$-rh+O(h^2)$であり、恒等的に零ではない。
よって$p$は関数体上で無限位数である。
\end{proof}

\begin{corollary}[固定したHesse係数でのgenericな位数]\label{er:cor:nontorsion}
固定した$t^3\ne-1$に対しても、$p_{h,t}$は$\CC(h)$上で無限位数である。
各固定整数$r\ge1$に対し、$[r]p_{h,t}=o$となる$h\in\mathscr U_t$は有限個である。
従って非常に一般的な$h$で$p_{h,t}$は無限位数である。
\end{corollary}
\begin{proof}
滑らかな楕円曲線族の零切断に沿う形式群にも$[r](v)=rv+O(v^2)$が成り立つので、
前の証明と同じ一次係数が非torsion性を与える。
二つの切断$[r]p$と$o$の一致集合は底の閉集合であり、恒等ではない。
一次元の有限型既約底の真の閉集合は有限である。
全$r$の例外集合の和は可算なので、最後の主張が従う。
\end{proof}

\section{相対中心商と、相対復元に残る一段階}\label{er:sec:relative}
この節では$t$を固定し、
$R=\CC[h,(U(h,t)b(h,t)\Delta(h,t))^{-1}]$とする。
$\mathscr E\to\Spec R$は$f_{\vartheta}=0$で定まる滑らかな楕円曲線族、
$\mathcal L=\OO_{\mathscr E}(1)$、$p=p_{h,t}$、$\sigma=t_p$とする。
$S_R$と$\Omega_R$は式\eqref{er:eq:relations}、\eqref{er:eq:center}で定義する。

\begin{proposition}[中心商の相対THCR表示]\label{er:prop:relative-THCR}
$\mathcal L_d=\mathcal L\otimes\sigma^*\mathcal L\otimes\cdots\otimes(\sigma^{d-1})^*\mathcal L$
と置くと、自然な次数付き同型
\begin{equation}\label{er:eq:relative-THCR}
 S_R/(\Omega_R)\simeq
 R\oplus\bigoplus_{d\ge1}H^0(\mathscr E,\mathcal L_d)
\end{equation}
が存在する。右辺の積は$\sigma$で捩った切断の積である。
この同型と両辺の各次数成分の形成は、任意の可換$R$-代数への底変更と両立する。
\end{proposition}
\begin{proof}
各$S_{R,d}$は有限生成$R$-加群であり、全閉fiberにおける次元は一定である。
$R$は$\CC[h]$の局所化であるからPIDであり、torsion成分があればその台でfiber次元が増える。
従って各次数は有限自由である。
同様に$\Omega_R$の乗法の余核も全fiberでHilbert級数
$(1-z^3)/(1-z)^3$を持ち、次数別有限自由である。

$d\ge1$では各fiberの$\mathcal L_d$の次数は$3d>0$である。
楕円曲線の$H^1$の消滅とcohomology and base changeにより、
$H^0(\mathscr E,\mathcal L_d)$は階数$3d$の有限局所自由$R$-加群で、その形成は任意の底変更と可換である。
$d=0$についても$\pi_*\OO_{\mathscr E}=\OO_{\Spec R}$で同じ主張が成立する。

三つの座標切断から$S_R$の右辺への写像を作る。
二次関係式が満たされることは各複素fiberで点曲線のgraphによる表示から従う。
右辺の各次数が$R$-torsion-freeなので、この恒等式は$R$上でも成立する。
同じ理由で$\Omega_R$の像は零である。
得られた商からの写像は$h\ne0$では古典的THCR表示、$h=0$では平面三次曲線の
projective normalityによる通常のsection ring表示であり、全fiberで同型である。
各次数の有限自由加群間の写像にNakayamaの補題を適用すると\eqref{er:eq:relative-THCR}を得る。
各次数が自由であり、切断の形成も底変更と可換なので、最後の主張も従う。
\end{proof}

\begin{remark}[証明できた相対命題と残る問題]\label{er:rem:relative-gap}
命題\ref{er:prop:relative-THCR}は、非可換三次元ブロックの中心商を実在する代数的底上で記述している。
しかし、それだけで
$\MF_R(S_R,\Omega_R)$と$\Perf(\mathscr E/R)$を
任意の底変更と両立して同一視したわけではない。
さらに$P(u)$を加えた残余圏について、対角のperfect性まで得たわけでもない。

正則な底についてのOrlov比較の候補はBurke--Stevenson
\cite[Definition 5.1, Theorem 6.4]{BurkeStevenson}である。
同定理には、全体の代数の左右noetherian性、有限injective dimension、
$\RHom_A(A_0,A)$のGorenstein型、$\chi$条件が必要である。
ファイバーごとのAS-regular性だけを、これらの相対的な条件の証明として扱ってはならない。
本稿はこの経路で相対同値とその底変更compatibilityを全て確認したとは主張しない。

明示的な滑らかな族$Y\to\Spec R$の相対Kuznetsov成分はsmooth properな可換模型として存在する。
これと実際の非可換量子化の相対残余圏を同定することが次の課題である。
一般の中心混合項および複数の無限位数ブロックの全生成は、さらに別の問題として残る。
\end{remark}
\section{Sklyaninブロックと次数シフト}\label{sec:elliptic}

ここからは、前節の相対的な存在定理とは異なり、
個々の有限係数の代数を具体的に扱う。
中心三次元と楕円点曲線を同定した後、
次数シフトに含まれる平行移動を追跡し、
可換側とのrotation付き比較を構成する。

\subsection{中心元と点曲線}

\begin{lemma}\label{lem:center}
 \(G_h=x^3+y^3+z^3-2h^2\sym(xyz)\) は、すべての \(h\) で
 \(S_h\) の非零な中心正則元である。
\end{lemma}
\begin{proof}
 \(D=x^3+y^3+z^3\) とする。関係式から
 \[
 [D,x]=h(zy^2z-yz^2y)=h[zy,yz]=h^2[zy,x^2].
 \]
 また六つの順列を同じ順序へ移すと
 \[
 \sym(xyz)=xyz+\frac h2(-x^3+y^3-z^3)
 \]
 となる。この式を \(x\) と交換させ、再び二次関係式を使えば
 \(2[\sym(xyz),x]=[zy,x^2]\) を得る。
 従って \([G_h,x]=0\)。巡回置換で \(y,z\) についても従う。

 非零性は特殊化で失われない形で確認できる。
 自由代数の次数三部分の線形汎関数 \(\ell\) を
 \[
 \ell(x^3)=\ell(y^3)=\ell(z^3)=1,\quad
 \ell(xyz)=\ell(yzx)=\ell(zxy)=h/2,
 \]
 \[
 \ell(xzy)=\ell(zyx)=\ell(yxz)=-h/2
 \]
 とし、ちょうど二文字が等しい語では零とする。
 三つの二次関係式の左右に一文字を掛けた18個の式を \(\ell\) はすべて殺す。
 例えば \(\ell(z(xy-yx-hz^2))=h/2+h/2-h=0\) である。
 従って \(\ell\) は \((S_h)_3\) に降り、\(\ell(G_h)=3\) となる。
 \((a,b,c)=(1,-1,-h)\) はSklyanin代数の12個の非正則パラメータに属さず、
 \(S_h\) はAS-regular domainである \cite[Introduction]{CGWZ}。
 \(h=0\) は多項式環である。よって \(G_h\) は左右正則である。
\end{proof}

標準Sklyanin記法では \((a,b,c)=(1,-1,-h)\) である。
標準中心元の式 \cite[Equation (2.16)]{Walton}
\[
 g=c(a^3-c^3)x^3+a(b^3-c^3)xyz+b(c^3-a^3)yxz+c(c^3-b^3)y^3
\]
に代入し、\(yxz=xyz-hz^3\) を用いると \(g=-hG_h\) となる。
\(h\ne0\)、\(h^6\ne-27\) では点曲線の公式
\cite[Theorem 2.14]{Walton}から
\[
 E_h=V(x^3+y^3+z^3+h^2xyz)\subset\PP^2
\]
を得る。これは滑らかな楕円曲線であり、
\[
 S_h/(G_h)\simeq B(E_h,L,\sigma_h),\qquad L=\OO_{E_h}(1),
\]
と書ける。ここで \(\sigma_h=t_p\) は平行移動である。
原点をinflection point \(o=[1:-1:0]\) に取ると、
多重線形化した関係式を第一成分 \(o\) で評価した核は
\(p=[1:-1:-h]\) である。従ってこの規約で \(\sigma_h=t_p\) となる。
この中心商と楕円曲線の比較は \cite[Sections 3.5--3.6]{EG} も参照。
\(h=0\) は通常のFermat平面三次曲線の斉次座標環である。
この場合は \(p=o\)、\(\sigma_0=\Id\) とする。

\subsection{rotationに含まれる平行移動を消す}

\begin{lemma}\label{lem:elliptic-conjugacy}
 \(E\) を滑らかな楕円曲線、\(L\) を次数三の直線束、
 \(\sigma=t_p\) とする。
 \(C=B(E,L,\sigma)\) のgraded singularity categoryを \(D^b(E)\) と同一視したとき、
 その次数シフトは可換平面三次曲線のrotation
 \(\rho_0=T_{\OO_E}\circ(-\otimes L)\) と共役である。
\end{lemma}
\begin{proof}
 右加群の規約を
 \[
 \Gamma_*(F)_d=H^0(E,F\otimes L_d),\qquad
 L_d=L\otimes\sigma^*L\otimes\cdots\otimes(\sigma^{d-1})^*L
 \]
 と取る。tailsの次数シフトは \(H(F)=\sigma_*(F\otimes L)\) である。
 Gorenstein parameter零のOrlov windowを比較すると
 \[
 \rho_{\mathrm{nc}}=T_{\OO_E}\circ H=t_{-p}^*\rho_0.
 \tag{2}\label{eq:rho-nc}
 \]
 このwindow公式を確認する。特異圏への商を \(q\) とし
 \(\Psi_i=qR\Gamma_{\ge i}\) と置くと、
 \((1)\Psi_1=\Psi_0H\)、\(q(k)=\Psi_0(\OO_E)[1]\)、
 \[
 \Psi_1F\longrightarrow\Psi_0F
 \longrightarrow\RHom(\OO_E,F)\otimes q(k)
 \]
 がある。楕円曲線のSerre双対から
 \(\Psi_1=\Psi_0T_{\OO_E}^{-1}\) となり、\eqref{eq:rho-nc}を得る。

 原点を \(o\)、\(P_q=\OO_E(q-o)\)、\(A_q=(-\otimes P_q)\) とする。
 spherical twistの核は \(I_\Delta[1]\) なので、
 \[
 K_{T_{P_q}}=(P_q^{-1}\boxtimes P_q)\otimes I_\Delta[1].
 \]
 see-sawを両方向のファイバーへ適用すると
 \[
 (P_q^{-1}\boxtimes P_q)\otimes I_\Delta
 \simeq I_{\{y=x-q\}}.
 \]
 よって \(T_{P_q}=t_q^*T_{\OO_E}\) であり、
 \[
 A_q\rho_0A_q^{-1}=t_q^*\rho_0,\qquad
 \rho_{\mathrm{nc}}=A_{-p}\rho_0A_{-p}^{-1}.
 \tag{3}\label{eq:conjugacy}
 \]
 これらは核とwindowによるdg同値として実現される。
 三次周期性の二つの同型の差は \(\HH^0(D^b(E))^\times=\CC^\times\) のscalarであり、
 共役の自然同型をその三乗根で調整すれば周期性も両立させられる。
 \(\HH^{<0}(D^b(E))=0\) なので、この自然同型の等式に追加の高次ホモトピー障害はない。
 後で用いる有限次数差の作用は、この周期性から構成する。
\end{proof}

\begin{corollary}\label{cor:block}
 \(h^6\ne-27\) のとき、三次周期性と次数シフトを保つdg Morita同値
 \[
 \MF^{\ZZ}(S_h,G_h)
 \simeq
 \MF^{\ZZ}\bigl(\CC[x,y,z],x^3+y^3+z^3+h^2xyz\bigr)
 \]
 が存在する。
\end{corollary}
\begin{proof}
 \cite[Theorem 3.6.2]{EG}、Orlov window、
 \cite[Theorem 5.8]{CCKM}によって両辺を \(D^b(E_h)\) と比較し、
 補題\ref{lem:elliptic-conjugacy}を使う。
 MFと特異圏の比較は、factorizationを商環上の両側無限の完全周期分解へ送る
 dg関手で実現する。中心正則元の商はGorensteinであり、complete resolutionは
 MCM加群の安定圏のdg模型を与える。
 CCKMの完全忠実性をすべてのshiftに適用すると、Hom複体の比較がquasi-isomorphismとなる。
 周囲がAS-regularであることによる高次syzygyのfactorization表示が本質的全射性を与える。
 特異圏からwindowへのdg比較は \cite[Section 5, Theorem 5.8]{CT} による。
 \(h=0\) では恒等の可換比較を用いる。
\end{proof}

\section{中心三次式を加える比較と通常次数への降下}\label{sec:TS}

単独の楕円ブロックの比較から中心変数を加えた模型へ進む。
外部積の完全忠実性と生成性を分けて示し、
最後に多重次数から通常の整数次数への降下を行う。

\subsection{二つの次数ブロックと外部積}

中心変数側の三次式をFermat型に制限する必要はない。
この節では \(T=\CC[u_1,\ldots,u_m]\)、\(m\ge1\)、とし、
\(P\in T_3\) は
\[
 J(P)=\left(\frac{\partial P}{\partial u_1},\ldots,
                 \frac{\partial P}{\partial u_m}\right),\qquad
 \dim_\CC T/J(P)<\infty                                      \tag{4a}\label{eq:central-isolated}
\]
を満たすものとする。これは斉次三次式 \(P\) のアフィン臨界点が
原点だけであるという条件である。
\[
 \Gamma=(\ZZ e_0\oplus\ZZ e_1)/\langle3e_0-3e_1\rangle,
 \qquad w=3e_0=3e_1                                       \tag{4b}\label{eq:central-gamma}
\]
と置く。\(S_d\) の次数を \(de_0\)、すべての \(u_i\) の次数を \(e_1\) とする。
したがって中心変数どうしの混合項を含む \(P\) も次数 \(w\) である。

\begin{lemma}[任意の孤立中心三次式に対する外部積]\label{lem:TS}
 \(S\) を連結非負次数付きnoetherian AS-regular \(\CC\)-代数、
 \(g\in S_3\) を中心正則元とする。
 \eqref{eq:central-isolated}の下で、外部積はpretriangulated hullと
 冪等完備化の後にdg Morita同値
 \[
 \MF^\ZZ(S,g)\otL_\CC\MF^\ZZ(T,P)
 \simeq
 \MF^\Gamma(S\otimes_\CC T,g+P)                             \tag{4c}\label{eq:central-TS}
 \]
 を与える。この同値は二つの因子の次数シフトと、
 両者の三次内部シフトを \([2]\) と同一視する周期性を保つ。
\end{lemma}
\begin{proof}
 \(R=S\otimes_\CC T=S[u_1,\ldots,u_m]\)、\(W=g+P\)、
 \(A=R/(W)\)、\(C=S/(g)\) と置く。
 まず \(W\) は中心正則元である。
 実際、\(u\) に関する単項式順序を固定すると、\(P\ne0\) の最高項の係数は
 非零のscalarであり、任意の非零多項式との積の最高項は消えない。
 以下では中心正則元の商に対するMFとgraded singularity categoryの比較
 \cite[Theorem 5.8]{CCKM}を用いる。
 \(\Gamma\to\ZZ\)、\(e_0,e_1\mapsto1\)、は有限の核を持ち、
 正の整数次数を与えるので、有限次数付き自由分解による同じ比較は
 \(\Gamma\)-次数についても適用できる。

 \paragraph{外部積と完全忠実性.}
 factorizationの規約を
 \[
 Q=(Q_1\xrightarrow{\varphi}Q_0
                  \xrightarrow{\psi}Q_1(w)),\qquad
 \psi\varphi=g\Id,\quad\varphi\psi=g\Id
 \]
 とする。これを \(Q^{i+2}=Q^i(w)\) を満たす曲がった複体に展開すると、
 \(Q^0=Q_0\)、\(Q^1=Q_1(w)\) である。
 二つのfactorizationの外部積は、二つの奇微分をKoszul符号付きで足すことにより定める。
 奇微分の外部和の平方は \(g+P\) である。
 二つの \([2]\) を同じ内部シフトに対応させる関係が
 \(3e_0=3e_1\) であり、得られるfactorizationは有限自由である。

 有限自由対象 \(Q,Q'\) と \(K,K'\) に対して、
 係数を \(S\)-次数と \(T\)-次数に分解すると
 \[
 \Hom_{\MF^\Gamma}(Q\boxtimes K,Q'\boxtimes K')
 \simeq
 \operatorname{Tot}\!\left(
   \Hom_{\MF^\ZZ(S,g)}(Q,Q')\otimes_\CC
   \Hom_{\MF^\ZZ(T,P)}(K,K')\right)                         \tag{4d}\label{eq:central-hom}
 \]
 を得る。具体的には、左辺の次数 \(r\) の斉次写像は、
 \(a+b=r\) を満たす右辺の次数 \((a,b)\) の写像の和に一意に分解する。
 二つの内部次数の差を移す操作は \((3,-3)\) の整数倍だけであり、
 これは二つのcohomological degreeを \((2,-2)\) だけ変える操作に対応する。
 微分と合成の符号も外部和のKoszul符号に一致する。
 各因子のHom複体は、有限個の自由加群の次数シフトと係数代数の非負次数性から
 下に有界である。従って各全cohomological degreeに寄与する項は有限個であり、
 \eqref{eq:central-hom}に完備化は必要ない。
 これが外部積のdg完全忠実性を与える。

 \paragraph{中心変数を零とした商への生成性の帰着.}
 \(\MF^\Gamma(R,W)\) の有限自由対象の奇微分を \(D\) とする。
 \(u_i\) は中心変数なので、係数行列の \(u_i\)-微分は \(R\) 上の導分であり、
 \[
 D\,\frac{\partial D}{\partial u_i}
 +\frac{\partial D}{\partial u_i}\,D
 =\frac{\partial P}{\partial u_i}\Id
 \]
 が成り立つ。内部シフトを付けると、これは
 \(\partial_iP:N\to N(2e_1)\) の零ホモトピーである。
 斉次イデアル \(J(P)\) は \((u_1,\ldots,u_m)\)-primaryなので、
 ある整数 \(d\ge2\) に対してすべての \(u_i^d\) が \(J(P)\) に属する。
 斉次係数でこれらを偏微分の線形結合に書けば、
 \(u_i^d:N\to N(de_1)\) も特異圏で零となる。

 \(N\) を有限生成 \(A\)-加群の有界複体で表し、
 \(u_1^d,\ldots,u_m^d\) に沿うKoszul複体を \(K(u^d;N)\) と書く。
 零写像のconeは元の対象を直和因子に持つため、\(N\) は
 \(K(u^d;N)\) の直和因子である。
 より正確には、一変数ごとのKoszul cone関手は完全関手でretractを保つので、
 各 \(u_i^d\) の零作用から得られるretractを順に合成すればよい。
 このKoszul複体のcohomologyは \((u_1^d,\ldots,u_m^d)\) に殺され、
 \[
 (u_1,\ldots,u_m)^{m(d-1)+1}\subset(u_1^d,\ldots,u_m^d)
 \]
 である。従って各cohomology加群は、有限個の \(C=A/(u_1,\ldots,u_m)\)-加群を
 商とする有限filtrationを持つ。標準切断の三角形も使うと、
 右辺の特異圏は \(C\)-加群の制限とその次数シフトで厚く生成される。

 ここで制限が特異圏へ降りることを確認する。
 \(u_1,\ldots,u_m\) は \(A\) の中心正則列である。
 一つの変数 \(u_i\) について、\(R\) で \(u_i a=Wb\) とする。
 \(u_i=0\) とした商では、\(\overline W=g+P|_{u_i=0}\) は正則である。
 特殊化した \(P\) が非零なら最高項のscalar係数を使い、
 零なら \(g\) の正則性を使えばよい。
 従って \(\overline b=0\)、すなわち \(b=u_i b_1\) であり、
 多項式環 \(R\) で \(u_i\) を消去して \(a=Wb_1\) を得る。
 先行する変数を順に零とした同じ議論が正則列を与える。
 よって \(C\) は有限自由Koszul分解を持つ \(A\)-perfect加群であり、
 制限は \(\Dsg(C)\to\Dsg(A)\) を誘導する。

 \paragraph{制限加群の安定化の明示.}
 Euler公式から
 \[
 q_i=\frac13\frac{\partial P}{\partial u_i},\qquad
 P=\sum_{i=1}^{m}u_iq_i
 \]
 と書ける。\(T\) 上の通常のKoszul分解を
 \[
 K_j=T(-je_1)\otimes_\CC\Lambda^j\CC^m,
 \qquad d_K=\sum_i u_i\iota_i
 \]
 とする。\(\iota_i\) は外積代数の縮約、\(\varepsilon_i\) は対応する基底による左外積である。
 \[
 h_K=\sum_iq_i\varepsilon_i:K_j\longrightarrow K_{j+1}(w)
 \]
 と置くと
 \[
 d_Kh_K+h_Kd_K=P\Id,\qquad h_K^2=0.                         \tag{4e}\label{eq:central-koszul}
 \]
 ここで \(q_i\) の次数は \(2e_1\) であり、
 \(K_{j+1}(w)\) のシフトによりすべての表示写像は内部次数零である。

 有限Koszul複体を偶奇で折り畳み、
 \[
 (K_P)_0=\bigoplus_{j\ge0}K_{2j}(jw),\qquad
 (K_P)_1=\bigoplus_{j\ge0}K_{2j+1}(jw)
 \]
 と置く。\(d_K+h_K\) は
 \[
 (K_P)_1\longrightarrow(K_P)_0
             \longrightarrow(K_P)_1(w)
 \]
 を与え、\eqref{eq:central-koszul}によってこれは \(P\) の有限自由factorizationである。

 \(M\) を \(C\)-加群とし、その特異圏での像を表すfactorizationを
 \(Q=(Q_1\xrightarrow\varphi Q_0\xrightarrow\psi Q_1(w))\) とする。
 有限回のsyzygyとcohomological shiftを取れば
 \(M=\operatorname{coker}\varphi\) としてよい。
 実際、元の加群との差に現れる \(C\)-perfect対象は制限後も \(A\)-perfectである。
 \(\psi\varphi=g\Id\) と \(g\) の正則性から \(\varphi\) は単射であり、
 \[
 0\longrightarrow Q_1\xrightarrow\varphi Q_0\longrightarrow M\longrightarrow0
 \]
 は \(S\)-自由分解である。
 これをhomological degree \(1,0\) に置き、
 \(d_Q=\varphi\)、\(s_Q:Q_0\to Q_1(w)\) を \(\psi\) とする。
 \(s_Q\) は \(Q_1\) 上では零とする。

 \(u_iM=0\) として \(M\) を \(R\)-加群に制限したときの有限自由分解は
 \[
 F_a=\bigoplus_{b+c=a}Q_b\otimes_\CC K_c,
 \qquad d=d_Q\otimes1+(-1)^b\otimes d_K
 \]
 である。\(b\in\{0,1\}\) は \(Q\) 側のhomological degreeである。
 \[
 s=s_Q\otimes1+(-1)^b\otimes h_K:F_a\longrightarrow F_{a+1}(w)
 \]
 と置くと、異なる因子の奇作用の反可換性から
 \[
 ds+sd=W\Id,\qquad s^2=0.                                  \tag{4f}\label{eq:central-shamash}
 \]
 が成り立つ。

 このデータの周期化が実際に \(A\)-自由分解を与えることも確認しておく。
 \(F_a=0\) for \(a<0\) とし、
 \[
 \mathcal F_n=
 \bigoplus_{j\ge0}(F_{n-2j}\otimes_R A)(-jw),\qquad n\ge0
 \]
 と置く。微分は同じ \(j\) 内の \(d\) と、\(j\ge1\) について
 \[
 s:(F_{n-2j}\otimes_R A)(-jw)
 \longrightarrow(F_{n-2j+1}\otimes_R A)(-(j-1)w)
 \]
 の和である。\eqref{eq:central-shamash}により平方は零になる。
 各 \(n\) で和は有限である。

 \(j\) による増大濾過を使うと、第一ページは
 \(H_*(F\otimes_R A)\) の内部・homological shiftの和となる。
 \(W\) が正則で \(WM=0\) なので、
 \[
 H_i(F\otimes_R A)=
 \begin{cases}
 M&i=0,\\
 M(-w)&i=1,\\
 0&i\ne0,1.
 \end{cases}
 \]
 次の微分は、\(j\ge1\) の \(H_0\) を \(j-1\) の \(H_1\) へ送る同型である。
 これを \(ds+sd=W\Id\) から直接見られる。
 \(a\in F_0\) が \(M\) の元を表すなら、\(s(a)\) の \(A\) 上の像はcycleであり、
 \(d s(a)=Wa\) である。
 短完全列
 \[
 0\longrightarrow F(-w)\xrightarrow{W}F
      \longrightarrow F\otimes_R A\longrightarrow0
 \]
 のconnecting mapはそのcycleを元の \(a\) に送る。
 従ってこの微分は、上の \(H_1\simeq M(-w)\) の標準同定の下で恒等である。
 隣接する二項はすべて相殺し、\(j=0\) の \(H_0=M\) だけが残る。
 濾過は各全homological degreeで有限なので収束し、
 \(\mathcal F\) は \(M\) の \(A\)-自由分解となる。

 \(F\) は有限複体なので、この分解の十分先は二周期的である。
 そのfactorizationは
 \[
 \widehat F_0=\bigoplus_{j\ge0}F_{2j}(jw),\qquad
 \widehat F_1=\bigoplus_{j\ge0}F_{2j+1}(jw),
 \qquad d+s
 \]
 による折り畳みであり、Koszul符号と内部シフトを含めて
 \(Q\boxtimes K_P\) に一致する。
 実際、\(\widehat F_0\) は
 \(Q_0\otimes(K_P)_0\oplus Q_1\otimes(K_P)_1(w)\)、
 \(\widehat F_1\) は
 \(Q_1\otimes(K_P)_0\oplus Q_0\otimes(K_P)_1\) である。
 従って制限加群の安定化は
 \[
 \operatorname{stab}_A(M)\simeq Q\boxtimes K_P.               \tag{4g}\label{eq:restriction-mf}
 \]
 ここで全体のshiftにも曖昧さはない。
 十分大きい \(\ell\) について
 \(\mathcal F_{2\ell}=(\widehat F_0\otimes_R A)(-\ell w)\)、
 \(\mathcal F_{2\ell+1}=(\widehat F_1\otimes_R A)(-\ell w)\) であり、
 そのcokernelは \(\Omega_A^{2\ell}M\) である。
 \(\Omega_A^{2\ell}M\simeq M[-2\ell]\) と \((w)\simeq[2]\) によって
 この内部シフトはちょうど打ち消される。

 最後に、\(\Gamma/\ZZ e_0\simeq\ZZ/3\) であるため、
 任意の \(\Gamma\)-次数付き \(C\)-加群は、
 三つの次数類 \(0,e_1,2e_1\) に属する部分に分解する。
 各部分は通常の \(\ZZ\)-次数付き \(C\)-加群の内部シフトであり、
 \eqref{eq:restriction-mf}の右辺の \(T\) 側を \((j)\)、\(j=0,1,2\)、
 だけシフトすればすべて得られる。
 前段のKoszul filtrationと合わせて、外部積の像は右辺全体を厚く生成する。
 これで本質的全射性が従う。
 周期性および二つの次数シフトとの両立は、ここで用いた外部和の構成そのものから従う。
\end{proof}

\begin{remark}[一変数の場合の明示行列]\label{rem:central-one-variable}
 \(m=1\)、\(P=u^3\) では \(K_P=(u,u^2)\) である。
 上の分解とホモトピーは
 \[
 Q_1(-e_1)\xrightarrow{(-u,\varphi)^t}
 Q_1\oplus Q_0(-e_1)\xrightarrow{(\varphi,u)}Q_0,
 \]
 \[
 s_0=(\psi,u^2)^t,\qquad s_1=(-u^2,\psi)
 \]
 となる。折り畳んだfactorizationは
 \[
 \alpha=
 \begin{pmatrix}\varphi&u\\-u^2&\psi\end{pmatrix}:
 Q_1\oplus Q_0(-e_1)\longrightarrow Q_0\oplus Q_1(w-e_1),
 \]
 \[
 \beta=
 \begin{pmatrix}\psi&-u\\u^2&\varphi\end{pmatrix}:
 Q_0\oplus Q_1(w-e_1)\longrightarrow
 Q_1(w)\oplus Q_0(w-e_1).
 \]
 中心性により \(\alpha\beta=\beta\alpha=(g+u^3)\Id\) である。
 これらの式は、一般の \(K_P\) による安定化の符号と次数シフトの具体的な検算になる。
\end{remark}

\subsection{有限群による整数次数の回復}

\(\Gamma\to\ZZ\)、\(e_0,e_1\mapsto1\)、の核は
\(T_\Gamma\simeq\ZZ/3\) であり、\(e_1-e_0\) で生成される。
分裂 \(1\mapsto e_0\) の下で、その双対群
\(H_\Gamma=\Hom(T_\Gamma,\CC^\times)\simeq\mu_3\) は
\(S\) 上恒等、すべての \(u_i\) を同時に三乗根倍する作用である。
\(P\) は斉次三次式なのでこの作用は \(W=g+P\) を保ち、
\[
 \MF^\Gamma(R,W)
 \simeq\bigl(\MF^\ZZ(R,W)\bigr)^{H_\Gamma}.                \tag{5}\label{eq:eq}
\]
有限群の位数は可逆である。
同変な有限次数付き自由加群は、連結な整数次数と有限群表現の半単純性により
同変な斉次基底を持つので、この比較は有限自由factorizationのdg模型で実現する。
さらに双対群 \(T_\Gamma\) で同変化すると、元の通常次数の圏が回復される
\cite[Theorem 4.2, Section 8, Lemma 8.6, Theorem 8.9]{Elagin}。
冪等完備化はこの比較の一部として取る。

外部積 \eqref{eq:central-TS} の左辺で双対作用の生成元は
\[
 \rho_T\rho_S^{-1}
\]
である。ここで \(\rho_S,\rho_T\) はそれぞれの因子の内部次数一のシフトである。
その三乗を恒等と同一視する自然同型は、両因子の
\((3)\simeq[2]\) から得られる。
二つのperiodicityを含めて \(\rho_S\) を保つdg同値はこの作用と両立し、
\eqref{eq:eq}と双対同変化により通常次数へ降りる。
この際 \(e_0\) のシフトも保存されるので、通常次数のdegree shift \((1)\) が保存される。
有限群作用のcoherenceはfactorizationの内部シフトのcoherenceから引き継ぎ、
単に自己同値の三乗が同型になるという存在命題だけから群作用を作るのではない。

\subsection{Sklyaninブロックへの適用}

\begin{proof}[定理\ref{thm:global}の証明]
 補題\ref{lem:TS}を \((S_h,G_h)\) と
 \((\CC[x,y,z],x^3+y^3+z^3+h^2xyz)\) のそれぞれに適用する。
 系\ref{cor:block}は第一因子の次数シフトと三次周期性を保ち、
 \(\MF^\ZZ(T,P)\) 側では恒等を用いる。
 前節の降下によって通常次数のdegree shiftを保つdg Morita同値
 \[
 \MF^\ZZ(S_h[u],G_h+P)
 \simeq
 \MF^\ZZ\!\left(\CC[x,y,z,u],x^3+y^3+z^3+h^2xyz+P\right)
 \]
 を得る。体上のOrlov同値と \cite{CCKM} により、これは残余圏の比較である。

 可換方程式の \(u_i\)-偏微分がすべて零なら、
 \eqref{eq:central-isolated}によって全 \(u_i=0\) である。
 \(h^6\ne-27\) の下では残る平面三次式のアフィン臨界点も原点だけである。
 よって射影超曲面 \(Y_{h,P}\) は滑らかである。
 次元は \(m+1\ge3\) なので、一意性は \cite[Theorem 1.7]{LZ} による。
\end{proof}



\begin{corollary}[平行移動を伴う一般の楕円ブロック]\label{cor:elliptic-central-cubic}
 \(S\) をHilbert series \((1-t)^{-3}\) を持つnoetherianな
 二次三次元AS-regular \(\CC\)-代数、\(g\in S_3\) を中心正則元とする。
 滑らかな平面三次曲線 \(E\)、\(L=\OO_E(1)\)、および
 ある原点に関する平行移動 \(\sigma=t_p\) に対して、次数付き同型
 \[
 S/(g)\simeq B(E,L,\sigma)
 \]
 が存在すると仮定する。
 \(f_E\in\CC[x,y,z]_3\) を \(E\subset\PP^2\) の方程式とし、
 \(P\in\CC[u_1,\ldots,u_m]_3\)、\(m\ge2\)、は
 \eqref{eq:central-isolated}を満たすとする。
 このとき \(S[u]/(g+P)\) のrotation付き残余圏の滑らかな可換モデルは
 \[
 V(f_E(x,y,z)+P(u))\subset\PP^{m+2}
 \]
 であり、射影同型を除いて一意である。
\end{corollary}
\begin{proof}
 Gorenstein parameter零のOrlov同値により
 \(\MF^\ZZ(S,g)\) は \(D^b(E)\) と同値であり、
 次数シフトは補題\ref{lem:elliptic-conjugacy}の
 \(T_{\OO_E}\sigma_*(-\otimes L)\) に対応する。
 同補題の直線束による共役とperiodicityのscalar調整により、
 \(\MF^\ZZ(\CC[x,y,z],f_E)\) との、次数シフトと三次周期性を保つdg同値を得る。
 あとは補題\ref{lem:TS}と有限次数差による降下を適用すればよい。
 滑らかさと一意性は定理\ref{thm:global}と同じ証明による。
 方程式 \(f_E\) の非零scalar倍による違いは、
 分離された \(x,y,z\) の同時三乗根倍によって射影同型に吸収される。
\end{proof}

\begin{remark}[この拡張の範囲]
 定理\ref{thm:global}では中心変数間の任意の三次混合項を許すが、
 非可換ブロックの生成元と中心変数の混合項は含めない。
 別の非可換ブロックについては、その生成元による通常の偏微分が代数の導分とは限らず、
 本節の \(J(P)\)-primaryな支持への帰着は適用できない。
 この点は独立した生成性の問題として残る。
 本節の拡張の先行研究上の新規性は未確認である。
\end{remark}


\section{複数ブロックの比較と生成性の帰着}\label{sec:multiblock}

複数の非可換ブロックでは、外部積関手の完全忠実性と
その像が全MF圏を生成することは別の主張である。
本節では、まず完全忠実性を構成し、
本質的全射性のために必要な支持と有限性の条件を明示する。

\subsection{生成性を仮定しない外部積}

\(A_1,\ldots,A_s\) を、\((A_i)_0=\CC\) を満たす局所有限な
非負次数付き代数、\(f_i\in Z(A_i)_3\) を中心的な三次斉次元とする。
この小節の \(\MF\) は、有限自由次数付きmatrix factorizationのdg圏の
pretriangulated hullを取り、さらに冪等完備化したものを表す。
この定義にはtensor product algebraのnoetherian性を仮定しない。
\[
 \Gamma_s=\ZZ^s/\langle3e_i-3e_1:2\le i\le s\rangle,
 \qquad w=3e_1,
 \qquad A=\bigotimes_{i=1}^s A_i,
 \qquad W=\sum_{i=1}^s f_i
\]
とする。\(A_i\) の次数 \(d\) の斉次元には \(de_i\) を割り当てる。

\begin{lemma}[外部積の完全忠実性]\label{lem:multiblock-ff}
 外部積はdg完全忠実関手
 \[
 \Phi:\Perf\!\left(\bigotimes_{i=1}^s\MF^{\ZZ}(A_i,f_i)\right)
 \longrightarrow\MF^{\Gamma_s}(A,W)
 \]
 を定める。\(i\) 番目の因子のdegree shiftは右辺の \((e_i)\) に対応し、
 \((3e_i)=[2]\) の標準同型とも両立する。
\end{lemma}
\begin{proof}
 まず有限自由factorization \(P_i,Q_i\) を取る。
 tensor productの奇微分はKoszul符号付きの外部和であり、
 その二乗は \(W\) となる。
 各因子のHom複体を \(C_i^\bullet\) と書くと、
 自由加群のHomを行列係数ごとに分解することにより
 \[
 \Hom_{\MF^{\Gamma_s}(A,W)}^\bullet
       (\boxtimes_iP_i,\boxtimes_iQ_i)
 \cong
 \operatorname{Tot}\!\left(\bigotimes_i C_i^\bullet\right)
 \tag{\(\dagger\)}\label{eq:multiblock-hom}
 \]
 を得る。ここでは次数群の関係 \(3e_i=3e_1\) が、
 各因子の二つのcohomological shiftを他の因子へ移す操作に正確に対応する。
 従って右辺の次数 \(q\) の和は
 \(q_1+\cdots+q_s=q\) にわたる通常のtensor product複体の和であり、
 追加の完備化や直積を取らない。

 この和の有限性を確認する。各 \(P_i,Q_i\) は有限個の内部次数shiftからなり、
 \(A_i\) は非負次数なので、\(C_i^q=0\) は十分小さい \(q\) で成り立つ。
 したがって固定した全次数 \(q\) に寄与する \((q_1,\ldots,q_s)\) は有限個である。
 同様に、行列係数を元の \(\ZZ^s\)-次数へ分解したときも、
 固定した \(\Gamma_s\)-次数と非負次数の条件の下では有限個の成分しか現れない。
 これにより \eqref{eq:multiblock-hom} は実際のHom複体の同型である。
 differentialとcompositionの一致はtensor productのKoszul符号から従う。
 完全忠実性は有限個のcone、shift、直和因子を取っても保たれるので、主張を得る。
 内部次数shiftとperiodicityとの両立はこの構成自体に含まれる。
\end{proof}

\subsection{任意個の楕円THCRの積と周囲の正則性}

Chanの多重Serre同値 \cite[Theorem 3.4]{ChanMulti} は、代数のnoetherian性を
仮定せず、right ampleな可逆bimoduleの組に適用できる。
同論文のCorollaries 5.8--5.9は二因子のTHCRと三次元AS-regular代数の積を扱う。
以下では、その証明で用いる面と帯の議論をfactor数について帰納し、任意個の楕円因子を扱う。
Keeler \cite[Theorem 2.7, Corollary 2.8]{KeelerMulti} によっても、
この積のbimoduleの組の左右NC-amplenessを確認できる。
各translationはnumerical Picard groupに自明に作用し、全ての座標が正の外部積はampleだからである。
ただし多重amplenessだけではnoetherian性は自動的に従わず、以下で帯の有限性を別に扱う。

\begin{lemma}[楕円THCRの有限テンソル積]
\label{lem:elliptic-thcr-tensor-noeth}
 \(E_1,\ldots,E_r\) を滑らかな複素射影楕円曲線、\(L_i\) を正次数の可逆層、
 \(\sigma_i\) をtranslationとする。このとき
 \[
 C=\bigotimes_{i=1}^r C_i,\qquad C_i=B(E_i,L_i,\sigma_i)
 \]
 は左右noetherianである。
\end{lemma}
\begin{proof}
 \(X=\prod_iE_i\) とし、各factorから引き戻した可逆bimoduleを
 \(\mathcal L_i\) と書く。異なるfactorの作用は可換であり、
 \(\mathcal L^{\mathbf d}=\mathcal L_1^{d_1}\cdots\mathcal L_r^{d_r}\) と置けば
 \[
 C_{\mathbf d}=H^0(X,\mathcal L^{\mathbf d})\qquad(\mathbf d\in\mathbb N^r)
 \]
 は \(C\) のテンソル積多重次数と乗法を与える。
 可逆bimoduleの逆を用いて \(\mathcal L^{\mathbf d}\) は
 \(\mathbf d\in\mathbb Z^r\) に対しても定義する。

 \emph{Ampleness.}
 外部積のvery ample可逆層 \(H=\boxtimes_iH_i\) を固定する。
 任意のcoherent層 \(\mathcal F\) を、有限個の \(H\) の負冪の直和からの全射で
 \(\dim X\) 段まで近似する。
 各 \(\mathcal L_i^{d_i}\) の台となる可逆層の次数は \(d_i\deg L_i\) なので、
 その各近似項を \(\mathcal L^{\mathbf d}\) でtwistした高次cohomologyは、
 \(d_i\) が全て十分大きいと消える。
 ここでは楕円曲線上の正次数の可逆層の \(H^1\) の消滅とK\"unneth公式を用いる。
 最後の核にはGrothendieck消滅を適用し、dimension shiftingによって
 \[
 H^q(X,\mathcal F\mathcal L^{\mathbf d})=0
 \quad(q>0,\ \mathbf d\gg0)
 \]
 を得る。従ってこの組はright ampleである。
 \(\mathcal F\otimes H^{-q}\) にもこの消滅を適用すると、
 Castelnuovo--Mumford regularityにより
 \(\mathcal F\mathcal L^{\mathbf d}\) の大域生成も \(\mathbf d\gg0\) で従う。
 自動同型によるpushforwardはcohomologyと大域生成を変えないので、
 この議論は可逆bimoduleの左右の規約に依存しない。

 \emph{面と有限幅の帯.}
 代数 \(C\) の右noetherian性と、全てのcoherent \(\mathcal F\)、\(q\ge0\) に対する
 \[
 Q^q(\mathcal F)=\bigoplus_{\mathbf d\in\mathbb N^r}
 H^q(X,\mathcal F\mathcal L^{\mathbf d})
 \]
 の有限生成性を、\(r\) について同時に示す。
 \(r=0\) の場合は有限次元ベクトル空間の主張である。
 \(r-1\) 個の因子まで成立すると仮定し、
 \(C_{\widehat i}=\bigotimes_{j\ne i}C_j\)、\(X_{\widehat i}=\prod_{j\ne i}E_j\) と置く。
 座標 \(d_i=a\) を固定した \(Q^q(\mathcal F)\) の面は、\(C_{\widehat i}\)-加群である。
 射影 \(\pi_i:X\to X_{\widehat i}\) と
 \(\mathcal F_i=\mathcal F\mathcal L_i^a\) に対するLerayスペクトル系列の
 全多重次数の和は
 \[
 E_2^{p,b}=Q^p_{X_{\widehat i}}
       (R^b\pi_{i*}\mathcal F_i)
 \ \Longrightarrow\
 \bigoplus_{\mathbf d_{\widehat i}\ge0}
 H^{p+b}(X,\mathcal F_i\mathcal L_{\widehat i}^{\mathbf d_{\widehat i}})
 \]
 となる。射影公式と残るtranslationに関する自然性により、これは
 \(C_{\widehat i}\)-加群のスペクトル系列である。
 帰納法で各 \(E_2\) 項はnoetherianな \(C_{\widehat i}\) 上有限生成である。
 従って、有限個の核・商・拡張で得られる各面もnoetherianである。

 次の帯に関する観察を使う。
 \(\mathbb N^r\)-次数付き \(C\)-加群 \(M\) の各面 \(M_{d_i=a}\) が
 \(C_{\widehat i}\) 上noetherianであり、\(M_{\ge\mathbf b}=0\) とする。
 \[
 M^{(i)}=\bigoplus_{d_1\ge b_1,\ldots,d_i\ge b_i}M_{\mathbf d}
 \quad(0\le i\le r)
 \]
 は \(M=M^{(0)}\supset\cdots\supset M^{(r)}=0\) という \(C\)-部分加群のfiltrationである。
 各 \(M^{(i-1)}/M^{(i)}\) をさらに \(d_i\ge a\) でfiltrationすると、
 有限個の商は固定面の \(C_{\widehat i}\)-部分加群になる。
 その商では \((C_i)_{>0}\) の作用は零なので、これらは \(C\)-加群としてもnoetherianである。
 よって \(M\) はnoetherianである。この観察は \(C\) 自身のnoetherian性を使わない。
 各面がnoetherianな加群の部分商についても、同じ観察を適用できる。

 \emph{切断加群.}
 \(q>0\) のとき、amplenessにより \(Q^q(\mathcal F)_{\ge\mathbf b}=0\) となる
 \(\mathbf b\) がある。従って前段階から
 \(Q^q(\mathcal F)\) は\emph{noetherian}である。
 特に、後で用いる \(Q^1\) の任意の部分加群も有限生成である。

 \(q=0\) について、\(\mathcal F\mathcal L^{\mathbf a}\) を大域生成する
 \(\mathbf a\ge0\) を選び、逆bimoduleをtensorして
 \[
 0\longrightarrow\mathcal K\longrightarrow
       (\mathcal L^{-\mathbf a})^{\oplus t}
       \longrightarrow\mathcal F\longrightarrow0
 \]
 を得る。\(\mathcal K\) はcoherentである。
 この積の場合、\(\mathbf v\) の一つでも成分が負なら
 \(H^0(X,\mathcal L^{\mathbf v})=0\) である。
 実際、外部積のそのfactorの次数が負であり、K\"unneth公式を適用できる。
 従って切断を全非負多重次数で取った列は
 \[
 C(-\mathbf a)^{\oplus t}\longrightarrow Q^0(\mathcal F)
       \longrightarrow Q^1(\mathcal K)
 \]
 と正確である。最初の写像の像は有限生成であり、余核は既にnoetherianと示した
 \(Q^1(\mathcal K)\) の部分加群なので、\(Q^0(\mathcal F)\) は有限生成である。
 また任意の \(\mathbf a\ge0\) について
 \[
 Q^0(\mathcal F)_{\ge\mathbf a}(\mathbf a)
 \cong Q^0(\mathcal F\mathcal L^{\mathbf a})
 \]
 であるから、第一象限によるtruncationも有限生成である。

 \emph{イデアルの有限生成.}
 多重斉次右イデアル \(I\subset C\) を取る。
 \cite[Theorem 3.4]{ChanMulti} により、sheafificationはexactであり、
 \(\widetilde I\subset\widetilde C=\OO_X\) となる。
 \(X\) はnoetherianなので、このquasi-coherent部分層はcoherentである。
 随伴のunit \(\eta_I:I\to Q^0(\widetilde I)\) の自然性と
 \(\eta_C:C\xrightarrow{\sim}Q^0(\OO_X)=C\) により、\(\eta_I\) は単射である。
 その余核は多重tailsに関するtorsionであり、\(Q^0(\widetilde I)\) の商だから有限生成である。
 有限生成torsion加群は一つの第一象限上で零となるため、ある \(\mathbf a\ge0\) に対して
 \[
 I_{\ge\mathbf a}=Q^0(\widetilde I)_{\ge\mathbf a}
 \]
 となる。右辺の有限生成性は前段階で証明した。
 一方、\(I/I_{\ge\mathbf a}\) は \(C/C_{\ge\mathbf a}\) の部分加群である。
 \(C\) の固定面は有限次元空間 \((C_i)_a\) と \(C_{\widehat i}\) のtensor productなので
 \(C_{\widehat i}\)-noetherianである。
 従って帯の観察により \(C/C_{\ge\mathbf a}\) はnoetherianであり、
 \(I/I_{\ge\mathbf a}\) も有限生成となる。これで \(I\) の有限生成性を得た。

 加法と両立する \(\mathbb N^r\) の整列順序、例えば全次数付き辞書式順序で
 任意の右イデアルの先頭多重斉次項を取る。
 その斉次イデアルの有限生成元を元のイデアルへ持ち上げ、
 先頭次数を下げる帰納法を行えば、任意の右イデアルも有限生成である。
 これで通常の右noetherian性と帰納法が完了する。
 opposite algebraには逆translationを伴う正次数の可逆bimoduleが対応するので、
 同じ証明により左noetherian性も得る。
\end{proof}

\begin{lemma}[正次数中心正則元からのnoetherian性の持ち上げ]
 \label{lem:central-noetherian-lift}
 \(D\) を連結局所有限な非負次数付き代数、
 \(g\in Z(D)\) を正次数の斉次正則元とする。
 \(D/(g)\) がnoetherianなら \(D\) もnoetherianである。
\end{lemma}
\begin{proof}
 \(g\)-adic associated gradedは \((D/(g))[t]\) でありnoetherianである。
 任意の斉次右イデアル \(J\) に誘導された \(g\)-adic filtrationの
 associated gradedはこの環の右イデアルなので有限生成である。
 その二重斉次生成元を \(J\) に持ち上げる。
 内部次数 \(d\) の元のinitial termをこれらの持ち上げで消すと、
 \(g\)-adic orderが真に増える。
 しかし \(g^nD\) の内部次数は少なくとも \(n\deg(g)\) なので、
 固定した \(d\) ではこの操作は有限回で終わる。
 よって全ての斉次右イデアルが有限生成である。
 任意の右イデアルについては通常の内部次数の最高次項を取り、
 生成元を持ち上げて次数に関する帰納法を行えば有限生成性を得る。
 左側も同じである。
\end{proof}

\begin{proposition}[周囲のテンソル積]
\label{prop:tensor-sklyanin-ambient}
 \(h_i^6\ne-27\)、\(S_i=S_{h_i}\)、\(g_i=G_{h_i}\) とし、
 \[
 A=(\bigotimes_{i=1}^rS_i)[u_1,\ldots,u_m],\qquad
 N=3r+m,\qquad W=\sum_i g_i+P(u),\qquad R=A/(W)
 \]
 と置く。ここで \(P\) は任意の中心斉次三次式であり、\(m=0\) なら \(P=0\) とする。
 \(A\) は左右noetherianなKoszul AS-regular代数で、
 \[
 \operatorname{gldim}A=N,\qquad H_A(t)=(1-t)^{-N},\qquad
 \operatorname{Ext}^j_A(\CC,A)=
 \begin{cases}\CC(N)&j=N,\\0&j\ne N\end{cases}
 \]
 を満たす。\(W\) は中心正則元であり、\(R\) はinjective dimension \(N-1\)、
 Gorenstein parameter \(N-3\) のAS-Gorenstein代数である。
 この主張に \(\sigma_i\) の有限位数や \(P\) の孤立性は必要ない。
\end{proposition}
\begin{proof}
 \(S_i/(g_i)\) の斉次基底の持ち上げ \(V_i\subset S_i\) を選ぶと、
 \(g_i\) の正則性と次数に関する帰納法により
 \(S_i\cong V_i\otimes\CC[g_i]\) は中心 \(\CC[g_i]\)-加群として自由である。
 \(g_1,\ldots,g_r\) は \(\bigotimes_iS_i\) 上の中心正則列をなし、その商は
 補題\ref{lem:elliptic-thcr-tensor-noeth}の \(C\) である。
 補題\ref{lem:central-noetherian-lift}により、正次数中心正則元について
 商のnoetherian性を元の連結次数付き代数へ持ち上げられる。
 これを反復し、中心変数についてのHilbert基底定理を用いると \(A\) は左右noetherianとなる。

 各 \(S_i\) の長さ三のKoszul双加群分解と、多項式環のKoszul双加群分解をtensorすると、
 \(A\) の対角双加群の長さ \(N\) の有限射影分解を得る。
 よって \(\operatorname{gldim}A\le N\) である。
 各factorの \(\CC\) のKoszul分解とそのdualのK\"unneth公式は、
 表示したExt公式と左側の同じ式を与える。
 従ってglobal dimensionはちょうど \(N\) であり、AS-Gorenstein条件も成立する。
 Koszul性とHilbert seriesはtensor分解から従い、GK dimensionは \(N\) である。

 \(A\) は \(\CC[g_1,\ldots,g_r,u_1,\ldots,u_m]\) 上自由である。
 この中心多項式環で \(W\) は \(g_1\) に関してmonic一次式なので、\(A\) 上でも左右正則である。
 中心正則元によるchange of ringsにより、商のinjective dimensionは一つ減り、
 Gorenstein parameterは三減る。
\end{proof}

\subsection{多ブロックHesse模型からの埋込み}

\(r\ge1\)、\(m\ge0\) とし、独立した変数 \(x_j,y_j,z_j\) に対して
\[
 S_{h_j}=\CC\langle x_j,y_j,z_j\rangle/
 ([x_j,y_j]-h_jz_j^2,[y_j,z_j]-h_jx_j^2,[z_j,x_j]-h_jy_j^2),
\]
\[
 G_j=x_j^3+y_j^3+z_j^3-2h_j^2\sym(x_jy_jz_j),\qquad
 F_j=x_j^3+y_j^3+z_j^3+h_j^2x_jy_jz_j
\]
と置く。異なるブロックの元は可換とし、
\[
 A_{\mathbf h,m}=(S_{h_1}\otimes\cdots\otimes S_{h_r})[u_1,\ldots,u_m],
 \qquad W_{\mathbf h,m}=\sum_{j=1}^rG_j+\sum_{a=1}^m u_a^3,
\]
\[
 F_{\mathbf h,m}=\sum_{j=1}^rF_j+\sum_{a=1}^m u_a^3,
 \qquad Y_{\mathbf h,m}=V(F_{\mathbf h,m})\subset\PP^{3r+m-1}
\]
とする。

\begin{proposition}[多ブロックの完全忠実比較]\label{prop:multiblock-embedding}
 全ての \(j\) について \(h_j^6\ne-27\) と仮定する。
 degree shiftと三次periodicityを保つdg完全忠実関手
 \[
 \MF^{\ZZ}(\CC[x_j,y_j,z_j,u_a],F_{\mathbf h,m})
 \longrightarrow
 \MF^{\ZZ}(A_{\mathbf h,m},W_{\mathbf h,m})
 \]
 が存在する。\(n=3r+m-2\ge3\) なら、左辺は
 \((\Ku(Y_{\mathbf h,m}),\rho_{Y_{\mathbf h,m}})\) である。
 この命題は、右辺の全対象が像に属するとは主張しない。
\end{proposition}
\begin{proof}
 \(s=r+m\) とし、各非可換ブロックと各一変数ブロックに別の \(e_i\) を割り当てる。
 系\ref{cor:block}は各 \(j\) について
 \[
 \MF^{\ZZ}(\CC[x_j,y_j,z_j],F_j)
 \simeq\MF^{\ZZ}(S_{h_j},G_j)
 \]
 をdegree shiftとperiodicityを保って与える。
 これらをtensorし、\(\MF^{\ZZ}(\CC[u_a],u_a^3)\) には恒等を用いる。
 非可換側へは補題\ref{lem:multiblock-ff}の完全忠実関手を用いる。

 可換側の外部積は本質的全射でもある。
 これは可換なsmooth affine ambientに対するequivariant Thom--Sebastiani定理
 \cite[Theorem 5.15, Corollary 5.18]{BFK}を適用したものである。
 二因子の場合は \(G=H=\mathbb G_m\)、potentialのcharacterを双方 \(t\mapsto t^3\)
 と取る。異なる二因子上でのcharacterの差 \((-3,3)\) は非torsionであり、
 \(G\times_{\mathbb G_m}H\) のcharacter groupは
 \(\ZZ^2/\langle3e_2-3e_1\rangle\) となる。
 各potentialは斉次なので、Euler公式により和のcritical locusは
 個々の零点の積に含まれる。従って同定理の支持条件を除くことができる。
 反復して \(\Gamma_s\) を得る。
 ここで引用している定理のambientは全て可換であり、非可換ambientへ拡張して引用してはいない。

 これで \(\Gamma_s\)-次数付きの完全忠実関手を得た。
 \(\Gamma_s\to\ZZ\)、\(e_i\mapsto1\) の核は
 \(T_s\cong(\ZZ/3)^{s-1}\) である。
 細分次数の圏から通常次数の圏へ戻す双対同変化では、
 \(T_s\) の生成元は因子のdegree shiftの差
 \(\rho_i\rho_1^{-1}\) として作用する。
 全ての比較はこれらと三次periodicityを保つので、
 有限群関係とそのcoherenceを保つ同変関手となる。
 同変対象間のHomは元のHomの有限群不変部分であり、標数零では不変部分を取る操作は完全である。
 従って完全忠実性は同変化で保たれる。
 \cite[Theorem 4.2 and Section 8]{Elagin}の双対同変化を用いて通常次数へ降ろせば、
 求める関手を得る。\((e_1)\) との両立は降下後の \((1)\) との両立になる。

 最後に \(Y_{\mathbf h,m}\) のsmoothnessを確認する。
 特異点では全ての \(u_a=0\) であり、各三変数ブロックの偏微分が同時に零となる。
 \(F_j\) の射影零点がsmoothなので、そのaffine critical locusは原点だけである。
 従って全座標が零となり、射影特異点は存在しない。
 Orlov同値により最後の主張を得る。
\end{proof}

\begin{remark}
 命題\ref{prop:tensor-sklyanin-ambient}により、この独立ブロック族の周囲は
 平行移動の位数によらずnoetherian AS-regularである。
 中心正則元の商についてのMF--特異圏比較とOrlov windowにより、
 右辺は \(A_{\mathbf h,m}/(W_{\mathbf h,m})\) の残余成分と同一視できる。
 この周囲の正則性は、残余圏そのもののpropernessや比較の本質的全射性とは別である。
\end{remark}

\subsection{Euler導分と中心支持への帰着}

\begin{lemma}[ブロックpotentialの零ホモトピー]\label{lem:block-euler}
 任意の有限自由factorization \(P\) of \(W_{\mathbf h,m}\) 上で、
 各 \(G_j:P\to P(w)\) と各 \(u_a^2:P\to P(2e_{r+a})\) は零ホモトピックである。
 この主張は \(\Gamma_s\)-次数付きでも通常の整数次数付きでも成り立つ。
\end{lemma}
\begin{proof}
 \(j\) 番目のブロックの次数を数え、他のブロックと \(u_a\) では零となるEuler導分を
 \(\delta_j\) と書く。Sklyanin関係式はそのブロックについて斉次なので、
 \(\delta_j\) は自由代数から商代数へ降りる。
 \(\delta_j(W_{\mathbf h,m})=3G_j\) である。
 factorizationの有限自由加群に斉次基底を選び、その微分行列 \(d\) の係数に
 \(\delta_j\) を作用させる。積のLeibniz則から
 \[
 d\,\delta_j(d)+\delta_j(d)\,d
 =\delta_j(d^2)
 =3G_j\,\Id_P.
 \]
 従って \(\frac13\delta_j(d)\) が所要のホモトピーである。
 \(\delta_j\) は代数の元の \(\Gamma_s\)-次数を変えない。
 \(3e_j=w\) という次数群の関係があっても、
 元のtensor factorsから定義した導分の存在は失われない。
 同じ計算を中心変数の導分 \(\partial_{u_a}\) に対して行えば
 \[
 d\,\partial_{u_a}(d)+\partial_{u_a}(d)\,d=3u_a^2\Id_P
 \]
 を得る。基底に依存するのはホモトピーの選択であり、零ホモトピックという結論ではない。
\end{proof}

\begin{proposition}[生成性の中心商への帰着]\label{prop:block-support}
 全ての \(h_j^6\ne-27\) とし、\(A=A_{\mathbf h,m}\) と置く。
 命題\ref{prop:tensor-sklyanin-ambient}により、この \(A\) はnoetherian AS-regularである。さらに
 \[
 R=A/(W_{\mathbf h,m}),\qquad
 C_j=S_{h_j}/(G_j),\qquad C=\bigotimes_{j=1}^r C_j
 \]
 と置く。\(C\) は \(R\) の商とみなす。
 次の等式が、冪等完備化した \(\Gamma_s\)-次数付き特異圏で成り立つ。
 \[
 \Dsg^{\Gamma_s}(R)
 =\thick\!\left(
 \operatorname{Res}\bigl(D^b(\operatorname{gr}^{\Gamma_s}C)\bigr)
 \right).
 \tag{\(\ddagger\)}\label{eq:block-support}
 \]
 また、\(C\) は \(R\)-perfectなので、restrictionは特異圏の間の関手にも降りる。
\end{proposition}
\begin{proof}
 各 \(G_j\) は中心正則元である。\(S_{h_j}/G_jS_{h_j}\) の斉次基底を持ち上げた
 ベクトル部分空間を \(V_j\subset S_{h_j}\) とすると、次数に関する帰納法と正則性から
 \[
 V_j\otimes_\CC\CC[t_j]\xrightarrow{\sim}S_{h_j},\qquad
 v\otimes t_j^b\longmapsto vG_j^b
 \]
 は中心 \(\CC[t_j]\)-加群の同型である。
 従って \(A\) は中心多項式環
 \(\CC[t_1,\ldots,t_r,u_1,\ldots,u_m]\) 上自由である。
 \(W\) はこの環の \(\sum_jt_j+\sum_au_a^3\) に対応するので正則であり、
 \(R\) は \(\CC[t_1,\ldots,t_{r-1},u_1,\ldots,u_m]\) 上自由となる。
 よって
 \[
 u_1,\ldots,u_m,G_1,\ldots,G_{r-1}
 \]
 は \(R\) の中心正則列で、その商は \(C\) である。
 Koszul分解は \(C\) の有限自由 \(R\)-分解を与える。
 これでrestrictionがperfect対象をperfect対象へ送ることも従う。

 まず冪等完備化前の \(\Dsg^{\Gamma_s}(R)\) の対象 \(N\) を取る。
 中心正則元によるAS-regular代数の商に対するMF比較
 \cite[Theorem 5.8]{CCKM}を使えば、\(N\) はfinite-free factorizationのcokernelのshiftで表せる。
 この比較は通常の整数次数で記述されているが、
 \(\Gamma_s\simeq\ZZ\oplus\Gamma_{s,\mathrm{tors}}\) と標数零を用い、
 各有限自由分解とホモトピーを斉次成分ごとに取る同じ構成で
 \(\Gamma_s\)-次数を保った比較を得る。
 補題\ref{lem:block-euler}より、\(N\) は
 \[
 K(G_1,\ldots,G_{r-1},u_1^2,\ldots,u_m^2;N)
 \]
 のshiftと内部twistの直和因子となる。
 ここで \(K\) は通常の中心Koszul複体である。
 この複体の各cohomologyは \(G_1,\ldots,G_{r-1},u_1^2,\ldots,u_m^2\) に殺される。
 さらに \((u_1,\ldots,u_m)^{m+1}\subset(u_1^2,\ldots,u_m^2)\) なので、
 各cohomologyの有限filtrationの商は全 \(u_a\) と全 \(G_j\) に殺される \(C\)-加群である。
 \(G_r\) の消滅は \(W=0\) から従う。
 noetherian性によりこれらは有限生成であり、有限個のtruncation triangleとfiltrationから
 未完備の特異圏の各対象が右辺に入る。
 最後に直和因子閉包を取れば、\eqref{eq:block-support} を得る。
\end{proof}

\subsection{必要な生成性の正確な内容}

命題\ref{prop:block-support}は、全ての \(C\)-加群のrestrictionが
\(C_1,\ldots,C_r\) の外部積から生成されることまでは述べていない。
十分条件の一つは
\[
 D^b(\operatorname{gr}^{\Gamma_s}C)
 =\thick\left\{
 \operatorname{coarsen}(M_1\boxtimes\cdots\boxtimes M_r)(\gamma)
 :M_j\in D^b(\operatorname{gr}^{\ZZ}C_j),\ \gamma\in\Gamma_s
 \right\}
 \tag{G}\label{eq:block-generation}
\]
である。\(\gamma\) の全shiftは、中心変数が零となった後にも残る次数類を含めるために必要である。
以下で用いるsyzygyへの置換を正当化するため、
\(A_{\ne i}=(\bigotimes_{j\ne i}S_{h_j})[u_1,\ldots,u_m]\) とし、
他因子の外部積に全ての \(u_a\) が零で作用する \(A_{\ne i}\)-加群を \(M_{\ne i}\) とする。
\(S_{h_i}\) の \(\CC[G_i]\) 上の自由性と \(W=0\) による \(G_i\) の消去から、
\(R\) は \(A_{\ne i}\) 上自由であり、
\[
 \operatorname{Res}(C_i\boxtimes M_{\ne i})
 \simeq M_{\ne i}\otL_{A_{\ne i}}R.
\]
各 \(S_{h_j}\) 上の有限自由分解と \(u_a\) のKoszul分解を外部tensorすると、
\(M_{\ne i}\) は \(A_{\ne i}\)-perfectであり、上のrestrictionは \(R\)-perfectである。
従って一因子が \(C_i\)-perfectである項はrestriction後に零となり、
各因子を有限syzygyへ置き換えてよい。有界複体の場合もtruncation triangleから従う。
各因子の有限syzygyを取りMFで表した後、外部tensor productの有限ambient分解を
\(W\) の零ホモトピーで周期化すると、そのrestrictionの安定化は外部積factorizationになる
（全体のcohomological shiftは生成性に影響しない）。
したがって (G) が成立すれば、命題\ref{prop:block-support}と
補題\ref{lem:multiblock-ff}により非可換側の外部積もMorita同値となり、
命題\ref{prop:multiblock-embedding}を圏同値へ強められる。
より弱く、(G)の等式をrestriction後の \(\Dsg^{\Gamma_s}(R)\) だけで示しても十分である。

しかし (G) を、multigraded tailsの可換な積の幾何だけから推論することはできない。
\(C\) は元来 \(\ZZ^r\)-次数を持つが、ここで必要なのは
\(\Gamma_s\)-次数付き加群である。\(\ZZ^r\to\Gamma_s\) の核は
\(3(e_i-e_1)\) からなる無限格子を含み、有限群降下ではない。
\(\Gamma_s\)-次数付き加群は一般には \(\ZZ^r\)-次数へ持ち上がらない。
最も単純には、\(\CC[a,b]\) に \(\deg a=e_1\)、\(\deg b=e_2\)、
\(3e_1=3e_2\) を入れると、\(\CC[a,b]/(a^3-b^3)\) は \(\Gamma_2\)-次数付きだが
\(\ZZ^2\)-次数へ持ち上がらない。そのannihilatorが \(\ZZ^2\)-斉次でないためである。
この例は (G) 自体を反証するものではなく、個々の対象を細分次数へ持ち上げる証明法を反証する。

さらにmultigraded tailsでは各座標軸上のtorsionが消える。
生成性を完成させるには、これらのtorsionと上記の非liftable加群を、
restriction後も含めて制御する必要がある。
第\ref{sec:finite-center}節では、高々一つを除く因子の有限中心性でこの両方を扱い、
一つの任意位数ブロックを許して本質的全射性を証明する。
二つ以上の無限位数因子の一般の場合は未解決として保持する。

\begin{remark}[別の十分条件]
 \(\mathcal A\to\cD\) をsmooth proper圏からのdg完全忠実関手とする。
 \(\cD\) がproperでSerre関手を持ち、埋込みがSerre関手と両立すれば、
 \(\mathcal A\) のsaturated性から左右随伴が存在し、その直交成分は左右一致する。
 したがって \(\cD\) に非自明な直交直和分解がなければ埋込みは同値である。
 これは上の生成性とは別の完成条件となる。
 ただし複数ブロックの全MFについて、必要なpropernessと直交分解不能性を
 ここで証明したわけではない。
\end{remark}

\section{一方有限中心の外部積生成性と多ブロックの可換モデル}\label{sec:finite-center}

前節で分離した生成性の問題について、一方の係数代数に
有限中心性を課さない外部積生成定理を証明する。
これを有限位数の楕円ブロックと組み合わせることで、
高々一つの無限位数ブロックを許す具体模型を得る。

この節では生成条件 \eqref{eq:block-generation} を、高々一つを除くブロックが
有限位数である場合に証明する。任意位数の一因子は係数代数として残し、
有限中心の因子を一つずつ追加する。
非可換係数付きBeilinson分解、中心のNoether正規化、分離的な稠密開集合と
その補集合の帰納法を用いる。
\(\ZZ^r\)-次数へ持ち上がらない対象と軸上のtorsionも初めから含む。
有限中心の場合の基礎的な局所化・生成法については \cite{ELS} と比較されたい。

\subsection{斉次中心局所化と稠密開集合}

\(\operatorname{gr}^{\Lambda}D\) は有限生成の \(\Lambda\)-次数付き右 \(D\)-加群の圏とする。
ここで \(\Lambda\) は有限生成アーベル群であり、torsionを持っていてよい。
全てのthick閉包は、有限直和、shift、cone、直和因子に関する閉包を表す。

\begin{lemma}[斉次中心局所化]\label{lem:graded-localization}
 \(D\) をnoetherianな \(\Lambda\)-次数付き代数、
 \(f\in Z(D)\) を斉次元とする。
 有界導来圏の局所化は
 \[
 D^b(\operatorname{gr}^{\Lambda}D)\big/
 D^b_{f\text{-tor}}(\operatorname{gr}^{\Lambda}D)
 \simeq D^b(\operatorname{gr}^{\Lambda}D_f)
 \]
 を与える。左辺の核は、各cohomologyが \(f\) の冪に殺される対象からなる。
 この核は \(D/(f)\) 上の有限生成 \(\Lambda\)-次数付き加群のrestrictionで生成される。
\end{lemma}
\begin{proof}
 局所化は中心元についての平坦な局所化であるからcohomologyと可換である。
 有限生成加群が局所化後に零なら、その有限個の生成元を同時に殺す \(f\) の冪がある。
 \(f\)-adic filtrationと有界truncation triangleにより核についての主張を得る。

 局所化後の有限生成次数付き加群は、有限個の斉次生成元の
 \(D\)-spanを取ることによって、局所化前の有限生成次数付き加群から得られる。
 射については有限生成加群の有限生成自由分解の必要な有限部分を使う。
 中心局所化と \(\operatorname{Ext}^q\) は可換であり、
 次数零の射は \(f\) の冪を分母としてclearできる。
 同じことは有界複体にも有限個のtruncationによって成立する。
 次数を保つ分母の射は \(f^N:M(-N\deg f)\to M\) と書く。
 そのconeは \(f\)-torsionである。
 従って上記Verdier商は局所化後の圏へ完全忠実に写り、
 加群とその有限個のconeを全て含むので本質的全射である。
 この議論で次数を細分する必要はない。
\end{proof}

\begin{lemma}[有限中心代数の良い稠密開集合]\label{lem:generic-separable}
 \(k\) を標数零の体とする。
 \(D\) を連結非負次数付き \(k\)-代数で、有限生成の可換次数付き中心部分代数
 \(Z\subset Z(D)\) 上有限なものとする。
 \(Z\) がreducedであると仮定する。
 ある斉次元 \(f\in Z\) と斉次両側イデアル \(I\subset D\) が存在し、次を満たす。
 \begin{enumerate}
 \item \(f\) は \(Z\) の全てのminimal primeを避ける。
       特に \(Z/(f)\ne0\) なら \(\dim Z/(f)<\dim Z\) である。
 \item \(Z_f\) は \(k\) 上smoothな有限型可換環である。
 \item \(I_f\) は冪零であり、\(Q_f=(D/I)_f\) は
       \(Z_f\) 上有限射影かつ分離的な代数である。
 \end{enumerate}
 \(\dim Z=0\) の場合には \(Z=k\)、\(f=1\) と取れる。
\end{lemma}
\begin{proof}
 \(K\) を \(Z\) の全商環とする。
 \(Z\) はreduced noetherianなので、\(K\) はminimal primeの分数体の有限積である。
 \(D_K\) は \(K\) 上有限次元であり、そのJacobson radical \(J\) は冪零である。
 \(I\) を \(D\to D_K\) による \(J\) の逆像とする。
 次数を与える \(\mathbb G_m\)-作用はminimal primeの集合と \(J\) を保つので、
 \(I\) は斉次である。
 \(I^N\) は十分大きい \(N\) で \(K\) 上零となり、
 有限生成性から、ある全商環の分母を反転すれば \(I^N=0\) となる。

 \(Q_K=D_K/J\) は半単純である。
 標数零なので、各simple factorの中心は対応する成分体の有限分離拡大であり、
 \(Q_K\) は \(K\) 上分離的である。
 generic freenessを用いて、ある稠密開集合上で \(Q\) を有限射影とする。
 分離性は、積の写像 \(Q\otimes_ZQ^{\mathrm{op}}\to Q\) の
 \(Q\)-双加群としてのsection、すなわち分離冪等元の存在で表せる。
 generic fiberにある分離冪等元と、それが満たす有限個の等式の分母を払えば、
 さらに小さい稠密開集合上でこの分離冪等元が存在する。
 また \(Z\) のsmooth locusは全てのgeneric pointを含む。

 以上の良い開集合は \(\mathbb G_m\)-作用で不変に取れる。
 実際、\(I^N\) の消滅、局所自由性、分離性、smoothnessは全て
 次数作用によって保たれる開条件である。
 補集合のイデアルは斉次で、どのminimal primeにも含まれない。
 \(Z_0=k\) と \(k\) の無限性により、そのイデアル内から
 全てのminimal primeを避ける一つの斉次元 \(f\) を選べる。
 必要なら候補を冪にして次数を揃え、一般の線型結合を取ればよい。
 \(\dim Z>0\) のとき零次元の既約成分はない
 （全ての既約成分は連結な錐の頂点を含む）ので、この操作に例外はない。
 \(Z=k\) の場合は有限次元代数のradicalと半単純商を用いればよい。
\end{proof}

\begin{lemma}[分離代数と射影分解]\label{lem:separable-perfect}
 \(Z\) をsmooth affineな有限型 \(k\)-代数、
 \(Q\) を \(Z\) 上有限射影かつ分離的な代数とする。
 このとき \(Q\) は有限global dimensionを持つ。
 特に \(Q\) に任意のアーベル群次数を与えたとき、
 \(D^b(\operatorname{gr}^{\Lambda}Q)\) は \(Q(\lambda)\)
 \((\lambda\in\Lambda)\) で生成される。
\end{lemma}
\begin{proof}
 任意の右 \(Q\)-加群 \(M\) は、分離冪等元により
 \(M\otimes_ZQ\to M\) の \(Q\)-線型sectionを持つ。
 \(Z\) は有限global dimensionを持ち、\(Q\) は \(Z\)-flatなので、
 \(M\) の \(Z\)-射影分解を \(-\otimes_ZQ\) で送ると
 \(M\otimes_ZQ\) の有限 \(Q\)-射影分解を得る。
 従ってその直和因子 \(M\) も有限射影次元を持つ。

 次数付き有限生成加群に有限生成斉次自由分解を作ると、十分高いsyzygyは
 ungradedには射影的である。
 自由加群からこのsyzygyへの次数零の全射のungraded sectionを取り、
 sectionの次数零成分を抜き出すと次数付きsectionになる。
 よってsyzygyは次数付き有限自由加群の直和因子であり、
 元の加群は有限個の \(Q(\lambda)\) から生成される。
\end{proof}

\subsection{非可換係数の多項式拡大と一方有限中心の生成定理}

ここでは係数代数の有限中心性を仮定しない。
これが一つの無限位数ブロックを許すための要点である。
以下、\(k\) は標数零の体、\(\Lambda\) は有限生成アーベル群とする。

\begin{lemma}[非可換係数を持つ多項式拡大]\label{lem:polynomial-coefficient}
 \(A\) を左右noetherianな \(\Lambda\)-次数付き \(k\)-代数とする。
 \(P=k[t_1,\ldots,t_d]\) の各変数に任意の \(\Lambda\)-次数を与え、
 \(A\otimes_kP\) には和の次数を与える。このとき
 \[
 D^b(\operatorname{gr}^{\Lambda}(A\otimes_kP))
 =\thick\{M\otimes_kP:\ M\in\operatorname{gr}^{\Lambda}A\}.
 \]
 \(P\) の斉次元 \(f\) を反転した \(P_f\) についても同じ結論が成立する。
 \(A\) が有限global dimensionを持つという仮定はない。
\end{lemma}
\begin{proof}
 中心多項式拡大のHilbert基底定理により \(A\otimes P\) は左右noetherianである。
 \(d=0\) は自明なので \(d>0\) とする。
 有限生成 \(\Lambda\)-次数付き \(A\otimes P\)-加群 \(T\) の有限斉次表示を取る。
 射影用の独立な整数次数を導入し、\(t_0,t_1,\ldots,t_d\) の射影次数を全て1、
 \(A\) の射影次数を0とする。\(\Lambda\)-次数は
 \(\deg_\Lambda(t_0)=0\)、\(\deg_\Lambda(t_i)=\deg_\Lambda(t_i\in P)\) とする。
 表示の各関係を \(t_0\) で斉次化すると、\(\PP^d_k\) 上の
 \(\mathcal A=A\otimes_k\OO_{\PP^d}\)-加群 \(\cF\) が得られる。
 これは \(\mathcal A\) 上有限表示で、\(t_0\ne0\) への制限が \(T\) である。
 各関係の斉次化では \(\Lambda\)-次数を変えていない。
 以後の \(\Lambda\)-次数は対角化可能群 \(D(\Lambda)\) による線型化としても解釈できる。

 必要なcohomologyの有限性を確認する。
 \(\mathcal A\) 上のcoherent加群には有限個の
 \(\mathcal A(-a)(\lambda)\) からの全射があり、核もcoherentである。
 これは中心変数による有限斉次表示とnoetherian性から従う。
 標準affine被覆の\v Cech複体は長さ \(d\) なので、cohomologyは次数 \(0,\ldots,d\)
 に限られる。一方、\(H^q(\PP^d,\mathcal A(a))
 =A\otimes_kH^q(\PP^d,\OO(a))\) は有限生成 \(A\)-加群である。
 上記全射とその核の長完全列を用い、\(q=d,d-1,\ldots,0\) の順に帰納すれば、
 全てのcoherent \(\mathcal A\)-加群の \(H^q\) が有限生成 \(A\)-加群と分かる。
 この議論は線型化と両立するので、\(\cF\) に有限階数線束をテンソルしたものの
 導来大域切断は \(D^b(\operatorname{gr}^{\Lambda}A)\) に属する。
 ここでcoherentは \(\mathcal A\) 上の意味であり、\(\OO\) 上有限とは仮定しない。

 Beilinsonの対角分解 \cite{Beilinson} を、因子の順序を適切に選んで
 \[
 0\longrightarrow\Omega^d(d)\boxtimes\OO(-d)
 \longrightarrow\cdots\longrightarrow
 \Omega^1(1)\boxtimes\OO(-1)
 \longrightarrow\OO\boxtimes\OO
 \longrightarrow\OO_\Delta\longrightarrow0
 \]
 と書く。これは \(D(\Lambda)\)-線型化を保つ分解である。
 \(\mathcal A\) の係数を保ったまま \(\cF\) に作用させると、\(\cF\) は
 \[
 R\Gamma(\PP^d,\cF\otimes_{\OO}\Omega^i(i))\otimes_k\OO(-i),
 \qquad 0\le i\le d,
 \]
 の有限個のconeから作られる。
 対角のKoszul分解は各射影に関してflatであるから、\(\cF\) の
 \(\OO\)-flatnessは必要ない。非可換な \(A\)-作用も全ての操作と可換である。
 係数複体は前段落により \(D^b(\operatorname{gr}^{\Lambda}A)\) に属する。
 \(t_0\ne0\) に制限すると、\(t_0\) の \(\Lambda\)-次数が0であるため
 \(\OO(-i)\) は次数を保って自明化される。
 よって \(T\) は \(M\otimes P\) から生成される。
 係数複体を有限個のcohomology加群に分けるには通常のtruncationを用いればよい。
 有界複体全体についても同じtruncationから結論する。

 \(P_f\) の場合は補題\ref{lem:graded-localization}による局所化を適用する。
 有限生成次数付き加群の有限表示は局所化前に持ち上がるので、
 局所化後の生成対象はちょうど \(M\otimes P_f\) である。
\end{proof}

\begin{lemma}[多項式中心上の分離的な稠密開集合]\label{lem:polynomial-separable}
 \(D\) を連結局所有限な非負次数付き \(k\)-代数で、有限型可換斉次中心部分代数
 \(Z\) 上有限なものとする。\(d=\dim Z\) と置く。
 正次数の斉次変数を持つ多項式部分代数 \(P=k[t_1,\ldots,t_d]\subset Z\)、
 非零斉次元 \(f\in P\)、斉次両側イデアル \(J\subset D\) を選んで、
 \(D\) が \(P\) 上有限、\(J_f\) が冪零、\(Q=(D/J)_f\) が
 \(P_f\) 上有限射影かつ分離的となるようにできる。
 分離冪等元は元の整数次数で次数0に取れる。
\end{lemma}
\begin{proof}
 斉次Noether正規化によって \(Z\) は正次数変数の多項式部分代数 \(P\) 上有限となる。
 \(Z\) のreducednessは必要ない。\(K=\operatorname{Frac}P\) とすると、
 \(D_K\) は有限次元 \(K\)-代数である。
 そのradicalの \(D\) への逆像を \(J\) とする。
 次数のスカラー作用は \(P\)、\(K\) とradicalを保つので \(J\) は斉次である。
 十分大きい \(a\) で \((J^a)_K=0\) だから、有限生成性によって
 ある非零元を反転すると \(J^a=0\) となる。
 半単純商 \(D_K/\operatorname{rad}(D_K)\) は標数零の体上分離的である。
 generic freenessと分離冪等元の有限個の等式の分母払いにより、
 ある稠密開集合上で \(D/J\) は有限射影かつ分離的となる。

 これらの条件の成立する開集合は次数作用に関して不変である。
 その補集合のイデアルは斉次であり、\(P\) は整域なので、そこから
 非零斉次元 \(f\) を一つ選べる。
 \(d=0\) の場合は \(P=k\)、\(f=1\) とし、有限次元代数のradicalを使う。
 最後に、分離冪等元を元の整数次数で分解し、次数0の成分を取る。
 \(\mu(e)=1\) と \((q\otimes1)e=e(1\otimes q)\) は斉次 \(q\) ごとに次数に分解されるから、
 次数0成分も同じ条件を満たす。これが必要な次数付き分離冪等元である。
\end{proof}

\begin{theorem}[一方だけが有限中心である外部積生成]\label{thm:one-sided-boxes}
 \(A\) を左右noetherianな \(\Lambda\)-次数付き \(k\)-代数とする。
 \(D\) を連結局所有限な非負次数付き \(k\)-代数で、
 有限型可換斉次中心部分代数 \(Z\) 上有限なものとする。
 任意の準同型 \(\beta:\ZZ\to\Lambda\) により \(D\) の次数を粗化し、
 \(A\otimes_kD\) には和の次数を与える。このとき
 \begin{equation}\label{eq:one-sided-boxes}
 \begin{split}
 D^b(\operatorname{gr}^{\Lambda}(A\otimes_kD))
 =\thick\{\,M\boxtimes\operatorname{coarsen}_\beta N:\qquad\qquad\qquad\quad\\
 M\in\operatorname{gr}^{\Lambda}A,\quad
 N\in\operatorname{gr}^{\ZZ}D\,\}.
 \end{split}
 \end{equation}
 特に \(A\) の有限中心性、有限global dimension、あるいは
 左辺の対象の細分次数へのliftは仮定しない。
\end{theorem}
\begin{proof}
 補題\ref{lem:polynomial-separable}の \(P\subset Z\) に対して
 \(A\otimes D\) は \(A\otimes P\) 上左右有限である。
 後者はnoetherianなので前者もnoetherianである。
 右辺のthick部分圏を \(\mathcal B\) とし、\(d=\dim Z\) に関して帰納する。
 \(d=0\) では \(D\) は有限次元であり、連結性から \(D_{>0}\) は冪零で商は \(k\) である。
 \(D_{>0}\)-adic filtrationの各層は有限生成 \(\Lambda\)-次数付き \(A\)-加群であるから、
 \(M\boxtimes k\) からの生成を得る。

 \(d>0\) とし、同補題の \(f,J,Q\) を取る。
 \(f=1\) と取れる場合も以下の開集合の議論はそのまま適用でき、閉集合側は空になる。
 \(J_f\)-adic filtrationにより、\(A\otimes D_f\) 上の全加群は
 \(A\otimes Q\) 上の有限生成加群から生成される。
 そのような加群 \(T\) は \(A\otimes P_f\) 上も有限生成である。
 次数0の分離冪等元 \(e=\sum_\nu x_\nu\otimes y_\nu\) により、乗法写像
 \[
 \operatorname{Res}T\otimes_{P_f}Q\longrightarrow T
 \]
 は \(A\otimes Q\)-線型かつ \(\Lambda\)-次数0のsection
 \(t\mapsto\sum_\nu tx_\nu\otimes y_\nu\) を持つ。
 \(A\) の作用は \(Q\) と可換なので、このsectionは \(A\)-線型でもある。

 補題\ref{lem:polynomial-coefficient}により \(\operatorname{Res}T\) は
 \(M\otimes P_f\) から生成される。
 \(Q\) は \(P_f\) 上射影なので \(-\otimes_{P_f}Q\) はexactであり、
 上のsectionと合わせて \(T\) は \(M\otimes Q\) から生成される。
 \(M\otimes Q\) は \(M\boxtimes(D/J)\) の局所化であるから、
 開集合上の全圏は \(\mathcal B\) の像で生成される。
 ここで \(A\otimes Q\) が有限global dimensionを持つとは主張していない。

 局所化の核は補題\ref{lem:graded-localization}により、
 \(A\otimes(D/(f))\) 上の有限生成次数付き加群から生成される。
 \(D/(f)\) は \(P/(f)\) の像という有限型可換斉次中心上有限である。
 \(f\) が非単元ならこの中心の次元は \(d-1\) 以下であるから、帰納法を適用できる。
 得られる外部積の第二因子は \(D/(f)\)-加群であり、元の \(D\)-加群としても
 元の整数次数を保っている。従って局所化の核は \(\mathcal B\) に含まれる。
 開集合と核の生成をVerdier局所化によって合わせれば \eqref{eq:one-sided-boxes} を得る。
 全てのconeと直和因子は粗化後の圏で取っており、対象の次数のliftは用いない。
\end{proof}

\subsection{任意の粗化次数と有限中心因子の反復}

\begin{theorem}[有限中心代数の外部積生成]\label{thm:finite-center-boxes}
 \(D_1,\ldots,D_r\) を連結局所有限な非負次数付き \(k\)-代数で、
 各 \(D_i\) は有限型可換斉次中心部分代数上有限とする。
 任意の準同型 \(\alpha:\ZZ^r\to\Lambda\) に対し、
 \(D=\bigotimes_iD_i\) に粗化次数を与えると
 \[
 D^b(\operatorname{gr}^{\Lambda}D)
 =\thick\{\operatorname{coarsen}_{\alpha}
 (M_1\boxtimes\cdots\boxtimes M_r)(\lambda):
 M_i\in\operatorname{gr}^{\ZZ}D_i,\ \lambda\in\Lambda\}.
 \tag{\(G_{\mathrm{fin}}\)}\label{eq:finite-center-generation}
 \]
 \(\alpha\) は単射・全射である必要はなく、無限の核を持っていてもよい。
\end{theorem}
\begin{proof}
 \(A=k\) から出発し、\(\beta_i(1)=\alpha(e_i)\) として
 定理\ref{thm:one-sided-boxes}を順に適用する。
 各段階の係数代数は有限型可換中心上有限なので左右noetherianである。
 \(k\) の有限生成 \(\Lambda\)-次数付き加群は有限個の \(k(\lambda)\) の和であり、
 反復によって表示の外部積を得る。
 複体の外部積は体 \(k\) 上exactなので、前段階のthick生成を次段階へ送れる。
\end{proof}

\begin{remark}[nonliftable対象と軸torsion]
 外部積のconeには粗化後の全ての次数零の射を用いる。
 例えば \(k[a,b]/(a^3-b^3)\) は \(3e_1=3e_2\) の次数では
 \(a^3-b^3\) による二つのfree shiftの間のconeで生成される。
 同じ加群が \(\ZZ^2\)-次数へ持ち上がる必要はない。
 また \(f\)-adic filtrationは正次数の中心元に殺される軸上の加群も保持する。
 multigraded tailsへの移行によってこれらを捨てる操作は用いない。
\end{remark}

\subsection{高々一つの無限位数ブロックへの帰結}

次の有限中心性は \cite[Corollary 2.3]{SmithTate} に含まれる既知の結果である。
必要な斉次中心を明示するため、証明を記す。

\begin{lemma}[有限位数THCRの有限中心性]\label{lem:finite-order-thcr}
 \(E\) をsmooth elliptic curve、\(L\) を正次数の線束とし、
 \(\sigma\in\operatorname{Aut}(E)\) が有限位数 \(q\) を持つとする。
 このとき \(B=B(E,L,\sigma)\) は有限型可換斉次中心部分代数上有限である。
\end{lemma}
\begin{proof}
 \(L=\mathcal O_E(D)\) となるdivisor \(D\) を選び、
 \(D_d=D+\sigma^*D+\cdots+(\sigma^{d-1})^*D\) と置く。
 通常のskew Laurent表示により、\(B\) を
 \[
 B_d=H^0(E,\mathcal O_E(D_d))t^d
 \subset k(E)[t,t^{-1};\sigma^*]
 \]
 と書ける。\(\sigma^q=1\) なので \(D_q\) は \(\sigma\)-不変であり、
 \(q\)-Veroneseは可換なsection ring
 \[
 R=B^{(q)}=\bigoplus_{n\ge0}H^0(E,\mathcal O_E(nD_q))t^{qn}
 \]
 になる。\(R\) 上では \(\sigma\) を係数に作用させ、\(t^q\) を固定する。
 \(Z=R^{\langle\sigma\rangle}\) は \(B\) の中心に含まれる。
 実際、\(a\in H^0(E,\mathcal O_E(nD_q))\) が \(\sigma\)-不変なら、
 任意の \(bt^d\in B_d\) に対して
 \[
 (at^{qn})(bt^d)=ab\,t^{qn+d}
 =b\sigma^{d*}(a)t^{qn+d}=(bt^d)(at^{qn})
 \]
 である。

 \(D_q\) は正次数なのでampleである。従って \(R\) は有限型可換代数であり、
 各 \(0\le d<q\) について
 \[
 \bigoplus_{n\ge0}H^0(E,\mathcal O_E(nD_q+D_d))t^{qn+d}
 \]
 は \(R\) 上有限なsection moduleである。
 これらの有限直和が \(B\) なので \(B\) は \(R\) 上有限である。
 有限群の不変式環について \(R\) は有限型環 \(Z\) 上有限であり、
 従って \(B\) も \(Z\) 上有限である。
\end{proof}




\begin{corollary}[一つの任意位数因子を含む生成条件]\label{cor:PI-generation}
 全ての \(j\) で \(h_j^6\ne-27\) とし、並べ替えた後の
 \(\sigma_j=t_{p_j}\ (j\ge2)\) のみ有限位数と仮定する。
 このとき \(C_j=S_{h_j}/(G_j)\)、\(C=\bigotimes_jC_j\) に対して、
 条件 \eqref{eq:block-generation} は成立する。
 中心変数を零とした後に残る \(\Gamma_s\) の全shiftも含む。
\end{corollary}
\begin{proof}
 \(C_1\) はnoetherianであり、その \(\Gamma_s\)-次数は
 単射 \(\ZZ e_1\hookrightarrow\Gamma_s\) によって与えられる。
 任意の有限生成 \(\Gamma_s\)-次数付き \(C_1\)-加群は
 \(\Gamma_s/\ZZ e_1\) の剰余類ごとの部分加群に分解する。
 有限個の生成元の次数が関わる剰余類は有限個だけであるから、
 各剰余類の代表によるshiftを戻せば、元の整数次数を持つ \(C_1\)-加群の有限和となる。
 ここで細分しているのは一因子上の加群だけであり、積の一般の加群ではない。

 \(j\ge2\) では補題\ref{lem:finite-order-thcr}によって \(C_j\) は有限中心上有限である。
 定理\ref{thm:one-sided-boxes}を
 \(A=C_1\)、\(D=C_2\)、\(\beta(1)=e_2\) に適用し、次に係数代数を
 \(C_1\otimes C_2\) として反復する。
 各段階で左右noetherian性は同定理の有限中心多項式拡大によって保証される。
 外部積は体上exactなので、前段階で得たthick生成をそのまま次段階へ送れる。
 これで全ての \(C_j\) について条件 \eqref{eq:block-generation} が得られる。
\end{proof}

以下では定理\ref{thm:multiblock-model}の仮定の下で
\[
 \begin{gathered}
 N=3r+m,\qquad A=\left(\bigotimes_{j=1}^rS_{h_j}\right)[u_1,\ldots,u_m],\\
 W=\sum_jG_{h_j}+P(u),\qquad R=A/(W)=B_{\mathbf h,P}.
 \end{gathered}
\]
と置く。\(m=0\) の場合は \(P=0\) とし、中心変数の因子を省く。

\begin{proof}[定理\ref{thm:multiblock-model}の証明]

 \emph{周囲の代数と中心元の正則性.}
 命題\ref{prop:tensor-sklyanin-ambient}により、\(A\) は左右noetherianな
 Koszul AS-regular代数で、global dimensionとGorenstein parameterは \(N\) である。
 \(W\) は中心正則元である。以下の支持帰着のために、中心上の自由性も具体的に用いる。

 \(C_j\) の斉次基底の持ち上げを \(V_j\subset S_{h_j}\) とすると、
 正則性と次数に関する帰納法により
 \[
 V_j\otimes\CC[t_j]\xrightarrow{\sim}S_{h_j},\qquad
 v\otimes t_j^a\longmapsto vG_j^a
 \]
 は中心 \(\CC[t_j]\)-加群としての同型である。
 従って \(A\) は中心多項式環
 \(Z_0=\CC[t_1,\ldots,t_r,u_1,\ldots,u_m]\) 上自由である。
 \(W\) は \(\sum_jt_j+P(u)\in Z_0\) の作用なので、左右正則である。
 その商 \(R\) は
 \(\CC[t_1,\ldots,t_{r-1},u_1,\ldots,u_m]\) 上自由であり、
 \[
 u_1,\ldots,u_m,G_1,\ldots,G_{r-1}
 \]
 は \(R\) の中心正則列で、商は \(C\) である。

 \emph{全生成性.}
 \(m>0\) では中心多項式環全体 \((T,P)=(\CC[u],P)\) を一つの因子とし、
 \[
 \Gamma=\ZZ^{r+1}/\langle3e_j-3e_0:1\le j\le r\rangle
 \]
 を用いる。\(\deg(u_a)=e_0\)、Sklyanin因子の次数は \(e_j\) である。
 \(m=0\) では \(e_0\) を省き、\(\Gamma=\Gamma_r\) とする。
 finite-free factorizationの微分を \(d\)、第 \(j\) ブロックのEuler導分を
 \(\delta_j\) とすれば、補題\ref{lem:block-euler}と同じ計算により
 \[
 d\delta_j(d)+\delta_j(d)d=3G_j\operatorname{id}.
 \]
 中心変数についても行列係数を微分することで
 \[
 d\,\partial_{u_a}(d)+\partial_{u_a}(d)\,d
 =\partial_{u_a}P\operatorname{id}
 \]
 を得る。これらは斉次のnull-homotopyである。

 \(J_P=(\partial_{u_1}P,\ldots,\partial_{u_m}P)\) とする。
 孤立性より、ある \(b\) について \((u_1,\ldots,u_m)^b\subset J_P\) である。
 命題\ref{prop:block-support}の証明で \(u_a^2\) を \(\partial_{u_a}P\) に置き換えると、
 各対象は \(G_1,\ldots,G_{r-1},\partial_{u_1}P,\ldots,\partial_{u_m}P\)
 に関するKoszul複体のshift・内部twistの直和因子になる。
 各cohomologyの \((u)\)-adic有限filtrationの商は全ての \(u_a\) に殺される。
 それらは \(G_1,\ldots,G_{r-1}\) にも殺され、\(W=0\) と \(P(0)=0\) から
 \(G_r\) にも殺される。従って
 \[
 \Dsg^\Gamma(R)
 =\thick\operatorname{Res}D^b(\operatorname{gr}^\Gamma C).
 \]
 \(m=0\) では中心変数についてのKoszul複体とfiltrationを省けばよい。
 全ての \(\Gamma\)-次数付き \(C\)-加群が
 系\ref{cor:PI-generation}により外部積で生成される。

 この外部積のrestrictionのstabilizationを説明する。
 各 \(C_j=S_{h_j}/(G_j)\) の加群を十分高いsyzygyに置き換えて二項ambient自由分解で表し、
 \(m>0\) では \(T/(P)\) 上の残余体 \(k\) についても同じ操作をする。
 一つのhypersurface因子においてperfectな項は、restriction後に \(R\)-perfectとなる。
 この点を右加群の規約で明示する。選んだambientを \(A_i\)、potentialを \(f_i\)、
 その商を \(D_i\)、他のambientの積を \(A'\)、potentialの和を \(g'\) と書く。
 \(A_i\cong V_i\otimes\CC[f_i]\) と \(f_i=-g'\) の消去により、
 \(R=(A_i\otimes A')/(f_i+g')\) は \(A'\) 上左右とも自由である。
 他因子の外部積 \(M'\) は \(M'g'=0\) を満たし、各ambient上の有限自由分解の
 tensor productによって \(A'\)-perfectである。従って
 \[
 \begin{aligned}
 M'\otL_{A'}R&\simeq M'\otimes_{A'}R
       \xrightarrow{\sim}\operatorname{Res}(D_i\boxtimes M'),\\
 m'\otimes[a\otimes b]&\longmapsto\bar a\otimes m'b
 \end{aligned}
 \]
 は次数を保つ同型であり、右辺は \(R\)-perfectである。
 有限cone、shift、直和因子を取ることで、\(D_i\) を任意の \(D_i\)-perfect複体に
 置き換えても成立する。
 中心因子についても \(m>0\) では孤立性から \(P\ne0\) であり、
 \(T\cong V_0\otimes\CC[P]\) となるので同じ議論を適用できる。
 従ってsyzygyへの置換で失う項は特異圏で零である。

 二項ambient分解の微分を \(d_i\)、potential倍のnull-homotopyを \(s_i\) とする。
 外部tensor積上のKoszul符号付きの和 \(d=\sum d_i\)、\(s=\sum s_i\) は
 \(ds+sd=W\)、\(s^2=0\) を満たす。
 そのShamash周期化のtailはfactorizationの外部積に一致する。
 よって各boxのrestrictionのstabilizationは外部積関手のthick imageに入る。
 全 \(\Gamma\)-shiftも各因子のdegree shiftで実現されるので、
 補題\ref{lem:multiblock-ff}の完全忠実関手はMorita同値である。

 \emph{可換側との比較と整数次数への降下.}
 各Sklyanin因子に系\ref{cor:block}を用い、中心因子には恒等を用いる。
 可換側の外部積には \cite[Theorem 5.15, Corollary 5.18]{BFK} を適用する。
 命題\ref{prop:multiblock-embedding}と同様に、\(\Gamma\to\ZZ\) の有限核の双対による
 同変化と双対同変化を用いる。因子ごとのdegree shiftと \((3)=[2]\) のcoherenceは
 各比較で保存されるので、通常の整数次数へ降下すると
 \[
 \MF^{\ZZ}(A,W)\simeq
 \MF^{\ZZ}\!\left(\CC[x_j,y_j,z_j,u_a],
 \sum_j(x_j^3+y_j^3+z_j^3+h_j^2x_jy_jz_j)+P(u)\right)
 \tag{\(*\)}\label{eq:PI-model}
 \]
 を得る。

 \emph{射影模型と一意性.}
 中心正則元の商 \(R\) はnoetherian AS-Gorensteinで、injective dimensionは \(N-1\)、
 Gorenstein parameterは \(N-3\) である。
 \cite[Theorem 5.8]{CCKM} のMF--特異圏比較と
 \cite[Theorem 2.5]{Orlov} の残余windowを適用できる。
 \(N>3\) ならそのwindowは
 \[
 D^b(\qgr R)=\langle\mathcal K_R,\OO,\OO(1),\ldots,\OO(N-4)\rangle
 \]
 を与え、\(N=3\) なら例外対象はなく \(\mathcal K_R=D^b(\qgr R)\) である。
 degree shiftはwindowを通してrotationに対応する。
 可換potentialの勾配は各変数群に分離しており、各 \(F_{h_j}\) と \(P\) の
 アフィン臨界点は原点だけであるから、その和も原点以外に臨界点を持たない。
 従って \(Y_{\mathbf h,P}\) は滑らかで、同じOrlov比較は可換側を
 \((\Ku(Y_{\mathbf h,P}),\rho_{Y_{\mathbf h,P}})\) と同定する。
 \(N-2\ge3\) における射影同型類の一意性は
 \cite[Theorem 1.7]{LZ} による。
\end{proof}

\begin{lemma}[位数五の明示的パラメータ]\label{lem:order-five}
 \(h^6=5\) とする。原点 \(o=[1:-1:0]\) を持つ
 \(E_h=V(x^3+y^3+z^3+h^2xyz)\) 上で、\(p=[1:-1:-h]\) はちょうど位数五である。
 従って対応するSklyaninの平行移動 \(\sigma=t_p\) も位数五である。
\end{lemma}
\begin{proof}
 \(o\) における接線は \(3x+3y-h^2z=0\) であり、
 この直線上で \(F_h=(1+h^6/27)z^3\) となるので、\(o\) はinflection pointである。
 逆元は \([x:y:z]\mapsto[y:x:z]\) である。実際、この二点と \(o\) は一直線上にある。
 \(a=h^3\) とし、\(q=[a+1:a-1:-2h]\) と置く。
 \(F_h(q)=0\) であり、\(p\) における接線
 \[
 (3+a)x+(3-a)y+2h^2z=0
 \]
 は \(-q=[a-1:a+1:-2h]\) を通る。
 \(h\ne0\) では \(-q\ne p\) であるから、接線の加法則により \(q=2p\) である。
 \(q\) における接線の係数は
 \[
 [a^2+8a+3:a^2-8a+3:h^2(a^2+11)]
 \]
 であり、これを \(p\) で評価すると \(a(5-a^2)=h^3(5-h^6)\) となる。
 \(p\ne q\) は \(q=2p\) と \(p\ne o\) から従う。
 \(h^6=5\) では \(p\) がこの接線上にあるので \(2q+p=0\)、従って \(5p=0\) である。
 \(p\ne o\) かつ五は素数なので、位数はちょうど五である。
\end{proof}

\begin{remark}[二・三ブロックのsixfold・sevenfold]
 全ての \(h_j^6=5\) とする。
 \(r=2,m=2,P=u_1^3+u_2^3\) は二つの非可換ブロックを持つcubic sixfoldの
 rotation付き残余圏の完全な可換模型を与える。
 \(r=3,m=0\) は三つの非可換ブロックを持つcubic sevenfoldの完全な可換模型を与える。
 両者の方程式は定理\ref{thm:multiblock-model}の \(Y_{\mathbf h,P}\) であり、
 元の対象の一部分ではなく全残余圏を実現する。
\end{remark}

\begin{lemma}[無限位数の明示的パラメータ]\label{lem:order-infinite}
 \(h=2\) とする。楕円曲線
 \[
 E_2=V(x^3+y^3+z^3+4xyz),\qquad o=[1:-1:0]
 \]
 上の点 \(p_2=[1:-1:-2]\) は無限位数である。
\end{lemma}
\begin{proof}
 曲線と点は \(\mathbb Q\) 上定義されている。
 \(2^6+27=91\) は5,11のどちらでも割れないので、Hesseの勾配による判定から
 5と11で良い還元を持つ。
 \(z=1\)、\(z=0,y=1\)、\([1:0:0]\) の三つのchartに分けて列挙すると
 \[
 \#E_2(\mathbb F_5)=9,\qquad \#E_2(\mathbb F_{11})=12.
 \]
 実際、\(z=1\) の点数はそれぞれ8と11、\(z=0\) にはどちらも一点だけ存在する。
 この列挙は整数の合同計算で再現でき、\texttt{verify\_algebra.py} に全座標を記録する。
 良い還元では剰余標数と素なtorsionの還元写像は単射である
 \cite[Reduction of torsion points]{SnowdenDVR}。
 従って \(p_2\) の位数 \(N\) が有限なら、\(N\) の5と素な部分は9を割り、
 11と素な部分は12を割る。後者により5-primary部分はなく、前者により11-primary部分もなく、
 残りの素数について両方を比較すれば \(N\mid3\) となる。

 \(p_2\ne o\) である。一方、接線による倍点公式から
 \[
 2p_2=[9:7:-4],\qquad -p_2=[1:-1:2]
 \]
 であり、これらは異なる射影点である。
 従って \(3p_2\ne0\) であり、\(N\mid3\) は不可能である。
 \(p_2\) は無限位数である。
\end{proof}

\begin{remark}[真に無限位数を含むsixfoldとsevenfold]
 \(h_1=2\)、\(h_j=5^{1/6}\ (j\ge2)\) と取る。
 第一ブロックの平行移動は無限位数、他のブロックは位数5である。
 \((r,m)=(2,2)\)、\(P=u_1^3+u_2^3\) では可換sixfoldは
 \[
 V\bigl(x_1^3+y_1^3+z_1^3+4x_1y_1z_1
       +x_2^3+y_2^3+z_2^3+5^{1/3}x_2y_2z_2+u_1^3+u_2^3\bigr)
 \subset\PP^7
 \]
 となる。\((r,m)=(3,0)\) では第二・第三ブロックの係数を \(5^{1/3}\) とした
 cubic sevenfoldを得る。いずれも全残余圏とrotationの同値である。
 さらに他の有限位数パラメータを固定すれば、\(h_1^6\ne-27\) を満たす
 全ての \(h_1\in\CC\) について同じ結論が成り立つ。
 これは各ファイバーでの主張であり、同値の族をこのパラメータ直線全体に
 代数的に降下させたという主張ではない。
\end{remark}

\subsection{外部積生成の十分条件とその限界}

有限位数条件は \(j\ge2\) のみに用いる。
従ってより一般に、\(C_1\) がnoetherian、\(j\ge2\) の各 \(C_j\) が
有限型可換斉次中心上有限で、周囲が命題\ref{prop:block-support}の仮定を満たせば、
同じ生成と安定化の議論を適用できる。
定理\ref{thm:one-sided-boxes}では任意のnoetherianな係数代数を保てる一方、
新たに加える因子の有限中心性は必要である。
従って二つ以上の独立な無限位数THCRを持つ場合は、この証明だけでは解けない。
これはその場合の非存在を意味せず、条件 (G) が未解決であるという意味である。

二つ以上の無限位数因子を含む場合に、ここでの生成法を拡張するには
条件 (G) が必要である。その十分な条件として、
THCRに対して任意のnoetherian係数を許す外部積生成、
または粗化次数を保った対角の有限構成が考えられる。
条件 (G) が得られれば、残りの有限中心因子を
定理\ref{thm:one-sided-boxes}で追加できる。
可換ambientのThom--Sebastiani定理を非可換ambientへ無条件に適用したわけではない。
先行研究上の新規性は未確認である。


\section{例外パラメータと相対復元の適用範囲}\label{sec:scope}

本節では、例外係数で可換復元が成立しない理由をpropernessによって説明し、
中心三次量子化族に形式復元・étale局所実現を適用するための
仮定をまとめる。これらは一般点主定理の仮定を変更するものではない。

\subsection{例外パラメータでpropernessが失われる}

\begin{proof}[命題\ref{prop:bad-intro}の証明]
 \(h_0^6=-27\) とする。\(E_{h_0}\) は三本の直線からなる三角形である。
 まず
 \[
 S_{h_0}/(G_{h_0})\simeq B(E_{h_0},L,\sigma_{h_0})                 \tag{6}\label{eq:triangle}
 \]
 を確認する。楕円曲線の場合だけの公式を、この特殊値へ無条件に適用しない。
 \(h_0\) の近傍で、関係式を多重線形化した行列
 \[
 M_h(x,y,z)=
 \begin{pmatrix}-y&x&-hz\\-hx&-z&y\\z&-hy&-x\end{pmatrix}
 \]
 は \(E_h\) 上でrank二である。実際その行列式は
 \(-h(x^3+y^3+z^3+h^2xyz)\) であり、rank一なら小行列式から
 \(yz+hx^2=yz-hx^2=0\) と巡回的な式を得て、\(x=y=z=0\) となる。
 核は相対的な射 \(\sigma_h\) を与える。グラフを転置すると \(h\) を \(-h\) に
 変える同じ関係式になるので、この射は相対自己同型である。

 \(\mathcal L_d=\bigotimes_{j=0}^{d-1}(\sigma^j)^*\OO_E(1)\) は、
 滑らかなファイバーでは次数 \(3d\)、三角形では各成分次数 \(d\) を持つ。
 \(d>0\) なら \(H^1(E_h,\mathcal L_d)=0\)、\(h^0=3d\) であり、
 相対切断加群はlocally freeで底変更と両立する。
 自然な次数付き写像 \(S_h\to\bigoplus_dH^0(E_h,\mathcal L_d)\) の下で、
 \(G_h\) は一般の滑らかなファイバーで零なので、局所自由性から恒等的に零となる。
 三角形のtwisted homogeneous coordinate ringが次数一で生成されることを確認する。
 \(L\) のbasepoint-free pencilに対するpencil trickを使うと、
 多次数 \((d,d,d)\) の \(M\) について
 \(H^0(L)\otimes H^0(M)\to H^0(LM)\) は
 \(H^1(ML^{-1})=0\) なら全射である。
 \(d\ge2\) では各成分上の正値性からこの消滅が従う。
 \(d=1\) で \(ML^{-1}\) が非自明なら、多次数零の直線束の切断は
 三角形の三成分の接合条件により零で、Euler数も零なので同じ消滅が成り立つ。
 \(M\simeq L\) なら、平面三次曲線の斉次イデアルに二次式がないことから全射である。
 従って \(S_{h_0}/(G_{h_0})\to B(E_{h_0},L,\sigma_{h_0})\) は全射である。
 補題\ref{lem:center}の正則性から左辺のHilbert係数は \(1,3d\) であり、
 右辺の係数と一致する。よってこの全射は同型で、\eqref{eq:triangle}が従う。
 \(\mathcal L_d\) の各成分上の次数が \(d\) とともに増大するので
 \(\sigma\)-amplenessも成り立ち、tailsは \(\operatorname{coh}E_{h_0}\) と同値になる。

 Gorenstein parameter零のOrlov同値と \cite[Theorem 5.8]{CCKM}から
 \[
 \MF^\ZZ(S_{h_0},G_{h_0})\simeq D^b(\operatorname{coh}E_{h_0}).
 \]
 三角形の節点 \(p\) のskyscraper sheafは、任意に高い次数で非零の自己Extを持つ。
 節点の完備局所環 \(\CC[[s,t]]/(st)\) の剰余体の最小自由分解は無限であり、
 局所大域Extスペクトル系列では高次Extの支持が \(p\) だけなので、
 この非零性は大域自己Extにも残る。従ってこのdg圏はproperでない。

 \(m=0\) の場合の非propernessはこれで示された。以下では \(m\ge1\) とする。
 \(P=\sum_i u_i^3\) とし、第\ref{sec:TS}節の二ブロック次数 \(\Gamma\) と
 Koszul factorization \(K_P\) を使う。
 \(K_P\) の全微分行列は \((u_1,\ldots,u_m)\) に属する。
 もし恒等写像が零ホモトピーなら、係数を原点で評価して非零自由加群上の恒等写像が
 零となるので矛盾する。従って \(H^0\REnd(K_P)\ne0\) である。
 節点の対象を \(N\) とすると、外部積の完全忠実性と体上のKünneth公式により
 \[
 H^r\REnd(N)\otimes H^0\REnd(K_P)
 \subset H^r\REnd(N\boxtimes K_P)
 \]
 が直和因子として残る。よって \(\Gamma\)-次数の全体圏はnonproperである。
 もし通常の整数次数付き圏がproperなら、その有限群 \(H_\Gamma\simeq\mu_3\) による
 同変圏も、Hom複体の不変部分が直和因子となることからproperとなる。
 これは式\eqref{eq:eq}に反する。

 さらに、この議論は単なるHom総次元の非有限性より強い。
 \(\Gamma\)-次数の一つの対象に、任意に高い正次数の自己Homが残っている。
 有限群同変化から忘却すると、同変Homは通常のHomの直和因子なので、
 通常次数でも同じ正次数非消滅を持つ一つの対象が得られる。
 \(m>0\) ではOrlovの完全忠実windowによりこれを \(D^b(\qgr B_{h_0,m})\) へ送れ、
 \(m=0\) では上の三角形との同値を用いる。
 もし \(\qgr B_{h_0,m}\) のhomological dimensionが有限なら、
 その有界導来圏の任意の固定された二対象の間の正次数Homは十分高次で消える。
 これは有界複体のcohomologyとExtからなる有限スペクトル系列で分かる。
 今得た対象はこれに反するから、homological dimensionは無限である。

 一方、周囲の \(S_{h_0}[u]\) はAS-regularであり、三次の中心元の正則性は
 \(m>0\) なら最後の中心変数についてmonicであること、\(m=0\) なら
 補題\ref{lem:center}から従う。従ってHilbert seriesは依然として
 \((1-t^3)/(1-t)^{m+3}\) であり、degree shiftの三乗も \([2]\) である。
 滑らかな射影多様体のadmissible成分はproperなので、求める非存在が従う。
\end{proof}

この例は、量子化を構成する開集合をさらに縮めて例外値を除くことを妨げない。
従って、零を含む十分小さな開集合での存在問題に対する反例ではない。

\subsection{中心三次量子化族への適用}

三次の中心元によるAS-regularな周囲の商
\[
 B_\hbar=A_\hbar/(\Omega_\hbar),\qquad B_0=\CC[x]/(F)
\]
を、次数ごとにflatな族として構成できていれば、
その形式完備化には定理\ref{thm:main}が適用できる。
この適用にはPoisson構造の分類を再度用いず、各Artin次数でのgraded flatnessだけを使う。
従って、量子化の構成自体の証明と、本稿の相対残余圏・復元の証明は独立に検証できる。
論文Aの代数的量子化定理への適用では、そこで構成する次数ごとのflatnessを用いる。
代数の量子化と圏の復元は異なる段階であり、それぞれの仮定を区別する。

形式的な可換モデルは、半古典極限の固定された \(X\) と同じとは限らない。
定理\ref{thm:global}の代表式では \(h^2xyz\) が現れる。
このモデルを初期の \(F\) と取り違えない。

\subsection{一般の有限パラメータへの適用範囲}

定理\ref{thm:etale}により、smooth properな代数的pairの族が与えられれば、
可換cubic点のétale近傍で存在問題は解ける。
一般の量子化族へこのétale局所定理を適用するには、
必要な相対saturationを別途確立する必要がある。
各Artin切断のsmooth proper性だけから、有限表示の代数的模型の存在を結論しない。

独立な複数ブロックでは、命題\ref{prop:multiblock-embedding}で得た完全忠実な像が
全MF圏を生成するかが、二つ以上の無限位数因子の場合の核心である。
命題\ref{prop:block-support}と条件\eqref{eq:block-generation}は、その具体的な帰着を与える。
高々一つの無限位数因子を持つ場合は定理\ref{thm:multiblock-model}で解決した。
中心三次式内の混合項は定理\ref{thm:global}と定理\ref{thm:multiblock-model}に含まれるが、
非可換ブロックと中心変数の混合項は含まれない。
また、étale局所モデルの元の底への降下と、実現開部分の外のファイバーの分類は
別の大域問題として残る。


\section{可換cubic模型の存在を妨げるrotationの不変量}
\label{sec:rotation-obstruction}

本節では、次数付き代数の条件と、可換cubicのrotationを持つという条件を区別する。
三次中心正則元を持つAS-regular代数の商というだけでは、可換模型の存在は従わない。
さらに、商の射影圏が滑らかな可換cubicのcoherent圏と同値であり、周囲の代数が
Calabi--Yauであっても、与えられたrotationを保つ可換模型が存在しない例がある。
以下の二つの反例は有限位数のZhang twistで構成する。
これらが本稿で構成する特定のPoisson量子化族のファイバーであることは主張しない。

\subsection{Serre関手とrotationの必要条件}

\begin{proposition}[Serre--rotation障害]
\label{prop:serre-rotation-obstruction}
 \(\cC\) をSerre関手 \(S_{\cC}\) を持つproperなdg圏、
 \(\rho\) をその自己同値とする。\(N=n+2\ge5\) とし、
 滑らかなcubic \(Y^n\subset\PP^{N-1}\) とのenhanced pair同値
 \[
 (\cC,\rho)\simeq(\Ku(Y),\rho_Y)
 \]
 が存在するなら、
 \begin{equation}\label{eq:serre-rotation-necessary}
 S_{\cC}\simeq\rho^{\,3-N}[N-2],\qquad
 \Theta_N(\cC,\rho):=S_{\cC}\rho^{N-3}[2-N]\simeq\Id
 \end{equation}
 である。特に \(\rho^3\simeq[2]\) だけでは十分でない。
\end{proposition}
\begin{proof}
 可換斉次座標環 \(R_Y=\CC[z_1,\ldots,z_N]/(F_Y)\) のKrull次元は
 \(N-1\)、標準加群は \(R_Y(3-N)\) である。
 滑らかさによって原点以外の錐は正則であり、次数付きAuslander--Reiten双対性は
 \(\Dsg^{\ZZ}(R_Y)\) のSerre関手を
 \((3-N)[N-2]\) と与える
 \cite[Theorem 7, p.~161]{UeyamaIsolated}。
 Orlovのwindowで内部次数シフト \((1)\) は \(\rho_Y\) に対応するので、
 \(S_{\Ku(Y)}\simeq\rho_Y^{3-N}[N-2]\) を得る。
 Serre関手は同値によって保存されるから、pair同値を通してこの式を移せばよい。
\end{proof}

\subsection{有限位数Zhang twistと次数シフト}

以後 \(\zeta\) は \(\zeta^3=1\)、\(\zeta\ne1\) を満たす。
\(S=\CC[x_0,\ldots,x_{N-1}]\)、\(F=\sum_i x_i^3\)、
\(X=V(F)\subset\PP^{N-1}\) とする。
\(q_i\in\{1,\zeta,\zeta^2\}\) を選び、
\(\tau(x_i)=q_i x_i\) と置く。\(\tau^3=\Id\)、\(\tau(F)=F\) である。
\(f:X\to X\) を \(f^*=\tau\) に対応する射影自己同型とする。
右Zhang twist \(A=S^{\tau}\) の積を、斉次元に対して
\begin{equation}\label{eq:right-zhang-product}
 a*b=a\,\tau^{\deg a}(b)
\end{equation}
で定義する。右辺の積は元の可換多項式環の積である。

\begin{lemma}\label{lem:zhang-cubic-structure}
 \(A\) は左右noetherianなKoszul AS-regular domainで、
 global dimensionは \(N\)、Hilbert級数は \((1-t)^{-N}\) である。
 同じ次数付きベクトル空間内の元 \(\Omega=F\in A_3\) は中心正則元であり、
 \[
 B=A/(\Omega)\simeq (S/(F))^{\tau},\qquad
 H_B(t)=\frac{1-t^3}{(1-t)^N}.
 \]
 さらに
 \begin{equation}\label{eq:zhang-qgr-pair}
 \qgr B\simeq\operatorname{coh}X,\qquad
 (\mathcal K_B,\rho_B)\simeq
 (\Ku(X),f_*\rho_X)
 \end{equation}
 である。従って \(\mathcal K_B\) はsmooth properで、
 \(\rho_B^3\simeq[2]\) を満たす。
\end{lemma}
\begin{proof}
 積\eqref{eq:right-zhang-product}は結合的であり、生成元の関係式は
 \[
 x_i*x_j=(q_j/q_i)x_j*x_i\qquad(i<j)
 \]
 となる。従って \(A\) はscalarによる可逆なOre拡大の反復である。
 順序付き単項式 \(x_0^{a_0}*\cdots*x_{N-1}^{a_{N-1}}\) は
 元の可換単項式の非零scalar倍なので基底をなし、domain性とHilbert級数が従う。
 Ore拡大のHilbert基底定理から左右noetherian性を得る。
 多項式環のKoszul分解をtwistすると、\(k\) の長さ \(N\) の線形自由分解になる。
 その双対もtwistされたKoszul分解であるから
 \[
 \operatorname{Ext}^j_A(k,A)=
 \begin{cases}k(N)&j=N,\\0&j\ne N\end{cases}
 \]
 となり、左側にも同じ式が成り立つ。
 同じtwistを双加群Koszul分解に施した有限分解、またはOre拡大のglobal dimensionの公式により、
 global dimensionは \(N\) である。これはZhang twistがAS-regular性を保つことの
 この場合の具体的な確認である \cite{ZhangTwist,RRZ}。

 任意の斉次 \(a\in A_d\) に対して
 \[
 a*F=a\tau^d(F)=aF,
 \qquad F*a=F\tau^3(a)=Fa
 \]
 であるから、\(F\) は中心的である。
 domain性から正則であり、\((F)\) の次数付きベクトル空間は元の多項式環の
 \((F)\) と一致する。このため商も右twistであり、表示したHilbert級数を得る。
 なお \(x_i*x_i*x_i=q_i^3x_i^3=x_i^3\) なので、
 \(\Omega\) はtwist後の生成元の三乗和としても同じ元である。

 次数付き右 \(R=S/(F)\)-加群 \(M\) に、
 \(m*a=m\tau^{\deg m}(a)\) という作用を入れた加群を \(\mathsf T(M)\) と書く。
 \(\mathsf T\) は逆twistを逆関手に持つ完全同値で、有限生成性とtorsionを保つ。
 従って \(\operatorname{gr}R\simeq\operatorname{gr}B\) および
 \(\operatorname{coh}X\simeq\qgr B\) を与える。

 次数シフトの規約を \(M(1)_d=M_{d+1}\) とする。
 \(\mathsf T(M)(1)\) をuntwistすると、次数 \(d\) の元 \(m\in M_{d+1}\) に対して
 \[
 m\cdot a=m*\tau^{-d}(a)=m\tau(a)
 \]
 となる。従って、untwist後の \((1)\) は \(M(1)_{\tau}\) であり、
 coherent圏では \(f_*(-\otimes\OO_X(1))\) に対応する。
 特に \(\OO_B(i)\) は \(\OO_X(i)\) に対応する。
 \(f_*\OO_X\simeq\OO_X\) なので、\(f_*\) は \(\OO_X\) によるleft mutationと可換である。
 よって残余圏とrotationについて\eqref{eq:zhang-qgr-pair}が従う。
 \(f_*\) は \(\rho_X\) と可換で、\(f_*^3\simeq\Id\) であるため、
 \((f_*\rho_X)^3\simeq\rho_X^3\simeq[2]\) となる。
\end{proof}

\begin{remark}[右twistにおけるNakayamaの規約]
\label{rem:zhang-nakayama}
 上の右twistのNakayama自己同型は
 \begin{equation}\label{eq:right-twist-nakayama}
 \mu_A=\tau^{-N}\xi_{\det\tau},\qquad
 \mu_A(x_i)=\frac{\prod_jq_j}{q_i^N}x_i,
 \end{equation}
 と取れる。ここで \(\xi_c\) は次数 \(d\) 上で \(c^d\) 倍する自己同型である。
 \cite[Theorem 5.4]{RRZ}は左twist
 \(a\circ b=\sigma^{\deg b}(a)b\) に対して
 \(\mu=\sigma^N\xi_{\det\sigma}^{-1}\) を与える。
 \(a\mapsto\tau^{-\deg a}(a)\) は右twistから \(\sigma=\tau^{-1}\) の左twistへの同型であり、
 この公式を移すと\eqref{eq:right-twist-nakayama}になる。
 このように左右の規約を固定してから用いる。
\end{remark}

\subsection{三次元の滑らかな反例}

\begin{theorem}[Serre--rotation条件を破るsmoothな非可換cubic]
\label{thm:zhang-threefold-counterexample}
 \(N=5\)、\(\tau=\operatorname{diag}(\zeta,1,1,1,1)\) とし、
 補題\ref{lem:zhang-cubic-structure}の \(B\) を取る。
 この \(B\) はnoetherianなKoszul AS-regular五次元代数の三次中心正則商であり、
 \(\qgr B\simeq\operatorname{coh}X\) は滑らかなFermat cubic threefoldのcoherent圏と同値である。
 しかし、いかなる滑らかな可換cubic threefold \(Y\) に対しても
 \[
 (\mathcal K_B,\rho_B)\not\simeq(\Ku(Y),\rho_Y).
 \]
\end{theorem}
\begin{proof}
 代数の諸性質とsmooth proper性は補題\ref{lem:zhang-cubic-structure}で示した。
 \(T=f_*|_{\Ku(X)}\) と書くと、同補題により
 \(\rho_B=T\rho_X\)、\(T^3\simeq\Id\)、\(T\rho_X\simeq\rho_XT\) である。
 underlying圏は \(\Ku(X)\) なので、そのSerre関手は
 \(S_X\simeq\rho_X^{-2}[3]\) である。
 もし可換pair模型が存在すれば、命題\ref{prop:serre-rotation-obstruction}により
 \[
 \rho_X^{-2}[3]\simeq (T\rho_X)^{-2}[3]
       \simeq T^{-2}\rho_X^{-2}[3]
 \]
 となり、\(T^2\simeq\Id\) を得る。\(T^3\simeq\Id\) と合わせると \(T\simeq\Id\) となる。

 これはHochschild homologyへの作用に反する。
 例外対象のHochschild homologyは次数零だけにあるため、加法性とHKR同型から
 \[
 \HH_1(\Ku(X))\simeq\HH_1(X)\simeq H^{2,1}(X)
 \]
 である。Griffithsのresidue同型\cite{Griffiths}は
 \[
 J(F)_1\xrightarrow{\sim}H^{2,1}(X),\qquad
 a\longmapsto\operatorname{Res}_X\frac{a\Omega_{\mathrm{proj}}}{F^2},
 \quad
 \Omega_{\mathrm{proj}}=
 \sum_{i=0}^4(-1)^i x_i\,dx_0\wedge\cdots\wedge\widehat{dx_i}\wedge\cdots\wedge dx_4
 \]
 と書ける。\(\tau(F)=F\)、\(\tau^*\Omega_{\mathrm{proj}}=(\det\tau)\Omega_{\mathrm{proj}}
 =\zeta\Omega_{\mathrm{proj}}\) である。
 \(J(F)=\CC[x_0,\ldots,x_4]/(x_0^2,\ldots,x_4^2)\) だから、
 \(f^*\) の \(H^{2,1}(X)\) 上の固有値は
 \[
 \zeta^2\ \text{が1個},\qquad \zeta\ \text{が4個}
 \]
 となる。\(f_*\) の作用はその逆であり、いずれにしても恒等ではない。
 従って \(T\not\simeq\Id\) であり、矛盾する。
\end{proof}

\begin{remark}
 この例では \(\rho_B^3\simeq[2]\) とsmooth proper性の両方が成立する。
 一方、\eqref{eq:right-twist-nakayama}のNakayama自己同型は非自明であり、
 周囲の \(A\) はCalabi--Yauではない。
 次の例は、この最後の条件を加えても一般的な存在命題が成立しないことを示す。
\end{remark}

\subsection{Calabi--Yauな周囲を持つ四次元の反例}

\begin{theorem}[Calabi--Yauな非可換cubicのrotation障害]
\label{thm:zhang-cy-fourfold-counterexample}
 \(N=6\)、\(\tau=\operatorname{diag}(\zeta,\zeta^2,1,1,1,1)\) とし、
 補題\ref{lem:zhang-cubic-structure}の \(A\)、\(B\) を取る。
 このとき \(A\) は六次元Calabi--Yauなnoetherian Koszul AS-regular代数であり、
 \(\Omega\) は三次中心正則元である。
 \(\qgr B\) は滑らかなFermat cubic fourfold \(X\) のcoherent圏と同値であり、
 \(\mathcal K_B\) はsmooth properな2-Calabi--Yau圏である。
 \(\rho_B^3\simeq[2]\) および\eqref{eq:serre-rotation-necessary}も満たす。
 しかし、いかなる滑らかな可換cubic fourfold \(Y\) に対しても
 \[
 (\mathcal K_B,\rho_B)\not\simeq(\Ku(Y),\rho_Y).
 \]
 より具体的に、Hochschild \emph{homology}の次数零部分について
 \begin{equation}\label{eq:cy-fourfold-hh0-obstruction}
 \dim\HH_0(\mathcal K_B)^{\rho_B}=8,
 \qquad
 \dim\HH_0(\Ku(Y))^{\rho_Y}=20
 \end{equation}
 である。
\end{theorem}
\begin{proof}
 \(\det\tau=1\)、\(\tau^6=\Id\) だから、
 \eqref{eq:right-twist-nakayama}によって \(\mu_A=\Id\) である。
 AS-regular性とNakayamaの自明性はCalabi--Yau性を与える
 \cite[Lemma 1.2, Theorem 5.4]{RRZ}。
 残りの代数的性質と \(\mathcal K_B\simeq\Ku(X)\) は
 補題\ref{lem:zhang-cubic-structure}による。
 cubic fourfoldのKuznetsov成分は2-Calabi--Yauなので
 \(S_{\mathcal K_B}\simeq[2]\) であり、
 \(\rho_B^{-3}[4]\simeq[2]\) からSerre--rotation条件も成立する。

 Hochschild homologyの作用を計算する。
 Mukai--HKR同定と残余成分の射影によって
 \begin{equation}\label{eq:cubic-fourfold-hh0-split}
 \HH_0(\Ku(X))\simeq
 H^{2,2}(X)_{\mathrm{prim}}\oplus (A_2\otimes_{\ZZ}\CC)
 \end{equation}
 と書ける。ここで前者は20次元、後者は2次元である。
 \(A_2\) はcubic fourfoldのKuznetsov成分に標準的に含まれる二つのMukai classの張る格子である。
 \cite[Proposition 3.12]{Huy}により、\(\rho_X\) はprimitive部分上で恒等、
 \(A_2\) 上で三つの根を巡回置換するので、後者の固有値は \(\zeta,\zeta^2\) である。
 この記述は一般の滑らかなcubic fourfoldにも成り立つため、
 \(\dim\HH_0(\Ku(Y))^{\rho_Y}=20\) を得る。

 射影自己同型 \(f\) の作用は分解\eqref{eq:cubic-fourfold-hh0-split}を保ち、
 \(A_2\) 上で恒等である。
 実際、その二つの標準classは直線 \(\ell\subset X\) 上の
 \(\OO_\ell(1)\)、\(\OO_\ell(2)\) の残余射影から作られる。
 \(f\) は直線を直線に送り、全ての直線のcohomology classは同じである。
 残余射影のcohomology上の作用は \(\OO_X,\OO_X(1),\OO_X(2)\) から定まり、
 これらを \(f\) が固定するので、二つのMukai classも固定される。
 Mukai--HKR同定に入るTodd classは射影自己同型で不変であるから、
 これは \(\HH_0\) 上の作用についても成立する。

 primitive部分ではresidue表示
 \[
 J(F)_3\xrightarrow{\sim}H^{2,2}(X)_{\mathrm{prim}},\qquad
 a\longmapsto\operatorname{Res}_X\frac{a\Omega_{\mathrm{proj}}}{F^3}
 \]
 を用いる。ここでresidue表示はpole filtrationの商を取った
 \(J(F)_3\simeq\operatorname{Gr}_F^2H^4(X)_{\mathrm{prim}}\) と、
 Hodge分解による \(\operatorname{Gr}_F^2\simeq H^{2,2}_{\mathrm{prim}}\) の合成を意味する。
 射影自己同型は両方の構造を保つ。
 \(\det\tau=1\)、\(\tau(F)=F\) なので、ここで \(f^*\) の作用は
 \(J(F)_3\) 上の通常の置換 \(a\mapsto\tau(a)\) に一致する。
 FermatのJacobian環は \(J(F)=\CC[x_0,\ldots,x_5]/(x_0^2,\ldots,x_5^2)\) であり、
 次数三の基底は平方自由な三変数単項式20個である。
 その固有値を分類すると
 \[
 \begin{array}{c|c|c}
 \text{単項式に含まれる変数}&\text{個数}&\tau\text{の固有値}\\\hline
 x_0,x_1\text{の両方}&4&1\\
 x_0,x_1\text{のどちらも含まない}&4&1\\
 x_0\text{だけ}&6&\zeta\\
 x_1\text{だけ}&6&\zeta^2
 \end{array}
 \]
 となる。従って \(f_*\) のprimitive部分の固定空間も8次元である。

 補題\ref{lem:zhang-cubic-structure}から \(\rho_B=f_*\rho_X\) である。
 primitive部分で \(\rho_X\) は恒等なので、\(\rho_B\) の固定空間は8次元である。
 \(A_2\) 上では \(f_*\) が恒等であり、\(\rho_X\) に固有値1はない。
 よって \eqref{eq:cy-fourfold-hh0-obstruction}が成立する。
 enhanced pair同値は \(\rho\) の \(\HH_0\) 上の作用を共役するので、
 固定空間の次元の不一致が同値の不存在を示す。
\end{proof}

\subsection{smooth properな連結族で保存される不変量}

前の反例と可換cubicのpairとの違いは、連結な圏の族を通じても検出できる。
Hochschild homologyの各次数の次元が個別に一定であることを仮定せず、
perfect複体のEuler標数を用いる。

\begin{proposition}[rotation不変Euler標数の局所一定性]
\label{prop:rotation-euler-locally-constant}
 \(S\) を複素スキームとし、\(\mathcal C_S\) をsmooth properな代数的dg圏の族とする。
 \(\rho_S\) は可逆perfect bimoduleで与えられる自己同値の族とし、
 整数 \(d>0\)、\(a\in\ZZ\) と自然同型 \(\rho_S^d\simeq[2a]\) を仮定する。
 幾何学的点 \(s\) に対して
 \[
 \chi_{\rho}(\mathcal C_s)
   =\sum_{i\in\ZZ}(-1)^i
        \dim_{\kappa(s)}\HH_i(\mathcal C_s)^{\rho_s}
 \]
 と定義する。この和は有限であり、\(s\mapsto\chi_{\rho}(\mathcal C_s)\) は局所一定である。
 特に \(S\) が連結なら一定である。
\end{proposition}
\begin{proof}
 相対Hochschild複体 \(H=\HH(\mathcal C_S/S)\) は \(S\)-perfectであり、
 導来底変更と両立する \cite[Lemmas 3.4, 3.6]{PerryIHC}。
 実際、étale局所的に \(\mathcal C_S=\Perf(E)\) と書けば、
 \(H=E\otimes^{\mathbf L}_{E^{\mathrm{op}}\otimes^{\mathbf L}_S E}E\) である。
 smoothnessによって対角は有限cell双加群の直和因子として表され、propernessによって
 各cellをtraceした複体は \(S\)-perfectとなる。perfect性は直和因子で保たれる。
 これらの構成から底変更も従う。

 \(\rho_S\) が誘導する \(H\) の自己同型を \(r\) と書く。
 Hochschild homologyの加法性により \([1]\) の誘導作用は \(-1\)、従って
 \([2a]\) の作用は恒等である \cite[Section 4, Example 5.6]{CisinskiTabuada}。
 よって \(r^d=1\) であり、
 \[
 e=\frac1d\sum_{j=0}^{d-1}r^j\in\operatorname{End}_{D(S)}(H)
 \]
 は冪等元である。\(\Perf(S)\) は冪等完備なので、その像 \(H^{\mathrm{inv}}\) はperfectである。
 幾何ファイバーで \(e\) の作用は有限巡回群の不変部分への射影であるから、
 \[
 H^{-i}(H^{\mathrm{inv}}\otimes_S^{\mathbf L}\kappa(s))
     =\HH_i(\mathcal C_s)^{\rho_s}.
 \]
 perfect複体のファイバーEuler標数は、局所的な有限階数ベクトル束複体の
 階数の交代和であるため局所一定である。これが主張である。
\end{proof}

\begin{corollary}[Calabi--Yau反例と可換pairの連結成分の分離]
\label{cor:zhang-cy-component-obstruction}
 定理\ref{thm:zhang-cy-fourfold-counterexample}のpairは、
 \(\rho^3\simeq[2]\) を相対的に備えた連結smooth proper代数的dg圏の族によって、
 滑らかな可換cubic fourfoldの標準pairに結ぶことができない。
 実際、両者の不変Euler標数はそれぞれ
 \[
 \chi_{\rho_B}(\mathcal K_B)=10,
 \qquad
 \chi_{\rho_Y}(\Ku(Y))=22
 \]
 である。
\end{corollary}
\begin{proof}
 cubic fourfoldのKuznetsov成分のHochschild homologyは次数 \(-2,0,2\) にのみあり、
 次数 \(-2,2\) は各一次元である。
 標準rotationはsymplecticなので、この二つの部分に恒等として作用する
 \cite[Proposition 3.12]{Huy}。
 定理\ref{thm:zhang-cy-fourfold-counterexample}の \(f\) も、
 \(\det\tau=1\)、\(\tau(F)=F\) により \(H^{3,1}(X)\) のresidueを固定し、
 その複素共役 \(H^{1,3}(X)\) も固定する。
 従って \(f_*\rho_X\) についても \(\HH_{\pm2}\) の不変部分は各一次元である。
 次数零の不変部分は\eqref{eq:cy-fourfold-hh0-obstruction}で計算したので、
 Euler標数はそれぞれ \(1+8+1=10\)、\(1+20+1=22\) となる。
 命題\ref{prop:rotation-euler-locally-constant}を適用する。
\end{proof}

\begin{remark}[量子化族との論理的関係]
\label{rem:counterexample-quantization-scope}
 定理\ref{thm:zhang-cy-fourfold-counterexample}は、
 「多項式環と同じHilbert級数を持つCalabi--Yau AS-regular代数の三次中心正則商であり、
 射影圏がsmooth properである」という条件による一般的な可換pair復元を反証する。
 またSerre--rotation条件を追加するだけでも十分でない。
 underlyingなKuznetsov成分には既に可換模型 \(X\) があるが、指定されたrotationが異なるのである。

 一方、この代数が本稿の特定のPoisson構造から作る、零パラメータで可換な
 中心三次量子化族のファイバーであることは示していない。
 \(\tau\) の有限位数性は中心元の構成に本質的であり、
 \(\zeta\) を任意の複素パラメータに置き換えても三次中心元が残るわけではない。
 従って本例から、その特定の量子化族に対する存在問題の否定は導かない。
 量子化族について復元の存在を問う際には、
 対象とする量子化族と残余pairの相対構成を指定し、
 その上での存在問題を述べる必要がある。
\end{remark}

\section{完全忠実比較を同値へ強める条件}
\label{sec:serre-generation}

多ブロック族では、既に構成した完全忠実関手の像が全体を生成するかが問題になる。
外部積による直接的な生成証明とは別に、properness、Serre関手および連結性から
本質的全射性を判定できる。ここでの \(\HH^0\) はcohomologyであり、
前節の反例に用いたhomologyの \(\HH_0\) とは異なる。

\begin{proposition}[Serreと連結性による全生成判定]
\label{prop:serre-connected-generation}
 \(\mathcal A\) を非零のsmooth properな、冪等完備pretriangulated dg圏、
 \(\mathcal D\) をproperな、冪等完備pretriangulated dg圏とする。
 \(\mathcal D\) はSerre関手 \(S_{\mathcal D}\) を持つとし、
 dg完全忠実関手 \(\Phi:\mathcal A\to\mathcal D\) が
 \[
 \Phi S_{\mathcal A}\simeq S_{\mathcal D}\Phi
 \]
 を満たすと仮定する。
 さらに環 \(\HH^0(\mathcal D)\) が \(0,1\) 以外の冪等元を持たないなら、
 \(\Phi\) はdg Morita同値である。
 特に \(\HH^0(\mathcal D)=\CC\) なら、この連結性条件は満たされる。
 \(\mathcal D\) のsmoothnessは仮定しなくてよい。
\end{proposition}
\begin{proof}
 まず \(\Phi\) の像がadmissibleであることを示す。
 \(\mathcal A\) はsmooth properなので、compact generator \(G\) と
 smooth proper dg代数 \(E=\REnd_{\mathcal A}(G)\) を用いて
 \(\mathcal A\simeq\Perf(E)\) と書ける。
 \(D\in\mathcal D\) に対し、\(\RHom_{\mathcal D}(\Phi G,D)\) は
 \(\mathcal D\) のpropernessにより \(\CC\)-perfectである。
 smoothなdg代数上では、このようなpseudo-perfect加群はperfectである。
 実際、対角 \(E\) を有限cell双加群の直和因子として表し、この加群をtensorすると、
 \(\CC\)-perfectな係数を持つ有限cell \(E\)-加群の直和因子として表される。
 したがって \(\mathcal A\) の対象で表現され、\(\Phi\) の右随伴を与える。
 opposite圏で同じ議論を行えば左随伴も得る。
 この構成はdg加群の表現可能性として行われるので、随伴もenhancedである。

 \(\mathcal B=\Phi(\mathcal A)^\perp\) とする。
 \(D\in\mathcal B\)、\(A\in\mathcal A\) に対し、Serre双対性から
 \[
 \Hom_{\mathcal D}(D,\Phi A[i])
 \simeq
 \Hom_{\mathcal D}(S_{\mathcal D}^{-1}\Phi A[i],D)^*
 \simeq
 \Hom_{\mathcal D}(\Phi S_{\mathcal A}^{-1}A[i],D)^*=0
 \]
 となる。従って右直交と左直交は一致する。
 admissibilityによる各対象の分解三角形は、二方向のHom消滅によって分裂する。
 よって \(\mathcal D\) は \(\Phi(\mathcal A)\) と \(\mathcal B\) の直交直和である。
 全てのcohomological shiftのHomが消えるため、これはdg Moritaの意味でも直和分解となる。

 この直和分解は対角双加群の分解を与え、第一成分への射影が
 \(\HH^0(\mathcal D)=\operatorname{End}(\Delta_{\mathcal D})\) の冪等元になる。
 \(\mathcal A\ne0\) なのでこの冪等元は零でない。
 もし \(\mathcal B\ne0\) なら恒等でもないため、仮定に反する。
 従って \(\mathcal B=0\) であり、\(\Phi\) は本質的全射である。
\end{proof}

\begin{corollary}[任意位数の独立ブロックに対する十分条件]
\label{cor:arbitrary-order-conditional-model}
 \(N=3r+m\ge5\) とし、全ての \(h_j^6\ne-27\) とする。
 \(P(u)\) は原点のみを臨界点に持つ中心三次式とし、\(m=0\) では省く。
 独立ブロックの代数
 \[
 B=\left(\bigotimes_{j=1}^rS_{h_j}\right)[u]
     \Big/\left(\sum_jG_{h_j}+P(u)\right)
 \]
 の通常次数付きMF圏を \(\mathcal K_B\)、次数シフトを \(\rho_B\) とする。
 次の三条件を仮定する。
 \begin{enumerate}
 \item \(\mathcal K_B\) はproperである。
 \item Serre関手が存在し、\(S_{\mathcal K_B}\simeq\rho_B^{3-N}[N-2]\) である。
 \item \(\HH^0(\mathcal K_B)\) は非自明な冪等元を持たない。
 \end{enumerate}
 このとき、平行移動の有限位数性を仮定せずに、
 \[
 (\mathcal K_B,\rho_B)\simeq
 \left(\Ku\!\left(V\!\left(\sum_j
 (x_j^3+y_j^3+z_j^3+h_j^2x_jy_jz_j)+P(u)\right)\right),\rho_Y\right)
 \]
 が成り立つ。
\end{corollary}
\begin{proof}
 各楕円ブロックのdegree shiftを保つ同値、外部積の完全忠実性、
 および有限群による通常次数への降下から、
 表示した滑らかな可換cubic \(Y\) について
 \(\Phi:\Ku(Y)\to\mathcal K_B\) が得られ、\(\Phi\rho_Y\simeq\rho_B\Phi\) である。
 中心三次式 \(P\) は一つの可換因子として扱えばよく、
 可換側のequivariant Thom--Sebastianiと非可換側の外部積の完全忠実性は
 Fermat型への分解を必要としない。
 \(Y\) の滑らかさは各変数群の勾配の孤立性から従う。
 \(\Ku(Y)\) は非零のsmooth proper圏であり、
 \eqref{eq:serre-rotation-necessary}と第2条件はSerre関手との互換性を与える。
 命題\ref{prop:serre-connected-generation}を適用すればよい。
\end{proof}

\begin{remark}[十分条件の適用範囲]
 この系は、二つ以上の無限位数ブロックの全生成性を無条件に証明したものではない。
 その場合に残る課題を、三つの明示した圏論的条件へ帰着する結果である。
 与えられた量子化族の全ファイバーについて、properness、Serre公式、
 および \(\HH^0\) の連結性を確認することは依然必要である。
 前節の反例が示すように、一般のAS-regular中心商に対しては、
 比較する完全忠実関手の存在も自動的には得られない。
\end{remark}

\section{本質的混合項の量子中心と相対Picardの制約}
\label{sec:mixed-quantum-center}
分離型を越える量子化では中心変数一次の補正が最初の課題となる。
以下の中心元表示はEtingof--Ginzburg \cite[\S9]{EG} のRains型計算および
本研究の\href{https://github.com/uedakazushi/nc-cubic2/blob/main/papers/nc_cubic_threefold_moduli.tex}{現行の中心三次対モジュライ稿}
の中心式の節にある結果のFermat規格化であり、新しい中心公式とは主張しない。
従来の独立計算は有限体での次数四イデアル所属検査に留まったが、
現在は論文Aの
\texttt{papers/support/papers/sections/fermat\_mixed\_universal\_center.tex}
で全係数の普遍的な中心性を厳密に証明した。

\subsection{全15混合方向の中心三次元}
$f=(x^3+y^3+z^3)/3$、$U=\langle u,v\rangle$とし、
$A,B,C\in U$、$D,E,H\in\Sym^2U$、$g\in\Sym^3U$を任意に取る。
\begin{equation}\label{eq:full-mixed-fermat}
 F=f+Axy+Byz+Czx+Dx+Ey+Hz+g
\end{equation}
とおき、$u,v$を中心変数として
\begin{equation}\label{eq:full-mixed-brackets}
 [x,y]=hF_z,\quad[y,z]=hF_x,\quad[z,x]=hF_y
\end{equation}
の右辺を完全対称化して読む。$T=xyz+xzy+yxz+yzx+zxy+zyx$、
$\operatorname{sym}(xy)=(xy+yx)/2$とする。

\begin{remark}[既存原稿の全混合中心元と検算状態]\label{rem:full-mixed-center}
既存の技術稿の中心元の表示は、現在の規約で
\begin{align}\label{eq:full-mixed-center}
\Omega_h={}&\frac{x^3+y^3+z^3}{3}-\frac{h^2}{9}T\notag\\
&+A\left(\operatorname{sym}(xy)+\frac{h^4}{18}z^2\right)
 +B\left(\operatorname{sym}(yz)+\frac{h^4}{18}x^2\right)
 +C\left(\operatorname{sym}(zx)+\frac{h^4}{18}y^2\right)\notag\\
&+\left(D+\frac{h^2B^2}{6}-\frac{h^4AC}{18}\right)x
 +\left(E+\frac{h^2C^2}{6}-\frac{h^4AB}{18}\right)y\notag\\
&+\left(H+\frac{h^2A^2}{6}-\frac{h^4BC}{18}\right)z+g(u,v)
\end{align}
となり、$h=0$で$F$に特殊化する。
この有限多項式は現在、論文Aで普遍的な中心元として証明済みである。
\texttt{computations/verify\_universal\_fermat\_mixed\_center.py}
では$A,B,C,D,E,H$を独立な中心不定元とし、
$\hbar$も変数のまま$\mathbb Q(\hbar)$上の次数四関係式行列を厳密に還元する。
行列の列数339、関係式の倍数162、一般階数と$\hbar=0$での階数はいずれも144、
全三交換子の剰余は零である。これにより原点を含む開集合での厳密な中心性が従う。
新たな \texttt{computations/verify\_mixed\_center.py} では
$p=1009$、$h=0,1,2,5,13$、各$h$に3組ずつ混合係数を代入して、
次数四の自由語625個における二次関係式の生成部分空間
（階数555）に全ての交換子が属することを確認した。
以下の有限特殊化は追加の検算であり、現在の普遍的証明の根拠ではない。
$A\ne0$だけの模型で$h^2A^2z/6$を省くと
$h=1$の次数四検査で中心性が失われることも確認した。
\end{remark}

\subsection{一次中心持上げの一意性}
\begin{proposition}\label{prop:mixed-center-jet-unique}
$S_h=\CC\langle x,y,z\rangle/
([x,y]-hz^2,[y,z]-hx^2,[z,x]-hy^2)$とし、
$Z(S_h)_2=0$を仮定する。
$R=\CC[u,v]/(u,v)^2$上の\eqref{eq:full-mixed-brackets}の
量子化族$\mathcal B_R$が次数ごとに$R$-自由なら、
固定された$S_h$の次数三の中心元$\Omega_{h,0}$への
次数三の中心持上げ$\Omega_R\in Z(\mathcal B_R)_3$は、
存在すれば一意である。
\end{proposition}
\begin{proof}
二つの持上げの差$\Delta$は$(u,v)(\mathcal B_R)_3$に属し、
次数別自由性と平方零条件により
$\Delta=uQ+vR'$、$Q,R'\in(S_h)_2$と一意に書ける。
$\Delta$の中心性を各生成元$a=x,y,z$との交換子に適用すると
$u[Q,a]+v[R',a]=0$を得る。
$u,v$が独立な平方零係数なので、
$Q,R'\in Z(S_h)_2=0$。従って$\Delta=0$。
\end{proof}
\begin{remark}[一次混合6方向の明示的持上げ]
\eqref{eq:full-mixed-center}を$(u,v)^2=0$に制限すると
\[
\begin{aligned}
 \Omega_h^{(1)}
 ={}&\frac{x^3+y^3+z^3}{3}-\frac{h^2}{9}T
 +A\left(\operatorname{sym}(xy)+\frac{h^4}{18}z^2\right)\\
 &+B\left(\operatorname{sym}(yz)+\frac{h^4}{18}x^2\right)
 +C\left(\operatorname{sym}(zx)+\frac{h^4}{18}y^2\right).
\end{aligned}
\]
これが全ての一次混合6方向の持上げの表示である。
命題\ref{prop:mixed-center-jet-unique}の一意性自体は、
既存の中心公式を仮定しない独立した証明である。
$h=0$では$Z(S_0)_2\ne0$なので適用しない。
\end{remark}

\subsection{相対Picardに対するスカラー方向の障害}
\begin{proposition}\label{prop:relative-scalar-picard}
$X\subset\PP^4$を滑らかなcubic threefoldとし、
$T_0\subset\operatorname{Gr}(2,H^0(\PP^4,\OO(1)))$を、
指定する平面$\Lambda_U=V(U)$との切断$C_U=X\cap\Lambda_U$が
滑らかな楕円曲線になる非空開部分とする。
$\mathscr C\to T_0$を楕円曲線族、
$\mathscr P=\operatorname{Pic}^0(\mathscr C/T_0)$を相対Picardとする。
非零Poisson錐の$\mathbb G_m$-束
$\mathscr L^\times\to T_0$から$\mathscr P$への
$T_0$-上の任意の有理写像は、稠密開部分において
$\mathscr L^\times\to T_0$を通じて有理的に因子化する。
すなわちPoissonスカラー方向には定数である。
幾何学的一般ファイバーで連結な有限\emph{étale}被覆を許しても同様である。
\end{proposition}
\begin{proof}
$K=\CC(T_0)$とする。$\mathscr L^\times$の一般ファイバーは
$\mathbb G_{m,K}$である。
$\overline K$へ移すと、与えられた有理写像は
$\mathbb G_m\dashrightarrow E$、$E$は種数1曲線となる。
固有性によって$\PP^1\to E$に延長されるが、
Riemann--Hurwitz公式から定数である。
従ってもとの写像は一般ファイバーに沿って定数であり、
$T_0$上の有理sectionを通じて因子化する。
また代数閉体上の$\mathbb G_m$の連結有限étale被覆は
$t\mapsto t^n$型で、その源も$\mathbb G_m$である。
同じ議論が適用できる。
\end{proof}
\begin{remark}[復元への含意と限界]
元の$X$とPoisson pencil $U$を幾何側でも固定したまま、
非零スカラーを楕円Picardの非自明な点へ送る大域有理模型は存在しない。
しかし量子化後に別の可換cubic $Y$や平面が動く比較、
分岐被覆を介した多価対応、形式局所同値の不存在は結論しない。
一般の非零パラメータでのrotation付き可換復元に必要な比較は未完成である。
有限Artin底での線束生成残余pairと相対MFの比較は、
論文BのArtin残余比較の節で扱う。
\end{remark}


\part{中心三次対の形式MF比較と混合jet}
\section{一般次元の局所代数的模型：仮定と主定理}
\label{lp:sec:main}
本節では、既に構成された次数別flatな量子化から、射影化のproperなdg模型を作る。
Poisson成分のブロック数、純ブロックの楕円性、有限中心性は仮定しない。
ただし、量子化代数そのものの全次数flatnessは入力である。
基底は完備化やArtin切断ではなく、可換点における実在の正則局所環とする。

\begin{definition}[局所入力]\label{lp:def:input}
$(R,\mathfrak m,k)$を有限生成複素代数の局所化である正則局所環とし、
$d=\dim R$、$N\ge5$とする。
$A=\bigoplus_{j\ge0}A_j$は中央$R$-代数、$A_0=R$で、各$A_j$は有限自由$R$-加群、
$A\otimes_R k\simeq k[z_1,\ldots,z_N]$を次数付き代数として満たすとする。
$\Omega\in Z(A)_3$が与えられ、その剰余$F$は滑らかな
$cubic\ X=V(F)\subset\mathbb P^{N-1}_k$を定めるとする。
$B=A/(\Omega)$、$n=N-2$、$a=N-3$と置く。
内部shiftは$M(j)_q=M_{q+j}$、複体はcohomological conventionを用いる。
\end{definition}

\begin{theorem}[局所proper模型と有界Orlov分解]\label{lp:thm:main}
定義\ref{lp:def:input}の下で以下が成立する。
\begin{enumerate}[label=(\arabic*)]
\item $A,B$は左右noetherianである。$A$の$R$上の対角には長さ$N$の有限自由双加群分解があり、
$\operatorname{gldim}A\le N+d$である。$\Omega$は中心正則で、各$B_j$は有限自由、
その形成は任意の通常の底変更と両立する。
\item 左右で
\[
 \operatorname{R\underline{Hom}}_A(R,A)\simeq R(N)[-N],\qquad
 \operatorname{R\underline{Hom}}_B(R,B)\simeq R(a)[-(N-1)].
 \tag{G}\label{lp:eq:G}
\]
$B$は有限左右injective dimensionを持ち、Burke--Stevensonの意味でGorensteinである。
従って通常の有界三角圏について、$D^b(\operatorname{qgr}B)$は
$D_{\rm sg}^b(\operatorname{gr}B)$と$a$個の$\operatorname{Perf}(R)$成分による
Orlov型半直交分解を持つ。特異圏では自然な三次周期性$(3)\simeq[2]$が成立する。
\item $\mathcal O(j)=\pi B(j)$とすると
\[
 \operatorname{Ext}^p_{\operatorname{qgr}B}(\mathcal O(i),\mathcal O(j))\simeq
 \begin{cases}
 B_{j-i}&p=0,\\
 B_{i-j-a}^{\vee}&p=n,\\
 0&\text{その他}
 \end{cases}\qquad (B_q=0\text{ for }q<0).
 \tag{H}\label{lp:eq:H}
\]
ここで$(-)^\vee=\operatorname{Hom}_R(-,R)$である。
この計算の自然な導来Homの底変更写像は、命題\ref{lp:prop:basechange}の通常の底変更で擬同型となる。
\item $\mathcal C_B=\operatorname{thick}_R\{\mathcal O(j):j\in\mathbb Z\}$を
$D^b(\operatorname{qgr}B)$の標準dg enhancement内の冪等完備な部分と定めると、
\[
 E=\bigoplus_{i=0}^{N-1}\mathcal O(i),\qquad
 \mathcal C_B=\operatorname{thick}_R(E)\simeq\operatorname{Perf}(\Lambda_R),
 \quad\Lambda_R=\operatorname{REnd}(E)\in\operatorname{Perf}(R).
 \tag{P}\label{lp:eq:proper-model}
\]
\item $\mathcal K_B^{\rm lin}=\mathcal C_B\cap
\langle\mathcal O,\ldots,\mathcal O(a-1)\rangle^\perp$は
有限個のmutationにより生成対象$G$を持ち、
$\mathcal K_B^{\rm lin}\simeq\operatorname{Perf}(\Gamma_R)$、
$\Gamma_R=\operatorname{REnd}(G)\in\operatorname{Perf}(R)$である。
特殊ファイバーでは$\mathcal C_B\otimes_R^{\mathbf L}k\simeq\operatorname{Perf}(X)$、
$\mathcal K_B^{\rm lin}\otimes_R^{\mathbf L}k\simeq\operatorname{Ku}(X)$を得る。
\end{enumerate}
\end{theorem}

\begin{remark}[主張しないこと]\label{lp:rem:limits}
$\Gamma_R$の底上proper性と$\Gamma_R^e$上の対角perfect性は異なる。
本定理は後者、一般の可換cubic模型の存在、有限型開集合へのdg構造のspread、
任意の導来基底上のOrlov同値を結論しない。
$\mathcal C_B=D^b(\operatorname{qgr}B)$や
$\mathcal K_B^{\rm lin}=D_{\rm sg}^b(\operatorname{gr}B)$も一般には主張しない。
上付き${\rm lin}$は、線束と有限操作から作る部分圏を全coherent残余圏と区別するための記号である。
本稿の結論は条件付きgeneric復元の入力のうち、proper模型・有界Orlov分解を具体化するものである。
\end{remark}

\subsection{先行研究と本節の寄与}
非可換射影圏のsection functorはArtin--Zhangの枠組み\cite{LP-AZ}に属する。
次数零部が体でないGorenstein環の有界Orlov分解には
Burke--Stevenson\cite[Definition 5.1, Lemmas 6.1--6.2, Theorem 6.4]{LP-BS}を用いる。
dg生成対象によるMorita表示はKeller\cite[Theorem 3.8]{LP-Keller}の一般理論である。
これらの定理自体を新しい結果とは扱わない。
以下では、与えられた量子化に対してnoetherian性・Gorenstein型・Homの底変更・有限生成を導き、
既存原稿で別途仮定していた相対proper模型の条件を正則局所基底で検証する。
この組合せの新規性は未確認であり、外部査読や形式証明は未実施である。

\section{局所基底でのnoetherian性と有限Koszul持上げ}
\label{lp:sec:lift}
\begin{lemma}[次数別有限性を用いたnoetherian性]\label{lp:lem:noetherian}
$(R,\mathfrak m)$をnoetherian局所環、$D=\bigoplus_{j\ge0}D_j$を中央$R$-代数とし、
各$D_j$は有限$R$-加群とする。
$D/\mathfrak mD$が右noetherianなら$D$は右noetherianである。左にも同様である。
\end{lemma}
\begin{proof}
$\mathfrak m$の有限生成元を選ぶと、
$(D/\mathfrak mD)[t_1,\ldots,t_e]\twoheadrightarrow\operatorname{gr}_{\mathfrak m}D$
という中央変数の多項式拡大からの全射がある。
従って右辺は右noetherianである。
斉次右イデアル$I\subset D$に誘導filtration $I\cap\mathfrak m^qD$を入れる。
その随伴次数イデアルの有限個の二重斉次生成元を、$f_1,\ldots,f_r\in I$へ持ち上げ、
$J=\sum f_iD$とする。
$x\in I_j$の最低$\mathfrak m$-次数記号を消す操作を繰り返すと、任意の$q$について
$x\in J_j+\mathfrak m^qD_j$を得る。
$D_j/J_j$は有限$R$-加群だからKrull交叉定理より$x\in J_j$である。
従って全ての斉次右イデアルは有限生成である。
任意の右イデアルには内部次数の上昇filtrationを入れ、最高次数項のイデアルの
有限生成元を持ち上げる。内部次数の有限降下により元のイデアルも有限生成である。
この証明は$R$の完備性を使わず、各内部次数における有限性を使う。
\end{proof}

\begin{lemma}[次数別分裂する自由分解の持上げ]\label{lp:lem:resolution}
$D$の各次数が有限自由$R$-加群で、$M$も各次数有限自由とする。
$\overline M$が$\overline D=D/\mathfrak mD$上、有限階数の次数付き自由加群による
有限分解$\overline P_\bullet$を持つなら、同じ項の階数とshiftを持つ
$D$-自由分解$P_\bullet\to M$を持ち上げられる。
分解は各内部次数で$R$-加群の分裂完全列であり、任意の通常の底変更後も完全である。
\end{lemma}
\begin{proof}
$\overline P_0$の生成元の像を$M$の同じ次数へ持ち上げる。
各次数の余核は有限$R$-加群で剰余が零なのでNakayamaの補題で零である。
$P_0\twoheadrightarrow M$は各次数で有限自由加群への全射だから分裂し、
核は各次数有限自由となる。剰余体への底変更後の核も正しく計算される。
核の生成元を$\overline P_1$から同様に持ち上げ、帰納する。
最後の核は各次数有限自由かつ剰余が零であるから零である。
各段階の内部次数別分裂により、底変更後にも完全性が残る。
\end{proof}

\begin{proposition}[全次数の相対Koszul性とambientの対角]\label{lp:prop:koszul}
定義\ref{lp:def:input}の$A$は次数一で生成され、二次関係で定まる。
$R$の左右$A$-加群としての分解と、$A$の$A^e=A\otimes_RA^{\rm op}$-加群としての分解は
それぞれ
\[
 P_j=A(-j)^{\binom Nj},\qquad
 Q_j=A^e(-j)^{\binom Nj}\quad(0\le j\le N)
 \tag{K}\label{lp:eq:koszul}
\]
という項を持ち、全ての通常の底変更と両立する。
特に$A$は$R$上smoothで、$\operatorname{gldim}A\le N+d$である。
\end{proposition}
\begin{proof}
$A_1$の剰余は$N$変数の空間であり、そのliftによる自由代数からの写像は
各次数のNakayamaにより全射となる。
多項式環の$k$のKoszul分解に補題\ref{lp:lem:resolution}を適用すると$P_\bullet$を得る。
第一次関係が次数二で生成されることもこの構成から従う。
同様に$A$を$A^e$-加群として見て、多項式環の対角Koszul分解を持ち上げると$Q_\bullet$を得る。
$A$は$R$上の自由加群（次数別の直和）なので$Q_\bullet$は右$A$-加群としても分裂完全である。
任意の左$A$-加群$M$を右からテンソルすると、項$A\otimes_R M$の有限直和からなる長さ$N$の分解になる。
$\operatorname{gldim}R=d$より各項の$A$-射影次元は高々$d$。
従って$\operatorname{pd}_A M\le N+d$で、右側も同様である。
\end{proof}

\begin{proposition}[相対Gorenstein型と中心商]\label{lp:prop:gorenstein}
式\eqref{lp:eq:G}が成立する。$B$は左右noetherian、各次数有限自由で、
左右injective dimensionは高々$N+d-1$である。
\end{proposition}
\begin{proof}
補題\ref{lp:lem:noetherian}を$A$に適用する。
$P_\bullet$を$A$で双対化した有界複体は、各内部次数で有限自由$R$-複体である。
剰余上ではKoszul双対性により、内部次数$-N$・cohomological degree $N$以外で完全。
有限自由有界複体の剰余が完全なら、単元行列成分の消去とNakayamaにより複体自身も分裂完全である。
内部次数$-N$では最後の項$R$のみ残る。従って$\operatorname{R\underline{Hom}}_A(R,A)=R(N)[-N]$を得る。
左側も同じである。

各次数での$\Omega:A_{j-3}\to A_j$の剰余は非零多項式$F$の乗法で単射。
行列の列数と同じ大きさの小行列式の一つが単元なので、この写像は局所基底変換で分裂単射となる。
よって$\Omega$は正則、$B_j$は自由で、全底変更と両立する。
中央正則元に対するchange-of-ringsは
\[
 \operatorname{R\underline{Hom}}_A(M,A)
 \simeq\operatorname{R\underline{Hom}}_B(M,B)(3)[-1]
 \quad(M\in D(\operatorname{Gr}B))
 \tag{CR}\label{lp:eq:change-rings}
\]
を与える。これは$0\to A(-3)\xrightarrow{\Omega}A\to B\to0$から
$\operatorname{R\underline{Hom}}_A(B,A)=B(3)[-1]$を計算し、導来随伴を使った式である。
$M=R$を代入して$B$の式\eqref{lp:eq:G}を得る。
一般の$M$に適用すると、$A$の有限大域次元から$B$の有限injective dimensionが従う。
商なのでnoetherian性も成り立つ。
\end{proof}

\section{射影化の局所コホモロジーと底変更}
\label{lp:sec:cohomology}
$\tau$を$B_{\ge1}$-torsion functor、$\pi:\operatorname{Gr}B\to\operatorname{QGr}B$を商関手とする。
Burke--Stevenson\cite[Section 4, Lemmas 6.1--6.2]{LP-BS}により、
上記のGorenstein型と$R$-flatnessは$\chi$と必要な有限cohomological dimensionを与える。
従って$D^b(\operatorname{qgr}B)$を$D(\operatorname{QGr}B)$で計算でき、導来sectionの三角
\[
 R\tau(B)\longrightarrow B\longrightarrow R\omega\pi(B)\longrightarrow
 \tag{LC}\label{lp:eq:torsion}
\]
を使用できる。

\begin{lemma}[局所コホモロジーの各次数]\label{lp:lem:local-cohomology}
$d_B=N-1$とすると
\[
 H^q(R\tau B)_t=
 \begin{cases}
  B_{-t-a}^{\vee}&q=d_B,\\
  0&q\ne d_B
 \end{cases}
 \tag{LC'}\label{lp:eq:localcohom}
\]
である。負添字は零とする。
\end{lemma}
\begin{proof}
$M_\ell=B/B_{\ge\ell}$の次数filtrationの商は
$B_j\otimes_R R(-j)$（$0\le j<\ell$）である。
式\eqref{lp:eq:G}と$B_j$の自由性により、各商の$B$への導来Homは次数$d_B$に集中し、
その内部次数は$-a-j$のみである。
長完全列によって$\operatorname{Ext}_B^q(M_\ell,B)=0$（$q\ne d_B$）。
内部次数$t$では$j=-t-a$のただ一層だけが寄与するので、その層が含まれれば群は$B_j^\vee$となる。
$\ell>j$で推移写像は同型である。
$R\tau B=\varinjlim_\ell\operatorname{R\underline{Hom}}_B(M_\ell,B)$のcohomologyを取ればよい。
\end{proof}

\begin{proof}[式\eqref{lp:eq:H}の証明]
$\operatorname{RHom}_{\operatorname{QGr}B}(\pi B,\pi B(t))$
は導来sectionの内部次数$t$に一致する。
式\eqref{lp:eq:torsion}と補題\ref{lp:lem:local-cohomology}を用いると
$H^0=B_t$、$H^{N-2}=B_{-t-a}^\vee$で、それ以外は零となる。
$a>0$なので同じ$t$に対して二群が同時に非零となることはない。
従ってこれはcohomology次元だけでなく、各Hom複体が有限自由加群の一つのshiftであることを示す。
$i$のtwistを外せば式\eqref{lp:eq:H}である。
\end{proof}

\begin{proposition}[線束Homの底変更]\label{lp:prop:basechange}
$R\to R'$を有限大域次元のnoetherian正則環への通常の準同型で、$A'=A\otimes_RR'$が左右noetherianであるものとする。
$B'=B\otimes_RR'$と置けば自然な写像
\[
 \operatorname{RHom}_{\operatorname{qgr}B}(\mathcal O(i),\mathcal O(j))
 \otimes_R^{\mathbf L}R'
 \longrightarrow
 \operatorname{RHom}_{\operatorname{qgr}B'}(\mathcal O'(i),\mathcal O'(j))
 \tag{BC}\label{lp:eq:bc}
\]
は擬同型である。
特に任意の剰余体$R\to\kappa(\mathfrak p)$、および正則noetherian局所環への局所準同型には適用できる。
\end{proposition}
\begin{proof}
$A,B$の有限相対分解と$B$の正則中心商表示は底変更と両立するので、
$A',B'$も対応するGorenstein型、有限injective dimension、$\chi$を持つ。
剰余体については中央の商・局所化によりnoetherian性を得る。
局所準同型の場合は剰余が多項式環なので補題\ref{lp:lem:noetherian}を再適用できる。

まず$R$の有限大域次元により$-\otimes_R^{\mathbf L}R'$は有限Tor振幅を持ち、
有界coherent graded複体を有界coherent複体へ送り、torsionを保つ。
従って射影商の有界導来圏にも底変更が定義される。
自然な底変更写像そのものを確認するため、$M_\ell$を有限階数graded自由$B$-加群の分解$P_\bullet$で計算する。
$M_\ell$は内部次数別$R$-自由なので、そのsyzygyを各次数で分裂させて、分解の全底変更を保てる。
固定した内部次数$t$で$C^\bullet=\operatorname{\underline{Hom}}_B(P_\bullet,B)_t$は
下に有界な有限自由$R$-加群の複体であり、前補題の証明により、
cohomologyは有限自由加群の一つの次数にしか現れない。
$R$の大域次元は有限なので、cohomologyが全て射影なこの複体は、そのcohomology複体にhomotopy同値である。
実際、射影加群からなる完全複体のcycleについてdimension shiftingを$d+1$回適用すると、
全cycleが射影となり複体が分裂する。
従って$C^\bullet$の導来テンソルは通常のテンソルで計算できる。
有限階数自由分解上ではHomと底変更が項ごとに可換であるため、$M_\ell$のHomの底変更を得る。
各$t$で安定するfiltered colimitと式\eqref{lp:eq:torsion}を通せば\eqref{lp:eq:bc}となる。
\end{proof}

\section{有限生成proper部分圏と残余圏}
\label{lp:sec:proper}
\begin{proof}[有界Orlov分解と有限生成の証明]
命題\ref{lp:prop:gorenstein}の結果はBurke--StevensonのDefinition 5.1の仮定を満たす。
$R$は有限大域次元で、$A,B$は中央$R$-flatなので、同論文Lemma 6.2によって$\chi$を満たす。
Theorem 6.4を$A$へ適用すると、$D_{\rm sg}^b(\operatorname{gr}A)=0$であるため
任意の連続した$N$個のtwistが$D^b(\operatorname{qgr}A)$を生成する。
$B$への導来制限は$0\to A(-3)\to A\to B\to0$により有限Tor振幅を持ち、
torsionをtorsionに送るので射影商へ降りる。
この関手は$\mathcal O_A(j)$を$\mathcal O_B(j)$へ送る。
従って全ての$\mathcal O_B(j)$は$E$から有限回のcone・直和因子で作られる。
これが式\eqref{lp:eq:proper-model}の生成性である。

$B$へのTheorem 6.4は、$a>0$より$a$個のtwistとgraded singularity categoryによる半直交分解を与える。
この有界分解の向きは引用定理の規約に従う。以下では線束生成部分における右直交を別に構成する。
この有界三角圏の同値を、任意の導来基底変更と両立するdg同値へ強めたとは主張しない。
\end{proof}

\begin{proposition}[properな残余dg代数]\label{lp:prop:residual}
$\mathcal E=\langle\mathcal O,\ldots,\mathcal O(a-1)\rangle$とする。
各成分は$\operatorname{Perf}(R)$であり、$\mathcal C_B$で
\[
 \mathcal C_B=\langle\mathcal K_B^{\rm lin},\mathcal O,\ldots,\mathcal O(a-1)\rangle
 \tag{SOD}\label{lp:eq:sod}
\]
を持つ。$p=L_{\mathcal O}\circ\cdots\circ L_{\mathcal O(a-1)}$と置くと
$G=p(E)$は$\mathcal K_B^{\rm lin}$を生成する。
$\Gamma_R=\operatorname{REnd}(G)$の底上の複体はperfectであり、
命題\ref{lp:prop:basechange}の底変更でこの構成はMorita同値まで可換である。
\end{proposition}
\begin{proof}
式\eqref{lp:eq:H}から$\operatorname{RHom}(\mathcal O(i),\mathcal O(i))=R$で、
$0\le i<j\le a-1$なら$\operatorname{RHom}(\mathcal O(j),\mathcal O(i))=0$である。
全$\mathcal C_B$の対象間のHomは、$E$間のperfect Homから有限cone・直和因子で作られるので$R$-perfectである。
従って各例外成分の左右随伴はHomとその$R$-双対で定義でき、mutationは有限操作となる。
右端の例外対象から順に左mutationして$p$を作ると、$p(\mathcal C_B)$は正確に$\mathcal E^\perp$である。
$E$の生成性より$G$が生成し、$\operatorname{REnd}(G)$も$R$-perfectとなる。
標準dg enhancementと生成対象のdg Yonedaを用いればMorita表示を得る\cite[Theorem 3.8]{LP-Keller}。

式\eqref{lp:eq:bc}は生成対象間の全導来Homと積に自然なので、有限coneと冪等完備化へ延長すると
線束生成部分圏の底変更で完全忠実性を得る。像は全線束を生成するのでMorita同値である。
評価写像とそのconeも底変更に可換で、$p,G$についても同じである。
特殊ファイバーでは$X$は滑らかで、全twistは$\operatorname{Perf}(X)=D^b(\operatorname{coh}X)$を生成する。
右直交は通常のKuznetsov成分なので主定理(5)が従う。
\end{proof}

\begin{lemma}[局所proper圏の完全忠実性の持上げ]\label{lp:lem:ff-lift}
$\mathcal C,\mathcal D$を$R$-properな冪等完備dg圏とし、対象間のHomが剰余体への導来底変更と可換であるとする。
$\Phi:\mathcal C\to\mathcal D$を底変更と両立するexactなdg関手とする。
$\Phi_k$が完全忠実なら$\Phi$も完全忠実である。
さらに$\Phi_k$が同値で、$\Phi$が右または左随伴を持つなら$\Phi$は同値である。
\end{lemma}
\begin{proof}
任意の対象$M,N$に対し、Hom上の自然な写像のconeは$R$-perfectで、その剰余は零である。
局所環上の導来Nakayamaによりconeは零となるので完全忠実性を得る。
また$Q\in\mathcal D$について$Q_k=0$なら、その自己Homの剰余も零である。
従って$\operatorname{RHom}(Q,Q)=0$、特に恒等射が零なので$Q=0$となる。
$\Phi$の像が片側admissibleであれば、任意の対象は像とその直交成分によるtriangleを持つ。
直交成分の対象$Q$は剰余上で$\Phi_k(\mathcal C_k)$と直交し、$\Phi_k$の同値性から$Q_k=0$。
したがって$Q=0$であり、像は全体である。
\end{proof}

\begin{proposition}[標準rotation候補の完全忠実性]\label{lp:prop:rotation-candidate}
$\widetilde\rho=L_{\mathcal O}\circ(1)$は$\mathcal K_B^{\rm lin}$を保つ完全忠実なexact dg自己関手であり、
特殊ファイバーでは通常のrotationとなる。
上記の底変更と両立する。
本節だけでは全基底上の本質的全射性と$\widetilde\rho^3\simeq[2]$を主張しない。
\end{proposition}
\begin{proof}
$M\in\mathcal E^\perp$なら
$\operatorname{RHom}(\mathcal O(j),M(1))=0$（$1\le j\le a$）。
$L_{\mathcal O}$は$\mathcal O$からのHomを零にする。
$1\le j<a$について$\operatorname{RHom}(\mathcal O(j),\mathcal O)=0$だから、
このmutationは他の必要な直交条件を壊さない。
twistは線束生成部分圏を保ち、mutationも有限操作なので自己関手を得る。
基底変更は前命題の証明と同じである。
特殊ファイバーでは滑らかなcubicの標準rotationであり自己同値なので、
補題\ref{lp:lem:ff-lift}によって$R$上で完全忠実になる。
\end{proof}
\begin{remark}[自己同値化を残す理由]
補題\ref{lp:lem:ff-lift}により、$\widetilde\rho$に有限対象を保つ片側随伴を構成すれば、
本質的全射性も従う。しかしproper性だけから任意のexact関手が随伴を持つとは限らない。
本稿ではこの随伴の存在を仮定なしに付け加えない。
同様に、$\widetilde\rho^3\to[2]$という自然変換を実際に構成し、特殊ファイバーで同型と示せれば、
対象ごとのconeに導来Nakayamaを適用して全体で同型にできる。
自然変換の持上げそのものを、その特殊ファイバーでの存在で代用してはならない。
\end{remark}

\subsection{中心商の導来自己交叉と特異圏の周期性}
\begin{proposition}[三次周期性の代数的な構成]\label{lp:prop:sing-period}
本節の仮定で、graded singularity categoryには自然な同型
\[
 (3)\simeq[2]
 \qquad\text{on }D_{\rm sg}^b(\operatorname{gr}B)
 \tag{PER}\label{lp:eq:period}
\]
が存在する。ここでは射影残余圏のsmoothnessを使わない。
\end{proposition}
\begin{proof}
$Q=B\otimes_A^{\mathbf L}B$をgraded $B$-双加群とする。
中央正則元の長さ一分解から、$H^0(Q)=B$、
$H^{-1}(Q)=(\Omega)/(\Omega)^2\simeq B(-3)$で、他のcohomologyは零である。
この同定は$B$-双加群として成立する。導来双加群の標準truncationは
\[
 B(-3)[1]\longrightarrow Q\longrightarrow B
 \xrightarrow{\eta}B(-3)[2]
\]
を与える。
任意の$M\in D^b(\operatorname{gr}B)$に$M\otimes_B^{\mathbf L}-$を作用させると
\[
 M(-3)[1]\longrightarrow M\otimes_A^{\mathbf L}B\longrightarrow M
 \xrightarrow{\eta_M}M(-3)[2]
\]
となる。$A$はnoetherianで有限大域次元なので$M$は$A$上perfectであり、
中項は$B$上perfectである。従って$\eta_M$はgraded singularity categoryで同型となる。
$\eta$が双加群射から来ることにより$M$に自然である。
内部shiftを三つ戻すと式\eqref{lp:eq:period}を得る。
\end{proof}
\begin{remark}[二種類のrotationを混同しない]
有界Orlov同値で内部shift $(1)$を移せば、全coherent残余圏には
三乗が$[2]$となる自己同値$\rho_{\rm sg}$を定義できる。
他方、命題\ref{lp:prop:rotation-candidate}は線束生成部分に
$\widetilde\rho=L_{\mathcal O}\circ(1)$を具体的に構成する。
この二つの関手を指定したwindowで同定し、全体での底変更とcoherenceを確認する作業は別である。
本稿は通常の中心商から得る特異圏の周期性そのものを未解決とはしないが、
上の同定を省略して線束生成残余部分の相対rotationが完成したとは主張しない。
また$Q$を$B\oplus B(-3)[1]$と双加群として分裂させてはならない。
その拡張類$\eta$こそが周期性を与えている。
\end{remark}

\section{非楕円成分への適用と有限計算}
\label{lp:sec:application}
前稿で構成した二パラメータ非楕円量子化の係数環は
$R_{\mathrm C}^0=\CC[r,s]$、$R_{\mathrm N}^0=\CC[r,s,(4-r^2)^{-1}]$である。
可換点$(r,s)=(0,0)$で局所化すると、どちらも正則局所環
$R=\CC[r,s]_{(r,s)}$になる。
各次数の自由性（局所化前は局所自由性）、七生成元の多項式環への特殊化、
中心三次元は前稿の全次数PBW構成から得られる。
従って主定理は$N=7,n=5,a=4$で適用できる。
この適用は節点・尖点の純ブロックの滑らかさを仮定しない。
既存の中心元・PBW性の証明と再現コードは本bundleの別論文に保持している。

\begin{corollary}[二つの成分の局所proper残余模型]\label{lp:cor:nonelliptic}
上の二つの成分に対し、$R=\CC[r,s]_{(r,s)}$上の実在するproperな残余dg代数
$\Gamma_R$と完全忠実な標準rotation候補を構成できる。
その特殊ファイバーは前稿の滑らかな$cubic\ fivefold$のKuznetsov成分とrotationである。
$B$は有界Orlov分解を持つ。
ただしこの残余dg代数の相対smoothnessは未証明である。
\end{corollary}

\begin{proposition}[線束生成子のcohomology階数]\label{lp:prop:ranks}
$b_q=\binom{N-1+q}{N-1}-\binom{N-4+q}{N-1}$（負次数の項は零）とする。
$\Lambda_R=\operatorname{REnd}(\bigoplus_{i=0}^{N-1}\mathcal O(i))$について
\begin{align*}
 \rank_R H^0(\Lambda_R)&=\sum_{q=0}^{N-1}(N-q)b_q
 =\binom{2N}{N-1}-\binom{2N-3}{N-4},\\
 \rank_R H^{N-2}(\Lambda_R)&=3+2N+\binom{N+1}{2}.
\end{align*}
それ以外のcohomologyは零である。$N=7$では階数は$2838$と$45$である。
\end{proposition}
\begin{proof}
第0次は差$j-i=q\ge0$を持つ組が$N-q$個あることから得る。
二つのbinomial和にhockey-stick公式を二回使うと閉公式になる。
上の次数は$i-j\ge a=N-3$だけが寄与し、差$a,a+1,a+2$を持つ組はそれぞれ3,2,1個である。
$B_0=R$、$\rank B_1=N$、$\rank B_2=\binom{N+1}{2}$を代入する。
\end{proof}

\begin{remark}[cohomologyの表はdg積の代用品ではない]\label{lp:rem:ainfty}
上の数値から$\Lambda_R$や$\Gamma_R$のdg積・高次積を復元したとは主張しない。
特に$H^0(\Lambda_R)\oplus H^{N-2}(\Lambda_R)[-(N-2)]$を
自動的に同じdg代数のsquare-zero模型とみなしてはならない。
計算コードはbinomial恒等式、Homの三角的消滅、例外成分の数を整数で確認するだけで、
noetherian性・Orlov定理・dg Morita同値を形式化したものではない。
\end{remark}

\section{残余smoothnessに残る正確な問題}
\label{lp:sec:remaining}
構成した$\Gamma_R\in\operatorname{Perf}(R)$について、なお必要なのは
\[
 \Gamma_R\in\operatorname{Perf}(\Gamma_R^{\rm op}\otimes_R^{\mathbf L}\Gamma_R).
 \tag{SM}\label{lp:eq:smoothness}
\]
周囲$A$の対角の有限分解は主定理で得たが、式\eqref{lp:eq:smoothness}ではない。
$\Gamma_R\otimes_Rk$がsmoothであることだけから、代数的な有限対角分解の持上げを結論しない。
実際のperfect複体$P$と$P\to\Gamma_R$を構成して一ファイバーで擬同型とすれば、
coneの底上perfect性からその消滅を開近傍へ広げられる。
この判定は有効だが、$P$と射を構成する仕事を代替しない。

同様に、通常の有界Orlov分解があることから、任意の導来基底でのdg同値、
標準rotationの可逆bimodule表示、三次周期性のcoherenceを無条件に結論してはならない。
本稿ではproperな線束生成部分を先に作り、その自己関手候補までを構成した。
全coherent残余圏の比較、標準rotationの本質的全射性と相対周期性の同定は次の課題である。

\begin{remark}[局所環と有限型開集合を区別する]
$R$は可換点の局所環なので、全次数にわたるNakayamaは同一の$R$上で使える。
しかしこれだけで、元の有限型基底の一つのZariski開集合上に
全ての無限次数条件が広がったとはいえない。
有限の自由分解・有限個の行列の個別のspreadと、proper dg代数の全高次積のspreadも区別する。
実際に代数的なsmooth proper pairを得た後に、既存原稿Cのétale局所実現原理(ER)を適用する。
(ER)の圏論的入力は本稿で再証明していない。
\end{remark}

\begin{remark}[同じHom階数でもsmoothとは限らない]
本稿の線束Homの計算は、非零三次式$F$が特異であっても成り立つ。
例えば$F=z_1^2z_2+z_3^3+\cdots+z_N^3$は$[0:1:0:\cdots:0]$で特異だが、
同じHilbert級数と命題\ref{lp:prop:ranks}の階数を持つ。
可換射影超曲面の$\operatorname{Perf}$はこの場合smoothでない。
実際、smooth properならcohomologically finiteなmoduleはperfectであり、
特異点のskyscraperもperfectとなって、その局所環の正則性に矛盾する。
この標準的な対角判定は\cite[Proposition 3.31]{LP-Orlov}のsmooth/properの枠組みに属する。
従って階数表だけをsaturationの証明書として使うことはできない。
\end{remark}
\section{線束生成子の相対最小$A_\infty$模型と対角smoothness}
\label{sec:minimal-ainfty}
本節は、残余圏の対角perfect性(SM)を実際に検証するため、
線束生成子のdg代数へ問題を移し、その相対最小模型の高次演算を制限する。
相対半直交分解の一般定理は既知であり
\cite[Lemma~4.15(3)]{AinftyPerry2019}を参照する。
最小模型の存在には\cite{PetersenTransfer,BurkeTransfer}を用いる。
以下はその適用と整数的な配置の絞り込みであり、SM自体の証明ではない。

\begin{theorem}[残余対角と線束生成子の対角]
\label{thm:ainfty-smooth-equiv}
定義\ref{lp:def:input}の正則局所底$R$上で、
既存の局所proper模型
$\mathcal C_B=\thick_R(E)\simeq\Perf(\Lambda_R)$、
$\mathcal K_B^{\mathrm{lin}}\simeq\Perf(\Gamma_R)$
を考える。ただし
\[
 E=\bigoplus_{i=0}^{N-1}\mathcal O(i),\quad
 \Lambda_R=\REnd(E),\quad N=n+2,\quad n\ge3.
\]
このとき、相対smoothnessについて
\[
 \boxed{\ \Gamma_R\in\Perf(\Gamma_R^e)
 \quad\Longleftrightarrow\quad
 \Lambda_R\in\Perf(\Lambda_R^e)\ },\qquad
 A^e:=A^{\mathrm{op}}\otimes_R^{\mathbf L}A.
 \tag{SM$'$}\label{eq:ainfty-sm-prime}
\]
従ってSMの有限対角証明書を$\Lambda_R$について求めてもよい。
\end{theorem}
\begin{proof}
既存の有限$R$-例外列により
\[
 \mathcal C_B=
 \langle\mathcal K_B^{\mathrm{lin}},
 \mathcal O,\mathcal O(1),\ldots,\mathcal O(n-2)\rangle
\]
は$R$-線形半直交分解である。各相互Homは$R$-perfectであり、
評価・余評価から定める左右mutationが実在するため、各成分はadmissibleである。
$\mathcal O(i)$の生成する成分は$\Perf(R)$と同値であり、$R$上smoothである。

相対半直交分解に関するPerryの補題
\cite[Lemma~4.15(3)]{AinftyPerry2019}により、
admissibleな各成分がsmoothであることと全体がsmoothであることは同値である。
この場合の具体的な理由も記しておく。
全体の対角を残余成分の左右両変数へ制限すると残余対角となり、
包含の両随伴によりrepresentable双加群はrepresentableへ送られるので、
全体の対角perfect性は残余対角perfect性を含意する。
逆に、残余成分と各例外成分がsmoothなら、各貼合せ双加群は
$R$上perfect（全体がproper）であり、
smoothな両成分の包絡dg圏上でもperfectである。
三角dg代数の対角を与える有限三角を使って帰納的に貼り合わせれば
全体の対角がperfectになる。Morita不変性から式\eqref{eq:ainfty-sm-prime}を得る。
\end{proof}

\begin{proposition}[二段階のsquare-zeroコホモロジー環]
\label{prop:ainfty-squarezero}
$H=H^*(\Lambda_R)$、$U=H^0$、$M=H^n$とおく。
通常の積について$H=U\ltimes M[-n]$かつ$M^2=0$である。
各対象指数において
\[
 U(i,j)=
 \begin{cases}B_{j-i}&j\ge i,\\0&j<i,\end{cases}\qquad
 M(i,j)=
 \begin{cases}B_{i-j-(n-1)}^\vee&i-j\ge n-1,\\
 0&\text{otherwise},\end{cases}
 \tag{SQ}\label{eq:ainfty-sq}
\]
である。$U$の積は$B$の積、$M$への作用は線束の合成と
Serre双対表示から定まる。
\end{proposition}
\begin{proof}
既存の線束Hom公式\eqref{lp:eq:H}により$H$は次数$0,n$だけにある。
$m_2$は次数零なので$H^nH^n\subset H^{2n}=0$。
残る成分の式は\eqref{lp:eq:H}そのものである。
\end{proof}

\begin{proposition}[相対最小模型の高次演算の間隙]
\label{prop:ainfty-arity}
$R$-線形でstrictly unitalな最小$A_\infty$模型を$H$上に取れる。
その積$m_k$と、$f_1=\mathrm{id}_H$である最小模型の比較同型の成分$f_k$は
\[
 m_k=0\quad(k\not\equiv2\!\!\!\pmod n),\qquad
 f_k=0\quad(k\not\equiv1\!\!\!\pmod n).
 \tag{AR}\label{eq:ainfty-arity}
\]
特に$m_3,\ldots,m_{n+1}$は零であり、
最初の可能な高次積は$m_{n+2}$、最初の可能な非線形な
比較同型は$f_{n+1}$である。
\end{proposition}
\begin{proof}
Homの各コホモロジーは有限自由$R$-加群で、
Hom複体は$R$-perfectである。従って射影複体模型を選べば
cycle、boundaryの列を分裂でき、単位を保存する
homotopy transferによって相対最小模型を得る。
$m_k$のcohomological degreeは$2-k$である。
次数$n$の入力を$p$個、出力次数を$q\in\{0,n\}$とすれば
$np+2-k=q$、従って$k\equiv2\pmod n$。
比較同型$f_k$の次数は$1-k$であり、同じ議論から
$k\equiv1\pmod n$。strict unital性によって
単位を含む$k\ge3$の入力では$m_k$が消える。
\end{proof}

\begin{proposition}[最初のHochschild類と次の二次障害]
\label{prop:ainfty-hoch}
$H$を$m_2$による次数付き代数とすると、
$m_{n+2}$は内部次数$-n$のHochschild cocycleを定め、
\[
 [m_{n+2}]\in\HH^{n+2}(H,H)_{-n}
 \tag{HC}\label{eq:ainfty-hh}
\]
は$f_1=\mathrm{id}$のstrictly unitalな最小模型同型に依らない。
次に現れ得る積は$m_{2n+2}$で、suspended cochainの規約では
\[
 \delta m_{2n+2}
   =-\frac12[m_{n+2},m_{n+2}].
 \tag{OB}\label{eq:ainfty-ob}
\]
従って任意の候補$m_{n+2}$について
$\frac12[m_{n+2},m_{n+2}]$の
$\HH^{2n+3}(H,H)_{-2n}$での類の消滅は延長可能性の必要条件である。
実在する$\Lambda_R$の模型ではこの類は消えるが、
$m_{2n+2}=0$とは結論しない。
\end{proposition}
\begin{proof}
$m_3,\ldots,m_{n+1}$が零なので、
長さ$n+3$のStasheff恒等式は$\delta m_{n+2}=0$である。
同型条件の最初の非自明な式では、$m_{n+2}$の差が
$\delta f_{n+1}$となるので類は不変である。
長さ$2n+3$の式は$m_2,m_{2n+2}$の合成と
$m_{n+2}$同士の挿入からなり、表示したGerstenhaber関係を得る。
後者の内部次数は$-2n$である。
\end{proof}

\begin{proposition}[最初の高次積の対象指数制限]
\label{prop:ainfty-indices}
$m_{n+2}$の非零候補は、次数$n$の入力を一個持ち
次数零を出力するもの、または次数$n$の入力を二個持ち
次数$n$を出力するものだけである。
前者が$\operatorname{End}^0(\mathcal O(i))=R$を出力するなら、
指数列は
\[
0\to1\to\cdots\to N-1\to0
\tag{CY}\label{eq:ainfty-cyclic}
\]
の$N$個の巡回移動に限る。
後者では$n$個の正の指数上昇の合計は
$n,n+1,n+2,n+3$のいずれかであり、
各上昇幅は高々$4$である。
\end{proposition}
\begin{proof}
式\eqref{eq:ainfty-sq}より、非単位な次数零の射は指数が
少なくとも一つ上昇し、次数$n$の射は指数が
$n-1,n,n+1$だけ下降する。
$m_{n+2}$の次数条件から高次入力数は一個または二個となる。
出力が同じ指数の場合、$n+1$個の上昇幅の和は
一個の下降幅に等しい。各上昇幅が$\ge1$、
下降幅が$\le n+1$なので全ての上昇幅は一、
下降は$N-1$から$0$に限られる。
二個の下降幅を$d_1,d_2$、出力の下降幅を$d_0$とすると
$n$個の上昇幅の和は$d_1+d_2-d_0$で、
各$d_i\in\{n-1,n,n+1\}$より$n\le u\le n+3$。
一つでも上昇幅が$5$以上なら$u\ge n+4$となり矛盾する。
\end{proof}

\begin{remark}[cubic fivefoldの有限計算の範囲]
$n=5,N=7$なら高次積候補は$m_7,m_{12},m_{17},\ldots$。
$H^0$と$H^5$の全$R$-階数は既存の命題\ref{lp:prop:ranks}で
それぞれ$2838,45$である。
strictly unitalな指数列に限ると、
$m_7$の対象指数配置は出力$H^0$側$2478$通り、
出力$H^5$側$303$通りで、前者の対角指数への出力は$7$通り。
これはHomのベクトル空間基底を数えた値ではない。
計算は\texttt{computations/verify\_minimal\_ainfty\_arity.py}で再現できる。
縮約$(i,p,h)$を実際に選べば、$m_7$は七葉の平面二分木
$C_6=132$個による有限和として計算できる。
しかし高次積の係数・相対bar右逆をまだ得ていない。

$m_7$を二度挿入する許容指数列として
\[
 0\to1\to2\to3\to4\to0\to1\to2\to3
 \to4\to0\to1\to2\to3
\]
がある。二つの$4\to0$が次数$5$で、それ以外は次数零。
最初の七射の$m_7$出力は$0\to2$の次数零に、
その出力を含む外側の$m_7$は$0\to3$の次数零に
入り得る。従って$[m_7,m_7]$は指数条件だけでは消えない。
非零の値やHochschild障害類を計算した主張ではない。
特に、ここから$\Gamma_R\in\Perf(\Gamma_R^e)$を結論してはならない。
\end{remark}

% Current local-window proof. Compatible with the general comparison in the following section.
\section{三つのstabilizationによる局所残余圏の表示}
\label{tw:sec:window}

本節は、主定理で用いる局所比較(L)を、代数的入力から導く。
形式的MF同値(CF)、相対対角のperfect性、特殊ファイバーからの生成性の持上げは使わない。
完備化はHom複体の有限性を証明するためだけに用い、
関手自体は代数的なsection--windowと余核のdg比較から作る。
したがって、両立するArtin同値の列を逆極限にして関手を得たとは主張しない。

$(R,\mathfrak m,k)$を有限Krull次元の正則noetherian局所$\mathbb C$-代数、
$k=\mathbb C$とし、$\widehat R$をその完備化とする。
$A=\bigoplus_{j\ge0}A_j$は中央$R$-代数、$A_0=R$、
各$A_j$は有限自由で、$A\otimes_R k=k[x_1,\ldots,x_N]$、$N\ge5$とする。
$\Omega\in Z(A)_3$の剰余$F$は滑らかな射影三次超曲面を定めるものとし、
$B=A/(\Omega)$、$a=N-3$、$i=1-a$と置く。
完備化後も$A_{\widehat R}=\bigoplus_j(A_j\otimes_R\widehat R)$という
\emph{次数別}の代数を使い、内部次数についての無限級数は導入しない。

\subsection{完備Hom複体の縮約と代数的基底への降下}

\begin{lemma}[次数別完備な複体の有限縮約]\label{tw:lem:contraction}
$S$を剰余体$\mathbb C$の完備noetherian局所$\mathbb C$-代数とし、
$C^\bullet$を各項が有限自由$S$-加群である複体とする。
$C^\bullet\otimes_S\mathbb C$の全cohomologyの次元の和が有限なら、
$C^\bullet$は有界有限自由複体と$S$-線形にhomotopy同値である。
複体の全長が有限であることは仮定しない。
\end{lemma}
\begin{proof}
$C_0=C\otimes_S\mathbb C$のcohomologyを$H$とする。
体上の分裂から、次数を保つ$i_0:H\to C_0$、$p_0:C_0\to H$と
次数$-1$の$s_0$を
\begin{equation}\label{tw:eq:retract}
 p_0i_0=1,\qquad d_0s_0+s_0d_0=1-i_0p_0,
 \qquad s_0^2=p_0s_0=s_0i_0=0
\end{equation}
となるように選べる。各次数の基底を選び、これらを$S$へ延長する。
元の微分は$d=d_0+\delta$、$\delta\in\mathfrak m_S\operatorname{Hom}^1(C,C)$と書ける。
各固定次数は有限自由なので、次の逆行列はその次数内で収束する：
\begin{equation}\label{tw:eq:perturbation}
 \begin{aligned}
 J&=(1+s_0\delta)^{-1},& L&=(1+\delta s_0)^{-1},\\
 i'&=Ji_0,&p'&=p_0L,\\
 s'&=s_0L,&d_H&=p_0\delta Ji_0.
 \end{aligned}
\end{equation}
これは標準のhomological perturbation公式である
\cite[Theorem~2.3]{TWCrainic}。
同文献のhomotopyの符号は本稿の$-s_0$に対応する。
$d^2=0$と式\eqref{tw:eq:retract}を用いると
$d_H^2=0$、$di'=i'd_H$、$p'd=d_Hp'$、$p'i'=1$、
$ds'+s'd=1-i'p'$となる。
各積は固定次数での有限行列の積であり、逆行列の級数だけが$\mathfrak m_S$進級数である。
$H$は有界有限次元なので$(H\otimes_{\mathbb C}S,d_H)$が所要の有限複体である。
\end{proof}

\begin{lemma}[有限rank MFのproper Homと底変更]\label{tw:lem:proper-hom}
上の$R,A,\Omega$に対し、有限自由graded MFの任意の二対象$P,Q$の
Hom複体は$R$-perfectである。その通常の複体は有界有限自由複体とhomotopy同値であり、
$R\to k,\widehat R,\operatorname{Frac}(R)$について
\begin{equation}\label{tw:eq:hom-bc}
 \operatorname{Hom}_{\mathrm{MF}}(P,Q)\otimes_R^{\mathbf L}S
 \simeq\operatorname{Hom}_{\mathrm{MF}_S}(P_S,Q_S)
\end{equation}
となる。ここで右辺は有限rank行列の通常の底変更から作るHomである。
\end{lemma}
\begin{proof}
周期変数の二重次数を$(2,-3)$とする。
固定cohomological degreeのHom行列成分では周期変数の冪と内部次数が固定される。
従ってHom複体$C$の各項は有限個の$A_j$の直和、すなわち有限自由$R$-加群であり、
各項の形成と通常の底変更は可換である。
剰余のHomは滑らかな可換cubicのgraded MFのHomなので、
古典Orlov比較により有界有限次元cohomologyを持つ。
補題\ref{tw:lem:contraction}を$C\otimes_R\widehat R$へ適用すると、
完備化後の複体は有界有限自由複体へ縮約する。

$R$はnoetherianなので$H^q(C)$は各$q$で有限生成である。
平坦な完備化はcohomologyと可換であり、忠実平坦性から
$H^q(C)=0$は$q$の有限区間の外で成り立つ。
正則$R$の大域次元は有限なので$C$は$R$-perfectである。
さらに有界有限自由複体$D$から$C$への準同型を選ぶと、そのconeは
各項が射影である完全複体になる。
$e=\operatorname{gldim}R$とすれば、完全複体のcycleに対して
$\operatorname{Ext}^1_R(Z^q,M)=\operatorname{Ext}^{e+1}_R(Z^{q+e},M)=0$
というdimension shiftingが適用できる。よって全cycleは射影で、このconeは可縮である。
従って$C$自体が$D$とhomotopy同値、特にK-flatである。
項ごとの底変更の同定が式\eqref{tw:eq:hom-bc}を与える。
有限coneと直和因子にも同じ結論が延長される。
\end{proof}

\subsection{stabilizationとsection--windowの実際のenhancement}

\begin{lemma}[正則局所基底上のdg stabilization]\label{tw:lem:stabilization}
$A,B$は左右noetherian、$A$の大域次元は有限、$\Omega$は正則である。
余核による比較は、内部shiftと中心元に由来する$(3)\simeq[2]$を保つ
\begin{equation}\label{tw:eq:stabilization}
 \mathrm{MF}^{\mathbb Z}(A,\Omega)
 \simeq D^b_{\mathrm{sg,dg}}(\operatorname{gr}B)^{\mathrm{idem}}
\end{equation}
というdg Morita同値を与える。
本節の有限windowに現れる対象では、$k$および$\widehat R$への底変更と整合する。
\end{lemma}
\begin{proof}
noetherian性、相対Koszul分解、$\operatorname{gldim}A\le N+\dim R$、
$\Omega$の正則性、$B_j$の自由性と有限左右injective dimensionは、
補題\ref{lp:lem:noetherian}、\ref{lp:lem:resolution}と
命題\ref{lp:prop:koszul}、\ref{lp:prop:gorenstein}による。
これらの次数別の証明は基底の完備性や有限型性を要求せず、$\widehat R$でも通る。

Chenの余核定理\cite[Theorem~7.2]{AMChen}は三角圏の定理である。
そのままdg同値を引用するのではなく、以下のenhancementを用いる。
有限自由graded MFの成分別split exact圏を$\mathcal F$、
有限生成graded Gorenstein-projective $B$-加群のexact圏を$\mathcal G$とする。
$A$の大域次元が有限なので、$\mathcal G$の全対象は有限$A$-射影次元を持つ。
各有限graded projective $A$-加群は、$A_0=R$が局所であることから自由である。
余核はexact関手$\mathcal F\to\mathcal G$であり、projective-injectiveを保つ。
Chenの証明は生成元と持上げを斉次に選ぶことでこのgraded版にも適用できる。
右加群については$A^{\mathrm{op}}$へ同じ結果を適用する。
従って、その誘導するstable categoryの関手は同値である。

exact圏の有界導来dg圏をprojective-injectiveの有界複体で割ることで、
実際のdg関手のzigzagを得る。
左側とMFの比較は有限totalizationである。
成分別splitな完全複体に横方向のcontraction $s$を取ると、
縦MF微分$\partial$に対する補正$s(1+\partial s)^{-1}$は横方向の有限性により有限和である。
中心曲率$w$の残差は$sws=ws^2=0$なので、totalizationは可縮である。
右側と特異圏の比較は$\mathcal G$の包含から誘導される。
完全射影分解がstable Homを計算し、有限injective dimensionによる高いsyzygyが
本質的全射性を与える。各zigzagの$H^0$は上記の三角同値である。
全てのshiftにも同じ議論を適用することでHomの全cohomologyの同型となり、
式\eqref{tw:eq:stabilization}を得る。
これは補題\ref{am:lem:stabilization}のdg構成の正則基底版である。

底変更には余核構成を使う。有限MFの二つの行列の積は$\Omega$であり、
剰余では非零多項式$F$の乗法になるため、各行列は剰余上単射である。
固定内部次数で単元小行列式を選ぶと、余核は有限自由$R$-加群となる。
従って余核の周期分解、有限totalization、windowで使う次数別自由な有限長加群は
通常の剰余底変更と整合する。完備化はさらに平坦である。
Morita逆であるstabilizationの比較は、この余核の自然な比較の逆として選ぶ。
任意のcoherent対象の任意の非正則底変更まで主張してはいない。
\end{proof}

\begin{theorem}[三対象window定理：局所比較(L)の確立]\label{tw:thm:window}
本節の代数的入力だけを仮定する。$K_R=\operatorname{stab}(R)$、
$T_R=K_R\oplus K_R(1)\oplus K_R(2)$と置くと、
\begin{equation}\label{tw:eq:main}
 \mathcal K_{B_R}^{\mathrm{lin}}
 \simeq\operatorname{thick}\{K_R,K_R(1),K_R(2)\}
 \simeq\operatorname{Perf}(\Gamma_R),\qquad
 \Gamma_R=R\operatorname{End}(T_R)\in\operatorname{Perf}(R)
\end{equation}
である。最初の同値は
\begin{equation}\label{tw:eq:window}
 \Psi_R=\operatorname{stab}\circ\bigl((R\omega -)_{\ge i}\bigr),
 \qquad i=1-a
\end{equation}
の制限で与えられる。
これは$k,\widehat R,\operatorname{Frac}(R)$への底変更と整合し、
$L_{\mathcal O}\circ(1)$をMFの内部shiftへ移す。
特にこのrotationは自己同値であり、中心元に由来する三次周期同型を持つ。
\end{theorem}
\begin{proof}
$E_R=\bigoplus_{j=0}^{N-1}\mathcal O(j)$は線束生成部分を生成し、
その全対象対のHomは$R$-perfectである。
例外列$\mathcal E=\langle\mathcal O,\ldots,\mathcal O(a-1)\rangle$を除くmutationを$p$とすると、
$G_R=p(E_R)$が$\mathcal K^{\mathrm{lin}}$を生成する
（命題\ref{lp:prop:residual}）。
必要なsection、有限cohomological dimension、$\chi$の入力は
Burke--Stevenson\cite[Proposition~4.18, Lemmas~6.1--6.2]{LP-BS}による。

\emph{関手の定義。}
線束Homの公式から、$0\le j\le N-1$では
$(R\omega\mathcal O(j))_{\ge i}=B(j)_{\ge i}$である。
$T_j=B(j)_{<i}$とすると、$T_j$の各次数は有限自由で、
$R$のshiftによる有限filtrationを持つ。
$R$の相対Koszul分解と$B$の長さ一の$A$-分解により、
$B(j)_{\ge i}$は$A$-perfectである。
従ってderived section、内部切断、dg商、補題\ref{tw:lem:stabilization}のMorita逆を合成して
式\eqref{tw:eq:window}が実際のdg quasi-functorとして定義される。
$0\le j<a$では$T_j=0$なので$\Psi_R(\mathcal O(j))=0$であり、
$\Psi_Rp\simeq\Psi_R$である。

\emph{完全忠実性。}
この関手の剰余は、古典的なsection--windowを特異圏へ送る関手である。
Burke--Stevenson\cite[Theorem~6.4とその証明]{LP-BS}において$i=1-a$を取ると、
そのkernelは$\mathcal E_k$であり、Verdier商を特異圏と同定する。
本稿の右直交と同文献の半直交分解の向きは異なり得るが、
双方が同じVerdier商を表示するため、$\mathcal E_k^\perp$への制限は同値である。

sourceの線束Homの底変更は命題\ref{lp:prop:basechange}による。
sectionの内部切断は生成線束上で上の加群の表示になり、
stabilizationの底変更は補題\ref{tw:lem:stabilization}による。
有限cone、直和因子、evaluationを通じ、$\Psi_R$のHom上の\emph{自然な写像}は
剰余で古典windowのHom写像へ移る。
targetのHomは補題\ref{tw:lem:proper-hom}により$R$-perfectかつ剰余底変更と可換である。
従って任意の二対象についてHom比較のconeはperfectで剰余が零となり、
局所環上の導来Nakayamaにより零である。これがenhancedな完全忠実性を与える。
単にHomの階数だけを比較しているのではない。

\emph{像の同定。}
残る$j=a,a+1,a+2$について、$B(j)$は特異圏で零なので
\begin{equation}\label{tw:eq:triangular}
 \Psi_R(\mathcal O(j))\simeq\operatorname{stab}(T_j)[-1],\qquad
 \operatorname{gr}(T_j)=\bigoplus_{r=0}^{j-a}B_r\otimes_RR(j-r).
\end{equation}
$B_0=R$である。まず$T_a=R(a)$、次に$T_{a+1}$の高次数部分は
$R(a)$の有限直和、最後に$T_{a+2}$の高次数部分は
$R(a),R(a+1)$の有限extensionである。
その商はそれぞれ$R(a+1),R(a+2)$なので、三角的に
\begin{equation}\label{tw:eq:three-thick}
 \operatorname{thick}(\Psi_R(G_R))
 =\operatorname{thick}\{K_R(a),K_R(a+1),K_R(a+2)\}
 =\operatorname{thick}(T_R).
\end{equation}
最後はMFの$(3)=[2]$による。
完全忠実なexact関手の冪等完備sourceの像はthickなので、これが式\eqref{tw:eq:main}の同値となる。
$\Gamma_R$のproper性は補題\ref{tw:lem:proper-hom}からも従う。
生成性の議論は有限長加群の実際のfiltrationによるものであり、
特殊ファイバーの生成性だけを持ち上げたものではない。

\emph{rotationと三次周期性。}
線束Homの消滅とevaluation三角により、$L_{\mathcal O}(M(1))$は再び残余圏に属する
（命題\ref{lp:prop:rotation-candidate}）。
$M\in\mathcal K^{\mathrm{lin}}$なら
$(R\omega M)_i=R\operatorname{Hom}(\mathcal O(a-1),M)=0$である。
切断の自然な三角により
\begin{equation}\label{tw:eq:rotation-cut}
 ((R\omega M)(1))_{\ge i}
 =(R\omega M)_{\ge i+1}(1)
 \simeq(R\omega M)_{\ge i}(1).
\end{equation}
$\Psi_R(\mathcal O)=0$とevaluationの三角を使うと
$\Psi_R(L_{\mathcal O}M(1))\simeq\Psi_R(M)(1)$を得る。
式\eqref{tw:eq:three-thick}の像は$(1)$と$(-1)$の両方で安定なので、
rotationは自己同値となる。
周期同型はMFで中心元$\Omega$から定まる$(3)=[2]$をこの同値で移したものとする。
この選択は切断、余核、evaluationの自然な比較から成り、底変更に整合する。
任意に別のスカラー倍を付けた周期同型を保存するという主張ではない。

完備化・分数体への底変更は平坦であり、有限MF行列のHomと積に項ごとに可換する。
sourceも有限線束生成とmutationから作った模型の底変更に一致する。
剰余底変更は上で確認した。これで必要な三対象とrotationの比較をすべて得る。
\end{proof}

\subsection{固定三対象で十分となる比較原理}

\begin{definition}[固定三対象の形式的比較]\label{tw:def:CF3}
条件\textup{(CF$_3$)}とは、$W\in\widehat R[x_1,\ldots,x_N]_3$、$W\bmod\mathfrak m=F$と、
可換側の有限rank graded MF $E$があり、
$E\otimes k\simeq K_F$で、$K_{\widehat R}(j)$から$E(j)$へ送る
三対象のfull dg部分圏の$\widehat R$-線形unital $A_\infty$ quasi-equivalenceがあることをいう
（$j=0,1,2$）。さらに、その有限cone・冪等包絡への延長は内部shiftと、
中心元からの三次周期比較を保つものとする。
これは全Hom複体と高次積の比較であり、cohomologyの次元の一致ではない。
\end{definition}

\begin{corollary}[固定三対象の比較による復元]\label{tw:cor:CF3}
定理\ref{gdr:thm:generic}の代数的入力に対して、
局所比較(L)は定理\ref{tw:thm:window}の結論であり、追加仮定ではない。
完備MF圏全体の比較を用いない場合にも、\textup{(CF$_3$)}の比較があれば
同主定理の結論が成立する。
全MF比較\textup{(CF)}からは、三対象への制限によって\textup{(CF$_3$)}が従う。
\end{corollary}
\begin{proof}
定理\ref{tw:thm:window}の$\Gamma_R$を使い、\textup{(CF$_3$)}を固定三対象上で
$L=\operatorname{Frac}(\widehat R)$へ平坦局所化する。
得られる像は$E_L,E_L(1),E_L(2)$のthick閉包である。
補題\ref{tw:lem:proper-hom}によってHomはproperで特殊底変更と可換なので、
命題\ref{rog:prop:euler}の非零Eulerベクトルが一般点でも保存され、$E_L\ne0$となる。
以後は既存の代数閉包上のBFK生成、有限降下、coherent飽和、基礎体上の勾配積分を使う。
これらは全MF圏上の形式functorを要求せず、固定三対象の比較だけを使っている。
\end{proof}

\begin{remark}[比較原理の適用範囲]\label{tw:rem:scope}
局所windowの構成は完備MF比較(CF)を入力とせず、
正則多変数局所基底にも適用できる。
その証明で使うperturbation lemma、MFの余核定理、
Gorenstein windowは既知の結果であり、ここでは
enhancement、有限filtration、底変更およびrotationとの整合を明記した。
固定三対象の比較(CF$_3$)を与える実際の曲率付き関手と逆関手は
第\ref{cft:sec:transport}節で構成する。
抽象的な非可換三次圏に中心三次超曲面表示が存在することは
本節の仮定にも結論にも含めない。
\end{remark}

% Complete formal comparison and its three-object restriction.
\section{次数別完備な曲率付き移送と形式比較の構成}
\label{cft:sec:transport}

本節では、完備係数上のMF比較を実際の関手と逆関手によって構成し、
その固定三対象への制限を明らかにする。
既知のKontsevich formalityを\emph{曲率を入れる前}に使い、
その後のゲージ経路は有限Artin次数ごとに直接積分する。
曲率の大きさ、全Hochschild cochain、単位、逆写像、有限rank、
二つの次数、指定された周期比較をそれぞれ確認する。
形式同値の存在を、解析的収束や全Laurent対象への関手と混同しない。

\subsection{主張と既知結果の使用範囲}
$(R,\mathfrak m)$を剰余体$\mathbb C$の完備noetherian局所$\mathbb C$-代数とする。
ここでは$R$の正則性は不要である。
$S=\mathbb C[x_1,\ldots,x_N]$、$N\ge3$、$\deg x_i=1$とし、
$A^{\rm gr}=\bigoplus_{j\ge0}A_j$をunitalな中央$R$-代数とする。
各$A_j$は有限自由、$A_0=R$、$A^{\rm gr}/\mathfrak m A^{\rm gr}=S$と仮定する。
$\Omega\in Z(A^{\rm gr})_3$の剰余$F$は滑らかな射影三次超曲面を定めるものとする。
次数別に剰余の単項式基底を持ち上げ、単位を保つ加群同定
$A_j\simeq S_j\otimes_{\mathbb C}R$を選ぶ。

\begin{theorem}[完備係数上の実際のgraded MF比較]
\label{cft:thm:complete}
上の入力に対し、$W\in R[x_1,\ldots,x_N]_3$、$W\bmod\mathfrak m=F$と、
有限rank graded MFの実際の$R$-線形unital $A_\infty$同値
\begin{equation}\label{cft:eq:equivalence}
 \operatorname{MF}^{\mathbb Z}(A^{\rm gr},\Omega)
 \simeq\operatorname{MF}^{\mathbb Z}(R[x_1,\ldots,x_N],W)
\end{equation}
が存在する。これは内部shiftと、中心元による$(3)\simeq[2]$を保ち、
剰余では恒等であり、すべての$R/\mathfrak m^{q+1}$への切断と両立する。
有限自由模型上で互いに逆な関手を構成し、有限coneと冪等完備化へ延長する。

曲率付きbar写像のTaylor成分$F_r$と逆写像の成分$H_r$は
\begin{equation}\label{cft:eq:bound}
 F_0=H_0=0,\quad F_1,H_1=1+O(\mathfrak m),\quad
 F_{2r}=H_{2r}=0,\quad F_{2r+1},H_{2r+1}\in\mathfrak m^r\quad(r\ge1)
\end{equation}
を満たすように取れる。次数付き代数の厳密同型は、この写像の前後に合成する。
\end{theorem}

Kontsevichの既知結果は、affine formality、HKR、Maurer--Cartan変形問題の比較
\cite[Theorem~4.6, Section~4.6.2, Theorem~6.3]{PKontsevich}である。
曲率付きtwistによるgraded MFの記述には
\cite[Section~2, 特にRemark~2.8とTheorem~2.12]{CT}を使う。
後者を形式的な無限arity関手へ無条件に適用せず、以下で評価式の有限性と収束を証明する。
これらの一般理論の新規性は主張せず、本定理の組合せの初出は未確認である。

\subsection{通常の結合変形と中心元の有限次元逆像}
\begin{lemma}[次数を保つunitalな星積模型]\label{cft:lem:model}
formal Poisson tensor
$\pi\in\mathfrak m(S_2\otimes\bigwedge^2\langle\partial_i\rangle)\widehat\otimes R$
と、単位・内部次数を保つ厳密な代数同型
\begin{equation}\label{cft:eq:phi}
 \phi:A^{\rm gr}\xrightarrow{\sim}(S\otimes R,*_\pi),\qquad
 \phi\equiv1\pmod{\mathfrak m}
\end{equation}
が存在する。右辺も内部次数ごとの直和とする。
\end{lemma}
\begin{proof}
単位を入力すると零になるnormalized Hochschild cochainを使う。
多項式環の対角のgraded Koszul分解とnormalized bar分解の比較は
内部次数を保つように選べる。従ってpolydifferential cochainから
\emph{全ての}内部次数指定Hochschild cochainへの包含はHKR同型を誘導する。
既存の積がbidifferentialであるという仮定は不要である。
Kontsevichのグラフ写像はaffine線形変換に同変なので、内部weightを保つ。
weight零に制限したformalityと上の包含の合成は準同型であり、
次数別flatな結合積の変形をPoisson変形に対応させる。

この対応を各$R/\mathfrak m^{q+1}$で行う。
一段の核$\mathfrak m^q/\mathfrak m^{q+1}$では、持上げ障害の比較はHKRの
次数二での単射性、積の差の補正は次数一での全射性による。
一入力gaugeを同時に持ち上げるので、すでに選んだ低次jetを保てる。
これは\cite[Theorem~4.6]{PKontsevich}の冪零変形問題の比較の逐次形である。
全Hochschild cochain側のdegree零gaugeは一入力の写像であり、
厳密な代数同型を与える。normalized性から単位を保つ。
極限は各有限自由$A_j$内で存在し、逆写像も同様に存在する。
weight零のbivectorの係数次数は二である。
\end{proof}

\begin{lemma}[中心逆像を有限行列で取り出す]\label{cft:lem:center}
formalityを固定すると、$W\in S_3\otimes R$が一意に存在して
\begin{equation}\label{cft:eq:center}
 C_\pi(W)=\phi(\Omega),\qquad [\pi,W]=0,\qquad W\bmod\mathfrak m=F,
 \quad
 C_\pi(f):=\sum_{r\ge0}\frac1{r!}\mathcal U_{r+1}(f,\pi,\ldots,\pi)
\end{equation}
となる。
\end{lemma}
\begin{proof}
関数一個とbivector $r$個のグラフの出力arityは零であり、内部weightを保つ。
よって$C_\pi$は有限自由加群$S_3\otimes R$の$R$-線形自己写像で、
$\mathcal U_1(f)=f$により$C_\pi=1+O(\mathfrak m)$である。
その逆は各固定次数内で収束する幾何級数であり、$W$を一意に定める。

接写像のchain恒等式は
\begin{equation}\label{cft:eq:chain-center}
 [*_\pi,C_\pi(W)]=\mathcal U^\pi([\pi,W])
\end{equation}
である。左辺は中心性から零。
右辺の入力は係数次数四のvector場であり、接写像の最低次は
vector場を導分として見るHKRの単射である。
$\mathfrak m^q/\mathfrak m^{q+1}$で順にこの単射性を使うと
$[\pi,W]\in\bigcap_q\mathfrak m^q=0$となる。
曲率付きの通常のHochschild cohomologyをJacobian環と同一視する操作は使わない。
\end{proof}

\subsection{三次式の次数制約とゲージ経路}
\begin{lemma}[唯一の定数三重vectorと二回目の作用]\label{cft:lem:primitive}
一意な$\Xi\in\mathfrak m\bigwedge^3\langle\partial_i\rangle\otimes R$が存在し、
\begin{equation}\label{cft:eq:primitive}
 \pi=[\Xi,W],\qquad [\Xi,\pi]=0
\end{equation}
となる。
\end{lemma}
\begin{proof}
滑らかな$F$の全偏微分は多項式環の正則列である。
勾配Koszul contractionは、係数次数とrankの和「係数次数$+2\,\mathrm{rank}$」を保つ。
各斉次部分は有限の有限自由$R$-複体である。
剰余の完全性を使い、$\mathfrak m^q/\mathfrak m^{q+1}$で原像を逐次持ち上げると、
完備化と分離性から正のrankで完全となる。
これは補題\ref{fo:lem:regular-gradient}の多変数・完備局所版でもある。

$[\pi,W]=0$なので、係数次数二・rank二のcycle $\pi$の原像は
係数次数零・rank三に取れる。二つの原像の差の原像は係数次数$-2$を要求するため零。
これで$\Xi$の存在・一意性と$\Xi\bmod\mathfrak m=0$を得る。
$D=[\Xi,\pi]$は係数次数一・rank四であり、graded Jacobiと$[\pi,\pi]=0$から
$[W,D]=-[\pi,\pi]=0$となる。
そのKoszul原像は係数次数$-1$を要求するので$D=0$である。
ここでJacobi条件は必須であり、Casimir条件だけから第二の等式は従わない。
\end{proof}

周期変数$t$の二重次数を$(2,-3)$とし、$t$を微分せず、係数として多重線形に扱う。
$\eta=\Xi/t$はcohomological degree零・内部weight零である。
\begin{equation}\label{cft:eq:path}
 \alpha_s=-tW+(1-s)\pi,\qquad
 \partial_s\alpha_s=[\eta,\alpha_s]=-\pi
\end{equation}
はMaurer--Cartan経路で、始点は非可換模型、終点は可換模型に対応する。
$W$は$\mathfrak m$の倍数でないので、通常の冪零MC定理をこの経路へ直接適用してはいけない。

\begin{lemma}[大きい曲率でも有限なformality評価]\label{cft:lem:arity}
$\alpha_s$の像はcurved associative構造
$b_s=(*_{(1-s)\pi},-tC_{(1-s)\pi}(W))$である。
$\gamma_s=\mathcal U^{\alpha_s}(\eta)$はarity一と三だけを持ち、
$\gamma_s\in\mathfrak m$で、
\begin{equation}\label{cft:eq:lax}
 \partial_sb_s=[\gamma_s,b_s]
\end{equation}
となる。$\gamma_s$の正のarity成分はnormalizedである。
\end{lemma}
\begin{proof}
グラフの辺数と配置空間の次元$2n+m-2$を等しくすると、
関数$k$個とbivector任意個の出力arityは$2-2k$、
さらにtrivector一個を入れた場合は$3-2k$となる。
従って関数入力は高々一個で、残る入力は$\mathfrak m$の倍数である。
mod $\mathfrak m^{q+1}$では必要な全Taylor和が有限となり、
$L_\infty$恒等式も有限和の恒等式として適用できる。
chain恒等式に$\partial_s\alpha_s=-[\alpha_s,\eta]$を代入すると
$\partial_sb_s=-[b_s,\mathcal U^{\alpha_s}(\eta)]$、すなわち式\eqref{cft:eq:lax}である。

正のarityで単位をboundary入力へ置く。
そこへ辺が入れば定数の微分が零であり、入らなければ積分形式はその頂点を忘れた
一つ低次元の配置空間からの引戻しなので最高次積分が零となる。
internal頂点が一個の場合はHKR写像であり、同じ消滅は直接分かる。
この議論は\cite[Sections~6.2--6.4]{PKontsevich}のグラフ表示に適用する。
$t$への多重線形性と内部weightは全ての項で保たれる。
\end{proof}

\subsection{bar写像と逆写像の一様なarity評価}
\begin{lemma}[順序付きflowと逆写像]\label{cft:lem:flow}
$\partial_s\mathcal F_s=\gamma_s\mathcal F_s$、$\mathcal F_0=1$を解くと、
各$R/\mathfrak m^{q+1}$上で実際のbar coalgebra同型を得る。
$\mathcal F_1b_0=b_1\mathcal F_1$であり、写像・逆写像は式\eqref{cft:eq:bound}を満たす。
\end{lemma}
\begin{proof}
mod $\mathfrak m^{q+1}$で
\begin{equation}\label{cft:eq:ordered}
 \mathcal F_s=1+\sum_{j=1}^{q}
 \int_{0\le u_j\le\cdots\le u_1\le s}
 \gamma_{u_1}\cdots\gamma_{u_j}\,du_j\cdots du_1
\end{equation}
とする。各係数は$s$の多項式で、積分は標数零での代数的積分である。
$\gamma$がcoderivationなので、$\Delta\mathcal F_s$と
$(\mathcal F_s\otimes\mathcal F_s)\Delta$は同じ微分方程式・初期値を持つ。
よってcoalgebra写像である。
$\mathcal F_s-1\in\mathfrak m$から逆は
$\sum_{j=0}^{q}(1-\mathcal F_s)^j$で与えられる。
式\eqref{cft:eq:lax}により$\mathcal F_s^{-1}b_s\mathcal F_s$の微分は零である。

零項成分はなく、三入力coderivationを一回使うごとに出力長は二つ減る。
一入力成分は長さを変えない。従ってarity $2r+1$には少なくとも$r$回の
coderivationが必要で、係数は$\mathfrak m^r$に入る。
偶数arityは現れず、逆flowにも同じ評価が成立する。
出力tensor長$\ell$に寄与する成分の積の位数を足すと、mod $\mathfrak m^{q+1}$では
入力長は高々$\ell+2q$である。従ってtensor長完備barへの延長も各剰余で連続である。
単位はnormalized性からflow中固定され、高arity成分の単位入力は零となる。
同じ形式的係数を切断しているので、商$R/\mathfrak m^{q+1}\to R/\mathfrak m^q$と両立する。
\end{proof}

\subsection{冪零でないMF行列への評価と周期比較}
\begin{proof}[定理\ref{cft:thm:complete}の証明]
上のbar写像を有限個のshiftを持つ自由加群の行列成分に拡張する。
合成可能な添字列に沿う和とsuspensionのKoszul符号を用いるので、
元のbar恒等式は行列上でも成立する。
MF行列$Q$は$\mathfrak m$の倍数でも上三角でもなくてよい。
mod $\mathfrak m^{q+1}$では、像は有限式
\begin{equation}\label{cft:eq:object}
 Q'=\sum_{r=0}^{q}F_{2r+1}(Q,\ldots,Q)
\end{equation}
で与えられる。$k$個の射に対するTaylor成分は、suspended表示で
\begin{equation}\label{cft:eq:insertions}
 \sum_{r_0,\ldots,r_k\ge0}
 F_{k+r_0+\cdots+r_k}
 (Q_k^{\otimes r_k},a_k,Q_{k-1}^{\otimes r_{k-1}},\ldots,a_1,Q_0^{\otimes r_0})
\end{equation}
である。$k+\sum r_i\le2q+1$の項だけが寄与するため有限である。

suspensionを$\varsigma$と書くと、tensor長完備bar上で$\mathcal F(e^{\varsigma Q})=e^{\varsigma Q'}$である。
$\mathcal F b_0=b_1\mathcal F$から対象のMF方程式と全functor恒等式が従う。
同じ構成を逆bar写像に施すと、挿入公式はbar写像の合成と可換するので、
二つのtwisted functorの合成は恒等となる。
したがって各Artin切断で実際の互いに逆な関手がある。

固定した有限自由加群、固定したcohomological次数と内部次数では、
各行列成分の$t$の冪と$x$の次数が固定される。
従って出力は有限個の$S_j\otimes R$の直和に属する。
式\eqref{cft:eq:bound}から、式\eqref{cft:eq:object}と\eqref{cft:eq:insertions}は
この有限自由加群内で$\mathfrak m$進収束する。
すべての恒等式は各Artin商で成立しているため、分離性から$R$上でも成立する。
これは同値の列の形式的な逆極限を宣言するのではなく、同じ基礎行列上の
関手と逆関手の明示的な係数wise極限である。

内部shiftは基礎加群のラベルを変えるだけであり、行列上の公式は不変である。
各$F_r,H_r$は$R[t,t^{-1}]$-多重線形なので、$t$が与える
$(3)\simeq[2]$の自然比較とも可換する。
従って周期同型が未知のスカラー倍だけずれる、という曖昧さはない。
別に選んだ任意の周期同型については主張しない。
正規化された単位も行列上で保たれる。
補題\ref{cft:lem:model}の厳密同型$\phi$と合成し、有限cone・冪等完備化へ延長して
式\eqref{cft:eq:equivalence}を得る。
\end{proof}

\begin{corollary}[固定三対象の比較条件の充足]\label{cft:cor:CF3}
定理\ref{tw:thm:window}の代数的入力について、完備化後の(CF$_3$)は成立する。
従って一般点復元の主定理において(L)、(CF)、(CF$_3$)を別の追加仮定とする必要はない。
正則多変数局所基底でも同じである。
\end{corollary}
\begin{proof}
$K_{\widehat R}=\operatorname{stab}(\widehat R)$を有限rank模型に取り、
定理\ref{cft:thm:complete}の像を$E$とする。
剰余で関手は恒等なので$E\otimes\mathbb C\simeq K_F$である。
三対象$K,K(1),K(2)$への制限と、内部shift・周期比較の整合が定義\ref{tw:def:CF3}を満たす。
一般点へ移るときは、この固定した有限対象集合のHomと高次積を先に構成してから
平坦な体拡大を行う。任意のLaurent係数MFへ形式級数を評価する必要はない。
\end{proof}

\begin{remark}[射程と検証の境界]\label{cft:rem:scope}
本節は指定した代数的入力の形式比較を証明するが、任意の抽象的nc cubic圏に
そのような表示があることを証明しない。$W$の係数の代数性・収束や、同値核の
元の体への降下、特殊点を含む代数的相対同値、全モジュライへの接着も別問題である。
中心元を変えれば$W$も変わり、一次Poisson tensorだけから復元先は決まらない。
特定の混合族についての明示二次jetおよび具体的可換模型の計算は、
第\ref{kj:sec:second-jet}節と第\ref{gdr:cor:mixed}節の仮定の下で扱う。
係数の有限次数計算は符号・arityの照合には有用だが、
Kontsevichの全グラフ積分や、この節のすべての高次関手恒等式の証明に
代わるものではない。
\end{remark}

% 2026-11-10 / v22
\section{有限curved gaugeからgraded MF同値への比較}
\label{cg:sec:comparison}

本節では、有限Artin基底上のcurved $A_\infty$ bar同型を、
有限rankのgraded matrix factorizationへ実際に延長する。
必要なのは、特殊ファイバーでのquasi-isomorphismだけではなく、
零項成分を持たず、写像と逆写像のarityが有限であるbar同型である。
この条件の下では、twisting行列が冪零でなくても全ての評価式が有限になる。

Kontsevichのformality定理
\cite[Sections~5--6, 特にSection~6.4]{PKontsevich}と、
curved twisted complexの枠組み
\cite[Theorem~2.12]{CT}\cite[Sections~2.2, 6.1]{CGTu}
を用いる。
Tuの結果は曲率・内部次数を扱う先例であり、
以下の非可換Artin基底上の有限rank比較を、そのまま引用して済ませるものではない。
本節では有限性、逆写像、内部シフトの保存をそれぞれ確認する。
この組合せの新規性は未確認である。

\subsection{周期変数と有限rank模型}

$\CC$-代数$R$上の内部次数付き代数を考え、
内部次数三の中心元を$W$とする。
周期変数$t$を加え、その二種類の次数を
\begin{equation}\label{cg:eq:bidegree}
 \deg_{\rm coh}(t)=2,\qquad \deg_{\rm int}(t)=-3,\qquad
 \deg_{\rm coh}(z_i)=0,\quad \deg_{\rm int}(z_i)=1
\end{equation}
と定める。$h$の両次数は零とする。
$t$は可逆な係数であり、以下の多重微分作用素で$t$を微分しない。

曲率の符号を$m_0=-tW$とする。
有限個のcohomological・内部shiftを持つ自由加群$E$上で、
cohomological degree一、内部次数零の行列$Q$が
\[
 m_0+Q^2=-tW+Q^2=0
\]
を満たすことが、graded MFの条件である。
この周期表示では$t$の作用が指定された同型$(3)\simeq[2]$を与える。
有限自由模型に有限coneと直和因子を加えた圏を
$\MF^{\ZZ}$と書く。一般のcurved $A_\infty$構造では、
同じ記法の対象条件を$\sum_{k\ge0}m_k(Q,\ldots,Q)=0$とする。

\subsection{有限bar同型のtwisted complexへの延長}

以下では$sE=E[1]$を用い、$sQ$の次数が零となるsuspension規約を採用する。
全ての符号はsuspended tensor coalgebraのKoszul規約で固定する。

\begin{lemma}[有限arityのbar同型とtwisted complex]
\label{cg:lem:tw-extension}
$m,m'$をstrictly unitalなcurved $A_\infty$構造とし、
それらのarityは有界であるとする。
対応するbar coderivationを$b,b'$とする。
零項成分を持たないbar coalgebra同型$\mathcal F$で
\[
 \mathcal F b=b'\mathcal F,\qquad F_0=0
\]
を満たすものがあり、$\mathcal F$と$\mathcal F^{-1}$の
Taylor成分がともに有限arityであると仮定する。
さらに両写像はstrictly unitalとする。
このとき$\mathcal F$は、有限rankのcurved twisted complexの間の
実際の$A_\infty$同値へ延長される。
対象に対する式は
\begin{equation}\label{cg:eq:object-map}
 Q'=\sum_{k\ge1}F_k(Q,\ldots,Q)
\end{equation}
である。内部次数と$t$の作用を保つ$\mathcal F$の場合、
この同値は内部シフトと$(3)\simeq[2]$を保つ。
\end{lemma}
\begin{proof}
まず有限個のshiftの直和に対して、$F_k$を行列成分の
合成可能な添字列に沿って作用させる。
各$k$について有限和なので、同じbar恒等式が行列上でも成立する。

tensor長について完備化したcoalgebra
$\widehat T^c(sE)=\prod_{r\ge0}(sE)^{\otimes r}$を用いる。
$F_k=0$（$k>L$）なら、出力長$r$の成分に寄与する入力長は
$r$から$Lr$までである。
従って$\mathcal F$はこの完備化上で定義され、連続である。
逆写像についても同じである。
$b,b'$もarityが有界なので連続であり、曲率の零項成分はtensor長を一つ増やすだけである。
ここでは$Q$の冪零性を仮定していない。

$q=sQ$と置き、
\[
 e^q=\sum_{r\ge0}q^{\otimes r}
\]
を考える。
bar写像がcoalgebra写像で$F_0=0$であることから
\[
 \mathcal F(e^q)=e^{q'},\qquad
 q'=\sum_{k\ge1}\overline F_k(q,\ldots,q)
\]
となる。上線はsuspension後のTaylor成分を表す。
$Q$のMaurer--Cartan方程式は$b(e^q)=0$と同値である。
従って
\[
 b'(e^{q'})=b'\mathcal F(e^q)=\mathcal F b(e^q)=0
\]
であり、式\eqref{cg:eq:object-map}は実際に対象を与える。

対象$Q_0,\ldots,Q_n$とsuspendedな射$a_1,\ldots,a_n$に対し、
functorの$n$項成分を
\begin{equation}\label{cg:eq:morphism-insertions}
 \sum_{r_0,\ldots,r_n\ge0}
 \overline F_{n+r_0+\cdots+r_n}
 (q_n^{\otimes r_n},a_n,q_{n-1}^{\otimes r_{n-1}},
 \ldots,a_1,q_0^{\otimes r_0})
\end{equation}
と定める。
arity上界により、この和も有限である。
これは
$e^{q_n}\otimes a_n\otimes e^{q_{n-1}}\otimes\cdots
\otimes a_1\otimes e^{q_0}$
への$\mathcal F$の作用を出力長一に射影したものである。
同じtensorへの$\mathcal F b=b'\mathcal F$の作用を展開すると、
twisted structureと式\eqref{cg:eq:morphism-insertions}の
$A_\infty$ functor恒等式を得る。
曲率に由来する項は各$q_j$のMaurer--Cartan方程式で消える。

$\mathcal F^{-1}$に同じ構成を行う。
twistingの挿入公式はbar写像の合成と可換なので、
二つのtwisted functorの合成は
$\mathcal F^{-1}\mathcal F=\Id$および
$\mathcal F\mathcal F^{-1}=\Id$のtwistingに等しい。
従って得られる二つのfunctorは実際に互いに逆である。
strict unitは、正のarity成分のnormalized性から保たれる。
有限rankは基礎となる有限自由加群を変更しないことから保たれる。

内部次数零の写像は内部shiftと可換し、
$R[t,t^{-1}]$-多重線形性は$t$の作用と可換する。
従って内部シフトと指定された周期同型を保つ。
有限cone・冪等完備化へ延長することで、対応するdg Morita同値も得る。
\end{proof}

\begin{lemma}[任意の有限Artin次数でのflowのarity]
\label{cg:lem:artin-flow}
$R_N=\CC[h]/(h^{N+1})$とする。
$\gamma_s$を、補助変数$s$に関して多項式であるdegree零のcoderivationとし、
$\gamma_s\in h$、その非零Taylor成分はarity一と三だけとする。
このとき
\[
 \partial_s\mathcal F_s=\gamma_s\mathcal F_s,\qquad
 \mathcal F_0=\Id
\]
の解は有限の順序付き積分で与えられるbar coalgebra同型である。
$\mathcal F_1$とその逆のTaylor成分は
\begin{equation}\label{cg:eq:general-arity}
 \begin{aligned}
 F_0&=0,\qquad F_1=\Id+O(h),\\
 F_{2r+1}&=O(h^r)\quad(1\le r\le N),\\
 F_{2r}&=0,\qquad F_k=0\quad(k>2N+1).
 \end{aligned}
\end{equation}
を満たす。
$\gamma_s$がnormalizedで内部次数と$t$の作用を保つなら、
解と逆写像もstrictly unitalで、同じ次数と作用を保つ。
\end{lemma}
\begin{proof}
解を
\[
 \mathcal F_s=\Id+
 \sum_{j=1}^{N}
 \int_{0\le u_j\le\cdots\le u_1\le s}
 \gamma_{u_1}\cdots\gamma_{u_j}\,
 du_j\cdots du_1
\]
と置く。$j$個のcoderivationの合成は$h^j$の倍数なので、
$j>N$の項は零である。
積分は$\CC$上の多項式の積分であり、解析的な収束を用いない。
この有限式を微分するとflow方程式を満たす。
同じ方程式の解の一意性は$h$の次数ごとの帰納法から従う。

$\gamma_s$のcoderivation恒等式により、
$\Delta\mathcal F_s$と
$(\mathcal F_s\otimes\mathcal F_s)\Delta$は同じflow方程式と初期条件を満たす。
従って$\mathcal F_s$はcoalgebra写像である。
$\mathcal F_s-\Id\in h$なので
\[
 \mathcal F_s^{-1}
 =\sum_{j=0}^{N}(\Id-\mathcal F_s)^j
\]
も有限式で定義される。

零項成分がないので$F_0=0$である。
coderivationの三項成分を一回使うごとに、
一つのTaylor成分を作る入力数は二つ増える。
従ってarity $2r+1$には少なくとも$r$回のcoderivationが必要であり、
その$h$進位数は少なくとも$r$である。
高々$N$回しか合成できないためarityは$2N+1$以下で、偶数arityは現れない。
逆flowまたは上の有限逆級数にも同じ数え上げが適用できる。

normalizedなcoderivationのflowは単位を固定し、
二項以上で単位を含む成分を生じさせない。
内部次数と$t$に関する多重線形性は合成・積分・逆写像で保存される。
\end{proof}

\subsection{曲率を含むKontsevich写像の有限性}

\begin{lemma}[arity、曲率入力、単位]
\label{cg:lem:arity}
Kontsevichのaffine formality写像を$\mathcal U$とする。
関数を$k$個、bivectorを$\ell$個入力したTaylor成分の出力arityは
$2-2k$である。
一個のtrivectorに加えて関数を$k$個、bivectorを$\ell$個入力した場合は
$3-2k$である。
負のarityの成分は零である。
従って前者はarity零または二、後者はarity一または三に限られる。
各bivectorとtrivectorが$h$の倍数であれば、固定した有限Artin次数で
必要となるTaylor和は全て有限である。
正のarityの成分はnormalizedである。
\end{lemma}
\begin{proof}
$n$個のinternal vertexと$r$個のboundary vertexを持つ
Kontsevichの配置空間の次元は$2n+r-2$である。
関数には出る辺がなく、bivectorには二本、trivectorには三本の辺がある。
関数$k$個とbivector $\ell$個の場合、辺数は$2\ell$、$n=k+\ell$なので
\[
 2\ell=2(k+\ell)+r-2,\qquad r=2-2k.
\]
一個のtrivectorを加えると
\[
 3+2\ell=2(1+k+\ell)+r-2,\qquad r=3-2k.
\]
従ってどちらも関数入力は高々一個である。
曲率の関数が$h$の倍数でなくても、このarity制約が関数入力数を制限する。
他の入力は$h$の倍数なので、固定した$h$次数ではその入力数も有界である。

単位を一つのboundary入力へ置く。
その頂点へ辺が入れば、定数の微分を含むので零となる。
辺が入らない場合、積分する微分形式はそのboundary頂点を忘れる写像で
一つ低い次元の配置空間から引き戻される。
$n\ge2$なら忘却後も配置空間は安定であり、最高次元の積分は零になる。
$n=1$の成分はHKR写像そのものなので、正のarityでは単位を微分して零になる。
これがnormalized性を与える。

全ての式を係数環$R[t,t^{-1}]$へ多重線形に拡張する。
$t$を微分する辺は導入しないため、この拡張は$t$の作用を保つ。
\end{proof}

\subsection{二次Artin基底での実際のMF同値}

$S=\CC[x,y,z,u_1,\ldots,u_m]$とし、
\[
 R_2=\CC[h]/(h^3),\qquad
 \xi=\partial_x\wedge\partial_y\wedge\partial_z,\qquad
 \pi=[\xi,F]
\]
と置く。$F,D\in S_3$、$W=F+h^2D$とする。
$\pi$は三つのactive変数のJacobian Poisson tensorであり、
passive変数は係数として扱う。
$\pi=[\xi,F]$という符号規約を用いる。
Kontsevichの星積を$*$、関数のtangent mapを
\[
 C_h(f)=\sum_{j=0}^{2}\frac{h^j}{j!}
       \mathcal U_{j+1}(f,\pi,\ldots,\pi)
\]
と書く。

\begin{theorem}[有限curved gaugeによるgraded MF比較]
\label{cg:thm:mf-comparison}
上のデータに対して、$h=0$で恒等となる$R_2$-linearな同値
\begin{equation}\label{cg:eq:MF-comparison}
 \MF^{\ZZ}\bigl((S\otimes R_2,*),C_h(W)\bigr)
 \simeq
 \MF^{\ZZ}(S\otimes R_2,W)
\end{equation}
が存在する。
これは有限rank対象の$A_\infty$同値から得られるdg Morita同値であり、
内部シフト$(1)$と周期同型$(3)\simeq[2]$を保つ。
この比較自体には$V(F)$の滑らかさを仮定しない。
\end{theorem}
\begin{proof}
周期多重vector場の側で
\[
 \alpha=-tW+h\pi,\qquad \eta=h\xi/t
\]
と置く。
$\alpha$のcohomological degreeは一、$\eta$のそれは零であり、
両者の内部次数は零である。
$[\pi,F]=[\pi,\pi]=0$であり、
$[h\pi,W]=0$は$h^3=0$の下で成立するため、$\alpha$はMaurer--Cartan元である。
$\xi,\pi$の脚は三つのactive変数だけを走るので
$[\xi,\pi]=0$である。従って
\begin{equation}\label{cg:eq:gauge-calculation}
 [\eta,-tW]=-h\pi,\qquad [\eta,h\pi]=0,\qquad
 \exp(\operatorname{ad}\eta)(-tW+h\pi)=-tW.
\end{equation}
ここで$h^3[\xi,D]$の項は零である。

$\alpha_s=-tW+(1-s)h\pi$とすると
$\partial_s\alpha_s=[\eta,\alpha_s]$である。
Kontsevich写像をこの経路へ適用する。
補題\ref{cg:lem:arity}により全ての必要なTaylor和は有限であり、
曲率が$h$の倍数でないことによる収束問題は生じない。
始点のcurved構造は
\[
 b_0\quad\longleftrightarrow\quad (*,-tC_h(W))
\]
であり、終点は
\[
 b_1\quad\longleftrightarrow\quad (\mu,-tW)
\]
である。$\mu$は通常の可換積である。
関数入力を含む項は零項曲率だけであり、
二個以上の関数入力はarityにより消えるので、これ以外の演算は現れない。

formalityによるgauge生成元を
\[
 \gamma_s=\mathcal U^{\alpha_s}(\eta)
 =\sum_{n\ge0}\frac1{n!}
   \mathcal U_{n+1}(\eta,\alpha_s,\ldots,\alpha_s)
\]
と書く。
gaugeの符号は$\partial_s b_s=[\gamma_s,b_s]$となる規約に従う。
補題\ref{cg:lem:arity}により$\gamma_s$のarityは一と三に限られ、
$\gamma_s=O(h)$である。
正のarity成分はnormalizedであり、内部次数と$t$の作用を保つ。

補題\ref{cg:lem:artin-flow}を$N=2$で適用する。
得られる実際のbar同型$\mathcal F$は
\[
 \mathcal F b_0=b_1\mathcal F,\qquad
 F_0=0,\qquad F_1=\Id+O(h),\quad F_3=O(h),\quad F_5=O(h^2)
\]
を満たし、その他のarity成分は零である。
逆写像にも同じ上界がある。
従って補題\ref{cg:lem:tw-extension}を適用できる。
対象の像は特に
\[
 Q'=F_1(Q)+F_3(Q,Q,Q)+F_5(Q,Q,Q,Q,Q)
\]
という有限式で与えられる。
写像と逆写像が有限rank、内部シフト、$t$の作用を保つことから結論を得る。
\end{proof}

\begin{corollary}[指定された中心元との照合]
\label{cg:cor:matching}
次数付き代数同型
$\Phi:A_2\xrightarrow{\sim}(S\otimes R_2,*)$があり、
中心三次元$\Omega_2$について
\[
 \Phi(\Omega_2)=C_h(F)+h^2D,\qquad D\in S_3
\]
を満たすとする。
このとき、内部シフトと三次周期性を保つ同値
\[
 \MF^{\ZZ}(A_2,\Omega_2)
 \simeq \MF^{\ZZ}(S\otimes R_2,F+h^2D)
\]
が存在する。
\end{corollary}
\begin{proof}
$C_h(F+h^2D)=C_h(F)+h^2D$である。
実際、$D$に対する一次HKR項以外にはさらに$h$が掛かるので、
全て$h^3$以上となる。
従って$\Phi$による厳密な同型と
定理\ref{cg:thm:mf-comparison}を合成すればよい。
\end{proof}

\subsection{射影残余模型との比較を区別する}

\begin{remark}[MF比較の適用範囲]
\label{cg:rem:scope}
上の定理は、指定された周囲の代数と中心元から作るMF pairについての主張である。
既存の局所残余模型$\Gamma_R$へ適用するには、別途
\begin{equation}\label{cg:eq:residual-bridge}
 \operatorname{Perf}(\Gamma_R\otimes_R^{\mathbf L}R_2)
 \simeq
 \MF^{\ZZ}(A_R\otimes_R R_2,\Omega_R\bmod h^3)
\end{equation}
というenhancedな比較と、rotationの同定が必要となる。
この比較は、相対windowとstabilizationを構成し、
生成対象、Homの導来底変更、特殊ファイバーでの同定を確認して行う。
通常の有界Orlov同値があるというだけでは、式\eqref{cg:eq:residual-bridge}の証明にならない。
第\ref{lp:sec:main}節の元の底変更命題は正則なtarget baseを対象としているため、
特異なArtin基底を扱う際には、その追加の比較を明示しなければならない。
この比較は第\ref{am:sec:comparison}節で構成する。
定理\ref{am:thm:comparison}が線束生成残余pairとの同定を与え、
命題\ref{am:prop:Gamma-basechange}が既存$\Gamma_R$への適用を保証する。

また有限Artin基底上の同値から、実在する代数的局所基底上の相対smoothness、
非零の複素数$h$への特殊化、あるいは収束する族への代数化は直接には従わない。
これらは別の主張として扱う。
\end{remark}

\begin{example}[Artin基底で絶対特異圏を使えない例]
\label{cg:ex:artin-singularity}
$R=\CC[\varepsilon]/(\varepsilon^2)$、$A=R[x]$、$W=x$とする。
任意の有限自由$\ZZ/2$-graded MF行列$Q$について$Q^2=xI$を$x$で微分すると
\[
 Q(\partial_xQ)+(\partial_xQ)Q=I
\]
となる。従って$\partial_xQ$が恒等射のcontracting homotopyであり、
有限自由MF圏は零である。

一方、$A/(x)=R$の絶対特異圏$D_{\rm sg}(R)$は零ではない。
剰余体$k=R/(\varepsilon)$は
\[
 \cdots\xrightarrow{\varepsilon}R
 \xrightarrow{\varepsilon}R
 \xrightarrow{\varepsilon}R\longrightarrow k\longrightarrow0
\]
という無限分解を持ち、
$\operatorname{Tor}^R_i(k,k)\simeq k$は全ての$i\ge0$で非零である。
従って$k$はperfectでなく、$D_{\rm sg}(R)$で非零である。
これはcubicについての反例ではなく、周囲の環が正則でないときに
有限自由MFと絶対特異圏を同一視できないことを示す例である。
Artin基底では、周囲の$A$上perfectとなる対象を用いた相対特異圏が必要となる。
\end{example}

% One-parameter formal MF comparison and explicit second-order formula.
\section{一変数の形式的MF比較と二次復元公式}
\label{fo:sec:formal-cubic}

第\ref{cft:sec:transport}節では完備局所基底上のMF比較を構成した。
本節では、その一変数形式的変形での具体的な構成を取り出すとともに、
二次の明示復元公式との関係を述べる。
基本となるのは、滑らかな三次式に対する勾配Koszul複体と
Jacobi恒等式である。これらによりPoisson bivectorの定数三重vectorによる
原像を構成し、全ての有限Artin次数で可換MFとの比較を得る。
特定のactive変数や中心一次変数を選ぶ必要はない。
第\ref{kj:sec:second-jet}節の混合族の二次明示計算は、
本節の形式的存在定理とは別の計算的結果として扱う。

formalityとその変形functorへの適用は
\cite[Sections~4.3--4.6, 6]{PKontsevich}の既知の結果である。
以下では、全Hochschild cochainの使用、内部次数、中心元の実際の逆像、
および有限rank行列因子分解への適用を明示する。
この形式的復元定理と同じ一般性を持つ結果の先行研究上の初出は未確認である。

\subsection{入力と結論の範囲}

$d\ge3$とし、全変数の内部次数を一として
\[
 S=\CC[z_1,\ldots,z_d],\qquad
 R_N=\CC[h]/(h^{N+1})
\]
と置く。$h$の内部次数は零であり、$F\in S_3$は
滑らかな射影cubic $V(F)\subset\PP^{d-1}$を定めるとする。

入力は次の性質を持つ対$(A_h,\Omega_h)$である。
\begin{enumerate}
 \item $A_h$は$h$進完備な次数付きunital $\CC[[h]]$-代数であり、
 次数を保つ$h$進加群同定$A_h\simeq S[[h]]$を持つ。
 $h=0$で積は$S$の通常の積になる。
 \item $\Omega_h\in Z(A_h)$は内部次数三で、$\Omega_h\equiv F\pmod h$である。
\end{enumerate}
加群の同定は積の同定ではない。
特に$A_h$の積の各係数がbidifferentialであることは仮定しない。
この入力には、markedな次数付きflat formal変形を、その自由な各次数成分の
基底を持ち上げて同定した場合が含まれる。
以下、$A_N=A_h/(h^{N+1})$、$\Omega_N=\Omega_h\bmod h^{N+1}$と書く。
行列因子分解の圏には第\ref{cg:sec:comparison}節の
有限自由模型とそのpretriangulated・冪等完備化を用いる。

\begin{theorem}[中心三次対の形式的MF復元]
\label{fo:thm:formal-reconstruction}
上の入力に対して、内部次数三の形式的多項式
\[
 W_h=F+\sum_{r\ge1}h^rW_r,\qquad W_r\in S_3
\]
が存在し、各$N\ge0$について内部シフト$(1)$と周期同型$(3)\simeq[2]$を保つ
\begin{equation}\label{fo:eq:MF-equivalence}
 \boxed{\quad
 \MF^{\ZZ}(A_N,\Omega_N)
 \simeq
 \MF^{\ZZ}(S\otimes R_N,W_h\bmod h^{N+1})
 \quad}
\end{equation}
を得る。同値は有限rank対象上の$A_\infty$同値から誘導されるdg Morita同値であり、
$h=0$で恒等である。
形式的データを一度選んでから切断することで、これらの同値は
$R_N\to R_{N-1}$による基底変換と両立するように取れる。
可換側の射影cubicは$R_N$上で相対的に滑らかである。

この定理単独では非零の複素数$h$での同値や収束する族への代数化を結論しない。
Artin底での線束生成残余pairと上のMF pairの同定は、
定理\ref{am:thm:comparison}と命題\ref{am:prop:Gamma-basechange}で別途証明する。
$W_h$の表示はformalityおよび同定の選択に依存し、標準的な方程式とは主張しない。
\end{theorem}

証明は、通常の結合代数の形式変形をPoisson変形へ移す段階と、
その後に曲率を加えてPoisson項を消す段階に分ける。
前段では周期変数$t$を導入しない。
この順序により、中心元の逆像をcohomology classだけでなく実際の関数として求められる。

\subsection{全Hochschild変形と三次Casimirの逆像}

$V=\langle\partial_{z_1},\ldots,\partial_{z_d}\rangle_\CC$を
定数係数の導分の空間とする。

\begin{lemma}[一般の結合変形を扱うformality]
\label{fo:lem:full-formality}
$A_h$の積は、次数付き$\CC[[h]]$-代数同型
\begin{equation}\label{fo:eq:Phi}
 \Phi_h:A_h\xrightarrow{\sim}(S[[h]],*_{\pi_h}),\qquad
 \Phi_h\equiv\Id\pmod h
\end{equation}
によって、formal Poisson tensor
\begin{equation}\label{fo:eq:formal-pi}
 \pi_h\in h\bigl(S_2\otimes_\CC\textstyle\bigwedge^2 V\bigr)[[h]]
\end{equation}
のKontsevich星積へ移せる。
また$W_h\in S_3[[h]]$が一意に存在して
\begin{equation}\label{fo:eq:central-inverse}
 [\pi_h,W_h]=0,\qquad
 \mathcal U^{\pi_h}(W_h)=\Phi_h(\Omega_h),\qquad
 W_h\equiv F\pmod h
\end{equation}
となる。ここで一意性は$\pi_h$、$\Phi_h$とformality morphismを固定した上での一意性である。
\end{lemma}
\begin{proof}
初めにmarkedな加群同定を単位$1$を保つように選ぶ。
結合変形と一入力gaugeを扱う際には、単位を入力すると零になる
normalized Hochschild complexを用いる。
通常のHochschild complexから退化成分を除く標準の縮約は次数を保ち、
Kontsevichの正のarity成分もnormalizedなので、この制限でquasi-isomorphismは保たれる。
degree $-1$の関数成分は、そのまま残す。

多重vector場、polydifferential cochain、全てのHochschild cochainをそれぞれ
\[
 T_{\mathrm{poly}}(S),\qquad
 D_{\mathrm{poly}}(S),\qquad
 C^{\bullet+1}(S,S)
\]
と書く。Kontsevichのaffine formalityと包含の合成
\begin{equation}\label{fo:eq:all-cochain-formality}
 T_{\mathrm{poly}}(S)
 \xrightarrow{\ \mathcal U\ }
 D_{\mathrm{poly}}(S)
 \lhook\joinrel\longrightarrow
 C^{\bullet+1}(S,S)
\end{equation}
は$L_\infty$ quasi-isomorphismである。
最後の包含を省略してはならない。
実際、$S\otimes_\CC S$上の対角のKoszul resolutionに$\Hom(-,S)$を適用すると、
全Hochschild cohomologyは多重vector場となり、HKR反対称化がその同定を与える。
polydifferential側のHKR quasi-isomorphismと比較して、包含もquasi-isomorphismとなる。
これは既存の積が微分作用素で書けることを要求しない。

全ての写像は内部次数を保つ。
bar resolutionとKoszul resolutionの比較は斉次に選べるため、
各内部weightに制限してもHKRはquasi-isomorphismである。
特にweight零の変形をweight零のformal Poisson tensorで表せる。
これは、全cochainを有限weightの直和に属すると仮定する操作ではなく、
斉次入力と出力についてweightを指定した複体の比較である。

$h$進filteredなMaurer--Cartan対応を
\eqref{fo:eq:all-cochain-formality}へ適用する
\cite[Section~4.4, Theorem]{PKontsevich}。
両側の変形は$h$の倍数であり、各有限Artin切断では通常の冪零変形問題となる。
二入力の結合積の変形は$\pi_h$の像とgauge同値になり、
Hochschild degree零のgaugeは一入力の線形同型なので、
これは\eqref{fo:eq:Phi}の厳密な代数同型を与える。
bivectorの内部weightが零であることは、その係数が二次式であることを意味する。
従って\eqref{fo:eq:formal-pi}を得る。

この構成は、既に指定された有限jetの$\pi_h$と$\Phi_h$を保って逐次に行える。
実際、一段持ち上げる際のPoisson側の障害は、その段の$h$係数を除けば
$T_{\mathrm{poly}}^2(S)$の元である。
実際の結合積は次の段へ持ち上がっているため、比較したそのHKR像は
Hochschild coboundaryとなる。
HKRのdegree二での単射性と多重vector場側の零微分から、障害自体が零となる。
残る積の差はHochschild degree一のcocycleであり、
HKRのこのdegreeでの全射性によって、新しいbivector係数と
一入力gaugeのcoboundaryに分けられる。
補正は新しい$h$次数だけに置けるので、既指定jetを変えずに帰納できる。
normalized複体を使っているため、このgaugeは単位を保つ。

次に$\pi_h$でtwistしたformality morphismを考える。
$h$進filtrationの次数零部分は上記のHKR比較なので、
twist後もquasi-isomorphismである。
周期変数をまだ加えていないため、shifted degree $-1$のcohomologyは
それぞれPoisson Casimirと星積の中心であり、degree $-2$以下の項は存在しない。
従って、このdegreeでのcohomologyの同型は
\[
 \mathcal U^{\pi_h}:
 \{f:[\pi_h,f]=0\}\xrightarrow{\sim} Z(S[[h]],*_{\pi_h})
\]
という実際の関数の同型である。
中心元$\Phi_h(\Omega_h)$の逆像を$W_h$とする。
$\pi_h$のweightが零なので、twisted differentialもtangent mapも各weightを保つ。
ここでは変形のweight零部分とは別に、中心のweight三部分を使う。
そのquasi-isomorphismから$W_h\in S_3[[h]]$を得る。
一次成分のHKR写像は関数上では恒等なので$W_h\equiv F\pmod h$となる。
ここでは、関数の同定を高次cochainを法とする同値類で置き換えていない。
\end{proof}

\subsection{三次式のKoszul原像と二回目の作用の消滅}

多重vector場のSchouten次数を、rank $r$に対して$r-1$とする。
以下の括弧は$[\partial_x\wedge\partial_y\wedge\partial_z,W]$
が通常のJacobian bivectorになる規約である。

\begin{lemma}[形式的勾配Koszul complexと係数次数]
\label{fo:lem:regular-gradient}
$F$が滑らかな射影cubicを定め、$W_h\equiv F\pmod h$なら、
$g=(\partial_1W_h,\ldots,\partial_dW_h)$によるKoszul contraction
\[
 \delta_{W_h}=\iota_{dW_h}:
 S[[h]]\otimes\textstyle\bigwedge^r V
 \longrightarrow
 S[[h]]\otimes\textstyle\bigwedge^{r-1}V
\]
は正のrankで完全である。
係数が斉次次数$q$のcycleの原像は、係数次数$q-2$に取れる。
特に負の係数次数しか許されない原像は零である。
同じ完全性は各有限Artin切断でも成立する。
\end{lemma}
\begin{proof}
全$d$個の偏微分$\partial_iF$の共通affine零点は原点だけなので、
それらが生成するイデアルの高さは$d$である。
多項式環はCohen--Macaulayであるから、全偏微分は正則列となる。
従って$h=0$のKoszul complexは正のrankで完全である。

各$R_N$上のKoszul complexを有限の$h$進filtrationで濾過する。
各項は$R_N$上自由なので、associated graded complexは
特殊fiberの勾配Koszul complexと$h^jR_N/h^{j+1}R_N$とのtensorの直和である。
そのhomologyはrank零にしか存在しない。
有限filtrationのspectral sequenceから、元のcomplexも正のrankで完全となる。
形式的なcomplexについても、各内部次数では有限自由な$\CC[[h]]$-加群の有限complexであり、
完備な$h$進filtrationまたは両立するcycleの逐次持上げにより同じ完全性を得る。
各$g_i$は内部次数二なので、Koszul微分は係数次数を二つ増やしrankを一つ減らす。
従ってKoszul complexは「係数次数$+2\,\mathrm{rank}$」でgradedであり、
この斉次部分ごとに完全性を適用すれば、原像の係数次数は$q-2$となる。
\end{proof}

\begin{lemma}[定数三重vectorによる二次Poissonの消去]
\label{fo:lem:constant-trivector}
$W_h\in S_3[[h]]$が上の条件を満たし、
$\pi_h\in h(S_2\otimes\bigwedge^2V)[[h]]$が
$[\pi_h,W_h]=[\pi_h,\pi_h]=0$を満たすとする。
このとき一意な定数係数三重vector
\[
 \Xi_h\in h\bigl(\textstyle\bigwedge^3V\bigr)[[h]]
\]
が存在し、
\begin{equation}\label{fo:eq:constant-trivector}
 \pi_h=[\Xi_h,W_h],\qquad [\Xi_h,\pi_h]=0
\end{equation}
となる。$\Xi_h$の分解可能性は仮定しない。
\end{lemma}
\begin{proof}
Casimir条件は、rank二、係数次数二の$\pi_h$が勾配Koszul cycleであることを意味する。
補題\ref{fo:lem:regular-gradient}により、rank三、係数次数零の原像$\Xi_h$が存在する。
従ってその係数は変数に依存しない。
二つの原像の差はrank三、係数次数零のcycleであり、
そのKoszul原像はrank四、係数次数$-2$でなければならない。
そのような多項式はないので$\Xi_h$は一意である。
$\pi_h\bmod h=0$に同じ一意性を適用すると$\Xi_h\bmod h=0$を得る。

$D_h=[\Xi_h,\pi_h]$と置く。
$\Xi_h$は定数係数で$\pi_h$の係数は二次式だから、
$D_h$はrank四、係数次数一である。
Schouten次数は$|W_h|=-1$、$|\Xi_h|=2$、$|\pi_h|=1$なので、graded Jacobiにより
\begin{align*}
 [W_h,D_h]
 &=[[W_h,\Xi_h],\pi_h]+[\Xi_h,[W_h,\pi_h]]\\
 &=-[\pi_h,\pi_h]=0.
\end{align*}
従って$D_h$も勾配Koszul cycleである。
その原像はrank五、係数次数$-1$でなければならず、再び零である。
よって$D_h=0$を得る。
rankが変数の数を超える場合も、その外積空間自体が零なので同じ結論となる。
\end{proof}

\begin{remark}[三次という次数を使う箇所]
補題\ref{fo:lem:constant-trivector}の二つの次数計算が本節の要点である。
最初のKoszul持上げでは「二次のbivectorから定数のtrivector」を得て、
次の持上げでは「一次係数のrank四tensorには多項式の原像がない」と結論する。
滑らかさは勾配Koszul complexの完全性に用いる。
従って、中心元が別の次数であったり、古典cubicが特異であったりする場合に、
同じ結論をそのまま適用してはいけない。
\end{remark}

\subsection{全ての有限Artin次数でのcurved比較}

\begin{proof}[定理\ref{fo:thm:formal-reconstruction}の証明]
補題\ref{fo:lem:full-formality}によって$\Phi_h,\pi_h,W_h$を形式的に選び、
補題\ref{fo:lem:constant-trivector}の$\Xi_h$を取る。
周期変数$t$に式\eqref{cg:eq:bidegree}の次数$(2,-3)$を与え、
curved Maurer--Cartan元とgauge生成元を
\[
 \alpha_h=-tW_h+\pi_h,\qquad \eta_h=\Xi_h/t
\]
と置く。$\Xi_h$はSchouten次数二、内部次数$-3$なので、
$\eta_h$はcohomological degree零、内部次数零である。
また$\eta_h$は$h$の倍数であり、
\[
 [\eta_h,-tW_h]=-\pi_h,\qquad [\eta_h,\pi_h]=0
\]
となる。従って全ての次数で
\begin{equation}\label{fo:eq:all-order-gauge}
 \boxed{\quad
 \exp(\operatorname{ad}\eta_h)(-tW_h+\pi_h)=-tW_h.
 \quad}
\end{equation}

この等式を$R_N$へ切断する。
補題\ref{cg:lem:arity}のarity計数は$N$にも変数の数にも依存しない。
bivectorの各入力が$h$の倍数であり、関数としての曲率入力は高々一個なので、
Kontsevich写像の必要なTaylor和は全て有限である。
式\eqref{fo:eq:central-inverse}から、$\alpha_h$の像は
\[
 \bigl(*_{\pi_h},-t\Phi_h(\Omega_h)\bigr)
\]
というcurved associative構造であり、右辺$-tW_h$の像は
通常の可換積と曲率$-tW_h$である。

gauge経路を$\alpha_s=-tW_h+(1-s)\pi_h$と取る。
そのKontsevich像のgauge生成元は
$\mathcal U^{\alpha_s}(\eta_h)$であり、符号は
$\partial_s\alpha_s=[\eta_h,\alpha_s]$に従って定める。
一個のtrivectorを含むTaylor成分のarityは、関数入力が零個なら三、
一個なら一である。
従ってこの生成元はarity一と三に限られ、$h$の倍数であり、
各$R_N$上で$s$の多項式となる。
正のarityの成分はnormalizedである。
補題\ref{cg:lem:artin-flow}から、対応するbar同型とその逆は
\[
 F_1=\Id+O(h),\qquad
 F_{2r+1}=O(h^r),\qquad F_k=0\ (k>2N+1)
\]
を満たす有限式で与えられる。
全成分は内部次数零であり、$t$について線形である。

補題\ref{cg:lem:tw-extension}をこのbar同型へ適用する。
有限rank行列因子分解のtwisting行列$Q$は$h$の倍数である必要はない。
像$\sum_kF_k(Q,\ldots,Q)$と射の挿入和は上のarity上界によって有限だからである。
得られる同値は有限rankを保ち、内部シフトと周期同型を保つ。
これに厳密な次数付き代数同型$\Phi_h\bmod h^{N+1}$を合成すると
式\eqref{fo:eq:MF-equivalence}を得る。

最後に、$\pi_h,\Phi_h,W_h,\Xi_h$は初めに一つの形式的データとして選んでいる。
gauge経路とその順序付き積分も同じ形式的式の切断であるから、
構成した同値と逆同値は$R_N\to R_{N-1}$による基底変換と両立する。
可換超曲面は滑らかな特殊fiberを持つflatな冪零変形なので、
各$R_N$上で相対的に滑らかである。
\end{proof}

\begin{corollary}[次数別完備基底上の有限rank比較]
\label{fo:cor:complete-graded}
$\widehat R=\CC[[h]]$とし、$A_h$の各斉次成分を保った代数
$A_h^{\mathrm{gr}}=\bigoplus_{j\ge0}(A_h)_j$を用いる。
定理\ref{fo:thm:formal-reconstruction}の同値は、$\widehat R$上でも
\[
 \MF^{\ZZ}(A_h^{\mathrm{gr}},\Omega_h)
 \simeq
 \MF^{\ZZ}(S\otimes_\CC\widehat R,W_h)
\]
という有限rank graded MFの$A_\infty$同値、従ってその包絡・冪等完備化のdg Morita同値を与える。
これは内部シフトと周期同型を保ち、各有限Artin同値へ還元される。
ここで全代数の次数を忘れた完備化や、非零の複素数への$h$の特殊化は主張しない。
\end{corollary}
\begin{proof}
固定した有限rank graded free moduleと固定したcohomological degreeについて、
その行列の各成分は内部次数が固定された
$(A_h)_j\simeq S_j\otimes_\CC\widehat R$に属する。
周期変数$t$の冪もcohomological degreeから固定されるため、
これらは有限個の有限自由$\widehat R$-加群の直和であり、$h$進完備かつ分離である。
定理の構成によるbar写像とその逆は形式的に
$F_{2r+1}=O(h^r)$、$F_{2r}=F_0=0$を満たす。
従って$Q$が$h$の倍数でなくても
$Q'=\sum_{r\ge0}F_{2r+1}(Q,\ldots,Q)$は上の固定された成分内で収束する。
固定個数の射に$Q$を挿入するfunctorの式も、
arityが増えると$h$進位数が増えるため同様に収束する。
これは各行列成分の固定内部次数内での極限であり、
変数の無制限な次数の級数を導入しない。
得られた対象がMF方程式を満たすこと、$A_\infty$恒等式、
および逆写像との合成が恒等であることは、
各$h^{N+1}$を法として既に成立している。
成分の分離性から、全ての恒等式が$\widehat R$上でも成立する。
従って有限rank対象上の実際の互いに逆な$A_\infty$ functorを得る。
内部次数、$t$の作用、単位も各切断で保存されるため極限で保存される。
\end{proof}

\begin{remark}[pure Poisson量子化のcanonicalな中心元]
\label{fo:rem:canonical-center}
$\pi$が二次Poisson tensorで$[\pi,F]=0$なら、
星積$*_{h\pi}$とcanonicalな中心元$\mathcal U^{h\pi}(F)$の対は、
全ての有限Artin次数で元の可換式$F$のMF pairと同値である。
実際、補題\ref{fo:lem:constant-trivector}で
$\pi=[\Xi,F]$、$[\Xi,\pi]=0$となる定数$\Xi$を取り、
$\eta=h\Xi/t$による同じcurved gaugeを適用すればよい。

一般に与えられた対$(A_h,\Omega_h)$では、formalityで対応する
$\pi_h$にも中心元の逆像$W_h$にも高次補正があり得る。
最初のPoisson tensorだけから$W_h$を決めることはできない。
例えば混合試験族の対称二次関係とpure Kontsevich星積$*_{h\pi}$
との生成元による同定は、第\ref{kj:sec:second-jet}節では$h^3$を法として証明したものであり、
その同定を全次数で同じ形に保てるとは主張していない。
\end{remark}

\subsection{対称二次量子化の明示的二次復元公式}

\begin{corollary}[一般の対称二次関係からの二次式]
\label{fo:cor:universal-second-jet}
$R_2=\CC[h]/(h^3)$とし、$F\in S_3$は滑らかな射影cubicを定めるとする。
$\pi=\frac12\pi^{ij}\partial_i\wedge\partial_j$は係数が二次式のPoisson tensorで、
$[\pi,F]=0$とする。
次数付き$R_2$-代数$A_2$は
\[
 [z_i,z_j]=h\,\sym_{A_2}(\pi^{ij})
\]
という対称二次関係を持ち、$S\otimes R_2$と同じ次数別自由加群を持つPBW変形とする。
中心三次元を
\[
 \Omega_2=\sym_{A_2}(F+h^2C),\qquad C\in S_3
\]
と仮定する。この表示には$h$の一次項を含めない。
このとき、有限の偏微分と積だけで計算できる三次式
\begin{equation}\label{fo:eq:universal-second-jet}
 \boxed{\begin{aligned}
 W_2&=F+h^2\left(C-\frac1{48}\sum_{i,k}Q_{ik}F_{ik}\right),\\
 Q_{ik}&=\sum_{j,l}(\partial_l\pi^{ij})(\partial_j\pi^{kl}).
 \end{aligned}}
\end{equation}
に対し、内部シフトを保つ同値
\[
 \MF^{\ZZ}(A_2,\Omega_2)\simeq
 \MF^{\ZZ}(S\otimes R_2,W_2)
\]
が存在する。この系には、入力の二次jetが全次数へ延長できることを仮定しない。
\end{corollary}
\begin{proof}
角度型Kontsevich星積の二次作用素は対称なので、
一次生成元の交換子は$h\pi^{ij}$である。
対称化された右辺の補正は$h^3$以上となり消えるため、
生成元を同名の一次式へ送る写像は$A_2$から星積への代数準同型を定める。
PBW性と$h=0$での恒等性によりこれは同型である。
補題\ref{kj:lem:symmetrization}は変数の数によらないので、
その同型は$\Omega_2$を式\eqref{fo:eq:universal-second-jet}の$W_2$へ送る。
補題\ref{kj:lem:curvature-wheel}により
$\mathcal U^{h\pi}(W_2)=W_2$でもある。

$[h\pi,W_2]=0$は$h^3=0$の下で成り立つ。
補題\ref{fo:lem:regular-gradient}と
\ref{fo:lem:constant-trivector}の有限Artin版を用いると、
$h\pi=[\Xi_2,W_2]$、$[\Xi_2,h\pi]=0$となる
$\Xi_2\in h(\bigwedge^3 V)\otimes R_2$を得る。
定理のcurved gaugeの証明をこの切断で直接行い、
補題\ref{cg:lem:artin-flow}と\ref{cg:lem:tw-extension}を適用すればよい。
\end{proof}

この式は、入力の二次Poisson tensorとWeyl中心式が与えられている場合の、
全変数に共通する明示的二次復元公式である。
三変数Jacobian型では
$\sum Q_{ik}F_{ik}=-6\det(\mathsf H)$となる。
さらに$C=-\det(\mathsf H)/12$なら、括弧内は$\det(\mathsf H)/24$となり、
第\ref{kj:sec:second-jet}節の混合族の式を再現する。

\subsection{混合試験族の二次式と射影残余への適用}

第\ref{mt:sec:family}節の混合試験族は定理の入力を満たす。
全Artin切断のPBW基底は両立しているので、逆極限を取ると
$A_h\simeq\CC[x,y,z,u,v][[h]]$という次数付き加群同定を得る。
中心元の閉公式も同じ逆極限で定義され、
$\lambda^3\ne-3$なら古典cubicは滑らかである。

第\ref{kj:sec:second-jet}節は、この族について
\[
 \widehat F_\lambda
 =F_\lambda+
 h^2\left(\frac{xyz}{3}-\frac{\lambda^2u^2z}{12}\right)
 \pmod{h^3}
\]
との実際のMF比較を与える。
これは単なる中心元の可換化から得た式ではなく、
Kontsevich星積とCasimirの規格化を指定した上での二次式である。
本節の形式的MC持上げは、与えられた有限jetの同定に合わせて選ぶことができる。
そのように選んだ$W_h$の二次Jacobian classは
$\bigl(1/3+\lambda^3/36\bigr)[xyz]$である。
高次の$W_r$は本節で計算していない。
それらを求めるには、formal Poisson tensorと代数同型を各次数で求め、
中心tangent map\eqref{fo:eq:central-inverse}を逆に解く必要がある。
本節はその形式的可解性を示すが、全次数の閉公式や収束を与えない。

定理\ref{am:thm:comparison}により、式\eqref{fo:eq:MF-equivalence}は
各有限Artin切断で線束生成残余pairの同値へ移る。
可換側は相対的に滑らかな$V(W_h)$のperfect Kuznetsov成分である。
既存の正則局所模型$\Gamma_R$からの底変更は
命題\ref{am:prop:Gamma-basechange}で確認する。
この比較の合成を系\ref{am:cor:formal-residual}に記載する。
非可換族の非零ファイバーのAS性、非零$h$での復元、および代数化は、
この有限Artin同値の系には含めない。

% Current finite-Artin comparison; extracted with complete proofs.
\section{Artin基底における線束生成残余圏と行列因子分解}
\label{am:sec:comparison}

本節では絶対特異圏を用いず、周囲の代数上perfectな複体だけを用いて、
有限階数graded MFと線束から生成される射影残余圏を比較する。
ここで構成するのは射影残余圏とMF圏の比較であり、
非可換MF圏と可換MF圏の間の別の同値の存在は仮定しない。
非正則Artin基底上で通常の有界coherent圏全体へOrlovの定理を適用することはしない。
相対的なstabilizationの加群論的入力は
Chen\cite[Theorem 2.10, Corollary 2.11, Theorem 7.2]{AMChen}である。
同論文の定理は三角圏についてのものであるため、以下ではexact圏とdg商を通じて
実際のenhanced関手を構成する。
本節の組合せと同じ一般性での先行研究上の初出は未確認である。

\subsection{入力と相対的な特異圏}

変数の数を$d\ge5$、Artin切断の次数を$N\ge0$とし、
\[
 R_N=\CC[h]/(h^{N+1}),\qquad k=R_N/(h),\qquad
 a=d-3,\qquad i=1-a=4-d
\]
と置く。$A=\bigoplus_{j\ge0}A_j$は中央$R_N$-代数で、$A_0=R_N$、
各$A_j$は有限自由$R_N$-加群、
$\overline A=A\otimes_{R_N}k\simeq k[z_1,\ldots,z_d]$と仮定する。
$\Omega\in Z(A)_3$の剰余$F$は滑らかな射影cubicを定め、$B=A/(\Omega)$とする。
以下の構成は内部次数を保つ。本節のMFと厚いdg部分圏は冪等完備化して扱う。

この入力から、左右noetherian性、$R_N$の左右$A$-加群としての有限自由Koszul型分解、
$\Omega$の正則性、$B_j$の自由性が従う。
これらは正則基底上の有限global dimensionと区別する。
実際、$h$による有限filtrationの随伴次数環がnoetherianなので$A,B$はnoetherianである。
この持上げは、基底の正則性ではなく次数別自由性を使う。
剰余の自由分解の零次項の生成元の像を同じ次数に持ち上げると、
次数別Nakayamaにより全射を得る。
各次数で有限自由加群への全射なので分裂し、核も次数別有限自由である。
この核の剰余を正しく計算できるため、次の自由項を同様に持ち上げて帰納する。
最後の核は剰余が零であり、再び次数別Nakayamaにより零となる。
従って多項式環のKoszul分解から
\begin{equation}\label{am:eq:relative-koszul}
 P_j=A(-j)^{\binom dj}\quad(0\le j\le d)
\end{equation}
という分解を得る。各内部次数で$R_N$-分裂し、全てのArtin商への底変更後も完全である。
$\Omega:A_{j-3}\to A_j$は剰余上で単射なので、有限自由加群の行列の単元小行列式を用いて
$R_N$上でも分裂単射となる。従って$B_j$は自由である。
この議論から$\operatorname{gldim}A<\infty$は結論しない。

\begin{definition}[周囲でperfectな特異圏]\label{am:def:relative-singularity}
標準dg enhancement内で
\[
 \mathcal D_{B/A}
 =\{M\in D^b(\operatorname{gr}B):
       \operatorname{Res}_A M\in\operatorname{Perf}(\operatorname{gr}A)\}
\]
を取り、
\[
 \mathcal S_{B/A}
 =\bigl(\mathcal D_{B/A}/\operatorname{Perf}(\operatorname{gr}B)\bigr)^{\mathrm{idem}}
\]
をそのdg商の冪等完備化とする。
$\operatorname{Perf}(\operatorname{gr}B)\subset\mathcal D_{B/A}$は
$0\to A(-3)\xrightarrow{\Omega}A\to B\to0$から従う。
\end{definition}

ここでのperfect性は周囲の\emph{$A$上}のperfect性である。
$B(j)$やその内部切断は、$R_N$上では一般に無限階数であり、
$\operatorname{Perf}(R_N)$の対象だとは述べない。
また$\mathcal S_{B/A}$を絶対特異圏$D_{\rm sg}(B)$と同一視しない。

\begin{lemma}[相対dg stabilization]\label{am:lem:stabilization}
上の入力に対し、内部shiftと三次周期性を保つdg Morita同値
\begin{equation}\label{am:eq:stabilization}
 \operatorname{cok}:
 \MF^{\ZZ}(A,\Omega)\xrightarrow{\sim}\mathcal S_{B/A}
\end{equation}
が存在する。そのMorita逆を$\operatorname{stab}$と書く。
この同値は有限自由MFの余核とその周期自由分解から誘導される。
Artin商への底変更に関して、これらの構成は自然である。
\end{lemma}

\begin{proof}
まず、次数付き有限射影$A$-加群は次数付き有限自由である。
実際、そのaugmentation商は局所環$R_N$上の次数付き有限射影加群なので自由であり、
その基底を持ち上げた自由加群からの写像に次数付きNakayamaを適用すればよい。

$\mathcal E$を、有限生成次数付きGorenstein-projective $B$-加群のうち
$A$上有限射影次元を持つもののexact圏とする。
Chenのchange-of-rings定理により、非零のこのような加群は
$A$上の射影次元が一である。
正則中心元の場合の同論文の証明は、全ての射と分解を斉次に取れば次数付き圏でも成立する。
特に$\mathcal E$はFrobenius exact圏であり、そのprojective-injective対象は
有限次数付き射影$B$-加群である。
syzygyとcosyzygyで有限$A$-射影次元が保たれること、およびその次元が再び一になることが、
このFrobenius部分圏への制限を保証する。

$\mathcal F$を有限自由graded MFのexact圏とする。
exact列は二つの成分で$A$-分裂する列であり、projective-injective対象は
$(1,\Omega)$、$(\Omega,1)$型の可縮factorizationの直和因子である。
余核はexact関手
\[
 \operatorname{cok}:\mathcal F\longrightarrow\mathcal E
\]
を与える。余核加群の$A$-自由分解の長さは一である。
逆に$M\in\mathcal E$の斉次分解
$0\to P_1\xrightarrow{\varphi}P_0\to M\to0$を取ると、
$\Omega M=0$より$\Omega\Id_{P_0}$は$\varphi$を通って一意にfactorする。
$\varphi$の単射性から他方の積も$\Omega\Id_{P_1}$となり、MFを得る。
射の持上げとhomotopyの同定を含め、これはChenのTheorem 7.2の余核同値である。
従って余核関手はstable categoryの三角同値を与える。

この三角同値だけでdg同値を宣言せず、次のdg商を取る。
exact関手から実際のdg関手
\begin{equation}\label{am:eq:exact-dg-quotients}
 \frac{D^b_{\rm dg}(\mathcal F)}{K^b_{\rm dg}(\operatorname{proj\!-\!inj}\mathcal F)}
 \longrightarrow
 \frac{D^b_{\rm dg}(\mathcal E)}{K^b_{\rm dg}(\operatorname{proj}B)}
\end{equation}
を得る。Frobenius exact圏の右辺・左辺の$H^0$は対応するstable categoryである。
これはprojective-injective分解を用い、有限複体をその高いsyzygyへ移すことで確認できる。
従って\eqref{am:eq:exact-dg-quotients}は$H^0$で同値であり、
shiftを含む任意の対象対へ同じ事実を適用すると全Hom cohomologyで同型となる。

左側のdg商から元のMF dg圏へは、有限複体の有限totalizationを取る。
成分別split exactな複体のtotalizationは可縮である。
具体的に、横微分$d$と横contraction$s$を
$ds+sd=1$、$s^2=0$とし、符号を含む縦MF微分を$\delta$とする。
$d\delta+\delta d=0$、$\delta^2=w$で、$w$は中央曲率である。
$s$は横次数を下げるので、横方向の有限性から
$1+\delta s$、$1+s\delta$は有限和で逆を持つ。
\[
 q=s(1+\delta s)^{-1}=(1+s\delta)^{-1}s
\]
と置けば、直接展開により
\[
 (1+s\delta)\bigl(1-(d+\delta)q-q(d+\delta)\bigr)
 (1+\delta s)=-sws=0
\]
となる。最後は中央性と$s^2=0$を用いた。
従って$(d+\delta)q+q(d+\delta)=1$であり、
totalizationは実際に可縮である。この計算は曲率の冪零性を必要としない。
可縮MFも零へ送られるため、totalizationは左側のdg商を通る。
$H^0$では通常のMF homotopy categoryへの同値であり、同様に全Homで準同型となる。

右側から$\mathcal S_{B/A}$へのdg関手は、$\mathcal E$の$B$-加群圏への包含から得る。
完全射影分解を用いると、Gorenstein-projective加群間のstable Homは特異圏のHomに一致する。
同じ分解のshiftを用いることで全Hom cohomologyについて一致する。
本質的全射性も具体的に確認できる。
$M\in\mathcal D_{B/A}$の有界上の有限生成$B$-射影分解を十分低い位置で切断すると、
特異圏で$M$はある$B$-加群$L$のshiftに同型となり、$L$は$A$上有限射影次元を持つ。
ChenのTheorem 2.10より$L$は$B$上有限Gorenstein-projective次元を持つので、
さらに高い$B$-syzygyを取れば$\mathcal E$に属する。
これで右側のdg商と$\mathcal S_{B/A}$の同値を得る。
三つの実際のdg関手のzigzagにより\eqref{am:eq:stabilization}が得られる。

最後に底変更を確認する。
MFの余核$M$の各内部次数の分解は有限自由$R_N$-加群間の列である。
剰余上で最初の射は単射である。実際、二つのfactorization mapの積は
非零多項式$F$の乗法なので、どちらのmapも単射である。
単元小行列式により各次数の余核$M_j$は$R_N$-自由となる。
従って余核、exact列、有限totalizationはArtin商への底変更と両立する。
周期分解とdg商から得た比較は、Morita localizationの普遍性を通じて自然となる。
選択したcontractionそのものが全ての底変更と厳密に可換することは必要としない。
内部shiftと周期性は全ての段階で保たれ、Morita逆もこれらと両立するように取れる。
\end{proof}

\subsection{Artin基底上の線束Homと有限window}

$\pi:\operatorname{Gr}B\to\operatorname{QGr}B$を射影商とし、
$\mathcal O(j)=\pi B(j)$と置く。
導来sectionを$\mathbf R\omega$と書く。
この段階では$D(\operatorname{QGr}B)$内で線束から生成するdg部分圏を用い、
全ての有界coherent複体についての底変更定理は要求しない。

\begin{lemma}[線束のsectionと底変更]\label{am:lem:line-bundles}
自然なderived Homの底変更は各Artin商と両立し、
\begin{equation}\label{am:eq:line-hom}
 \operatorname{Ext}^q(\mathcal O(r),\mathcal O(s))=
 \begin{cases}
 B_{s-r}&q=0,\\
 B_{r-s-a}^{\vee}&q=d-2,\\
 0&\text{その他}
 \end{cases}
\end{equation}
である。ここで$(-)^\vee=\operatorname{Hom}_{R_N}(-,R_N)$、負添字は零とする。
また$0\le j\le d-1$について自然に
\begin{equation}\label{am:eq:window-sections}
 (\mathbf R\omega\mathcal O(j))_{\ge i}\simeq B(j)_{\ge i}
\end{equation}
となり、この同定もArtin商への底変更と両立する。
\end{lemma}

\begin{proof}
式\eqref{am:eq:relative-koszul}を$A$で双対化する。
各内部次数は有界有限自由$R_N$-複体であり、剰余は通常のKoszul双対性を満たす。
単元成分の消去によって
\[
 \mathbf R\underline{\operatorname{Hom}}_A(R_N,A)
 \simeq R_N(d)[-d]
\]
となる。中心正則元のchange-of-ringsから
\[
 \mathbf R\underline{\operatorname{Hom}}_B(R_N,B)
 \simeq R_N(a)[-(d-1)]
\]
を得る。この二つの計算は基底の有限global dimensionを用いない。

標準の導来section--torsion三角を用いる。
ここで必要な$\chi$はArtin基底でも確認できる。
$N>0$では$B/(h)$の$B$-自由分解は、中央の$h,h^N$による周期分解である。
有限生成$B$-加群$M$について、そのHom cohomologyは有限生成$B/(h)$-加群となる。
change-of-rings spectral sequenceと可換noetherian環$B/(h)$の$\chi$から
$\operatorname{Ext}_B^q(k,M)$は各$q$で内部次数について有限となる。
$R_N$の有限な$k$-filtrationにより、$\operatorname{Ext}_B^q(R_N,M)$も同じ性質を持つ。
$N=0$では通常の可換計算である。
従ってArtin--Zhangのsectionとlocal cohomologyの枠組を適用できる
\cite{LP-AZ}。

$M_\ell=B/B_{\ge\ell}$の内部次数filtrationの商は
$B_j\otimes_{R_N}R_N(-j)$である。
直前のrelative Gorenstein計算と$B_j$の自由性から
\[
 H^q(\mathbf R\tau B)_t=
 \begin{cases}
 B_{-t-a}^{\vee}&q=d-1,\\
 0&q\ne d-1
 \end{cases}
\]
を得る。固定した$t$では$j=-t-a$の一層のみが寄与し、$\ell>j$で安定する。
この計算では、次数filtrationの各商の導来Homがcohomological degree
$d-1$と内部次数$-a-j$に集中することを用いた。
三角$\mathbf R\tau B\to B\to\mathbf R\omega\pi B$から
\eqref{am:eq:line-hom}を得る。

自然な底変更写像についても確認する。
$M_\ell$の有限階数$B$-自由分解を、syzygyが各内部次数で$R_N$-自由となるように取る。
固定した内部次数でのHom複体は、各項が有限自由$R_N$-加群である。
その剰余はcohomologyが一つの次数にしかない複体であり、
体上のsplit contractionをnilpotent perturbationによって持ち上げると、
有限自由cohomologyの一項複体への実際のhomotopy同値を得る。
したがってこの複体はK-flatであり、通常のtensorが導来tensorを計算する。
項ごとのHomとtensorの同定、filtered colimit、上のsection三角を通じて
自然な底変更写像そのものが準同型であることが従う。
filtered colimitはK-flat複体をK-flat複体へ送るので、ここでunboundedな自由複体の
K-flat性を無条件に用いてはいない。

最後に$\mathbf R\omega\mathcal O(j)$の正のcohomologyは内部次数
$t\le-j-a$にしかない。
$i=1-a>-j-a$なので、$j\ge0$では内部切断後に全て消える。
次数零のsectionは$B(j)$であるから\eqref{am:eq:window-sections}を得る。
\end{proof}

\begin{lemma}[線束生成部分の有限模型]\label{am:lem:linear-category}
\[
 E=\bigoplus_{j=0}^{d-1}\mathcal O(j),\qquad
 \mathcal C_B=\operatorname{thick}_{R_N}\{\mathcal O(j):j\in\ZZ\}
\]
とすると$\mathcal C_B=\operatorname{thick}_{R_N}(E)$である。
全ての対象対のHomは$R_N$-perfectであり、Artin商への底変更と両立する。
例外成分$\langle\mathcal O,\ldots,\mathcal O(a-1)\rangle$を除く右直交を
\[
 \mathcal K_B^{\rm lin}
 =\mathcal C_B\cap\langle\mathcal O,\ldots,\mathcal O(a-1)\rangle^\perp
\]
とすると、有限mutationの射影$p$により$G=p(E)$がこれを生成する。
\end{lemma}

\begin{proof}
式\eqref{am:eq:relative-koszul}を$B$へ導来tensorする。
そのcohomologyは$R_N$と$R_N(-3)$のshiftだけであり、いずれも射影商で零となる。
従って線束の長さ$d$のKoszul型複体は射影商で完全である。
両端の係数の階数が一であるため、任意の連続した$d$個のtwistから全てのtwistを
前後に生成できる。
式\eqref{am:eq:line-hom}と有限cone・直和因子によりproper性と底変更が従う。
同式から$\mathcal O,\ldots,\mathcal O(a-1)$は$R_N$-例外列であり、
Homおよびその$R_N$-双対を用いて左右随伴が作れる。
従って有限個のevaluation coneが射影$p$を与え、$p(E)$が右直交を生成する。
\end{proof}

\subsection{MFのHomのproper性と生成子の持上げ}

\begin{lemma}[nilpotent perturbationによるMFの底変更]\label{am:lem:mf-hom}
有限自由graded MFの任意の二対象$P,Q$に対し、
$\operatorname{Hom}_{\MF}(P,Q)$は$R_N$-perfectである。
自然なHomの導来底変更は通常のMFの底変更Homと準同型である。
\end{lemma}

\begin{proof}
周期変数$t$にcohomological degree $2$、内部次数$-3$を与え、
代数$A$をcohomological degree零に置く。
この二重次数で内部次数零のMF Homを書けば、
各cohomological degreeは有限個の$A_j$の直和であり、有限自由$R_N$-加群である。
負のcohomological方向では必要な内部次数が負となるため零である。
剰余のHom複体$C_0$は、滑らかな古典cubicのgraded MFにおけるHomであり、
通常のOrlov同値により有限次元の有界cohomologyを持つ。

体上で$C_0$から有限cohomology複体$H$へのsplit contractionを選ぶ。
各次数の基底を持ち上げて、元のHom微分を$d_0+\delta$、$\delta\in h\operatorname{End}$と書く。
contractionのhomotopyを$s$とすると、
$(1-\delta s)^{-1}$および$(1-s\delta)^{-1}$は$h^{N+1}=0$により有限和である。
通常のperturbation恒等式をこの有限和に適用すると、
元のHom複体と$H\otimes_kR_N$上の変形微分との間のhomotopy同値を得る。
後者は有界有限自由複体である。
この同値からK-flat性が従い、項ごとのMF Homの底変更が導来底変更と一致する。
\end{proof}

\begin{lemma}[Artin基底上の生成子の持上げ]\label{am:lem:generator}
$P\in\MF^{\ZZ}(A,\Omega)$の剰余$\overline P$が古典graded MF圏を生成するとする。
このとき$P$は$\MF^{\ZZ}(A,\Omega)$を生成する。
\end{lemma}

\begin{proof}
補題\ref{am:lem:mf-hom}により、カテゴリーの底変更から古典MFへの関手は
Hom上で完全忠実である。像は$\overline P$を含むので、冪等完備化後にはMorita同値となる。
生成性を正確に持ち上げるため、$\mathcal D=\operatorname{Ind}(\MF^{\ZZ}(A,\Omega))$へ移る。
$P$はcompactであり、$\operatorname{RHom}(P,-)$はcolimitと可換である。
従って$M\in\mathcal D$が$P$と右直交なら
\[
 \operatorname{RHom}(\overline P,M\otimes_{R_N}^{\mathbf L}k)
 \simeq\operatorname{RHom}(P,M)\otimes_{R_N}^{\mathbf L}k=0.
\]
古典fiberでの生成性から$M\otimes_{R_N}^{\mathbf L}k=0$となる。
$R_N$はその$h$-filtrationによって$D(R_N)$内で$k$の有限extensionから作られる。
従って$R_N\in\operatorname{thick}(k)$であり、$M\otimes^{\mathbf L}k=0$は$M=0$を含意する。
よって$P$は$\mathcal D$のcompact generatorであり、compact部分を取って主張を得る。
ここではsmallなproper圏における右直交の消滅だけから生成性を結論していない。
\end{proof}

\subsection{有限windowによるenhanced比較}

\begin{theorem}[Artin線束残余pairのMF表示]\label{am:thm:comparison}
本節の入力に対し、有限window
\begin{equation}\label{am:eq:window-functor}
 \Psi=\operatorname{stab}\circ
       \bigl((\mathbf R\omega -)_{\ge\,4-d}\bigr)
\end{equation}
は$\mathcal K_B^{\rm lin}$上でdg Morita同値
\begin{equation}\label{am:eq:residual-mf}
 \Psi:\mathcal K_B^{\rm lin}\xrightarrow{\sim}\MF^{\ZZ}(A,\Omega)
\end{equation}
を与える。これはArtin商への底変更と両立する。
さらに$\rho=L_{\mathcal O}\circ(1)$に対し自然に
\begin{equation}\label{am:eq:rotation}
 \Psi\rho\simeq(1)\Psi
\end{equation}
である。特に$\rho$は自己同値となり、MFの三次周期性は
$\rho^3\simeq[2]$を与える。
\end{theorem}

\begin{proof}
まず\eqref{am:eq:window-functor}をdg quasi-functorとして定義する。
derived sectionと内部切断はenhanced関手であり、
stabilizationは補題\ref{am:lem:stabilization}で構成したdg Morita同値の逆である。
この関手が$\mathcal C_B$上で定義されることを生成子で確認する。
補題\ref{am:lem:line-bundles}から、$0\le j\le d-1$について
\[
 (\mathbf R\omega\mathcal O(j))_{\ge i}=B(j)_{\ge i}.
\]
これは$A$上perfectである。
実際、$B(j)$は$A$上長さ一の分解を持ち、有限長商
$T_j=B(j)_{<i}$の次数filtrationの各商は有限自由$R_N$のshiftである。
$R_N$は式\eqref{am:eq:relative-koszul}により$A$上perfectなので、$T_j$も同様である。
よって短完全列
\[
 0\to B(j)_{\ge i}\to B(j)\to T_j\to0
\]
の三項はいずれも$\mathcal D_{B/A}$に入る。
有限coneと直和因子へ延長して$\mathcal C_B$全体に$\Psi$を定義できる。

$0\le j<a$なら$-j\ge i$であり、$B(j)_{\ge i}=B(j)$は$B$上perfectである。
従って$\Psi$は例外成分$\mathcal O,\ldots,\mathcal O(a-1)$を零へ送る。
残る三つの$j=a,a+1,a+2$では
\begin{equation}\label{am:eq:three-images}
 \Psi(\mathcal O(j))\simeq\operatorname{stab}(T_j)[-1],\qquad
 \operatorname{gr}(T_j)=
 \bigoplus_{r=0}^{j-a}B_r\otimes_{R_N}R_N(j-r).
\end{equation}
従って三つの像は、三角的に
$\operatorname{stab}(R_N(a)),\operatorname{stab}(R_N(a+1)),
\operatorname{stab}(R_N(a+2))$を生成する。

剰余上では$\Psi$は通常のOrlovのsection--window関手である。
window $i=1-a$のkernelは$\langle\mathcal O,\ldots,\mathcal O(a-1)\rangle$であり、
その右直交への制限は古典Kuznetsov成分とgraded MFの同値である
\cite[Theorem 6.4 and its proof]{LP-BS}。
右直交と、同定されたwindowの反対側の直交の向きは混同しない。
両者は例外成分による同じVerdier商を表示し、$\Psi$の制限はその商関手である。

全ての自然なHom写像を剰余へ移すと、この古典同値のHom写像となる。
sourceのHomのproper性・底変更は補題\ref{am:lem:linear-category}、
targetは補題\ref{am:lem:mf-hom}による。
$\Psi$の底変更自然性はsectionの計算、内部切断、
補題\ref{am:lem:stabilization}の余核構成から従う。
従ってHom写像のconeは$R_N$-perfectで剰余が零であり、導来Nakayamaにより零である。
これで$\Psi|_{\mathcal K_B^{\rm lin}}$の完全忠実性を得る。
$G=p(E)$の像の剰余は古典MFのgeneratorである。
補題\ref{am:lem:generator}により$\Psi(G)$はArtin上でも全MFを生成するので、
本質的全射性が従う。

最後にrotationの自然性を確認する。
まず$\rho$は$\mathcal K_B^{\rm lin}$を保つ。
実際、$1\le j<a$では
$\operatorname{RHom}(\mathcal O(j),M(1))
=\operatorname{RHom}(\mathcal O(j-1),M)=0$であり、
式\eqref{am:eq:line-hom}より$\operatorname{RHom}(\mathcal O(j),\mathcal O)=0$である。
従ってevaluationの三角から$j\ge1$の直交性が従い、$j=0$の直交性はleft mutationの定義による。
$M\in\mathcal K_B^{\rm lin}$なら
\[
 (\mathbf R\omega M)_i
 \simeq\operatorname{RHom}(\mathcal O(-i),M)
 =\operatorname{RHom}(\mathcal O(a-1),M)=0.
\]
従って内部切断の自然な三角から
\[
 ((\mathbf R\omega M)(1))_{\ge i}
 =(\mathbf R\omega M)_{\ge i+1}(1)
 \simeq(\mathbf R\omega M)_{\ge i}(1)
\]
を得る。stabilizationは内部shiftを保つので、
$\Psi(M(1))\simeq\Psi(M)(1)$となる。
また$\Psi(\mathcal O)=0$であるから、evaluationの三角を$\Psi$で送ると
$\Psi(L_{\mathcal O}M(1))\simeq\Psi(M(1))$を得る。
これらの自然同型を合成して\eqref{am:eq:rotation}を得る。
ここでの周期同型は、指定された中心元のMFにおける周期変数$t$から移したものである。
\end{proof}


\section{正則局所基底からのArtin底変更}\label{ar:sec:regular}

$(R,\mathfrak m,\mathbb C)$を有限生成複素代数の局所化である正則局所環とする。
$A_R=\bigoplus_{j\ge0}(A_R)_j$は中央$R$-代数、$(A_R)_0=R$、
各次数有限自由、剰余が$d\ge5$変数の多項式環であるとする。
$\Omega_R\in Z(A_R)_3$の剰余は滑らかな三次式とし、
$B_R=A_R/(\Omega_R)$、$a=d-3$と置く。
$E_R=\bigoplus_{j=0}^{d-1}\mathcal O(j)$の例外線束成分を除く射影を$G_R$とし、
$\Gamma_R=\operatorname{REnd}(G_R)$と定める。
以下ではこの定義が可能になる有限性と、必要な底変更だけを確認する。

\subsection{局所模型の線束Hom}

\begin{lemma}[正則局所模型に必要な有限性]\label{ar:lem:regular-hom}
$A_R,B_R$は左右noetherianで、$(B_R)_j$は有限自由である。
線束のHomは$R$-perfectであり、
\[
 \operatorname{Ext}^q(\mathcal O(r),\mathcal O(s))=
 \begin{cases}(B_R)_{s-r}&q=0,\\
 (B_R)_{r-s-a}^{\vee}&q=d-2,\\
 0&\text{その他}
 \end{cases}
\]
となる。例外列$\mathcal O,\ldots,\mathcal O(a-1)$による射影が存在し、
$G_R$は線束生成残余部分を生成する。
固定次数のtorsion計算に現れる自由Hom複体は有限自由cohomologyに
homotopy同値であり、任意の通常のtensorで導来tensorを計算できる。
\end{lemma}
\begin{proof}
まずnoetherian性を次数別に確認する。
$A_R$の$\mathfrak m$進随伴次数環は、内部次数ごとの自由性より
$\operatorname{gr}_{\mathfrak m}R\otimes_{\mathbb C}(A_R/\mathfrak m A_R)$であり、
noetherianである。斉次右イデアル$I$の初期形式イデアルの有限生成元を持ち上げ、
それらが生成する右イデアルを$J$とする。
最低$\mathfrak m$次数を順に消すことで、各内部次数$j$と各$q$について
$I_j\subset J_j+\mathfrak m^q(A_R)_j$を得る。
$(A_R)_j/J_j$は有限$R$-加群なので、Krull交叉定理より$I_j=J_j$である。
一般の右イデアルには内部次数filtrationを用い、最高次数の有限降下を行う。
左側も同じであり、$B_R$は商なのでnoetherianである。

第\ref{am:sec:comparison}節の次数別分裂する持上げを、
$R$の$A_R$-加群としてのKoszul分解と、$A_R$の対角Koszul分解に適用する。
前者は項$A_R(-j)^{\binom dj}$、後者は項$(A_R^e)(-j)^{\binom dj}$を持つ。
対角分解を任意の加群へtensorし、$R$の有限大域次元を使えば
$A_R$も有限大域次元を持つ。
また$\Omega_R$による次数別の行列は剰余上単射なので、
単元小行列式により分裂単射であり、各$(B_R)_j$は自由である。

Koszul分解の双対に単元行列成分の消去を適用し、正則中心元のchange-of-ringsを用いると
\[
 \mathbf R\underline{\operatorname{Hom}}_{A_R}(R,A_R)=R(d)[-d],\qquad
 \mathbf R\underline{\operatorname{Hom}}_{B_R}(R,B_R)=R(a)[-(d-1)]
\]
を得る。左右のGorenstein型と$R$-flatnessにより、section--torsion三角に必要な
$\chi$と有限cohomological dimensionを使える
\cite[Section~4, Lemmas~6.1--6.2]{LP-BS}。
$M_\ell=B_R/(B_R)_{\ge\ell}$のfiltrationの商は$(B_R)_j\otimes_R R(-j)$である。
上の双対性から固定内部次数$t$では$j=-t-a$だけが寄与し、
$\ell>j$で安定する。従って
$H^{d-1}(\mathbf R\tau B_R)_t=(B_R)_{-t-a}^{\vee}$、他は零となる。
section--torsion三角から表示した線束Hom公式を得る。
$a>0$なので、固定した内部次数で二つのcohomology群が同時に非零になることはない。

底変更については、$M_\ell$の自由分解を各次数のsyzygyが$R$-自由となるように取る。
固定次数の自由Hom複体は下に有界で、cohomologyは有限自由加群の一つの次数に集中する。
$R$の有限大域次元を用い、完全な射影複体のcycleにdimension shiftingを行うと、
全cycleが射影となり、その複体は分裂する。
従ってcohomology複体へのhomotopy同値があり、通常のtensorが導来tensorを計算する。
項別Hom、filtered colimit、section三角を通じて自然な底変更を得る。

最後にKoszul型の線束完全複体の両端は階数一なので、
連続した$d$本から全twistを前後に生成できる。
Hom公式は例外列とその左右随伴を与え、有限evaluation coneによる射影$p_R$が存在する。
この射影を$E_R$へ適用した$G_R$が線束生成残余部分を生成する。
\end{proof}

\subsection{残余dg代数の底変更}

\begin{proposition}[既存の$\Gamma_R$模型のArtin底変更]\label{am:prop:Gamma-basechange}
この節の正則局所基底$R$上の入力から、
剰余体が$\CC$である局所準同型$R\to R_N$によって$A,B$を底変更したとする。
このとき自然に
\[
 \operatorname{Perf}(\Gamma_R\otimes_R^{\mathbf L}R_N)
 \simeq\mathcal K_{B\otimes_RR_N}^{\rm lin}
 \simeq\MF^{\ZZ}(A\otimes_RR_N,\Omega\otimes1)
\]
であり、mutationで定めたrotation候補はMFの内部shiftへ移る。
\end{proposition}

\begin{proof}
正則な変更先についての底変更だけでは十分でないため、
線束に必要な底変更をArtin基底上で確認する。
正則$R$上で$M_\ell=B/B_{\ge\ell}$の自由分解から作った、固定内部次数のHom複体は、
補題\ref{ar:lem:regular-hom}の証明により有限自由cohomologyへhomotopy同値である。
従って$R_N$への通常のtensorが導来tensorを計算する。
syzygyの次数別$R$-分裂により、底変更後の自由分解も完全である。
そのHom、filtered colimit、section三角を用いると、
変更先について補題\ref{am:lem:line-bundles}で計算したHomと自然に一致する。
よって線束間のHomとその積について必要な底変更を得る。

$R$は正則なので、$-\otimes_R^{\mathbf L}R_N$は有界複体に対して有限Tor振幅を持ち、
射影商への関手が定義される。
線束間のHomの底変更を有限cone・直和因子・evaluationへ延長すれば、
$E_R,G_R$とそのendomorphism dg代数の底変更を得る。
これが第一の同値である。第二の同値とrotationの記述は
定理\ref{am:thm:comparison}による。
この証明は、$D^b(\operatorname{qgr}B)$全体の任意の底変更についての主張を含まない。
\end{proof}


\section{形式MF存在から全有限切断の残余復元へ}
\begin{corollary}[全有限Artin切断における可換残余復元]
\label{am:cor:formal-residual}
$d\ge5$とし、定理\ref{fo:thm:formal-reconstruction}の入力を取る。
$B_N=A_N/(\Omega_N)$、
$Y_N=V(W_h\bmod h^{N+1})\subset\PP^{d-1}_{R_N}$と置く。
このとき同じ形式三次式$W_h$を用いて、全ての$N$で
\[
 (\mathcal K_{B_N}^{\rm lin},\rho)
 \simeq
 (\Ku_{\rm perf}(Y_N/R_N),\rho_{Y_N})
\]
を得る。これは指定された内部shift、三次周期性、Artin商への底変更と両立する。
ここで$\Ku_{\rm perf}$は$\operatorname{Perf}(Y_N)$内で
$\mathcal O,\ldots,\mathcal O(d-4)$を除いた右直交である。
既存の正則局所模型からの入力には、命題\ref{am:prop:Gamma-basechange}も適用できる。
\end{corollary}
\begin{proof}
定理\ref{am:thm:comparison}を非可換側と可換側に適用し、
その間に定理\ref{fo:thm:formal-reconstruction}のMF同値を挟む。

可換側では、$\mathcal O,\mathcal O(-1),\ldots,\mathcal O(-(d-1))$の直和が
$D_{\rm QCoh}(Y_N)$をcompactに生成する
\cite[Tag 0BQT, Remark 36.16.4]{AMStacks}。
具体的には、周囲の射影空間の相対対角分解がその線束生成子を与え、
閉埋込みのpushforwardがconservativeであることと随伴により、
そのpullbackが$Y_N$の生成子となる。
qcqs schemeでcompact対象はperfect対象である
\cite[Tag 09M1, Proposition 36.17.1]{AMStacks}から、
これらの線束の厚い閉包は$\operatorname{Perf}(Y_N)$である。
補題\ref{am:lem:linear-category}で用いた全twistの厚い閉包と一致する。
各関手のrotationと底変更の両立から、主張するpairの同値を得る。
\end{proof}

\section{最小混合試験族の三次中心元の閉公式}
\label{mt:sec:family}

一楕円ブロックの分離型から一般混合型へ進むため、
五変数のcubic
\begin{equation}\label{mt:eq:classical}
 F_\lambda=\frac{x^3+y^3+z^3}{3}+u^3+v^3+\lambda uxy
\end{equation}
を固定する。$u,v$を中心、$h,\lambda$を次数零のパラメータとし、
\begin{equation}\label{mt:eq:relations}
 [y,z]=h(x^2+\lambda uy),\qquad
 [z,x]=h(y^2+\lambda ux),\qquad
 [x,y]=hz^2
\end{equation}
で周囲の代数$A_{h,\lambda}$を定める。
本節ではこの族の三次中心元を多項式の閉公式で構成し、
自由代数内の三つの恒等式によって中心性を証明する。
これは可換復元先の二次係数$G_2$の計算ではなく、
それを計算するための非可換代数側の入力である。
中心式の特殊形は既存の混合中心式の有限体検算に現れるが、
ここでは特殊化を用いず、全パラメータに対する証明書を与える。
閉公式の先行研究上の初出は未確認である。

\begin{lemma}[古典特殊ファイバーの滑らかさ]
\label{mt:lem:classical-smooth}
$V(F_\lambda)\subset\PP^4$は$\lambda^3\ne-3$のとき、かつそのときに限り滑らかである。
\end{lemma}
\begin{proof}
全偏微分が零なら$z=v=0$であり、残る方程式は
\[
 x^2+\lambda uy=y^2+\lambda ux=3u^2+\lambda xy=0.
\]
$u=0$なら$x=y=0$となるので射影点は存在しない。
$u=1$とすると、$\lambda=0$では解がなく、$\lambda\ne0$なら$x,y$は非零である。
最初の二式を掛けると$xy=\lambda^2$を得て、第三式は$3+\lambda^3=0$となる。
逆にこの等式が成立すれば、$x^3=3$を満たす$x$と
$y=-x^2/\lambda$を選ぶことで解を得る。
三次のEuler等式により、これらの解は$F_\lambda=0$も満たす。
\end{proof}

\subsection{中心三次式と自由代数の証明書}

$a=\lambda u$と略記し、$T=xyz+xzy+yxz+yzx+zxy+zyx$と置く。

\begin{proposition}[混合試験族の中心元]
\label{mt:prop:center}
\begin{equation}\label{mt:eq:center}
 \boxed{\begin{aligned}
 \Omega_{h,\lambda}={}&
 \frac{x^3+y^3+z^3}{3}-\frac{h^2}{9}T
 +\frac a2(xy+yx)\\
 &+\frac{h^2a^2}{6}z+\frac{h^4a}{18}z^2+u^3+v^3.
 \end{aligned}}
\end{equation}
は$A_{h,\lambda}$の三次中心元である。
これは$\CC[h,\lambda]$上の恒等的な中心性であり、
$h\ne0$への局所化も、有限$h$-jetへの制限も必要としない。
$h=0$で$\Omega_{0,\lambda}=F_\lambda$である。
\end{proposition}
\begin{proof}
まず$a$を独立な中心係数とし、自由代数
$\mathbb Q[h,a]\langle x,y,z\rangle$で計算する。
関係式を
\[
 r_1=xy-yx-hz^2,\quad
 r_2=xz-zx+hy^2+hax,\quad
 r_3=yz-zy-hx^2-hay
\]
と書く。二つの演算を
\[
 \mathcal S_p(r)=rp^2+prp+p^2r,\qquad
 \mathcal D_{p,q}(r)=rpq-prq-qrp+qpr
\]
とし、$\{p,q\}_+=pq+qp$とする。
$\omega=\Omega_{h,\lambda}-u^3-v^3$について、以下の三式が
\emph{関係式で割る前の自由代数}において成立する：
\begin{align}
 [\omega,x]={}&
 -\frac13\mathcal S_y(r_1)-\frac13\mathcal S_z(r_2)
 +\frac{h^2}{18}
   \bigl(\mathcal D_{z,x}(r_1)+\mathcal D_{y,x}(r_2)\bigr)\notag\\
 &+\frac{ah^3}{18}[x,r_1]-\frac a2\{x,r_1\}_+
  +\frac{ah}{6}[z,r_2]+
  \frac h6[yz+zy,r_3],\label{mt:eq:certificate-x}\\
 [\omega,y]={}&
 \frac13\mathcal S_x(r_1)-\frac13\mathcal S_z(r_3)
 -\frac{h^2}{9}\mathcal D_{z,y}(r_1)
 +\frac{ah^3}{18}[y,r_1]+\frac a2\{y,r_1\}_+\notag\\
 &+\frac{ah^2}{18}[r_2,x]+\frac h9[r_2,2xz+zx]
 +\frac{h^2}{9}[r_2,y]y+\frac{ah}{6}[r_3,z],
 \label{mt:eq:certificate-y}\\
 [\omega,z]={}&
 \frac{a(2h^3+9)}{18}[r_1,z]
 +\frac h9[2xy+yx,r_1]+\frac{h^2}{9}[r_1,z]z\notag\\
 &+ayr_2+\frac13\mathcal S_x(r_2)
 -\frac{h^2}{9}\mathcal D_{y,z}(r_2)
 +ar_3x+\frac13\mathcal S_y(r_3).
 \label{mt:eq:certificate-z}
\end{align}
各式は語と$h,a$の冪ごとに展開して確認でき、全係数は有理数である。
右辺の各項は$r_i$を含むので、商代数で全交換子が零となる。
$a=\lambda u$を代入し、中心スカラー$u^3+v^3$を戻せば中心性を得る。
零パラメータでの表示は直接代入による。
\end{proof}

\subsection{Artin切断でのsorted PBW表示}

閉公式の中心性と、周囲の代数のflatnessを別々に確認する。
$u,v$の中心性を先に使い、active語を$z<y<x$の順に並べる。
還元規則は
\begin{equation}\label{mt:eq:rewrite}
 xy\longmapsto yx+hz^2,\qquad
 xz\longmapsto zx-hy^2-hax,\qquad
 yz\longmapsto zy+hx^2+hay
\end{equation}
である。これは一般の非零$h$で通常の語順序に関して停止する規則とは主張しない。

\begin{lemma}[全Artin切断でのPBW性]
\label{mt:lem:artin-pbw}
任意の$N\ge0$について、$h^{N+1}=0$を課した代数は
$\CC[\lambda,u,v,h]/(h^{N+1})$上で
$\{z^iy^jx^k:i,j,k\ge0\}$を基底に持つ。
従って各次数成分は$\CC[\lambda,h]/(h^{N+1})$上自由であり、
周囲のHilbert級数は$(1-t)^{-5}$である。
\end{lemma}
\begin{proof}
$h^r w$の還元に、辞書式に並べた二つの非負整数
$(N-r,\operatorname{inv}(w))$を用いる。
$h$を含まない交換項では反転数が減り、補正項では$h$の冪が増える。
従って$h^{N+1}=0$のもとで全還元が停止する。
active語の重なりは$xyz$の一つだけである。
最初に$xy$を還元した場合と、最初に$yz$を還元した場合は、いずれも有限回の置換で
\begin{equation}\label{mt:eq:overlap}
 zyx+h(x^3-y^3+z^3)
\end{equation}
になる。$a$を含む項は両経路でそれぞれ相殺する。
包含の曖昧さはない。係数に逆元を導入しないmonicな合流的還元なので、
既約語$z^iy^jx^k$が基底となる。
中心係数$a$を$\lambda u$へ置き換え、独立な中心変数$v$を加えても同じ証明が成り立つ。
\end{proof}

この補題は全ての$h$-Artin切断およびその逆極限についてのformal PBW性を示す。
非零$h$での全ファイバーのAS-regular性や代数的族全体のflatnessを、
Artin切断の計算だけから結論していない。
また、中心商は$v$に関してmonicな三次式であるため、
各Artin切断で$1,v,v^2$を、$v$を含まない周囲の代数上の基底に持つ。
従って中心三次商も次数別に平坦で、Hilbert級数は
$(1-t^3)/(1-t)^5$である。


\section{混合試験族のjetと同変制約}
\subsection{二次jetと基底表示の区別}

sorted PBW表示に直すと\eqref{mt:eq:center}は
\begin{align}\label{mt:eq:second-jet}
 \Omega_{h,\lambda}\equiv{}&
 \frac{x^3+y^3+z^3}{3}+a\,yx+u^3+v^3
 +\frac{ha}{2}z^2\notag\\
 &+h^2\left(-\frac23zyx+\frac{a^2}{6}z\right)
 \pmod{h^3}.
\end{align}
対称表示では偶数冪だけに見える式が、sorted表示では一次項を持つ。
これは積と基底表示の違いであり、可換復元先の一次変形が非零だという意味ではない。
$h^2a^2z/6$の補正は必要であり、これを省いた元の$x$との交換子は
\[
 -\frac{h^3a^2}{6}(y^2+ax)\pmod{h^4}
\]
となる。従って古典混合項を対称化して加えるだけでは中心にならない。

$\lambda=0$では既存の分離型に戻る。
一方、古典式の混合方向$uxy$は
\[
 J(F_0)_3=\bigl(\CC[x,y,z,u,v]/(x^2,y^2,z^2,u^2,v^2)\bigr)_3
\]
で非零であり、
active変数の中心シフトが作る$u\langle x^2,y^2,z^2\rangle$にも属さない。
従って中心部分空間を固定したactive変数の線形シフトだけで、
この混合方向を分離型へ戻すことはできない。

\subsection{対角対称性による二次復元候補の一方向への制限}

固定した$\lambda$について、有限対角群
\[
 H=\{(\alpha,\beta,\gamma,\delta,\varepsilon)\in\mu_3^5:
       \alpha\beta\gamma=1\},\qquad
 \chi(\alpha,\beta,\gamma,\delta,\varepsilon)=\alpha\beta\delta
\]
を考える。作用は
\[
 (x,y,u,z,v,h)\longmapsto
 (\alpha x,\beta y,\gamma u,\delta z,\varepsilon v,\chi h)
\]
で与える。$\lambda$は固定する。

\begin{lemma}[量子族の同変性とJacobian固有空間]
\label{mt:lem:character-space}
この作用は\eqref{mt:eq:relations}の関係式イデアルを保ち、
$\Omega_{h,\lambda}$および$F_\lambda$を固定する。
$J_\lambda=\CC[x,y,u,z,v]/(\partial F_\lambda)$とすると、
次数三の固有空間は
\begin{equation}\label{mt:eq:character-spaces}
 (J_\lambda)_3^\chi=\CC\cdot[xyz],\qquad
 (J_\lambda)_3^{\chi^2}=0.
\end{equation}
ここで固有空間の計算は、$\lambda^3=-3$を含む全ての$\lambda$について成り立つ。
\end{lemma}
\begin{proof}
例えば最初の関係式は全体が$\beta\delta$倍になる。
実際、$\chi\alpha^2=\beta\delta$、
$\chi\gamma\beta=\beta\delta$である。
他の二関係式も同様である。
中心式の各項の重みはそれぞれ一である。
特に$h^2T$の重みは$\chi^3=1$、
$h^2a^2z$の重みは$\chi^2\gamma^2\delta=1$、
$h^4az^2$の重みは$\chi^4\gamma\delta^2=1$となる。

固有空間は単項式の列挙で求まる。
$\gamma=(\alpha\beta)^{-1}$を消去し、独立な文字を
$(\alpha,\beta,\delta,\varepsilon)$とすると、
$(x,y,u,z,v)$の重みはそれぞれ
\[
 (1,0,0,0),\ (0,1,0,0),\ (2,2,0,0),\
 (0,0,1,0),\ (0,0,0,1)
 \quad\text{in }(\ZZ/3)^4
\]
である。多項式環の次数三の$\chi=(1,1,1,0)$固有空間は
$\langle xyz,u^2z\rangle$だけであり、
Jacobianイデアルの同じ固有空間は
\[
 \langle z\partial_uF_\lambda\rangle
 =\langle\lambda xyz+3u^2z\rangle
\]
だけである。後者の$u^2z$の係数は非零なので、
商は$[xyz]$で生成される一次元空間となる。
同様に$\chi^2=(2,2,2,0)$固有空間は
多項式環では$\langle uz^2\rangle$だけであり、
$u\partial_zF_\lambda=uz^2$がこれを消す。
有限群の標数零の表現は半単純なので、この計算は商の固有空間を正確に与える。
\end{proof}

\begin{corollary}[marked同変復元を仮定した二次係数の制約]
\label{mt:cor:second-character}
$\lambda^3\ne-3$とする。
非可換族に対するmarkedな可換復元が$\CC[h]/(h^3)$上に存在し、
上記の$H$作用が復元先の偏極付きcubic族に持ち上がって、
超平面束の特殊ファイバー上の所定の線形化も持ち上がり、
底の$h\mapsto\chi h$と両立すると仮定する。
その場合、復元先の一次変形類は零であり、
一次項を除去した後の二次変形類は
\begin{equation}\label{mt:eq:unknown-scalar}
 [G_2]=c(\lambda)[xyz]\in (J_\lambda)_3
\end{equation}
という一つの係数で表される。
本系は復元の存在や、その同変性を証明するものではない。
\end{corollary}
\begin{proof}
有限群について平均を取ることにより、
超平面束の切断加群の基底を、特殊ファイバーの上記の対角表現を保つように持ち上げられる。
方程式の生成元も、$F_\lambda$への留数を保って$H$不変に選べる。
従って方程式を
$F_\lambda+hG_1+h^2G_2$と書いたとき、
$G_i$の重みは$\chi^{-i}$である。
補題\ref{mt:lem:character-space}より$[G_1]=0$である。
実際、同変座標では$G_1$は$uz^2$のスカラー倍であり、
$z\mapsto z-hc u$という同変な座標変換で取り除ける。
この変換の同変性は$\chi\gamma=\delta$による。
その後も$G_2$の重みは$\chi^{-2}=\chi$なので、
同じ補題により\eqref{mt:eq:unknown-scalar}を得る。
\end{proof}

量子代数側の同変性は補題で直接確認したが、
可換復元への同変性の持上げは系の明示的な仮定である。
この区別を保ったうえで、残余圏の計算が求める二次変形類を
一つの係数$c(\lambda)$に絞れる。

第\ref{kj:sec:second-jet}節では星積・指定中心元・曲率を照合し、
$G_2=xyz/3-\lambda^2u^2z/12$、$c(\lambda)=1/3+\lambda^3/36$を実際に求める。
同節の積表はMF比較によって得られ、
第\ref{am:sec:comparison}節の相対windowにより線束生成残余pairにも移る。
本節の中心式の$h^2$係数をそのまま$G_2$と同一視するものではない。
一般の非零パラメータへの代数化と残余対角のperfect性は、別の未解決問題である。

証明書の再現は
\texttt{python3 computations/verify\_mixed\_test\_family.py}
で行う。同コードは
\eqref{mt:eq:certificate-x}--\eqref{mt:eq:certificate-z}を
自由代数内の係数ごとに厳密検算する。
別に$h^3,h^4,h^5,h^9$を法とするsorted還元、唯一の重なり、
二次補正を省いた負例も確認し、
\texttt{mixed\_test\_family\_certificate.json}を生成する。
自由代数恒等式による全パラメータの中心性と、有限jetの独立検算を区別して記録する。

\section{具体的な可換模型と中心式の可換化の相違}
\label{em:sec:comparison}

可換復元には、候補となる方程式を取り出すことと、その方程式の
Kuznetsov成分へのenhanced同値を作ることの二段階がある。
一楕円ブロックの分離型では、前者は二次関係式の多重線形化から計算でき、
後者には定理\ref{er:thm:main}のブロック比較を使える。
本節ではこの具体性を明示し、三次中心元の表示から積の順序を消す操作が、
その可換模型を与えないことを証明する。
一般混合型への圏同値は本節の結論に含めない。
Hessianによる点曲線公式はMori--Smithの既知の結果であり、
本節では量子化の規約に合わせた証明を示し、既存の圏比較への接続を明示する。
以下の二次比較と混合項の障害も初等的に証明できる。
それらを先行研究に対して新規なものとする確認は未完了である。

\subsection{三変数対称Jacobian量子化のHessian公式}

三変数の任意の斉次三次式$f\in\CC[x,y,z]_3$に対し、
$\operatorname{Hess}(f)=\det(\partial_i\partial_j f)$と書く。
これはHessian\emph{行列}ではなく、その行列式である。
以下の多項式恒等式は$f$の滑らかさを仮定しない。
点曲線の記述はMori--Smith
\cite[Lemma 7.1, Theorem 7.2, pp.~23--25]{PMoriSmith}にある。
同論文の$A_{f,\lambda}$の規約では、行列式は
$\lambda^3\operatorname{Hess}(f)/8+3\lambda f$である。
ここでは関係式の符号を合わせて、その式を$-3h$で正規化する。

\begin{proposition}[Mori--SmithのHessian公式の量子化表示]
\label{em:prop:hessian-formula}
対称Jacobian関係式
\[
 [y,z]=h\sym(\partial_xf),\qquad
 [z,x]=h\sym(\partial_yf),\qquad
 [x,y]=h\sym(\partial_zf)
\]
を多重線形化すると、$v=(x,y,z)^{\mathsf T}$について係数行列は
\[
 M_f(v)=K(v)-\frac h2 H_f(v),\qquad
 K(v)=\begin{pmatrix}0&-z&y\\z&0&-x\\-y&x&0\end{pmatrix},
 \quad H_f=(\partial_i\partial_j f)
\]
となる。その行列式は
\begin{equation}\label{em:eq:hessian-identity}
 \det M_f=-3h\left(f+\frac{h^2}{24}\operatorname{Hess}(f)\right).
\end{equation}
従って$h\ne0$では、多重線形化の退化三次式は
\begin{equation}\label{em:eq:hessian-candidate}
 f_h^{\mathrm{pt}}=f+\frac{h^2}{24}\operatorname{Hess}(f)
\end{equation}
で与えられる。この式は、退化三次式が恒等的に零になる場合も含む恒等式である。
実際の点曲線との同定には、対応するAS-regular代数と点加群についての追加仮定が必要である。
\end{proposition}
\begin{proof}
二次式$q$の対称化の多重線形化は
$\frac12\sum_j(\partial_jq)(v)v'_j$である。
$q=\partial_i f$と置けば行列表示を得る。
任意の対称$3\times3$行列$H$と$\varepsilon$について、
$K^{\mathsf T}=-K$より
\[
 \det(K-\varepsilon H)=-\det(K+\varepsilon H).
\]
左辺は$\varepsilon$の三次以下の奇多項式であり、
$\operatorname{adj}K=vv^{\mathsf T}$なので一次係数は
$-\operatorname{tr}(\operatorname{adj}K\,H)=-v^{\mathsf T}Hv$である。
三次係数は$-\det H$であり、従って
\[
 \det(K-\varepsilon H)=-\varepsilon v^{\mathsf T}Hv-
 \varepsilon^3\det H.
\]
三次のEuler等式を二回用いると$v^{\mathsf T}H_fv=6f$である。
$\varepsilon=h/2$を代入すれば\eqref{em:eq:hessian-identity}を得る。
\end{proof}

この公式は座標選択とも整合する。実際、
$g(x)=\lambda f(Lx)$、$L\in\mathrm{GL}_3(\CC)$、$\lambda\ne0$なら
\begin{equation}\label{em:eq:hessian-covariance}
 \operatorname{Hess}(g)(x)
 =\lambda^3\det(L)^2\operatorname{Hess}(f)(Lx),\qquad
 g_h^{\mathrm{pt}}(x)=\lambda f_{h'}^{\mathrm{pt}}(Lx),
 \quad h'=\lambda\det(L)h.
\end{equation}
第一式は$H_g=\lambda L^{\mathsf T}H_f(Lx)L$の行列式を取れば従い、
第二式は代入で得られる。$h'$はJacobian Poisson括弧の座標変換則と一致する。
固定した体積形式の下での式であり、Hessianの数値係数を座標変換と独立に扱わない。

\begin{corollary}[一般の滑らかな一ブロックの候補式]
\label{em:cor:hessian-reconstruction}
$g$は滑らかな平面三次曲線を定め、$\Omega_h^g$は
系\ref{er:cor:ternary}の正規化された中心元とする。
同系のパラメータ開集合と$P(u)$の仮定の下では、
定理\ref{er:thm:main}の圏比較を入力とすることにより、
\begin{equation}\label{em:eq:general-hessian-output}
 V\left(g(x)+\frac{h^2}{24}\operatorname{Hess}(g)(x)+P(u)\right)
\end{equation}
が$S_h^g[u]/(\Omega_h^g+P)$のrotation付き可換模型となる。
これは同系の模型の射影同型な表示であり、一般混合型の復元定理ではない。
\end{corollary}
\begin{proof}
$g=\lambda f_t(Lx)$とHesse形に移し、$h'=\lambda\det(L)h$とする。
直接計算で
\[
 \operatorname{Hess}(f_t)=
 2(t^3+4)xyz-2t^2(x^3+y^3+z^3),\qquad
 f_t+\frac{h'^2}{24}\operatorname{Hess}(f_t)
 =b(h',t) f_{\vartheta(h',t)}
\]
を得る。従って\eqref{em:eq:hessian-covariance}により
\eqref{em:eq:general-hessian-output}の方程式は
$\lambda b(h',t) f_{\vartheta(h',t)}(Lx)+P(u)$となる。
$b(h',t)\ne0$なので、三変数を一斉に$b(h',t)^{-1/3}$倍すれば、
系\ref{er:cor:ternary}の方程式
$\lambda f_{\vartheta(h',t)}(Lx)+P(u)$になる。
この座標変換と同系の圏同値により結論を得る。
ここで用いた三乗根は複素ファイバーごとに選べる。
族上の同型を選ぶ際には、必要な有限\'etale被覆を別途指定する。
\end{proof}

\subsection{二次関係式から取り出せる分離型の方程式}

$S_{h,t}$、$\Omega_{h,t}$、$P(u)$および$\mathscr U_t$は
第\ref{er:sec:intro}節のものとする。
$t^3\ne-1$、$h\in\mathscr U_t$、$m\ge2$とし、$P$のアフィン臨界点は原点のみとする。
三次中心商は
\[
 B_{h,t,P}=S_{h,t}[u_1,\ldots,u_m]/(\Omega_{h,t}+P(u))
\]
である。既存のSklyanin中心商と楕円曲線の比較
\cite[pp.~20--21]{SmithTate}\cite[Sections 3.5--3.6]{EG}だけでなく、
内部シフトとの整合性および中心変数を加えた後の生成性を使う点が重要である。

\begin{proposition}[分離型の具体的復元手順]
\label{em:prop:determinant}
$\alpha=1-ht/2$、$\beta=-1-ht/2$、$\gamma=-h$と置く。
関係式を第二の三変数について線形にした行列を
\[
 M(x,y,z)=
 \begin{pmatrix}
 \gamma x&\beta z&\alpha y\\
 \alpha z&\gamma y&\beta x\\
 \beta y&\alpha x&\gamma z
 \end{pmatrix}
\]
とすると
\begin{equation}\label{em:eq:determinant}
 -\det M=\alpha\beta\gamma(x^3+y^3+z^3)
 -(\alpha^3+\beta^3+\gamma^3)xyz.
\end{equation}
$h\ne0$では、$b=1-h^2t^2/4$を用いて
\begin{equation}\label{em:eq:explicit-output}
 F_{\mathrm{rec}}(h)
 =-\frac{\det M}{3hb}+P(u)
 =\frac{x^3+y^3+z^3}{3}+\vartheta(h,t)xyz+P(u)
\end{equation}
を得る。右辺は$h=0$まで正則に延長する。
定理\ref{er:thm:main}の圏比較を入力とすれば、
$Y_h=V(F_{\mathrm{rec}}(h))$は$B_{h,t,P}$のrotation付き可換模型である。
\end{proposition}
\begin{proof}
多重線形化された関係式は
$M(x,y,z)(x',y',z')^{\mathsf T}=0$である。
行列式を六項に展開すると\eqref{em:eq:determinant}を得る。
さらに
\[
 \alpha\beta\gamma=hb,\qquad
 \alpha^3+\beta^3+\gamma^3
 =-3h\left(t+\frac{h^2}{3}(1+t^3/4)\right)
\]
を代入すれば\eqref{em:eq:explicit-output}となる。
零パラメータでは行列式全体が零になるため、共通因子$h$を除いてから延長する。
これは$S_{0,t}$のpoint schemeを楕円曲線と同一視する操作ではない。
最後の圏同値は定理\ref{er:thm:main}そのものであり、
行列式の計算だけから従うものではない。
\end{proof}

一般の三次元Sklyanin関係式
$\alpha yz+\beta zy+\gamma x^2=0$とその巡回置換で与えられた場合も、
滑らかな点曲線および標準中心三次商であることを確認したうえで、
\eqref{em:eq:determinant}が点曲線の方程式を直接与える。
分離型では方程式の非零スカラー倍はブロック変数の一斉拡大で吸収できる。
その拡大に必要な三乗根は$\CC$上で存在するが、一般の基礎体上では底拡大の指定が必要である。
これは元の古典極限の方程式を再利用する手順ではなく、
与えられた量子関係式から変形した楕円曲線を取り出す手順である。

\subsection{中心式の積を可換にすると二次で違う模型になる}

$S_3=x^3+y^3+z^3$と置き、中心元の指定された対称表示
$\Omega_{h,t}=U(h,t)S_3/3+W(h,t)\sym(xyz)$の積の順序を消した多項式を
\[
 F_{\mathrm{sym}}(h)=U(h,t)S_3/3+W(h,t)xyz+P(u)
\]
と定義する。これは選んだ表示に依存する操作であって、
代数の可換化とは異なる。可換化との相違は次の小節で扱う。

\begin{proposition}[二次Jacobian類の不一致]
\label{em:prop:second-order-failure}
$F_0=f_t+P$、$J(F_0)=\CC[x,y,z,u]/(\partial F_0)$とする。
$F_{\mathrm{rec}}(h)$と$F_{\mathrm{sym}}(h)$はいずれも$h$の一次項を持たず、
二次係数のJacobian類はそれぞれ
\begin{equation}\label{em:eq:two-classes}
 \frac{1+t^3}{3}[xyz],\qquad
 -\frac{2(1+t^3)}{3}[xyz]
 \quad\text{in }J(F_0)_3
\end{equation}
である。差$(1+t^3)[xyz]$は非零である。
従って二つの偏極付きcubic族は、$\CC[h]/(h^3)$上で
特殊ファイバーの恒等写像を保つ射影座標変換によって同型にならない。
特に、中心式の積の順序を消す操作は、分離型でもmarkedな可換復元を与えない。
\end{proposition}
\begin{proof}
式\eqref{er:eq:coefficients}を展開すると
\begin{align*}
 \vartheta(h,t)&=t+\frac{1+t^3}{3}h^2+O(h^4),\\
 U(h,t)&=1+\frac{t^2}{2}h^2+O(h^4),\\
 W(h,t)&=t-\frac{t^3+4}{6}h^2+O(h^4).
\end{align*}
$J(f_t)$では
$x^2+tyz=y^2+tzx=z^2+txy=0$なので
$S_3=-3txyz$である。従って$F_{\mathrm{sym}}$の二次係数の類は
\[
 \left[\frac{t^2}{6}S_3-\frac{t^3+4}{6}xyz\right]
 =-\frac{2(1+t^3)}{3}[xyz]
\]
となる。

$[xyz]\ne0$も有限な線形代数で直接確認できる。
次数三の単項式を
\[
 (x^3,x^2y,x^2z,xy^2,xyz,xz^2,y^3,y^2z,yz^2,z^3)
\]
の順に並べる。列を
$(x\partial_xf_t,y\partial_xf_t,z\partial_xf_t,
x\partial_yf_t,y\partial_yf_t,z\partial_yf_t,
x\partial_zf_t,y\partial_zf_t,z\partial_zf_t,xyz)$
とする係数行列の行列式は
\[
 -(1+t^3)^2.
\]
従って$t^3\ne-1$のとき、最初の九列は独立であり、$xyz$はその張る空間に含まれない。
変数の分離から$J(F_0)=J(f_t)\otimes J(P)$であり、
$J(P)_0=\CC$なので$[xyz]\otimes1$も非零である。

最後に、$F_0$は滑らかなcubicなので、その偏微分は二次の正則列をなす。
これらに一次係数のsyzygyはなく、Euler等式を除いて無限小射影自己同型はない。
従って一次項が零である二つのmarked族の間の一次座標変換は射影的に零である。
二次座標変換と方程式の単元倍で変えられる二次項は$(\partial F_0)_3$に限られる。
差のJacobian類が非零なので、求める同型は存在しない。
\end{proof}

\begin{example}[Fermatブロック]
\label{em:ex:fermat}
$t=0$、$P(u,v)=u^3+v^3$では
\begin{align*}
 F_{\mathrm{rec}}(h)&=\frac{x^3+y^3+z^3}{3}+u^3+v^3
                       +\frac{h^2}{3}xyz,\\
 F_{\mathrm{sym}}(h)&=\frac{x^3+y^3+z^3}{3}+u^3+v^3
                       -\frac{2h^2}{3}xyz.
\end{align*}
両式の差は正確に$h^2xyz$である。
ここで$J(F_0)=\CC[x,y,z,u,v]/(x^2,y^2,z^2,u^2,v^2)$なので、
差の非零性は平方自由単項式$xyz$だけで確認できる。
この判定はmarked変形の判定であり、
有限自己同型群で割った粗モジュライの局所座標について、同じ次数の主張を加えていない。
\end{example}

\subsection{代数そのものの可換化は次元さえ保たない}

\begin{proposition}[文字通りの可換化の失敗]
\label{em:prop:abelianization}
$t=0$、$h\ne0$とする。$B_{h,0,P}$を、全ての交換子が生成する両側イデアルで割った
可換代数を$B_{h,0,P}^{\mathrm{ab}}$と書く。このとき
\[
 B_{h,0,P}^{\mathrm{ab}}
 \simeq
 \CC[x,y,z,u_1,\ldots,u_m]
 /(x^2,y^2,z^2,P(u)-2h^2xyz/3).
\]
その被約環は$\CC[u]/(P)$であり、$\operatorname{Proj}$の次元は$m-2$である。
従って、次元$m+1$の復元先cubicをこの可換化から得ることはできない。
\end{proposition}
\begin{proof}
関係式$[y,z]=hx^2$とその巡回置換に可換性を課すと、$h\ne0$から
$x^2=y^2=z^2=0$が従う。
中心方程式の三乗項は消え、表示した最後の関係式が残る。
$x,y,z$は冪零なので、被約化後に残る関係は$P=0$である。
$P$は原点以外に臨界点を持たない非零三次式であり、特に被約である。
よって被約環と次元の主張を得る。
\end{proof}

\subsection{非楕円混合型のモデルへの適用障害}

節点型と尖点型で用いる三変数三次式は
\[
 f_{\mathrm N}=y^3+z^3+xyz,\qquad
 f_{\mathrm C}=xz^2+y^3
\]
であり、全体の可換特殊ファイバーには$ux^2$という混合項がある。
この項は、一ブロック定理で許される中心変数方向への線形平行移動に吸収できない。

\begin{lemma}[混合項$ux^2$は中心シフトで除去できない]
\label{em:lem:essential-mixing}
$\epsilon=\mathrm N,\mathrm C$のいずれについても
$x^2\notin\langle\partial_xf_\epsilon,\partial_yf_\epsilon,
\partial_zf_\epsilon\rangle\subset\CC[x,y,z]_2$である。
従って$f_\epsilon+ux^2$から$u$の一次混合項を
$(x,y,z)\mapsto(x+a u,y+b u,z+c u)$によって除去できない。
\end{lemma}
\begin{proof}
点$(x,y,z)=(1,0,0)$では両方の$f_\epsilon$の三つの偏微分が全て零であり、
$x^2=1$である。よって表示した線形空間に$x^2$は属さない。
中心シフトで$f_\epsilon$の$u$の一次係数に加わるものは
$a\partial_x f_\epsilon+b\partial_y f_\epsilon+c\partial_z f_\epsilon$だけなので、結論を得る。
\end{proof}

従って非楕円二例には、単に三変数のpoint curveを取り出して
定理\ref{er:thm:main}を適用することはできない。
候補式を得るためには、全七変数の量子関係式と中心三次元から残余圏の
rotation付き乗法データを計算するか、混合項を保つ具体的な行列因子分解同値を作る必要がある。
候補cubicが滑らかであることや、有限個のHochschild次元が一致することだけでは同値は証明されない。
また中心変数を通常の偏微分で扱う生成性の議論は、非楕円混合型では
その偏微分が代数導分であることから確かめ直す必要がある。

\subsection{圏論的入力と有限計算の役割}

系\ref{em:cor:hessian-reconstruction}と命題\ref{em:prop:determinant}の圏同値部分は
定理\ref{er:thm:main}と、その証明で用いるdg window、
periodicityを保つブロック比較、中心変数に対する外部積生成、
有限群の双対同変化\cite[Theorem 4.2, Section 8]{Elagin}を入力としている。
行列式や中心係数の一致は、これらの圏論的入力の独立監査を代替しない。
命題\ref{em:prop:second-order-failure}、命題\ref{em:prop:abelianization}、
補題\ref{em:lem:essential-mixing}の証明は上記の初等的計算で完結する。
付属の\texttt{verify\_explicit\_model\_comparison.py}は、
一般の十係数三次式に対するHessian恒等式、Hesse形への特殊化、
多重線形化行列の行列式、二次係数、十次元係数行列の行列式、
Fermat例のJacobian剰余、混合項の評価を有理数上で厳密に検算する。
この検算は圏同値、一般混合型の可換復元、残余smoothnessの計算機証明ではない。

% Explicit second-order MF model and rotation products in the mixed family.
\section{混合試験族の明示的な二次MF復元とrotation積}
\label{kj:sec:second-jet}

本節では第\ref{mt:sec:family}節の非可換中心対について、二次の可換MF模型を計算する。
$h^3=0$で周囲の代数とKontsevich星積の同型および指定中心元の像を求め、
定理\ref{cg:thm:mf-comparison}によって内部シフト付きMF同値を得る。
MF pairの次数一・二のrotation積も決定する。
第\ref{am:sec:comparison}節の相対windowにより、
滑らかな古典fiberを持つ線束生成射影残余pairにも適用できる。

二次星積にはDitoの明示式\cite[式~(3)、arXiv版の印刷頁5]{KJDito}を使い、
曲率の量子化にはShoikhetのwheel消滅\cite[Section~1.3]{KJShoikhet}を使う。
これらは異なるgraphの重みである。
混合族への適用と明示的積表の先行研究上の初出は未確認であり、
formalityやgraph積分自体を本稿の新しい結果とは扱わない。

\subsection{星積の規格化とWeyl対称化}

$S=\CC[x,y,z,u,v]$、$R=\CC[\lambda,h]/(h^3)$とする。
$F=F_\lambda$を式\eqref{mt:eq:classical}の三次式として
\begin{equation}\label{kj:eq:pi}
 \pi=F_x\partial_y\wedge\partial_z+
 F_y\partial_z\wedge\partial_x+F_z\partial_x\wedge\partial_y
 =\frac12\pi^{ij}\partial_i\wedge\partial_j
\end{equation}
と置く。添字は五変数を走り、$u,v$を含む$\pi^{ij}$は零である。
$\{x_i,x_j\}=\pi^{ij}$という規格化で、$F,u,v$はCasimirである。
角度型Kontsevich星積は
\begin{align}
 f*g={}&fg+\frac h2\pi^{ij}f_i g_j\notag\\
 &+h^2\left\{\frac18\pi^{ij}\pi^{kl}f_{ik}g_{jl}
 +\frac1{12}\pi^{ij}\partial_j\pi^{kl}(f_{ik}g_l-f_k g_{il})
 -\frac1{24}Q_{ik}f_i g_k\right\},\label{kj:eq:star}\\
 Q_{ik}:={}&\sum_{j,l}(\partial_l\pi^{ij})(\partial_j\pi^{kl}).
 \label{kj:eq:eye-tensor}
\end{align}
同じ添字を足し合わせる。
Ditoの一次項は$\hbar\pi_{\rm Dito}^{ij}f_i g_j$なので、
$\hbar=h$、$\pi_{\rm Dito}=\pi/2$とした。
その二次式の最後の係数$-1/6$は本稿では$-1/24$になる。
この最後の項をeye項と呼ぶ。二次作用素全体は$f,g$について対称である。

星積によるWeyl対称化を
\[
 T(x_{i_1}\cdots x_{i_d})
 =\frac1{d!}\sum_{\sigma\in\mathfrak S_d}
 x_{i_{\sigma(1)}}*\cdots*x_{i_{\sigma(d)}}
\]
で定める。
\begin{lemma}[対称化写像の二次項]\label{kj:lem:symmetrization}
任意の多項式$f$について
\begin{equation}\label{kj:eq:T}
 T(f)=f-\frac{h^2}{48}\sum_{i,k}Q_{ik}\partial_i\partial_kf
\end{equation}
が成り立つ。これはJacobianイデアルを法としない多項式の等式であり、
任意の変数の数と一般のPoisson tensorについて有効である。
\end{lemma}
\begin{proof}
定数係数の一次式$\ell$について$B_1(\ell^m,\ell)=0$である。
二次項では$\ell$の二階微分を含む項は消え、
残るtree項も$\partial_j\pi^{kl}\ell_k\ell_l=0$で消える。
従って
\[
 B_2(\ell^m,\ell)=-\frac m{24}Q_{ik}\ell_i\ell_k\ell^{m-1}.
\]
既に$h^2$を含む補正に一次星積を作用させても$h^3$以上にしか寄与しない。
帰納法により
\[
 T(\ell^m)=\ell^m-\frac{h^2m(m-1)}{48}
 Q_{ik}\ell_i\ell_k\ell^{m-2}.
\]
各斉次次数の多項式は一次式の冪で張られるから、偏極により主張を得る。
\end{proof}

active Hessianを
\begin{equation}\label{kj:eq:hessian}
 \mathsf H=(F_{ij})_{i,j\in\{x,y,z\}},\qquad
 \mathcal H=\det\mathsf H=8xyz-2\lambda^2u^2z
\end{equation}
とする。
\begin{lemma}[Jacobian tensorとHessian]\label{kj:lem:hessian}
$Q$のactive部分は$-2\operatorname{adj}(\mathsf H)$、他の成分は零であり、
\begin{equation}\label{kj:eq:T-F}
 T(F)=F+\frac{h^2}{8}\mathcal H
\end{equation}
となる。tensor恒等式は、passive変数を係数とする任意の三変数多項式にも成り立つ。
\end{lemma}
\begin{proof}
交代tensorを$\epsilon_{ijk}$とすると$\pi^{ij}=\epsilon_{ijr}F_r$である。
従って余因子の表示から
\[
 Q_{ik}=\epsilon_{ijr}\epsilon_{kls}\mathsf H_{rl}\mathsf H_{sj}
       =-2\operatorname{adj}(\mathsf H)_{ik}.
\]
可逆性を仮定せず成立する恒等式である。$F_{ik}$と縮約すると
$\sum Q_{ik}F_{ik}=-6\det\mathsf H$となり、式\eqref{kj:eq:T}を使えばよい。
$u,v$を含む成分は$\pi$の脚がactive方向にしかないため零である。
\end{proof}

\subsection{指定された非可換中心元の像}

\begin{proposition}[周囲の代数と中心元の二次同定]
\label{kj:prop:algebra-and-center}
$A_R=A_{h,\lambda}/(h^3)$を式\eqref{mt:eq:relations}の代数とする。
生成元を同名の一次式へ送る写像は次数付き$R$-代数同型
\begin{equation}\label{kj:eq:Phi}
 \Phi:A_R\xrightarrow{\sim}(S\otimes R,*)
\end{equation}
を定め、式\eqref{mt:eq:center}の中心元について
\begin{equation}\label{kj:eq:effective}
 \begin{split}
 \Phi(\Omega_{h,\lambda})&=\widehat F_\lambda:=F_\lambda+h^2G_2,\\
 G_2&=\frac{\mathcal H}{24}
     =\frac{xyz}{3}-\frac{\lambda^2u^2z}{12}
 \end{split}
\end{equation}
となる。この命題は$\lambda^3\ne-3$を仮定しない。
\end{proposition}
\begin{proof}
二次作用素の対称性から$[x_i,x_j]_* =h\pi^{ij}$である。
関係式の右辺の二次式を星積で評価した際の補正には$h^2$が掛かり、
関係式の係数$h$と合わせると消える。例えば
\[
 [y,z]_* =h(x^2+\lambda uy)=h(x*x+\lambda u*y).
\]
他の関係も同様で、$u,v$との星積は通常の積である。
従って$\Phi$はwell-definedである。
補題\ref{mt:lem:artin-pbw}により各次数の定義域と値域は同階数の自由$R$-加群である。
$h=0$で恒等なので、冪零イデアル$(h)$について逆行列を有限展開でき、$\Phi$は同型となる。

通常の多項式から$A_R$へのWeyl対称化を$\sym_A$と書く。
中心元の閉公式は
\begin{equation}\label{kj:eq:omega-weyl}
 \Omega_{h,\lambda}=\sym_A(F_\lambda-h^2\mathcal H/12)
\end{equation}
と書ける。$\sym_A(xyz)$は六語の和の$1/6$倍で、
$\sym_A(\lambda uxy)=\lambda u(xy+yx)/2$である。
$\Phi\circ\sym_A=T$なので
\[
 \Phi(\Omega)=T(F)-\frac{h^2}{12}T(\mathcal H)
 =F+h^2\left(\frac18-\frac1{12}\right)\mathcal H
 =F+\frac{h^2}{24}\mathcal H .
\]
\end{proof}

中心元のWeyl記号$F-h^2\mathcal H/12$と、
Kontsevich側の曲率$F+h^2\mathcal H/24$は異なる。
差は対称化の$+h^2\mathcal H/8$であり、中心元の語を単に可換に置き換えても得られない。

\subsection{曲率のwheelと有限階数MF比較}

Kontsevichのaffine formality morphismを$\mathcal U$とし、
一次成分を通常の反対称化HKRに規格化する\cite[Sections~4--6]{PKontsevich}。
関数$f$に対するtangent mapは
\begin{equation}\label{kj:eq:tangent-map}
 \mathcal U^{h\pi}(f)=f+h\mathcal U_2(\pi,f)
 +\frac{h^2}{2}\mathcal U_3(\pi,\pi,f).
\end{equation}
\begin{lemma}[曲率側の二輪は消える]\label{kj:lem:curvature-wheel}
角度型formalityについて、任意の関数$f$に対し
$\mathcal U^{h\pi}(f)=f$が$h^3=0$で成り立つ。
\end{lemma}
\begin{proof}
出力は零項演算でありboundary vertexはない。
$\mathcal U_2(\pi,f)$では$\pi$からの二本の辺の終点は$f$しかなく、
自己辺と重複辺が許されないためgraphが存在しない。
$\mathcal U_3(\pi,\pi,f)$では、各$\pi$頂点から他方の$\pi$頂点と$f$へ
一本ずつ辺を出す二輪だけが可能である。
これは\cite[Section~8.3.3.1]{PKontsevich}のtangent-wheelで、
重み$W_2$は\cite[Section~1.3]{KJShoikhet}により零である。
重みはgraphだけで決まる定数なので$\pi$が線形である必要はない。
\end{proof}
星積のeye項には二つのboundary vertexがあり、その重みは零ではない。
曲率のgraphでは$f$がinternal vertexに置かれ、boundary vertexは零個である。
同じ$\partial\pi\,\partial\pi$というtensorが現れることだけから重みを同一視してはいけない。

\begin{theorem}[混合試験族の二次MF復元]
\label{kj:thm:curved-MC}\label{kj:thm:mf-reconstruction}
$\lambda$を任意の複素数に固定し、$R_\lambda=\CC[h]/(h^3)$とする。
第\ref{cg:sec:comparison}節の有限階数・冪等完備なMF圏について、
$h=0$で恒等となる$R_\lambda$-linearなdg Morita同値
\begin{equation}\label{kj:eq:mf-reconstruction}
 \bigl(\MF^{\ZZ}(A_R,\Omega_{h,\lambda}),(1)\bigr)
 \simeq
 \bigl(\MF^{\ZZ}(S\otimes R_\lambda,\widehat F_\lambda),(1)\bigr)
\end{equation}
が存在し、指定された三次周期性$(3)\simeq[2]$も保つ。
\end{theorem}
\begin{proof}
命題\ref{kj:prop:algebra-and-center}が生成元を固定する$\Phi$と
$\Phi(\Omega)=\widehat F_\lambda$を与える。
曲率wheelの消滅により$C_h(F)=F$、
$C_h(F+h^2G_2)=F+h^2G_2$なので、系\ref{cg:cor:matching}の条件を満たす。
同系の有限curved gaugeとtwisted complexへの延長を適用すればよい。
bar同型とその逆のTaylor成分はarity一・三・五に限られ、有限rankと内部shiftを保存する。
従って中心性やHochschild類の一致だけからの推論ではない。
\end{proof}

\subsection{低次数rotation積の完全な二次表}

$\mathcal M_R=\MF^{\ZZ}(A_R,\Omega_{h,\lambda})$、$\rho_R=(1)$と置く。
\begin{corollary}[MF pairの低次数積]\label{kj:cor:conditional-Nat}
$\lambda^3\ne-3$とする。
$R_\lambda$-linearな次数一・二の導来自然変換と積は
\[
 V=R_\lambda\langle x,y,z,u,v\rangle,\qquad
 W=\bigl(R_\lambda[x,y,z,u,v]/(\partial\widehat F_\lambda)\bigr)_2
\]
の積で与えられる。核の基底は
\begin{align}
 q_x&=x^2+\lambda uy+\frac{h^2}{3}yz,&
 q_y&=y^2+\lambda ux+\frac{h^2}{3}xz,\notag\\
 q_z&=z^2+h^2\left(\frac{xy}{3}-\frac{\lambda^2u^2}{12}\right),&
 q_u&=3u^2+\lambda xy-\frac{h^2\lambda^2}{6}uz,\notag\\
 q_v&=3v^2.&&\label{kj:eq:gradient-kernel}
\end{align}
$W$の順序付き基底を
\begin{equation}\label{kj:eq:W-basis}
 (xy,xz,xu,xv,yz,yu,yv,zu,zv,uv)
\end{equation}
とすれば、異なる二文字の積は対応する基底vectorで、平方の積は
\begin{align}
 x\cdot x&=-\lambda uy-\frac{h^2}{3}yz,&
 y\cdot y&=-\lambda ux-\frac{h^2}{3}xz,\notag\\
 z\cdot z&=-h^2\left(\frac13+\frac{\lambda^3}{36}\right)xy,&
 u\cdot u&=-\frac{\lambda}{3}xy+\frac{h^2\lambda^2}{18}uz,\notag\\
 v\cdot v&=0.&&\label{kj:eq:multiplication-table}
\end{align}
これは$10\times15$の対称積行列を完全に指定する。
\end{corollary}
\begin{proof}
定理\ref{kj:thm:mf-reconstruction}により可換側で計算できる。
可換超曲面は滑らかな特殊fiberを持つ冪零変形で、相対的に滑らかである。
命題\ref{prop:ej-geometric-low-degree}の対角Koszul計算を相対Artin基底で行う。
偏微分列は特殊fiberで正則列であり、有限な$h$進filtrationを用いると
正のKoszul次数の完全性が持ち上がる。
五変数では非恒等sectorはrotation次数一・二のcohomological degree零には寄与しない。
従って$V,W$と積は上の相対Jacobianデータである。

核は$\widehat F$の直接微分である。
平方$\{x^2,y^2,z^2,u^2,v^2\}$の係数小行列の行列式は$9$なので、
これら五関係は平方を一意に消去し、残る十単項式が$W$の自由基底となる。
$z^2$の式に$u^2=-\lambda xy/3\pmod{h^2}$を代入すれば表示した$z\cdot z$を得る。
他の積も同様である。
\end{proof}

\begin{corollary}[線束生成射影残余pairの二次復元]
\label{kj:cor:conditional-Ku}
$\lambda^3\ne-3$とし、非可換側の線束生成残余pairを
第\ref{am:sec:comparison}節の$\mathcal K_{B_R}^{\rm lin}$で定める。
このとき
\[
 (\mathcal K_{B_R}^{\rm lin},\rho_R)\simeq
 (\Ku_{\rm perf}(V(\widehat F_\lambda)/R_\lambda),\rho_{\widehat F_\lambda})
\]
であり、上の低次数積表はこのpairの積でもある。
可換側は$\operatorname{Perf}(V(\widehat F_\lambda))$内の残余成分を意味する。
\end{corollary}
\begin{proof}
定理\ref{am:thm:comparison}を非可換側と可換側に適用し、
定理\ref{kj:thm:mf-reconstruction}と合成する。
可換側の線束生成圏とperfect圏の一致は
系\ref{am:cor:formal-residual}の証明で確認する。
\end{proof}
Artin底の全有界coherent圏を、perfect圏と同一視していない。

古典Jacobian環での二次項は
\begin{equation}\label{kj:eq:scalar}
 [G_2]=c(\lambda)[xyz],\qquad
 c(\lambda)=\frac13+\frac{\lambda^3}{36}.
\end{equation}
実際、多項式恒等式
\begin{equation}\label{kj:eq:scalar-certificate}
 G_2-\left(\frac13+\frac{\lambda^3}{36}\right)xyz
 =-\frac{\lambda^2z}{36}\partial_uF_\lambda
\end{equation}
が証明書である。
さらに$u'=u-h^2\lambda^2z/36$という線形座標変換により
\begin{equation}\label{kj:eq:marked-normal-form}
 \widehat F_\lambda(x,y,z,u,v)
 =F_\lambda(x,y,z,u',v)+h^2\left(\frac13+\frac{\lambda^3}{36}\right)xyz
 \pmod{h^3}.
\end{equation}
従って$\lambda^3=-12$では、このMF pairは二次まで元の可換cubicのpairと同値になる。
この点でも古典cubicは滑らかであるが、三次以上の変形の消滅は従わない。
$\lambda$自体は切断していない。
$\lambda=0$の$c(0)=1/3$との一致は独立な整合確認であり、
係数$1/36$はその特殊化から推定したものではなく二次graphの係数から求めた値である。

\subsection{厳密計算の再現範囲}

\texttt{computations/verify\_kontsevich\_second\_jet.py}は有理数上で、
一般対称$3\times3$行列の$Q=-2\operatorname{adj}(\mathsf H)$、
混合族の三関係、三次の全35単項式の星積対称化、
生成元の125通りの結合性、中心元の像、
$5\times15$勾配行列と$10\times15$積行列の核・像を検算する。
平方小行列の行列式$9$、Euler積分、式\eqref{kj:eq:scalar-certificate}も確認し、
指定された出力先へJSON証明書を再生成する。

これらは多項式恒等式の検算であり、formality、wheel積分、
有限curved gaugeからのMF同値は計算機検証していない。
積行列のMF解釈は本文の比較定理と相対sector計算に依存する。
有限計算の再現範囲と、組版および他の証明の検証範囲は
区別して記録する。

% 2026-11-10 / v23
\section{混合中心対の有限群同変形式復元と全次数正規形}
\label{mnf:sec:equivariant-normal-form}

本節では第\ref{mt:sec:family}節の混合中心対を固定し、
形式復元定理\ref{fo:thm:formal-reconstruction}で得られる可換三次式の
全次数の選択則を証明する。
従来の対角群に$x,y$の交換と$h$の符号反転を加えることで、
二つの形式級数を$h^3$ではなく$h^6$の級数として表せる。
高次係数の明示値、収束、非零複素パラメータへの代数化、
残余対角のperfect性は依然として未解決である。

$H=\{(\alpha,\beta,\gamma,\delta,\varepsilon)\in\mu_3^5:
\alpha\beta\gamma=1\}$、
$\chi=\alpha\beta\delta$とする。$H$は
$(x,y,u,z,v,h)\mapsto
(\alpha x,\beta y,\gamma u,\delta z,\varepsilon v,\chi h)$
によって作用する。補題\ref{mt:lem:character-space}において、
この作用が非可換関係式と指定中心元を保つことを確認した。

\begin{lemma}[交換対合と非可換中心対]
\label{mnf:lem:swap-involution}
$\tau$を
\[
 \tau(x)=y,\quad \tau(y)=x,\quad
 \tau(z)=z,\quad \tau(u)=u,\quad \tau(v)=v,\quad
 \tau(h)=-h,\quad \tau(\lambda)=\lambda
\]
と定める。この作用は混合量子化$A_{h,\lambda}$の関係式
イデアルを保ち、$\Omega_{h,\lambda}$を固定する。
さらに$\tau^2=1$で、$\tau$は$H$を正規化する。
従って$G=H\rtimes\langle\tau\rangle$は中心対の有限群対称性である。
\end{lemma}
\begin{proof}
第\ref{mt:sec:family}節と同様に$a=\lambda u$とおき、
自由代数内の関係式を
\[
 r_1=xy-yx-hz^2,\quad
 r_2=xz-zx+h(y^2+ax),\quad
 r_3=yz-zy-h(x^2+ay)
\]
と書く。語の順序を変えずに$\tau$を代入すれば
\[
 \tau(r_1)=-r_1,\quad \tau(r_2)=r_3,\quad \tau(r_3)=r_2
\]
となる。式\eqref{mt:eq:center}の
$x^3+y^3$、六語の和$T$、$xy+yx$はすべて$x,y$の交換で不変で、
$h$は偶数冪でのみ現れるため$\tau(\Omega_{h,\lambda})=\Omega_{h,\lambda}$。
また$\tau$は$H$の$\alpha,\beta$を交換し、
$\chi=\alpha\beta\delta$を不変に保つので、$H$を正規化する。
以上は商を取る前の自由代数における恒等式である。
\end{proof}

\begin{lemma}[有限群同変なformalityデータ]
\label{mnf:lem:equivariant-formality}
$\lambda^3\ne-3$とする。
定理\ref{fo:thm:formal-reconstruction}に現れる
$\Phi_h,\pi_h,W_h,\Xi_h$は$G$同変に選べる。
これらの選択は、第\ref{kj:sec:second-jet}節で構成した
$h^3$を法とする生成元の同定と中心式を保って行える。
特に$W_h$は$G$不変である。
\end{lemma}
\begin{proof}
全Artin切断のPBW性は補題\ref{mt:lem:artin-pbw}にある。
まずそこで得られる次数付きの加群同定
$m:S[[h]]\to A_{h,\lambda}$を選ぶ。
$H$にはsorted語基底が適合するが、$\tau$は一般にsorted語を
sorted語に送らない。従ってsorted標識そのものが
$G$同変であると仮定してはならない。
代わりに$G$の全元による平均
\[
 m^G=\frac1{|G|}\sum_{g\in G}g_A\circ m\circ g_S^{-1}
\]
を取る。$g$の$h$への作用は半線形だが、各共役写像は
$\CC[[h]]$線形で次数を保つ。また$m^G\bmod h=\mathrm{id}$である。
各内部次数で$h$進Nakayamaを適用すると$m^G$は
$G$同変な次数付き加群同型となる。

Affine Kontsevich formalityは線形座標変換に同変であり
\cite[Section~6]{PKontsevich}\cite{MNFDolgushev}、
polydifferential cochainsから全normalized Hochschild cochainsへの
包含も同変である。
$G$は有限群、標数は零なので不変部分を取る操作は完全である。
そのため$G$不変な$h$進Maurer--Cartan変形と一入力gaugeは、
これらの$L_\infty$ quasi-isomorphismの不変部分を通して
逐次比較できる。これは非線形gaugeを単純に平均するのではなく、
各持上げ段階のcocycleとその原像を不変部分で選ぶ操作である。

二次までの指定データは、$h\pi$のKontsevich星積と、
一次生成元を同名の一次多項式へ送る同型である。
Poisson tensorの係数と$h$への作用を合わせると
これは$G$不変であり、
補題\ref{kj:prop:algebra-and-center}の中心同定も$G$同変となる。
次の段の補正を$h^3$以降に限れば、指定した二次jetを保てる。
中心のtangent mapも$G$同変であるため、
その一意な三次Casimir逆像$W_h$は$G$不変である。
さらに$\Xi_h$は補題\ref{fo:lem:constant-trivector}の
一意な定数Koszul原像だから$G$不変となる。
各Artin段階の選択を両立させて逆極限を取る。
\end{proof}

\begin{lemma}[三指標のJacobian商と交換対合]
\label{mnf:lem:character-quotients}
$S=\CC[x,y,z,u,v]$、
$J(F_\lambda)=S/(\partial F_\lambda)$とする。
$\lambda^3\ne-3$のとき
\[
\begin{array}{c|c|c}
\text{指標}&S_3\text{の指標空間}&J(F_\lambda)_3\text{の指標空間}\\ \hline
1&\langle x^3,y^3,z^3,u^3,v^3,uxy\rangle&
       \CC[uxy]\\
\chi&\langle xyz,u^2z\rangle&\CC[xyz]\\
\chi^2&\langle uz^2\rangle&0
\end{array}
\]
である。二つの非零な商空間における$\tau$作用は恒等である。
指標$1$の$S_3$の$\tau$反不変部分は
$\CC(x^3-y^3)$であり、これはJacobianイデアルに属する。
\end{lemma}
\begin{proof}
$\gamma=(\alpha\beta)^{-1}$を消去すると、
$(\alpha,\beta,\delta,\varepsilon)$について
$x,y,u,z,v$の重みはそれぞれ
$(1,0,0,0),(0,1,0,0),(2,2,0,0),(0,0,1,0),(0,0,0,1)$
である。三次単項式35個を列挙して中央列を得る。
\[
 F_x=x^2+\lambda uy,\quad F_y=y^2+\lambda ux,\quad
 F_z=z^2,\quad F_u=3u^2+\lambda xy,\quad F_v=3v^2 .
\]
指標$1$のJacobian部分空間は
$xF_x,yF_y,zF_z,uF_u,vF_v$で五次元。
指標$\chi$では$zF_u=\lambda xyz+3u^2z$、
指標$\chi^2$では$uF_z=uz^2$が関係を与える。
従って商は表示の通りである。
$uxy,xyz$は$x,y$の交換で不変であり、
指標$1$に属する六つの単項式の反不変部分は
$x^3-y^3=xF_x-yF_y$で生成される。
\end{proof}

\begin{theorem}[二つの$h^6$形式級数による正規形]
\label{mnf:thm:two-series}
$\lambda^3\ne-3$とし、混合中心対の
補題\ref{mnf:lem:equivariant-formality}による
形式的な可換三次式を$W_h$とする。
このとき$h=0$で恒等となる$G$同変な
$T_h\in\operatorname{GL}_5(\CC[[h]])$が存在し、
\begin{equation}\label{mnf:eq:normal-form}
 \boxed{\quad
 T_h^*W_h=
 \frac{x^3+y^3+z^3}{3}+u^3+v^3+
 \lambda_*(h^6)\,uxy+h^2c_*(h^6)\,xyz
 \quad}
\end{equation}
とできる。ここで
\[
 \lambda_*(q)\in\lambda+q\CC[[q]],\qquad
 c_*(q)\in\CC[[q]],\qquad
 c_*(0)=\frac13+\frac{\lambda^3}{36}.
\]
高次係数の一意性・代数性・解析的収束は主張しない。
\end{theorem}
\begin{proof}
$H$不変性と$h$の指標$\chi$から、
$W_h=\sum_{r\ge0}h^rW_r$の係数は
$W_r\in S_3^{\chi^{-r}}$となる。
$\tau$不変性からさらに
$(x\leftrightarrow y)W_r=(-1)^rW_r$が成り立つ。
補題\ref{mnf:lem:character-quotients}によれば、
$r\equiv1,5\pmod6$なら$W_r=0$、
$r\equiv3\pmod6$なら$W_r$は
$x^3-y^3$の倍数でありJacobianイデアルに属する。
また$r\equiv4\pmod6$ならJacobian類は零。
従ってJacobian商において許される非零な高次方向は
\[
 \begin{cases}
 \CC[uxy],&r\equiv0\pmod6,\\
 \CC[xyz],&r\equiv2\pmod6,\\
 0,&r\equiv1,3,4,5\pmod6.
 \end{cases}
\]
これが選択則である。

上式から形式正規化を帰納的に行う。
次数$r$未満で主張の形が得られたとする。
残余の三次式は、その次数の許された指標と$\tau$固有値を持つ。
Jacobian商への像を、$r\equiv0$なら$uxy$、
$r\equiv2$なら$xyz$の係数に吸収する。
残りは三次Jacobianイデアルに属し、
$\sum_iL_i\partial_iF_\lambda$と書ける（$L_i$は一次式）。
$G$の平均により、$h^r\sum_iL_i\partial_i$を
$G$不変な線形vector fieldとして選び直せる。
その次数$r$の形式的線形座標変換でJacobian項を消去しても、
既に固定した低次数係数は変わらない。
この操作を全次数で繰り返すと$h$進的に収束した
$G$同変の$T_h\equiv\mathrm{id}\pmod h$を得る。

初めの二次式は
$W_h=F_\lambda+h^2(xyz/3-\lambda^2u^2z/12)+O(h^3)$。
一次変換を恒等とし、
$u\mapsto u+h^2\lambda^2z/36$とすれば
$u^2z$が消え、
$xyz$係数は$1/3+\lambda^3/36$となる
（式\eqref{kj:eq:scalar-certificate}）。
この座標変換も$G$同変である。
それ以降の操作は次数三以上から始めるため、
$c_*(0)$は変化しない。
\end{proof}

\begin{remark}[最初の未決定係数と判別点]
式\eqref{mnf:eq:normal-form}で
\[
 \lambda_*(q)=\lambda+a_1q+O(q^2),\qquad
 c_*(q)=\left(\frac13+\frac{\lambda^3}{36}\right)+b_1q+O(q^2).
\]
最初の未決定係数は\emph{$h^6uxy$の$a_1$と
$h^8xyz$の$b_1$}である。
$h^3,h^4,h^5$のJacobian変形類は零であるが、
一般の非同変表示で生の係数が零とは限らない。
$\lambda^3=-12$では二次の非自明類も消え、
形式的に最初に現れ得るのは$h^6$方向となる。
これが実際に非零かどうかは計算していない。
\end{remark}

\noindent
\textbf{計算的照合：}
\texttt{python3 computations/verify\_mixed\_characters.py}は
三指標の全単項式、Jacobian関係行列の階数$(5,1,1)$、
交換対合による自由代数の三関係と中心式の不変性、
反不変三次式のJacobian恒等式、および$c_*(0)$の
多項式恒等式を$\mathbb Q[\lambda]$上で検算する。
同変formalityと全次数の正規化は数学的証明に依存し、
有限計算や外部形式検証による結論ではない。


\part{有限対角・bar積・Morita核と代数化}
% 2026-11-10 / finite-algebraic-recovery v1.0 / proposed nc-cubic v24
\section{有限の低次数積からの代数的な三次式の復元}
\label{far:sec:recovery}

本節では、代数的な有限データから可換三次式を取り出す問題を扱う。
あるformalityを選んで得た形式三次式の係数が代数的であることは要求しない。
代わりに、その勾配空間が代数的に与えられているとき、同じ勾配空間を持ち、
任意の指定された有限jetを保つ代数的三次式を構成する。
完成底上では、元の形式三次式との線形座標同値も得られる。
ただし、非可換の残余圏から必要な有限データを取り出すことと、
一般ファイバー上の圏同値を構成することは、ここでの結論とは別である。

既知の構造との関係を述べる。
可換超曲面のrotationの自然変換とJacobian環との比較には
Lin--Zhangのsector計算と積の比較を用いる
\cite[Theorem 3.5, Remark 3.6, Section 5]{FARLinZhang}。
同じ勾配空間を持つ三次式の自由度はsymmetrizerの問題であり、
その一般的背景はHwang\cite{FARHwang}にある。
完備化の忠実平坦性と有限加群の完備化は
\cite[Tags 00MC, 0BNH]{FARStacks}の標準結果である。
以下はこれらを有限復元へ適用するための直接的な代数的整理であり、
一般的な復元原理の初出を主張しない。

\subsection{勾配空間の積分と任意の有限jetの保持}

$(R,\mathfrak m)$をNoether局所$\mathbb Q$-代数、
$\widehat R$をその$\mathfrak m$進完備化、$k=R/\mathfrak m$とする。
$S_R=R[x_1,\ldots,x_N]$とし、添字の$2,3$は多項式次数を表す。
$K\subset (S_R)_2$は階数$N$の自由直和因子で、基底を
$q=(q_1,\ldots,q_N)^{\mathsf t}$とする。
次の有限加群を定める：
\begin{equation}\label{far:eq:linear-kernel}
 \mathcal L_K=\left\{A\in\operatorname{Mat}_N(R):
 \sum_a A_{ia}\partial_jq_a=\sum_a A_{ja}\partial_iq_a
 \quad(1\le i<j\le N)\right\}.
\end{equation}
これは$N^2$未知数について、高々$N\binom N2$本の斉次線形方程式の核である。

\begin{theorem}[勾配空間からの局所代数化]\label{far:thm:algebraization}
$\widehat W\in (S_{\widehat R})_3$の特殊化$F\in k[x]_3$が
幾何学的に滑らかな射影三次超曲面を定め、
\begin{equation}\label{far:eq:formal-gradient}
 K\otimes_R\widehat R
 =\langle\partial_1\widehat W,\ldots,\partial_N\widehat W\rangle
 \subset (S_{\widehat R})_2
\end{equation}
を満たすとする。このとき任意の$\ell\ge1$について、
$A\in\mathcal L_K\cap\operatorname{GL}_N(R)$を選び、
\begin{equation}\label{far:eq:euler}
 G=\frac13\sum_i x_i(Aq)_i
\end{equation}
とおくことで、次を満たす$G\in(S_R)_3$を構成できる。
\begin{enumerate}
\item $\nabla G=Aq$であり、その成分は$K$の基底である。
\item $G\equiv\widehat W\pmod{\mathfrak m^\ell\widehat R}$である。
\item $V(G)\subset\mathbb P^{N-1}_R$は$R$上滑らかである。
\item $T\in\operatorname{GL}_N(\widehat R)$、
$T\equiv I\pmod{\mathfrak m^\ell\widehat R}$が存在し、
$\widehat W(x)=G(Tx)$となる。
\end{enumerate}
この方程式の構成には、$\widehat W$の係数の代数性、解析的収束、
非可換圏の相対対角のperfect性を仮定しない。
\end{theorem}

\begin{proof}
滑らかさから$F$の偏微分は$k$上一次独立である。
実際、定数係数の線形関係があれば座標変換後に$F$は一変数に依存せず、
その変数軸の射影点は特異点となる。
仮定\eqref{far:eq:formal-gradient}とNakayamaの補題から、
$\nabla\widehat W=\widehat A q$を満たす
$\widehat A\in\operatorname{GL}_N(\widehat R)$がある。
混合偏微分が交換するので、$\widehat A$は
\eqref{far:eq:linear-kernel}の完成底上の解である。

$R\to\widehat R$は平坦だから、線形写像の核について
\[
 \mathcal L_K\otimes_R\widehat R=\mathcal L_{K\otimes\widehat R}
\]
が成り立つ。$R$はNoether環なので$\mathcal L_K$は有限生成である。
生成元を$A_1,\ldots,A_r$とし、
$\widehat A=\sum_\nu\widehat c_\nu A_\nu$と書く。
$R/\mathfrak m^\ell\simeq\widehat R/\mathfrak m^\ell\widehat R$より、
$c_\nu\in R$を$\widehat c_\nu$の$\ell$次未満の近似に取れる。
$A=\sum c_\nu A_\nu$は正確に$\mathcal L_K$に属し、
$A\equiv\widehat A\pmod{\mathfrak m^\ell}$である。
特に$\det A$は単元である。
ここで解加群の特殊化次元が一定であることは仮定していない。

$g=Aq$とおく。可積分性と二次のEuler等式から
\[
 \partial_j\left(\frac13\sum_i x_ig_i\right)
 =\frac13\left(g_j+\sum_i x_i\partial_jg_i\right)
 =\frac13\left(g_j+\sum_i x_i\partial_i g_j\right)=g_j.
\]
これが第一の主張を与える。三次のEuler等式を
$\widehat W$にも適用すると第二の主張が従う。

滑らかさは有限の次数で検出できる。
$F$の$N$個の偏微分は正則列で、その商のHilbert級数は$(1+t)^N$である。
従って$\overline q_i$の$k[x]_{N-1}$倍は$k[x]_{N+1}$を張る。
有限自由加群間の写像
\begin{equation}\label{far:eq:macaulay}
 (S_R)_{N-1}^{\oplus N}\longrightarrow(S_R)_{N+1},
 \qquad (b_i)\longmapsto\sum_i b_iq_i
\end{equation}
は特殊化で全射なので、Nakayamaにより全射である。
各ファイバーで全$q_i$が零となる射影点は存在しない。
$\nabla G=Aq$と$A$の可逆性から$G$にも同じことがいえ、
射影アフィンchartのJacobian判定で$V(G)\to\operatorname{Spec}R$は滑らかである。

最後に$B=\widehat A A^{-1}$とすると、
$\nabla\widehat W=B\nabla G$、$B\equiv I\pmod{\mathfrak m^\ell}$である。
$G(x)=\mathcal T(x,x,x)$と書く対称三重線形形式について、
$U=B^{\mathsf t}$は
\[
 \mathcal T(Ua,b,c)=\mathcal T(a,Ub,c)=\mathcal T(a,b,Uc)
\]
を満たす。これは$\widehat W$の偏極の対称性から従う。
$U$の多項式、およびその$\mathfrak m$進極限も同じ条件を満たす。
よって
\[
 T=\sum_{j\ge0}\binom{1/3}{j}(U-I)^j
\]
は定義でき、$T^3=U$である。
各引数間で$T$を移すと
\[
 G(Tx)=\mathcal T(Tx,Tx,Tx)
 =\mathcal T(Ux,x,x)=\widehat W(x).
\]
また$T\equiv I\pmod{\mathfrak m^\ell}$である。
\end{proof}

\begin{remark}[代数化されるもの]\label{far:rem:scope}
定理は$G$が$R$上に存在することを示すが、
座標変換$T$が$R$上に定義されるとは主張しない。
$\widehat W$自体も代数的である必要はない。
例えば$e^{3h}x_1^3+\sum_{i>1}x_i^3$は
$x_1\mapsto e^h x_1$によって定数族と形式同型である。
さらに、非可換pairと$\widehat W$のpairの形式同値を仮定すれば、
これを$G$のpairへの形式同値へ置き換えられるが、
その同値の非零パラメータへの代数化は別問題である。
\end{remark}

\subsection{低次数積からの適用条件}

\begin{corollary}[代数的な低次数積があれば十分]\label{far:cor:product}
$V_R,W_R$を有限自由$R$-加群、
$\mu:V_R\otimes_RV_R\to W_R$を$R$-双線形写像とする。
完成底上に、$V_R\otimes\widehat R\simeq(S_{\widehat R})_1$と
$W_R\otimes\widehat R\simeq\operatorname{Jac}(\widehat W)_2$という同定があり、
$\mu\otimes\widehat R$がJacobian環の積となると仮定する。
$\widehat W$の特殊ファイバーは定理\ref{far:thm:algebraization}と同様に滑らかとする。
このとき$\mu$は対称であり、
\[
 \overline\mu:\operatorname{Sym}^2V_R\longrightarrow W_R
\]
は全射である。その核$K$は階数$N=\operatorname{rank}V_R$の自由直和因子となり、
定理\ref{far:thm:algebraization}を適用できる。
完成底での一次部分の同定を用いて、$\widehat W$の座標も変更してよい。
\end{corollary}
\begin{proof}
$\mu-\mu^{\mathrm{op}}$は忠実平坦な底変更後に零だから零である。
$\operatorname{coker}\overline\mu$も底変更後に零だから零である。
$W_R$が自由なので全射は分裂し、核は有限射影、従って局所環上自由である。
核の底変更は勾配空間に一致し、その階数は$N$である。
\end{proof}

\begin{proposition}[有限複体からの降下]\label{far:prop:finite-model}
非可換pairに対し、次数$1,2$の導来自然変換の$H^0$を計算する
有限自由三項複体
\[
 C_j^{-1}\longrightarrow C_j^0\longrightarrow C_j^1\qquad(j=1,2)
\]
と、$H^0$上で合成を誘導する代数的な写像が与えられたとする。
これらの複体が実際の導来自然変換を計算すること、
完備化との比較が合成を保つことを仮定する。
完成底上でそれらの$H^0$と積が系\ref{far:cor:product}のJacobianデータに一致すれば、
$R$上で同系の入力を構成できる。
\end{proposition}
\begin{proof}
$H^0(C_j)$は有限加群で、平坦な底変更についてコホモロジーは可換である。
完成底上で有限自由となるので、忠実平坦降下により$R$上有限射影、従って自由となる。
合成を$H^0$に降ろして系\ref{far:cor:product}を使う。
\end{proof}

\begin{remark}[未構成の入力]
非可換圏がproperであることや、各Artin切断で可換模型を持つことだけから、
命題\ref{far:prop:finite-model}の有限複体と比較を得たことにはならない。
特に、bar複体を有限長で切った行列や、線束間のHomの階数表を、
導来自然変換の積そのものと同一視してはならない。
この命題は非可換残余対角のperfect性を仮定しない条件付き経路であり、
その有限模型を無条件に構成した定理ではない。
可換cubic $N\ge5$について次数$1,2$の同定は
\cite[Theorem 3.5, Remark 3.6]{FARLinZhang}のsector計算による。
非恒等sectorはこの二つの次数には寄与しないが、全次数の同定とは区別する。
\end{remark}

\subsection{一変数底での有限アルゴリズム}

以下$R=k[h]_{(h)}$、$k$は計算可能な標数零の体とする。
$k=\mathbb Q(\lambda)$等も許すが、この場合は$\lambda$についての一般点の計算であり、
分母が消える特殊値への一様な主張はしない。

\begin{lemma}[核の局所飽和]\label{far:lem:saturation}
$E$を$R$係数の有限行列とする。
$k(h)$上の核の基底を取り、各列の分母を払って多項式列を得る。
これを並べた行列を$B$とする。
$B(0)$の列が従属なら、$0\ne c\in\ker B(0)$と$c_j\ne0$を取り、
第$j$列を$Bc/h$で置き換える。
この操作は有限回で終了し、得られた列は$\ker_R E$の$R$-基底となる。
\end{lemma}
\begin{proof}
$Bc$は$h$で割り切れるので各段階の列は多項式である。
列の外積は$c_j/h$倍されるから、その全Pl\"ucker座標に共通する
$h$の最大冪の指数は一ずつ減る。
指数は非負整数なので停止する。停止時にはある最大minorが単元である。
従って列は$R^n$の直和因子を張る。
$k(h)$上の列空間は最初の核と同じであり、単元minorで座標を解けば、
その空間に属する任意の$R^n$の元の係数は$R$に属する。
これが核全体との一致を示す。
\end{proof}

\begin{proposition}[指定特殊ファイバーを保つ実行手順]\label{far:prop:algorithm}
$\overline\mu$の行列と滑らかな特殊三次式$F$が与えられ、
$F$の勾配が特殊化した核の基底になるとする。
次の有限手順は、$G\bmod h=F$を満たす局所可逆な積分解の存在を判定し、
存在すればその一つを有理関数係数で構成する。
\begin{enumerate}
\item $\overline\mu(0)$の非零最大minorを選び、その列をpivotとして核の基底$q$を作る。
\item $\nabla F=A_0q(0)$を満たす可逆行列$A_0$を線形代数で求める。
\item \eqref{far:eq:linear-kernel}の行列$E_K$に補題\ref{far:lem:saturation}を適用し、基底$B$を求める。
\item $B(0)c_0=\operatorname{vec}(A_0)$を解く。解がなければ指定特殊解の持上げは存在しない。
\item $\operatorname{vec}(A)=Bc_0$とし、\eqref{far:eq:euler}を用いる。
\end{enumerate}
定理\ref{far:thm:algebraization}の形式的入力が存在すれば第四段階は成功する。
\end{proposition}
\begin{proof}
最初のminorは$R$の単元なので、核の基底は$h=0$で正則である。
飽和基底は解加群全体の基底だから、第四段階は特殊化$A_0$を持つ解の存在と同値である。
最後の$A$は$A_0$に特殊化するため可逆であり、Euler積分の証明を適用できる。
\end{proof}

\begin{example}[一般点の核だけでは不十分]
$E=(h\ \ 1\ \ 1)$の核に対して、分母を払った二列
$(-1,h,0)^{\mathsf t}$、$(-1,0,h)^{\mathsf t}$は
$h=0$で従属となる。その差を$h$で割ると
$(0,-1,1)^{\mathsf t}$を得る。
この飽和により、例えば特殊解$(0,1,-1)^{\mathsf t}$の持上げを正しく検出できる。
一方$E=(h\ \ 0)$では特殊解$(1,0)^{\mathsf t}$は持ち上がらない。
特殊化した線形方程式を解くだけでは、この区別を検出できない。
\end{example}

\begin{remark}[検証器の範囲]
\texttt{computations/finite\_recovery.py}は以上の手順を実装する。
入力は\emph{既に与えられた}対称積の行列であり、非可換代数からその行列を作るコードではない。
出力時に核、可積分性、勾配、特殊化を厳密に検査する。
滑らかさは特殊勾配イデアルの零次元性で確認する。
理論上は\eqref{far:eq:macaulay}の単元minorも有限証明書となる。
$N=5$の場合、この行列のサイズは$210\times350$である。
一般アルゴリズムの最適な計算量は主張しない。
\end{remark}

\section{混合候補の二パラメータ族と厳密な積表}
\label{far:sec:mixed}

次の可換族を、非可換混合族との比較に用いる\emph{候補族}として考える：
\begin{equation}\label{far:eq:family}
 H_{L,C}=\frac{x^3+y^3+z^3}{3}+u^3+v^3+L uxy+C xyz.
\end{equation}
本節はこの可換族自身の計算であり、$L,C$を非可換入力から決定したものではない。
特に、以下の積表を実際の非可換rotation積と同定する比較は未構成である。

\subsection{全パラメータでの積と可積分性}

Jacobian商の次数二において
\begin{align*}
 x^2&=-Cyz-Lyu,& y^2&=-Cxz-Lxu,& z^2&=-Cxy,\\
 u^2&=-\tfrac L3xy,& v^2&=0.
\end{align*}
従って
\[
 (xy,xz,xu,xv,yz,yu,yv,zu,zv,uv)
\]
は次数二商の基底である。
始域を$(x^2,y^2,z^2,u^2,v^2)$とこの十単項式の順に取ると、
積は$[D\mid I_{10}]$であり、$D$の非零行は
\[
\begin{array}{c|rrrrr}
 &x^2&y^2&z^2&u^2&v^2\\\hline
xy&0&0&-C&-L/3&0\\
xz&0&-C&0&0&0\\
xu&0&-L&0&0&0\\
yz&-C&0&0&0&0\\
yu&-L&0&0&0&0
\end{array}
\]
である。残り五行は零である。この表示は全パラメータで階数十を持つ。

\begin{proposition}[混合領域の積分解空間]\label{far:prop:symmetrizer}
$q=\nabla H_{L,C}$とし、$LC\ne0$とする。
可積分行列の空間は正確に
\[
 \left\{\operatorname{diag}(a,a,a,a,b):a,b\in\mathbb C\right\}
\]
である。従って勾配空間を保つ三次式は
$a(H_{L,C}-v^3)+bv^3$の形である。
可逆な積分解の条件は$ab\ne0$である。
\end{proposition}
\begin{proof}
表示した二次元空間は直接代入で可積分である。
逆に$50\times25$の可積分性行列$E$は、次のminorによって階数二十三以上となる。
変数順を$(x,y,z,u,v)$とし、行は$(i,j)$を辞書式に並べ、
各対について係数変数を同じ順に並べる。列は$A_{ia}$を行優先で並べる。
\emph{零始まり}で行
\[
\begin{gathered}
0,1,2,5,6,7,8,10,11,12,13,15,16,17,19,22,23,34,37,38,42,44,49
\end{gathered}
\]
と列
\[
0,1,2,3,4,5,6,7,8,9,10,11,12,13,14,15,16,17,19,20,21,22,23
\]
を選んだ行列式は
\begin{equation}\label{far:eq:minor}
 5971968\,C^5L^5
\end{equation}
である。各成分は\eqref{far:eq:linear-kernel}の微分から指定されるため、
これは有限の多項式恒等式による証明書である。
添付コードは有理係数でこの行列式を展開して検査する。
核の次元は二以下であり、既に二次元の解を得ているので等号となる。
\end{proof}

$L=1,C=0$では可積分性行列の階数は二十二、$LC\ne0$では二十三である。
従って、混合候補の特殊点で解空間が跳ぶことは実際に起こる。
定理\ref{far:thm:algebraization}と局所飽和の手順は、この跳びを許している。
軸上でのすべての値を一括して一定次元とは扱わない。
例えば$L=0,C=2$では核の次元は五になる。

\subsection{滑らかさの必要十分条件}

\begin{proposition}[候補族の滑らかさ判別]\label{far:prop:discriminant}
複素数$L,C$について$V(H_{L,C})\subset\mathbb P^4$が滑らかであることと、
\begin{equation}\label{far:eq:discriminant}
 \Delta(L,C)=(L^3+3+3C^3)^2-12L^3C^3\ne0
\end{equation}
は同値である。ここで$\Delta$は滑らかさを判定する多項式であり、
五変数三次式の普遍判別式をこの族へ制限したものの重複度まで同定してはいない。
\end{proposition}
\begin{proof}
偏微分が同時に零なら$v=0$である。
さらに$xy=0$なら$x=y=z=u=0$となるので、射影特異点では$xy\ne0$である。
$xH_x-yH_y=x^3-y^3$より、$x=1$とし$y=\zeta$、$\zeta^3=1$とできる。
$z=\zeta^2z_0,u=\zeta^2u_0$とおくと、条件は
\[
 z_0^2=-C,\qquad u_0^2=-L/3,\qquad 1+Cz_0+Lu_0=0
\]
となる。$a=Lu_0,b=Cz_0$とおけば
$a^2=-L^3/3,b^2=-C^3,a+b=-1$である。
従って$4a^2b^2=(1+L^3/3+C^3)^2$から$\Delta=0$を得る。

逆に$LC\ne0$かつ$\Delta=0$とする。
\[
 a=\frac{L^3-3-3C^3}{6},\qquad
 b=\frac{-3-L^3+3C^3}{6}
\]
とおくと、$a+b=-1$であり、
$a^2+L^3/3=b^2+C^3=\Delta/36=0$である。
よって$u_0=a/L,z_0=b/C$が所要の解を与える。
$C=0$では条件は$L^3=-3$となり、$u_0=-1/L,z_0=0$が解である。
$L=0$では$C^3=-1$となり、$z_0=-1/C,u_0=0$が解である。
$L=C=0$では$\Delta=9\ne0$である。
Euler等式から偏微分の共通零点は超曲面上にあり、必要十分性が従う。
\end{proof}

\subsection{有限jetが決めないもの}

混合量子化で既に得られている二次式は、座標変換後に
\[
 H_{\lambda,\,h^2c_0(\lambda)}\pmod{h^3},
 \qquad c_0(\lambda)=\frac13+\frac{\lambda^3}{36}
\]
となる。これと形式的選択則だけから全次数の係数は決まらない。
実際、任意の$a,b\in\mathbb C$について
\[
 L(h)=\lambda+ah^6,\qquad
 C(h)=h^2\bigl(c_0(\lambda)+bh^6\bigr)
\]
は、同じ二次jetと同じ冪の選択則を持つ代数的な候補を与える。
$\lambda^3\ne-3$なら$\Delta(L(0),C(0))=(\lambda^3+3)^2\ne0$なので、
すべて零の近傍で滑らかである。
これは実際の非可換入力に複数の復元先があるという主張ではない。
未計算の低rotation積が$a,b$を区別する必要がある、という注意である。

検算では、人工的な可換入力
$L=1+2h^6$、$C=13h^2/36+3h^8$の積の終域基底をさらに変更し、
その行列だけから命題\ref{far:prop:algorithm}で両係数を全次数で回収した。
非可換由来の積を計算したという意味ではない。
実際の混合中心対に対する未解決の入力は、
命題\ref{far:prop:finite-model}の有限複体、合成、および完成底との比較である。

% 2026-11-11 / v27; finite curved diagonal and genuine noncommutative tests.
\section{曲率付き対角からの有限計算}
\label{cd:sec:finite-diagonal}

本節では、代数的非可換入力から有限の行列を実際に作る。
結論の範囲を最初に固定する。中心差を用いた対角の行列因子分解と、
その自己準同型の積は構成・計算する。一方、それが射影残余圏の
導来自然変換の積に一致することは本節の結論ではない。
この区別は、曲率付き双加群の導来圏と包絡代数上の導来圏が一般には異なるという
Guan--Holstein--Lazarev--Rasmussenの指摘\cite[Theorem~4.7, Theorem~6.14, Remark~6.16]{CDGHLR}
にも対応する。

\subsection{有界分解の有限曲率化}
$R$を可換$\mathbb Q$-代数、$A$を非負内部次数付き$R$-代数とする。
各$A_d$は有限自由、$A_0=R$と仮定する。
右$Q=A^{\mathrm{op}}\otimes_R A$-加群としての対角$A$に、
有限自由な線形分解
\[
 0\longrightarrow P_N\longrightarrow\cdots\longrightarrow P_0
 \xrightarrow{\epsilon} A\longrightarrow0,\qquad
 P_p=E_p\otimes_R Q,\quad \deg E_p=p
\]
が与えられたとする。$E_p$は有限自由である。
中心三次元$\Omega\in A_3$について
$f=\Omega^{\mathrm{op}}\otimes1-1\otimes\Omega$と置く。
これは$Q$の中心に属し、対角$A$を零化する。

\begin{theorem}[有限個の高次ホモトピー]
\label{cd:thm:finite-curvature}
内部次数$3j$、鎖次数$2j-1$の写像
$s_j:P_p\to P_{p+2j-1}$を、
\[
 ds_1+s_1d=f\,\mathrm{id},\qquad
 ds_j+s_jd=-\sum_{a=1}^{j-1}s_as_{j-a}\quad(j\ge2)
\]
となるように選べる。
$j>\lfloor(N+1)/2\rfloor$では$s_j=0$としてよい。
$s_j$の行列係数は$Q_{j+1}$に属する。
$P_p$をコホモロジー次数$-p$に置き、
$\deg(t)=(2,-3)$とすれば
\[
 D=P\otimes_R R[t,t^{-1}],\qquad
 \mathcal D=d+\sum_{j\ge1}t^js_j,
 \qquad \mathcal D^2=tf\,\mathrm{id}
\]
を得る。この和は有限で、$N=5$なら$s_1,s_2,s_3$だけである。
\end{theorem}
\begin{proof}
$\operatorname{End}_Q(P)$には鎖次数を一つ下げる微分
$\delta T=dT-(-1)^{|T|}Td$を入れる。
$f\,\mathrm{id}$が誘導する$A$上の写像は零なので、
有界射影分解の比較定理により$\delta s_1=f\,\mathrm{id}$を解ける。

$s_1,\ldots,s_{j-1}$が構成されたとする。
$O_j=\sum_{a=1}^{j-1}s_as_{j-a}$の鎖次数は$2j-2$である。
既存の式を代入して$\delta O_j$を展開すると、三重積は相殺し、
残る$f s_{j-1}-s_{j-1}f$も中心性によって零となる。
鎖次数が正であるため、対応する負次数の
$\operatorname{Ext}_Q(A,A)$は零であり、$\delta s_j=-O_j$を解ける。
内部次数ごとにこの議論を行えば、$s_j$を内部次数$3j$に選べる。
$p$次の生成元から$p+2j-1$次の生成元への係数次数は
$3j-(2j-1)=j+1$である。

$2j-1>N$ならその鎖次数の写像は存在しない。
このとき$O_j$の鎖次数も$N$を超えれば自明に零である。
残る可能性は$2j-2=N$の場合だけだが、そこでの閉元は境界であり、
境界を与える鎖次数$N+1$の写像空間が零なので$O_j=0$となる。
これで有限和の全恒等式が成立する。
\end{proof}

この構成は高次ホモトピーの標準的な持上げ議論であり、原理そのものの新規性を主張しない。
本研究で利用する点は、内部次数によって必要な係数空間と計算上限が固定されることである。
高次ホモトピーを特殊ファイバーでだけ作るのではなく、実際の$R$上で解く。

\subsection{小さい複体と、積を計算する大きい複体}
$A[t,t^{-1}]$には零微分を入れる。$f$はこの加群を零化するので、
\[
 C_r=\operatorname{Hom}_Q(D,A[t,t^{-1}])_r,
 \qquad d_C\phi=-(-1)^{|\phi|}\phi\mathcal D
\]
は複体である。ここで添字$r$は内部次数である。
\begin{proposition}[各双次数の有限性]
\label{cd:prop:small-complex}
\[
 C_r^q=\bigoplus_{\substack{0\le p\le N\\p\equiv q\ (2)}}
 A_{\,r+(3q-p)/2}\otimes_R E_p^\vee,
 \qquad A_d=0\ (d<0).
\]
一方、$\mathcal E_r=\operatorname{Hom}_Q(D,D)_r$の$(p,s)$ブロックは
\[
 Q_{\,r+(3q-p+s)/2}\otimes_R
 \operatorname{Hom}_R(E_p,E_s),\qquad q-p+s\equiv0\pmod2
\]
である。増大射は擬同型$\mathcal E_r\to C_r$を誘導する。
従って$H^0(C_r)$の元を$\mathcal E_r$の閉元へ持ち上げ、行列を合成して
再び$C$へ送ることで、次数付き代数$H^0(\mathcal E_\bullet)$の積を計算できる。
\end{proposition}
\begin{proof}
$P_p$上の写像が$t^a$を含むとき、全コホモロジー次数は$q=p+2a$である。
内部次数$r$の条件は係数の次数を$r+p+3a$にするので、最初の式を得る。
$P_s$を値域とする場合は$q=p-s+2a$となり、同様に第二の式を得る。
$p,s$の範囲は有界だから各項は有限自由である。

$t$の指数でフィルターすると、$d$による微分は指数を保ち、
各$t^js_j$は指数を正に増やす。
固定したコホモロジー次数・内部次数では、とり得る$t$の指数は有限集合である。
増大射のconeの随伴次数複体は
\[
 \operatorname{Hom}_Q(P,P)\longrightarrow
 \operatorname{Hom}_Q(P,A)
\]
のconeの内部次数成分と、その$t$倍からなる。
$P$が有界射影分解であるからこれは非輪状である。
各次数で有限なフィルトレーションに沿って閉元を順に消去することにより、
曲率補正後のconeも非輪状となる。
積は$\mathcal E$の通常の合成から運ぶので、持上げの選択によらず結合的である。
この議論は$C$上の係数を単純に掛けることを正当化するものではない。
\end{proof}

$N=5$、$\operatorname{rank}E_p=\binom5p$、
$\dim A_d=\binom{d+4}{4}$なら、必要な小複体は
\[
\begin{array}{c|ccc}
 r& C_r^{-1}&C_r^0&C_r^1\\\hline
 1&0&15&126\\
 2&5&70&330
\end{array}
\]
である。これらはコチェイン空間の次元であり、$H^0$の次元ではない。
$\mathcal E$による閉元の持上げ・合成も含め、必要な係数は$A$の内部次数六以下で扱える。
実際、$s_j$の構成に必要な係数次数は最大五であり、
$r=1,2$、$q=-1,0,1$の$\mathcal E_r^q$に現れる$Q$の係数次数は最大六である。
これは分解の全次数での完全性を有限次数の検査で置き換える主張ではない。

\subsection{混合非可換入力と相対的な実計算}
$u$を中心変数とし、$K=\mathbb Q(h)$上で
\[
 [y,z]=h(x^2+\lambda uy),\qquad
 [z,x]=h(y^2+\lambda ux),\qquad [x,y]=hz^2
\]
を課す。$a=\lambda u$、$T$を$xyz$の六語の和として
\[
 \Omega_{\mathrm{act}}
 =\frac{x^3+y^3+z^3}{3}-\frac{h^2}{9}T
 +\frac a2(xy+yx)+\frac{h^2a^2}{6}z+\frac{h^4a}{18}z^2
\]
を指定する。これは現行の混合試験族の入力である。
相対計算では$Q=A^{\mathrm{op}}\otimes_{K[u]}A$を使う。
三つの関係式を$r_x,r_y,r_z$とすれば、分解の階数は$(1,3,3,1)$で、
$d_1=(x_L-x_R,y_L-y_R,z_L-z_R)$、$d_2$は関係式の非可換差分、
$d_3=(x_L-x_R,y_L-y_R,z_L-z_R)^{\mathsf t}$となる。
非可換差分は、各語$w=a_1\cdots a_m$を
\[
 \partial^{\mathrm{nc}}_i(w)
 =\sum_{a_j=x_i}(a_1\cdots a_{j-1})_L
                  (a_{j+1}\cdots a_m)_R
\]
へ送る操作であり、$u$は係数として扱う。

この相対分解の完全性は、既存の全Artin PBW入力からも確認できる。
$\mathbb Q[h]_{(h)}$上の各内部次数成分は有限自由であり、
$h=0$では中心$u$を底に置いた通常のKoszul対角分解となる。
増大射を含む有界複体に次数別のNakayamaを適用して完全性を得た後、
$h$を反転する。PBW性を本節の有限数値検査だけから推論しない。

\begin{proposition}[相対対角の低次数積の厳密計算]
\label{cd:prop:relative-formula}
上の相対分解から構成した$C$について、一般の$h,\lambda$で
$\dim H^0(C_1)=4$、$\dim H^0(C_2)=7$である。
次数一の閉元をその$P_0$成分が$u,x,y,z$となるように規格化すると、
積は対称かつ全射であり、その核は次の三次式の
$x,y,z$偏微分によって張られる：
\begin{align*}
 G^{\mathrm{rel}}_{h,\lambda}={}&
 \frac{x^3+y^3+z^3}{3}+\frac{h^2}{3}xyz
 +\lambda\left(1+\frac{h^6}{54}\right)uxy\\
 &-\frac{h^4\lambda}{18}uz^2
 -\frac{h^2\lambda^2}{12}u^2z.
\end{align*}
これは特定の$h$での補間ではなく、$\mathbb Q(h)$上の有限行列計算である。
\end{proposition}
\begin{proof}[計算証明と再現手順]
$\lambda=1$について、付属コードは定義された関係式の次数別商から積を構成する。
$d^2=0$、$ds_1+s_1d=f$、
$ds_2+s_2d=-s_1^2$を有理関数体上で検査する。
小複体の次元は$r=1$で$(0,7,34)$、$r=2$で$(3,22,70)$であり、
微分の階数は順に$(0,3)$、$(3,12)$となる。
次数一の四閉元を実際の閉行列に持ち上げ、全十六個の合成を求める。
反対称積が境界となることと、対称積の階数七を検査した結果、その核は
\[
\begin{split}
 x^2+\frac{h^2}{3}yz+\left(1+\frac{h^6}{54}\right)uy,\
 y^2+\frac{h^2}{3}xz+\left(1+\frac{h^6}{54}\right)ux,\
 z^2+\frac{h^2}{3}xy-\frac{h^4}{9}uz-\frac{h^2}{12}u^2
\end{split}
\]
となる。これらは表示した三次式の偏微分である。
一般の$\lambda\ne0$では中心変数の置換$u\mapsto\lambda u$を用いる。
特殊層での次元の不変性は、この置換だけから主張しない。
証明書は分解・高次ホモトピー・閉元の持上げと積の有限データを保持する。
\end{proof}

\begin{remark}[中心の三次式は相対対角からは見えない]
$\Omega$に$g(u)$を加えても、相対包絡代数内では
$g(u)_L-g(u)_R=0$である。
従って相対計算だけでは$u^3$の係数を決められない。
特に$G^{\mathrm{rel}}+u^3+v^3$を、元の非可換cubicの復元先と宣言してはならない。
また表示の$h^6uxy$係数を、そのまま既存の同変正規形の未知係数と同一視しない。
絶対対角、中心方向の微分、座標規格化の比較が別に必要である。
\end{remark}

\subsection{絶対模型への接続と比較条件}
絶対模型では$u$を底に含めず、四生成元の二次関係の空間$W_2$から
\[
 W_3=(V\otimes W_2)\cap(W_2\otimes V),\qquad
 W_4=(V\otimes W_3)\cap(W_3\otimes V)
\]
を計算する。次元は$(1,4,6,4,1)$である。
中心元には$u^3$も含める。
独立な$v$と$v^3$は一変数のKoszul因子とのテンソル積で加えられる。
\begin{proposition}[絶対対角から得られる全次数の三次式]
\label{cd:prop:absolute-formula}
$\lambda=1$、$K=\mathbb Q(h)$とし、中心元を
$\Omega_{\mathrm{act}}+u^3$に固定する。
絶対模型の$H^0(C_1),H^0(C_2)$の次元はそれぞれ四、六である。
次数一の元を$P_0$成分$u,x,y,z$で規格化すると、その対称積の核は
\begin{align*}
 q_u&=\left(3+\frac{h^6}{54}\right)u^2-\frac{h^2}{6}uz+xy,\\
 q_x&=x^2+\frac{h^2}{3}yz+\left(1+\frac{h^6}{54}\right)uy,\\
 q_y&=y^2+\frac{h^2}{3}xz+\left(1+\frac{h^6}{54}\right)ux,\\
 q_z&=z^2-\frac{h^4}{18}uz-
              \left(\frac{13h^2}{12}+\frac{h^8}{162}\right)u^2
\end{align*}
で張られる。$B(h)=1+29h^6/648+h^{12}/4374$と置くと、
\begin{align}
G_h={}&\frac{x^3+y^3+z^3}{3}+\frac{h^2}{3}xyz
 +\left(1+\frac{h^6}{54}\right)uxy\notag\\
 &-\frac{h^4}{18}uz^2-\frac{h^2}{12}u^2z+B(h)u^3
 \label{cd:eq:absolute-polynomial}
\end{align}
はこの核を勾配空間とする。
\end{proposition}
\begin{proof}[計算証明および積分]
絶対Koszul空間は二次関係の交叉から計算する。各空間を、昇順の語の係数が
単位行列となるように規格化する。
付属コードは$\mathbb Q(h)$上で全ての曲率恒等式を検査した後、
四つの次数一閉元を実際の閉じた自己準同型へ持ち上げる。
持上げは$t$の指数と入力の鎖次数に沿う有限の線形方程式で行い、
構成後に$\mathcal DF-F\mathcal D=0$を全ブロックで再検査する。
全十六個の積を合成して増大射で$C_2$へ送り、境界で割る。
小複体の次元は$r=1$で$(0,10,56)$、$r=2$で$(4,35,120)$、
微分の階数はそれぞれ$(0,6)$、$(4,25)$である。
積の反対称部分が境界であること、対称積の階数が六であること、
表示した四つの二次式がその核の基底であることを厳密に検査する。
この計算には、三次式$G_h$を入力していない。
計算後の積分検査においてのみ右辺の候補を使う。

積分は手計算でも確認できる。$L=1+h^6/54$として
\[
 \partial_xG_h=q_x,\quad \partial_yG_h=q_y,\quad
 \partial_zG_h=q_z+\frac{h^2}{3}q_u,\quad
 \partial_uG_h=Lq_u-\frac{h^4}{18}q_z.
\]
この基底変換の行列式は
\[
 \det A=L+\frac{h^6}{54}=1+\frac{h^6}{27}\ne0
 \quad\text{in }\mathbb Q(h).
\]
従って表示した三次式の勾配空間は求める核に等しい。
これは有限個の有理値による補間ではない。
\end{proof}

\begin{corollary}[五変数の代数的候補]
上の絶対対角に、独立な$v$と曲率$v^3$の一変数因子をテンソルする。
低次数の自己準同型代数の次元は五、十となり、その対称積の核は
$G_h+v^3$の勾配空間である。
$V(G_h+v^3)\subset\mathbb P^4$は$h=0$を含む非空開集合上で滑らかである。
\end{corollary}
\begin{proof}
一変数の小複体は、可逆な周期変数を除いて$3v^2$によるKoszul複体であり、
そのコホモロジーは$K[v]/(v^2)$である。
四変数模型とのテンソル積とK\"unnethにより、次数一には$v$の一方向、
次数二には$v$と既存の四方向の積だけが加わる。
従って核には$v^2$が加わる。これは$G_h+v^3$の勾配空間である。
$h=0$は元の滑らかな$F_1$である。
実際、全偏微分が零なら$z=v=0$であり、$u=0$なら$x=y=0$、
$u\ne0$なら$x^2=-uy,y^2=-ux,xy=-3u^2$から
$9u^4=-3u^4$となって矛盾する。
射影族の滑らかさの開性から一般の$h$での滑らかさを得る。
\end{proof}

\begin{remark}[generic計算と局所正則化]
今回保存した行列の一部には$h=0$で極がある。
特に絶対模型の高次ホモトピー、閉元の持上げ、商の基底には
負の$h$冪が現れる。次数一の選んだ閉元と最終的な三次式が正則であることと、
選んだ計算模型全体が$h=0$で正則であることは別である。
本節の積計算の底は$\mathbb Q(h)$であり、これをそのまま
$\mathbb Q[[h]]$へテンソルして比較したとは主張しない。
局所正則な模型の選び直しと低次数コホモロジーの全Artin底変更は、
後続の局所正則化の節で実行する。残余圏および指定形式模型との比較は未完成である。
\end{remark}

\begin{remark}[何を決定したか]
式\eqref{cd:eq:absolute-polynomial}は、実際の非可換関係式から作った
\emph{曲率付き対角の自己準同型積}の核を積分したものである。
指定されたformalityが与える$W_h$との形式線形同値や、
非可換射影残余圏とのgeneric同値を証明したわけではない。
二次jetは既存の式と一致するが、その一致だけで全次数の比較を結論しない。
絶対計算は$\lambda=1$についての恒等式であり、任意の$\lambda$への公式は未確定とする。
\end{remark}

最終的に必要なのは、上の行列代数から得た積と
\[
 \bigoplus_{r=0}^2H^0\operatorname{RHom}_{\mathcal C^e}
        (\Delta,P^{\otimes^{\mathbf L}r})
\]
との合成を保つ比較である。内部shiftとの両立と、同変化に現れる非恒等sectorの扱いも必要である。
Guanほか\cite[Theorem~6.14]{CDGHLR}の比較は、包絡代数の有限ホモロジー次元、
coherence、外積に関するCondition~1を仮定する。
同論文Remark~6.16は最後の条件を対角対象に関する条件へ弱めるが、
本研究の一般混合非可換族でその条件を確認したことにはならない。
同論文Theorem~4.7が無条件に扱うHochschild複体と、
有限対角から得る複体を、名前が似ているという理由で同一視しない。
本節の有限計算はこの比較問題を残しつつ、代数的入力からの方程式抽出に
具体的な計算対象を与えるものである。

% 2026-11-11 / v27
\section{混合対角の局所正則化と全Artin底変更}
\label{rd:sec:regular}

第\ref{cd:sec:finite-diagonal}節の有理関数体上の計算を、
$R=\mathbb Q[h]_{(h)}$上の行列へ置き換える。
本節では高次ホモトピー、閉自己準同型、コホモロジーの基底および積を
すべて$h=0$で正則に選ぶ。最終的な三次式は
式\eqref{cd:eq:absolute-polynomial}と同じである。
これによって有限対角の低次数積を全Artin切断へ移せるが、
射影残余圏の導来自然変換との比較を証明したことにはならない。

\subsection{単元小行列を用いる局所線形方程式}
\begin{lemma}[一定階数の単元小行列解法]\label{rd:lem:unit-solve}
$R=k[h]_{(h)}$、$K=k(h)$とし、$M\in\operatorname{Mat}_{a\times b}(R)$の
$K$上の階数と$M(0)$の階数がともに$r$であるとする。
$M(0)$の独立な$r$列$J$を選び、その$r$列の独立な$r$行$I$を選ぶ。
$B\in\operatorname{Mat}_{a\times s}(R)$について$MX=B$が$K$上可解なら、
\[
 X_J=M_{I,J}^{-1}B_I,\qquad X_{J^c}=0
 \tag{RD1}\label{rd:eq:unit-solve}
\]
は$R$上の解である。この構成は$R$の任意の商への底変更と両立する。
\end{lemma}
\begin{proof}
$\det M_{I,J}$は$R$の単元であるから表示した$X$の係数は$R$に属する。
$K$上で$M_J$は$\operatorname{im}M$の基底であり、行$I$への射影は
この像上で単射である。$MX-B$はこの像に属し、行$I$で零なので零である。
分母は単元のみであり、逆行列と等式を商環へ移せる。
\end{proof}

一定階数を仮定しない任意の系では、この小行列だけでは全方程式を解けない。
実装は常に$MX=B$を全行について再検査し、不整合なら停止する。
一般の核の飽和とは用途が異なり、既知の完全な分解の像へ持ち上げる系では
この単元小行列法を優先できる。

\begin{proposition}[有界分解の正則な曲率化]\label{rd:prop:regular-existence}
第\ref{cd:thm:finite-curvature}の入力を$R$上に持ち、各内部次数において
増大射を含む分解が有限自由$R$-加群の分裂完全列であると仮定する。
このとき高次ホモトピー$s_j$を$R$上で選べる。
構成に現れる$d$による持上げを式\eqref{rd:eq:unit-solve}で実行できる。
小複体の正則な閉元から閉自己準同型への持上げについても、
増大射を固定した逐次比較の各段階に同じ解法を使える。
\end{proposition}
\begin{proof}
分裂完全性により、各次数の$d$の像と核は有限自由な直和因子であり、
$d$の階数は特殊点と一般点で等しい。
$f=\Omega_L-\Omega_R$の対角上の作用は零なので、初段の持上げが存在する。
以後の障害が閉じていること、負のExtの消滅、有限長による停止は
定理\ref{cd:thm:finite-curvature}の証明と同じである。
各段階の右辺は既に$R$上にあり、所要の像に入るので
補題\ref{rd:lem:unit-solve}を使える。
閉元の比較では$t$の指数と入力鎖次数の順に解く。
増大射により固定される最下段を除き、右辺は比較写像の恒等式から$d$の閉元となる。
同じ持上げを行い、最後の鎖次数で残る式も全行検査する。
\end{proof}

本研究の混合入力では、次数別自由性と分解の完全性に既存の全Artin PBW結果を用いる。
その仮定を有限個のHilbert関数の計算だけで置き換えない。
$R$はDVRなので、増大射の値域を含めて各次数が自由なら、
完全列を右端から順に分裂させて上の分裂完全性を得る。

\subsection{正則な四変数模型の実計算}
関係式と中心元は第\ref{cd:prop:absolute-formula}と同じ$\lambda=1$の入力とする。
$P$の階数は$(1,4,6,4,1)$、$\deg(t)=(2,-3)$である。
\begin{theorem}[局所正則な低次数積]\label{rd:thm:local-product}
この入力について、曲率付き対角$D$と四つの次数一閉自己準同型を
$R$上で構成できる。次数一の元の増大射の$P_0$成分を$u,x,y,z$に規格化すると、
\[
 H^0(C_1)\simeq R^4,\qquad H^0(C_2)\simeq R^6
\]
であり、後者には
\[
 ux,\ uy,\ uz,\ xy,\ xz,\ yz
 \tag{RD2}\label{rd:eq:class-basis}
\]
という実際の合成の類を基底に取れる。対称積は全射で、$D_h=h^6+162$、
$L_h=1+h^6/54$と置くと、四つの平方の積は
\begin{align}
 u^2&=\frac{9h^2}{D_h}\,uz-\frac{54}{D_h}\,xy,\notag\\
 x^2&=-L_huy-\frac{h^2}{3}yz,\qquad
 y^2=-L_hux-\frac{h^2}{3}xz,\label{rd:eq:product}\\
 z^2&=\frac{4h^{10}+675h^4}{36D_h}\,uz
       -\frac{2h^8+351h^2}{6D_h}\,xy.\notag
\end{align}
すべての係数は$R$に属する。
この対称積の核は第\ref{cd:prop:absolute-formula}の$q_u,q_x,q_y,q_z$であり、
$G_h$の勾配空間に等しい。
\end{theorem}
\begin{proof}[有限計算と局所自由性の証明]
\texttt{regular\_diagonal.py}は二次非可換関係から作った絶対Koszul分解に
単元小行列解法を適用する。
$d s_1+s_1d=f$と$d s_2+s_2d=-s_1^2$、およびすべての高次の残余式を
$\mathbb Q(h)$上の恒等式として検査する。
各解の分母は$h=0$で非零であることも検査する。
次数一の閉元を閉行列に持ち上げた後、全16個の順序付き合成を計算する。
これは係数の積でなく行列の合成である。

小複体の隣接項の次元と微分の階数は次のとおりである。
\[
\begin{array}{c|ccc|cc}
 r&\dim C_r^{-1}&\dim C_r^0&\dim C_r^1&
 \operatorname{rank}d^{-1}&\operatorname{rank}d^0\\\hline
 1&0&10&56&0&6\\
 2&4&35&120&4&25
\end{array}
\]
後二列は$\mathbb Q(h)$上でも$h=0$でも同じである。
特殊点での独立な小行列は$R$上の単元小行列となり、
一般点でより大きな階数を持たないので、各微分の像と核は直和因子となる。
従ってコホモロジーは自由で、その階数は四と六である。

$B=d^{-1}:C_2^{-1}\to C_2^0$の四列と、式\eqref{rd:eq:class-basis}の
六つの合成を$C_2^0$へ送った列を並べて、$35\times10$行列$T$を作る。
所定の基底順序における最初の10行の小行列の行列式は、$h=0$で$-3$である。
$\ker d^0$の階数は10なので、$T$はその$R$上の基底となる。
全16個の合成をこの基底に展開し、最初の四つの境界方向を除く。
反対称積が零となり、対称積の結果は式\eqref{rd:eq:product}である。

その核を求めると表示した四つの$q$を得る。
これらの平方単項式に関する小行列式は$3+h^6/54$で、$R$の単元である。
さらに$\nabla G_h=A(q_u,q_x,q_y,q_z)^{\mathsf t}$における
\[
 A=\begin{pmatrix}
 L_h&0&0&-h^4/18\\
 0&1&0&0\\0&0&1&0\\h^2/3&0&0&1
 \end{pmatrix},\qquad \det A=1+h^6/27
\]
も$R$上可逆である。最後の積分検査にのみ$G_h$を用い、
非可換行列の構成や積の計算には用いない。
\end{proof}

\begin{corollary}[全Artin切断と完備化]\label{rd:cor:basechange}
$r=1,2$と任意の$m\ge0$について、自然な写像
\[
 H^0(C_r)\otimes_R R/(h^{m+1})
 \longrightarrow H^0(C_r\otimes_R R/(h^{m+1}))
\]
は同型であり、合成と両立する。同じことが$\mathbb Q[[h]]$への底変更についても成り立つ。
特殊ファイバーの積は$G_0$のJacobian環の次数一・二の積である。
\end{corollary}
\begin{proof}
証明中の単元小行列による分解は、任意の$R$-代数への底変更後も分裂する。
隣接する微分の像・核とそれらの商の分解から、主張するコホモロジーの底変更を得る。
積の代表を与える閉行列はすべて$R$上で定義されているので、
合成も底変更と可換する。$h=0$では
$u^2=-xy/3$、$x^2=-uy$、$y^2=-ux$、$z^2=0$となる。
\end{proof}

この系の対象は\emph{有限曲率付き対角の自己準同型代数}である。
全Artin切断でこの代数が定義できることと、残余圏のrotation積との比較は別である。
独立な$v^3$の因子は第\ref{cd:sec:finite-diagonal}節の一変数計算により加えられる。

保存した分解、高次ホモトピー、閉行列および積表の全455種類の有理係数を調べると、
その分母の最小公倍数（非零有理数倍を除く）は
\[
 (h^3-1)(h^6+162)
\]
である。この事実も証明書の有限係数から厳密に再計算する。
これは選んだ行列表示の分母の記述であって、非可換射影圏のsmoothnessの判別式ではない。


\subsection{有限arity比較で実際に使えること}
普通のHochschildコチェインの積全体を用いる議論と、arityについて直和を取る議論を区別する。
後者と曲率付き双加群の導来圏の比較には条件が必要である
\cite[\S3.3, \S3.5]{RDPP}\cite[Theorem~6.14, Remark~6.16]{CDGHLR}。
一方、次の有限性は直接証明できる。

\begin{lemma}[有界arity同型による共役]\label{rd:lem:bounded-conjugation}
$c,c'$を有界arityの曲率付きbar coderivationとし、$c^2=(c')^2=0$とする。
$F$を$F_0=0$、$Fc=c'F$を満たすbar coalgebra同型とし、
$F$と$G=F^{-1}$のTaylor成分のarity上界をそれぞれ$L_F,L_G$とする。
arity上界$M$のcoderivation$X$に対し、$FXG$はarity上界
\[
 L_G(L_F+M-1)
 \tag{RD3}\label{rd:eq:arity-bound}
\]
を持つ。従って共役は、arityについて有限支持なHochschildコチェインの
微分付きLie代数の同型を与える。内部次数と周期変数を保つ場合は、この同型もそれらを保つ。
\end{lemma}
\begin{proof}
入力長$n$の語に$G$を作用させて長さ$k$の語が現れるためには、
各ブロックの長さが高々$L_G$なので$n\le L_Gk$である。
$X$のarity $m\le M$の成分は長さを$k-m+1$に変える。
その後$F$を作用させて出力長一へ射影した成分が非零なら
$k-m+1\le L_F$、従って$k\le L_F+M-1$である。
これが式\eqref{rd:eq:arity-bound}を与える。$m=0$の場合も同じ数え方が成り立つ。
共役はcoderivationと交換子を保ち、
$[c',FXG]=F[c,X]G$である。逆方向にも同じ議論を適用する。
\end{proof}

\begin{remark}[Lie比較から積比較を飛躍させない]
補題\ref{rd:lem:bounded-conjugation}だけではcup積の保存を証明していない。
cup積は単なるcoderivationの交換子とは異なる。
次に必要なのは、一つの入力を標識したbar双加群の比較を構成し、
有限対角のYoneda積とcompactly supported Hochschildのcup積との一致を示すことである。
さらに、それを残余圏の導来自然変換へ移す際には内部shiftと非恒等sectorを扱う必要がある。
この二つの比較は本節では未完成とする。
\end{remark}

\subsection{検算の範囲と再現}
新しい厳密計算は
\begin{verbatim}
python3 computations/verify_regular_diagonal.py \
  --output build/regular_diagonal_certificate.json \
  --full-output build/regular_diagonal_full.json
\end{verbatim}
で再現する。通常の証明書は積表、階数、単元小行列、検査範囲と
全行列データのSHA256を持つ。全行列も上の指定によって再生成できる。
容易に再生成できる大きな出力はGitで二重保守しない。

本節で実行した有限恒等式の検算、先行する全Artin PBW証明への依存、
本文の局所線形代数の証明、未完成の圏論的比較を区別する。
一般$\lambda$の絶対公式、指定形式模型との全次数同値、
genericなrotation付き圏同値および相対smoothnessは本節の結論ではない。

% 2026-11-11 / v30
\section{可換曲率付き対角の有限ホモトピーとkernel作用}\label{dki:sec:identity}

本節は、論文Bの相対Artin節における条件(KF)のうち、
\emph{stabilized diagonalが有限rank graded matrix factorization上で恒等関手として作用すること}を、
任意の可換基底環上で有限のホモトピーにより証明する。
この結果は、kernel間のHom複体とdg双加群の導来自然変換複体が
擬同型であるという\emph{全忠実性}とは別である。
後者が必要な(BC)および非可換残余圏のrotation積は未完成として区別する。

以下の曲率付き摂動の原理はTu\cite[Appendix~A, Lemma~A.1]{DKITu}、
体上の積分変換・kernel比較はBallard--Favero--Katzarkov
\cite[Theorem~5.22]{DKIBFK}とDyckerhoff\cite[Section~5]{DKIDyckerhoff}
を参照する。特にBFKの定理は体上のdg圏の内部Homについての結果であり、
非正則Artin環へ仮定なしで移すものではない。
ここで直接示すのは\emph{objectwiseな恒等作用}とその底変更である。

\subsection{差分商による有限対角}

$R$を任意の可換環、$S_x=R[x_1,\ldots,x_n]$、
$S_y=R[y_1,\ldots,y_n]$、$S_e=R[x,y]$、
$W\in R[z_1,\ldots,z_n]$を斉次三次式とする。
$s_i=y_i-x_i$と置く。
$W_i=W(x_1,\ldots,x_i,y_{i+1},\ldots,y_n)$、$W_0=W(y)$、
$W_n=W(x)$と定め、
\begin{equation}
 a_i=\frac{W_i-W_{i-1}}{y_i-x_i}\in S_e
 \quad\Longrightarrow\quad
 \sum_{i=1}^n s_i a_i=W(x)-W(y).
 \tag{DKI1}\label{dki:eq:diff}
\end{equation}
右辺の割算は多項式の差分商であり、$s_i$を局所化しない。
$E=\bigoplus_iRe_i$、$K=S_e\otimes_R\bigwedge E$とし、
$\deg_{\rm coh}(e_i)=-1$、$\deg_{\rm int}(e_i)=1$、
$\deg t=(2,-3)$とする。
\begin{equation}
 d_0=\sum_i s_i\iota_{e_i^*},\qquad
 \delta=t\sum_i a_i(e_i\wedge-),\qquad
 d_K=d_0+\delta.
 \tag{DKI2}\label{dki:eq:koszul}
\end{equation}
すると$d_K^2=t(W(x)-W(y))$であり、$d_K$はcohomological degree一・
内部次数零を持つ。$K$をstabilized diagonalの明示的有限模型とする。

$M$を$S_y[t,t^{-1}]$上の有限自由graded MF、$d_M^2=tW(y)$とする。
$S_x$への射影に沿った、核の外積との畳込みの\emph{大きな}模型は
\begin{equation}
 \mathscr C(M)=K\otimes_{S_y}M,
 \qquad D_M=d_K\otimes1+(-1)^{|e|}\otimes d_M,
 \qquad D_M^2=tW(x).
 \tag{DKI3}\label{dki:eq:convolution}
\end{equation}
ここで$(-1)^{|e|}$は外積次数のparityであり、$S_x$-加群としては
$y$の多項式変数を残している。したがってこれを有限rankの
$S_x$-MFと\emph{そのまま}同一視しない。

\subsection{対角評価の明示的な強ホモトピー同値}

$M(x)$を係数置換$y_i\mapsto x_i$で得たMFとする。
$\pi_M:\mathscr C(M)\to M(x)$を、外積次数零への射影と
$y=x$での評価を合成して定める。

\begin{theorem}[有限対角のobjectwise恒等作用]\label{dki:thm:identity}
任意の可換基底環$R$、斉次三次式$W$、有限自由MF $M$に対し、
$\pi_M$はgraded MFの閉じた射であり、
\emph{明示的な有限和}で与えられる逆とホモトピーを持つ。
具体的に、多項式Koszul分解$(K,d_0)$の標準的な強変形収縮
\[
 (S_x\otimes_{S_y}M,0)\;
 \xleftarrow[\pi_0]{\ \ }(K\otimes M,d_0)
 \xrightarrow[\ ]{i_0}\;(S_x\otimes_{S_y}M,0)
\]
を$h_0$で表し、
\[
 p_M=\delta\otimes1+(-1)^{|e|}\otimes d_M(y),\qquad D_M=d_0+p_M
\]
とする（左端は$M$の自由基底を用いる同次加群としての表示）。
このとき
\begin{equation}\label{dki:eq:finite-hpl}
 \begin{split}
 i_M&=\sum_{j=0}^{n}(-h_0p_M)^j i_0,\\
 H_M&=\sum_{j=0}^{n}(-h_0p_M)^j h_0
 \end{split}
 \qquad\text{および}\qquad \pi_M=\pi_0
 \tag{DKI4}
\end{equation}
が定義され、次の恒等式を満たす：
\begin{equation}\label{dki:eq:identities}
 \begin{gathered}
 D_Mi_M=i_Md_M(x),\qquad
 \pi_MD_M=d_M(x)\pi_M,\qquad \pi_Mi_M=1,\\
 D_MH_M+H_MD_M=1-i_M\pi_M.
 \end{gathered}
 \tag{DKI5}
\end{equation}
以上の構成は任意の$R\to R'$への底変更と両立する。
$\mathbb G_m$の斉次次数と内部shiftを保ち、
三次斉次模型の$\mu_3$スカラー作用とも両立する。
\end{theorem}
\begin{proof}
一変数のKoszul分解$R[x,y]e_i\xrightarrow{s_i}R[x,y]$には
\[
 h_i(f)=e_i\frac{f(y_i)-f(x_i)}{y_i-x_i},\qquad
 h_i(e_if)=0
\]
という$R[x_i]$-線形な収縮がある。
その射影は$y_i=x_i$での評価、包含は$x_i$の多項式を
$y_i$に依らない式と見る写像である。
この収縮のtensor trickにより、$n$変数の$(K,d_0)$について
\[
 d_0h_0+h_0d_0=1-i_0\pi_0,\quad
 h_0^2=\pi_0h_0=h_0i_0=0
\]
を満たす収縮が得られる。例えば$n=2$では、外積次数零の
$f(y_1,y_2)$を
\[
 h_0(f)=
 e_1\frac{f(y_1,y_2)-f(x_1,y_2)}{y_1-x_1}
 +e_2\frac{f(x_1,y_2)-f(x_1,x_2)}{y_2-x_2}
\]
へ送り、$e_2f$を
$e_1\wedge e_2(f(y_1,y_2)-f(x_1,y_2))/(y_1-x_1)$へ送る。
この表示には$h$-進分母も基底環の非単元の反転もない。

$p_M$は外積次数を減少させない（$d_M$項は次数を保持し、
$\delta$項は一つ上げる）。$h_0$は外積次数を必ず一つ上げるので
\[
 (h_0p_M)^{n+1}=(p_Mh_0)^{n+1}=0.
\]
曲率は$D_M^2=tW(x)$であり、これはすべての演算子と可換する。
共通の中心曲率を持つ形式への有限homological perturbation
\cite[Appendix~A, Lemma~A.1]{DKITu}を適用すると、
式\eqref{dki:eq:finite-hpl}と\eqref{dki:eq:identities}が従う。
移送した微分は
\[
 \pi_0p_M(1+h_0p_M)^{-1}i_0=d_M(x)
\]
である。実際、増大射$\pi_0$は正の外積次数を消し、
$\delta$は外積次数を上げるため、唯一残るのは
$d_M(y)$を外積次数零に作用させ、$y=x$に評価した項である。
$\pi_0p_Mh_0=0$なので摂動後の射影は$\pi_0$に一致する。

すべての式は多項式の差分商と長さ高々$n+1$の有限和で書かれ、
$R\to R'$への底変更と可換する。内部次数は上の$e_i,t$の規約で
保存される。$\mu_3$のスカラー作用では、差分商、評価、外積シフトは
いずれも同次である。
\end{proof}

\begin{corollary}[自然な対角評価と有限Artinでの恒等性]
\label{dki:cor:unit}
$R_m=k[h]/(h^{m+1})$とすると、対角kernelの畳込みは
有限自由graded MFの各対象に、内部shiftと両立して
\emph{自然なホモトピー同値}を与える。
その特殊化は$R_0=k$上の同じ差分商対角評価である。
特に、(KF)に必要な\emph{対角の恒等作用とその底変更}は成立する。
\end{corollary}
\begin{proof}
次数零のMF射$f:M\to N$に対し、$1\otimes f(y)$は
畳込み側の閉射で、評価射について
\[
 \pi_N\circ(1\otimes f(y))=f(x)\circ\pi_M
\]
が厳密に成り立つ。高次のdg自然性は同じ式にKoszul符号を付けて得る。
各$\pi_M$は定理\ref{dki:thm:identity}により明示的ホモトピー逆を持つ。
特殊化は定理の底変更性による。
\end{proof}

\subsection{(KF)の残り：核作用の全忠実性}

\begin{remark}[恒等作用とkernel全忠実性の相違]\label{dki:rem:kf}
上の系は、相対stabilized diagonalが有限rank MF上で
\emph{恒等関手として作用すること}を示す。
これは\emph{相対kernel Homと導来自然変換Homの全忠実性}、すなわち
論文Bの相対Artin節で考えた写像
\[
 \operatorname{Hom}(\mathscr D_m,\mathscr D_m(r))
 \longrightarrow
 \RHom_{\mathcal T_m^e}(\Delta,\Delta(r))
 \tag{DKI6}\label{dki:eq:kf-map}
\]
が擬同型であることとは異なる。体上ではBFKのkernel定理が
対応する全忠実性を与えるが、Artin底では、同一enhancement上に
式\eqref{dki:eq:kf-map}の\emph{導来}模型を作り、二側barの値が
その模型に一致することを確認する必要がある。
定理\ref{dki:thm:identity}はそのためのunitの明示的な比較を提供するが、
(BC)や非可換rotation積の全同定は結論しない。
\end{remark}

\begin{example}[項別平坦性はK-flat性を意味しない]
\label{dki:ex:kflat}
$R=k[\varepsilon]/(\varepsilon^2)$とし、両方向に無限な複体
\[
 P:\quad\cdots\xrightarrow{\varepsilon}R
 \xrightarrow{\varepsilon}R\xrightarrow{\varepsilon}R\xrightarrow{\varepsilon}\cdots
\]
を考える。各項は自由で、$\ker(\varepsilon)=\varepsilon R$より
$P$は非輪状である。しかし$P\otimes_Rk$は零微分で
各次数のコホモロジーが$k$である。従って$P$はK-flatでない。
特に、(KF)の二側barの\emph{各項が平坦}であるという情報だけでは、
それが必要な導来tensorや導来双加群Homを計算することは保証されない。
一方、定理\ref{dki:thm:identity}は\emph{強ホモトピー同値}なので
この障害を回避し、任意の底変更で保存される。
\end{example}

\paragraph{再開すべき核心.}
指定生成子$K_m\oplus K_m(1)\oplus K_m(2)$の
適切な相対cofibrant双加群模型を作り、kernel作用の写像
\eqref{dki:eq:kf-map}を、体上BFKの全忠実性写像へ\emph{同じ鎖レベル}で
特殊化できることを確認する。単にgeneratorの特殊ファイバーでの
生成性や、対象毎の対角恒等作用だけから(BC)を導かない。

\paragraph{有限計算の範囲.}
\texttt{computations/verify\_diagonal\_identity.py}は、
一変数三次式と、混合二変数三次式
$W=y_1^3+y_2^3+2y_1^2y_2$を持つそれぞれrank2、rank4の
明示的MFについて、$t=1$で式\eqref{dki:eq:identities}、
$D_M^2=W(x)$、元のKoszul収縮を有理係数の多項式恒等式として検査する。
これは一般の$n$についての証明でも、(KF)の全忠実性の計算でもない。

% 2026-11-11 / v28
\section{有限対角のYoneda積と有限支持Hochschild積}
\label{mb:sec:comparison}

本節の目的は、有限対角の計算と射影残余圏の間にある比較を分解することである。
以下で証明するのは、同次分裂bar比較を持つ有限対角の合成積と、
有限支持Hochschild cochainのcup積の同定、およびMFへの乗法的characteristic写像の構成である。
この写像の全Artin切断における同型性には後述の条件(BC)を残す。
生成された有限行列の恒等式から、この条件やgeneric圏同値を結論しない。

先行研究との関係を明記する。cap作用によるYoneda積との比較は
Negron\cite[Theorem~5.3]{MBNegron}の標準的方法に基づく。
同論文では受け手はdg代数であるため、ここで必要な曲率を持つ受け手・
bar長の完備化・有限支持という組合せを直接検証する。
曲率摂動にはTu\cite[Appendix~A, Lemma~A.1]{MBTu}を用いる。
有限支持のtwisting挿入はLowen\cite[Propositions~3.11, 3.13]{MBLowen}の
brace写像であり、これらの原理自体の新規性は主張しない。
本研究の有限行列模型への同時適用の初出は未確認である。

\subsection{入力と二種類の完備化}
$R$を可換$\mathbb Q$-代数、$A=R\oplus A_+$を連結内部次数付き$R$-代数とする。
$A$はコホモロジー次数零にあり、各内部次数成分は有限自由とする。
$w\in Z(A)_3$を固定し、$A^e=A\otimes_R A^{\rm op}$と置く。
正規化bar分解を
\[
 B_p=A\otimes_R(sA_+)^{\otimes p}\otimes_R A,
 \qquad |s a|_{\rm coh}=-1
\]
と書く。内部次数は多項式次数である。
$b$は通常のbar微分であり、鎖長を一つ下げる。
$\deg(t)=(2,-3)$、$t$を可逆とし、bar長について各二重次数内で完備化する。
完成barを$\widehat B$と書く。一方、Hochschild cochainは\emph{arityについて直和}を取り、
\[
 C_c(A,w)=\Hom^{\rm cont}_{A^e[t,t^{-1}]}(\widehat B,A[t,t^{-1}])
\]
とする。右辺の値域はbar長に関して離散であり、各同次cochainは有限arityである。
この$C_c$を通常のarity積totalizationと置き換えない。

$B$の元を$a_0[a_1|\cdots|a_p]a_{p+1}$と書く。
$\delta$はbarのすべての隙間に$t w$を交代符号で挿入する作用である。
$\mathcal D_B=b+\delta$とすると
\begin{equation}\label{mb:eq:bar-curvature}
 \mathcal D_B^2=t(w_L-w_R),\qquad \varepsilon_B\mathcal D_B=0.
\end{equation}
実際、$b^2=\delta^2=0$であり、$b\delta+\delta b$では内部の項が相殺し、
左右端の$tw$だけが残る。従って$C_c$の微分は平方零である。

次の入力を明示的に仮定する。
\begin{enumerate}
\item $P_p=0$（$p<0$または$p>N$）で、$P\to A$は同次有限生成射影$A^e$-分解である。
\item $P$と$B$の間に、鎖長を保つ$A^e$-線形な同次special deformation retraction（収縮を$\eta$と書く）
\begin{equation}\label{mb:eq:sdr}
 pi=1,\quad b\eta+\eta b=1-ip,\quad \eta^2=\eta i=p\eta=0
\end{equation}
があり、$\eta$は鎖長を一つ上げる。増大射とも両立する。
\end{enumerate}
例えばKoszul分解がbar分解の次数毎に分裂する部分複体であり、商の各項が射影なら、
商は非負鎖次数の完全な射影複体である。最下次数から順に全射を分裂させることで商を縮約でき、
これから\eqref{mb:eq:sdr}を得る。
本研究の適用では、Koszul完全性・次数毎自由性・交叉空間の分裂を入力として用いる。
有限次数のHilbert関数だけを、この全次数の入力の代わりにしない。

\subsection{bar曲率の有限対角への移送}
\begin{lemma}[有限支持に適合する曲率摂動]\label{mb:lem:perturbation}
上の入力の下で
\begin{align*}
 I&=(1+\eta\delta)^{-1}i,&
 \Pi&=p(1+\delta\eta)^{-1},&
 H&=(1+\eta\delta)^{-1}\eta,\\
 D_P&=d_P+p\delta(1+\eta\delta)^{-1}i
\end{align*}
は同次な曲率付きspecial deformation retractionを与える。すなわち
\[
 D_P^2=t(w_L-w_R),\quad \Pi I=1,\quad
 1-I\Pi=\mathcal D_BH+H\mathcal D_B,
 \quad \varepsilon_P\Pi=\varepsilon_B.
\]
$I,H$は連続で、$\Pi$はbar長$N$より大きい入力で零である。
$D_P$は$d_P+\sum_{j\ge1}t^j s_j$という有限和で、$s_j$は鎖長を$2j-1$上げる。
\end{lemma}
\begin{proof}
$\eta\delta$と$\delta\eta$はbar長を二つ上げるため、幾何級数は完成bar上で意味を持つ。
曲率$t(w_L-w_R)$は$A^e$-線形写像と可換し、Tuの曲率摂動補題を
\eqref{mb:eq:sdr}の符号規約で適用できる。
同補題は線形写像の恒等式なので、この分裂した$R$-線形状況でも同じ証明が通る。
$D_P$の$t^j$成分は鎖長を$2j-1$上げるから、$P$の有界性により有限個しか残らない。
$\Pi$の各項は、bar長を増やしてから$p$を適用する。
従って入力長が$N$を超えると零であり、残る場合にも和は有限である。
増大射の等式は、$\varepsilon_Pp=\varepsilon_B$と
正の長さを持つ項を増大射が消すことから従う。
\end{proof}

$P$に既に別の曲率付き微分$D'$が構成されていても、同じ比較を使える。
\begin{lemma}[有限曲率化の選択独立性]\label{mb:lem:uniqueness}
$D'$と$D_P$がともに$d_P+\sum_{j\ge1}t^js_j$の形で同じ曲率を持つなら、
\[
 U=1+\sum_{j\ge1}t^j u_j,\qquad D'U=UD_P
\]
となる有限な可逆写像がある。$u_j$は鎖長を$2j$上げ、$\varepsilon_PU=\varepsilon_P$である。
\end{lemma}
\begin{proof}
$t$の次数順に方程式を解く。$j$段目の障害は
$\operatorname{End}_{A^e}(P)$のコホモロジー次数$1-2j<0$の閉元である。
閉じていることは$D'^2=D_P^2$と、それ以前の比較式から従う。
$P$が対角の射影分解なので、この負次数のコホモロジーは
$\Ext^{1-2j}_{A^e}(A,A)=0$である。従って$u_j$を選べる。
$2j>N$なら$u_j=0$であり、最後に残り得る障害も同じ負Ext消滅から零となる。
$U-1$は鎖長を厳密に上げるので、その逆も有限幾何級数である。
増大射との両立も鎖長から従う。
\end{proof}

\subsection{cap作用と乗法性のホモトピー}
$\Delta_B$をbarのdeconcatenationとする。
内部の積と曲率挿入を切断の左右へ分ければ、
$\Delta_B\mathcal D_B=(\mathcal D_B\otimes1+1\otimes\mathcal D_B)\Delta_B$となる。
テンソル積では中間の右曲率と左曲率が相殺する。
cochain $\phi$の右cap作用を
\[
 R_\phi=(1\otimes\phi)\Delta_B
\]
で定める。ここでは$(\psi\otimes\phi)(x\otimes y)=
(-1)^{|\phi||x|}\psi(x)\otimes\phi(y)$のKoszul符号を用いる。
例えば偶次数代数上のarity $p$のcochainは、長さ$n$のbarの最後の$p$入力を取り、
符号$(-1)^{p(n-p)}$を伴って値を右端へ掛ける。

\begin{lemma}[曲率付きcap恒等式]\label{mb:lem:cap}
有限支持cochainについて
\[
 [\mathcal D_B,R_\phi]=R_{d\phi},\qquad
 R_\phi R_\psi=R_{\phi\smile\psi},\qquad
 \varepsilon_BR_\phi=\phi
\]
である。$R_\phi$は連続で、bar長を高々$\operatorname{arity}(\phi)$だけ下げる。
\end{lemma}
\begin{proof}
微分の式は$\Delta_B$の微分との両立、積の式はdeconcatenationの余結合性による。
符号はテンソル積のKoszul符号で一致する。曲率挿入についても、切断の左側に挿入する項と
右側に挿入する項の直和であり、同じ証明が成り立つ。
最後の式は左のbar長が零となる切断だけを増大射が残すことによる。
有限支持なので、各出力長に寄与する入力長は有限区間に限られる。
\end{proof}

\begin{theorem}[有限対角の乗法的bar比較]\label{mb:thm:yoneda-cup}
上の入力の下で、内部次数を保つ環同型
\[
 H^*\operatorname{End}_{A^e[t,t^{-1}]}(P,D_P)
 \ \simeq\ H^*C_c(A,w)
\]
がある。左辺の積は行列合成、右辺はcup積である。
一つの向きのcochain写像は
\begin{equation}\label{mb:eq:q-map}
 Q(a)=\varepsilon_Pa\Pi
\end{equation}
で与えられる。これを非可換入力から作った別の有限曲率化に用いる場合は、
補題\ref{mb:lem:uniqueness}の$U$で共役する。
\end{theorem}
\begin{proof}
$K=\Hom_{A^e[t,t^{-1}]}((P,D_P),A[t,t^{-1}])$と置く。
補題\ref{mb:lem:perturbation}を双対化すると
\[
 \alpha:C_c\longrightarrow K,\quad\phi\longmapsto\phi I,
 \qquad
 \beta:K\longrightarrow C_c,\quad v\longmapsto v\Pi
\]
は互いに逆なホモトピー同値である。
$H$は長さを増やすため、$\phi H$は再び有限支持であり、この双対化に積totalizationを要しない。

増大射からの写像
$e:\operatorname{End}(P,D_P)\to K$、$a\mapsto\varepsilon_Pa$も擬同型である。
実際、$t$の指数で濾過する。固定コホモロジー次数では$P$の二つの鎖次数が
$0,\ldots,N$に限られるため、可能な$t$指数は有限区間に入る。
従ってこの濾過の比較は収束し、随伴次数では通常の
$\Hom(P,P)\to\Hom(P,A)$の擬同型へ帰着する。
後者は$P$が有界射影分解であることから従う。
よって$Q=\beta e$は擬同型である。

乗法性を別に確認する。
\[
 J(\phi)=\Pi R_\phi I
\]
と置くと、$eJ=\alpha$なので$J$も擬同型であり、$H(J)$は$H(Q)$の逆である。
$\phi,\psi$が次数零の閉cochainなら
\begin{align}
 J(\phi)J(\psi)-J(\phi\smile\psi)
 &=\Pi R_\phi(I\Pi-1)R_\psi I\notag\\
 &=-d_{\operatorname{End}(P)}\bigl(\Pi R_\phi H R_\psi I\bigr).
 \label{mb:eq:product-homotopy}
\end{align}
一般次数では右辺に$(-1)^{|\phi|}$を加える。
$I$は無限和でも、$R_\psi$の有界な長さ減少と$\Pi$の長さ上界から、
この式の$P$上の値は有限和である。
従って$H(J)$、その逆$H(Q)$は環同型である。
単位は$\Pi I=1$から保たれる。
\end{proof}

この定理は、適用する分裂比較の入力を明示した有限対角の理論である。
文献\cite[Theorem~6.14, Remark~6.16]{MBGHLR}の
二つの第二種理論を無条件に同一視することはしていない。
同論文の外積生成条件や対角の本質像条件は、本節の環同型からは従わない。

\section{MFへのcharacteristic写像と低次数判定}
\label{mb:sec:characteristic}

\subsection{有限支持brace挿入}
$\mathcal T$を$(A,-tw)$の有限rank graded MFのdg圏とする。
その対象のtwisting行列を$q_X=sQ_X$とし、suspension後の次数は零とする。
有限支持cochain$\phi$を有限行列に拡張し、
\begin{equation}\label{mb:eq:embrace}
 (\operatorname{embr}_Q\phi)_n(a_n,\ldots,a_1)=
 \sum_{j_0,\ldots,j_n\ge0}
 \overline\phi_{n+\sum j_i}
 (q_{X_n}^{\otimes j_n},a_n,\ldots,a_1,q_{X_0}^{\otimes j_0})
\end{equation}
と置く。中間の$q_{X_i}$も各$a_{i+1},a_i$間へ挿入する。
内部次数$r$のcochainは$\Id$から$(r)$へのcoherent自然変換を与える。

\begin{proposition}[乗法的characteristic写像]\label{mb:prop:characteristic}
式\eqref{mb:eq:embrace}は有限和であり、微分とbraceを保つ。
従って有限対角について
\[
 \begin{gathered}
 \Theta_r:H^0\operatorname{End}(P,D)_r\longrightarrow
 \Hom_{D(\mathcal T^e)}(\Delta_{\mathcal T},\Delta_{\mathcal T}(r)),\\
 [a] \longmapsto [\operatorname{embr}_Q(\varepsilon_Pa\Pi)].
 \end{gathered}
\]
という、積と内部シフトに両立する写像を得る。
この写像の存在だけから同型性は従わない。
\end{proposition}
\begin{proof}
$\phi$のarity上界を$M$とすると、$n+\sum j_i\le M$の項しか現れない。
従って$Q_X$の冪零性を必要としない。
brace合成に現れる平面的な入力配置を、先に$q$を挿入するか、先にcochainを挿入するかで
数え直すと、各項が一対一に対応する。$q$のsuspended次数は零なので、
残りの入力のKoszul符号も一致する。これがLowenのbrace恒等式の有限支持版である。
構造cochain$m=m_0+m_2$はtwisting後にMF圏の構造へ移り、
対象方程式により曲率成分は零となる。従って微分とcup積を保つ。
通常のcategory barからのcochainは、増大射を導来圏で逆にすることで
表示した導来自然変換へ写る。ここではnaive cochain複体全体が
導来Hochschild複体を計算すると仮定する必要はない。
最後に定理\ref{mb:thm:yoneda-cup}を合成する。
\end{proof}

\begin{corollary}[低rotation次数に必要なarity]\label{mb:cor:arity}
$P_p=A^e\otimes_R E_p$が線形で、$E_p$の内部次数が$p$なら、
コホモロジー次数零・内部次数$r$の$\varepsilon_Pa$は
偶数$p\le\min(N,2r)$の成分しか持たない。
従って$\Theta_1$と$\Theta_2$には、それぞれarity高々二、高々四の
characteristic cochainを選べる。
\end{corollary}
\begin{proof}
$P_p$から$A[t,t^{-1}]$への次数零写像では$t$指数が$-p/2$であり、
値の多項式次数は$r-p/2$となる。これは非負でなければならない。
$\Pi$はbar長を増やしてから有限対角へ射影するので、cochainの入力arityも同じ上界を持つ。
\end{proof}

\subsection{全Artin同型性に残る条件}
本研究で計算した五変数模型について、$R=\CC[h]_{(h)}$、
$R_m=R/(h^{m+1})$とする。有限対角側の次数一・二は自由rank五・十で、
積と全Artin底変更は既存の正則行列計算で確認されている。
$\mathcal T_m=\MF^{\ZZ}(A_m,\Omega_m)$と置き、
\[
 T_{r,m}=\Hom_{D(\mathcal T_m^e)}(\Delta_m,\Delta_m(r))
 \quad(r=1,2)
\]
と書く。次の条件を独立に掲げる。

\noindent\textbf{(BC)} $T_{1,m},T_{2,m}$がそれぞれ自由rank五・十であり、
その自然な剰余写像が$T_{r,m}\otimes_{R_m}\CC\simeq T_{r,0}$を与え、
この同定と$\Theta_r$が両立する。

\begin{theorem}[条件付きArtin同定]\label{mb:thm:artin-criterion}
各$m$で(BC)が成り立つなら、$\Theta_1,\Theta_2$は全$m$で同型であり、
有限対角で計算した対称積は実際のMF内部shift積と同定される。
さらに既存のArtin残余比較を入力すると、同じ結論が線束生成残余pairのrotation積に成り立つ。
\end{theorem}
\begin{proof}
$h=0$では多項式$x_i$のcharacteristic写像は行列因子分解上の$x_i$倍である。
可換cubicの次数一・二の導来自然変換はJacobian環の同じ次数成分であり、
この多項式倍が同定を与える\cite[Theorem~3.5, Remark~3.6]{MBLinZhang}。
五変数では非恒等sectorは奇数コホモロジー次数を持ち、この次数零比較に寄与しない。
有限対角側の特殊積も同じJacobian積なので、$\Theta_r\bmod h$は同型である。
(BC)とsourceの自由性から$\Theta_r$は同じrankの自由$R_m$-加群間の写像で、
行列式の剰余が非零である。行列式は単元なので同型である。
乗法性は命題\ref{mb:prop:characteristic}、rotationへの移送は入力した残余比較による。
\end{proof}

\begin{remark}[未確認の適用条件]\label{mb:rem:remaining}
本稿は(BC)を証明していない。
既存の形式MF同値によってtargetを可換形式cubicへ移した後、
相対のsector/kernel模型とその底変更を確認することが次の具体的課題である。
体上の次元公式を、非正則Artin底へそのまま引用して済ませてはならない。
特に、導来Hom複体の底変更とその$H^0$の底変更は異なる。
sourceの正則性やrank計算だけではtarget側の条件は成立しない。
形式MF同値・Artin残余比較そのものの全証明の独立再監査も、本稿の作業では行っていない。
\end{remark}

(BC)がすべての$m$で両立して確認されれば、低次数積の核から
既存の形式模型$W_h$と計算済み多項式$G_h+v^3$の勾配空間を同定できる。
基底の同定を$h=0$で恒等に規格化し、形式的symmetrizerの三乗根を取れば、
両三次式は恒等に近い形式線形座標変換で同型になる。
これは有限代数的積分の定理の適用である。
しかし、その帰結も現段階では(BC)に条件付きであり、非零複素パラメータの圏同値は別に残る。

\section{計算済み混合候補の代数的正規形}
\label{mb:sec:normal-form}
ここでは圏論的な比較を使わず、有限対角から得た具体的多項式だけを扱う。
$q=h^6$、
\[
 A(q)=1+\frac{29q}{648}+\frac{q^2}{4374}
\]
として
\begin{align*}
 G_h+v^3={}&\frac{x^3+y^3+z^3}{3}+v^3+
 \frac{h^2}{3}xyz+(1+q/54)uxy\\
 &-\frac{h^4}{18}uz^2-\frac{h^2}{12}u^2z+A(q)u^3
\end{align*}
を固定する。これは$\lambda=1$に対する候補である。

\begin{proposition}[四本の代数方程式による正規化]\label{mb:prop:algebraic-normal-form}
$q=0$の\'etale近傍上に定義された代数的関数$a,b,c,d$で
\[
 (a,b,c,d)(0)=\left(1,\frac1{36},\frac{25}{432},1\right)
\]
を満たすものがあり、
\[
 u=a(q)U+h^2b(q)Z,\qquad z=h^4c(q)U+d(q)Z
\]
は上の候補を
\[
 \frac{x^3+y^3+Z^3}{3}+U^3+v^3+
 L(q)Uxy+h^2C(q)Zxy
\]
へ変換する。$a,b,c,d$は次の四方程式の指定局所解である：
\begin{align}
 A a^3-qa^2c/12-q^2ac^2/18+q^2c^3/3&=1,\notag\\
 3Aa^2b-(a^2d+2qabc)/12-q(2acd+qb c^2)/18+qc^2d&=0,\notag\\
 3Aab^2-(2abd+qb^2c)/12-(ad^2+2qbc d)/18+cd^2&=0,\label{mb:eq:algebraic-system}\\
 Aq b^3-qb^2d/12-qbd^2/18+d^3/3&=1/3.\notag
\end{align}
ここで
\[
 L=(1+q/54)a+qc/3,\qquad C=(1+q/54)b+d/3.
\]
最初の係数は
\begin{equation}\label{mb:eq:first-coefficients}
 \boxed{\quad L(q)=1+\frac{127}{5184}q+O(q^2),\qquad
 C(q)=\frac{13}{36}+\frac{1613}{1679616}q+O(q^2).\quad}
\end{equation}
\end{proposition}
\begin{proof}
代入後の$U^3,U^2Z,UZ^2,Z^3$の係数を比較すると
\eqref{mb:eq:algebraic-system}となる。二つの混合係数はそれぞれ$h^2,h^4$を
因子に持ち、それらを多項式として除いているので、$h=0$でも方程式が定義される。
指定した初期値は$q=0$での解であり、四方程式の$(a,b,c,d)$に関する
Jacobian行列の行列式はこの点で$9$である。
従ってその局所解空間から$q$直線への射影は相対次元零のsmooth射、すなわち\'etale射であり、指定形式解は一意で代数的である。
この標準のJacobian判定は\cite[Tag 00T7]{MBStacks}を参照する。
線形変換の行列式$ad-qbc$は初期値一なので可逆である。

四方程式を$q$で微分し、初期値で解くと
\[
 a'(0)=-\frac{23}{1728},\quad
 b'(0)=-\frac{139}{1679616},\quad
 c'(0)=-\frac{2323}{3359232},\quad
 d'(0)=\frac{37}{23328}.
\]
$L,C$の定義に代入すれば\eqref{mb:eq:first-coefficients}を得る。
すべて有理数上の線形代数であり、二次以上の項の推測や数値補間は用いていない。
\end{proof}

\begin{remark}[係数の意味]
これにより、\emph{計算済み候補}の$h^6Uxy$係数と$h^8Zxy$係数が決まった。
(BC)と形式比較が未確認なので、既存のformalityで選んだ$W_h$の係数が
無条件にこれらの値であるとは主張しない。
またこの命題は四本の有限代数方程式による明示化であり、
全域で一価な有理式を与えるものでもない。
\end{remark}

\subsection{再現と検証の範囲}
\texttt{computations/verify\_bar\_comparison.py}は
$M_2(\mathbb Q)[z]/(z^2)$、曲率$zI_2$に対し、有限語でbar曲率144件、
capと微分の両立288件、cap積360件を厳密演算で検査する。
$t=1$への偶周期特殊化による符号点検であり、無限barの証明の代用ではない。
さらに有限分裂複体で\eqref{mb:eq:product-homotopy}を30件点検し、
うち25件では乗法性の差が実際に非零で、そのホモトピー原像が一致する。
四本の正規化方程式・Jacobian行列式・一次係数も有理数上で再計算する。
このscriptは実際のcategorical Nat、generic圏同値、v27の全行列を検証していない。

% 2026-11-11 / v31
\section{Artin基底上のJacobian sectorと導来自然変換の底変更}
\label{ras:sec:artin-sector}

本節では標数零の体$k$、$R_m=k[h]/(h^{m+1})$、
$S_m=R_m[x_1,\ldots,x_5]$、次数三の
$W_m\in(S_m)_3$を考える。$F=W_m\bmod h$は
$\PP^4_k$の滑らかなcubic threefoldを定めるとする。
目的は、相対的な五変数のJacobian sectorと、graded MF圏の
$\Id\to(1),(2)$の導来自然変換を区別しながら底変更を調べることである。
以下の$\RHom$は常に\emph{$R_m$-線形}な導来Homである。
$R_m$の絶対的なHochschild理論を取り込まない。

体上のkernelの計算にはBallard--Favero--Katzarkov\cite{RASBFK}、
Lin--Zhang\cite[Theorem~3.5, Remark~3.6]{RASLZ}を用いる。
残余体のstabilizationが生成元となる体上の結果は
Dyckerhoff\cite{RASDyck}を用いる。
以下では、非正則Artin底へ適用するときに必要な
平坦複体の持上げと、可換対角の有限模型を明示する。
これら既知の結果そのものの新規性は主張しない。

\subsection{勾配Koszul複体の相対完全性}

\begin{lemma}[勾配の全Artin完全性]\label{ras:lem:gradient}
$K_m=K_\bullet(\partial_1W_m,\ldots,\partial_5W_m;S_m)$を
通常のKoszul鎖複体とする。このとき
\[
 H_i(K_m)=0\quad(i>0),\qquad
 H_0(K_m)=J_m:=S_m/(\partial_1W_m,\ldots,\partial_5W_m).
\]
$J_m$は内部次数別に有限自由な$R_m$-加群で、Hilbert級数は
\[
 \sum_{j\ge0}\operatorname{rank}_{R_m}(J_m)_jt^j=(1+t)^5.
 \tag{RAS1}\label{ras:eq:hilbert}
\]
特に$(J_m)_1\simeq R_m^5$、$(J_m)_2\simeq R_m^{10}$である。
全ての商$R_m\to R_a$（$a\le m$）への底変更は
$J_m\otimes_{R_m}R_a\simeq J_a$と両立する。
\end{lemma}
\begin{proof}
$F$の偏微分は原点以外に共通零点を持たない。
$k[x_1,\ldots,x_5]$はCohen--Macaulayなので、五つの二次式は
正則列であり、$K_0$は次数零以外で完全である。
従って$J_0$のHilbert級数は
$(1-t^2)^5/(1-t)^5=(1+t)^5$である。

$m\ge1$のとき、各成分は$R_m$-自由なので
\[
 0\longrightarrow K_0\xrightarrow{h^m}K_m
 \longrightarrow K_{m-1}\longrightarrow0
 \tag{RAS2}\label{ras:eq:short}
\]
は鎖複体の短完全列である。左端の微分は
$h^m\partial W_m=h^m\partial F$なので、左端は実際に$K_0$である。
長完全列による帰納法で$H_i(K_m)=0$（$i>0$）を得て、
\[
 0\longrightarrow J_0\xrightarrow{h^m}J_m
 \longrightarrow J_{m-1}\longrightarrow0
\]
も得る。各内部次数$j$の$J_m$の$k$-長さは
$(m+1)\dim_k(J_0)_j$である。
$(J_0)_j$の基底を同次数の多項式に持ち上げると、Nakayamaにより
$R_m^{d_j}\twoheadrightarrow(J_m)_j$を得る。
両辺の$k$-長さは$(m+1)d_j$なので、これは同型である。
同じ持上げ基底を使えば任意の$R_a$への底変更が従う。
\end{proof}

\begin{lemma}[有限nilpotent比較の検出]\label{ras:lem:nilpotent}
$f:C^\bullet\to D^\bullet$を、各項が$R_m$-平坦な複体の射とする。
$f\otimes_{R_m}k$が擬同型なら、$f$は擬同型である。
有界性は仮定しないが、\emph{各項の平坦性}は必要である。
\end{lemma}
\begin{proof}
$E=\operatorname{Cone}(f)$も項別平坦である。
有限フィルトレーション
$E\supset hE\supset\cdots\supset h^mE\supset0$の各商は、
平坦性により$E/hE$と複体として同型である。
$E/hE$は非輪状であるから、有限回の短完全列を用いて$E$も
非輪状となる。
\end{proof}

\begin{example}[特殊ファイバーの次元だけでは足りない]
\label{ras:ex:basechange-fails}
$R_1=k[h]/(h^2)$上の二項複体
$C=[R_1\xrightarrow{h}R_1]$を次数$0,1$に置く。
$H^0(C)=hR_1\simeq k$だが、
$H^0(C)\otimes_{R_1}k\to H^0(C\otimes_{R_1}k)=k$は\emph{零写像}である。
両コホモロジーの$k$-次元が一であっても底変更の同型は従わない。
従って、後述の標的の底変更を体上の次元公式だけで済ませてはならない。
\end{example}

\subsection{体上のkernel比較をArtin底へ持ち上げる際の構成}

$(\mathcal T_m,(1))$を、有限rank自由な次数付き
matrix factorizationとその冪等完備化から作った$R_m$-線形dg圏とする。
そのdg双加群による自然変換を
\[
 T_{r,m}:=H^0\RHom_{\mathcal T_m^e}
 (\Delta_{\mathcal T_m},\Delta_{\mathcal T_m}(r))
 \qquad(r=1,2)
 \tag{RAS3}\label{ras:eq:target}
\]
とする。

\begin{lemma}[導来bar評価の条件付き持上げ]\label{ras:lem:generator}
$\mathcal C_m$を項別$R_m$-平坦なHomを持つdg圏、
$G_m\in\mathcal C_m$とする。各対象$X,Y$について、
$G_m$を介した\emph{導来二側bar評価}の鎖模型があり、
その始域と終域が項別$R_m$-平坦で、$h=0$への剰余が
体上のMorita評価と厳密に一致すると仮定する。
体上で$G_0$がMorita生成元なら、この評価は擬同型である。
従って$G_m$は、この導来bar模型でのMorita生成元となる。
\end{lemma}
\begin{proof}
特殊ファイバーの評価はMorita生成性から擬同型である。
与えられた導来bar評価のconeは項別平坦で、
その剰余は特殊ファイバーで非輪状だから
補題\ref{ras:lem:nilpotent}によりcone自身も非輪状である。
\end{proof}
\begin{remark}[一般の特殊生成子持上げは偽]
指定した導来bar評価の存在・導来性・特殊化整合性は
単なるHom複体の項別平坦性からは従わない。
具体的な反例と隣接sectorの障害は
後続のArtin障害の節に記した。
\end{remark}

五変数の場合、$W_m=\sum_i x_i(\partial_iW_m/3)$から
残余点のstabilization $K_m$をKoszul型の有限rank因子分解として
作れる。$K_m,(1)K_m,(2)K_m$の直和を$G_m$とする。
体上ではこれがgraded MFのMorita生成元となる\cite{RASDyck,RASBFK}。
しかしこのことだけからArtin底の生成性は従わない。
補題\ref{ras:lem:generator}の導来bar評価・剰余整合性を別途確認する。
上のbarの各項は同次多項式の有限自由$R_m$-加群の直和である。
自然変換を計算するbar \emph{cochain}では直積が現れるが、
$R_m$はNoetherian（従ってcoherent）であり、自由加群の直積は平坦である。
$R_m\to k$への剰余は、各係数を成分毎に剰余する操作と可換する。

差$W_m(x)-W_m(y)$を
\[
 W_m(x)-W_m(y)=\sum_i(x_i-y_i)D_iW_m(x,y)
\]
と分解し、$\sum_i(x_i-y_i)\iota_i+\sum_iD_iW_m\,\wedge e_i$を
微分とするKoszul因子分解$\mathscr D_m$を取る。
これは$\mathbb G_m$-同変な\emph{stabilized diagonal kernel}である。
$x_i-y_i$は任意の$R_m$上で正則列なので、有限rank因子分解への畳込みにより
恒等関手を与えることは、通常のKoszul contractionで確認できる。
内部shiftを持つ$\mathscr D_m(r)$も同様である。

\begin{proposition}[相対kernel作用のnilpotent検出]\label{ras:prop:kernel-action}
$\mathscr D_m$によるkernel作用から得られる、相対的な自然な複体の射
\[
 \kappa_{r,m}:\Hom_{\mathrm{MF}^{\rm gr}(W_m(x)-W_m(y))}
 (\mathscr D_m,\mathscr D_m(r))
 \longrightarrow
 \RHom_{\mathcal T_m^e}(\Delta,\Delta(r))
 \tag{RAS5}\label{ras:eq:kappa}
\]
を、右辺で$G_m$の二側bar模型を使って表す。
この比較の特殊ファイバーが体上のkernel定理による擬同型なら、
$\kappa_{r,m}$は全ての$m$で擬同型である。
\end{proposition}
\begin{proof}
左辺は有限自由因子分解の同次Homなので項別$R_m$-自由である。
右辺の$G_m$-bar模型の各項は上記のように$R_m$-平坦で、
特殊化は体上の同じbar模型となる。kernelの作用は
$R_m$上の畳込み・合成で作るため、特殊化と可換である。
従って特殊化が擬同型なら、補題\ref{ras:lem:nilpotent}を適用できる。
\end{proof}

\begin{remark}[比較の適用範囲]
体上のkernel定理\cite{RASBFK}は、そのdg enhancementと
連続kernel functorを指定する。式\eqref{ras:eq:kappa}について
\emph{同じ}enhancementを用い、相対畳込みと$G_m$-bar評価を
整合させる必要がある。異なる第二種Hochschild totalizationを
同じ名前だからと同一視してはならない。
以下の計算は、有限自由のstabilized diagonalの相対模型について行う。
\end{remark}

\subsection{恒等sectorと非恒等sector}

\begin{proposition}[五変数の相対sector計算]\label{ras:prop:sectors}
標準の$\mu_3$-sector分解を相対stabilized diagonalに適用すると、
$\mathscr D_m$の内部shift付き自己準同型の次数零は
\[
 H^0\Hom(\mathscr D_m,\mathscr D_m(r))
 \simeq (J_m)_r\quad(r=1,2)
 \tag{RAS6}\label{ras:eq:sectors}
\]
となり、積は$J_m$の通常の多項式積である。
この同型は全$R_m\to R_a$と両立する。
\end{proposition}
\begin{proof}
$\mu_3$は標数零で線型簡約だから、sectorへの射影は
$R_m$上の直和分解である。恒等sectorでは、二つの対角の
Koszul homotopyを$y=x$で縮約すると、偏微分
$(\partial_1W_m,\ldots,\partial_5W_m)$のKoszul複体を得る。
補題\ref{ras:lem:gradient}により、次数零は$J_m$で、
正のKoszulホモロジーは零である。

非恒等$\zeta\in\mu_3$では、twisted対角との交点に
$((1-\zeta)x_1,\ldots,(1-\zeta)x_5)$のKoszul複体が現れる。
$1-\zeta$は単元なので、この五つは$S_m$上の正則列である。
その縮約後に残るnormal determinantは\emph{奇数}（階数五）の
cohomological parityを持つ。$(3)=[2]$によるシフトは偶数なので、
どの内部shift$r$でもcohomological degree零には寄与しない。
この計算は体上のsector formula\cite[Theorem~3.5]{RASLZ}の
明示的Koszul証明を$R_m$上で行ったものであり、
$1-\zeta$の可逆性以外に体の性質を使わない。
恒等sectorの積はdiagonalへの多項式倍であり、積の同定も従う
\cite[Remark~3.6]{RASLZ}。
\end{proof}

\begin{theorem}[相対kernel比較の下での(BC)]\label{ras:thm:bc}
各$R_m$について、\emph{指定生成子の導来二側bar評価が
補題\ref{ras:lem:generator}の条件を満たし、かつ相対kernel作用の射
\eqref{ras:eq:kappa}が体上のBFK kernel比較と同じenhancementにおいて
定義され、その$h=0$剰余が擬同型である}とする。
このとき、論文Bの条件(BC)がすべての$m$について成り立つ。
特に
\[
 T_{1,m}\simeq R_m^5,\quad T_{2,m}\simeq R_m^{10},\quad
 T_{r,m}\otimes_{R_m}k\xrightarrow{\sim}T_{r,0}\quad(r=1,2)
\]
であり、低次数の積は形式可換模型$W_m$のJacobian積である。
\end{theorem}
\begin{proof}
補題\ref{ras:lem:generator}の追加仮定により、導来bar模型で
実際の導来自然変換を計算できる。
命題\ref{ras:prop:kernel-action}のnilpotent持上げと
命題\ref{ras:prop:sectors}を合成すれば、標的の同定と底変更を得る。
内部shiftはdiagonalのequivariant characterのshiftであり、
比較の畳込みに組み込まれている。積の比較はkernelの合成であり、
恒等sectorでは多項式積となる。
\end{proof}

\begin{corollary}[混合族に関する条件付き正式]\label{ras:cor:formal-G}
既存の有限Artin MF同値・線束残余比較と、
上の相対kernel作用の仮定が与えられているとする。
$\lambda=1$で有限曲率付き対角から計算した多項式候補
$G_h+v^3$と、既存のformalityで選んだ可換形式三次式$W_h$は、
$\CC[[h]]$上、恒等に近い形式線形座標変換で同型である。
非零複素数$h$での圏同値はこの結論に含まれない。
\end{corollary}
\begin{proof}
定理\ref{ras:thm:bc}と既存の乗法的characteristic写像の
Artin判定から、全Artin次数で低次数積が一致する。
対応を逆極限に取り、次数二の核として勾配空間を同定する。
$h=0$では座標の同定を恒等に規格化できる。
有限代数的積分の定理とsymmetrizerの形式三乗根により、
両三次式を形式的な線形座標変換で同定できる。
\end{proof}

\subsection{残る検証の境界}

相対kernel作用の\emph{特殊ファイバー擬同型}は体上の既知定理だが、
非正則Artin上の有限自由enhancement、畳込み、二側barを
同じモデルで整合させた完全な鎖レベルの実装は、
この原稿では独立の有限行列証明書として提供していない。
従って上の定理は必要な比較モデルを指定した\emph{条件付き}の定理である。
前節の有限曲率付き対角と相対stabilized diagonalとの比較を
省略して、実際の非可換rotation積と無条件に同定してはならない。
最優先課題は、そのchain-level kernel作用の自然性・$R_m$-平坦性・
特殊化を一つの有限自由模型で具体的に検証することである。

具体的な五変数の古典特殊ファイバー
\[
 F_1=(x^3+y^3+z^3)/3+u^3+v^3+uxy
\]
では$F_x=x^2+uy$, $F_y=y^2+ux$, $F_u=3u^2+xy$から
$u\ne0$を仮定すると$(xy)^2=9u^4$かつ
$x^2y^2=-3u^4$となり矛盾する。従って唯一の勾配共通零点は原点。
厳密なGr\"obner証明書は
\texttt{computations/verify\_artin\_sectors.py}で再生成でき、
Hilbert関数は$(1,5,10,10,5,1)$である。
さらに計算済み候補$G_{1/2}+v^3$も同じ関数を持つことを検査する。
これは(BC)の鎖レベルの実装を検証したものではない。

% 2026-11-11 / v32
\section{低rotation次数のArtin障害と特殊生成子の持上げ}\label{aro:sec:obstructions}

五変数三次の体上の高次導来自然変換を計算し、
Artin底変更条件(BC)を次数一と二に分離する。
特殊生成子のMorita生成性を任意のArtin底へ持ち上げる旧補題は
以下の反例のため一般には成立しない。
既存の有限対角・形式MF復元結果は維持する。

\subsection{次数零に隣接するsector群}
\begin{proposition}[五変数三次の隣接群]\label{aro:prop:sector-neighbors}
$k$を標数零の代数閉体、$F\in k[x_1,\ldots,x_5]_3$を
滑らかな射影cubicを定める三次式、$\mathcal T_0$を
そのgraded MF dg圏、$\rho=(1)$とする。
\[
\begin{array}{c|ccc}
r&H^{-1}\RHom(\Delta,\Delta(r))&
H^0\RHom(\Delta,\Delta(r))&H^1\RHom(\Delta,\Delta(r))\\\hline
1&0&k^5&k^2\\2&0&k^{10}&0
\end{array}\tag{ARO1}\label{aro:eq:neighbors}
\]
$r=1$の$H^1$の二次元は二つの非恒等$\mu_3$-sectorである。
これは既知公式の特殊化であり、新しいHochschild分類とは主張しない。
\end{proposition}
\begin{proof}
Ballard--Favero--KatzarkovのTheorem~5.39
\cite{AROBFK}またはLin--Zhang\cite{AROLZ}を使う。
恒等sectorは偶数cohomological次数$m=2j$に
$\operatorname{Jac}(F)_{r+3j}$を与える。
勾配が五つの二次正則列なのでHilbert級数は$(1+t)^5$。
$m=0$では次元$5,10$、$m=\pm1$では寄与しない。
非恒等$g\in\mu_3$は原点のみを固定し、normal rankは五、
行列式重みは$-5$である。そのsectorは$m$が奇数で
\[
r+5+\frac{3(m-5)}2=0
\quad\Longleftrightarrow\quad2r+3m-5=0
\tag{ARO2}\label{aro:eq:weight}
\]
の場合に寄与する。$r=1,m=1$には二つの非恒等元が各一次元、
$r=2,m=-1,1$には寄与しない。
\end{proof}

\subsection{nilpotent底の正確な判定}
\begin{lemma}[隣接群消滅による自由な持上げ]\label{aro:lem:neighbor-vanishing}
$R_m=k[h]/(h^{m+1})$、$C_m$を項別$R_m$-平坦な
$\mathbb Z$-次数付き複体で、
$C_m\otimes_{R_m}R_a=C_a$が全ての$a\le m$で
鎖レベルで成立するとする。
$H^{-1}(C_0)=H^1(C_0)=0$、
$H^0(C_0)=k^d$なら
\[
H^{-1}(C_m)=H^1(C_m)=0,\quad
H^0(C_m)\simeq R_m^d,\quad
H^0(C_m)\otimes_{R_m}R_a\simeq H^0(C_a).
\tag{ARO3}\label{aro:eq:flat}
\]
有界性は不要である。
\end{lemma}
\begin{proof}
各項の平坦性から
$0\to C_0\xrightarrow{h^m}C_m\to C_{m-1}\to0$
は短完全列。長完全列で帰納すると両隣群が消え
$0\to H^0(C_0)\to H^0(C_m)\to H^0(C_{m-1})\to0$。
$H^0(C_m)$の$k$-長さは$(m+1)d$。
剰余基底の持上げからNakayamaで
$R_m^d\twoheadrightarrow H^0(C_m)$、長さ一致で同型。
\end{proof}

\begin{proposition}[次数一の最初の障害行列]\label{aro:prop:first-obstruction}
$C_1$を項別$R_1$-自由で$C_0=C_1/hC_1$とする。
$[x]\in H^0(C_0)$の閉代表を$\widetilde x$へ持ち上げ
$d\widetilde x=h\widetilde y$とすると
\[
\operatorname{ob}_1([x])
=[\widetilde y\bmod h]\in H^1(C_0)
\tag{ARO4}\label{aro:eq:obstruction}
\]
はよく定義され、
$\operatorname{im}(H^0(C_1)\to H^0(C_0))
=\ker\operatorname{ob}_1$。
従って表\eqref{aro:eq:neighbors}と両立する相対Nat模型を
構成できれば、$r=1$の一次障害は最大$2\times5$行列である。
$r=2$は補題\ref{aro:lem:neighbor-vanishing}の範囲になるが、
正しい標的複体の構成は未完成である。
\end{proposition}
\begin{proof}
$d^2\widetilde x=h\,d\widetilde y=0$から
項別自由性で$d\widetilde y\bmod h=0$。
持上げを$hz$変えると$\widetilde y\bmod h$は境界だけ変わる。
障害が零なら持上げを境界の原像で補正して閉じた元を得る。
\end{proof}
\begin{example}[非零障害]\label{aro:ex:obstruction}
$R_1=k[h]/h^2$上で$C^0=R_1^5,C^1=R_1^2$、
$d=h\begin{psmallmatrix}1&0&0&0&0\\0&1&0&0&0\end{psmallmatrix}$
とする。特殊ファイバーの$H^0,H^1$は$k^5,k^2$だが、
$H^0(C)=R_1^3\oplus k^2$、剰余写像の像は三次元。
体上のrankだけで(BC)は従わない。
\end{example}

\subsection{特殊生成子の持上げに関する訂正}
\begin{proposition}[項別平坦性だけでは生成性は持ち上がらない]
\label{aro:prop:generator-counterexample}
$R=k[\epsilon]/(\epsilon^2)$上の有限自由二周期複体のdg圏
$\mathcal C$を考える。
$G=R$を偶次数・零微分に置く。
$G_0=k$は体上のMorita生成元で、
全Hom複体は項別自由であるが、
$G$は$\mathcal C$のMorita生成元でない。
\end{proposition}
\begin{proof}
$T$を偶・奇各次数に$R$、両方向の微分を$\epsilon$倍とする。
$\ker\epsilon=\epsilon R=\operatorname{im}\epsilon$ゆえ
$H^*(T)=H^*\Hom(G,T)=0$。
一方$T$が可縮なら$1=\epsilon(a+b)$となり、
$\bmod\epsilon$で矛盾。従って$T$は非零対象で$G$は生成しない。
体上の有限自由二周期複体はそのコホモロジーと可縮部分に分解し、
$k$と$k[1]$が生成するため$G_0$はMorita生成元。
有限自由複体間のHomも項別$R$自由。
\end{proof}
\begin{remark}[条件付き(KF)から(BC)への追加仮定]
これは一般の二周期dg圏の反例であり、正の内部次数の
graded cubic MFの生成性が偽とは示していない。
しかし相対Artin節の「項別平坦Homと特殊生成子だけで
Morita生成性が持ち上がる」という旧補題は誤り。
正しい導来二側bar評価とその剰余・導来tensorとの両立を
(KF)とは別に検証する必要がある。実際のrotation積、
一般混合族のgeneric同値は未解決である。
\end{remark}

% 2026-11-11 / v32
\section{独立中心三次変数の対称性によるArtin障害の消滅}
\label{csl:sec:symmetry}

ここでは論文Bの混合中心対において、$v$が独立中心変数であることを使う。
前節の五変数三次の体上の隣接sector計算を、\emph{正しい導来自然変換複体が構成された場合}の全Artin持上げ定理へ強化する。
これにより次数一の$2\times5$障害行列を逐次計算する必要がなくなるが、
導来bar複体の正当性やkernel--functor比較(KF)そのものは証明しない。

\subsection{体上の障害sectorへの補助対称性}
$k$を三乗根$\zeta$を含む標数零の体とし、
\[
 F_\lambda=(x^3+y^3+z^3)/3+u^3+v^3+\lambda uxy,
 \qquad \lambda^3\ne-3
\]
とする。$\Gamma=\mu_3^{(v)}$を$v\mapsto\zeta v$、
$x,y,z,u\mapsto x,y,z,u$で作用させる。
混合量子化の二次関係は$v$を含まず、指定中心元の$v$への依存は$v^3$のみである。
従って$\Gamma$は量子パラメータ$h$を固定し、非可換中心対へ作用する。

\begin{proposition}[障害sectorの非自明な$\Gamma$-指標]
\label{csl:prop:characters}
$\mathcal T_0$を$F_\lambda$のgraded MF圏とする。
$\Gamma$の作用について
\[
 H^0\RHom(\Delta,\Delta(1))
 =\langle x,y,z,u\rangle\oplus\langle v\rangle,
 \qquad
 \left(H^1\RHom(\Delta,\Delta(1))\right)^\Gamma=0.
 \tag{CSL1}\label{csl:eq:characters}
\]
最初の直和の左項は自明表現で、右項は非自明な一次元表現である。
後者の$H^1$は二つの非恒等scalar-$\mu_3$ sectorの和であり、
各sectorは$\Gamma$の非自明な一次元表現となる。
\end{proposition}
\begin{proof}
前節の命題\ref{aro:prop:sector-neighbors}から
$H^0\simeq\operatorname{Jac}(F_\lambda)_1$、
$H^1\simeq k^2$である。$H^0$の表現は変数の作用から従う。
後者の各非恒等sectorは原点のみを固定し、
normal空間は五つの座標方向全部である。
Ballard--Favero--Katzarkovのsector formula
\cite[Theorem~5.39, proof]{AROBFK}では、
normal方向のtop exterior power（双対規約ではその双対）の
行列式characterが現れる。$\Gamma$はnormalの$v$方向だけを
$\zeta$倍するので、その行列式characterは$\zeta^{\pm1}$であり
自明でない。$\Gamma$はscalar-$\mu_3$と可換で各sectorを保つから、
二つのsectorの和にも不変元はない。
この判定にはcharacterの双対規約を選ぶ必要はない。
\end{proof}

\subsection{有限Artin環上の一般持上げ補題}
$R$を剰余体$k$を持つ局所Artin $k$-代数とし、$\Gamma$は$R$を固定する。
以下の$R$-線形cochain複体$C_R$は各項が$R$-平坦で、
すべてのArtin商への縮小と\emph{鎖レベルで}両立すると仮定する。
実際のdg圏の導来自然変換を正しく表すという仮定は、この条件と別に必要である。

\begin{lemma}[対称性と一つの恒等的閉元による全Artin自由性]
\label{csl:lem:equivariant-lift}
$C_R$を上記の$\Gamma$-同変な複体の族とし、特殊ファイバー$C_k$について
\[
 H^{-1}(C_k)=0,\quad H^0(C_k)=V_0\oplus kv,
 \quad V_0=H^0(C_k)^\Gamma\simeq k^4,
 \quad H^1(C_k)^\Gamma=0
\]
を仮定する。$v$の類はすべての$R$で互換な閉元$v_R\in Z^0(C_R)$へ
持ち上がるとする。このとき、
\[
 H^{-1}(C_R)=0,\qquad H^0(C_R)\simeq R^5,
 \qquad H^0(C_R)\otimes_R R/I\simeq H^0(C_{R/I})
 \tag{CSL2}\label{csl:eq:artin-free}
\]
が任意のArtin商$R/I$に対して成り立つ。
\end{lemma}
\begin{proof}
$R$の長さについて帰納し、長さ一の小拡大
$0\to I\to R'\to R\to0$を取る。ただし$\mathfrak m_{R'}I=0$。
項別平坦性から
\[
 0\to I\otimes_k C_k\to C_{R'}\to C_R\to0
\]
は複体の短完全列となる。$H^{-1}(C_k)=0$より、
帰納法で$H^{-1}(C_{R'})=0$を得る。

$V_0$の基底の不変閉代表を$C_R$に持ち上げてあるとする。
$\Gamma$は有限群で$\operatorname{char}k=0$だから、
不変部分を取る操作は完全であり、$C_{R'}^0$に不変な持上げを選べる。
その微分は$I\otimes C_k^1$に属する不変なcocycleである。
仮定$H^1(C_k)^\Gamma=0$より、
$I\otimes C_k^0$の不変な元の微分として消せる。
従って四基底は閉元のまま持ち上がる。
$v_R$には初めから互換な閉元がある。
よって$H^0(C_{R'})\to H^0(C_R)$は全射。
長完全列と$H^{-1}(C_R)=0$により
\[
 0\to I\otimes_k H^0(C_k)\to H^0(C_{R'})
 \to H^0(C_R)\to0
\]
が成り立つ。従って$H^0(C_{R'})$の$k$-長さは
$5\operatorname{length}(R')$であり、剰余体の五基底の持上げから
$R'^5\twoheadrightarrow H^0(C_{R'})$が得られる。
Nakayamaと長さの一致により同型である。
小拡大について繰り返せば任意の$R$とすべての商で結論する。
\end{proof}

\begin{theorem}[混合族の低rotation底変更に関する条件付き対称性定理]
\label{csl:thm:conditional-bc}
$F_\lambda$を特殊ファイバーとする$\Gamma$-同変な
Artin変形のgraded MFあるいは非可換MFについて、
$r=1,2$の\emph{実際の導来自然変換}を計算する$R$-線形cochain複体$C_{r,R}$が
次の条件を満たすとする：
(i) 項別$R$-平坦、任意の商への剰余が鎖レベルで一致する；
(ii) $\Gamma$同変で、特殊ファイバーでは既知のsector同定と作用に一致する；
(iii) 中心変数$v$による積が$C_{1,R}$の互換な閉元を与える。
このとき、任意の$R$に対し
\[
 H^0(C_{1,R})\simeq R^5,\qquad
 H^0(C_{2,R})\simeq R^{10},
 \tag{CSL3}\label{csl:eq:bc}
\]
かつ両者のすべてのArtin底変更が同型である。
特に論文Bの乗法的characteristic写像$\Theta_r$の
自然な特殊化が既知の多項式倍との同定であるなら、
$\Theta_1,\Theta_2$は局所環上の単元行列式判定によって同型となり、
条件(BC)が満たされる。
\end{theorem}
\begin{proof}
$r=1$には命題\ref{csl:prop:characters}と
補題\ref{csl:lem:equivariant-lift}を使う。
$r=2$では前節の\ref{aro:prop:sector-neighbors}から
$H^{-1}(C_{2,k})=H^1(C_{2,k})=0$、
$H^0(C_{2,k})\simeq k^{10}$であり、
補題\ref{aro:lem:neighbor-vanishing}の小拡大版を適用する。
積と$\Theta$の条件付き同定は、標的とsourceの自由性および特殊ファイバーでの同型から、$\det(\Theta_r)$が単元となることによる。
\end{proof}

\begin{remark}[解決したこと・残ること]\label{csl:rem:scope}
この定理は\emph{仮定した正しい複体の内部での}全次数の障害消滅であり、
任意の量子化や一般混合パラメータで条件(i)--(iii)が成立するとは主張しない。
特に非正則Artin上の導来二側bar、相対cofibrant模型、
$\Theta$の実際の導来自然変換への写像、kernel全忠実性(KF)の検証は
別に残る。これらを体上の次元公式や項別平坦性から推論しない。
一方、$\Gamma$が保存される部分族については、
以前の$2\times5$障害行列の高次係数を個別に計算する優先度は下がり、
正しいArtin互換のenhancementの構成が最優先となる。
\end{remark}

% 2026-11-11 / v33
\section{射影核による低rotation積の相対底変更}\label{gr:sec:basechange}

本節では、最終目標に必要な比較を見直し、一般の非正則基底上のMF核定理を
先に完成させる経路を迂回する。可換側のKuznetsov成分は滑らかな射影族の
admissible部分圏なので、その\emph{真の導来自然変換}はperfectな幾何学的核の
導来Homで計算できる。これと座標切断の持上げを組み合わせると、
次数一・二の全Artin底変更が従う。独立中心変数も有限群対称性も不要である。

To\"enの相対導来Morita定理\cite[Theorems~8.9, 8.15]{GRToen}、
Kuznetsovの射影核によるHochschild理論\cite{GRKuznetsov}、
perfect直接像の底変更\cite[Tag~0B91]{GRStacks}を既知の入力とする。
To\"enの基底$k$は原論文\S2で任意の可換環として固定されており、体に限られない。
最近のLin--Zhang\cite[\S2.1]{GRLZ26}にもprojected kernelを用いる相対的な枠組みがある。
本節はこれらのMorita原理の新規性を主張しない。
体上の低次数同定には\cite[Theorem~5.39]{GRBFK}と
\cite[Theorem~3.5, Remark~3.6, \S5]{GRLZ}だけを用い、
体上の定理を非正則基底へ直接引用することはしない。

\subsection{持ち上がった基底があれば$H^1$の消滅は不要}

\begin{lemma}[閉基底持上げによる$H^0$底変更]\label{gr:lem:lifts}
$(R,\mathfrak m,k)$をArtin局所環、$C\in\operatorname{Perf}(R)$とする。
$H^{-1}(C\otimes_R^{\mathbf L}k)=0$、
$H^0(C\otimes_R^{\mathbf L}k)\simeq k^d$と仮定する。
$\alpha:R^d[0]\to C$が$D(R)$の射で、その剰余の$H^0$が同型なら、
\[
 H^0(\alpha):R^d\xrightarrow{\sim}H^0(C)
 \tag{GR1}\label{gr:eq:lift-isom}
\]
は同型である。任意のArtin商$R\to R'$について自然な写像
\[
 H^0(C)\otimes_R R'\longrightarrow
 H^0(C\otimes_R^{\mathbf L}R')
\]
も同型である。$H^1(C\otimes_R^{\mathbf L}k)$の消滅は仮定しない。
\end{lemma}
\begin{proof}
$C$を有界有限自由複体で表し、各成分の閉代表を選んで$\alpha$を鎖写像とする。
$E=\operatorname{Cone}(\alpha)$も有界有限自由である。
特殊ファイバーの長完全列から$H^{-1}(E_k)=H^0(E_k)=0$。
$E$の有限$\mathfrak m$-進フィルトレーションの各商は
$(\mathfrak m^j/\mathfrak m^{j+1})\otimes_k E_k$なので、
短完全列の長完全列を繰り返すと$H^{-1}(E)=H^0(E)=0$を得る。
$R^d\to C\to E$の長完全列が\eqref{gr:eq:lift-isom}を与える。
同じ鎖写像を$R'$へテンソルし、同じ議論を行うと、両側が
$R'^d$に同定される。これが\emph{自然な}底変更写像と一致する。
\end{proof}

\begin{remark}
$[R\xrightarrow{h}R]$を次数$0,1$に置いた例では、
特殊ファイバーの次数零基底を閉元として持ち上げられない。
上の補題はperfect性や特殊次元だけで底変更を推論しているのではない。
逆に、全基底の閉持上げが既にあれば、次数一の障害を別途計算する必要はない。
\end{remark}

\subsection{相対Kuznetsov成分のprojected kernel}

$R$をArtin局所$\mathbb C$-代数、$N\ge5$、
$X=V(W)\subset\mathbb P_R^{N-1}$を滑らかな相対三次超曲面とする。
$D=\operatorname{Perf}(X)$、$a=N-3$、$E_j=\mathcal O_X(j)$と置き、
\[
 \mathcal A_R=\langle E_0,\ldots,E_{a-1}\rangle^\perp\subset D
\]
をperfect Kuznetsov成分とする。$D^b(\operatorname{coh}X)$全体とは同一視しない。

射影空間の相対線束コホモロジーと
$0\to\mathcal O_{\mathbb P}(j-3)\to\mathcal O_{\mathbb P}(j)\to\mathcal O_X(j)\to0$
により、$E_0,\ldots,E_{a-1}$は$R$-例外列である。
$R\operatorname{Hom}(E_j,E_i)=0$（$j>i$）と
$R\operatorname{Hom}(E_i,E_i)=R$は任意の底変更に両立する。

核の規約を$\Phi_K(M)=Rp_{2*}(p_1^*M\otimes^{\mathbf L}K)$とし、
$K\star L$は$\Phi_K\Phi_L$を表す畳込みとする。
\[
 \mathsf M_j=\operatorname{Cone}(E_j^\vee\boxtimes E_j\longrightarrow\mathcal O_\Delta),
 \qquad P=\mathsf M_0\star\cdots\star\mathsf M_{a-1}
 \tag{GR2}\label{gr:eq:projector}
\]
は順に左mutationと左射影$p=i\ell$を表す。
$i:\mathcal A_R\hookrightarrow D$、$\ell\dashv i$、$\ell i\simeq\mathrm{Id}$である。
右端のmutationから作用することに注意する。

$T=\Delta_*\mathcal O_X(1)$、
\[
 \rho=L_{E_0}\circ(-\otimes\mathcal O_X(1))|_{\mathcal A_R},
 \qquad Q=P\star\mathsf M_0\star T\star P
 \tag{GR3}\label{gr:eq:rotation-kernel}
\]
とする。$M\in\mathcal A_R$なら、$j\ge1$に対し
$R\operatorname{Hom}(E_j,M(1))=R\operatorname{Hom}(E_{j-1},M)=0$であり、
$R\operatorname{Hom}(E_j,E_0)=0$でもある。従って$\rho$は$\mathcal A_R$を保ち、
$Q$は$i\rho\ell$を表す。$Q^{\star r}$は$i\rho^r\ell$を表す。

\begin{theorem}[真の導来Natのperfect模型]\label{gr:thm:perfect-target}
上の状況で、$q:X\times_R X\to\operatorname{Spec}R$として
\[
 \mathcal C_r=Rq_*R\mathcal Hom(P,Q^{\star r})
 \tag{GR4}\label{gr:eq:nat-complex}
\]
は$R$-perfectであり、自然な擬同型
\[
 \mathcal C_r\simeq R\operatorname{Hom}_{\mathcal A_R^e}
       (\Delta_{\mathcal A_R},P_{\rho^r})
\]
を持つ。任意の$R\to R'$に対する導来底変更は、
$X_{R'}$とその同じ射影核から作る複体への自然な擬同型を与える。
比較は自然変換の合成と両立する。
\end{theorem}
\begin{proof}
$X/R$は滑らかなので相対対角は正則埋込みであり、$\mathcal O_\Delta$と$T$はperfect。
従って\eqref{gr:eq:projector}の核はperfectである。
畳込みは滑らかな三重積の射影に沿う導来引戻し・テンソル・proper直接像である。
\cite[Tag~0B91]{GRStacks}により、perfect性と任意の導来底変更との両立が保たれる。
同じ結果を$P^\vee\otimes Q^{\star r}$の直接像へ適用すると、
$\mathcal C_r$のperfect性と底変更が従う。

導来自然変換との同定は体上の次元比較ではない。
\cite[Theorem~8.15]{GRToen}を基底環$R$に適用すると、$D$のenhanced自己関手間の
導来Homは幾何学的核間の導来Homに同定される。
さらに$F\mapsto iF\ell$はenhanced関手圏上で完全忠実である。
実際、$i$上の制限と$\ell i\simeq\mathrm{Id}$から逆向きの写像を得る。
任意の$M\in D$に対する単位$M\to i\ell M$は、
$iF\ell$へ作用させると同型になるため、自然性により
変換は$\mathcal A_R$上の制限から一意に復元される。
同じ単位と三角恒等式はcoherent自然変換複体上の互いに逆な比較を与える。
従って$P$と$Q^{\star r}$の間のHomは、
$\mathrm{Id}_{\mathcal A_R}$と$\rho^r$の間の真の導来自然変換を計算する。

核と変換は全て評価射・mutation・畳込みから作られる。
これらの導来底変更と射の合成は可換であるので、最後の両立性も得る。
無界で項別自由なbar複体をK-flatとみなす段階はない。
\end{proof}

\subsection{座標切断からの全基底持上げ}

任意の$s\in H^0(X,\mathcal O_X(1))$に対し、
\[
 \alpha_s(M):M\xrightarrow{\,s\,}M(1)
       \longrightarrow L_{E_0}(M(1))=\rho(M)
 \tag{GR5}\label{gr:eq:sections}
\]
はenhancedな自然変換である。後半はmutationの単位であり、
核の射としても$R$上で定義される。
$x_1,\ldots,x_N$はこうして$H^0(\mathcal C_1)$の閉持上げを与える。
次数二には$\rho(\alpha_{x_j})\circ\alpha_{x_i}$を用いる。
この段階では順序付き積を使えば十分で、可換性を先取りしない。

\begin{lemma}[特殊ファイバーの低次数入力]\label{gr:lem:field}
$F=W\bmod\mathfrak m$に対し、$r=1,2$では
\[
 H^{-1}(\mathcal C_{r,k})=0,\qquad
 H^0(\mathcal C_{r,k})=\operatorname{Jac}(F)_r.
 \tag{GR6}\label{gr:eq:field}
\]
\eqref{gr:eq:sections}の特殊化はJacobian環の$x_i$に対応し、
その合成は多項式積に対応する。従って次数一の$N$個の座標と、
次数二で適当に選んだ$\binom N2$個の二次単項式が基底を与える。
\end{lemma}
\begin{proof}
体上のBFK公式とLin--Zhangの乗法同定を使う
\cite[Theorem~5.39]{GRBFK}\cite[Theorem~3.5, Remark~3.6, \S5]{GRLZ}。
恒等sectorは偶数コホモロジー次数$j=2b$に$\operatorname{Jac}(F)_{r+3b}$を与える。
非恒等sectorは、$j\equiv N\pmod2$かつ
\[
 2r+3j-N=0
\]
のときだけ寄与する。$j=-1$なら$N=2r-3\le1$、
$j=0$なら$N=2r\le4$なので、$N\ge5$ではいずれも起こらない。
特殊ファイバーの勾配は$N$個の二次正則列なので、
次数一・二の次元は$N,\binom N2$である。
座標による変換とその合成は同公式の恒等sectorの多項式作用である。
ここではrotation自体を入力するため$\gcd(N,3)=1$を要求しない。
\end{proof}

\begin{theorem}[全Artin低rotation底変更]\label{gr:thm:bc}
任意のArtin局所$\mathbb C$-代数$R$と、滑らかな
$X=V(W)\subset\mathbb P_R^{N-1}$（$N\ge5$）について
\[
 \operatorname{Nat}_{dg}^0(\mathrm{Id},\rho)\simeq R^N,
 \qquad
 \operatorname{Nat}_{dg}^0(\mathrm{Id},\rho^2)\simeq R^{\binom N2}.
 \tag{GR7}\label{gr:eq:ranks}
\]
自然な$H^0$の底変更は全Artin商で同型で、合成と両立する。
独立な中心変数・有限群対称性・$H^1$の消滅は仮定しない。
\end{theorem}
\begin{proof}
定理\ref{gr:thm:perfect-target}により真の標的をperfect複体$\mathcal C_r$で計算できる。
$r=1$では$\alpha_{x_i}$を用い、$r=2$では特殊ファイバーの基底となる
$\binom N2$個の順序付き合成を選ぶ。
これらはそれぞれ$R^N[0]\to\mathcal C_1$、
$R^{\binom N2}[0]\to\mathcal C_2$という射を与える。
補題\ref{gr:lem:field}と補題\ref{gr:lem:lifts}を適用する。
全ての射が核上で底変更と可換なので、得られる$H^0$の同型も自然なものであり、
積と底変更の両立が従う。
\end{proof}

\begin{remark}[障害群を消すより短い理由]
五変数の$r=1$では特殊ファイバーの$H^1$は二次元であり、
七変数の$r=2$でも二つの非恒等sectorが$H^1$に寄与する。
それでも上の定理は成立する。幾何学的座標によって必要な全基底が
既に閉元として持ち上がっているからであり、障害群そのものを零にする必要はない。
これは以前の中心対称性による消滅結果を否定するものではなく、
本問題の可換標的には不要な追加条件を取り除くものである。
\end{remark}

\subsection{既存Artin比較に使う生成性の仮定の点検}

\begin{lemma}[導来特殊化とcompact生成子]\label{gr:lem:derived-generation}
$R$をArtin局所$\mathbb C$-代数、$\mathcal B$を小さな冪等完備pretriangulated
$R$-線形dg圏とする。全Hom複体が$R$-K-flatであり、通常の係数特殊化が
導来Morita特殊化$\mathcal B\otimes_R^{\mathbf L}k$を表すとする。
$G\in\mathcal B$の剰余がその導来特殊ファイバーを生成すれば、$G$は$\mathcal B$を生成する。
とくに、\emph{項別平坦性だけ}をこの仮定の代わりに使うことはできない。
\end{lemma}
\begin{proof}
$\mathcal D=\operatorname{Ind}(\mathcal B)$を取る。$G$はcompactであり、
$R\operatorname{Hom}(G,-)$は全colimitと可換する。
compact対象のHomの導来底変更は表現可能対象で成立し、
両辺がcolimitと可換することから任意の$M\in\mathcal D$へ延長される。
$R\operatorname{Hom}(G,M)=0$なら、特殊ファイバーの生成性から
$M\otimes_R^{\mathbf L}k=0$。
$R$は$D(R)$内で$k$の有限extensionから作られるので、
$M=M\otimes_R^{\mathbf L}R=0$である。
従って$G$は$\mathcal D$のcompact generatorであり、compact部分を取ると結論を得る。
\end{proof}

本研究の正の内部次数付き有限rank MFでは、固定cohomological次数のHomは
有限自由$R_m$-加群であり、特殊ファイバーのHomの全cohomologyは有限次元かつ有界である。
特殊Hom複体の体上の分裂収縮を選ぶと、$h$の冪零性により
homological perturbationの逆は有限和となる。
従って実際のHom複体は有界有限自由$R_m$-複体へホモトピー同値であり、K-flatである。
これは既存Artin残余比較の\texttt{am:lem:mf-hom}の入力である。
特殊化関手はHomで導来底変更を表し、その像が古典生成子を含むことを
windowの三つのstabilizationで確認してから、上の補題を適用できる。
零曲率の任意の二周期圏で項別自由性だけを使う旧一般補題の反例とは矛盾しない。
本節ではこの必要なK-flat性・導来特殊化・生成子の像を分離しており、
形式MF復元やstabilizationの全証明を再証明したとはしない。

\subsection{非可換入力への適用と依存する既存結果}

以下で用いる既存入力を明記する。
\begin{description}
\item[(F)] 全有限Artin切断で、指定rotationと底変更を保つ
$\mathcal K^{\rm lin}_{B_m}\simeq\operatorname{Ku}_{\rm perf}(Y_m/R_m)$がある。
これは現行論文Bの形式MF復元とArtin残余比較を合成した結果である。
本節の幾何学的定理は(F)に依存しないが、非可換側への適用は依存する。
\item[(D)] 計算した有限対角の低次数$H^0$は自由で、全Artin底変更に両立し、
実際の残余pairへの乗法的characteristic写像$\Theta$がある。
これは局所正則対角、乗法的bar比較、および(F)のwindowとの両立を入力とする。
\end{description}
ここに一般のMF kernel--functor比較(KF)を追加してはならない。
(F)と(D)は既存証明への依存であり、本節だけで全既存証明を独立再監査したとの意味ではない。

\begin{corollary}[既存形式比較の下での(BC)の充足]\label{gr:cor:nc}
(F)が成立する中心三次対について、標的側の条件(BC)は
定理\ref{gr:thm:bc}から従う。(D)も成立し、特殊ファイバーの$\Theta_1,\Theta_2$が
Jacobianの標準同定なら、両写像は全有限Artin切断で同型である。
\end{corollary}
\begin{proof}
(F)をdg双加群の導来自然変換へ適用し、標的を\eqref{gr:eq:nat-complex}で計算する。
Morita同値はrotation双加群とその合成を保つ。
標的は定理\ref{gr:thm:bc}により所定の階数の自由加群で底変更に両立する。
(D)のsourceも同じ階数の自由加群である。
$\Theta_r$の行列式は剰余で非零なので単元であり、同型である。
乗法性と全切断の整合性は入力した$\Theta$のものを保つ。
\end{proof}

可換側で低積を厳密にJacobian積と同定するには、既存の相対window比較を用いる。
可換curved Hochschild複体では導分$\partial_i$の微分が、符号を除いて$t\partial_iW$である。
従って$t^{-1}\partial_i$は$\partial_iW$の有限arityの原像を与える。
characteristic写像により、偏微分による自然変換は境界になる。
windowは多項式倍と内部shiftを保ち、mutationの単位は商で同型になるので、
これは\eqref{gr:eq:sections}の多項式作用と一致する。
相対勾配Koszulの完全性から$\operatorname{Jac}(W)_1,\operatorname{Jac}(W)_2$は所定の階数で自由。
これらから$H^0(\mathcal C_r)$への写像は特殊化で同型なので、全Artinで同型である。
ここでもMF核の全忠実性は用いない。

\begin{corollary}[計算済み三次式の形式的同定]\label{gr:cor:formal-polynomial}
既存入力(F),(D)を満たす$\lambda=1$の混合族について、
有限対角から得た多項式$G_h+v^3$と形式復元先$W_h$は、
$h=0$で恒等な$\operatorname{GL}_5(\mathbb C[[h]])$変換で同型である。
従って全ての$m$で、非可換線束残余pairは
$V(G_h+v^3\bmod h^{m+1})$のperfect Kuznetsov pairと同値である。
これは非零複素パラメータのgeneric圏同値を主張しない。
\end{corollary}
\begin{proof}
系\ref{gr:cor:nc}により低次数の積が同定され、次数二の核として勾配空間が一致する。
次数一の基底の対応は特殊ファイバーで恒等に規格化でき、
全切断の互換性により逆極限の線形対応を得る。
この座標変更後、$\nabla W_h=A_h\nabla(G_h+v^3)$、$A_h\equiv1\pmod h$と書ける。
三次式の対称三重形式を$T$とすると、$A_h^{\mathsf t}$は
$T(A_h^{\mathsf t}u,v,w)=T(u,A_h^{\mathsf t}v,w)=T(u,v,A_h^{\mathsf t}w)$を満たす。
$A_h^{\mathsf t}$の形式三乗根も同じ性質を持つため、
$C_h=(A_h^{\mathsf t})^{1/3}$によって
$(G_h+v^3)(C_hx)=W_h(x)$となる。
この形式座標同型と(F)を各切断で合成する。
\end{proof}

\subsection{generic復元までの最短経路と残る証明}

本節により、既存の形式残余比較を利用できる範囲では、
一般の相対MF核比較(KF)、無界barの全面的な特殊化理論、
中心対称性による次数一障害の逐次処理を、最終目標の前提にする必要がなくなった。
有限代数的抽出という目標は維持し、標的の処理だけを幾何学的射影核へ置き換える。

しかし\emph{式の代数性}と\emph{圏同値の代数性}は別である。
(F)から得た全Artin同値や上の形式座標変換から、
非可換側の指定生成子がgenericでも生成するとは推論できない。
次に必要なのは、代数的な残余模型と$\operatorname{Ku}(V(G_h+v^3))$を結ぶ
双加群、逆双加群、両合成の単位・余単位、rotation intertwinerを
\emph{共通の有限cell/retractデータ}で構成することである。
その有限データがあれば、構造方程式・逆写像・ホモトピーの恒等式を
有限の代数方程式として検証し、存在する開集合へ広げられる。
現時点でその有限証明書、一般混合型の相対対角perfect性、generic生成性は未構成である。
滑らかなproper dg代数の有限$A_n$決定性\cite{GRToenFinite}を利用する場合も、
非可換側のsmoothnessを先に確認する必要があり、ここを循環させない。

本節の計算コードは、体上sector重みと閉基底持上げ補題の有限自由複体例を検査する。
幾何学的核定理、既存の形式MF同値、または非可換対角の全行列を計算機で検証したものではない。

% 2026-11-11 / v34
\section{片側compact核からのMorita同値と代数化}
\label{os:sec:algebraization}

本節では、一般ファイバーへの移行のために必要な有限データを減らす。
非可換側のproper dg代数を$\Gamma$、既に構成した可換cubicの
perfect Kuznetsov成分のsmooth proper生成代数を$\mathsf C$と書く。
前節までの周囲の次数付き代数$A$、中心商$B$とは異なる記号である。
関手の向きを\emph{可換側から非可換側へ}取ることが重要である。

主張は二段階からなる。第一に、右$\Gamma$-perfectな一つの
$(\mathsf C,\Gamma)$-双加群が特殊点で同値を与えれば、底上でも同値になる。
逆核と二つの逆性を独立の未知量にする必要はない。
第二に、そのような\emph{ordinaryな導来圏のcompact核}が完備底上に存在すれば、
perfect加群の有限表示性とPopescuの定理により、指定有限jetを保って
\'etale近傍に代数化できる。
これはArtin切断ごとのperfect性からcompact核の存在を結論する定理ではない。

使用する既知入力はTo\"en--Vaqui\'eの相対pseudo-perfect性とperfect性の比較、
およびperfect加群のfiltered-colimit降下である
\cite[Lemma~2.8, Lemma~2.10]{OSTV}。
後者とPopescuの定理を使うため、非可換側のsmoothnessを先に仮定したり、
未検証の有限$A_n$決定性を使ったりしない。

\subsection{smoothなのは片側だけでよい}

全てのテンソル積・Hom・包絡代数は導来のものとする。
通常のdg模型を使う場合は、基底環上K-flatな置換を固定する。
右加群の規約を用い、$P$が$(\mathsf C,\Gamma)$-双加群なら
$F_P(M)=M\otimes_{\mathsf C}^{\mathbf L}P$とする。
properとは基底環上perfectであること、smoothとは対角が導来包絡代数上perfectであることをいう。

\begin{lemma}[片側perfectから双加群perfectへ]\label{os:lem:compact}
$R$を可換環、$\mathsf C$を$R$上smoothなdg代数、$\Gamma$を任意の$R$-dg代数とする。
$(\mathsf C,\Gamma)$-双加群$P$が右$\Gamma$-加群としてperfectなら、
\[
 P\in\operatorname{Perf}(\mathsf C^{\mathrm{op}}\otimes_R^{\mathbf L}\Gamma).
 \tag{OS1}\label{os:eq:compact}
\]
また$R$上perfectな右$\mathsf C$-加群は$\mathsf C$上perfectである。
\end{lemma}
\begin{proof}
$\mathsf C$の対角を、$\mathsf C\otimes_R\mathsf C$から有限coneと直和因子で作る。
右側に$P$をテンソルすると、$P$は$\mathsf C\otimes_R P$から同じ有限操作で得られる。
$P$の右$\Gamma$-perfect性により、$\mathsf C\otimes_R P$は双加群としてperfectである。
後半も、対角の有限表示に$R$-perfectな加群をテンソルすれば従う。
これは\cite[Lemma~2.8(2)]{OSTV}の一対象版である。
\end{proof}

\begin{theorem}[片側核の特殊ファイバー判定]\label{os:thm:local}
$(R,\mathfrak m,k)$をNoether局所環、$\Gamma,\mathsf C$を$R$上properなdg代数とし、
$\mathsf C$は$R$上smoothとする。
$(\mathsf C,\Gamma)$-双加群$P$が右$\Gamma$-perfectであり、
\[
 F_{P_k}:\operatorname{Perf}(\mathsf C_k)
       \xrightarrow{\sim}\operatorname{Perf}(\Gamma_k)
 \tag{OS2}\label{os:eq:special-equivalence}
\]
が同値であると仮定する。添字$k$は導来特殊化を表す。
このとき$F_P$は$R$上のdg Morita同値である。
特に$\Gamma$のsmoothnessは仮定ではなく結論となる。
逆核は
\[
 Q=\operatorname{RHom}_{\Gamma}(P,\Gamma)
 \tag{OS3}\label{os:eq:inverse}
\]
で与えられ、自然な射
\[
 \mathsf C\longrightarrow P\otimes_\Gamma^{\mathbf L}Q,
 \qquad Q\otimes_{\mathsf C}^{\mathbf L}P\longrightarrow\Gamma
 \tag{OS4}\label{os:eq:unit-counit}
\]
はともに擬同型である。
\end{theorem}
\begin{proof}
$P$は右$\Gamma$-perfectで、$\Gamma$は$R$-perfectだから、$P$も$R$-perfectである。
有限cellと直和因子による帰納法により、任意の$N\in\operatorname{Perf}(\Gamma)$について
$\operatorname{RHom}_{\Gamma}(P,N)$は$R$-perfectで、任意の導来底変更と両立する。
$\mathsf C$のsmoothnessと補題\ref{os:lem:compact}から、これは右$\mathsf C$-perfectでもある。
従って
\[
 F_P:\operatorname{Perf}(\mathsf C)\rightleftarrows
 \operatorname{Perf}(\Gamma):G_P,
 \qquad G_P(N)=\operatorname{RHom}_{\Gamma}(P,N)
\]
は両側のperfect部分に制限できる随伴である。
$P$の双対可能性から$G_P(N)=N\otimes_\Gamma^{\mathbf L}Q$となる。

単位の$\mathsf C$における値は
$\mathsf C\to\operatorname{RHom}_\Gamma(P,P)$であり、両項は$R$-perfectである。
その底変更は\eqref{os:eq:special-equivalence}の単位なので同型である。
よってconeは特殊化で零となる$R$-perfect複体であり、導来Nakayamaにより零である。
ここで必要なNakayamaは有界有限自由複体に対するもので、
単元を持つ微分成分を順に消去しても証明できる\cite[Tag~068U]{OSStacksSplit}。
単位は全$\operatorname{Perf}(\mathsf C)$上で同型となり、$F_P$は完全忠実である。

$Q=G_P(\Gamma)$は$\mathsf C$-perfectなので、$F_P(Q)$は$\Gamma$-perfect、従って$R$-perfectである。
余単位$F_P(Q)\to\Gamma$も底変更と両立し、その特殊化は同値の余単位である。
同じNakayamaによりこの余単位は同型となる。
像は$\Gamma$を含み、有限coneと直和因子に閉じているので本質的全射である。
随伴から単位・余単位の三角恒等式も自動的に成立する。
最後にsmoothnessのMorita不変性により$\Gamma$はsmoothとなる。
\end{proof}

\begin{corollary}[開近傍と有限な二つのcone]\label{os:cor:open}
上のデータがNoether基底$S$上で定義され、点$s$で\eqref{os:eq:special-equivalence}が成立するとする。
$\mathsf C/S$のsmooth proper性、$\Gamma/S$のproper性、$P$の右perfect性を仮定すれば、
$s$の開近傍で$F_P$は同値である。
\end{corollary}
\begin{proof}
\eqref{os:eq:unit-counit}の二つのconeは$S$-perfectである。
特殊ファイバーで零のperfect複体はその点のある開近傍で零である。
有限個の単元小行列による消去、あるいはcohomologyの閉supportから従う。
二つの開近傍の共通部分に前定理の議論を適用する。
\end{proof}

\subsection{rotationと三次周期性も片側で処理する}

$L_{\mathsf C},L_\Gamma$を、それぞれperfect部分上の自己同値を表す可逆双加群とする。
$\mathsf C$がsmoothなので、右$\Gamma$-perfectな双加群の間の比較
\[
 \theta:L_{\mathsf C}\otimes_{\mathsf C}^{\mathbf L}P
       \longrightarrow P\otimes_\Gamma^{\mathbf L}L_\Gamma
 \tag{OS5}\label{os:eq:rotation}
\]
は、\eqref{os:eq:compact}のperfect双加群の間の射である。
$\theta_k$が同型なら、そのconeの$R$-perfect性から$\theta$も同型となる。
従ってrotation比較の\emph{逆射}を別に未知量として用意する必要はない。

\begin{proposition}[周期比較のscalar補正]\label{os:prop:period}
定理\ref{os:thm:local}の状況で$R$は$\mathbb Q$-代数とし、
\[
 HH^0(\mathsf C/R)=R,
 \quad L_{\mathsf C}^{\otimes3}\simeq\mathsf C[2],
 \quad L_\Gamma^{\otimes3}\simeq\Gamma[2]
\]
を指定する。\eqref{os:eq:rotation}が同型で、その特殊化が指定周期性と両立するとする。
このとき、剰余体を変えない\'etale近傍で$\theta$をscalar倍することにより、
三次周期性とも両立させられる。
\end{proposition}
\begin{proof}
$\theta$を三回合成して二つの周期同型と比較すると$P[2]$の双加群自己同型を得る。
$F_P$は同値だから、その次数零自己準同型環は$HH^0(\mathsf C/R)=R$である。
従って差はある$u\in R^\times$の乗法であり、$u\bmod\mathfrak m=1$。
$\theta$を$c\theta$へ置換すると差は$c^3u$となる。
$c^3u=1$、$c\equiv1$を満たす分枝は、初期点で微分$3$が可逆なので\'etale局所に存在する。
補正後は指定された三次周期同型の図式が可換となる。
これはこの周期図式の可換性の主張であり、別途指定した任意の高次群作用の
全coherenceを自動的に構成するという主張ではない。
\end{proof}

\subsection{完備底のcompact核から\'etale代数化へ}

\begin{theorem}[片側compactな形式核の代数化]\label{os:thm:algebraization}
$R$をNoether局所G-ring、$\widehat R$をその完備化とする。
$\Gamma,\mathsf C$は$R$上proper、$\mathsf C$はsmoothとし、必要ならrotation双加群を指定する。
\emph{通常の導来圏}に双加群
\[
 \widehat P\in D(\mathsf C_{\widehat R}^{\mathrm{op}}
                    \otimes_{\widehat R}^{\mathbf L}\Gamma_{\widehat R})
 \tag{OS6}\label{os:eq:formal-datum}
\]
があり、右$\Gamma_{\widehat R}$-加群としてperfectで、剰余で同値を与えると仮定する。
rotationを扱う場合は\eqref{os:eq:rotation}の形式版の射$\widehat\theta$を与え、
その剰余が同型であると仮定する。

このとき任意の整数$\ell\ge1$について、剰余体を変えない\'etale近傍$R\to R'$上に
$P',\theta'$を構成できる。指定した点の近傍で$P'$はMorita同値を与え、
$\theta'$はrotationを保つ同型である。
$R'/\mathfrak m^\ell R'\simeq R/\mathfrak m^\ell$の下で、これらのデータは
$\widehat P,\widehat\theta$の$\mathfrak m^\ell$剰余と同型である。
一般に完備化全体で\emph{同じ形式分枝}を再現するとは主張しない。
\end{theorem}
\begin{proof}
$R\to\widehat R$はregularなので、Popescuの定理により
\[
 \widehat R=\mathop{\mathrm{colim}}_i R_i
 \tag{OS7}\label{os:eq:popescu}
\]
と書ける。ここで$R_i$はsmoothな有限表示$R$-代数である\cite[Tag~07GC]{OSPopescu}。
補題\ref{os:lem:compact}により$\widehat P$は\eqref{os:eq:formal-datum}の包絡代数上perfectである。
K-flatなdg模型を用いると、包絡代数の$R_i$への底変更のfiltered colimitは完備底の包絡代数となる。
perfect dg加群とその間の射はfiltered colimitに降下する
\cite[Lemma~2.10]{OSTV}ので、ある$i$に$P_i$が存在する。
これは有限cell表示とhomotopy idempotentを有限段階に降ろすことに相当し、
無限個の$A_\infty$係数を有限個だと見なす操作ではない。

$\mathsf C_{R_i}$は$R_i$-properだから、包絡代数上perfectな$P_i$を
右$\Gamma_{R_i}$へ制限してもperfectである。
rotationを施した二つの核も右perfectであり、$\mathsf C_{R_i}$のsmoothnessから双加群perfectである。
従って$i$を大きくすれば$\widehat\theta$も射$\theta_i$として降下する。

$R_i\to\widehat R\to R/\mathfrak m^\ell$を取る。
smooth代数の剰余切断の\'etale持上げ\cite[Tag~07M7]{OSEtale}により、
$R'/\mathfrak m^\ell R'\simeq R/\mathfrak m^\ell$と
$R_i\to R'$を得る。この写像で$P_i,\theta_i$を底変更する。
指定剰余点で$P'$は元と同じ同値を与え、右$\Gamma_{R'}$-perfectである。
定理\ref{os:thm:local}と系\ref{os:cor:open}により、その点の開近傍では同値となる。
$\theta'$の可逆性も、そのperfectなconeと剰余での可逆性から従う。
必要なら命題\ref{os:prop:period}のscalar補正を行う。
形式周期性が初めから整合し、$\ell$次まで保存した場合は、
同じ\'etale根を$1$ modulo $\mathfrak m^\ell$に選ぶことでそのjetも保持できる。
\end{proof}

\begin{corollary}[一般点への適用]\label{os:cor:generic}
$R$が既約な有限型$\mathbb C$-基底の一点の局所環で、前定理の入力が与えられれば、
その基底の幾何学的一般点でrotation付き同値を得る。
\end{corollary}
\begin{proof}
得られた\'etale近傍の指定点を含む成分は、基底の非空開へ支配的である。
その一般点の体は基底の関数体の有限分離拡大であり、同値をそこで底変更できる。
両側の随伴逆性は核の擬同型であるため、さらに代数閉包へ移しても保たれる。
\end{proof}

\subsection{本研究への適用範囲と残る一つの構成}

$\mathsf C_R$は明示した滑らかな可換候補$Y_G/R$のperfect Kuznetsov成分の生成代数、
$\Gamma_R$は非可換の指定残余生成代数とする。
既存の全Artin同値を逆向きに取り、可換側生成子$H$の像を$P_m$とする。
次に確認すべきことは、これらが\emph{同じ完備底上の核}
$\widehat P$とrotation射として実現され、
\[
 \boxed{\quad\widehat P\in\operatorname{Perf}(\Gamma_{\widehat R})\quad}
 \tag{OS8}\label{os:eq:remaining}
\]
となることである。左$\mathsf C_{\widehat R}$作用とrotation射の導来的coherenceも入力に含める。
この条件を確認すれば、逆双加群・両合成の逆性・代数化・一般点での生成性は
上の定理から従い、それらを独立した未知問題にする必要はない。

\eqref{os:eq:remaining}は\emph{本混合族についてまだ証明していない}。
固定rankの自由MFとして表示できること、$R$上perfectなHomを持つこと、
各$P_m$がそれぞれperfectであることは、この条件とは異なる。
たとえば体上で$\Gamma=k[e]/(e^2)$、$P=k$なら、$P$は$k$-perfectだが
$\Gamma$-perfectではない。$e$による無限自由分解を$k$へtensorすると
全次数のTorが$k$となるからである。
本研究の一般的生成子持上げの反例や全Artinでsmoothな反例とも矛盾しない。

具体的な次の着手は、可換側の\emph{一つの}生成子の形式逆像を追跡し、
$\widehat\Gamma$の有限個のshiftのconeの直和因子として書くことである。
各Artin段階ごとの表示ではなく、同じ有限filtrationを持つ表示が必要である。
その表示がないままPopescuやArtin近似を適用しない。
同値の順方向の像だけでは、上の非対称な判定の仮定を満たしたことにならない。

\subsection{検証と新規性の境界}

本節はsmooth対角、随伴、perfect Nakayama、perfect加群の有限表示性、
Popescuの定理を組み合わせた研究用の判定である。
これらの一般原理自体の新規性は主張しない。
独立に検算するコードは、ねじった行列積による通常のMorita模型で
随伴評価と双加群恒等式を確認し、scalar周期補正とcompact性の注意例を扱う。
非可換cubicの$\widehat P$を構成・検証したコードではない。


\part{一般点への適用と飽和の別構成}
\section{問題、既知の入力、主結果}\label{RSRsec:intro}
本研究の目標は、任意の$n\ge3$と滑らかな複素三次超曲面$X^n$について、
そのPoissonスキームの各被約既約成分の幾何学的一般点で、構成した量子化の残余dg圏を
可換な滑らかな三次超曲面のrotation付きKuznetsov成分として実現することである。
全てのAS-regular中心三次商の全点に同じ存在を要求する問題ではない。

以下の全ての圏、関手、自然変換にはdg enhancementを保ち、底変更はMorita底変更を意味する。
圏の同値を単なる三角圏の同値へ弱めない。
$R$上properとは対象間の全導来Homが$R$-perfectであること、smoothとは対角が
包絡dg圏上perfectであることをいう。
$V^\vee=\RHom_R(V,R)$とし、直交条件は全cohomological shiftで課す。

\begin{hypothesis}[局所圏論的入力]\label{RSRhyp:categorical}
$(R,\mathfrak m,k)$はnoetherian局所環である。
$\cC$は$R$-properな冪等完備pretriangulated dg圏で、dg自己同値$T$と対象$\OO$を持つ。
$\OO_i=T^i\OO$と書き、有限個の$\OO_i$の直和$E$が$\cC$を古典的に生成する。
ある滑らかな三次超曲面$X\subset\mathbb P^{n+1}_k$、$n\ge3$に対し、指定された同値
\[
 (\cC_k,T_k,\OO_k)\simeq(\Perf(X),-\otimes\OO_X(1),\OO_X)
 \tag{L}\label{RSReq:local-input}
\]
がある。ここで$\cC_k$は導来Morita底変更の冪等完備化であり、対象間のHomの
底変更写像はこの同定と両立する。
\end{hypothesis}

$a=n-1$と置く。相対例外列$\OO_0,\ldots,\OO_{a-1}$と、その右直交
\[
 \cK=\langle\OO_0,\ldots,\OO_{a-1}\rangle^\perp\subset\cC
\]
を得ることは以下で証明する。

\begin{theorem}[rotationの自己同値性]\label{RSRthm:rotation}
仮定\ref{RSRhyp:categorical}の下で
\[
 \rho=\LL_{\OO}\circ T:\cK\longrightarrow\cK
\]
はdg自己同値であり、その逆関手は
\[
 \boxed{\ \rho^{-1}=T^{-1}\circ\RR_{\OO}\big|_{\cK}\ }.
 \tag{R}\label{RSReq:inverse}
\]
評価・余評価のconeから、両合成と恒等関手との自然な同型が得られる。
この主張には$\cK$のsmoothnessを仮定しない。
\end{theorem}

\begin{theorem}[saturation後に従う二つの帰結]\label{RSRthm:summary}
第\ref{RSRsec:projective}節の局所中心三次商$B$とそのproper模型
$\cC_B=\thick\{\OO(j)\}$、$\cK_B=\Perf(\Gamma_R)$を考える。
追加条件
\[
 \Gamma_R\in\Perf(\Gamma_R^{\mathrm{op}}\otL_R\Gamma_R)
 \tag{SM}\label{RSReq:SM}
\]
を仮定すると、次が成立する。
\begin{enumerate}[label=(\arabic*)]
\item $\cC_B=D^b(\qgr B)$であり、線束生成残余圏は全coherent残余圏と一致する。
\item $R=(R_0)_{\mathfrak m}$、$R_0$が有限型複素整域なら、
$\mathfrak m$を含むある主開集合上へ$(\Perf(\Gamma_R),\rho)$は
smooth properな代数的dg圏と可逆perfect双加群のpairとして降下する。
\end{enumerate}
従って、これら二項を\eqref{RSReq:SM}とは独立な未解決仮定として重ねて課す必要はない。
\end{theorem}

この二定理と、定理\ref{er:thm:ER}の局所実現定理を合わせたgeneric復元を
定理\ref{RSRthm:generic}に記す。\eqref{RSReq:SM}自体の証明とは区別する。

\subsection{先行研究との関係と主張の状態}
mutationの基本操作は既知である。ここでの要点は、特殊ファイバーのSerre traceから
局所底上の必要な双対性を作り、逆rotationが残余圏を保つことを実際に証明する点にある。
「properだから任意の関手に随伴がある」とは主張しない。

有界非可換射影導来圏についてはBurke--Stevenson
\cite[Proposition 4.18, Lemmas 6.1--6.2, Theorem 6.4]{RSRBS}を使う。
smooth proper dg代数の有限型底への降下はToën\cite[Théorème 1]{RSRToen}の定理である。
perfect加群のfiltered colimitに関する性質は同論文F4、
Toën--Vaquié\cite[Lemma 2.10]{RSRTV}による。
これらの既知の定理を本研究の新結果として数えない。
Orlov\cite[Theorem 3.25]{RSROrlov}のgluing判定は体上の記述なので、
第\ref{RSRsec:fullness}節では必要な方向を一般の可換底上の双加群三角で証明する。

Lin--Zhang\cite[\S2.1, Theorem 2.2, Proposition 2.4]{ERLinZhang}は、幾何的な相対admissible成分について、
同値核の変形障害の消去と代数的correspondenceを扱う。
その入力には相対perfect projectorと底変更の両立がある。
第\ref{er:sec:etale}節では、同じ障害消去原理を実際のperfect双加群のmoduliで定式化し、
非幾何的なtargetとrotation intertwinerを含めて証明する。
同論文をそのまま非幾何的なdg族へ適用したとは扱わない。

本稿のrotation定理と、その後の還元命題には証明を記す。
具体的量子化への適用は第\ref{RSRsec:projective}節の入力の検証に依存する。
全Poisson成分での局所量子化の存在と\eqref{RSReq:SM}は未解決である。
本稿の組合せとしての新規性は未確定であり、外部査読・証明支援系による形式検証は行っていない。

\section{一つのSerre traceの持上げ}\label{RSRsec:trace}
\begin{lemma}[perfect複体に対する局所Nakayama]\label{RSRlem:nakayama}
$V\in\Perf(R)$が$V\otL_Rk=0$を満たせば$V=0$である。
$V\otL_Rk\simeq k[0]$なら$V\simeq R[0]$である。
\end{lemma}
\begin{proof}
有限自由加群からなる有界複体を取り、微分行列の単元成分を順次消去する。
残った最小複体の微分は全て$\mathfrak m$に値を取り、剰余上では零になる。
剰余が零なら全ての項が零であり、剰余が$k[0]$なら次数零の階数一の項だけが残る。
\end{proof}

\begin{lemma}[例外列と残余生成子]\label{RSRlem:exceptional}
$\OO_0,\ldots,\OO_{a-1}$は相対例外列である。
$\cC=\langle\cK,\OO_0,\ldots,\OO_{a-1}\rangle$を持ち、
\[
 p=\LL_{\OO_0}\circ\cdots\circ\LL_{\OO_{a-1}},\qquad G=p(E)
\]
は$\cK$の生成子である。$\Gamma_R=\REnd(G)$は$R$-perfectである。
\end{lemma}
\begin{proof}
単位写像$R\to\RHom(\OO_i,\OO_i)$のconeと、$j>i$について
$\RHom(\OO_j,\OO_i)$は$R$-perfectで、その剰余は零である。
特殊ファイバーでの消滅は三次超曲面の線束コホモロジーから従う。
補題\ref{RSRlem:nakayama}を適用する。
proper性により、各例外対象への左右のHomはperfectであり、評価・余評価は$\cC$の対象を保つ。
右端から左mutationを繰り返すと指定した右直交への射影$p$を得る。
$E$の生成性から$p(E)$の生成性が従い、proper性から$\Gamma_R$は$R$-perfectである。
\end{proof}

\begin{proposition}[持ち上げたSerre pairing]\label{RSRprop:trace}
ある射
\[
 \operatorname{tr}:\RHom_{\cC}(\OO_a,\OO[n])\longrightarrow R
\]
を選べて、任意の$M\in\cC$について合成とtraceが与える写像
\[
 \beta_M:\RHom_{\cC}(M,\OO[n])
 \longrightarrow\RHom_{\cC}(\OO_a,M)^\vee
 \tag{D}\label{RSReq:duality}
\]
は擬同型である。従って
\[
 M\in{}^\perp\OO\quad\Longleftrightarrow\quad M\in\OO_a^\perp.
 \tag{O}\label{RSReq:orthogonals}
\]
\end{proposition}
\begin{proof}
$L=\RHom(\OO_a,\OO[n])$は$R$-perfectである。
$\omega_X=\OO_X(-a)$より、$L\otL_Rk\simeq k[0]$であり、
非零射$L_k\to k$はスカラー倍を除いて$X$のSerre traceである。
補題\ref{RSRlem:nakayama}から$L\simeq R$なので、そのtraceを還元にもつ同型を一つ選ぶ。

合成$\RHom(M,\OO[n])\otL\RHom(\OO_a,M)\to L$とtraceから
\eqref{RSReq:duality}を作る。両辺が$R$-perfectなのでconeもperfectである。
導来底変更、合成、perfect双対は両立する。剰余上の$\beta_M$は
$S_X(\OO_X(a))=\OO_X[n]$に対するSerre dualityそのものなので擬同型である。
再び補題\ref{RSRlem:nakayama}を使う。
双対はperfect複体上保守的なので\eqref{RSReq:orthogonals}を得る。
\end{proof}

\begin{remark}
これは対角加群の持上げではない。底上perfectであることが既知のHom複体に対して
一つのtrace射を選び、そのconeにNakayamaを適用したものである。
$\Gamma_R$が包絡代数上perfectであるという主張は含まれない。
\end{remark}

\section{逆rotationの構成}\label{RSRsec:rotation}
$V\in\cC$についてmutationを次のfunctorial triangleで定める：
\begin{align}
 \RHom(\OO,V)\otL_R\OO&\longrightarrow V
       \longrightarrow\LL_{\OO}V,\label{RSReq:left}\tag{LM}\\
 \RR_{\OO}V&\longrightarrow V
       \longrightarrow\RHom(V,\OO)^\vee\otL_R\OO.
       \label{RSReq:right}\tag{RM}
\end{align}
proper性により全ての項は$\cC$に属する。

\begin{lemma}[直交圏間のmutation]\label{RSRlem:mutation}
$\LL_{\OO}$と$\RR_{\OO}$は互いに逆な同値
\[
 {}^\perp\OO\ \underset{\RR_{\OO}}{\stackrel{\LL_{\OO}}{\rightleftarrows}}\ \OO^\perp
\]
を与える。
\end{lemma}
\begin{proof}
\eqref{RSReq:left}に$\RHom(\OO,-)$を適用すると、評価により最初の写像は同型となる。
従って$\LL_{\OO}V\in\OO^\perp$である。
\eqref{RSReq:right}に$\RHom(-,\OO)$を適用し、perfect二重双対性を使うと
$\RR_{\OO}V\in{}^\perp\OO$となる。
両mutationは$\OO\otL_R\Perf(R)$を零に送る。
$V\in{}^\perp\OO$なら$\RR_{\OO}V\simeq V$なので、
\eqref{RSReq:left}へ$\RR_{\OO}$を適用して
$\RR_{\OO}\LL_{\OO}V\simeq V$を得る。
同様に$W\in\OO^\perp$について
$\LL_{\OO}\RR_{\OO}W\simeq W$である。
全ての射は評価・余評価から作るため、対象に自然である。
\end{proof}

\begin{proof}[定理\ref{RSRthm:rotation}の証明]
まず$M\in\cK$なら$1\le i<a$に対し
$\RHom(\OO_i,TM)=\RHom(\OO_{i-1},M)=0$である。
$\RHom(\OO_i,\OO)=0$なので、\eqref{RSReq:left}によるmutation後もこの消滅は保たれる。
$i=0$の消滅は左mutationの定義から従い、$\rho M\in\cK$を得る。
さらに
\[
 \RHom(\OO_a,TM)=\RHom(\OO_{a-1},M)=0
\]
と命題\ref{RSRprop:trace}から$TM\in{}^\perp\OO$である。

次に$N\in\cK$として$Q=\RR_{\OO}N$と置く。
$Q\in{}^\perp\OO$なので命題\ref{RSRprop:trace}から$\RHom(\OO_a,Q)=0$。
$1\le i<a$について\eqref{RSReq:right}に$\RHom(\OO_i,-)$を適用すると、
$N\in\cK$と$\RHom(\OO_i,\OO)=0$により$\RHom(\OO_i,Q)=0$を得る。
従って$0\le j<a$に対し
\[
 \RHom(\OO_j,T^{-1}Q)=\RHom(\OO_{j+1},Q)=0.
\]
これは$\psi=T^{-1}\RR_{\OO}$が$\cK$を保つことを示す。

補題\ref{RSRlem:mutation}より、$TM\in{}^\perp\OO$と$N\in\OO^\perp$に対して
\[
 \psi\rho M=T^{-1}\RR_{\OO}\LL_{\OO}TM\simeq M,
 \qquad
 \rho\psi N=\LL_{\OO}\RR_{\OO}N\simeq N.
\]
これらはdg自然同型であり、定理が従う。
\end{proof}

\begin{corollary}[底変更]\label{RSRcor:rotation-bc}
Homの導来底変更が成立し、properな模型と$T,\OO$が保たれる底変更について、
上のrotation、逆rotation、両合成の同型は底変更と両立する。
\end{corollary}
\begin{proof}
perfect双対、評価、余評価、有限cone、$T^{\pm1}$はこの底変更と両立する。
逆関手の同型はそれらのfunctorial triangleから作ったので、独立に同型を持ち上げる必要はない。
\end{proof}

\begin{remark}[周期性とは異なる]
本定理は$\rho^3\simeq[2]$を証明しない。
中心三次商のgraded singularity categoryでの$(3)\simeq[2]$と、射影残余圏上の
$\LL_{\OO}(1)$の比較にはwindowの同定が必要である。
特殊ファイバーで同型である二つの関手を、自然変換を構成せずに局所底上で同一視してはならない。
\end{remark}

\section{Serre関手と有限の開条件}\label{RSRsec:serre}
\begin{proposition}[局所Serre関手]\label{RSRprop:serre}
仮定\ref{RSRhyp:categorical}の$\cC$には相対Serre自己同値$S_{\cC}$が存在し、対象として
\[
 S_{\cC}(\OO_i)\simeq\OO_{i-a}[n]\qquad(i\in\mathbb Z)
\]
を満たす。$\cK$は$\cC$で両側admissibleであり、相対Serre自己同値を持つ。
\end{proposition}
\begin{proof}
\eqref{RSReq:duality}を$T$で移すと
$M\mapsto\RHom(\OO_i,M)^\vee$は$\OO_{i-a}[n]$で表現される。
この表現可能性を持つ第一変数の集合は、有限cone、shift、直和因子で閉じている。
$E$が生成子なので全ての第一変数で表現可能であり、dg YonedaによりSerre準関手を得る。
二重双対性からこれは完全忠実である。その像は生成子$T^{-a}E[n]$を含むので本質的全射である。

有限相対例外列の生成部分は左右随伴を持つ。
Serre関手は左直交を右直交へ移すので、その右直交$\cK$も両側admissibleとなる。
包含$i:\cK\to\cC$の右随伴を$i^!$とすれば、$i^!S_{\cC}i$がそのSerre関手である。
\end{proof}

\begin{remark}
対象$\OO_i$上の同型から、関手として$S_{\cC}=T^{-a}[n]$とは結論しない。
生成対象上の自己Homへの作用を比較する必要があり、Nakayama型のtwistを見落としてはならない。
\end{remark}

\begin{proposition}[代数的proper族の有限trace証明書]\label{RSRprop:open}
有限型noetherianスキーム$S$上に、properなdg圏の族$\cC_S$、$T$、$\OO$、
有限個のtwistからなる生成子$E$が実際に存在するとする。
点$s$で\eqref{RSReq:local-input}が成立すれば、$s$のZariski開近傍$U$上で
$\LL_{\OO}T$と$T^{-1}\RR_{\OO}$は残余圏の互いに逆な自己同値となる。
$U$は有限個のperfect複体の消滅条件で定められる。
\end{proposition}
\begin{proof}
有限個の単位写像のconeと逆順Homの消滅開集合で縮小し、相対例外列を得る。
$L=\RHom(\OO_a,\OO[n])$は$s$で階数一・振幅零なので、さらに縮小して
次数零の線束となる。$s$の近傍で自明化してtraceを選ぶ。
有限のperfect複体
\[
 W=\Cone\!\left(\RHom(E,\OO[n])
              \xrightarrow{\beta_E}\RHom(\OO_a,E)^\vee\right)
 \tag{FC}\label{RSReq:finite-certificate}
\]
は$s$で零である。そのコホモロジーのsupportを除くと$W=0$になる。
$\beta_M$が擬同型となる$M$の全体はthickなので、$E$の生成性から全ての$M$で成り立つ。
あとは第\ref{RSRsec:rotation}節の証明をそのまま適用できる。
\end{proof}

\begin{remark}
\eqref{RSReq:finite-certificate}は有限の圏論的証明書の\emph{形式}であり、
具体的$\Gamma_R$の行列を書き下して数値計算したという意味ではない。
この命題はproper族が有限型底上に既にある場合の主張であり、その族自身の降下を代替しない。
\end{remark}

\section{局所中心三次商からの入力}\label{RSRsec:projective}
ここでは本稿で使う既存局所模型を、具体的な代数的仮定まで遡って記す。
本節は親リポジトリ\cite{RSRRepo}の\texttt{local\_projective\_model.tex}と同じ範囲の結果であり、
本稿の新しいrotation論証と区別する。必要な推論を以下に再構成する。

\begin{hypothesis}[次数別flatな局所量子化]\label{RSRhyp:algebra}
$R$は有限型複素代数の局所化である正則局所環、$d=\dim R$、$N=n+2\ge5$とする。
$A=\bigoplus_{j\ge0}A_j$は中央$R$-代数、$A_0=R$で、各$A_j$は有限自由であり、
$A_k\simeq k[z_1,\ldots,z_N]$である。
$\Omega\in Z(A)_3$の剰余$F$は滑らかな三次超曲面を定める。
$B=A/(\Omega)$、$a=N-3$とする。
\end{hypothesis}

\begin{proposition}[局所proper模型の入力検証]\label{RSRprop:local-model}
仮定\ref{RSRhyp:algebra}の下で$A,B$は左右noetherianで、$\Omega$は中心正則元である。
$\cC_B=\thick_R\{\OO(j):j\in\mathbb Z\}\subset D^b(\qgr B)$は仮定\ref{RSRhyp:categorical}を満たす。
生成子は$E=\bigoplus_{i=0}^{N-1}\OO(i)$でよい。
また全ての$M\in D^b(\qgr B)$について$\RHom(E,M)$は$R$-perfectであり、$E$はこの大きい圏で弱生成子である。
\end{proposition}

\begin{proof}
\emph{(i) noetherian性と自由分解。}
一般に$D_j$が有限$R$-加群で$D/\mathfrak mD$が右noetherianなら、
$\operatorname{gr}_{\mathfrak m}D$は$(D/\mathfrak mD)[t_1,\ldots,t_e]$の商なので右noetherianである。
斉次右イデアル$I$の随伴次数付きイデアルの有限個の二重斉次生成元を持ち上げ、
それらの生成するイデアルを$J$とする。最低$\mathfrak m$次数を順に消すと
$I_j\subset J_j+\mathfrak m^qD_j$を全ての$q$で得る。
有限$R$-加群$D_j/J_j$にKrull交叉を適用して$I_j=J_j$となる。
任意の右イデアルは内部次数の最高項を用いて同様に有限生成となる。左側も同じである。

$D,M$が次数別有限自由で、剰余上の$M_k$が有限階数自由加群による有限graded分解を持つとする。
第零項の生成元の像を持ち上げると、次数別Nakayamaから全射となる。
その核は各次数で分裂直和因子なので有限自由であり、剰余上の核も正しく得られる。
これを反復すると同じshiftと階数の分解を持ち上げられる。最後の核は次数別Nakayamaで零である。
この操作を多項式環のKoszul分解と対角Koszul分解に適用し、
$R$の$A$-分解および$A$の$A^e$-分解を得る。項はそれぞれ
$A(-j)^{\binom Nj}$と$A^e(-j)^{\binom Nj}$、$0\le j\le N$である。
全分解は次数別$R$-splitなので任意の通常の底変更で完全性を保つ。
対角分解を加群とテンソルすれば$\operatorname{gldim}A\le N+d$となる。
次数一生成と二次関係も、剰余上からの同じ持上げで従う。

\emph{(ii) Gorenstein性と中心商。}
Koszul分解を$A$で双対化すると、各内部次数は有界有限自由$R$-複体である。
剰余上の完全性と補題\ref{RSRlem:nakayama}から、内部次数$-N$、cohomological degree $N$以外は零となる。
この唯一の箇所では$R$が残る。従って左右で
\[
 \operatorname{R\underline{Hom}}_A(R,A)\simeq R(N)[-N].
\]
$\Omega:A_{j-3}\to A_j$の剰余は単射で、列数大の小行列式に単元がある。
従って各次数で分裂単射であり、$\Omega$は正則、$B_j$は有限自由となる。
長さ一の$A$-分解からchange-of-rings
\[
 \operatorname{R\underline{Hom}}_A(M,A)
 \simeq\operatorname{R\underline{Hom}}_B(M,B)(3)[-1]
\]
を得る。$M=R$を代入すると
$\operatorname{R\underline{Hom}}_B(R,B)=R(a)[-(N-1)]$である。
$A$の有限大域次元から$B$の左右injective dimensionは高々$N+d-1$となる。
$B$は商なのでnoetherianである。

\emph{(iii) 線束Homとその底変更。}
Burke--Stevenson\cite[Lemmas 6.1--6.2]{RSRBS}により$B$は$\chi$と有限cohomological dimensionを満たす。
$M_\ell=B/B_{\ge\ell}$を内部次数で濾過すると、その商は$B_j\otimes_RR(-j)$である。
上のGorenstein型から、これらの$B$へのExtは次数$N-1$、内部次数$-a-j$にのみ存在する。
各固定内部次数$t$では$j=-t-a$という唯一の層が寄与し、$\ell>j$で安定する。
従って局所コホモロジーは
\[
 H^q(R\tau B)_t=
 \begin{cases}B_{-t-a}^\vee&q=N-1,\\0&q\ne N-1.\end{cases}
\]
三角$R\tau B\to B\to R\omega\pi B$から
\[
 \operatorname{Ext}^p(\OO(i),\OO(j))=
 \begin{cases}
 B_{j-i}&p=0,\\ B_{i-j-a}^\vee&p=n,\\0&\text{その他}
 \end{cases}
 \tag{H}\label{RSReq:line-hom}
\]
となる。$a>0$なので二群が同じ差で同時に非零となることはない。

底変更はコホモロジー階数の一致だけでなく、自然な導来Hom写像について確認する。
$M_\ell$の有限階数graded自由分解を取り、各内部次数で$R$-splitにする。
$\underline{\operatorname{Hom}}_B(P_\bullet,B)_t$は下に有界な有限自由複体で、
そのコホモロジーは一つの次数の有限自由加群である。
$R$は有限大域次元なので、完全な射影複体のcycleにdimension shiftingを適用すると
完全部分は分裂する。従って通常のテンソルが導来テンソルを計算する。
項ごとのHom底変更と、各$t$で安定するfiltered colimitを使えば、線束Homの自然な底変更が従う。
この議論は$R'\!$が正則noetherian・有限大域次元で$A\otimes_RR'$がnoetherianである
通常の底変更に適用され、特に剰余体と分数体に適用される。

\emph{(iv) 有限生成と特殊ファイバー。}
Burke--Stevenson\cite[Theorem 6.4]{RSRBS}を$A$に適用すると、$\Dsg(\gr A)=0$より
任意の連続した$N$個のtwistが$D^b(\qgr A)$を生成する。
$A\to B$は有限Tor振幅を持ち、導来制限は$\OO_A(j)$を$\OO_B(j)$へ送る。
よって全twistは$E$のthick閉包に入る。
\eqref{RSReq:line-hom}から$\REnd(E)$は$R$-perfectで、有限cone・直和因子によって$\cC_B$全体がproperとなる。
自然なHom底変更は積と両立し、有限cone・冪等完備化へ延長できる。
特殊ファイバーで全twistは$\Perf(X)$を生成するため、\eqref{RSReq:local-input}を得る。

\emph{(v) 全coherent圏へのHom。}
\cite[Proposition 4.18]{RSRBS}の有界導来sectionは$D^b(\qgr B)$を$D^b(\gr_{\ge i}B)$へ送る。
各内部次数の成分は有限$R$-加群である。従って固定した$j$について
$\RHom(\OO(j),M)$は有界かつ有限生成コホモロジーを持つ。
$R$の正則性よりこれは$R$-perfectである。
$\RHom(E,M)=0$なら全twistからのHomが零であり、導来sectionの各次数が零になる。
同命題の完全忠実性から$M=0$となる。
\end{proof}

\begin{corollary}[局所量子化のrotation]\label{RSRcor:local-rotation}
仮定\ref{RSRhyp:algebra}の残余部分圏$\cK_B$上で
$\rho=\LL_{\OO}(1)$は自己同値であり、$\rho^{-1}=(-1)\RR_{\OO}$である。
残余smoothnessは仮定しない。
\end{corollary}

\section{smoothnessと全coherent残余圏}\label{RSRsec:fullness}
\begin{lemma}[smooth代数上のpseudo-perfect加群]\label{RSRlem:pseudoperfect}
$H$は可換環$R$上smoothなdg代数とする。
$H$-加群$V$の基礎$R$複体がperfectなら、$V$は$H$-perfectである。
\end{lemma}
\begin{proof}
対角$H$は$H^e$の有限cone・shift・直和因子から作られる。
これに$V\otL_H-$を作用させると、自由双加群の像は$V\otL_RH$である。
$V$が$R$-perfectなのでこれは$H$-perfectであり、その有限操作で$V$が得られる。
これは\cite[Lemma 2.8(2)]{RSRTV}、\cite[F6]{RSRToen}の性質の直接の証明である。
\end{proof}

\begin{lemma}[弱生成子による全体の回収]\label{RSRlem:fullness}
$\cD$を冪等完備dg圏、$E\in\cD$を弱生成子とし、$H=\REnd(E)$と置く。
全ての$M\in\cD$について$\RHom(E,M)$が$R$-perfectで、$H$が$R$上smoothなら
$\cD=\thick_R(E)$である。
\end{lemma}
\begin{proof}
$V=\RHom(E,M)$は補題\ref{RSRlem:pseudoperfect}により$H$-perfectである。
$\thick(E)\simeq\Perf(H)$から、それに対応する$N\in\thick(E)$を選ぶ。
任意の$N\in\thick(E)$について、自然な写像
\[
 \RHom_{\cD}(N,M)\longrightarrow
 \RHom_H(\RHom(E,N),\RHom(E,M))
\]
は擬同型である。これは$N=E$のdg Yonedaから有限操作で延長すればよい。
従って選んだ加群同型を射$N\to M$へ戻せる。
そのconeは$E$からの全Homが零なので、弱生成性により零となる。
\end{proof}

\begin{proposition}[相対gluingのsmoothness]\label{RSRprop:gluing}
$R$-properで有限生成なdg圏$\cC$が、両側admissibleな半直交分解
$\cC=\langle\cK,\OO_0,\ldots,\OO_{a-1}\rangle$を持つとする。
$\OO_i$が相対例外対象なら、$\cC$が$R$上smoothであることと$\cK$が$R$上smoothであることは同値である。
\end{proposition}
\begin{proof}
まず二成分のgluingを考える。生成子を取り、零である方向のHom複体を零に置き換えると、
次の三角dg代数にMorita同値な模型を得る：
\[
 H=\begin{pmatrix}A&M\\0&D\end{pmatrix},\qquad e+f=1.
\]
$A,D$がsmoothで$M$が$R$-perfectなら、$A\otL_RD^{\mathrm{op}}$もsmoothなので
補題\ref{RSRlem:pseudoperfect}から$M$は$A$-$D$双加群としてperfectである。
次の対角の三角を用いる：
\[
 He\otL_A M\otL_D fH\longrightarrow
 (He\otL_A eH)\oplus(Hf\otL_D fH)
 \longrightarrow H.
 \tag{G}\label{RSReq:gluing-diagonal}
\]
最初の射は$m$を左右どちら側に掛けるかの差、二番目は乗法である。
homotopy-flatな模型を取り各$e,f$成分を見れば、対角成分は恒等射、
交差成分は$M\to M\oplus M\to M$のsplit完全列であり、この三角が確認できる。

$A$の対角の有限perfect表示を両側から誘導すれば$He\otL_A eH$は$H^e$-perfectとなる。
$D$の項も同様である。$M$のperfect性から第一項も$H^e$-perfectなので、$H$の対角はperfectとなる。
各例外成分$\Perf(R)$を一つずつ加える帰納法により、$\cK$のsmoothnessから$\cC$のsmoothnessを得る。
この議論は体を仮定せず、Orlov\cite[Theorem 3.25]{RSROrlov}の必要な方向を相対化したものである。

逆に$\cC$がsmoothなら、その有限perfect対角を$\cK$の両変数へ制限する。
左右随伴により、制限された表現可能加群は再び表現可能である。
したがって制限はperfect性を保ち、$\cK$の対角もperfectとなる。
\end{proof}

\begin{proof}[定理\ref{RSRthm:summary}(1)の証明]
命題\ref{RSRprop:serre}により$\cK_B$は両側admissibleである。
\eqref{RSReq:SM}と命題\ref{RSRprop:gluing}から$\cC_B$がsmoothとなる。
命題\ref{RSRprop:local-model}の最後の主張と補題\ref{RSRlem:fullness}により
$\cC_B=D^b(\qgr B)$となる。
同じ例外列の右直交を取れば、線束生成残余圏と全coherent残余圏が一致する。
既存の有界Orlov分解との同定も、この右直交を用いるwindowで得られる。
ただし任意の導来底変更に関するOrlov核のcoherenceまで主張しているのではない。
\end{proof}

\section{smooth proper pairの有限型底への降下}\label{RSRsec:spread}
\begin{proposition}[局所pairのspread]\label{RSRprop:spread}
$R=(R_0)_{\mathfrak m}$を有限型複素整域の局所化とする。
$\Gamma_R$が$R$上smooth properで、$\rho$が$\Perf(\Gamma_R)$の$R$-linear dg自己同値なら、
ある$f\in R_0\setminus\mathfrak m$と$U=\operatorname{Spec}R_0[f^{-1}]$上の
smooth proper dg代数$\Gamma_U$、可逆perfect双加群$P_U$が存在し、
局所化後に$(\Perf(\Gamma_R),\rho)$を与える。
\end{proposition}
\begin{proof}
$R=\varinjlim_{f\notin\mathfrak m}R_0[f^{-1}]$である。
Toën\cite[Théorème 1]{RSRToen}により$\Gamma_R$は、ある段階のsmooth proper dg代数へ降下する。
$\rho$とその逆関手を表すMorita双加群を$P,Q$とする。
その基礎$R$複体は、生成子とその像のHomで計算できるのでperfectである。
$\Gamma_R^e$もsmoothであるから、補題\ref{RSRlem:pseudoperfect}により
$P,Q\in\Perf(\Gamma_R^e)$である。

perfect加群のfiltered colimit性\cite[F4]{RSRToen}\cite[Lemma 2.10]{RSRTV}により
$P,Q$と逆同型
\[
 P\otL_{\Gamma_R}Q\simeq\Gamma_R,
 \qquad Q\otL_{\Gamma_R}P\simeq\Gamma_R
\]
を有限の段階へ降下できる。射、逆射、および有限個の合成が恒等となる関係も、
同じfiltered colimit性を使ってさらに一段階先で成立させられる。
従って降下した$P_U$は可逆なperfect双加群である。
元の$\rho$との同定は$R$上で保たれる。
\end{proof}

\begin{proof}[定理\ref{RSRthm:summary}(2)の証明]
系\ref{RSRcor:local-rotation}により$\rho$は実在するdg自己同値である。
\eqref{RSReq:SM}と既存のproper性の下で命題\ref{RSRprop:spread}を適用する。
\end{proof}

\begin{remark}
この証明は、$R$-perfectなdg代数を有限階数のstrict dg代数と同一視しない。
個々の高次$A_\infty$積の分母を無限に列挙して一つの開集合を得ようとする議論も不要である。
smoothnessが与えるhomotopical finite presentationが有限型降下を可能にする。
\end{remark}

étale局所実現は第\ref{er:sec:etale}節の証明を用いる。
Hochschild計算は第\ref{ho:sec:parity}節を用いる。

\section{generic可換復元への正確な還元}\label{RSRsec:reconstruction}
\begin{theorem}[局所saturationからgeneric復元へ]\label{RSRthm:generic}
$S$を滑らかな既約有限型複素スキーム、$s\in S(\CC)$、$R=\OO_{S,s}$とする。
$R$上に仮定\ref{RSRhyp:algebra}の量子化があり、残余対角条件\eqref{RSReq:SM}が成立するとする。
$K=\CC(S)$と置くと、ある有限分離拡大$L/K$、滑らかな三次超曲面$Y_L$、
rotationを保つenhanced同値
\[
 (\cK_{B_L}^{\mathrm{coh}},\LL_{\OO}(1))
       \simeq (\Ku(Y_L),\rho_{Y_L})
 \tag{CR}\label{RSReq:reconstruction}
\]
が存在する。ここで左辺は全coherent残余圏である。
\end{theorem}
\begin{proof}
局所proper模型と定理\ref{RSRthm:rotation}から$(\Perf(\Gamma_R),\rho)$を得る。
\eqref{RSReq:SM}によりsmooth properとなるので、命題\ref{RSRprop:spread}から$s$を含む開集合$U\subset S$へ
pairを降下する。その$s$-ファイバーは元の$(\Ku(X),\rho_X)$である。
定理\ref{er:thm:ER}を適用してpointed étale近傍$V\to U$と$Y/V$を得る。
指定点を通る既約成分を選ぶと、その像は$U$の非空開集合を含む。
étale射は準有限なので、その関数体$L$は$K$の有限分離拡大である。

線束Homの自然な分数体底変更と有限mutationにより、
$\Gamma_R\otL_RL$は実際の$B_L$の線束生成残余圏を表す。
smoothnessは底変更で保存されるので、第\ref{RSRsec:fullness}節を$L$上で適用して
線束生成部分と全coherent残余圏を同定する。
従ってgenericファイバーで\eqref{RSReq:reconstruction}を得る。
\end{proof}

\begin{corollary}[Poisson成分に対する適用範囲]\label{RSRcor:component}
Poisson既約成分$Z$を支配する滑らかな既約chart $S$と可換点$s$があり、
$\CC(S)/\CC(Z)$が有限分離であるとする。
$s$の局所環上で仮定\ref{RSRhyp:algebra}と\eqref{RSReq:SM}を満たせば、
$Z$の関数体の有限分離拡大上に\eqref{RSReq:reconstruction}の可換模型が存在する。
\end{corollary}
\begin{proof}
定理\ref{RSRthm:generic}を適用し、二つの有限分離拡大を合成する。
\end{proof}

\begin{remark}[残る仮定の量化]
各Poisson成分に一つ、この局所入力を実現すればよい。
全成分での量子化自体が本稿から生じるわけではない。
定理\ref{er:thm:ER}は自己同値一つを指定したpairについての定理である。
可換模型から周期性同型を移すことはできるが、特定のmatrix factorization表示の
周期性とのcoherenceまで自動的に同定したとは扱わない。
\end{remark}

\section{非楕円型への適用と未解決の核心}\label{RSRsec:remaining}
親プロジェクトの二つの非楕円型成分では、可換点の局所環は
$R=\CC[r,s]_{(r,s)}$、$N=7$、$n=5$、$a=4$である。
その全次数flatなambientと中心三次商の構成を入力すれば、系\ref{RSRcor:local-rotation}により
\[
 \rho^{-1}=(-1)\RR_{\OO}
\]
が局所残余圏上に存在する。節点型・尖点型の純三次ブロックの滑らかさは使わない。
具体的な係数と中心式は研究ノートに記す。

この二成分に対する残る核心は、$\Gamma_R$の対角が$\Gamma_R^e$上perfectであることにある。
これが得られれば、全coherent比較、pairのspread、定理\ref{er:thm:ER}により
一般点のrotation付き可換模型が有限分離拡大上に存在する。
本稿はこの対角perfect性を証明したものではない。

\subsection{非楕円二成分に対する具体的な可換変形slice}
変数の順を$(x,y,z,a,b,c,u)$とし、
\[
 F_\epsilon=f_\epsilon+ux^2+a^3+b^3+c^3+u(a^2+b^2+c^2)+u^3,
 \quad f_C=xz^2+y^3,\quad f_N=xyz+y^3+z^3
\]
と置く。これらの滑らかさは研究ノートの勾配消去で確認されている。
\begin{proposition}[35係数の横断slice：厳密線形代数の証明書]\label{RSRer:prop:slice}
次の単項式集合はそれぞれ$J(F_C)_3$と$J(F_N)_3$の基底を与える。
\begin{align*}
\mathcal M_C=\{&axy,bxy,cxy,abx,acx,bcx,yz^2,ayz,byz,cyz,uyz,\\
 &aby,acy,auy,bcy,buy,cuy,u^2y,az^2,bz^2,cz^2,uz^2,\\
 &abz,acz,bcz,buz,cuz,u^2z,abc,abu,acu,bcu,bu^2,cu^2,u^3\},\\
\mathcal M_N=\{&abx,acx,bcx,ay^2,by^2,cy^2,ayz,byz,cyz,uyz,\\
 &aby,acy,bcy,buy,cuy,u^2y,az^2,bz^2,cz^2,uz^2,abz,acz,\\
 &auz,bcz,buz,cuz,u^2z,abc,abu,acu,au^2,bcu,bu^2,cu^2,u^3\}.
\end{align*}
従って$V(F_\epsilon+\sum_{m\in\mathcal M_\epsilon}t_m m)$の滑らかな開部分は、
原点で定理\ref{er:thm:criterion}の幾何的接写像が同型となる35次元のsource族である。
\end{proposition}
\begin{proof}
84個の三次単項式を、変数番号の非減少三つ組の辞書式順で並べる。
84行49列の整数行列$M_\epsilon$の列を$z_j\partial_iF_\epsilon$の係数ベクトルとする。
列は$i$を先、$j$を後に増加させて並べる。
上記$\mathcal M_\epsilon$に対応する35行を除いた49行の小行列を$D_\epsilon$とすると、
整数のfraction-free消去により
\[
 \det D_C=2^{15}3^{26},\qquad \det D_N=-2^6 3^{25}
\]
を得る。独立な二つの決定的アルゴリズムによる検算を
\texttt{computations/verify\_etale\_inputs.py}に収める。
これは有限体の標本検査ではなく、上で完全に指定した整数行列の行列式計算である。
残りの行を$U$、採用した行を$P$とすると、
$v\mapsto v_U-(M_\epsilon)_U D_\epsilon^{-1}v_P$は$M_\epsilon$を零へ送り、
$\mathcal M_\epsilon$の単項式上では恒等となる。この二恒等式も有理数上で検算する。
よって84次元の三次形式空間は、49次元の座標変換方向と表示した35次元の空間の直和である。
Euler等式により$F_\epsilon$自身も前者に入る。補題\ref{er:lem:cubic}で接写像の主張を得る。
\end{proof}

\begin{remark}
この基底は局所実現のsourceを具体化するものであり、量子化パラメータ$(r,s)$から
$t_m$への閉公式を計算したものではない。各小拡大の実際の障害類が得られれば、
その消去をこの35係数で行える。smoothness証明書を代替する計算でもない。
\end{remark}

\begin{proposition}[Serre自己同値はsmoothnessを含意しない]\label{RSRprop:warning}
体$k$上の$H=k[\epsilon]/(\epsilon^2)$に対し、$\Perf(H)$はproperでSerre関手が恒等であるがsmoothではない。
\end{proposition}
\begin{proof}
$\epsilon$の係数を取るtraceは非退化な対称Frobenius pairingを与え、$H^\vee\simeq H$は双加群同型である。
従ってperfect加群上のSerre関手は$-\otL_HH^\vee\simeq\mathrm{id}$である。
一方、単純加群$k$は
$\cdots\xrightarrow{\epsilon}H\xrightarrow{\epsilon}H\to k\to0$
という無限最小分解を持ち、$H$-perfectでない。
$H$がsmoothなら補題\ref{RSRlem:pseudoperfect}に反する。
\end{proof}

この例は局所量子化が非smoothであることを示す反例ではない。
今回構成したSerre関手・逆rotationだけでは対角perfectnessの証明にならないことを示している。
対角の持上げ、形式的Artin切断、底上perfectなHom複体という三種類の有限性を区別する必要がある。


% Rotation orbit generation and generic descent: alternative proof of the generator statement.
\section{次数シフト軌道による一般ファイバーの可換復元}
\label{rog:sec:generic}

完備底全体でcompactな逆核を先に構成することは、一般点での復元には過剰な十分条件である。
本節では、既存の次数別完備な有限rank MF同値を\emph{固定した三対象}へ制限し、
先にパラメータを可逆にする。その像の非零性をEuler形式で検出し、
孤立超曲面の特異圏の分類から三つの次数シフトが生成することを示す。
逆核のcompact性は、この体上の生成性から得る。
従って全Artinの生成性を完備底へ持ち上げる誤った推論は使わない。

既知の入力は、Stevensonの特異圏のthick部分圏分類
\cite[Theorem~6.13]{ROGStevenson}、行列因子分解と特異圏の比較および
残余体stabilizationのendomorphism計算\cite[\S2.4, \S5.6]{ROGDyckerhoff}、
compact対象の局所化原理\cite[Proposition~13.37.6]{ROGStacks}、
perfect双加群の有限表示性\cite[Lemmas~2.8, 2.10]{ROGTV}である。
これらは既知の原理であり、非可換中心三次対への適用に必要な
局所window比較(L)と形式MF比較(CF)は、
それぞれ定理\ref{tw:thm:window}と定理\ref{cft:thm:complete}で構成する。
具体的な可換方程式を指定する場合には、別途(D)を要する。

\subsection{孤立斉次超曲面の次数シフト軌道}

$K$を標数零の体、$S=K[x_1,\ldots,x_N]$、$W\in S_d$、$d\ge2$とする。
$S/(W)$の特異点集合は原点だけであると仮定する。
有限生成射影加群上のgraded MFのホモトピー圏を冪等完備化して$\mathcal T^{\rm gr}$と書く。
内部shiftの規約は$M(r)_q=M_{q+r}$、周期性は$(d)=[2]$とする。

\begin{theorem}[非零対象のシフト軌道は生成する]\label{rog:thm:orbit}
任意の非零対象$E\in\mathcal T^{\rm gr}$について
\[
 \operatorname{thick}\{E(j):j\in\mathbb Z\}=\mathcal T^{\rm gr}.
 \tag{RO1}\label{rog:eq:orbit}
\]
従って$E\oplus E(1)\oplus\cdots\oplus E(d-1)$は生成対象である。
単独の$E$がgraded圏を生成するとは主張しない。
\end{theorem}
\begin{proof}
次数を忘れる関手を$U$とする。$U$は有限rank対象で保守的である。
実際、$U(E)$の恒等射が非斉次の多項式ホモトピーの境界なら、
その行列の各成分を内部次数に分解し、恒等射に寄与する次数だけを取ると
gradedな縮約を得る。冪等因子についても、対応する閉冪等射の境界に同じ操作を行う。
従って$U(E)\ne0$である。

$S/(W)$の特異点集合は一点なので、\cite[Theorem~6.13]{ROGStevenson}により
次数を忘れた特異圏の冪等完備化には非零の真のthick部分圏がない。
正則環$S$上のMFと特異圏との比較から、任意の有限graded MF $Z$について
$U(Z)\in\operatorname{thick}(U(E))$となる。
この事実だけでgraded生成性を宣言せず、次の次数の再付与を使う。

非次数付き$S$-加群$M$に対し、
\[
 (\mathcal R M)_j=M\quad(j\in\mathbb Z),\qquad
 a_\ell:M_j\longrightarrow M_{j+\ell}
\]
と定め、後者は元の$a_\ell$の作用とする。
非次数付きMFの第一微分はコピーの添字を保ち、第二微分は添字を$d$上げるように定める。
すると$\mathcal R$は、任意rankのgraded projective MFのホモトピー圏への完全関手となる。
非次数付き射はコピーの添字を保つ射へ送り、ホモトピーの二つの成分には
内部次数$-d,0$を与える。これにより境界の式を保つ。
coneとの両立は、coneの第二成分にあるコピーの添字を$d$ずらして確認できる。
また自由加群は、全shiftを直和したgraded自由加群へ移り、射影加群はその直和因子へ移る。

有限graded MFについて自然な同型
\[
 \mathcal R U(E)\simeq\bigoplus_{j\in\mathbb Z}E(j)
 \tag{RO2}\label{rog:eq:regrading}
\]
がある。元の内部次数とコピーの添字の差を固定した成分が、右辺の各直和成分である。
第一微分は両次数を保ち、第二微分は両次数を$d$上げるので、この分解はMFの分解である。
同じことは$Z$についても成り立つ。
従って$\mathcal R$を有限thick構成へ適用すると、
$Z$は$\{E(j)\}$の生成する局所化部分圏に属する。

有限rank graded MFから直和への射は有限個の行列係数で決まり、有限個の直和成分を通る。
従ってこれらは大きなgraded MFホモトピー圏のcompact対象である。
compact対象の局所化原理\cite[\S13.37]{ROGStacks}により、
局所化部分圏に属するcompact対象はそのcompact生成対象のthick閉包に属する。
この原理は$\{E(j)\}$が生成するdg局所化部分圏に適用すればよく、
大きなMF圏全体の生成性を別に仮定する必要はない。
実際、この局所化部分圏で全$E(j)$と右直交する$Y$に対して、
$\operatorname{Hom}(-,Y)$が直和を直積へ送ることから$Y$自身も$Y$と直交し、$Y=0$となる。
よって\eqref{rog:eq:orbit}を得る。最後は$(d)=[2]$から従う。
\end{proof}

\subsection{残余体stabilizationのEuler形式}

今度は$d=3$とし、$F\in k[x_1,\ldots,x_N]_3$を滑らかな射影cubicを定める三次式とする。
$K_F$を原点の残余体のstabilizationとする。
\begin{proposition}[非零のEulerベクトル]\label{rog:prop:euler}
$r=0,1,2$について
\[
 \chi_r:=\sum_p(-1)^p\dim_k\operatorname{Hom}(K_F,K_F(r)[p])
 =\sum_{\substack{0\le s\le N\\s\equiv-r\pmod3}}(-1)^s\binom Ns.
 \tag{RO3}\label{rog:eq:euler}
\]
この三成分ベクトルは零でない。より正確に
\[
 \chi_0+\chi_1+\chi_2=0,\qquad
 \chi_0^2+\chi_1^2+\chi_2^2=2\cdot3^{N-1}.
 \tag{RO4}\label{rog:eq:norm}
\]
$N=5$では$(\chi_0,\chi_1,\chi_2)=(-9,9,0)$である。
\end{proposition}
\begin{proof}
$F=\sum_i x_if_i$、$f_i=(\partial_iF)/3$、
$f_i=\sum_jx_ja_{ij}$、$a_{ij}=(\partial_jf_i)/2$と置く。
外積作用$e_i\wedge-$の二重次数を$(-1,1)$、
収縮$\iota_i$を$(1,-1)$、周期変数$t$を$(2,-3)$とする。
$K_F$の微分は
\[
 Q=\sum_i x_i\iota_i+t\sum_i f_i(e_i\wedge-),\qquad Q^2=tF.
\]
そのendomorphism代数のsuper交換子微分について
\[
 d(e_i\wedge-)=x_i,\qquad
 \xi_i:=\iota_i-t\sum_j a_{ij}(e_j\wedge-),\qquad d\xi_i=0.
\]
順序を固定した$e_J\xi_I$はendomorphism加群の基底である。
実際、標準の$e_J\iota_I$からの変換は三角的である。
この基底では微分は$x_1,\ldots,x_N$のKoszul微分となる。
また$\xi_i\xi_j+\xi_j\xi_i=-t(a_{ij}+a_{ji})$であり、右辺は原点で零となる一次式なので境界である。
従って、周期変数を残した二重次数付きコホモロジーは
\[
 \bigwedge(\xi_1,\ldots,\xi_N)\otimes_k k[t,t^{-1}],
 \qquad\deg\xi_i=(1,-1)
 \tag{RO5}\label{rog:eq:exterior}
\]
となる。これは\cite[\S5.6]{ROGDyckerhoff}のClifford計算の三次の場合でもある。

$\xi_I t^j$が内部shift $r$のHomに寄与する条件は
$-|I|-3j=r$であり、cohomological次数は$|I|+2j$である。
$0\le|I|\le N$なので各$r$について有限個しか寄与せず、符号は$(-1)^{|I|}$。
これで\eqref{rog:eq:euler}を得る。

$A(z)=(1-z)^N$を$z^3-1$で割った余りの係数が、順序を除いて$\chi_r$である。
原始三乗根$\zeta$において$A(\zeta)=(1-\zeta)^N\ne0$なので三成分は同時には零でない。
$A(1)=0$、$|1-\zeta|^2=3$と三点のFourier公式から\eqref{rog:eq:norm}を得る。
これは整数恒等式なので、基礎体に依らない。$N=5$の値は直接代入する。
\end{proof}

\begin{remark}
一般$N$で自己Euler数$\chi_0$だけを使うのは誤りである。
例えば$N=9$では$\chi_0=0$となるが、三成分ベクトルは非零である。
また上の群は\emph{生成対象間のHom}であり、$\operatorname{Nat}(\mathrm{Id},\rho^r)$ではない。
\end{remark}

\subsection{固定した有限対象上では形式functorを先に局所化できる}

$R=\mathbb C[h]_{(h)}$、$K=\mathbb C(h)$、
$\widehat R=\mathbb C[[h]]$、$L=\mathbb C((h))$とする。
代数的中心三次対$(A_R,\Omega_R)$について、次の既存入力を使用する。
\begin{description}
\item[\textup{(L)}] 次数別flatな局所量子化と有限Koszul分解、
線束生成残余圏のproper模型、regular底のstabilization/window比較がある。
原点の$R$-加群の有限rank stabilizationを$K_R$とすると
\[
 T_R=K_R\oplus K_R(1)\oplus K_R(2),\quad
 \Gamma_R=R\operatorname{End}(T_R),\quad
 \mathcal K^{\rm lin}_{B_R}\simeq\operatorname{Perf}(\Gamma_R),
 \qquad \Gamma_R\in\operatorname{Perf}(R).
 \tag{RO6}\label{rog:eq:source}
\]
この模型とHomは$\widehat R,L$への平坦底変更に両立する。
\item[\textup{(CF)}] 次数別完備底上の有限rank graded MFに、
内部shiftと指定周期性を保つ実際の$\widehat R$-線形$A_\infty$同値
\[
 \Phi:\operatorname{MF}^{\mathbb Z}(A^{\rm gr}_{\widehat R},\Omega)
 \simeq\operatorname{MF}^{\mathbb Z}(\widehat R[x],W_h)
 \tag{RO7}\label{rog:eq:completed}
\]
がある。$W_h\bmod h=F$は滑らかで、$\Phi$は剰余で恒等。
単なる両立したArtin同値の列ではなく、
定理\ref{cft:thm:complete}で構成した固定内部次数内で収束する
実際の関手を用いる。
\end{description}
記号(CF)は完備底の実際のfunctorを指し、先行節で(F)と呼んだArtin残余同値の列と区別する。
本節では、これらの比較定理を入力として、
固定三対象の局所化から一般点の生成性へ進む部分を独立に記述する。

\begin{remark}[\textup{(L)}の三つのstabilization]
三つの対象を選ぶのは、特殊生成性からの持上げではない。
$a=N-3$とすると、windowは線束$\mathcal O(a),\mathcal O(a+1),\mathcal O(a+2)$を
有限長切断$B(j)_{<1-a}$のstabilizationのshiftへ送る。
そのfiltrationの商は$B_r\otimes_RR(j-r)$（$0\le r\le j-a$）である。
従って三角的に$K_R(a),K_R(a+1),K_R(a+2)$と同じthick部分圏を生成する。
三次周期性でこれを\eqref{rog:eq:source}の三つの対象に換えられる。
proper性は線束生成模型のHom公式から継承する。
この主張は全coherent残余圏と線束生成部分の同一性を含まない。
\end{remark}

\begin{lemma}[有限対象への制限とLaurent局所化]\label{rog:lem:localization}
\textup{(L),(CF)}の下で$E=\Phi(K_{\widehat R})$と置くと、
$E$は通常の$\widehat R[x]$上の有限rank graded MFであり、
\[
 \operatorname{Perf}(\Gamma_R\otimes_RL)
 \simeq\operatorname{thick}\{E_L,E_L(1),E_L(2)\}
 \tag{RO8}\label{rog:eq:image}
\]
というrotationと周期性を保つ同値がある。
また$\operatorname{Hom}(E,E(r))$は$r=0,1,2$で$\widehat R$-perfectであり、
その特殊・一般ファイバー間のEuler数は\eqref{rog:eq:euler}に等しい。
\end{lemma}
\begin{proof}
固定した有限rankの基礎加群と固定cohomological次数では、MF行列の各成分の
内部次数と$t$の指数は固定される。従って各成分は有限自由な
$\widehat R$-加群$\widehat R[x]_j$に属する。
(CF)の$F_{2q+1}=O(h^q)$によるtwistingの和は、その有限自由成分内で収束する。
結果は$x$の次数を無制限に増やす級数ではなく、$h$の形式係数を持つ有限次数の多項式行列である。
固定個数の射に対する各Taylor成分も同じ理由で通常の$\widehat R$-多重線形写像となる。
恒等式と逆functorは全Artin剰余で成立し、成分の分離性によって$\widehat R$上で成立する。

従って$T_{\widehat R}$の有限対象部分圏では、(CF)は\emph{通常の}endomorphism複体の
$A_\infty$準同型を与える。$\widehat R\to L$は平坦なので、これをテンソルすると
通常の準同型である。固定cohomological次数のHomの形成はこの平坦底変更と可換する。
$\Phi$を任意のLaurent係数対象へ新たに拡張することは使わず、
この固定部分圏のpretriangulated・冪等包絡を取って\eqref{rog:eq:image}を得る。
有限twisted complexのtwisting行列は厳密上三角に取れるため、包絡への延長では有限個のTaylor成分だけが使われる。
したがって$L$への移行後に新しい無限級数の収束を要求しない。

Homのperfect性は(L)と(CF)から従い、$K_R(r)$は周期性により$T_R$のshiftの直和因子である。
有界有限射影複体のEuler数は特殊化で不変なので、剰余の$K_F$について計算した
\eqref{rog:eq:euler}が一般ファイバーでも保たれる。
内部shiftと周期性の自然比較も、この固定部分圏への制限と平坦底変更で保たれる。
\end{proof}

\begin{theorem}[Laurent体上の生成性と可換復元]\label{rog:thm:laurent}
\textup{(L),(CF)}の下で、$N\ge5$に対し
\[
 (\mathcal K^{\rm lin}_{B_R}\otimes_RL,\rho)
 \simeq\bigl(\operatorname{Ku}(V(W_h)_L),\rho_{W_h}\bigr).
 \tag{RO9}\label{rog:eq:laurent}
\]
特に$\Gamma_R\otimes_RL$はsmooth properである。
完備底上の逆核のcompact性、または非可換側のsmoothnessを先に仮定しない。
\end{theorem}
\begin{proof}
$W_h$の判別式は特殊ファイバーで非零なので単元であり、$V(W_h)_L$は滑らか。
$E_L=0$なら全$r$について$\operatorname{Hom}(E_L,E_L(r))$のEuler数は零となる。
補題\ref{rog:lem:localization}と命題\ref{rog:prop:euler}に矛盾するので$E_L\ne0$。
定理\ref{rog:thm:orbit}により、$E_L\oplus E_L(1)\oplus E_L(2)$は
可換graded MF全体を生成する。
\eqref{rog:eq:image}の像は従って全MFである。
体上のOrlov比較を合成すると\eqref{rog:eq:laurent}を得る。
可換側のKuznetsov成分はsmooth properなので、Morita不変性から最後の結論を得る。
\end{proof}

\begin{remark}[既存の反例との相違]
$[\mathcal O_{\mathbb P^1}(-1)\xrightarrow{hs}\mathcal O_{\mathbb P^1}]$は一般点で点の層となり、
全Artinでの生成性から完備底の生成性を導けない。
本節はこの推論を行わない。一般ファイバーの\emph{非零性}と、
孤立斉次超曲面の\emph{全次数シフト軌道}という追加構造を使う。
Euler数が非零という条件だけでは、一般の導来圏で生成性は従わない。
\end{remark}

\subsection{体拡大上の同値から有限代数拡大上の同値へ}

\begin{lemma}[固定された可換模型へのgeneric降下]\label{rog:lem:descent}
$K$を標数零の体、$L/K$を任意の体拡大、$\Gamma,\mathsf C$を$K$上properなdg代数とし、
$\mathsf C$をsmoothとする。両側にrotationを表す可逆双加群と、必要なら三次周期同型を指定する。
$L$上でrotation・周期性を保つMorita同値があれば、ある有限代数拡大$K'/K$上でも
同じ二つの代数の間にそれらを保つMorita同値がある。
\end{lemma}
\begin{proof}
$L$上の同値を可換側から非可換側への核$P_L$で表す。
これは右$\Gamma_L$-perfectであり、$\mathsf C_L$のsmoothnessにより
$\mathsf C_L^{\rm op}\otimes_L\Gamma_L$上でもperfectである
\cite[Lemma~2.8]{ROGTV}。
$L$を有限生成$K$-部分代数のfiltered colimitとして表し、
\cite[Lemma~2.10]{ROGTV}により核とrotation比較を一つの有限生成整域$D\subset L$へ降ろす。
$\mathsf C_D$は$D$-properなので、包絡代数上のperfect加群を右$\Gamma_D$へ制限してもperfectである。

$Q_D=R\operatorname{Hom}_{\Gamma_D}(P_D,\Gamma_D)$と置く。
片側smoothnessにより$Q_D$は右$\mathsf C_D$-perfectであり、随伴の単位・余単位のconeは
$D$-perfectである。$L$への底変更後は零なので、
有限生成性と$D\subset L$により、ある$0\ne f\in D$を反転すれば両coneは零となる。
rotation比較のconeも同様に消せる。
指定周期性の図式はperfect双加群間の有限個の射の等式であるから、
必要なら段階を拡大し、さらに一つの元を反転して等式を保つ。
任意の高次群作用を有限の周期図式と同一視してはいない。

$\operatorname{Spec}D_f$の閉点を一つ取る。その剰余体$K'$はZariskiの補題により$K$の有限拡大であり、
標数零なので分離的である。そこでの底変更が所要の同値を与える。
これは元の$L$上の形式分枝を全次数で再現する主張ではない。
\end{proof}

\begin{theorem}[一変数量子化の幾何学的一般点]\label{rog:thm:geometric-generic}
\textup{(L),(CF)}を満たす一変数中心三次量子化に対し、
ある有限拡大$K'/\mathbb C(h)$と滑らかな三次超曲面$Y/K'$が存在して
\[
 (\mathcal K^{\rm lin}_{B_K}\otimes_KK',\rho)
 \simeq(\operatorname{Ku}(Y),\rho_Y)
 \tag{RO10}\label{rog:eq:generic}
\]
となる。$\mathcal K^{\rm lin}$と全coherent残余圏との同一性は、この定理に含めない。
\end{theorem}
\begin{proof}
定理\ref{rog:thm:laurent}で$L=\mathbb C((h))$上の同値を得た。
$W_h$の有限個の三次係数と判別式の逆を含む有限生成$K$-部分代数$D_0\subset L$を取ると、
$Y_0=V(W_{D_0})$は$D_0$上の滑らかな射影cubic族である。
その線束の有限和を残余成分へ射影した生成子を$H_{D_0}$、
生成dg代数を$\mathsf C_{D_0}=R\operatorname{End}(H_{D_0})$とする。
これはsmooth properで、底変更と両立する。
ここで生成子は相対ampleな全線束が生成することと、座標のKoszul完全列によって
$N$個の連続したtwistへ減らすことから得られる。
射影は有限個のexceptional mutationによる。
この有限線束生成の基底環に関する一般性は\cite[\S36.16]{ROGStacks}を参照する。

$L$上の同値核を、$D_0$を含む有限生成$K$-部分代数へ降ろす。
以後は補題\ref{rog:lem:descent}の証明と同じであり、
properな二つのconeを消して閉点を選ぶ。
その点では$Y_0$も滑らかなcubicへ特殊化し、\eqref{rog:eq:generic}を得る。
この一般存在定理は、$Y$の係数の閉公式や同値核の行列を明示するものではない。
\end{proof}

\begin{corollary}[計算済み混合三次式への復元]\label{rog:cor:mixed}
\textup{(L),(CF)}に加えて、既存の有限対角・characteristic比較\textup{(D)}を使う。
$\lambda=1$の混合族に対し、$\widetilde G_h$を
\[
\begin{aligned}
 \widetilde G_h={}&(x^3+y^3+z^3)/3+v^3+h^2xyz/3+(1+h^6/54)uxy\\
 &-h^4uz^2/18-h^2u^2z/12+
 (1+29h^6/648+h^{12}/4374)u^3
\end{aligned}
\tag{RO11}\label{rog:eq:explicit}
\]
とする。このとき、ある有限拡大$K'/\mathbb C(h)$上で
\[
 (\mathcal K^{\rm lin}_{B_K}\otimes_KK',\rho)
 \simeq(\operatorname{Ku}(V(\widetilde G_h)_{K'}),\rho_{\widetilde G}) .
\]
従って、この線束生成残余pairの\emph{幾何学的一般点}は、計算済みの代数的三次式から復元される。
\end{corollary}
\begin{proof}
既存の幾何学的低rotation底変更と\textup{(D)}により、
$W_h$と$\widetilde G_h$は$\widehat R$上で形式線形同型である。
その行列を$L$へ移せば通常の線形座標変換になる。
定理\ref{rog:thm:laurent}の可換模型をこの固定された$K$上のcubicへ換え、
補題\ref{rog:lem:descent}を適用する。
特殊式$\widetilde G_0$は滑らかなので、$\widetilde G_h$の一般点も滑らかである。
\end{proof}

\subsection{生成性と有限降下の射程}

以上の議論は、局所window比較(L)と完備MF比較(CF)を用いて、
一般点での生成性とcompactな逆核の存在を示す。
これらの比較の構成は、それぞれ定理\ref{tw:thm:window}と
定理\ref{cft:thm:complete}で与えた。
完備底全体の右$\widehat\Gamma$-perfectな逆核や、$h=0$を含むétale近傍の同値は結論しない。
体上で得た有限cellの存在は証明したが、そのcell数の明示上界、
双加群の具体的行列、必要な有限拡大の次数は求めていない。

本節の独立した内容は、次数シフト軌道の生成、Euler形式による非零性の検出、
固定対象を局所化する順序、および体拡大からの有限降下にある。
線束生成部分と全coherent残余圏の同定は
定理\ref{crs:thm:coherent}で行う。
正則多変数局所基底についても個々の一般点には同じ議論が適用できるが、
異なるPoisson成分を含む全モジュライ上の接着、任意の$\lambda$の明示式、
同値核の具体的行列と有限拡大次数の評価はここでは扱わない。

% 2026-11-11 / v36
\section{線束生成部分の飽和と全coherent残余圏}
\label{crs:sec:saturation}

本節は、一般点復元で用いた線束生成部分と全coherent導来圏の差を解消する。
鍵は、一般点で既に得た線束生成残余圏のsmoothnessを、
全coherent対象に対する有限性と組み合わせることである。
特殊点の生成性を一般点へ持ち上げる議論は使わない。

既知の入力は、Burke--Stevensonの導来section、Gorenstein環の
$\chi$条件と有限cohomological dimension、およびOrlov型分解
\cite[Proposition~4.18, Lemmas~6.1--6.2, Theorem~6.4]{CRSBS}、
smooth dg圏上のpseudo-perfect加群のperfect性
\cite[Lemma~2.8(2)]{CRSTV}、admissibleな半直交分解とsmoothness
\cite[Lemma~4.15(3)]{CRSPerry}である。
これらの一般原理自体の新規性は主張しない。
本研究での寄与は、その仮定を一般点の中心三次商について確認し、
線束生成残余pairの復元を全coherent残余pairへ拡張する組合せにある。
この組合せの包括的新規性調査・外部査読は未実施である。

\subsection{抽象的な飽和判定}

$k$を体、$\mathcal D$を小さな冪等完備pretriangulated $k$-線形dg圏とし、
$G\in\mathcal D$、$\mathcal C=\operatorname{thick}(G)$、
$\Lambda=\operatorname{RHom}_{\mathcal D}(G,G)$と置く。
以下の直交は全cohomological shiftに関するものである。

\begin{theorem}[有限性とsmoothnessによる全体への飽和]
\label{crs:thm:abstract}
次の三条件を仮定する。
\begin{enumerate}
\item $\Lambda$は$k$上smooth、すなわち
$\Delta_\Lambda\in\operatorname{Perf}(\Lambda^{\mathrm{op}}\otimes_k\Lambda)$である。
\item 任意の$M\in\mathcal D$について
$\operatorname{RHom}(G,M)\in\operatorname{Perf}(k)$である。
\item $\mathcal C^\perp=0$である。
\end{enumerate}
このとき$\mathcal C=\mathcal D$であり、
$\operatorname{Perf}(\Lambda)\to\mathcal D$はdg Morita同値である。
\end{theorem}
\begin{proof}
右$\Lambda$-加群$N_M=\operatorname{RHom}(G,M)$を考える。
smoothな対角の有限cell/retract表示へ
$N_M\otimes_\Lambda^{\mathbf L}(-)$を適用する。
自由な包絡加群の像は$N_M\otimes_k\Lambda$であり、
$N_M\in\operatorname{Perf}(k)$だから右$\Lambda$-perfectである。
従って$N_M\in\operatorname{Perf}(\Lambda)$となる。
これは\cite[Lemma~2.8(2)]{CRSTV}の一対象版の直接証明である。

$G$を用いるMorita関手
$F:\operatorname{Perf}(\Lambda)\xrightarrow{\sim}\mathcal C$
で$N_M$を移し、評価射$F(N_M)\to M$を取る。
ここで$F(N_M)$は有限coneと直和因子だけで作られ、
$\mathcal D$に無限直和があることを要求しない。
$F$と$\operatorname{RHom}(G,-)$の単位は$\Lambda$で同型なので、
有限操作による帰納法で全perfect加群に対して同型である。
従って評価射に$\operatorname{RHom}(G,-)$を施すと同型となる。
評価射のconeは$G$、従って$\mathcal C$と右直交する。
第三条件によりこのconeは零である。
全$M$が$\mathcal C$に属し、enhancedな完全忠実性は
$\operatorname{thick}(G)$におけるMorita表示から従う。
\end{proof}

\begin{corollary}[対角の有限表示からの一様な生成長]
\label{crs:cor:level}
上の対角が自由包絡加群から$q$段の拡大、有限直和、shift、直和因子で作れるなら、
任意の$M\in\mathcal D$も$G$から高々$q$段の同じ操作で作れる。
この系は$q$の存在から生成長を評価するもので、具体的な値は与えない。
\end{corollary}
\begin{proof}
前証明の各自由包絡加群の像は$N_M\otimes_k\Lambda$である。
体上では$N_M$はその有限次元有界cohomologyの直和に準同型だから、
この像は有限個の$\Lambda$のshiftの直和である。
テンソル関手と$F$は拡大段数を増やさない。
\end{proof}

\begin{example}[smoothnessを省けない]
$\Lambda=k[\epsilon]/(\epsilon^2)$、
$\mathcal D=D^b(\operatorname{mod}\Lambda)$、$G=\Lambda$とする。
$\operatorname{RHom}(G,M)$は全$M$について$k$-perfectで、$G^\perp=0$である。
しかし$k$の射影分解は$\epsilon$による無限周期分解であり、
$\operatorname{Tor}^{\Lambda}_i(k,k)=k$が全$i\ge0$に残る。
従って$\operatorname{thick}(G)=\operatorname{Perf}(\Lambda)$は$\mathcal D$より真に小さい。
単なるright-spanning性を有限生成性と同一視してはいけない。
\end{example}

\subsection{中心三次商のcoherent対象に対する有限性と検出}

$A$を連結局所有限な両側noetherian Koszul AS-regular $k$-代数とし、
global dimensionとGorenstein parameterをともに$N\ge5$とする。
$w\in Z(A)_3$を正則元、$B=A/(w)$、$a=N-3$と置く。
$\operatorname{qgr}B$は有限生成次数付き右加群を有限次元加群で割った圏、
$\mathcal O(j)=\pi B(j)$とする。標準のenhancementを用いる。

\begin{lemma}[必要なGorenstein・section条件]
\label{crs:lem:sections}
$B$は両側noetherian AS-Gorensteinで、parameterは$a$である。
$B$は$\chi$条件と有限cohomological dimensionを満たす。
任意の$M\in D^b(\operatorname{qgr}B)$と任意の$j\in\mathbb Z$について
\[
 \operatorname{RHom}_{\operatorname{qgr}B}(\mathcal O(j),M)
 \in\operatorname{Perf}(k).
 \tag{CS1}\label{crs:eq:finite-Hom}
\]
また、全$j,p$で$\operatorname{Hom}(\mathcal O(j),M[p])=0$なら$M=0$である。
\end{lemma}
\begin{proof}
$0\to A(-3)\xrightarrow{w}A\to B\to0$と導来change-of-ringsにより、
$B$-加群$M$について
\[
 R\underline{\operatorname{Hom}}_B(M,B)
 \simeq R\underline{\operatorname{Hom}}_A(M,A)(-3)[1]
\]
を得る。ここでは$w$の中央性を使い、normal元のtwistを省略したのではない。
左右でinjective dimensionは有限であり、$M=k$を代入すると
$R\underline{\operatorname{Hom}}_B(k,B)\simeq k(N-3)[-(N-1)]$となる。
$B$のnoetherian性は商であることから従う。
従って\cite[Lemmas~6.1--6.2]{CRSBS}を適用できる。

\cite[Proposition~4.18とその直前の導来随伴]{CRSBS}により、任意の整数$b$について
\[
 R\Gamma_{\ge b}:D^b(\operatorname{qgr}B)
 \longrightarrow D^b(\operatorname{gr}_{\ge b}B)
 \tag{CS2}\label{crs:eq:sections}
\]
は完全忠実で、各cohomologyは有限生成次数付き$B$-加群である。
固定内部次数の成分は有限次元であり、cohomologyは有界だから、
$\bigl(R\Gamma_*M\bigr)_{-j}=R\operatorname{Hom}(\mathcal O(j),M)$は
$k$-perfectとなる。任意の$j$に対し$b\le-j$を選べばよい。

全ての線束とのHomが零なら$R\Gamma_{\ge0}M$の各次数のcohomologyが零である。
この複体自体が零となり、完全忠実性から$M=0$を得る。
この検出は高次Extも含み、射の次数零だけを調べているのではない。
\end{proof}

\begin{lemma}[有限線束生成対象]
\label{crs:lem:finite-lines}
\[
 \mathcal C_B=\operatorname{thick}\{\mathcal O(j):j\in\mathbb Z\},
 \qquad G_B=\bigoplus_{j=0}^{N-1}\mathcal O(j)
\]
とすると$\mathcal C_B=\operatorname{thick}(G_B)$であり、$\mathcal C_B$はproperである。
\end{lemma}
\begin{proof}
$k$の$A$上の有限Koszul分解を$B$へ導来テンソルする。
二項の$A$-分解で計算すると$\operatorname{Tor}^A_i(k,B)$は
$i=0,1$にのみあり、それぞれ$k,k(-3)$である。
従ってこのKoszul複体を$qgr$へ移すと完全になる。
両端は一つの線束で、中間の項のshiftは$1,\ldots,N-1$にある。
任意の連続$N$個の線束から、一つ手前と一つ後をそれぞれ生成できる。
前後への帰納により全てのtwistを得る。
proper性は補題\ref{crs:lem:sections}を有限個の$\mathcal O(i),\mathcal O(j)$に適用し、
有限coneと直和因子へ延長して従う。
\end{proof}

\begin{theorem}[線束部分のsmoothnessで全coherent圏を得る]
\label{crs:thm:coherent}
上の中心三次商について、$\mathcal C_B$がsmoothなら
\[
 \boxed{\quad
 \mathcal C_B=D^b(\operatorname{qgr}B).
 \quad}\tag{CS3}\label{crs:eq:full}
\]
従って右辺はsmooth proper dg圏で、$G_B$が古典的生成対象になる。
\end{theorem}
\begin{proof}
任意のcoherent有界対象に対する$\operatorname{RHom}(G_B,-)$の有限性は
補題\ref{crs:lem:sections}による。線束による検出と
補題\ref{crs:lem:finite-lines}から$\mathcal C_B^\perp=0$である。
定理\ref{crs:thm:abstract}を適用する。
ここで全$D^b(\operatorname{qgr}B)$のsmoothnessを仮定していない。
\end{proof}

\subsection{残余部分のsmoothnessと有限拡大からの降下}

\begin{lemma}[有限体拡張からのsmoothness降下]
\label{crs:lem:descent}
有限体拡大$k'/k$と$k$-dg代数$\Gamma$について、
$\Gamma\otimes_k k'$が$k'$上smoothなら$\Gamma$は$k$上smoothである。
\end{lemma}
\begin{proof}
$E=\Gamma^{\mathrm{op}}\otimes_k\Gamma$と置く。
$E\otimes_k k'$は$E$上有限自由なので、そのperfect加群のrestrictionは$E$-perfectである。
従ってsmoothな対角$\Gamma\otimes_k k'$をrestrictionしたものは$E$-perfect。
$\lambda:k'\to k$を$\lambda(1)=1$となる$k$-線形写像とすれば、
$\Gamma\to\Gamma\otimes_k k'\xrightarrow{1\otimes\lambda}\Gamma$は
$E$-線形なretractである。従って$\Gamma$は$E$-perfectである。
\end{proof}

\begin{proposition}[残余のsmoothnessから全線束部分へ]
\label{crs:prop:residual-smooth}
$\mathcal K_B^{\rm lin}=\mathcal C_B\cap
\langle\mathcal O,\ldots,\mathcal O(a-1)\rangle^\perp$とする。
$\mathcal K_B^{\rm lin}$がsmoothなら$\mathcal C_B$もsmoothである。
従ってこの場合
\[
 \mathcal K_B^{\rm lin}=\mathcal K_B^{\rm coh}:=
 D^b(\operatorname{qgr}B)\cap
 \langle\mathcal O,\ldots,\mathcal O(a-1)\rangle^\perp.
 \tag{CS4}\label{crs:eq:residual-full}
\]
\end{proposition}
\begin{proof}
$\mathcal O,\ldots,\mathcal O(a-1)$は例外列である。
これはGorensteinの線束cohomology公式、または
\cite[Theorem~6.4]{CRSBS}のOrlov分解から従う。
$\mathcal C_B$はproperなので、例外列への評価・余評価は有限であり、
その生成部分はadmissibleである。残余部分もadmissibleとなり
\[
 \mathcal C_B=\langle\mathcal K_B^{\rm lin},
                 \mathcal O,\ldots,\mathcal O(a-1)\rangle
\]
を得る。smoothnessの貼合せ\cite[Lemma~4.15(3)]{CRSPerry}を適用する。
具体的には各例外成分は$\operatorname{Perf}(k)$でsmooth、
相互Homはproperであり、smoothな成分の間の貼合せ双加群は
\cite[Lemma~2.8(2)]{CRSTV}によりperfectである。
上三角dg圏の対角の有限三角を順に貼り合わせれば全体の対角もperfectとなる。
定理\ref{crs:thm:coherent}から\eqref{crs:eq:residual-full}を得る。
\end{proof}

\subsection{一般点復元を全coherent残余pairへ拡張する}

\begin{theorem}[全coherent残余pairの幾何学的一般点復元]
\label{crs:thm:generic}
$K=\mathbb C(h)$とし、既存の一般点復元定理の入力\textup{(L),(CF)}を満たす
一変数量子化$B_K=A_K/(\Omega_K)$を考える。
$A_K$は上のnoetherian Koszul AS-regular仮定を満たすとする。
この仮定は本研究の局所量子化入力の一般ファイバーで成立する。
既存定理が与える有限拡大$K'/K$と滑らかな三次式$W'/K'$について
\[
 \boxed{\quad
 (\mathcal K^{\rm coh}_{B_{K'}},\rho_B)
 \simeq(\operatorname{Ku}(V(W')),\rho_{W'}).
 \quad}\tag{CS5}\label{crs:eq:generic}
\]
さらに$D^b(\operatorname{qgr}B_K)=\mathcal C_{B_K}$であり、
このdg圏と$\mathcal K^{\rm coh}_{B_K}$は$K$上smooth properである。
混合$\lambda=1$で既存入力\textup{(D)}も用いる場合、$W'$は既に計算した
多項式$G_h+v^3$に取れる。
\end{theorem}
\begin{proof}
既存の次数シフト軌道による一般点定理を用いると、
$\mathcal K^{\rm lin}_{B_{K'}}$は滑らかな可換cubicのKuと同値でありsmooth properである。
線束生成子のHomの平坦底変更も入力(L)に含まれているため、
残余生成代数について$\Gamma_{K'}=\Gamma_K\otimes_KK'$と同定できる。
補題\ref{crs:lem:descent}で$\Gamma_K$のsmoothnessが従う。
命題\ref{crs:prop:residual-smooth}と定理\ref{crs:thm:coherent}を$K$上で適用する。

同じ議論は$K'$上でも成立する。有限拡大後のnoetherian性は
$A_K\otimes_KK'$が$A_K$上有限加群であることから従う。
標数零なので$K'/K$は分離的であり、任意の加群は
restriction後のscalar extensionの直和因子になる。
この分離idempotentと$A_K$上の有限射影分解によりglobal dimensionの上界を保つ。
Koszul分解とGorensteinの等式は平坦底変更で保たれる。
従って\eqref{crs:eq:residual-full}を$K'$にも適用でき、
既存の線束部分の同値がそのまま全coherent残余pairの同値となる。
rotationは同じ$L_{\mathcal O}\circ(1)$であり、新たな関手への変更はない。
(D)による既存の明示式同定を合成すると最後の結論を得る。
\end{proof}

\begin{remark}[主張の境界]
本定理は局所window比較(L)と完備MF比較(CF)を用いる。
これらは定理\ref{tw:thm:window}および定理\ref{cft:thm:complete}で構成されている。
特定の可換方程式を指定する場合には、さらに比較(D)を用いる。
全$D^b(\operatorname{qgr}B_{K'})$と$D^b(\operatorname{coh}V(W'))$の同値は主張しない。
残余成分が同値でも例外成分との貼合せは異なり得る。
またアフィン次数付き環$B_K$のglobal dimensionが有限とは言っていない。
特殊点を含む同値のspread、同値核の行列・有限拡大次数、
複数の基底チャートにわたる相対同値の接着は別の問題である。
\end{remark}

\subsection{既知の生成原理と適用条件の確認}

StevensonのTheorem~6.13\cite{CRSStevenson}は、affineな局所超曲面環について
\emph{全ての}thick部分圏を分類する。tensor idealに限定した定理ではない。
本研究で使う$K[x]/(W)$は原点のみが特異なので、適用範囲は合っている。
次数の再付与では有限対象から無限直和へ移るが、有限rank MFからの射は
有限個の直和成分を通るため、最後のcompact局所化によって有限thick構成へ戻る。
従って、次数の再付与に現れる無限直和は最終的な生成性の結論ではなく、
有限thick閉包に戻るための中間的な構成である。

完備MF同値の局所化では、先に固定有限対象へ制限してから$h$を反転する。
各固定cohomological次数は有限自由な斉次成分からなり、
平坦局所化はそのHomとcohomologyに可換する。
任意のLaurent係数対象へ新しく無限Taylor和を評価することは使わない。
曲率付き関手と局所windowの構成は、それぞれ
第\ref{cft:sec:transport}節と第\ref{tw:sec:window}節に与える。

全coherent残余圏への拡張には、上で示したsection有限性、
smooth対角、評価射および直交成分の検出を用いる。
有限係数のEuler計算は非零性を調べるものであって、
これらの圏論的議論を置き換えるものではない。

% 2026-11-12 / B-v41. Current proof, not a version-history appendix.
\section{一般点復元の基礎体降下と定量的生成}
\label{gcd:sec:descent}

本節では可換模型の\emph{幾何学的一般点での存在}を得た後の議論を整理する。
方程式の定義体と、圏同値の定義体は異なる。
前者は低rotation積の可積分性により元の体へ降下できるが、
後者について本稿が保証するのは有限拡大後の存在である。
可積分性の有限線形代数そのものは第\ref{sec:explicit-jacobian-recovery}節と
第\ref{far:sec:recovery}節の結果を用いる。

\subsection{次数別flat入力から従う一般ファイバーのAS正則性}

\begin{proposition}[一般点のambientに対する追加仮定の除去]
\label{gcd:prop:automatic-as}
$(R,\mathfrak m,k)$をnoetherian局所整域、$K=\operatorname{Frac}R$とする。
$A_R$は非負次数の中央$R$-代数、$(A_R)_0=R$、各$(A_R)_j$は有限自由で、
$A_R/\mathfrak mA_R\simeq k[x_1,\ldots,x_N]$とする。
このとき$A_K$は左右noetherianな二次Koszul AS正則代数であり、
大域次元とGorenstein parameterは$N$、Hilbert級数は$(1-z)^{-N}$である。
さらに$\Omega_R\in Z(A_R)_d$の剰余が非零なら、$\Omega_R$は左右正則で、
商の各次数は$R$-自由である。
\end{proposition}
\begin{proof}
補題\ref{lp:lem:noetherian}の次数別Krull交叉の議論を$A_R$に適用すると、
左右noetherian性を得る。中央集合$R\setminus\{0\}$による局所化で$A_K$もそうなる。
ここでnoetherian性を有限Hilbert級数だけから推論していない。

補題\ref{lp:lem:resolution}を、augmentation加群$R$と、
$A_R^e=A_R\otimes_RA_R^{\mathrm{op}}$上の対角$A_R$に適用する。
多項式環のKoszul分解と対角分解を持ち上げると、項がそれぞれ
$A_R(-j)^{\binom Nj}$、$A_R^e(-j)^{\binom Nj}$である長さ$N$の分解を得る。
持上げの各段で、剰余上の生成元を持ち上げた写像の全射性を各内部次数の
Nakayamaで示す。有限自由$R$-加群への全射は分裂するので核も各次数で自由であり、
最後の核は剰余が零のため零となる。この構成は任意の通常の底変更に可換である。
したがって$A_K$は二次Koszulで、対角分解から左右大域次元は高々$N$となる。

augmentation分解の$A_R$双対を固定内部次数ごとに取る。
その剰余複体は、内部次数$-N$・コホモロジー次数$N$以外で完全である。
有界有限自由複体の単元行列成分を順に消去すれば、他の内部次数の複体は分裂完全となる。
内部次数$-N$では最後の$R$のみが残る。よって左右で
\begin{equation}\label{gcd:eq:as}
 R\underline{\operatorname{Hom}}_{A_R}(R,A_R)\simeq R(N)[-N]
\end{equation}
となり、底変更後の非零な最上Extから大域次元は正確に$N$である。
次数別自由性はHilbert級数を与え、増大度も多項式的である。
最後に各次数の$\Omega_R$乗法は剰余上で単射である。
列数と同じ大きさの単元小行列式を使えば$R$上でも分裂単射となり、
余核も自由となる。中央性により左右で同じ結論を得る。
\end{proof}

\subsection{smooth properなpairから取り出す低次数積}

$k$を標数零の体、$\mathcal C$をsmooth properな$k$-線形dg圏、
$\rho$をdg Morita自己同値とする。対角を$\Delta$、$\rho$の双加群を$P$と書き、
\begin{equation}\label{gcd:eq:low}
 V=H^0R\operatorname{Hom}_{\mathcal C^e}(\Delta,P),\qquad
 U=H^0R\operatorname{Hom}_{\mathcal C^e}(\Delta,P\otimes_{\mathcal C}^{\mathbf L}P)
\end{equation}
と置く。積$\mu:V\otimes V\to U$は$\rho(\beta)\circ\alpha$である。
通常の三角圏上の自然変換と、この導来自然変換を区別する。

\begin{lemma}[低次数積の体拡大との両立]\label{gcd:lem:basechange}
$V,U$は有限次元であり、任意の体拡大$\ell/k$に対して、その形成と積は底変更と可換である。
\end{lemma}
\begin{proof}
smoothnessにより$\Delta$は$\mathcal C^e$上perfectである。
proper性と$P$の可逆性により、$P$と$P\otimes_{\mathcal C}^{\mathbf L}P$は
基礎体上perfectである。有限cellとretractで$\Delta$を表示すると、
式\eqref{gcd:eq:low}のHom複体も$k$-perfectとなる。
同じ有限表示を使えばHomと平坦な体拡大が可換であり、$H^0$も可換である。
双加群のテンソルと射の合成はこの底変更に自然なので、積も保たれる。
\end{proof}

\begin{theorem}[可換cubic模型の方程式は基礎体上に取れる]
\label{gcd:thm:field-model}
$n\ge3$、$N=n+2$とする。上のpairがある代数閉拡大体$\Omega/k$上で
滑らかなcubic $Y/\Omega$のrotation付きKuznetsov pairとMorita同値であると仮定する。
このとき次を満たす$G\in\operatorname{Sym}^3_k V$が存在する。
\begin{enumerate}
\item $X_G=V(G)\subset\mathbb P(V^*)$は$k$上幾何学的に滑らかで、$X_{G,\Omega}\simeq Y$である。
\item ある有限拡大$k'/k$上で$(\mathcal C_{k'},\rho)\simeq(\operatorname{Ku}(X_{G,k'}),\rho_G)$である。
元の同値に三次周期比較を指定していた場合、それも保存できる。
\item $\mathcal C$の任意の滑らかなcubic模型は、共通の代数閉拡大体上で射影同型である。
\end{enumerate}
従って方程式の定義体には有限拡大を必要としない。
一方、同値双加群の$k$への降下や、$k$上での模型の一意性は主張しない。
\end{theorem}
\begin{proof}
命題\ref{prop:ej-geometric-low-degree}に用いる
Lin--Zhangのsector計算\cite[Theorem~3.5, Remark~3.6]{LZ}を$\Omega$上で適用する。
三次式の非恒等sectorがコホモロジー次数零へ寄与する内部次数は、
$N$が偶数の場合の$N/2$のみである。$N\ge5$では内部次数$1,2$に干渉しない。
したがって、$\Omega$上で$V,U,\mu$はJacobian環の次数$1,2$と積に一致する。
この計算は標数零の代数閉体上の代数的なfactorization計算である。
補題\ref{gcd:lem:basechange}と忠実平坦性から、$k$上でも
\begin{equation}\label{gcd:eq:gradient-data}
 \dim V=N,\quad \dim U=\binom N2,\quad
 \overline\mu:\operatorname{Sym}^2V\twoheadrightarrow U,\quad
 Q:=\ker\overline\mu,\quad\dim Q=N
\end{equation}
となる。積の対称性も降下する。
$Q$には代数閉包上で共通射影零点がない。これは$\Omega$上で勾配空間となることと
$Y$の滑らかさから従う。

$V$の基底を$x_i$、$Q$の基底を$q_a$とする。可積分性の線形解空間
\begin{equation}\label{gcd:eq:integrability}
 \mathcal L_Q=\left\{A\in\operatorname{Mat}_N(k):
 \sum_a A_{ia}\partial_jq_a=\sum_a A_{ja}\partial_iq_a\quad(i<j)\right\}
\end{equation}
は体拡大と可換である。$Y$の方程式の勾配を$q_a$で展開すると、
$\mathcal L_Q\otimes_k\Omega$に可逆行列が存在する。
$k$-基底$A_1,\ldots,A_r$を選べば
$D(z)=\det(\sum_\nu z_\nu A_\nu)$は零多項式でない。
$k$は無限体なので$D(c)\ne0$となる$c\in k^r$が存在する。
実際、各変数の次数は高々$N$なので$\{0,1,\ldots,N\}^r$のいずれかで非零となる。
$A=\sum c_\nu A_\nu$として
\begin{equation}\label{gcd:eq:model}
 G=\frac13\sum_{i,a}x_iA_{ia}q_a
\end{equation}
と置く。可積分性とEuler等式から$\partial_iG=\sum_aA_{ia}q_a$である。
従って$G$の勾配空間は$Q$、その共通射影零点はなく、$X_G$は幾何学的に滑らかである。

$\Omega$上で$G$と$Y$の方程式は同じ勾配空間を持つ。
命題\ref{prop:ej-projective-uniqueness}の対称三重線形形式の証明は
任意の標数零代数閉体で通るので、両者は線形座標変換で一致する。
この証明はsymmetrizer代数の可換性と非零冪零元の不存在を用い、
固有空間ごとの非零係数を三乗根で吸収するものである。
その座標変換はgraded MFの内部shiftと三次周期性を保つ。
従って$\Omega$上で$(\mathcal C,\rho)$と$X_G$のpairの同値を得る。

$X_G$は既に$k$上にあり、そのKuznetsov成分はsmooth properである。
固定された二つのproper dg模型に補題\ref{rog:lem:descent}を適用する。
perfect同値核、rotation比較、指定周期図式を有限生成$k$-部分代数へ降ろし、
単位・余単位のconeを消す非空開の閉点を選べば、有限拡大$k'/k$上の同値を得る。
第三の主張も、任意の模型に式\eqref{gcd:eq:gradient-data}を適用し、
同じ勾配空間の射影的一意性を用いれば従う。
\end{proof}

\begin{remark}[存在証明と実計算の違い]\label{gcd:rem:limits}
この定理は幾何学的な可換模型の存在を\emph{仮定して}、その方程式の定義体を改善する。
任意の非可換圏で次元と積分条件を検査するだけでは圏同値は得られない。
また、$V,U,\mu$の有限性の証明は、その非可換入力での具体行列を計算したという意味ではない。
$k=\mathbb C(h)$のときも$G$の係数が$h=0$で正則、あるいは$G_0=F$とは限らない。
特殊点を含む代数化には第\ref{far:sec:recovery}節の局所入力を別途確認する必要がある。
\end{remark}

\begin{example}[同じ勾配空間でも基礎体上の同型類は一意でない]
\label{gcd:ex:forms}
$k=\mathbb C(h)$、$N\ge5$とし、
\[
 F=x_1^3+\cdots+x_N^3,\qquad
 H=h x_1^3+x_2^3+\cdots+x_N^3
\]
を考える。両方とも幾何学的に滑らかで、勾配空間は
$\langle x_1^2,\ldots,x_N^2\rangle$である。
$k(h^{1/3})$上では同型だが、$k$上では射影同型でない。
実際、両方のHessianの零因子は座標超平面の和である。
射影線形同型はこの超平面配置を保つのでmonomial行列で表される。
係数を比較すると$h$が$k$の三乗であることを要するが、
$h$進付値が1なのでこれは不可能である。
したがって、勾配空間から得る方程式は\emph{幾何学的な}同型類を定めるのであり、
基礎体上のformや同値双加群の降下障害を自動的に消すものではない。
\end{example}

\subsection{既知の生成定理による最短経路と一様な上界}

\begin{proposition}[三shiftの定量的生成：BFKの帰結]
\label{gcd:prop:bound}
$k$は標数零の代数閉体、$Y\subset\mathbb P^{N-1}_k$は滑らかなcubic、$N\ge5$とする。
任意の非零$E\in\operatorname{Ku}(Y)$について
\begin{equation}\label{gcd:eq:bound}
 \operatorname U\bigl(E\oplus\rho_YE\oplus\rho_Y^2E\bigr)
 \le 2N(N+1)-1
\end{equation}
となる。ここで$\operatorname U$は、有限直和・shift・直和因子を段数零として、
圏の全対象を得るためのcone段数の最小値である。
\end{proposition}
\begin{proof}
Orlov比較により$\rho_Y$はgraded singularity categoryの内部shiftとなる。
Ballard--Favero--Katzarkov\cite[Theorem~5.15]{GCDBFK}は、
$N$変数の次数$d$の孤立斉次超曲面について、任意の非零対象の$d$個のshiftの直和が
生成時間$2N\bigl(N(d-2)+1\bigr)-1$以下のstrong generatorとなることを証明している。
$d=3$を代入すればよい。滑らかな射影cubicのアフィン錐の特異点は原点だけである。
\end{proof}

上の生成性と上界はBFKの既知定理の帰結である。
第\ref{rog:sec:generic}節にはStevenson分類と次数再付与を用いる別証明も与える。
一般点復元の主証明では、上記のBFK定理を直接適用する。
BFKの論文は基礎体を代数閉・標数零としているので、
$\mathbb C((h))$へそのまま適用せず、まずその代数閉包へ底変更する。
一般点同値の有限拡大への降下はその後で行う。
この上界は同値核の行列サイズや必要な体拡大次数の上界ではない。

\begin{proposition}[固定三対象の比較による独立の十分条件]
\label{gcd:prop:three-object}
主定理の代数的入力および局所比較(L)を仮定する。
完備MF圏全体の比較(CF)を使用せず、以下の固定三対象の比較を与えるだけでも
定理\ref{gdr:thm:generic}の結論が成立する。
ある体拡大$\ell/K$、滑らかなcubic $Y/\ell$、非零$E\in\operatorname{Ku}(Y)$について、
\[
 \Gamma_K\otimes_K\ell\simeq
 R\operatorname{End}(E\oplus\rho_Y E\oplus\rho_Y^2 E)
\]
というenhancedなMorita比較があり、これが生成部分圏上のrotationと指定周期比較を保つ。
ここでは単に両側のcohomology群や低次数の積が一致するだけでは不十分である。
\end{proposition}
\begin{proof}
右辺のperfect加群圏は三対象のthick閉包である。
$\ell$の代数閉包上で命題\ref{gcd:prop:bound}を適用すれば、
それは$\operatorname{Ku}(Y)$全体となる。
以後は主定理の有限降下、coherent飽和、基礎体上の方程式抽出と同じである。
\end{proof}

この条件は既に構成された代数的sourceに対する\emph{固定三対象の実際の比較}であり、
全MF圏上の関手を先に構成することを要しない。
この命題は、主定理に用いた全MF比較よりも弱い十分条件を分離したものである。
他方、相対的な変形理論では全MF比較そのものが必要になるため、
その構成は第\ref{cft:sec:transport}節に詳述する。

\subsection{多変数局所基底への適用と接着の問題}

\begin{proposition}[任意次元の正則局所基底の一般点]
\label{gcd:prop:multibase}
$(R,\mathfrak m,\mathbb C)$を有限生成複素代数の正則局所環とし、
$K=\operatorname{Frac}R$とする。次数別有限自由な中心三次対$(A_R,\Omega_R)$の
特殊ファイバーを多項式環と滑らかなcubicとする。
条件(L)をこの$R$について、条件(CF)を$\mathfrak m$進完備化$\widehat R$について仮定する。
このとき定理\ref{gdr:thm:generic}の結論は$K$をこの分数体に置き換えて成立する。
\end{proposition}
\begin{proof}
正則局所環の完備化は整域で、$K\hookrightarrow L=\operatorname{Frac}\widehat R$である。
(CF)を固定三対象に制限してから$L$、次いでその代数閉包へ底変更する。
(L)のperfect Homにより特殊点と一般点のEulerベクトルは同じで、
そのノルムは$2\cdot3^{N-1}>0$である。
命題\ref{gcd:prop:bound}により像は全可換MF圏を生成する。
以後のperfect核の有限降下、有限拡大からのsmoothness降下、coherent飽和は
基礎体が$\mathbb C(h)$であることを使わない。
命題\ref{gcd:prop:automatic-as}と定理\ref{gcd:thm:field-model}で結論を得る。
\end{proof}

\begin{lemma}[多変数でも成立する形式Poisson消去の代数部分]
\label{gcd:lem:multi-koszul}
$\widehat R=\mathbb C[[t_1,\ldots,t_r]]$、$S=\mathbb C[x_1,\ldots,x_N]$とし、
$W\in\widehat R\otimes S_3$の剰余が滑らかな射影cubicを定めるとする。
$\pi\in\mathfrak m(\widehat R\otimes S_2\otimes\bigwedge^2\langle\partial_x\rangle)$が
$[\pi,W]=[\pi,\pi]=0$を満たせば、一意な
$\Xi\in\mathfrak m\widehat R\otimes\bigwedge^3\langle\partial_x\rangle$があり、
$\pi=[\Xi,W]$、$[\Xi,\pi]=0$である。
\end{lemma}
\begin{proof}
勾配Koszul複体は特殊ファイバーで正の外積次数において完全である。
固定した重み（係数次数と外積次数の二倍の和）ごとに有限自由複体を取り、
$\mathfrak m$進の次数ごとにcycleを持ち上げれば、完備化上でも同じ完全性を得る。
二次係数の二重vectorの原像は係数次数零の三重vectorである。
その差の原像は係数次数$-2$を要求するので一意である。
Jacobi恒等式により$[\Xi,\pi]$は係数次数一の四重vectorのKoszul cycleであり、
その原像は係数次数$-1$となるから零である。
\end{proof}

この補題から、周期変数$u$について
$\exp(\operatorname{ad}(\Xi/u))(-uW+\pi)=-uW$が従う。
実際の完備MF比較と相対windowの構成には、それぞれ
定理\ref{cft:thm:complete}と定理\ref{tw:thm:window}を用いる。
これらの局所的な比較をモジュライの異なるチャートにわたって接着し、
共通の相対模型を構成する問題は本節の結論に含めない。

% 2026-11-12 / B-v47: editorial revision; mathematical statements and proofs preserved.
\section{中心三次非可換変形の一般点可換復元}
\label{gdr:sec:main}

\subsection{代数的入力と残余圏の定義}
$n\ge3$、$N=n+2$、$a=N-3$、$R=\mathbb C[h]_{(h)}$、
$\widehat R=\mathbb C[[h]]$、$K=\mathbb C(h)$、$L=\mathbb C((h))$とする。
$A_R=\bigoplus_{j\ge0}(A_R)_j$は中央$R$-代数、$(A_R)_0=R$で、
各次数成分は有限自由、$A_R/hA_R\simeq\mathbb C[x_1,\ldots,x_N]$と仮定する。
$\Omega_R\in Z(A_R)_3$の剰余$F$は滑らかな射影三次超曲面を定めるものとし、
$B_R=A_R/(\Omega_R)$と置く。
以下の結果はこの代数的な組を対象とする。
抽象的な非可換三次圏に対して同様の表示が存在することは仮定せず、結論ともしない。

体上の$B=B_K$について$\mathcal O_B(j)=\pi B(j)$、
$\mathcal C_B=\operatorname{thick}\{\mathcal O_B(j):j\in\mathbb Z\}$とし、
\begin{equation}\label{gdr:eq:residual-defs}
\begin{aligned}
 \mathcal K_B^{\rm lin}
 &=\mathcal C_B\cap\langle\mathcal O_B,\ldots,\mathcal O_B(a-1)\rangle^\perp,\\
 \mathcal K_B^{\rm coh}
 &=D^b(\operatorname{qgr}B)\cap\langle\mathcal O_B,\ldots,\mathcal O_B(a-1)\rangle^\perp
\end{aligned}
\end{equation}
を区別する。rotationは
\begin{equation}\label{gdr:eq:rotation}
 \rho_B=L_{\mathcal O_B}\circ(1)
\end{equation}
とする。滑らかな$Y\subset\mathbb P^{n+1}_k$についても
\begin{equation}\label{gdr:eq:ku-def}
 \operatorname{Ku}(Y)=\langle\mathcal O_Y,\ldots,\mathcal O_Y(n-2)\rangle^\perp,
 \qquad \rho_Y=L_{\mathcal O_Y}\circ(1)
\end{equation}
と定める。enhancementと冪等完備化を伴う。
線束生成残余と全coherent残余の一致は後で証明する。

\subsection{主定理に用いる二つの比較定理}
局所比較(L)は定理\ref{tw:thm:window}により代数的入力から構成する。
$K_R=\operatorname{stab}(R)$と$T_R=K_R\oplus K_R(1)\oplus K_R(2)$について、
$\mathcal K^{\rm lin}_{B_R}\simeq\operatorname{Perf}(R\operatorname{End}(T_R))$であり、
Homは$R$-perfect、底変更、内部shift、中心元由来の周期比較と整合する。

形式比較(CF)も定理\ref{cft:thm:complete}で構成する。
すなわち、$W_h\in\widehat R[x]_3$、$W_h\bmod h=F$と、実際の
unital $A_\infty$同値
\begin{equation}\label{gdr:eq:cf}
 \operatorname{MF}^{\mathbb Z}(A^{\rm gr}_{\widehat R},\Omega_{\widehat R})
 \simeq\operatorname{MF}^{\mathbb Z}(\widehat R[x_1,\ldots,x_N],W_h)
\end{equation}
があり、内部shiftと$(3)\simeq[2]$を保つ。
その固定三対象への制限が定義\ref{tw:def:CF3}を満たす
（系\ref{cft:cor:CF3}）。以後、(L)、(CF)、(CF$_3$)は構成された比較の略称とし、
独立した追加仮定とはしない。既存の支持節でそれらを仮定する命題には、
この二つの比較定理を適用する。

\begin{theorem}[代数的中心三次変形の一般点可換復元]
\label{gdr:thm:generic}
上の代数的中心三次対に対し、幾何学的に滑らかな三次超曲面
$Y=V(G)\subset\mathbb P^{n+1}_K$を元の体$K=\mathbb C(h)$上に取れて、
ある有限拡大$K'/K$上で
\begin{equation}\label{gdr:eq:main}
 (\mathcal K^{\rm coh}_{B_{K'}},\rho_B)
 \simeq(\operatorname{Ku}(Y_{K'}),\rho_Y)
\end{equation}
となる。同値は中心元からの三次周期比較を保つMorita同値である。
さらに次が成立する。
\begin{enumerate}
\item $D^b(\operatorname{qgr}B_K)$はsmooth properで、
$\bigoplus_{j=0}^{N-1}\mathcal O_B(j)$が古典的生成対象である。
\item $Y$は任意の共通代数閉拡大体上で射影同型を除いて一意である。
\item 代数閉包$\overline K$上の任意の非零$M\in\mathcal K^{\rm coh}_{B_K}\otimes_K\overline K$に対し、
\begin{equation}\label{gdr:eq:uniform-bound}
 \operatorname U(M\oplus\rho_BM\oplus\rho_B^2M)\le2(n+2)(n+3)-1.
\end{equation}
段数零は有限直和、shift、直和因子のみを許す規約とする。
\end{enumerate}
\end{theorem}
\begin{proof}
\emph{一般点への移行。}
命題\ref{gcd:prop:automatic-as}により$A_K$は左右noetherian Koszul AS正則で
次元とparameterは$N$、$\Omega$は正則である。
定理\ref{tw:thm:window}の固定三対象へ定理\ref{cft:thm:complete}を制限し、
像を$E,E(1),E(2)$とする。固定した全Homと高次積を先に構成してから$L$へ局所化する。
任意のLaurent係数対象へ無限Taylor和を評価しない。
完全忠実性は元の比較と平坦性から従い、像は三対象のthick閉包である。

特殊ファイバーで命題\ref{rog:prop:euler}は
\begin{equation}\label{gdr:eq:euler}
 \chi_r=\sum_{\substack{0\le s\le N\\s\equiv-r\;({\rm mod}\,3)}}
 (-1)^s\binom Ns,\qquad \sum_{r=0}^2\chi_r^2=2\cdot3^{N-1}>0
\end{equation}
を与える。proper HomによるEuler数の不変性から$E_L\ne0$。
$W_h$の判別式は特殊化で非零なので単元である。
代数閉包$\overline L$上でBFKの既知定理による命題\ref{gcd:prop:bound}を適用すると、
三shiftの像は全可換graded MFを生成し、Orlov比較から
\begin{equation}\label{gdr:eq:laurent}
 (\mathcal K^{\rm lin}_{B_K}\otimes_K\overline L,\rho_B)
 \simeq(\operatorname{Ku}(V(W_h)_{\overline L}),\rho_{W_h})
\end{equation}
となる。

\emph{有限降下とcoherent飽和。}
三次係数、判別式の逆、perfect同値核、rotationと周期比較を
有限生成$K$-部分代数へ降ろす。
補題\ref{rog:lem:descent}と定理\ref{rog:thm:geometric-generic}の降下の証明により、
非空開の閉点を選び、有限拡大上の可換Kuznetsov pairとの同値を得る。
補題\ref{crs:lem:descent}でsmoothnessを$K$へ戻し、
命題\ref{crs:prop:residual-smooth}と定理\ref{crs:thm:coherent}でcoherent飽和を行う。
これが第一項と全coherent残余圏の比較を与える。

\emph{方程式の定義体と一意性。}
得られた$K$上のsmooth proper pairへ定理\ref{gcd:thm:field-model}を適用する。
低rotation積の底変更、勾配空間の可積分性と無限体上の可逆行列の選択により、
方程式$G$を$K$上に取れる。必要なら同値核の有限定義体を拡大して
式\eqref{gdr:eq:main}を得る。
勾配空間の幾何学的一意性が第二項、可換模型へのBFKの適用が第三項を与える。
\end{proof}

\begin{corollary}[全ての$n\ge3$における可換模型の一意性]
\label{gdr:cor:uniqueness}
上のrotation付きpairの可換三次模型は代数閉包後に射影同型を除いて一意である。
\end{corollary}
\begin{proof}
低rotation積はJacobian環の次数$1,2$と積を与え、同じ勾配空間を持つ
滑らかな三次式は射影同型である（定理\ref{gcd:thm:field-model}）。
これはLin--Zhang\cite[Theorem~1.7, Section~5]{LZ}とも整合する既知のTorelliの帰結であり、
一意性自体の新規性を主張しない。
\end{proof}

\begin{corollary}[正則多変数局所基底の一般点]\label{gdr:cor:multibase}
$R$が有限生成複素代数の局所化である正則局所環、剰余体が$\mathbb C$で、
各次数有限自由な中心三次対と滑らかな可換特殊ファイバーが与えられたとする。
$K=\operatorname{Frac}(R)$についても、主定理と同じ結論が成立する。
\end{corollary}
\begin{proof}
定理\ref{tw:thm:window}を$R$上で、定理\ref{cft:thm:complete}をその完備化上で用いる。
$R,\widehat R$は整域であり、$K\subset\operatorname{Frac}(\widehat R)$へ
固定三対象の比較を移せる。以後は命題\ref{gcd:prop:multibase}の証明における
非零Eulerベクトル、BFK、有限降下、飽和、基礎体上の勾配積分を適用する。
\end{proof}

\subsection{明示模型と結論の射程}
\begin{corollary}[五変数混合族の明示的復元]\label{gdr:cor:mixed}
$n=3$の既存の$\lambda=1$混合一変数量子化について、
構成した比較に対する有限対角・乗法的characteristic比較(D)を仮定する。
このとき可換模型は
\begin{equation}\label{gdr:eq:mixed}
\begin{aligned}
 W^{\rm mix}_h={}&(x^3+y^3+z^3)/3+v^3+h^2xyz/3\\
 &+(1+h^6/54)uxy-h^4uz^2/18-h^2u^2z/12\\
 &+(1+29h^6/648+h^{12}/4374)u^3
\end{aligned}
\end{equation}
の零点に取れる。圏同値は有限拡大後に得る。
\end{corollary}
\begin{proof}
(D)による系\ref{gr:cor:formal-polynomial}が形式模型とこの式の形式線形同型を与える。
それを$L$へ移し、系\ref{rog:cor:mixed}の有限降下とcoherent飽和を適用する。
(L),(CF)は本稿で構成した比較を使う。
\end{proof}

主定理の存在部分には(D)や低次数自然変換の積の明示行列を要しない。
一方で同値核の$K$への降下、$h=0$で正則な$G$、特殊点込み相対同値、
全Poisson moduli上の接着、有限拡大次数・同値核行列の明示は未解決である。
全射影化の導来圏を可換cubic全体の導来圏と同一視しない。
形式比較の証明を構成したこと、原稿全体の全補題を独立再監査したこと、
この結果の適用は、指定された次数別有限自由な代数的中心三次対に限られる。

\section{各Poisson成分の支配的チャートと条件付き一般点復元}
\label{sec:all-poisson-components}

本節では、任意の$n\ge3$の滑らかな複素三次超曲面$X^n$と
アフィンPoissonスキームの被約既約成分$Z$について、
その一般点を扱うための滑らかな支配的チャートを構成する。
得られる幾何学的なチャートと、そこからrotation付き可換復元を導くための
圏論的な追加条件を区別する。
この構成は三次三様体の17次元族や特定のPoisson階数に限定されない。
とくに、特異な錐の頂点への量子化族の降下を
一般点での存在問題の前提としない。

\subsection{各成分の零点と滑らかな支配的chart}

\begin{proposition}[Poisson錐の成分と零点]
\label{prop:all-components-cone}
$X$を滑らかな複素射影多様体とし、
\[
 E_X=H^0(X,\textstyle\bigwedge^2T_X),\qquad
 \operatorname{Pois}(X)=\{\pi\in E_X:[\pi,\pi]=0\}
\]
とおく。その被約化の各既約成分はスカラー作用で不変で、零を含む。
滑らかなcubicを動かす相対Poissonスキームでも、各成分は
その成分にある$(X,\pi)$に対応する$(X,0)$を含む。
\end{proposition}
\begin{proof}
Jacobi方程式は$\pi$について二次斉次である。
既約成分$Z$に対し、$\mathbb G_m\times Z$の作用像の閉包は既約で$Z$を含む。
$Z$の極大性からその閉包は$Z$であり、スカラー作用は$Z$を保つ。
$t\mapsto t\pi$を$\mathbb A^1$まで延長すると、閉性から$0\in Z$となる。
相対の場合は$t\mapsto(X,t\pi)$とすればよい。
\end{proof}

\begin{theorem}[滑らかなスカラーchart]
\label{cw:thm:chart}
$Z\ne\{0\}$を有限次元複素ベクトル空間$E$内の被約既約アフィン錐とする。
$\mathbb PZ$の滑らかな稠密アフィン開部分$T$と正則写像
$p:T\to Z\setminus\{0\}$を選べて、
\[
 S=T\times\mathbb A^1_h,\qquad
 c:S\longrightarrow Z,\quad (t,h)\longmapsto h p(t)
 \tag{C}\label{cw:eq:contraction}
\]
は次を満たす。
$S$は滑らかな既約スキーム、$c$は支配的双有理、
$c(t,0)=0$であり、$c|_{T\times\mathbb G_m}$は$Z$の稠密開部分への同型である。
とくに$\mathbb C(S)=\mathbb C(Z)$である。
\end{theorem}
\begin{proof}
$\mathbb PZ$は被約既約で、標数零なので滑らかな稠密開部分を持つ。
その非空開部分を、ある線形座標$z_0\ne0$の射影chartの中に取り、さらにアフィンに縮める。
代表を$z_0(p(t))=1$と規格化すれば、$p$は正則である。
$h\ne0$では逆写像は
\[
 z\longmapsto\bigl([z],z_0(z)\bigr)
\]
であり、像は$z_0\ne0$かつ$[z]\in T$という開集合である。
これが双有理性を含む全主張を証明する。
\end{proof}

\begin{remark}[頂点での降下はgenericな目標には不要]
\label{cw:rem:no-cone-descent}
$Z$が頂点で特異でも$S$は滑らかである。
$S$上の量子化族が$Z$へ降下するとは主張しない。
しかし一般点の模型を求めるだけなら、同じ関数体を持つ$S$上での構成で足りる。
$t$を一つ固定した曲線ではなく、\emph{$T$の全方向を含む族}を保つことが重要である。
射影Poisson空間そのものに零点はなく、式\eqref{cw:eq:contraction}の$h=0$が可換ファイバーを与える。
相対Poissonスキームについても、その射影化内の滑らかな稠密chartで
相対tautological lineを自明化すれば同じ構成ができる。
零切断内にのみ含まれる成分は、もともと可換な族として別に扱う。
\end{remark}

\begin{proposition}[ブロック正規形の有限被覆上の固定]
\label{cw:prop:normal-form-cover}
$K=\mathbb C(T)$とし、genericなPoisson構造の三重ベクトル$P$と三次式$F$が
$\overline K$上で、ある$N=3k+m$について
\[
 P=\sum_{s=1}^k\partial_{x_{s1}}\wedge\partial_{x_{s2}}\wedge\partial_{x_{s3}},
 \qquad F=\sum_{s=1}^k f_s(x_s,u)
 \tag{B}\label{cw:eq:block-form}
\]
という正規形を持つと仮定する。
$T$を縮めると、整な有限étale全射$T'\to T$が存在し、
この順序付きブロック、中心変数、座標変換を全て$T'$上で正則に選べる。
$S'=T'\times\mathbb A^1_h\to Z$は支配的でgenerically finiteであり、
$h=0$は可換なcubicの族を与える。
\end{proposition}
\begin{proof}
座標変換、各ブロックの基底、三次式の係数を未知数とすると、
\eqref{cw:eq:block-form}と座標変換の可逆性は有限個の多項式方程式・不等式である。
$\overline K$上の解の有限個の座標は$K$の有限拡大$L$に属する。
標数零なので$L/K$は分離的である。
原始元の最小多項式を取り、係数の分母と判別式を反転すると、
$T$の開部分上に関数体$L$の有限étale被覆$T'$を得る。
解の座標に現れる有限個の分母は、そのノルムをさらに反転して全て可逆にできる。
これで座標変換と式が$T'$全体へ広がる。
有限étale被覆と式\eqref{cw:eq:contraction}の合成は、$h\ne0$で有限étaleである。
\end{proof}
\begin{remark}
この命題の正規形の存在は論文AのPoisson正規形を入力とする。
その証明を本節で再検証したとは扱わない。
$k$は選んだ成分のgenericなブロック数であり、$k$が同じ成分を同一視しない。
また、各$f_s(x_s,0)$が滑らかな平面三次曲線を定めるという仮定は置いていない。
\end{remark}

\subsection{有限行列による中心三次元の代数化}

\begin{lemma}[genericな形式解の有限jetを保つ代数化]
\label{cw:lem:linear-algebraization}
$R$を有限生成整域、$K=\operatorname{Frac}(R)$、
$M(h)\in\operatorname{Mat}_{a\times b}(R[h])$とする。
$\widehat z\in K[[h]]^b$が$M\widehat z=0$を満たすなら、
任意の整数$r\ge0$について、$R$の非零元$s$と
$q(h)\in R_s[h]$、$q(0)\in R_s^\times$、および
\[
 z\in R_s[h,q^{-1}]^b,
 \qquad Mz=0,\qquad z\equiv\widehat z\pmod {h^{r+1}}
 \tag{L}\label{cw:eq:linear-approx}
\]
が存在する。合同は$K[[h]]^b$内で解釈する。
行列$M$のrankが$h=0$で一定であることは仮定しない。
\end{lemma}
\begin{proof}
$\mathcal O=K[h]_{(h)}$とし、$z_{\le r}\in K[h]^b$を$\widehat z$の$r$次切断とする。
$Mz_{\le r}$は$h^{r+1}$で割れるから
\[
 b_r=-h^{-r-1}Mz_{\le r}\in K[h]^a
\]
と置ける。形式解$w=(\widehat z-z_{\le r})/h^{r+1}$は$Mw=b_r$を満たす。
$\mathcal O\to K[[h]]$は忠実平坦である\cite[Tag 00MC]{Stacks}。
従って$\operatorname{coker}(M:\mathcal O^b\to\mathcal O^a)$における$b_r$の像は零で、
$Mv=b_r$を満たす$v\in\mathcal O^b$が存在する。
$z=z_{\le r}+h^{r+1}v$とすれば所要の合同と恒等式を得る。
$v$の共通分母を$q\in K[h]$、$q(0)\ne0$と取る。
$K$の係数に現れる有限個の分母と$q(0)$を$R$で反転すると、
式\eqref{cw:eq:linear-approx}を得る。
\end{proof}

\begin{corollary}[次数三の中心持上げの一様化]
\label{cw:cor:center-algebraization}
$\mathcal A$を$R[h]$上の次数一生成代数とし、
$\mathcal A_3,\mathcal A_4$が有限自由、$\mathcal A/(h)$が可換とする。
$\mathcal A_1$の生成元$x_1,\ldots,x_N$に対し
\[
 M_h:\mathcal A_3\longrightarrow\mathcal A_4^{\oplus N},
 \qquad z\longmapsto h^{-1}([z,x_1],\ldots,[z,x_N])
 \tag{Z}\label{cw:eq:center-matrix}
\]
は有限行列で表される。
$\mathcal A_K$の次数別$h$-進完備化に中心三次元$\widehat\Omega$があり、
$\widehat\Omega\bmod h=F$なら、前補題の局所化後に
$\Omega\in Z(\mathcal A_{R_s[h,q^{-1}]})_3$、$\Omega\bmod h=F$を取れる。
指定した任意の有限jetも保持できる。
Hilbert級数が$(1-t)^{-N}$のPBW族なら、未知数の数と方程式の数はそれぞれ
\[
 b=\binom{N+2}{3},\qquad a=N\binom{N+3}{4}
\]
であり、$h$の打切り次数によって無限に増やす必要はない。
\end{corollary}
\begin{proof}
$\mathcal A/(h)$が可換で$\mathcal A_4$が自由なので、各交換子を$h$で一意に割れる。
中心性は生成元との交換子が全て零であることと同値である。
前補題をこの行列に適用する。
\end{proof}
\begin{remark}[この代数化が証明しないこと]
\label{cw:rem:center-limit}
形式中心元の存在と、適用する周囲の代数$\mathcal A$の構成はこの系の入力である。
係数方向を縮めることは許すが、$q(0)$は可逆なので縮めた基底の零切断は全て残る。
この議論は中心元を\emph{代数的な有理式}にするものであり、
全ての形式係数を一致させること、全Poisson錐上の普遍中心元、
全次数PBW性、AS-regular性、相対saturationを証明するものではない。
\end{remark}

\begin{lemma}[正則中心商の次数別flatness]
\label{cw:lem:flat-quotient}
$D$を有限型基底、$\mathcal A$を$D$上の次数付き代数で各次数が有限局所自由とする。
全ての幾何学的ファイバーが整域で、$\Omega\in Z(\mathcal A)_3$が
各ファイバーで非零であれば、$\mathcal B=\mathcal A/(\Omega)$は各次数で有限局所自由である。
$H_{\mathcal A_s}(t)=(1-t)^{-N}$なら
$H_{\mathcal B_s}(t)=(1-t^3)/(1-t)^N$である。
\end{lemma}
\begin{proof}
各次数$d$で$\Omega:\mathcal A_{d-3}\to\mathcal A_d$は
全ての幾何学的ファイバーで単射である。
有限局所自由加群の行列の極大小行列式を用いると、この写像は局所的に分裂した単射で、
余核は有限局所自由である。従って主張とHilbert級数の式を得る。
\end{proof}

\subsection{一つの可換ファイバーから成分の一般点へ}

\begin{definition}[条件 (ER)]
\label{cw:def:ER}
本節で条件 (ER) とは次の局所実現原理を指す。
有限型複素スキーム$D$上のsmooth properな代数的dg圏と
可逆perfect bimoduleで表されるrotationの族$(\mathcal K_D,\rho_D)$が、
ある複素点$d$で滑らかなcubicの標準pairとenhanced同値なら、
$d$のpointed étale近傍で相対cubicとrotationを保つ同値が存在する。
これは現行論文Bの\texttt{rotation\_models.tex}にある
\texttt{thm:etale}を\emph{入力}とする条件である。
本節ではその形式変形論・実現スタックの証明全体を再証明しない。
Toën--Vaquié \cite[Theorem 3.6]{TV} はperfect対象のmoduliの代数性を与える入力だが、
それ単独が条件 (ER) の証明ではない。
\end{definition}

\begin{theorem}[支配的chart上の一可換点によるgeneric復元]
\label{thm:all-components-conditional}\label{cw:thm:generic}
$X^n$、$n\ge3$を滑らかなcubic、$Z\ne\{0\}$を
$\operatorname{Pois}(X)_{\mathrm{red}}$の任意の既約成分とする。
定理\ref{cw:thm:chart}と命題\ref{cw:prop:normal-form-cover}の
$S'=T'\times\mathbb A^1_h\to Z$を使う。
ある点$s_0=(t_0,0)$を含む開部分$V\subset S'$上に、
実際の次数別flatな量子化中心三次商と、これに対応する
smooth properな代数的残余pair族$(\mathcal K_V,\rho_V)$があり、
$s_0$で$(\mathcal K_{s_0},\rho_{s_0})\simeq(\Ku(X),\rho_X)$とする。
条件 (ER) が成立すれば、$V$の稠密開部分の各幾何学的ファイバーは
滑らかな可換cubicのrotation付きpairを持つ。
特に$\mathbb C(Z)$のある有限分離拡大上で、その構成したgenericな量子化の可換模型が存在する。
\end{theorem}
\begin{proof}
条件 (ER) のétale近傍の像$W\subset V$は$s_0$を含む開集合である。
$S'$は既約だから$W$は稠密で、$W\cap\{h\ne0\}$は非空である。
$S'\to Z$は$h\ne0$で$Z$の稠密開部分への有限étale射なので、その像も開で稠密である。
étale近傍の支配的な成分のgeneric residue fieldは
$\mathbb C(S')$の有限分離拡大であり、
$\mathbb C(S')/\mathbb C(Z)$も有限分離なので結論を得る。
\end{proof}
\begin{remark}[量化と降下]
\label{rem:all-components-gap}
定理は\emph{構成した}量子化族についての条件付き存在であり、
同じPoisson初項を持つ任意の非同値な量子化を同一視しない。
$Z$上に量子化そのものが既に降下している場合には、
$Z$の稠密開の各幾何学的ファイバーについて結論を読むこともできる。
そうでなければ、有限支配被覆上のgeneric模型の存在として読む。
また、多ブロック生成性はこの方法で独立の前提にする必要はないが、
その代わり全体の残余圏についてsmooth proper性を示す必要がある。
\end{remark}

\begin{proposition}[有限対角証明書によるsaturationの開近傍]
\label{cw:prop:diagonal-certificate}
$R$をNoether環、$\Lambda$を$R$上のdg代数で$\Lambda\in\Perf(R)$とする。
$\Lambda^e=\Lambda^{\mathrm{op}}\otimes_R^{\mathbf L}\Lambda$と置く。
実際に
\[
 P\in\Perf(\Lambda^e),\qquad \alpha:P\to\Lambda
 \tag{D}\label{cw:eq:diagonal-certificate}
\]
が構成され、ある点$s$で$\alpha_s$が擬同型なら、$s$のZariski開近傍で
$\Lambda$はsmooth properである。この結論は導来底変更と両立する。
\end{proposition}
\begin{proof}
$\Lambda^e$は$R$-perfectなので、$P$も$R$-perfectである。
従って$C=\operatorname{Cone}(\alpha)$は$R$-perfectで、$C\otimes_R^{\mathbf L}\kappa(s)=0$。
有限個の項からなる有限局所自由複体で表すと、ファイバーの非零コホモロジーの存在座は閉である。
よって$s$の開近傍で$C=0$であり、$\Lambda\simeq P\in\Perf(\Lambda^e)$を得る。
perfect性と擬同型の導来引戻しの両立性から最後の主張も従う
\cite[Tag 0BDK]{Stacks}。
\end{proof}
\begin{remark}
ここで必要なのは\emph{既に持ち上がった}式\eqref{cw:eq:diagonal-certificate}である。
特殊ファイバーの対角がperfectであるだけから、その持上げを結論してはいない。
また$R$-perfectとは底上の複体の性質であって、
有限自由$R$-加群だけからなる\emph{厳密なdg代数模型}の存在を仮定していない。
\end{remark}

\subsection{次元合同類と残る未解決条件}

\begin{proposition}[Serre関手からのrotationの復元]
\label{prop:all-components-rotation-serre}
smooth properなdg圏と自己同値$\rho$が
\begin{equation}\label{eq:all-components-serre}
 \rho^3\simeq[2],\qquad S\simeq\rho^{1-n}[n]
\end{equation}
を満たすとする。
$n\not\equiv1\pmod3$なら、$3a+(1-n)b=1$を満たす整数$a,b$に対して
\begin{equation}\label{eq:all-components-rotation}
 \rho\simeq S^b[2a-nb].
\end{equation}
特に$n\equiv0\pmod3$なら$\rho\simeq S[-n/3]$、
$n\equiv2\pmod3$なら$\rho\simeq S^{-1}[(n+4)/3]$である。
$n\equiv1\pmod3$ならこれらの式から$S\simeq[(n+2)/3]$を得るが、
rotationの同型類は一般には決定されない。
\end{proposition}
\begin{proof}
$\gcd(3,n-1)=1$のときBézout等式を用いて
$\rho=(\rho^3)^a(\rho^{1-n})^b$とする。
$n-1=3m$の場合は$S\simeq[n-2m]$である。
この後者の場合の追加の曖昧さは、論文Bの四次元Zhang twist反例で実際に現れる。
\end{proof}

\begin{example}[全体が滑らかでも純ブロックは特異になり得る]
\label{cw:ex:singular-block}
$\mathbb P^6$内のcubic fivefold
\[
 F=x^3+y^3+u z^2+a^3+b^3+c^3+u^3=0
\]
と、三重ベクトル
$P=\partial_x\wedge\partial_y\wedge\partial_z+
\partial_a\wedge\partial_b\wedge\partial_c$を考える。
$F$は滑らかで、$\iota_{dF}P$は一般階数4のPoisson構造である。
しかし第一ブロックの$u=0$での純三次式$x^3+y^3$は特異である。
\end{example}
\begin{proof}
勾配の零点では$x=y=a=b=c=0$、$uz=0$、$z^2+3u^2=0$なので、$z=u=0$である。
よってアフィン錐の特異点は原点のみである。
二ブロックは互いに可換で$u$はCasimir、各ブロックはJacobian括弧なのでJacobi恒等式を満たす。
$u\ne0$のchartで両ブロックが非零となる点は存在し、そのとき階数は4である。
$x^3+y^3=0$は$[0:0:1]$で特異である。
\end{proof}
\begin{remark}
この例が、そのPoisson既約成分の\emph{一般点}でも特異ブロックを強制するとは主張しない。
全点に対する議論で、全体の滑らかさから各純ブロックの楕円性を仮定してはならないという注意である。
\end{remark}

\begin{remark}[全非可換成分との区別]\label{rem:all-components-nc-scope}
任意のAS-regular中心三次商の全点での復元は既存のZhang twist反例により偽である。
個別の反例点だけではPoisson既約成分の一般点での復元は否定されない。
本節は$n\ge3$を対象とし、曲線・cubic surfaceへ圏論的入力を無条件に拡張しない。
\end{remark}

各Poisson成分の一般点を扱う際には、錐の頂点へ量子化族を降下させる必要はない。
ただし、本節のチャートから可換復元へ進むには、
その上で\emph{実際の}相対PBW代数と中心正則商を構成し、
残余圏のproper模型、有限対角証明書、rotationと相対Orlov比較を確立する必要がある。
形式中心の存在が分かれば中心元だけの代数化は補題\ref{cw:lem:linear-algebraization}で処理できるが、
これを残余圏のsaturationの証明と混同しない。

\subsection{有限計算と証明書の範囲}

添付コードは、$\operatorname{Gr}(2,5)$の標準chartと錐の縮約、
上の特異ブロック例の勾配・Jacobi恒等式、中心行列の線形代数化の例を有理数係数で確認する。
さらに、三変数Fermat量子化の中心元については、
有限特殊化ではなく\emph{$\mathbb Q[h]$上の自由代数恒等式}を証明書として収録した。
$r_1=xy-yx-hz^2$、$r_2=yz-zy-hx^2$、$r_3=zx-xz-hy^2$、
$D=x^3+y^3+z^3$、$T=xyz+xzy+yxz+yzx+zxy+zyx$とすると、
\begin{align*}
 [3D-h^2T,x]={}&-3r_1y^2+h^2r_1zx-h^2xr_1z-3yr_1y-h^2zr_1x
 +h^2xzr_1-3y^2r_1\\
 &+2h^2r_2x^2-3hr_2yz-h^2xr_2x-h^2x^2r_2+3hyzr_2\\
 &+3r_3z^2+3zr_3z+3z^2r_3.
\end{align*}
語ごとの係数比較でこの式を確認でき、巡回置換で$y,z$との交換子も得る。
従って$(D/3-h^2T/9)$はこの三変数商で中心的である。
この証明書は全混合項の中心公式、PBW性、圏同値の検証ではなく、
標準的なFermat入力について特殊値$h=0$を除外しない普遍的検証である。


\part{相対proper模型と有限対角の検出}
\section{相対Beilinson模型から得られるproperな族}
\label{sec:relative-proper-model}

本節は、一般の量子化族の相対saturationに向けた部分結果である。
以下の相対Beilinson仮定のもとで、超曲面圏と残余圏を
実在する基底上のproper dg代数で表す。
この結論にsmoothness、homotopical finite presentation、
または一般の有限パラメータでの可換復元を含めない。
相対Beilinson仮定そのものと、最後に残る対角加群の有限性は別の問題である。

\subsection{仮定と基底変更の意味}

\(R\) をNoetherな可換 \(\CC\)-代数、\(N=n+2\ge5\) とする。
\(A=\bigoplus_{d\ge0}A_d\) は \(A_0=R\) を満たす次数付き代数で、
有限生成の \(A_1\) により生成され、各 \(A_d\) は有限射影 \(R\)-加群とする。
中心元 \(\Omega\in A_3\) は非零因子で、
\(B=A/(\Omega)\) の各次数成分も有限射影 \(R\)-加群とする。
したがって
\[
 0\longrightarrow A(-3)\xrightarrow{\Omega}A
   \longrightarrow B\longrightarrow0
 \tag{Q}\label{eq:relative-quotient-sequence}
\]
は任意の可換環への底変更後にも完全である。
\(\QGr\) は全次数付き右加群の圏をtorsion加群で割ったGrothendieck圏、
\(\OO_A(i)\), \(\OO_B(i)\) は構造加群のシフトの像を表す。

\begin{definition}[本節の相対Beilinson仮定]\label{def:relative-beilinson-hyp}
 必要とする各Noetherな底変更 \(R\to R'\) に対し、次を仮定する。
 \begin{enumerate}[label=(RB\arabic*)]
 \item 次数付き加群の導来局所化は
 \[
 D(\QGr A_{R'})\simeq
 D(\operatorname{Gr}A_{R'})/
 D_{\operatorname{Tors}}(\operatorname{Gr}A_{R'}),
 \]
 および \(B_{R'}\) に対する同じ表示を持つ。
 これらのdg enhancementを用い、次数シフトと導来係数拡大が
 局所化へ降りるものとする。
 \item \(G_{R'}=\bigoplus_{i=0}^{N-1}\OO_{A_{R'}}(i)\) はcompact generatorで、
 \[
 \REnd(G_{R'})\simeq T_{R'},\qquad
 e_jTe_i=A_{j-i}\quad(0\le i,j\le N-1),
 \tag{B}\label{eq:relative-beilinson-algebra}
 \]
 である。ただし \(A_d=0\) for \(d<0\) とし、右辺は次数零の三角代数である。
 積はHomの合成によるものとし、この側のBeilinson代数を \(T=\nabla A\) と書く。
 oppositeに関する行列の規約はこの等式で固定する。
 同値は自然な \(T\otimes_RR'\to\REnd(G_{R'})\) によって与えられる。
 \end{enumerate}
\end{definition}

ここで(RB2)は、体上のAS-regular性を述べる条件より強い。
各fiberにおけるBeilinson型定理を引用するだけでは、
この普遍的な底変更を含む条件を検証したことにはならない。
通常の相対構成では、有限長の線形Koszul分解、
その双対複体、torsion局所化の計算を用いて(RB2)を検証する。
本節ではこの検証を、任意の量子化族について完了したとは主張しない。

\subsection{超曲面圏のproper模型}

\begin{proposition}[相対proper模型]\label{prop:relative-proper-model}
 定義\ref{def:relative-beilinson-hyp}の条件のもとで
 \[
 E=\bigoplus_{i=0}^{N-1}\OO_B(i),\qquad
 \Lambda=\REnd_{\QGr B}(E)
 \]
 とおくと、\(E\) は \(D(\QGr B)\) のcompact generatorであり、
 \[
 D(\QGr B)^c\simeq\Perf(\Lambda)
 \]
 を得る。\(\Lambda\) は \(R\)-perfectなdg代数である。
 また、仮定の範囲の各底変更に対して、dg代数として
 \[
 \Lambda\otL_RR'\simeq
 \REnd_{\QGr B_{R'}}(E_{R'})
 \tag{BC}\label{eq:relative-proper-basechange}
 \]
 が成り立つ。
\end{proposition}

\begin{proof}
 次数付き加群のrestrictionはtorsion加群をtorsion加群へ送り、
 \(B\)-加群について \(A\)-torsionと \(B\)-torsionは一致する。
 よってrestrictionはexact、conservativeかつ直和を保つ関手
 \[
 j_*:D(\QGr B)\longrightarrow D(\QGr A)
 \]
 を誘導する。その左随伴 \(j^*\) は導来係数拡大である。
 実際、\eqref{eq:relative-quotient-sequence}によって
 \(-\otL_A B\) はtorsion cohomologyを持つ複体を同種の複体へ送り、
 導来局所化へ降りる。随伴のunitとcounitも降りる。
 また、restriction後の \(j_*j^*\) には自然なtriangle
 \[
 X(-3)\xrightarrow{\Omega}X\longrightarrow j_*j^*X
 \tag{M}\label{eq:relative-quotient-monad}
 \]
 がある。\(j_*\) の直和保存とconservativityから、
 \(j^*G=E\) はcompact generatorとなる。

 \(T=\nabla A\) は有限射影 \(R\)-加群である。
 さらに \(S=R^{\oplus N}\subset T\) とすると、
 augmentation \(T\to S\) の核 \(J\) は狭義三角行列からなる。
 \(J^{\otimes_S q}=0\) for \(q\ge N\) であり、
 normalized relative bar resolution
 \[
 \cdots\longrightarrow
 T\otimes_SJ^{\otimes_S2}\otimes_ST
 \longrightarrow T\otimes_SJ\otimes_ST
 \longrightarrow T\otimes_ST\longrightarrow T
 \]
 は有限長で、各項は有限射影 \(T\)-bimoduleである。
 したがって \(T\) は \(R\) 上smooth properである。

 \(D(\QGr A)\simeq D(T)\) のもとで次数シフト \((-3)\) を表す
 invertible \(T\)-bimoduleを \(M_{-3}\) と書く。
 これは \(T\) のsmoothnessとMorita自己同値により
 \(\Perf(T^{\mathrm{op}}\otL_RT)\) に属し、特に \(R\)-perfectである。
 \eqref{eq:relative-quotient-monad}を \(G\) に適用すると、基底上の複体として
 \[
 \Lambda\simeq
 \operatorname{cofib}\bigl(M_{-3}\xrightarrow{\Omega}T\bigr).
 \tag{C}\label{eq:relative-perfect-cone}
 \]
 よって \(\Lambda\) は \(R\)-perfectである。

 このcofiberの表示だけから積を定義してはいけない。
 \(\Lambda\) の積と単位は実際の \(\REnd_B(E)\) のものを用いる。
 同じ内容をMorita模型で述べると、\(j_*j^*\) は
 unitとmultiplicationを備えたmonadであり、
 \eqref{eq:relative-perfect-cone}のcofiberはそのmonadを表す
 \(T\)-bimoduleの代数対象である。
 この代数構造が、随伴による
 \(\REnd_B(j^*G)\simeq\RHom_A(G,j_*j^*G)\) の合成と一致する。

 最後に、(RB2)により周囲のMorita同値は底変更と両立する。
 次数シフト、中心元 \(\Omega\)、随伴のunitとmultiplicationも
 導来係数拡大から作られるので底変更と両立する。
 したがって、monadの代数対象としての底変更から
 \eqref{eq:relative-proper-basechange}が得られる。
 これは基底複体だけの同一視より強い主張である。
\end{proof}

\begin{remark}[有限性の範囲]\label{rem:relative-model-finiteness}
 命題\ref{prop:relative-proper-model}は、各Artin切断を別々に構成して
 逆極限を取る主張ではない。\(R\) 上に存在するadjunctionとmonadから、
 \(R\)-perfectな実際のdg代数を与える。
 ただし、基底上のperfect複体に移した \(A_\infty\) 積が
 有限個の高次演算だけで記述できることや、\(\Lambda\) が
 homotopically finitely presentedであることは従わない。
 特に、この命題をsmoothnessの証明と読み替えることはできない。
\end{remark}

\subsection{局所cohomologyによる確認可能な数値条件}

(RB2)を検証する際の具体的な目標は、ある可逆 \(R\)-加群 \(L\) に対する
次の構造層cohomologyである：
\[
 H^pR\Gamma(\OO_A(d))\simeq
 \begin{cases}
 A_d,&p=0,\ d\ge0,\\
 A_{-N-d}^{\vee}\otimes_RL,&p=N-1,\ d\le-N,\\
 0,&\text{otherwise}.
 \end{cases}
 \tag{H}\label{eq:relative-ambient-cohomology}
\]
この式そのものにも底変更との両立を要求する。
有限線形Koszul分解の存在はshiftの有限生成性を、
\eqref{eq:relative-ambient-cohomology}はtiltingのExt消滅を与える。
一方、compact性と全生成性を得るには、torsion局所化の導来関手が
必要な直和を保つこと、または同等のcompact generationの証明が別途必要である。

\eqref{eq:relative-ambient-cohomology}が分かれば、
\(0\le i,j\le N-1\), \(d=j-i\) に対して
\[
 H^p\RHom_B(\OO_B(i),\OO_B(j))\simeq
 \begin{cases}
 B_d,&p=0,\ d\ge0,\\
 A_{3-N-d}^{\vee}\otimes_RL,
     &p=N-2,\ 0\le3-N-d\le2,\\
 0,&\text{otherwise}
 \end{cases}
 \tag{H$_B$}\label{eq:relative-hypersurface-cohomology}
\]
が \eqref{eq:relative-quotient-monad}から従う。
特にこの窓のHomは基底上有限射影cohomologyを持ち、その振幅は
\([0,N-2]=[0,n]\) に含まれる。
この式は積の全情報を代替しない。

\begin{lemma}[fiberwise Koszul性から得られる有限分解]
\label{lem:relative-koszul-resolution}
 上の次数別有限射影性を仮定し、各幾何学的fiber \(A_s\) が
 Koszulで、augmentation加群 \(k(s)\) の射影次元が \(N\)、
 Hilbert級数が \((1-t)^{-N}\) であるとする。
 このとき有限射影 \(R\)-加群 \(W_p\) を用いる普遍的に完全な線形分解
 \[
 0\to W_N\otimes_RA(-N)\to\cdots\to
 W_1\otimes_RA(-1)\to A\to R\to0
 \tag{K}\label{eq:relative-linear-resolution}
 \]
 が存在し、\(\rank W_p=\binom Np\) である。
\end{lemma}

\begin{proof}
 \(A_+\) の次数別成分は有限射影である。
 \(W_1=A_1\) として、射影加群 \(W_1\otimes_RA(-1)\) から \(A_+\) への射を取る。
 各fiberで全射なので、各次数でNakayamaの補題により全射である。
 像の次数別成分が射影であるため、核も次数別有限射影であり、
 核の形成は任意の底変更と両立する。
 同じ操作をsyzygyに反復し、\(p\) 次syzygyを生成する次数 \(p\) の成分を
 \(W_p\) とする。fiberのKoszul性により、この成分がsyzygy全体を生成する。
 したがって各段階で同じ射影性と底変更の主張が成り立つ。
 第 \(N\) 段階の次の核はすべてのfiberで零なので、次数ごとのNakayamaにより零である。
 fiberのKoszul Betti数はHilbert級数から \(\binom Np\) となる。
 各次数の短完全列は \(R\)-加群として分裂するので、分解は普遍的に完全である。
\end{proof}

この補題は有限なKoszul分解を構成するが、
各fiberのAS-regular性だけからrelative tailsの性質をすべて得たという主張ではない。
双対Koszul複体と導来torsion局所化を含めて
\eqref{eq:relative-ambient-cohomology}および(RB2)を検証する必要が残る。

\subsection{残余圏のpropernessとrotationの開性}

さらに \(\OO_B,\ldots,\OO_B(N-4)\) が相対例外列で、
そのHomと底変更が両立すると仮定する。
例えば \eqref{eq:relative-hypersurface-cohomology}はこの条件を与える。
\[
 \cK_R=\langle\OO_B,\ldots,\OO_B(N-4)\rangle^\perp
 \subset\Perf(\Lambda)
\]
とおく。

\begin{proposition}\label{prop:relative-residual-proper}
 \(\cK_R\) のinclusionは左随伴 \(p\) を持ち、
 \(G_{\cK}=p(E)\) はcompact generatorである。
 \(\Gamma_R=\REnd(G_{\cK})\) は \(R\)-perfectで、
 \(\cK_R\simeq\Perf(\Gamma_R)\) は底変更と両立する。
 また \(\rho=L_{\OO_B}\circ(1)\) は \(\cK_R\) の自己関手を与え、
 底変更と両立する。
\end{proposition}

\begin{proof}
 propernessにより各相対例外対象 \(F\) に対し、
 \(\Perf(R)\to\Perf(\Lambda)\), \(V\mapsto V\otimes_RF\) は
 左右随伴 \(\RHom(-,F)^\vee\), \(\RHom(F,-)\) を持つ。
 その左mutationは \(F^\perp\) への左projectionを与える。
 \(a=N-3\), \(F_i=\OO_B(i)\) とおくと
 \[
 p=L_{F_0}\circ L_{F_1}\circ\cdots\circ L_{F_{a-1}}
 \]
 は \(\cK_R\) への左projectionである。
 実際、右側から順に適用し、\(j>i\) に対する
 \(\RHom(F_j,F_i)=0\) を用いると、既に得た直交性が保存される。
 各evaluation triangleはcompact対象の有限なtriangleであり、
 \(p|_{\cK_R}=\Id\) なので、\(p(E)\) は残余圏を生成する。
 そのEnd複体のpropernessおよび底変更は、
 元のproper dg模型とこれらのevaluation triangleから従う。

 \(F\in\cK_R\) とする。\(1\le j\le a-1\) なら
 \(\RHom(\OO_B(j),F(1))=\RHom(\OO_B(j-1),F)=0\) であり、
 \(\RHom(\OO_B(j),\OO_B)=0\) である。
 \(L_{\OO_B}\) のevaluation triangleから \(\rho(F)\in\cK_R\) が従う。
 同じtriangleから底変更との両立も従う。
\end{proof}

\begin{remark}[両側admissibilityは追加条件]
 proper性だけから \(\cK_R\) のinclusionが右随伴も持つとは結論しない。
 同様に、\(L_{\OO_B}\) または \(\rho\) が右随伴を持つことも、
 上の証明だけからは従わない。
 例外対象への関手の左右随伴が存在することと、
 その右直交圏の両側admissibilityは別の主張である。
\end{remark}

\begin{lemma}[随伴が与えられた場合のrotationの開性]
\label{lem:relative-rotation-open-conditional}
 上の状況で、さらに \(\rho\) がcompact対象上の \(R\)-linear右随伴
 \(\rho^R\) を持ち、そのunitとcounitを含めて底変更と両立すると仮定する。
 ある点 \(s\in\Spec R\) で \(\rho_s\) が自己同値ならば、
 \(s\) のZariski近傍へ制限することで \(\rho\) は自己同値となる。
\end{lemma}

\begin{proof}
 unitとcounitを生成子 \(G_{\cK}\) 上で評価した二つのconeはcompact対象である。
 これらからの \(\RHom(G_{\cK},-)\) は \(R\)-perfectで、
 fiber \(s\) で零になる。
 perfect複体のsupportの閉性により、両方が零となる \(s\) の開近傍がある。
 そこで生成子性を用いると二つのconeは零であり、
 unitとcounitが生成子上、従ってそのthick閉包全体で同型となる。
\end{proof}

上の右随伴の追加仮定は、例えば底変更と両立するrelative Serre dualityを
周囲の圏に別途与えれば検討できるが、本節のproper模型からは導かない。
またrotationの可逆性が分かった場合でも、そのkernelが
\(\Gamma_R^{\mathrm{op}}\otL_R\Gamma_R\) 上perfectであることは別の条件である。
対角加群のperfect性が未証明なら、恒等関手についてさえこの条件は未証明である。

\subsection{smoothnessに十分な有限の証明書と残る問題}

\begin{lemma}[対角の有限証明書]\label{lem:relative-diagonal-certificate}
 \(\Gamma\) を \(R\)-perfectなdg代数とし、\(\Gamma^e=\Gamma^{\mathrm{op}}\otL_R\Gamma\)
 と書く。ある \(P\in\Perf(\Gamma^e)\) と射 \(P\to\Gamma\) が存在し、
 点 \(s\) での導来fiberがquasi-isomorphismならば、\(s\) の開近傍で \(\Gamma\) はsmoothとなる。
\end{lemma}

\begin{proof}
 \(\Gamma^e\) は \(R\)-perfectなので、有限cellと直和因子で作られる \(P\) も
 \(R\)-perfectである。\(\Cone(P\to\Gamma)\) も \(R\)-perfectであり、
 \(s\) でのfiberが零なので、ある開近傍で零である。
 従ってその近傍では対角 \(\Gamma\simeq P\) が \(\Gamma^e\) 上perfectである。
\end{proof}

\begin{remark}[現時点で未解決の点]\label{rem:relative-saturation-gap}
 fiber \(\Gamma_s\) のsmoothnessは、そのfiber上での有限cell対角分解を与える。
 しかし、その有限cell、微分、retract、および対角への射を
 \(R\) の近傍へ同時に持ち上げた証明書はまだ構成していない。
 各Artin切断でのsmoothnessだけではこの構成を代替できない。
 Toën \cite[Lemma 3.14]{ToenFinitude} が述べるのは、properなdg代数の
 「smoothかつcohomological dimension \(\le d\)」という部分前層の
 quasi-compactnessであり、一般の開性ではない。
 この補題をsmoothnessの開性として使わない。

 また、Toën--Vaquié \cite[Theorem 3.6]{TV} の有限型dg圏のobject moduli定理を、
 proper性だけが分かっている \(\Gamma\) に適用することはできない。
 特にpseudo-perfect加群とcompactなperfect加群の同一視にはsmoothnessが必要である。
 既知の飽和圏内で生成子を動かして比較する方法にも、
 生成性の開性および比較moduliの有限表示性の証明が別途必要となる。
 これらのmoduliによる迂回を完成した代数化定理として扱わない。
\end{remark}


\subsection{有限例外列による相対smoothnessの還元}
現行論文の「相対saturationへの二つの還元」節の命題により、
相対properなdg圏$\mathcal C_R=\Perf(\Lambda_R)$が
有限個の相対例外対象と残余圏$\mathcal K_R$に分解するとき、
$\mathcal C_R$と$\mathcal K_R$の相対smoothnessは同値となる。
gluing bimoduleは$R$-perfectで、残余圏がsmoothなら
対角のperfect性により実際にperfect bimoduleとなる。
従って求める証明書は$\Lambda_R$の対角について一つあればよい。
これは(RB1)--(RB2)を仮定する還元であって、相対smoothnessそのものの
証明ではない。閉ファイバーのBeilinson対角分解を全族へ
持ち上げる問題が残る。詳細な証明は現行論文
の同節に置く。

graded matrix factorizationの内部シフトは
中心元の次数が三ならば常に$(1)^3\simeq[2]$を満たす。
これは相対Orlov同値が与えられればrotationの自己同値性と
相対三次周期性を保証する。ただし相対Orlov同値自体は
依然として未構成であり、形式的周期性からその存在は導けない。


\section{PIの中心幾何から得られる有限生成子}
\label{sec:notes-pi-residue-generator}

有限位数のSklyaninブロックでは、本文の外部積生成定理とは別に、
中心の零sliceのAzumaya幾何から全MF圏のpropernessと有限生成子を直接得られる。
この証明は \(\gcd(|\sigma_i|,3)=1\) を仮定しない。

\paragraph{設定.}
\(r\ge1\)、\(m\ge0\)、\(h_i^6\ne-27\) とし、
\[
 \begin{aligned}
 A&=\bigotimes_{i=1}^rS_{h_i}\otimes_\CC T,
 &T&=\CC[u_1,\ldots,u_m],\\
 W&=\sum_{i=1}^r g_i+P,
 &R&=A/(W)
 \end{aligned}
\]
と置く。ここで \(g_i=G_{h_i}\)、\(P\in T_3\) であり、
\(m>0\) の場合は \(T/(\partial_1P,\ldots,\partial_mP)\) が有限次元、
\(m=0\) の場合は \(T=\CC,P=0\) とする。
全ての \(\sigma_i\) は有限位数 \(q_i\) を持つものとする。
\(h_i=0\) は多項式環のfactorで、\(q_i=1\) と解釈する。
周囲の代数 \(A\) はnoetherian AS-regularであり、\(W\) は中心正則元である
（本文のテンソル積の正則性の命題）。
従ってMFと特異圏の比較 \cite[Theorem 5.8]{CCKM} を用いられる。
以下ではこれらの圏を冪等完備化し、\(k=R/R_{>0}\) は特異圏での像、
またはその有限自由stabilizationも表す。

\(s=r\) (\(m=0\))、\(s=r+1\) (\(m>0\)) と置き、
\[
 \Gamma_s=\ZZ^s/\langle3e_i-3e_1:2\le i\le s\rangle,
 \qquad w=3e_i
\]
とする。各 \(S_{h_i}\) の生成元の次数を \(e_i\) とし、
中心変数があれば全ての \(u_a\) の次数を \(e_s\) とする。
通常次数は \(\Gamma_s\to\ZZ,\ e_i\mapsto1\) による粗化であり、このとき \(w=3\) である。
この写像は有限の核を持ち、正の通常次数を与えるため、
\(\Gamma_s\)-斉次自由分解と斉次射を用いる同じMF--特異圏比較が適用できる。

\begin{proposition}[PIの場合のpropernessと残余体生成]
\label{prop:notes-pi-residue-generator}
上の仮定の下で \(\MF^{\Gamma_s}(A,W)\) と \(\MF^\ZZ(A,W)\) はproperである。
特異圏において
\[
 \begin{aligned}
 \Dsg^{\Gamma_s}(R)&=\thick(K_{\Gamma_s}),\\
 K_{\Gamma_s}&=
 \bigoplus_{(a_1,\ldots,a_s)\in\{0,1,2\}^s}
 k(a_1e_1+\cdots+a_se_s)
 \end{aligned}
 \tag{N1}\label{eq:notes-pi-gamma-generator}
\]
および
\[
 \Dsg^\ZZ(R)=\thick(K_{\ZZ}),\qquad
 K_{\ZZ}=k\oplus k(1)\oplus k(2)
 \tag{N2}\label{eq:notes-pi-z-generator}
\]
が成立する。
\end{proposition}

\begin{proof}
\paragraph{一次資料から用いる中心の事実.}
\(Z_i=Z(S_{h_i})\)、\(Y_i=\Spec Z_i\) とし、\(o_i\) を斉次頂点とする。
\(h_i\ne0\) のとき、\(S_{h_i}\) が有限生成可換中心 \(Z_i\) 上有限であることは
\cite[Theorem 3.7(c), p.~38]{SmithTate} による。
さらに \cite[Theorem 1.3(1),(2a)]{WWY3} は
\[
 V(g_i)\setminus\{o_i\}\subset Y_i^{\mathrm{sm}}
       =\operatorname{Az}(S_{h_i}),\qquad
 V(g_i)\setminus\{o_i\}\ \text{は滑らか}
 \tag{N3}\label{eq:notes-pi-zero-slice}
\]
を与える。原論文の三次中心元 \(g\) と \(g_i\) は非零定数倍なので、
零sliceについて同じ主張になる。
従って \(y_i\in V(g_i)\setminus\{o_i\}\) において
\[
 \widehat{(Z_i)_{y_i}}\simeq\CC[[g_i,a,b]],\qquad
 S_{h_i}\otimes_{Z_i}\widehat{(Z_i)_{y_i}}
       \simeq M_{q_i}(\CC[[g_i,a,b]]).
 \tag{N4}\label{eq:notes-pi-matrix-completion}
\]
実際、滑らかな三次元中心と滑らかな零sliceから \(dg_i\ne0\) が従い、
\(g_i\) を正則パラメータに選べる。
完全局所環上のAzumaya代数は、剰余体 \(\CC\) 上でsplitし、
行列冪等元の持ち上げによって行列環になる。
\(h_i=0\) でも、行列sizeを一とし、Fermat三次式の勾配を用いれば
\eqref{eq:notes-pi-zero-slice}--\eqref{eq:notes-pi-matrix-completion} が成立する。

\paragraph{行列Homの支持.}
\(Z=\bigotimes_i Z_i\otimes T\) は有限型可換代数であり、\(A\) は \(Z\) 上有限である。
有限自由factorization \(Q,Q'\) の次数を忘れた二周期行列Hom複体を
\[
 \overline C=\Hom_A^{\ZZ/2}(Q,Q'),\qquad
 \partial(f)=d_{Q'}f-(-1)^{|f|}fd_Q
\]
と書く。その二つの項は有限生成 \(Z\)-加群なので、
\(\overline H^\epsilon=H^{\bar\epsilon}(\overline C)\) も有限生成である。
ブロックEuler導分と中心偏微分により
\[
 d_Q\delta_i(d_Q)+\delta_i(d_Q)d_Q=3g_i\Id_Q,\qquad
 d_Q\partial_{u_a}(d_Q)+\partial_{u_a}(d_Q)d_Q
       =(\partial_aP)\Id_Q.
\]
これらをHomへ合成することで、\(g_i\) と \(\partial_aP\) は
\(\overline H^\epsilon\) を殺す。
Jacobian idealのradicalは \((u_1,\ldots,u_m)\) なので
\[
 \operatorname{Supp}_Z\overline H^\epsilon
 \subset V(g_1,\ldots,g_r,u_1,\ldots,u_m).
 \tag{N5}\label{eq:notes-pi-support}
\]

右辺の閉点 \(y=(y_1,\ldots,y_r,0)\) で、ある \(y_i\ne o_i\) とする。
他のSklyanin factorのテンソル積をその中心点で完備化した代数を \(B\) とすると、
\eqref{eq:notes-pi-matrix-completion} と中心上の完備テンソル積により
\[
 \widehat A_y\simeq
 M_{q_i}\!\left(B[[g_i,a,b,u_1,\ldots,u_m]]\right),
 \qquad W=g_i+\sum_{j\ne i}g_j+P(u).
 \tag{N6}\label{eq:notes-pi-local-matrix}
\]
ここでは次数を忘れてよい。
右辺の形式変数 \(g_i\) に関する成分ごとの導分を \(\partial_{g_i}\) とすれば、
\(\partial_{g_i}W=1\) である。
任意の完備化したfactorizationの微分 \(d\) に対して
\[
 h=\partial_{g_i}(d),\qquad dh+hd=\Id
 \tag{N7}\label{eq:notes-pi-contraction}
\]
となる。従って局所的なungraded MFは実際にcontractibleである。
行列Homにも \(\mathcal H(f)=h_{Q'}f\) による
\(\partial\mathcal H+\mathcal H\partial=\Id\) がある。
中心完備化はflatなのでHom cohomologyと可換であり、
局所環上ではfaithfully flatなので
\((\overline H^\epsilon)_y=0\) が従う。

こうして \(\overline H^\epsilon\) の支持は唯一の頂点
\(o=(o_1,\ldots,o_r,0)\) に含まれる。
\(Z\) は有限型可換 \(\CC\)-代数であるから、
\(\overline H^0\oplus\overline H^1\) は有限次元 \(\CC\)-ベクトル空間である。
この議論には、完備化した他のfactorのglobal dimensionを求める必要はない。

\paragraph{周期性とproperness.}
二周期Homを、偶数部分の微分が内部次数零、奇数部分の微分が内部次数 \(w\)
となるように次数付ける。
具体的に \(Q=(Q_1\to Q_0\to Q_1(w))\) の規約では
\[
 \begin{aligned}
 \overline C^0&=\underline{\Hom}_A(Q_0,Q'_0)
       \oplus\underline{\Hom}_A(Q_1,Q'_1),\\
 \overline C^1&=\underline{\Hom}_A(Q_0,Q'_1(w))
       \oplus\underline{\Hom}_A(Q_1,Q'_0).
 \end{aligned}
\]
とする。ここで \(\underline{\Hom}\) は全ての内部次数のHomの直和である。
\(\overline H^\epsilon\) はこの内部次数を継承し、\(\Lambda=\Gamma_s,\ZZ\) のいずれでも
\[
 \begin{aligned}
 H^{2n}\Hom_{\MF^\Lambda}(Q,Q')
      &=(\overline H^0)_{nw},\\
 H^{2n+1}\Hom_{\MF^\Lambda}(Q,Q')
      &=(\overline H^1)_{nw}.
 \end{aligned}
 \tag{N8}\label{eq:notes-pi-hom-periodicity}
\]
\(w\) は無限位数であり、有限次元の \(\overline H^\epsilon\) は有限個の内部次数しか持たない。
従ってMFのHomの全cohomologyの直和は有限次元である。
有限個のconeと直和因子にもこの結論は保たれるので、冪等完備化したMF圏はproperである。

\paragraph{残余体による生成.}
任意の有限自由factorization \(Q\) を取る。
中心の正次数斉次生成元を \(z_1,\ldots,z_t\) とすると、
全内部次数を含むstable endomorphism加群が頂点に支持されることから、十分大きい \(d\) に対して
\(z_j^d\Id_Q\) は全て零ホモトピックである。
このclassの内部次数は \(d\deg z_j\) であり、
Homの対応する斉次成分を取ることで、
\[
 z_j^d:Q(-d\deg z_j)\longrightarrow Q
\]
のnull-homotopyを次数を保った形に選べる。
中心Koszul複体とのtensor積はこれらのhomotopyを保つので、
これらの乗法のconeを順次取った対象には \(Q\) が直和因子として入る。

\(\Dsg^\Lambda(R)\) で \(Q\) を表す有限生成moduleを \(M\) とする。
今のconeの反復は、boundedな中心Koszul複体
\(K_R(z_1^d,\ldots,z_t^d)\otimes_R M\) の特異圏での像である。
そのcohomologyは \((z_1^d,\ldots,z_t^d)\) に殺される。
\(R\) は \(Z\) 上有限で、\(Z/(z_1^d,\ldots,z_t^d)\) は有限次元なので、
これらのcohomologyも有限次元である。
有限次元の次数付き \(R\)-moduleは正の通常次数によるfiltrationを持ち、
各商は有限個の \(k(\lambda)\) の直和である。
従って全ての \(Q\) は \(\{k(\lambda)\}_{\lambda\in\Lambda}\) で生成される。
冪等完備化についても同じ結論になる。

最後に \((w)=[2]\) であり、
\[
 \Gamma_s/\ZZ w\simeq(\ZZ/3)^s,\qquad \ZZ/3\ZZ\simeq\ZZ/3
\]
なので、それぞれの剰余類の代表だけで十分である。
これが \eqref{eq:notes-pi-gamma-generator} と \eqref{eq:notes-pi-z-generator} を与える。
\end{proof}

\paragraph{次の計算への使い方.}
細分次数では \(3^s\) 個、通常次数では \(k,k(1),k(2)\) の三つのstabilizationを用意すればよい。
その直和 \(K_\Lambda\) に対して
\[
 B_\Lambda=\REnd_{\Dsg^\Lambda(R)}(K_\Lambda),\qquad
 \MF^\Lambda(A,W)\simeq_{\mathrm{Morita}}\Perf(B_\Lambda)
\]
となる。自己Homはstabilization後のMFのdg模型で取り、\(H^*(B_\Lambda)\) は有限次元である。
本文の有限位数多ブロック比較と組み合わせると、この模型はsmooth properでもある。
この生成子の記述は、外部積による全生成性の証明とは独立に、
全MF圏を具体的な有限個の対象の自己Homから調べる入口を与える。
既知の中心幾何からの上の帰結について、先行研究上の新規性は未確認である。


\section{設定と区別すべき三つのperfect性}
$R$を可換環、$\Gamma$を単位的な$R$-dg代数とし、
\[
 \Gamma^e=\Gamma^{\mathrm{op}}\otL_R\Gamma,\qquad
 \Delta_\Gamma=\Gamma\in D(\Gamma^e)
\]
とおく。$\Gamma\in\Perf(R)$をproper性の入力とする。
ここで\emph{相対smoothness}とは$\Delta_\Gamma\in\Perf(\Gamma^e)$をいう。
$\Gamma$のunderlying $R$-複体のperfect性、$\Gamma\otL_Rk$の対角perfect性、$\Delta_\Gamma$の$\Gamma^e$-perfect性は、論理的に異なる条件である。

非楕円型の節点・尖点量子化では$R=\mathbb C[r,s]_{(r,s)}$、$\Gamma_R=\operatorname{REnd}(G)$が実在するproper残余dg模型である。特殊ファイバー$\Gamma_R\otL_Rk$は滑らかなcubic fivefoldのKuznetsov成分を表しsmoothであるが、この事実だけで$\Gamma_R$の相対smoothnessを結論したことにはならない。

\section{有限bar段階による対角判定}
$\operatorname{Bar}_\bullet(\Gamma/R)$を、導来テンソルで形成する相対二側bar解消とする。必要ならcofibrantな模型を取る。$P_d$をその$d$-skeletonの有限部分実現とし、増大写像を$\epsilon_d:P_d\to\Delta_\Gamma$と書く。

\begin{proposition}[有限段階のretract判定]\label{prop:bar}
$\Gamma\in\Perf(R)$と仮定する。このとき$P_d\in\Perf(\Gamma^e)$で、
\[
 \operatorname{hocolim}_{d}P_d\simeq\Delta_\Gamma.
\]
従って次の条件は同値である。
\begin{enumerate}
\item $\Gamma$が$R$上homologically smoothである。
\item ある$d\ge0$に対し$D(\Gamma^e)$で$\epsilon_d:P_d\to\Delta_\Gamma$が右逆
$s_d:\Delta_\Gamma\to P_d$を持つ。
\end{enumerate}
\end{proposition}
\begin{proof}
barの各simplicial項は、導来双加群として
\[
 \Gamma\otL_R(\Gamma)^{\otL_R d}\otL_R\Gamma
 \simeq \Gamma^e\otL_R (\Gamma)^{\otL_R d}
\]
という形である。$\Gamma\in\Perf(R)$なので有限テンソル冪は$R$-perfectであり、その$\Gamma^e$への拡大は$\Gamma^e$-perfectである。有限skeletonはこれらの有限個のshift、cone、retractで作られるから$P_d$はperfectである。二側barのaugmentationは実現後に擬同型となる。

右逆が存在すれば$\Delta_\Gamma$は$P_d$のretractなのでperfect。逆に$\Delta_\Gamma$がperfectなら$D(\Gamma^e)$でcompactである。したがって
\[
 \operatorname{Hom}_{D(\Gamma^e)}
 (\Delta_\Gamma,\operatorname{hocolim}_d P_d)
 \cong \varinjlim_d\operatorname{Hom}_{D(\Gamma^e)}(\Delta_\Gamma,P_d)
\]
を用いると、恒等射はある有限段階を経由する。この因子化は$\epsilon_d$の右逆を与える。
\end{proof}

\begin{remark}[有限段階の選択は解決ではない]\label{rem:lifting}
特殊ファイバーの対角がsmoothなら、ある$d$で$\epsilon_{d,k}$に右逆がある。しかし、その右逆やそのhomotopyが$R$上に\emph{実際に}持ち上がることは別の課題である。必要なのは$\Gamma_R^e$上の有限データ
\[
 s_d:\Gamma_R\longrightarrow P_{d,R},\quad
 H:\epsilon_{d,R}s_d\simeq\mathrm{id}_{\Gamma_R}
\]
（またはこれと同値な有限perfect複体と増大射）である。各Artin切断で異なる有限段数のretractが存在することは、一様な$d$や代数的な右逆を与えない。
\end{remark}

\section{局所Artin底での検出と完備底の障害}
\begin{lemma}[Artin底でのperfectness検出]\label{lem:artin}
$R$が局所Artin環で剰余体$k$を持つとする。任意の$R$-dg代数$\Lambda$と$M\in D(\Lambda)$について、
\[
 M\in\Perf(\Lambda)
 \quad\Longleftrightarrow\quad
 M\otL_Rk\in\Perf(\Lambda\otL_Rk).
\]
特に$\Lambda\otL_Rk$がsmoothなら$\Lambda$は相対smoothである。
\end{lemma}
\begin{proof}
逆向きを述べればよい。$R$は$k$を組成因子とする有限filtrationを持つため$R\in\thick_{D(R)}(k)$である。$M\otL_Rk$のcompactnessにより、$X\mapsto\RHom_\Lambda(M,-\otL_RX)$が任意の直和を保つという条件を満たす$X$は$k$を含むthick部分圏をなす。ゆえに$R$も含み、$M$は$\Lambda$-compact。smoothnessの主張は$\Lambda^e$上の対角に適用する。
\end{proof}

\begin{remark}[完備底への誤った外挿]
$R=k[[h]]$では一般に$R\notin\thick_{D(R)}(k)$であり、上の議論は使えない。全ての$R/(h^n)$でperfectであっても、$R$上でperfectとは限らない。\texttt{notes/source\_recovery/dg\_perfectness\_counterexamples.tex}の明示例はその失敗を示すが、特殊ファイバーの代数自身はsmoothでないため、滑らかな特殊Kuznetsov成分からの相対対角持上げの直接反例ではない。
\end{remark}

\section{generic fiberのperfect性のspread}
\begin{proposition}[genericなperfectnessの局所化]\label{prop:spread}
$R$をnoetherian整域、$K=\operatorname{Frac}(R)$、$\Lambda$をunderlying $R$-perfectなdg代数、$M\in D(\Lambda)$をunderlying $R$-perfectなdg加群とする。$M_K\in\Perf(\Lambda_K)$なら、ある$0\neq s\in R$が存在して
\[
 M\otL_RR[s^{-1}]\in\Perf(\Lambda\otL_RR[s^{-1}]).
\]
\end{proposition}
\begin{proof}
$K=\varinjlim_{s\neq0}R[s^{-1}]$である。perfect dg加群のfiltered-colimit descent（有限cellとretractの表示、またはToën--VaquiéのLemma~2.10に基づく議論）により、ある局所化$R[s^{-1}]$上のperfect $\Lambda_s$-加群$Q_s$で$Q_s\otL_{R_s}K\simeq M_K$となるものを得る。

$Q_s$のcompactnessにより、generic上の同型は必要に応じてさらに局所化した後、射$Q_s\to M_s$で表せる。$Q_s$は$\Lambda_s$-perfect、$\Lambda_s$は底上perfectなので$Q_s$は$R_s$-perfectである。従って$C=\Cone(Q_s\to M_s)$は$R_s$-perfectであり、genericで消滅する。その有限個の有限生成cohomologyを殺す非零元を反転すると$C=0$となる。よって局所化後$M_s\simeq Q_s$となる。
\end{proof}

\begin{corollary}[相対smoothnessのgeneric spreading]\label{cor:spread}
$R$をnoetherian整域とし、$\Gamma\in\Perf(R)$をdg代数とする。$\Gamma_K$が$K$上smoothなら、ある$0\neq s\in R$で$\Gamma_s$は$R_s$上smooth properである。
\end{corollary}
\begin{proof}
命題\ref{prop:spread}を$\Lambda=\Gamma^e$、$M=\Delta_\Gamma$に適用する。どちらもunderlying $R$-perfectである。
\end{proof}

\begin{remark}[何を証明していないか]\label{rem:not-open}
系\ref{cor:spread}はgeneric fiberがすでにsmoothと分かっているときの結果である。特殊ファイバーがsmoothという情報だけではgeneric smoothnessは得られない。よってこの系を$\Gamma_R$の相対smoothness (SM) の証明にそのまま用いることはできない。
\end{remark}

\begin{corollary}[一変数での条件付きfiberwise結論]\label{cor:onevariable}
$R=k[h,\Delta(h)^{-1}]$、$\Delta(0)\neq0$とする。$\Gamma\in\Perf(R)$で必要なdg底変更が成立すると仮定する。$\Gamma_0$と$\Gamma\otL_Rk((h))$がそれぞれの体上smoothなら、$q(0)\neq0$を満たす多項式$q(h)$が存在し、$\Delta(h)q(h)\neq0$である全体値ファイバー$\Gamma_\lambda$はsmooth properとなる。
\end{corollary}
\begin{proof}
$k(h)\subset k((h))$の体拡大に対するfaithfully flatな下降により、generic $k(h)$上で対角がperfectとなる。系\ref{cor:spread}によってある$D(s)$上で相対smoothとなる。$s=h^a q(h)$（$q(0)\neq0$）と書けば、$h\neq0$では$D(q)$の全点が$D(s)$に含まれ、$h=0$は別途仮定でsmooth。これは\emph{全体のファイバーの}結論であって、$D(q)$全体で相対対角がperfectという主張ではない。
\end{proof}

\section{非楕円型cubic fivefoldに残る有限作業}
節点型・尖点型の二パラメータ族のための実在するproper残余dg代数を$\Gamma_R$とする。既存原稿\texttt{papers/sections/local\_projective\_model.tex}はその構成と線束Homの底変更を扱い、\texttt{papers/sections/nonelliptic\_quantization.tex}は量子化と中心三次商の具体式を扱う。いずれも今回のノートの相対対角右逆を作ったものではない。

優先的に構成すべきデータは、$\Gamma_R$の有限のdg（または$A_\infty$）表示を選び、ある$d$に対しbarの有限段階$P_{d,R}$、右逆$s_d$、増大後の恒等homotopyを\emph{有限個の整合式として}提示することである。$\Gamma_R$のunderlying cohomology階数（fivefoldの場合、線束生成dg代数の0次と5次にそれぞれ2838、45）や\(\HH^3(\operatorname{Ku}(X))=0\)だけでは、このデータは得られない。

\paragraph{出典と検証範囲}
本稿は\texttt{notes/pdf-only/finite\_smoothness\_notes.pdf}（2026年11月1日版）第18--21節の補題・命題・証明を現在の問題へ組み直したもの。原本の他の節は残るためPDFのGit追跡解除を認めない。Toën--Vaquié, \emph{Moduli of objects in dg-categories}, Ann. Sci. ENS \textbf{40} (2007), 387--444, Lemmas 2.8 and 2.10, \url{https://arxiv.org/abs/math/0503269} を参照。参考文献の適用と導来bar模型の細部は最終公表前に改めて監査する。

\section{規約と問題}
すべてcohomological gradingとし、微分の次数を$+1$とする。
$\Perf(\Lambda)$はunbounded derived category $D(\Lambda)$のcompact objects、
同値に有限coneとretractで$\Lambda$から生成される対象を表す。
基底環上のperfectnessとdg加群としてのperfectnessを混同しない。
次の含意は一般に偽である：
\[
\Lambda,M\in\Perf(R),\qquad
M\otL_Rk\in\Perf(\Lambda\otL_Rk)
\quad\Longrightarrow\quad M\in\Perf(\Lambda).
\tag{DG4.1}
\]

\section{特殊ファイバーだけperfectな二生成子の例}
\begin{theorem}[特殊ファイバーとすべてのArtin切断でperfect]\label{thm:two}
$R=k[[h]]$, $K=k((h))$とし
\[
\Lambda=R\oplus R\varepsilon\oplus R\eta,\qquad
|\varepsilon|=|\eta|=1,\qquad
(\varepsilon,\eta)^2=0,\quad d_\Lambda=0.
\]
右dg加群$M=e\Lambda$は$|e|=0$, $d_M(e)=he\varepsilon$とする。
このとき
\begin{enumerate}
\item $\Lambda,M$は$R$上の有界有限自由複体である。
\item $M_k=M\otL_Rk\simeq\Lambda_k$は$\Lambda_k$-perfectである。
\item $M_K=M\otL_RK$は$\Lambda_K$-perfectでない。
\item すべての$n\ge1$で$M_n=M\otL_RR/(h^n)$は$\Lambda_n$-perfectである。
\item $M$は$\Lambda$-perfectでない。
\end{enumerate}
\end{theorem}
\begin{proof}
$(\varepsilon,\eta)^2=0$からLeibniz則で$d_M^2=0$となる。
基礎の$R$複体は
\[
0\longrightarrow R\xrightarrow{(h,0)}
R\varepsilon\oplus R\eta\longrightarrow0
\]
であるから1が従い、$h=0$で微分が零となるので2が従う。
$h$を可逆にすると$H^0(M_K)=0$, $H^1(M_K)=Ke\eta$である。
augmentation $\Lambda_K\to K$を通して$K$を加群とみれば
\[
K[-1]\xrightarrow{\simeq}M_K,\qquad1\longmapsto e\eta
\]
は$\Lambda_K$-linear quasi-isomorphismである。

$K$が$\Lambda_K$-perfectでないことを示す。
$\varepsilon,\eta$のすべての有限語$w$（空語を含む）に
次数0の自由生成子$e_w$を対応させる左semifree加群$Q$をとり、
\[
d(e_\varnothing)=0,\qquad d(e_{aw})=ae_w
\quad (a\in\{\varepsilon,\eta\})
\]
とする。radical square zeroゆえ$d^2=0$。
語長filtrationはsemifree filtrationであり、
非空語とradicalの元が対応して消えるので$Q\to K$は
quasi-isomorphismである。augmentationの右$K$とtensorすると
全微分が零になるから
\[
K\otL_{\Lambda_K}K\simeq
\bigoplus_{w\in\{\varepsilon,\eta\}^{*}}Ke_w
\]
は無限次元のcohomologyを持つ。
もし$K\in\Perf(\Lambda_K)$なら
$-\otL_{\Lambda_K}K$はこれを$K$-perfectに送るため矛盾する。
従って3。5はperfectnessが導来基底変更で保存されることから従う。
4は系\ref{cor:artin}から従う。
\end{proof}

\begin{corollary}[fiberwise perfect locusの非開性]
$R=k[h]$上で同じ式を定義すると
\[
\{p\in\Spec R\mid M\otL_R\kappa(p)\in
\Perf(\Lambda\otL_R\kappa(p))\}=V(h).
\]
\end{corollary}
\begin{proof}
$h=0$では自由加群であり、$h\ne0$の剰余体では
定理\ref{thm:two}のaugmentationを用いた議論が成立する。
\end{proof}
\begin{remark}
本例は$H^{>0}(\Lambda_k)\ne0$であり、
$H^{>0}=0$というconnectivity仮定を伴う定理への反例ではない。
さらに$\Lambda_k$自体はsmoothではない。
従って$E_k$ smoothを仮定する相対対角$E$の
perfectnessの反例を構成したものではない。
\end{remark}

\section{Artin切断とthick detection}
\begin{lemma}[thick detection]\label{lem:thick}
$R$を可換環、$S$を可換$R$代数とし
$R\in\thick_{D(R)}(S)$と仮定する。
$R$-dg代数$\Gamma$と$P\in D(\Gamma)$について
\[
P\otL_RS\in\Perf(\Gamma\otL_RS)
\quad\Longrightarrow\quad P\in\Perf(\Gamma)
\]
が成り立つ。connectivityは不要である。
\end{lemma}
\begin{proof}
$X\in D(R)$に対し函手
$N\mapsto\RHom_\Gamma(P,N\otL_RX)$が
任意の直和を保つという条件を考える。
この条件を満たす$X$はthick部分圏をなす。
$X=S$ではextension--restriction adjunctionによって
\[
\RHom_\Gamma(P,N\otL_RS)
\simeq \RHom_{\Gamma\otL_RS}
(P\otL_RS,N\otL_RS)
\]
となり、仮定されたcompactnessと基底変更の
直和保存性から右辺は直和を保つ。
よって$S$はこのthick部分圏に入り、仮定から$R$も入る。
従って$P$はcompactである。
\end{proof}
\begin{corollary}\label{cor:artin}
$R_n=k[h]/(h^n)$上のdg moduleのperfectnessは
$k$への導来基底変更で検出される。特に定理\ref{thm:two}の
すべての$M_n$は$\Lambda_n$-perfectである。
\end{corollary}
\begin{proof}
$R_n$は$k$を組成因子に持つ有限filtrationを持ち、
$R_n\in\thick_{D(R_n)}(k)$である。
補題\ref{lem:thick}を適用する。
\end{proof}
したがって$M\simeq\mathbf R\varprojlim M_n$かつ各$M_n$が
perfectでも$M$自身の$\Lambda$-perfectnessは従わない。

\section{全ファイバーがperfectでも相対perfectでない一生成子の例}
\begin{proposition}
$R=k[[h]]$または$k[h]_{(h)}$とし
\[
\Gamma=R\oplus R\varepsilon,\quad
|\varepsilon|=1,\quad\varepsilon^2=0,\quad d_\Gamma=0,
\qquad P=e\Gamma,\quad d_P(e)=he\varepsilon.
\]
$P$は$R$-perfectで、$P_k\simeq\Gamma_k$および
$P_K\simeq0$はいずれもperfectだが
$P\notin\Perf(\Gamma)$である。
\end{proposition}
\begin{proof}
基礎の$R$複体は$R\xrightarrow{h}R$。
仮に$P$が$\Gamma$-perfectならaugmentation
$\Gamma\to R$により$N=P\otL_\Gamma R\in\Perf(R)$。
導来基底変更では$N\otL_Rk\simeq k$だが
$N\otL_RK\simeq0$である。
局所環上のperfect complexは有界有限自由複体で表されるため
Euler characteristicは各項のrankの交代和であり、
special/genericで一致すべきところ$1\ne0$となる。矛盾。
\end{proof}

\section{完備基底でなお成立するHom有限性}
\begin{proposition}\label{prop:hom}
$R=k[[h]]$, $\Gamma$を$R$-dg代数、
$P,N\in D(\Gamma)$とする。
\[
P\otL_Rk\in\Perf(\Gamma\otL_Rk),\qquad
N\in\Perf(R)
\]
ならば$\RHom_\Gamma(P,N)\in\Perf(R)$である。
\end{proposition}
\begin{proof}
$C=\RHom_\Gamma(P,N)$とする。
$k$は$R$-perfectで$-\otL_Rk$は$h$倍写像のconeだから、
$P$のcompactnessを仮定せず
\[
C\otL_Rk\simeq
\RHom_{\Gamma\otL_Rk}(P\otL_Rk,N\otL_Rk)
\]
であり、右辺は$k$-perfect。
$N$は$R$-perfectゆえderived $h$-completeで、
$\RHom_R(R[1/h],N)=0$。
tensor--Hom adjunctionから
\[
\RHom_R(R[1/h],C)
\simeq\RHom_\Gamma(P,\RHom_R(R[1/h],N))=0.
\]
従って$C$はderived $h$-completeであり
$C\simeq\mathbf R\varprojlim_n(C\otL_RR/(h^n))$。
可換完備環上のperfect complexのeffectivity
（Stacks Project, \href{https://stacks.math.columbia.edu/tag/09AW}{Tag 09AW}）
から$C\in\Perf(R)$を得る。
\end{proof}
\begin{remark}
上の結論は$P$の$\Gamma$-compactnessを意味しない。
$R[1/h]$への一般の$\RHom$基底変更は、
$k$への有限coneでの基底変更とは異なる。
\end{remark}
\begin{corollary}
$E$が$R$-perfectな相対dg代数として構成済みで
$E_k$がsmoothなら、
\[
\RHom_{E^{\mathrm{op}}\otL_RE}(E,E)\in\Perf(R).
\]
これはHochschild cochainの有限性であり
$E$の相対smoothnessを意味しない。
\end{corollary}

\section{非可換cubicの正則性研究への適用限界}
Artin切断で対角がperfectであることから
完備基底上の対角の一様な有限cell/retract模型を
無条件に得ることはできない。
上の反例は一般の加群補題の使用を否定するが、
特定の非可換cubicのtailsがsmoothでないという
結論は与えない。中心局所化と有限行列因子化など、
対角の追加構造を利用した証明と区別する。


\appendix
\part{低次数復元と一次変形の基礎}
\section{rotationの低次数積からの可換方程式の有限復元}
\label{sec:explicit-jacobian-recovery}

本節では、可換模型の\emph{方程式の抽出}を、その模型との
\emph{dg圏同値の存在}から分離する。
前者には有限次元線形代数による必要十分な判定がある。
これを非可換cubicの残余圏に適用するには、まずその圏について
以下の低次数積を実際に計算しなければならない。
本節は一般の非可換cubicの可換模型の存在を証明するものではない。

先行研究との関係を明確にしておく。
Huybrechts--Rennemoはrotationから得られるHochschild環と
Jacobian環を比較し\cite{HR}、Lin--Zhangはgraded matrix factorizationを用いて
一般のFano超曲面に対するrotation付きTorelliを証明している
\cite[Theorem~3.5, Remark~3.6, Section~5]{LZ}。
勾配空間から多項式を復元する問題はUeda--Yoshinagaの研究
\cite[Lemma~3]{ExplicitUedaYoshinaga}に現れ、
積分の非一意性はHwangのsymmetrizer群
\cite[Theorems~1.5--1.6]{ExplicitHwang}によって説明される。
以下の有限検査器はこれらの既知の構造を本研究に適用するための整理であり、
復元原理そのものの新規性を主張しない。
具体的な検査器および反例の初出は未確認である。

\subsection{計算すべき量は次数一と二の積である}

$\cC$を小さな$\CC$-線形dg圏、$\rho$をdg自己同値とする。
通常の三角圏上の自然変換ではなく、dg双加群による\emph{導来自然変換}を用いる。
対角双加群を$\Delta$、$\rho$の双加群を$P$と書き、
\[
 V:=H^0\RHom_{\cC^e}(\Delta,P),\qquad
 W:=H^0\RHom_{\cC^e}(\Delta,P\otimes^{\mathbf L}_{\cC}P)
 \tag{EJ1}\label{eq:ej-low-degree}
\]
とおく。$\alpha,\beta:\Id\to\rho$の積を
$\rho(\beta)\circ\alpha:\Id\to\rho^2$で定める。
これにより$\mu:V\otimes V\to W$を得る。
一般のpairではこの積の対称性は自動ではない。
従って最初に
\[
 \mu(v,w)-\mu(w,v)=0\quad(v,w\in V)
 \tag{EJ2}\label{eq:ej-commutativity}
\]
を確認し、その後でのみ
$\bar\mu:\operatorname{Sym}^2V\to W$および
$K:=\ker\bar\mu$を定義する。
以下、これらのベクトル空間は有限次元であるとする。

\begin{proposition}[可換cubicにおける低次数の同定]
\label{prop:ej-geometric-low-degree}
$F\in\CC[x_0,\ldots,x_{N-1}]_3$が
滑らかなcubic $X\subset\PP^{N-1}$を定め、$N=n+2\ge5$とする。
$(\cC,\rho)=(\Ku(X),\rho_X)$に対して
\[
 V\simeq\operatorname{Jac}(F)_1=\CC\langle x_0,\ldots,x_{N-1}\rangle,
 \qquad W\simeq\operatorname{Jac}(F)_2
\]
であり、$\mu$はJacobian環の積である。特に積は対称で、
\[
 \dim V=N,\quad \dim W=\binom N2,\quad
 \bar\mu\text{ は全射},\quad
 K=\langle\partial_0F,\ldots,\partial_{N-1}F\rangle
 \subset\operatorname{Sym}^2V.
 \tag{EJ3}\label{eq:ej-jacobian-kernel}
\]
\end{proposition}
\begin{proof}
Orlovのdg比較をrotationも含めて用いると、$\rho_X$はgraded matrix
factorizationの内部次数シフトとなる。
Lin--Zhangの\cite[Theorem~2.2, Theorem~3.5, Remark~3.6, Section~5]{LZ}
にある対角kernelの計算を適用する。
恒等sectorは$\operatorname{Jac}(F)_t$を与える。
三次式の非恒等sectorは全$N$変数を動かすので、
cohomological degree零の$\Hom(\Id,\rho^t)$に寄与できるのは
$N$が偶数で$t=N/2$の場合に限る。
$N$が奇数なら寄与はなく、$N$が偶数でも$N\ge6$なので$t=1,2$では寄与しない。
従って表示した二つの同定を得る。
このsector計算には$\gcd(N,3)=1$を必要としない。
同論文Proposition~3.7の互いに素という条件は、内部シフトを
Serre関手とcohomological shiftで表示してSerre代数へ埋め込むための条件であり、
ここでは$\rho$自体を入力している。

恒等sector上の合成は多項式の積である
\cite[Remark~3.6]{LZ}。
Jacobianイデアルは二次式$\partial_iF$で生成され、次数一部分は零である。
滑らかさにより$\partial_iF$は一次独立である。
実際、定数係数の関係があれば線形座標変換後に$F$は一つの変数に依らず、
その変数軸上の射影点が特異点になる。
従って次数二の核と全射性、および次元公式が従う。
\end{proof}

\begin{remark}[低次数だけを使う理由]
全ての$\Hom(\Id,\rho^t)$をそのままJacobian環と同一視してはならない。
例えばcubic fourfoldでは$t=3$に非恒等sectorが現れる。
命題\ref{prop:ej-geometric-low-degree}は、$n\ge3$に共通する
$t=1,2$だけを取り出すことでこの問題を避けている。
\end{remark}

\subsection{二次関係の積分は有限線形方程式である}

$V$の基底を$x_0,\ldots,x_{N-1}$とし、
$S=\operatorname{Sym}V=\CC[x_0,\ldots,x_{N-1}]$とする。
$N$次元部分空間$K\subset S_2$の基底を$q_0,\ldots,q_{N-1}$と書く。
次の線形部分空間を定義する：
\[
 \mathcal L_K:=\left\{A=(A_{ia})\in\operatorname{Mat}_N(\CC):
 \sum_a A_{ia}\partial_jq_a=\sum_a A_{ja}\partial_iq_a
 \text{ for all }i,j\right\}.
 \tag{EJ4}\label{eq:ej-integrability}
\]
各等式の両辺は一次式なので、係数比較により
高々$N\binom N2$本の斉次線形方程式を$N^2$個の未知数について解くだけでよい。

\begin{theorem}[方程式抽出の必要十分条件]
\label{thm:ej-effective-recovery}
$K\subset S_2$を$N$次元部分空間とする。
次の二条件は同値である。
\begin{enumerate}
 \item 滑らかなcubic $F\in S_3$が存在し、
 $K=\langle\partial_0F,\ldots,\partial_{N-1}F\rangle$である。
 \item $K$の共通零点は$\mathbb A^N$で原点だけであり、
 $\mathcal L_K$は可逆行列$A$を含む。
\end{enumerate}
第二条件が成り立つとき、任意の可逆$A\in\mathcal L_K$について
\[
 \boxed{\quad
 F_A=\frac13\sum_{i,a}x_i A_{ia}q_a,
 \qquad \partial_iF_A=\sum_aA_{ia}q_a
 \quad}
 \tag{EJ5}\label{eq:ej-euler-recovery}
\]
が求める滑らかなcubicを与える。
\end{theorem}
\begin{proof}
第一条件が成り立つなら、$\partial_iF$は$K$の別の基底なので
$\partial_iF=\sum_a A_{ia}q_a$となる可逆行列$A$がある。
混合偏微分の交換性が$A\in\mathcal L_K$を与える。
三次のEuler等式により、偏微分の共通射影零点は
$V(F)$の特異点に一致するから、$K$には射影共通零点がない。

逆に$g_i=\sum_aA_{ia}q_a$と置く。
式\eqref{eq:ej-integrability}と二次のEuler等式から
\[
 \partial_jF_A
 =\frac13\left(g_j+\sum_i x_i\partial_jg_i\right)
 =\frac13\left(g_j+\sum_i x_i\partial_i g_j\right)=g_j.
\]
$A$が可逆なので、勾配は$K$を張る。
その共通射影零点がないことから$V(F_A)$は滑らかである。
\end{proof}

\begin{corollary}[有限の判定手順]
\label{cor:ej-finite-algorithm}
$K$の係数が計算可能な標数零の体にあれば、
代数閉包上で定理\ref{thm:ej-effective-recovery}の条件を有限計算で判定できる。
\end{corollary}
\begin{proof}
まず線形方程式\eqref{eq:ej-integrability}の核の基底
$A_1,\ldots,A_r$を求める。
$\mathcal L_K$に可逆行列がある条件は
\[
 D_K(z_1,\ldots,z_r)=\det\left(\sum_{\nu=1}^r z_\nu A_\nu\right)
 \not\equiv0
 \tag{EJ6}\label{eq:ej-determinant-test}
\]
であり、有限の多項式恒等判定である。
非零なら、例えば$\{0,1,\ldots,N\}^r$の有限格子のいずれかで非零になる。
これは$D_K$の各変数の次数が高々$N$であることと一変数の零点数から帰納的に従う。
その値を使って式\eqref{eq:ej-euler-recovery}を適用すればよい。
共通零点条件は、斉次イデアル$(K)$の商が零次元であることと同値であり、
Gr\"obner基底で判定できる。
この方法は実行可能性を保証する有限手順であって、最適な計算量の主張ではない。
\end{proof}

\subsection{候補は射影同型を除いて一意だが、式は一意とは限らない}

\begin{proposition}[滑らかな候補の射影的一意性]
\label{prop:ej-projective-uniqueness}
同じ$K\subset S_2$を勾配空間に持つ二つの滑らかなcubicは射影同型である。
ただし固定座標での方程式がscalar倍に限られるとは限らない。
\end{proposition}
\begin{proof}
$F,G$をその二つのcubicとする。
それぞれの勾配は$K$の基底なので$\nabla G=A\nabla F$を満たす可逆行列$A$がある。
$F$を、$F(x)=T(x,x,x)$を満たす対称三重線形形式$T$で表す。
$a=A^{\mathsf t}$とすれば、$G$の対称性は
\[
 T(au,v,w)=T(u,av,w)=T(u,v,aw)
 \tag{EJ7}\label{eq:ej-symmetrizer}
\]
を意味する。一般にこの式を満たす線形写像全体を$\mathfrak s_T$とする。

$a,b\in\mathfrak s_T$について、三つの引数の間で写像を移すと
\[
 \begin{aligned}
 T(abu,v,w)&=T(bu,av,w)=T(u,av,bw)\\
 &=T(au,v,bw)=T(au,bv,w)=T(bau,v,w).
 \end{aligned}
\]
$F$が滑らかなので$u\mapsto T(u,-,-)$は単射であり、$ab=ba$が従う。
さらに式\eqref{eq:ej-symmetrizer}から$ab$もsymmetrizerである。
従って$\mathfrak s_T$は単位を持つ可換部分代数である。

この代数に非零nilpotentがあれば、その適当な冪を取って
$0\ne b\in\mathfrak s_T$かつ$b^2=0$とできる。
$u=bv\ne0$を選ぶと
\[
 T(u,u,w)=T(bv,bv,w)=T(v,b^2v,w)=0
 \quad\text{for every }w.
\]
これは$[u]$が$V(F)$の特異点であることを意味し、矛盾する。
Jordan分解のnilpotent部分は元の写像の多項式であるから、
従って$\mathfrak s_T$の全ての元は半単純で、特に$a$は対角化できる。
固有空間分解を$\CC^N=\bigoplus_\lambda U_\lambda$とすると、
式\eqref{eq:ej-symmetrizer}により異なる固有空間を含む混合項は零になる。
よって
\[
 F=\sum_\lambda F_\lambda,
 \qquad G=\sum_\lambda\lambda F_\lambda,
 \qquad F_\lambda\in\operatorname{Sym}^3 U_\lambda^*.
\]
$a$が可逆なので各$\lambda\ne0$である。
各$U_\lambda$の座標を$\lambda^{1/3}$倍すれば$F$を$G$へ移せる。
この線形変換が射影同型を与える。
\end{proof}

\begin{example}[Fermatの場合の積分解]
\label{ex:ej-fermat}
$K=\langle x_0^2,\ldots,x_{N-1}^2\rangle$では
式\eqref{eq:ej-integrability}は$A$が対角行列であることと同値である。
従って$\dim\mathcal L_K=N$であり、
\[
 F_A=\frac13\sum_i a_i x_i^3\qquad(a_i\ne0)
\]
は全て解である。これらは射影同型であるが、scalar倍だけの一意性はない。
特にFermat cubic threefoldでは積分解空間の次元は$5$であり、$1$ではない。
\end{example}

\subsection{正しいHilbert級数を持っても積分できない二次関係}

次の反例は、次数一・二の次元や完全交差性だけでは
可換cubicの候補すら作れないことを示す。
非可換cubicの残余圏からこの例が生じるとは主張しない。

\begin{proposition}[五変数の積分不可能な完全交差]
\label{prop:ej-nonintegrable-ci}
添字を$\ZZ/5\ZZ$で読み、
\[
 q_i=x_i^2+x_{i+1}x_{i+2},\qquad
 K=\langle q_0,\ldots,q_4\rangle
 \subset\CC[x_0,\ldots,x_4]_2
 \tag{EJ8}\label{eq:ej-cycle-example}
\]
とする。このとき以下が成り立つ。
\begin{enumerate}
 \item $K$には共通射影零点がなく、$S/(K)$はArtin完全交差で
 Hilbert級数$(1+t)^5$を持つ。
 \item $\mathcal L_K=0$である。
 特に$S/(K)$はどの滑らかなcubic threefoldのgraded Jacobian環とも同型でない。
\end{enumerate}
\end{proposition}
\begin{proof}
$q_i(x)=0$とする。ある$x_i=0$なら
$q_{i-2}=x_{i-2}^2+x_{i-1}x_i=0$から$x_{i-2}=0$である。
$2$と$5$は互いに素なので全ての座標が零になる。
全ての$x_i$が非零なら五つの式$x_i^2=-x_{i+1}x_{i+2}$の積から
$\prod_i x_i^2=-\prod_i x_i^2$となり矛盾する。
従って共通零点は原点だけである。
五つの$q_i$は各々固有の平方項を持つため一次独立であり、
多項式環のsystem of parametersなので正則列である。
よってHilbert級数は$(1-t^2)^5/(1-t)^5=(1+t)^5$である。

次に$A\in\mathcal L_K$とし、$g_i=\sum_aA_{ia}q_a$とおく。
直接微分すると
\[
 \partial_jg_i=2A_{ij}x_j+A_{i,j-1}x_{j+1}+A_{i,j-2}x_{j-1}.
 \tag{EJ9}\label{eq:ej-cycle-derivative}
\]
$j=i+1$の可積分性で$x_{i+2}$と$x_{i-1}$の係数を比較すると
$A_{ii}=0$および$A_{i+1,i-2}=0$である。
添字を動かせば$A_{i,i+2}=0$も得る。
$j=i+2$で$x_{i+3}$と$x_i$の係数を比較すると
$A_{i,i+1}=0$および$A_{i+2,i}=0$、従って$A_{i,i+3}=0$を得る。
残り得る成分を$a_i=A_{i,i-1}$とすると、
$j=i+1$で$x_i$の係数を比較して$a_i=2a_{i+1}$となる。
五回繰り返すと$a_i=32a_i$なので$a_i=0$である。従って$A=0$。

もしgraded Jacobian環と同型なら、その次数一部分の同型は
多項式環の線形座標変換に持ち上がる。
次数二の核から、この座標変換後の$K$は滑らかな三次式の勾配空間になるはずである。
しかし可積分性は座標変換に不変であり、$\mathcal L_K=0$と矛盾する。
\end{proof}

\subsection{非可換cubicに適用するときの結論と残る課題}

\begin{corollary}[唯一の可換候補を抽出する検査器]
\label{cor:ej-candidate}
非可換cubicのenhanced残余pair $(\cC,\rho)$に対して、
式\eqref{eq:ej-low-degree}の$V,W$および積$\mu$が計算されたとする。
このpairが滑らかな可換cubic $n$-foldと同値なら、$N=n+2$として
次の全てを満たす：
\begin{enumerate}
 \item $\dim V=N$、$\dim W=\binom N2$、積の対称性と全射性。
 \item $K=\ker\bar\mu$は$N$次元で、共通射影零点を持たない。
 \item $\mathcal L_K$には可逆行列が存在する。
\end{enumerate}
いずれかが失敗すれば、このpairの可換cubic模型は存在しない。
全て成功すれば、式\eqref{eq:ej-euler-recovery}は
\emph{そのような模型が存在するときの唯一の候補}を射影同型を除いて与える。
この成功だけから$(\cC,\rho)\simeq(\Ku(V(F_A)),\rho_{F_A})$を結論してはならない。
\end{corollary}
\begin{proof}
必要性は命題\ref{prop:ej-geometric-low-degree}と
定理\ref{thm:ej-effective-recovery}による。
候補の射影的一意性は命題\ref{prop:ej-projective-uniqueness}による。
これらの証明は低次数の導来自然変換の積しか使わず、
pair間の同値双加群やその逆を構成していない。
従って圏同値の存在は独立に証明すべき事項として残る。
\end{proof}

具体的な非可換圏への適用では、残余dg代数とrotation双加群から
導来自然変換の次数零コホモロジーと合成を求める必要がある。
線束Homの階数表、Poissonの一次項、ambientの二次関係を
そのまま低次数積に代入することはできない。
また、この有限判定は実際の同値双加群や相対対角のperfect性を構成しない。
これらは方程式候補の抽出とは独立の入力である。

\paragraph{計算の再現範囲.}
\texttt{computations/verify\_explicit\_jacobian\_recovery.py}は有理数上で、
Fermat三次式、混合項を持つ滑らかな五変数三次式、
および命題\ref{prop:ej-nonintegrable-ci}の完全交差を検査する。
Jacobianの次数二商の積から核を取り直して積分し、
勾配との一致、線形系の階数、共通零点、Hilbert関数を確認する。
反例では非零の最大小行列式も証明書に記録する。
入力は既知の可換多項式または明示的な二次関係であり、
非楕円量子化の残余圏の導来自然変換を計算したものではない。

% First-order cancellation, quadratic obstructions, and categorical tangent map.
\section{一次の非可換変形と可換模型の接空間}
\label{fpk:sec:first-order}

この節では一般点の全次数形式比較を前提とせず、
双対数環上で非可換積の一次摂動が曲率付きMFにどう作用するかを直接調べる。
これは全次数の定理\ref{cft:thm:complete}の一次部分の独立検証でもある。
曲率の一次変形を固定しなければ、任意の非可換変形が自明という意味ではない。

$S=\mathbb C[x_1,\ldots,x_N]$、$N\ge5$、
$D=\mathbb C[\varepsilon]/(\varepsilon^2)$ とする。
$F\in S_3$は滑らかな射影三次超曲面を定めるものとする。
$A_D$を単位付き・次数別有限自由な$D$-代数で
$A_D/\varepsilon A_D\simeq S$、
$\Omega_D\in Z(A_D)_3$の剰余が$F$となるものとする。

\begin{theorem}[一次Poisson方向の曲率付き消滅]
\label{fpk:thm:infinitesimal}
上の入力に対し、単位と内部次数を保つ$D$-加群標識を適切に選べば、
$A_D$の積と中心元は
\begin{equation}\label{fpk:eq:normal}
 a*b=ab+\varepsilon\,\Theta_\pi(a,b),\qquad
 \Omega_D=F+\varepsilon G,\qquad
 \Theta_\pi(a,b)=\frac12\{a,b\}_\pi
\end{equation}
と書ける。ここで$\pi$は二次係数のbivectorで$[\pi,F]=0$である。
さらに、この標識で有限rank graded MFの
$D$-線形unital $A_\infty$同値
\begin{equation}\label{fpk:eq:equivalence}
 \operatorname{MF}^{\mathbb Z}(A_D,\Omega_D)
 \simeq
 \operatorname{MF}^{\mathbb Z}(S\otimes D,F+\varepsilon G)
\end{equation}
が存在し、剰余では恒等、内部shiftおよび
中心元による三次周期同型$(3)\simeq[2]$を保つ。

二つの標識で得た$G$の差は
$J(F)_3=(\partial_1F,\ldots,\partial_NF)_3$に属する。
従って、一次の可換模型の座標同値類を定める写像は
\begin{equation}\label{fpk:eq:tangent}
 [G]\in(S/J(F))_3
\end{equation}
を通る。特に$G=0$と取れる正規化では、一次Poisson変形は
rotation付きgraded MF圏の一次変形を生じない。
\end{theorem}

\begin{proof}
次数別自由な加群標識を選ぶと、積は$\mu+\varepsilon\beta$となる。
単位を保つよう標識を取り、$\beta$をnormalizedな内部weight零の
Hochschild二cochainとする。結合律の一次部分は
$d_{\mathrm{Hoch}}\beta=0$である。
polynomial algebraのHKR同型により、そのクラスは
二次係数bivector$\pi$の反対称微分cochain$\Theta_\pi$で表せる。
差$\beta-\Theta_\pi$は内部weight零のHochschild coboundaryなので、
$\mathrm{id}+\varepsilon T$という一次の厳密な加群標識変更で
積を式\eqref{fpk:eq:normal}にできる。
この変更は元の代数の結合律を保ち、単位も保つ。
中心元の表示を$F+\varepsilon G$と書く。
中心性の一次部分は$[\pi,F]=0$である。

$F$が滑らかなので、全偏微分は正則列である。
$\iota_{dF}$による斉次Koszul複体をrank二・係数次数二で用いると、
$[\pi,F]=0$から定数係数三重vector
$\Xi\in\bigwedge^3\langle\partial_1,\ldots,\partial_N\rangle$を得て
\begin{equation}\label{fpk:eq:primitive}
 \pi=[\Xi,F]
\end{equation}
となる。原像の差を消すにはrank四・係数次数$-2$の原像が必要となり、
そのような多項式はないので$\Xi$は一意である。
ここでは$\pi$のJacobi恒等式を使わない。

$\Theta_\Xi=\operatorname{HKR}(\Xi)$を、標準の$1/3!$反対称化で
正規化した三入力Hochschild cochainとする。
三つの座標方向の単純三重vectorについては
\[
 \Theta_\Xi(a,b,c)
 =\frac16\det
 \begin{pmatrix}
  \partial_i a&\partial_i b&\partial_i c\\
  \partial_j a&\partial_j b&\partial_j c\\
  \partial_k a&\partial_k b&\partial_k c
 \end{pmatrix}.
\]
式\eqref{fpk:eq:primitive}と反対称性より
\begin{equation}\label{fpk:eq:insert}
 \Theta_\Xi(F,a,b)-\Theta_\Xi(a,F,b)+\Theta_\Xi(a,b,F)
 =\Theta_\pi(a,b).
\end{equation}
$\Theta_\Xi$は多重導分なので
$d_{\mathrm{Hoch}}\Theta_\Xi=0$である。

周期変数に$\deg_{\mathrm{coh}}t=2$、
$\deg_{\mathrm{int}}t=-3$を与え、曲率を$-tF$とする。
\[
 \eta=\varepsilon\,t^{-1}\Theta_\Xi
\]
はcohomological・内部次数ともに零、arity三のnormalized cochainである。
曲率付きGerstenhaber括弧で式\eqref{fpk:eq:insert}を使うと
\[
 [\eta,\mu]=0,\qquad
 [\eta,-tF]=-\varepsilon\Theta_\pi
\]
となる（符号は式\eqref{fpk:eq:insert}を採用する）。
従って、曲率付き構造
$\mu-tF+\varepsilon(\Theta_\pi-tG)$は、
$\eta$による一次ゲージで$\mu-t(F+\varepsilon G)$へ移る。
$\varepsilon^2=0$だから、そのbar写像は
$F_0=0,F_1=\mathrm{id},F_3=\varepsilon t^{-1}\Theta_\Xi$のみを持ち、
逆は$F_3$の符号を反転させればよい。

有限rank MFの行列$Q$は$\varepsilon$の倍数である必要はない。
しかし対象像$Q'=Q+F_3(Q,Q,Q)$および射のtwisting挿入公式は
双対数環上では有限和である。
bar恒等式がMF方程式と関手恒等式を与え、逆bar写像で互いに逆な
$A_\infty$関手が得られる。写像は内部weight零で
$t$を微分しないので内部shiftと$(3)\simeq[2]$を保つ。
これで式\eqref{fpk:eq:equivalence}が従う。

最後に、積を$\mu+\varepsilon\Theta_\pi$へ正規化する
二つの標識の差を$\mathrm{id}+\varepsilon T$とする。
同じ正規化された積を保つ条件は
$d_{\mathrm{Hoch}}T=0$であるから、$T$は内部weight零の導分である。
従って$T(x_i)=\sum_j a_{ij}x_j$であり、
中心元の一次係数は$G\mapsto G+T(F)$と変わる。
この差は$J(F)_3$に属する。
また、同じ$F$の一次座標変更による変形は
この形の$T(F)$で尽くされる。
従って式\eqref{fpk:eq:tangent}を得る。
\end{proof}

\begin{corollary}[Jacobian接空間の次元]
\label{fpk:cor:dimension}
滑らかな三次式$F$に対して
\begin{equation}\label{fpk:eq:dimension}
 \dim_{\mathbb C}(S/J(F))_3
 =\binom{N+2}{3}-N^2.
\end{equation}
特にcubic threefoldでは$10$、cubic fourfoldでは$20$である。
\end{corollary}
\begin{proof}
偏微分の二次式は正則列であり、一次係数のsyzygyを持たない。
従って$J(F)_3$は$N$個の二次偏微分に$N$次元の一次式を掛けた
$N^2$次元の部分空間である。
一方$\dim S_3=\binom{N+2}{3}$である。
\end{proof}

\begin{remark}[Jacobi条件が必要になる位置]
\label{fpk:rem:jacobi}
一次の曲率付き同値に必要なのは中心性からの$[\pi,F]=0$であり、
$[\pi,\pi]=0$は必要ない。
二次以降で同じゲージを延長するには
$[\Xi,\pi]=0$が必要になる。
滑らかな三次式の下では、Jacobi条件$[\pi,\pi]=0$から
勾配Koszulと係数次数により$[\Xi,\pi]=0$を導ける
（補題\ref{cft:lem:primitive}）。
従って一次消滅と全次数の比較は異なる論理的強さを持つ。
\end{remark}


\subsection{一次対変形の接空間と二次のPl\"ucker障害}

\begin{proposition}[次数付き中心三次対の一次接空間]
\label{fpk:prop:tangent-pair}
$N\ge5$とし、可換対$(S,F)$における次数別flatな中心三次対の
双対数変形を、剰余で恒等となる単位付き次数付き同型で割る。
中心元は次数三の実際の元として標識する。
この一次変形集合には自然なベクトル空間構造があり、
HKR正規形を選べば
\begin{equation}\label{fpk:eq:tangent-pair}
 T^1_{\mathrm{pair},(S,F)}
 \simeq
 \bigwedge^3\mathbb C^N\oplus\operatorname{Jac}(F)_3
\end{equation}
と表される。第一成分は$\pi=\iota_{dF}\Xi$、第二成分は
$\Omega=F+\varepsilon G$の$[G]$である。
特に$N=5$ではこの接空間は$20$次元である。
\end{proposition}
\begin{proof}
定理\ref{fpk:thm:infinitesimal}で用いたnormalized HKRにより、
積の一次項の同型類は二次係数bivector$\pi$で記述される。
中心性の一次条件は$[\pi,F]=0$であり、勾配Koszulの完全性と
係数次数から唯一の定数三重vector$\Xi$が存在して$\pi=[\Xi,F]$となる。
逆に任意の$\Xi$はこの中心性条件を満たし、
$\Theta_\pi$を一次積とする双対数上の結合代数を与える。
任意の三次式$G$は中心元の一次項に取れる。
一方、積のHKR正規形を保つ剰余恒等の次数付き標識変更は
一次の導分$T$であり、$G$を$G+T(F)$へ移す。
内部次数零の導分は一次形式の線形変換だから、
$T(F)$の張る空間は$J(F)_3$である。
よって表示を得る。
\end{proof}

以下では$N=5$とする。固定体積形式を使い、
$\bigwedge^3V\simeq\bigwedge^2V^*\otimes\det V$
で$\Xi$に対応する二形式を$\alpha_\Xi$と書く。
その五つのPfaffianを$q_0,\ldots,q_4$とし、
\[
 q(\Xi)=\alpha_\Xi\wedge\alpha_\Xi
 \in\bigwedge^4V^*
\]
で表す。$q(\Xi)=0$は$\Xi$の分解可能性と同値で、
$q_i$は$\operatorname{Gr}(3,V)$のPl\"uckerイデアルを生成する。

\begin{theorem}[二次Jacobi障害のスキーム論的Pl\"ucker表示]
\label{fpk:thm:plucker-obstruction}
$V=\mathbb C^5$、$F\in\operatorname{Sym}^3V^*$を
滑らかな射影三次超曲面の方程式とする。
定数三重vector$\Xi$について$\pi_\Xi=\iota_{dF}\Xi$と置くと
\[
 [\pi_\Xi,\pi_\Xi]=0
 \quad\Longleftrightarrow\quad
 \Xi\text{は分解可能（零を含む）}.
\]
さらに$[\pi_\Xi,\pi_\Xi]$の全係数を$\Xi$の座標の二次式と見て
生成されるイデアルは、五つのPl\"ucker二次式のイデアル
$(q_0,\ldots,q_4)$に\emph{スキームとして}一致する。
\end{theorem}
\begin{proof}
まず$\Xi$が分解可能なら、座標変換により
$\Xi=\partial_0\wedge\partial_1\wedge\partial_2$とできる。
$\pi_\Xi$は三つのactive変数のJacobian bracketであり、
任意の$F$についてJacobi恒等式を満たす。

逆に$\Xi$が分解不可能なら、五次元における二形式の交代標準形より、
座標変換と非零スカラー倍で
\[
 \Xi=\partial_0\wedge\partial_1\wedge\partial_2+
       \partial_0\wedge\partial_3\wedge\partial_4
\]
とできる。$F_i=\partial_iF$、$F_{0i}=\partial_0\partial_iF$
と書く。$\pi_{ij}=\sum_k\Xi_{ijk}F_k$を用いてJacobiatorを直接展開すると、
$[\pi_\Xi,\pi_\Xi]=0$はちょうど
\begin{equation}\label{fpk:eq:jacobi-normal}
 F_{0i}F_j-F_{0j}F_i=0
 \qquad(0\le i<j\le4)
\end{equation}
に等しい（成分の順序による符号と共通非零係数は無関係）。
式\eqref{fpk:eq:jacobi-normal}に$x_j$を掛けて$j$について加え、
Euler等式$\sum_jx_jF_j=3F$および
$\sum_jx_jF_{0j}=2F_0$を使えば
\begin{equation}\label{fpk:eq:euler-obstruction}
 3F F_{0i}=2F_0F_i
 \qquad(0\le i\le4)
\end{equation}
を得る。もし$F_0=0$なら$F$は$x_0$に依らず、
$[1:0:\cdots:0]$が特異点となる。
$F_0\ne0$なら、滑らかな射影三次式$F$は既約なので
$\gcd(F,F_0)=1$である。式\eqref{fpk:eq:euler-obstruction}から
$F\mid F_i$が全$i$について従うが、$\deg F_i=2<3$なので
全$F_i=0$となり矛盾する。従って$\Xi$は分解可能である。

次に二次式の\emph{イデアル}の一致を示す。
$[\pi_\Xi,\pi_\Xi]$の係数は$\Xi$について斉次二次式であり、
分解可能な$\Xi$で零になる。Grassmannianの斉次イデアルは
五つのPl\"ucker二次式で生成されるので、
係数全体は$q_0,\ldots,q_4$の線形結合である。
従ってある線形写像$L_F$について、Jacobiatorの全係数ベクトルは
$L_F(q(\Xi))$と表される。

ここで二次写像$q:\bigwedge^3V\to\bigwedge^4V^*$は全射である。
実際、五次元では任意の非零四形式は分解可能であり、
適当な基底で$c\,e^1\wedge e^2\wedge e^3\wedge e^4$と書ける。
$\alpha=e^1\wedge e^2+(c/2)e^3\wedge e^4$とすれば
$\alpha\wedge\alpha$はその四形式である。
もし$L_F$に非零の核があれば、$q(\Xi)$がその非零元となる$\Xi$を選べる。
すると$[\pi_\Xi,\pi_\Xi]=0$なのに$q(\Xi)\ne0$であり、
前半の結論と矛盾する。よって$L_F$は単射である。
その双対写像の全射性から、Jacobiatorの係数二次式は
Pl\"ucker二次式の五次元空間全体を張り、生成イデアルも一致する。
\end{proof}

\begin{corollary}[二次まで延長できる非可換一次方向]
\label{fpk:cor:second-order}
$N=5$の中心三次対の一次変形の非可換方向
$\Xi\in\bigwedge^3V$が、次数付き\emph{結合積}の二次変形
$\mathbb C[\varepsilon]/(\varepsilon^3)$へ延長可能であるためには、
五つのPl\"ucker二次式が零でなければならない。
結合積に関する二次障害の零スキームはPl\"ucker錐と一致する。
その錐の局所次元は$7$で、頂点のZariski接空間は$10$次元である。
これは中心元を含めた対の高次延長障害がすべて消えるという主張ではない。
\end{corollary}
\begin{proof}
結合積$\mu+\varepsilon\beta_1+\varepsilon^2\beta_2$の結合律の
二次項は
$d_{\rm Hoch}\beta_2+\frac12[\beta_1,\beta_1]=0$である。
HKR同型で二次障害の$HH^3$類は
$\frac12[\pi_\Xi,\pi_\Xi]$に対応し、
この類が零であることが結合積の二次延長の必要十分条件である。
前定理からその係数イデアルは五つのPl\"ucker関係式になる。
$\operatorname{Gr}(3,5)$の次元は$3(5-3)=6$なので
錐の次元は$7$、零点では全五式の一次微分が零で接空間は$10$次元である。
\end{proof}

\begin{remark}[Pl\"ucker障害の射程]
\label{fpk:rem:plucker-scope}
滑らかな三次三様体のPoisson構造とGrassmannianの対応は
既知の研究\cite{PMayanskiy}と整合する。
上の定理は、その対応を\emph{結合積の二次延長障害の係数イデアル}として
直接記述するものである。
この記述の先行研究における初出は未確認である。
障害の零スキームを同定したことだけから、
中心元も含む対の全高次形式変形関手を同定したとは結論しない。
\end{remark}


\subsection{Kuznetsov成分の圏論的接空間への射}
\begin{theorem}[可換点での圏論的忘却射の微分]
\label{fpk:thm:categorical-tangent}
$N=5$、$X=V(F)\subset\mathbb P^4$を滑らかなcubic threefoldとする。
次数別flatな中心三次対の一次変形を、剰余で恒等な次数付き同型で割った
$T^1_{\mathrm{pair},(S,F)}$を考える。
Artin基底上のenhancedなstabilization/window比較
（定理\ref{am:thm:comparison}）によって、
rotation付き残余圏の変形に付随する
\emph{基礎dg圏を忘れた}Kodaira--Spencer写像
\[
 \kappa_{\rm cat}:T^1_{\mathrm{pair},(S,F)}
 \longrightarrow HH^2(\operatorname{Ku}(X))
\]
を定めることができる。
命題\ref{fpk:prop:tangent-pair}のHKR正規形による分解では
\begin{equation}\label{fpk:eq:categorical-KS}
 \kappa_{\rm cat}(\Xi,[G])=\eta([G]),
 \qquad
 \eta:H^1(X,T_X)\xrightarrow{\ \sim\ }HH^2(\operatorname{Ku}(X)).
\end{equation}
特に次の完全列を得る。
\begin{equation}\label{fpk:eq:categorical-exact}
 \boxed{
 0\longrightarrow\bigwedge^3\mathbb C^5
 \longrightarrow T^1_{\mathrm{pair},(S,F)}
 \xrightarrow{\ \kappa_{\rm cat}\ }
 HH^2(\operatorname{Ku}(X))
 \longrightarrow0 .
 }
\end{equation}
したがって、可換点での圏論的忘却射の微分は全射で、
その核は$10$次元、像も$10$次元である。
この結論は非零Poisson一般点での微分の階数を主張しない。
\end{theorem}
\begin{proof}
滑らかな三次三様体の標準正規束完全列とEuler列によって、
三次式の変形の接空間は
\[
 H^1(X,T_X)\simeq
 S_3/\langle x_i\partial_jF:0\le i,j\le4\rangle
 =\operatorname{Jac}(F)_3
\]
と同定される。一次Poisson方向は命題\ref{fpk:prop:tangent-pair}の
$\Xi$で表される。定理\ref{fpk:thm:infinitesimal}は、
同じ双対数環上で、剰余が恒等でrotationと指定周期を保つ
MFの実際の$A_\infty$同値を
$(A_D,\Omega_D)$と$(S_D,F+\varepsilon G)$の間に構成する。
Artin基底上のenhanced windowと可換側の相対Orlov比較を合成すると、
この同値は両者の残余dg圏の一次変形を同一視する。
従って、その$HH^2$におけるKodaira--Spencer類は
非可換Poisson方向$\Xi$には依存せず、
滑らかな三次式$F+\varepsilon G$のKuznetsov成分の変形類
$\eta([G])$に等しい。

Jacovskis--Lin--Liu--Zhang
\cite[Theorem~4.4(2)]{JLLZ}は、任意の滑らかな
index-two・degree-three Fano threefoldについて
$\eta:H^1(T_X)\to HH^2(\operatorname{Ku}(X))$が同型であることを示す。
従って式\eqref{fpk:eq:categorical-KS}から
式\eqref{fpk:eq:categorical-exact}が従う。
ここで$HH^2$は\emph{基礎dg圏}の変形接空間であり、
rotationを忘れずに扱うモジュライstack全体の接空間と
自動的に同一視したものではない。
\end{proof}
\begin{remark}[圏論的微分と相対復元]
\label{fpk:rem:categorical-scope}
完全列\eqref{fpk:eq:categorical-exact}は\emph{可換点での}
基礎dg圏の変形接空間に関するものである。
非零Poisson一般点に沿う忘却射の微分を記述するには、
相対的な可換復元、中心元の標識とKodaira--Spencer写像の比較を
さらに必要とする。これは可換点の計算とは別の問題である。
\end{remark}

\begin{remark}[モジュライの局所次元と接空間]
\label{fpk:rem:moduli}
滑らかな三次三様体では$\dim\operatorname{Jac}(F)_3=10$である。
これに対し、\emph{非零Poisson一般点の次元}と
\emph{可換点でのZariski接空間}は異なる。
固定した滑らかな五変数三次式に対するPoisson錐はPl\"ucker錐であり、
非零分解可能な一般点では$1+\dim\operatorname{Gr}(2,5)=7$次元である。
一方、錐の頂点では全てのPl\"ucker方程式が二次式なので、
Zariski接空間は$\bigwedge^3\mathbb C^5$の全$10$次元となる。
元の可換cubicの変形接空間$10$次元を合わせると、
中心三次対の可換点における一次接空間は$20$次元である。
二次障害零スキームの局所次元$17$はこの接空間の次元と異なる。
本節の二次障害計算はその差を直接説明するが、
対の\emph{全高次}形式局所環の同定をこの計算だけから結論しない。

定理\ref{fpk:thm:categorical-tangent}の完全列は可換点での核が
$10$次元であることを示すが、非零Poisson一般点での核の次元は与えない。
後者については注意\ref{fpk:rem:categorical-scope}を参照されたい。
\end{remark}

% Todd-corrected HKR proof of Poisson first-order vanishing.
\section{Todd補正から見た三次三様体の一次Poisson消滅}
\label{pf:sec:first-order}
中心三次対の曲率付きHochschild cochainによる直接的な一次MF比較とは別に、
本節では\emph{滑らかな三次三様体}について
Todd補正付きHKRとKuznetsov成分への制限を用いてPoisson方向の一次消滅を示す。
純粋なPoisson方向はKuznetsov成分の一次変形では消えるため、
可換模型が存在するなら、その方程式の本質的な変化は二次以降に現れる。
ここで証明するのは一次接写像の消滅であって、可換模型の存在、
残余smoothness (SM)、全次数での局所自明性ではない。

用いる外部結果は、Calaque--Rossi--Van den BerghのTodd補正付きHKRと
加群構造の両立性\cite[Theorem 1.4]{PFCRVdB}、および
Jacovskis--Lin--Liu--ZhangのHochschild作用の可換図式と
cubic threefoldに対する単射性\cite[Theorems 3.10 and 4.4(2)]{JLLZ}である。
以下の消滅命題はこれらの既知結果の適用であり、その帰結としての新規性は未確認である。

\subsection{通常のHKRとTodd補正の比較}
$X\subset\mathbb P^4_{\mathbb C}$を滑らかな三次超曲面とし、
\[
 \operatorname{HT}^2(X)
 =H^0(X,\wedge^2T_X)\oplus H^1(X,T_X)\oplus H^2(X,\mathcal O_X)
\]
と書く。通常のHKR同型を
$I_{\mathrm{HKR}}:\operatorname{HT}^2(X)\xrightarrow{\sim}\operatorname{HH}^2(X)$、
加群構造と両立するTodd補正付き同型を$I^K$とする。
純粋なPoisson方向とは、一次変形類が
$I_{\mathrm{HKR}}(\pi)$、$\pi\in H^0(X,\wedge^2T_X)$、
$[\pi,\pi]=0$であることをいう。
同時に動かす可換複素構造の成分をこの条件に含めない。

\begin{lemma}[bivectorによる第一Chern類の縮約]
\label{pf:lem:chern}
滑らかな複素代数多様体$Z$と
$\pi\in H^0(Z,\wedge^2T_Z)$に対し、
\[
 \pi^\sharp c_1(K_Z)=0\quad\text{in }H^1(Z,T_Z).
\]
この主張にPoisson条件$[\pi,\pi]=0$は不要である。
\end{lemma}
\begin{proof}
$K_Z$を自明化するアフィン開被覆$\{U_i\}$と、
各$U_i$で消えない最高次微分形式$\mu_i$を取る。
重なりで$\mu_j=g_{ij}\mu_i$、$g_{ij}\in\mathcal O^\times(U_i\cap U_j)$と書く。
$c_1(K_Z)$は、定数倍と符号の規約を除き、
$\{d\log g_{ij}\}$で表される$H^1(\Omega_Z^1)$の類である。

局所ベクトル場$D_i$を
\[
 \iota_{D_i}\mu_i=d(\iota_\pi\mu_i)
\]
で定める。最高次微分形式との縮約は
$T_{U_i}\simeq\Omega^{\dim Z-1}_{U_i}$という同型なので、$D_i$は正則である。
積の微分により
\begin{align*}
 \iota_{D_j-D_i}\mu_i
 &=g_{ij}^{-1}d\bigl(g_{ij}\iota_\pi\mu_i\bigr)
       -d(\iota_\pi\mu_i)\\
 &=d\log g_{ij}\wedge\iota_\pi\mu_i.
\end{align*}
右辺は、縮約と$\pi^\sharp$の符号を固定すれば
$\pm\iota_{\pi^\sharp(d\log g_{ij})}\mu_i$である。
従って$\{\pi^\sharp(d\log g_{ij})\}$は
$\{\pm D_i\}$の\v{C}ech coboundaryであり、そのコホモロジー類は零となる。
符号の選択は消滅の結論に影響しない。
\end{proof}

\begin{lemma}[cubic threefoldの純bivector上ではTodd補正が消える]
\label{pf:lem:todd}
任意の$\pi\in H^0(X,\wedge^2T_X)$について
\[
 I^K(\pi)=I_{\mathrm{HKR}}(\pi).
\]
\end{lemma}
\begin{proof}
\cite[Theorem 1.4]{PFCRVdB}の規約では
$I^K=I_{\mathrm{HKR}}\circ\iota_{\sqrt{\operatorname{td}(X)}}$である。
\[
 \sqrt{\operatorname{td}(X)}
 =1+\tfrac14c_1(T_X)+\theta_2+\cdots,
 \qquad\theta_2\in H^2(X,\Omega_X^2)
\]
を$\pi$に縮約すると、第一補正は
$\tfrac14\pi^\sharp c_1(T_X)\in H^1(T_X)$であり、
$c_1(T_X)=-c_1(K_X)$と補題\ref{pf:lem:chern}から零である。
第二補正は$H^2(X,\mathcal O_X)=0$に値を取る。
三次以上の形式はbivectorには縮約できない。
よって補正は全て消える。逆Todd類を用いる規約でも同じ理由で結論は変わらない。
\end{proof}

\subsection{残余圏への制限写像の核}
$\mathcal K=\operatorname{Ku}(X)$とし、enhancementの完全忠実な包含から誘導される写像を
\[
 \operatorname{res}_X:\operatorname{HH}^2(X)\longrightarrow\operatorname{HH}^2(\mathcal K)
\]
と書く。

\begin{proposition}[Poisson方向の一次消滅]
\label{pf:prop:vanishing}
滑らかなcubic threefold $X$に対し、
\[
 \operatorname{res}_X\circ I_{\mathrm{HKR}}
 \big|_{H^0(X,\wedge^2T_X)}=0.
 \tag{PV}\label{pf:eq:vanishing}
\]
従って、実在する純Poisson一次変形からenhanced残余変形が誘導され、
そのKodaira--Spencer類が上の制限写像で与えられるなら、
残余変形の$\operatorname{HH}^2(\mathcal K)$における一次類は零である。
\end{proposition}
\begin{proof}
\cite[Theorem 3.10]{JLLZ}によれば、
$\alpha'=\operatorname{res}_X\circ I^K$と
\[
 \gamma:\operatorname{HH}^2(\mathcal K)\longrightarrow
 \operatorname{Hom}_{\mathbb C}
 \bigl(H^1(X,\Omega_X^2),H^2(X,\Omega_X^1)\bigr)
\]
の合成はpolyvectorの縮約作用である。
ここで同論文のHochschild homologyの符号では、
$H\Omega_{-1}(X)=H^1(X,\Omega_X^2)$、
$H\Omega_1(X)=H^2(X,\Omega_X^1)$となる。
これら以外の該当するHodge成分は、cubic threefoldのLefschetz定理と
$H^i(X,\mathcal O_X)=0$（$i>0$）により零である。

$\pi\in H^0(X,\wedge^2T_X)$の縮約は
\[
 H^1(X,\Omega_X^2)\xrightarrow{\ \iota_\pi\ }
 H^1(X,\mathcal O_X)=0
\]
なので$\gamma\alpha'(\pi)=0$である。
cubic threefoldについて$\gamma$は単射である
\cite[Theorem 4.4(2)]{JLLZ}。従って$\alpha'(\pi)=0$を得る。
補題\ref{pf:lem:todd}で$I^K$を通常のHKRに戻せば
\eqref{pf:eq:vanishing}となる。
最後の主張は、仮定されたKodaira--Spencer類の制限との両立性による。
\end{proof}

\subsection{明示的復元式に対する制約}
\begin{corollary}[可換復元先の方程式は一次で動かない]
\label{pf:cor:second-order}
$X=V(F)\subset\mathbb P^4$、$R=\mathbb C[[h]]$とし、
特殊ファイバーが$\operatorname{Ku}(X)$であるenhanced残余族$\mathcal K_R$を考える。
その$R/(h^2)$への還元が、命題\ref{pf:prop:vanishing}の意味で
純Poisson一次変形から誘導されていると仮定する。
さらに滑らかな可換cubic族$Y_R=V(F_h)$とenhanced同値
\[
 \mathcal K_R\simeq\operatorname{Ku}(Y_R/R),\qquad Y_0=X
\]
が\emph{既に存在する}と仮定する。
$F_h=F+hG+O(h^2)$と書けば
\[
 [G]=0\quad\text{in }
 \operatorname{Jac}(F)_3,
 \qquad \operatorname{Jac}(F)=\mathbb C[z_0,\ldots,z_4]/(\partial_0F,\ldots,\partial_4F).
 \tag{T0}\label{pf:eq:tangent-zero}
\]
従って$h=0$で恒等となる線形座標変換と方程式の単元倍により
$F_h=F+O(h^2)$とできる。
同じ結論は一次底$\mathbb C[h]/(h^2)$上だけでも成立する。
\end{corollary}
\begin{proof}
命題\ref{pf:prop:vanishing}より、$\mathcal K_R$の一次類は零である。
enhanced同値はこの類を$\operatorname{Ku}(Y_R/R)$の一次類に送るので、
幾何的Kodaira--Spencer類$\kappa_Y\in H^1(X,T_X)$について
\[
 \eta(\kappa_Y)=0,\qquad
 \eta:H^1(X,T_X)\longrightarrow\operatorname{HH}^2(\operatorname{Ku}(X))
\]
となる。$H^1(T_X)$上のTodd補正も$H^2(\mathcal O_X)=0$で消えるので、
ここでも通常のHKRと補正付きHKRは一致する。$\eta$は同型\cite[Theorem 4.4(2)]{JLLZ}であるから
$\kappa_Y=0$である。ここでは特殊ファイバーの同値が恒等である必要はなく、
Hochschild類の零は任意のenhanced同値で零に移る。

超曲面の法束完全列と射影空間のEuler列により、
$H^1(X,T_X)=\operatorname{Jac}(F)_3$で、$\kappa_Y$は$[G]$に対応する。
従って一次形式$\ell_i$を用いて
$G=\sum_{i=0}^4\ell_i\partial_iF$と書ける。
座標変換$z_i\mapsto z_i-h\ell_i(z)$を代入すると一次項が消える。
方程式のスカラー変化もEuler恒等式
$\sum_i z_i\partial_iF=3F$によって同じJacobian部分に含まれる。
\end{proof}

\begin{remark}[一次消滅の適用条件と高次項]
\label{pf:rem:scope}
量子代数の交換子の半古典極限が$\pi$であることだけでは、
射影化した圏の一次変形が純粋に$I_{\mathrm{HKR}}(\pi)$であるとは断定できない。
可換方程式やgluingを一次で同時に動かすと、$H^1(T_X)$成分が加わり得る。
具体的な量子化族に系\ref{pf:cor:second-order}を適用する際は、
この純粋性とenhanced残余変形への制限との両立性を確認する。

確認後は、可換模型の最初の本質的な係数を
$F_h=F+h^2G_2+h^3G_3+\cdots$の$[G_2]\in\operatorname{Jac}(F)_3$から探せばよい。
本節は$G_2$を計算しておらず、その非零性も、全ての奇数次項の消滅も主張しない。
また$\operatorname{HH}^2$での消滅はunderlying残余圏の一次類についての主張であり、
指定rotationを含む全coherentデータの自明化を結論しない。
存在・SM・高次同値の証明は、それぞれ独立に残る。
\end{remark}
\section{Poisson 錐の二次方程式}
$W$ を五次元複素ベクトル空間，
$S=\Sym W=\C[x_0,\ldots,x_4]$ とし，
$F\in S_3$ が滑らかな三次三様体
$X=V(F)\subset\PP(W^*)$ を定めるとする．
随伴公式と三次元における自然同型によって
\[
K_X\simeq\mathcal O_X(-2),\qquad
\bigwedge^2T_X\simeq\Omega_X^1(2),\qquad
H^0(X,\bigwedge^2T_X)\simeq\bigwedge^2W
\]
となる．最後の同型は Euler 列と conormal 列から従う．
$p=(p_{ij})\in\bigwedge^2W$ に対して
\[
\beta_p=\sum_{i<j}p_{ij}\,dx_i\wedge dx_j,\quad
E=\sum_{i=0}^4x_i\partial_{x_i},\quad
\alpha_p=\iota_E\beta_p,\quad
\Omega=dx_0\wedge\cdots\wedge dx_4
\]
とおく．定数ベクトル $v_p$ を
$\beta_p\wedge\beta_p=\iota_{v_p}\Omega$
で定義する．その成分 $q_i(p)$ は五つの Pl\"ucker 二次式である．

\begin{theorem}[二次方程式のスキームとしての Poisson 錐]
\label{thm:cone}
Poisson 条件を Schouten 括弧で与えたアフィンスキームは
\[
\operatorname{Pois}(X)
=\Spec\frac{\C[p_{ij}\mid0\le i<j\le4]}
{(q_0,\ldots,q_4)}
=\widehat{\Gr}(2,W)
\]
である．とくに非零部分の射影化は $\Gr(2,5)$，
錐は既約かつ七次元である．
\end{theorem}
\begin{proof}
$d\alpha_p=2\beta_p$ より
$\alpha_p\wedge d\alpha_p=\iota_E(\beta_p\wedge\beta_p)$ であり，
Euler 恒等式 $E(F)=3F$ から
\[
\alpha_p\wedge d\alpha_p\wedge dF
=v_p(F)\,\iota_E\Omega-3F\,\iota_{v_p}\Omega
\tag{1}
\]
を得る（体積形式の縮約の符号は上式で固定する）．
$X$ 上の Frobenius 条件は $v_p(F)$ が $S/(F)$ で零になること
と同値である．これは二次式なので $S_2$ で零になることと同値である．
線形写像 $W^*\to S_2,\ v\mapsto v(F)$ は単射である．
実際 $v(F)=0$ で $v\ne0$ なら $F$ は座標変換後に一変数に依存せず，
$X$ は特異な射影錐になる．
従って (1) の係数比較から五つの $q_i(p)$ 自身が
Poisson 方程式の生成元である．単なる集合論的同一視ではない．
\end{proof}

\section{頂点の二次無限小近傍}
$\mathfrak m=(p_{ij})$ とおいて
\[
R_2=
\frac{\C[p_{ij}]}
{(q_0,\ldots,q_4)+\mathfrak m^3}
\]
を考える．これは Poisson 錐の頂点の二次無限小近傍の座標環である．
同じ $F$ に対し五変数アフィン空間上の次数二の双ベクトル場
$\Pi_p$ を
\[
\iota_{\Pi_p}\Omega=dF\wedge\beta_p
\tag{2}
\]
で定める．係数は $p_{ij}$ に線形であり，$F$ は
$\Pi_p$ の Casimir である．$\Pi_p$ は Pl\"ucker 方程式の下で
Jacobi 恒等式を満たす．
添字の総和規約で $f_i=\partial_i f$ 等と書き，$S\otimes R_2$ 上の積を
\[
\begin{aligned}
f\star g={}&fg+\frac12\Pi^{ij}f_i g_j
+\frac18\Pi^{ij}\Pi^{k\ell}f_{ik}g_{j\ell}\\
&+\frac1{12}\Pi^{ij}(\partial_j\Pi^{k\ell})
  (f_{ik}g_\ell-f_k g_{i\ell})
\end{aligned}
\tag{3}
\]
で定める．これは Kontsevich の二次の星積を
$\mathfrak m^3=0$ で打ち切ったものである．
\begin{theorem}[普遍二次ジェットの平坦量子化]
$\star$ は結合的な次数付き $R_2$-代数構造を定め，
$F$ はその中心に属する．商
\[
\mathcal A_2=(S\otimes R_2,\star)/(F)
\]
は次数ごとに $R_2$-自由であり，
特別ファイバーは $S/(F)$ である．その一次の交換子は
$\Pi_p$ の Poisson 括弧に一致する．
\end{theorem}
\begin{proof}
(3) の結合性を二次まで計算した際の障害は
Jacobi 恒等式 $[\Pi_p,\Pi_p]=0$ の普遍的線形結合である．
定理\ref{thm:cone}によりこれは $R_2$ 上で零で，
$\mathfrak m^3=0$ より高次項も零である．
一次の交換子は $\Pi^{ij}f_i g_j$ である．
二次項は対称なので交換子の二次部分は消える．
特に (2) により $\Pi_p(dF,-)=0$ だから，
$F$ は $\star$-中心的である．
各次数で $F\star-$ は通常の単射
$S_{d-3}\to S_d$ の冪零摂動である．
その余核は $R_2$ 上の自由加群に同型である．
実際，多項式環の各次数でこの単射の像に
$\C$-線形補空間を選べば，
$\mathfrak m$ に関する冪零行列の消去により
$\operatorname{im}(F\star-)$ の補空間も同じ rank で選べる．
従って $\mathcal A_{2,d}$ は $R_2$-自由である．
\end{proof}

\begin{remark}[意味と限界]
頂点そのもののファイバーは可換な $X$ である．
非被約な二次近傍上の族は，頂点から出る一次方向に対して
$p\wedge p=0$ が二次障害になることを記録する．
この主張は，七次元の被約錐全体上で有限係数の普遍族が存在する
という主張ではない．一般の非零方向に沿う有限係数量子化と
中心三次超曲面および17次元の真の変形については，
リポジトリの別原稿で扱う．
\end{remark}

\section{既知の結果との関係と再生成}
本稿で用いる形式量子化の普遍的な二次積は
Kontsevich の既知の公式である．固定した滑らかな三次三様体の
非零 Poisson 構造の集合論的分類は
Loray--Pereira--Touzet と Mayanskiy に依存する．
本稿の二次方程式の計算はスキーム論的な等式を明記したものである．
既存の有限パラメータでの Sklyanin 型代数，
全混合項の中心三次元，および ambient 非可換 $\PP^4$ の復元は
\texttt{papers/nc\_cubic\_threefold\_moduli.tex} と
nc-cubicの論文A を参照する．
本稿はこれらの長い証明を重複しない．
原稿の再生成には LuaLaTeX と LuaTeX-ja を用いる．


\begin{thebibliography}{99}
\bibitem{Kontsevich}
M. Kontsevich, \emph{Deformation quantization of Poisson manifolds},
Lett. Math. Phys. \textbf{66} (2003), 157--216.
\url{https://arxiv.org/abs/q-alg/9709040}.
Section~4.6.2のHKR比較を参照する。

\bibitem{CT}
A. C\u{a}ld\u{a}raru and J. Tu,
\emph{Curved $A_\infty$ algebras and Landau--Ginzburg models},
New York J. Math. \textbf{19} (2013), 305--342.
Section~2, Remark~2.8, Theorem~2.12.
\url{https://nyjm.albany.edu/j/2013/19-16v.pdf}.

\bibitem{PMayanskiy}
E. Mayanskiy, \emph{Poisson cohomology of two Fano threefolds},
arXiv:1303.5916, Theorem~1およびSection~2の係数計算。
\url{https://arxiv.org/abs/1303.5916}.
同Theorem~1はLoray--Pereira--Touzetの既知結果として述べられている。

\bibitem{JLLZ}
A. Jacovskis, X. Lin, Z. Liu and S. Zhang,
\emph{Infinitesimal categorical Torelli theorems for Fano threefolds},
J. Pure Appl. Algebra \textbf{227} (2023), 107418.
\url{https://doi.org/10.1016/j.jpaa.2023.107418}.
著者版（2022年9月25日）のTheorem~3.10, Lemma~3.13およびTheorem~4.4(2)
（印刷pp.~9--11, 15）を参照する。
\url{https://ajacovskis.github.io/hochschild_only_2.pdf}.

\bibitem{AMChen}
X.-W. Chen, \emph{Module factorizations}, arXiv:2402.11613v3 (2025).
Theorem~2.10, Corollary~2.11, Theorem~7.2。
\url{https://arxiv.org/html/2402.11613v3}.
余核の三角同値を使用し、enhancementは本文でexact圏とdg商により構成する。

\bibitem{AZ}
M. Artin and J. J. Zhang, \emph{Noncommutative projective schemes},
Adv. Math. \textbf{109} (1994), 228--287.
\url{https://doi.org/10.1006/aima.1994.1087}.

\bibitem{Orlov}
D. Orlov, \emph{Derived categories of coherent sheaves and triangulated categories of singularities},
arXiv:math/0503632v3 (2009), Theorem~2.5。
\url{https://arxiv.org/abs/math/0503632}.

\bibitem{BS}
J. Burke and G. Stevenson, \emph{The derived category of a graded Gorenstein ring},
arXiv:1507.00830v1 (2015), Theorem~6.4およびその証明。
\url{https://arxiv.org/abs/1507.00830}.
付録では剰余体上のwindowとVerdier商の向きの同定に用いる。

\bibitem{PFCRVdB}
D. Calaque, C. A. Rossi and M. Van den Bergh,
\emph{C\u{a}ld\u{a}raru's conjecture and Tsygan's formality},
Ann. of Math. (2) \textbf{176} (2012), no.~2, 865--923.
\url{https://doi.org/10.4007/annals.2012.176.2.4}.
Theorem~1.4, p.~872（Todd補正付きHKRとHochschild加群構造）。
\url{https://annals.math.princeton.edu/wp-content/uploads/annals-v176-n2-p04-p.pdf}.

\bibitem{HR}
D. Huybrechts and J. V. Rennemo,
\emph{Hochschild cohomology versus the Jacobian ring, and the Torelli theorem for cubic fourfolds},
Algebraic Geometry \textbf{6} (2019), 76--99.
\href{https://doi.org/10.14231/AG-2019-005}{doi:10.14231/AG-2019-005}.
Sections~2--3。一般のrotation付きHochschild環とJacobian環の比較の背景として参照する。
本稿の低次数の同定には次の\cite{LZ}を用いる。

\bibitem{LZ}
X. Lin and S. Zhang,
\emph{Serre algebra, matrix factorization and categorical Torelli theorem for hypersurfaces},
\href{https://arxiv.org/abs/2310.09927v2}{arXiv:2310.09927v2} (2024).
Theorem~2.2（dg比較とshift）、Theorem~3.5、Remark~3.6（sectorと積）、Section~5。
Proposition~3.7の互いに素という仮定はSerre代数への表示のための条件と区別する。

\bibitem{ExplicitUedaYoshinaga}
K. Ueda and M. Yoshinaga,
\emph{Logarithmic vector fields along smooth divisors in projective spaces},
Hokkaido Math. J. \textbf{38} (2009), 409--415.
\href{https://doi.org/10.14492/hokmj/1258553970}{doi:10.14492/hokmj/1258553970};
\href{https://arxiv.org/abs/0802.2133}{arXiv:0802.2133}.
Lemma~3：勾配イデアルとSebastiani--Thom型の関係。

\bibitem{ExplicitHwang}
J.-M. Hwang,
\emph{Symmetrizer group of a projective hypersurface},
J. Math. Soc. Japan \textbf{78} (2026), 55--62.
\href{https://doi.org/10.2969/jmsj/93469346}{doi:10.2969/jmsj/93469346};
\href{https://arxiv.org/abs/2501.08633}{arXiv:2501.08633}.
Theorems~1.5--1.6（勾配空間のfiberとsymmetrizer群）、
Proposition~2.2（可換代数の構造）。

\bibitem{PSPSmithTate}
S.~P.~Smith and J.~Tate,
\emph{The center of the 3-dimensional and 4-dimensional Sklyanin algebras},
K-Theory \textbf{8} (1994), 19--63, Section~1 (pp.~20--21).
\href{https://sites.math.washington.edu/~smith/Research/tate.pdf}{著者公開版}.
本稿では中心三次元による楕円曲線のtwisted homogeneous coordinate ringへの商を用いる。
有限位数を仮定する同論文後半の中心の分類は用いない。

\bibitem{PSPDeLaet}
K.~De Laet,
\emph{On the center of 3-dimensional and 4-dimensional Sklyanin algebras},
J. Algebra \textbf{487} (2017), 244--268,
\href{https://doi.org/10.1016/j.jalgebra.2017.05.032}{doi:10.1016/j.jalgebra.2017.05.032}.
参照箇所は\href{https://arxiv.org/abs/1612.06158}{arXiv:1612.06158v1}の
Definition~4.1、Sections~4.1--4.2、Theorem~4.4。
点曲線、平行移動$\tau=[a:b:c]$、中心三次元による点加群の消去を用いる。

\bibitem{TV} B. Toën and M. Vaquié,
\emph{Moduli of objects in dg-categories},
Ann. Sci. École Norm. Sup. (4) \textbf{40} (2007), 387--444.
\url{https://doi.org/10.1016/j.ansens.2007.05.001}.
Theorem 3.6, Corollaries 3.15, 3.17, 3.19, 3.21。
\url{https://www.math.univ-toulouse.fr/~vaquie/marque/vaquie_AnnEns40.pdf}.

\bibitem{ERLowen} W. Lowen,
\emph{Hochschild cohomology, the characteristic morphism and derived deformations},
Compos. Math. \textbf{144} (2008), 1557--1580.
\url{https://arxiv.org/abs/0707.2602v1}.
Theorem 4.8, Section 4.4。任意の小拡大に対するdg加群の記述は本文で説明する。

\bibitem{BFKER} M. Ballard, D. Favero and L. Katzarkov,
\emph{A category of kernels for equivariant factorizations and its implications for Hodge theory},
Publ. Math. Inst. Hautes Études Sci. \textbf{120} (2014), 1--111.
\url{https://doi.org/10.1007/s10240-013-0059-9}.
Theorems 5.39, 6.5。

\bibitem{HRER} D. Huybrechts and J. V. Rennemo,
\emph{Hochschild cohomology versus the Jacobian ring and the Torelli theorem for cubic fourfolds},
Algebraic Geometry \textbf{6} (2019), 76--99.
\url{https://doi.org/10.14231/AG-2019-005}.
Lemma 2.19, Sections 2.5, 3.3。

\bibitem{JLLZER} A. Jacovskis, X. Lin, Z. Liu and S. Zhang,
\emph{Infinitesimal categorical Torelli theorems for Fano threefolds},
J. Pure Appl. Algebra \textbf{227} (2023), 107418.
\url{https://arxiv.org/abs/2203.08187}.
Theorem 4.4(2)。

\bibitem{HuyER} D. Huybrechts,
\emph{The K3 category of a cubic fourfold},
Compos. Math. \textbf{153} (2017), 586--620.
\url{https://doi.org/10.1112/S0010437X16008137}.
Proposition 3.12, Section 3.5。

\bibitem{StacksER} The Stacks Project Authors, \emph{The Stacks Project}.
\url{https://stacks.math.columbia.edu/tag/02HX};
\url{https://stacks.math.columbia.edu/tag/055U}.
2026年11月13日確認。smoothnessの無限小判定とsmooth射のétale局所section。

\bibitem{ERLinZhang} X. Lin and S. Zhang,
\emph{Categorical Torelli and Fourier--Mukai partners of hypersurfaces},
arXiv:2611.05915v1 (2026).
\url{https://arxiv.org/abs/2611.05915}.
Section 2.1, Theorem 2.2, Proposition 2.4（元原稿の参照箇所）。
本稿の判定法の証明の入力ではなく、関連研究として挙げる。

\bibitem{AMSasha} Sasha,
Answer to \emph{Is the property of being a dg generator open?},
MathOverflow, 5 December 2016.
\url{https://mathoverflow.net/questions/256391/is-the-property-of-being-a-dg-generator-open}.
射影直線上の生成対象がskyscraperへ変形して生成性を失う例。

\bibitem{LP-AZ}
M. Artin and J. J. Zhang, \emph{Noncommutative projective schemes},
Adv. Math. \textbf{109} (1994), 228--287.
\url{https://doi.org/10.1006/aima.1994.1087}.

\bibitem{LP-BS}
J. Burke and G. Stevenson, \emph{The derived category of a graded Gorenstein ring},
arXiv:1507.00830v1 (2015).
\url{https://arxiv.org/abs/1507.00830}.
Section~4, Lemmas~6.1--6.2, Theorem~6.4およびその証明。

\bibitem{AMStacks}
The Stacks Project Authors, \emph{The Stacks Project}.
Tags \href{https://stacks.math.columbia.edu/tag/0BQT}{0BQT}
（有限線束生成子）および
\href{https://stacks.math.columbia.edu/tag/09M1}{09M1}
（compact対象とperfect対象）。

\bibitem{PHartshorne}
R.~Hartshorne,
\emph{Algebraic Geometry}, Graduate Texts in Mathematics 52,
Springer, 1977.
\href{https://doi.org/10.1007/978-1-4757-3849-0}{doi:10.1007/978-1-4757-3849-0}.

\bibitem{PEtingofLie}
P.~Etingof,
\emph{Lie Groups and Lie Algebras}, MIT lecture notes, 2024 course version.
\href{https://live.ocw.mit.edu/courses/18-755-lie-groups-and-lie-algebras-ii-spring-2024/mit18_755_s24_lec_full.pdf}{MIT OpenCourseWare 公開原稿}.
本文では Proposition~16.9、\S19.4、Theorem~25.17、\S32.2 を参照する。

\bibitem{PMoriSmith}
I.~Mori and S.~P.~Smith,
The classification of 3-Calabi--Yau algebras with 3 generators and 3 quadratic relations,
\href{https://arxiv.org/abs/1502.07403v2}{arXiv:1502.07403v2} (2016).
Theorems~1.3, 1.4, 4.2 and Corollaries~1.6, 2.7, 2.9.
点曲線のHessian公式はLemma~7.1, Theorem~7.2, pp.~23--25。

\bibitem{PBergerTaillefer}
R.~Berger and R.~Taillefer,
Poincar\'e--Birkhoff--Witt deformations of Calabi--Yau algebras,
\emph{Journal of Noncommutative Geometry} \textbf{1} (2007), 241--270.
\href{https://arxiv.org/abs/math/0610112}{arXiv:math/0610112}.
本文で用いる potential による PBW 変形は Theorem~1.1(i)。

\bibitem{PKontsevich}
M.~Kontsevich,
Deformation quantization of Poisson manifolds,
\emph{Letters in Mathematical Physics} \textbf{66} (2003), 157--216.
\href{https://doi.org/10.1023/B:MATH.0000027508.00421.bf}{doi:10.1023/B:MATH.0000027508.00421.bf};
原稿名は \emph{Deformation quantization of Poisson manifolds, I},
\href{https://arxiv.org/abs/q-alg/9709040}{arXiv:q-alg/9709040} (1997).
参照箇所は原稿のSections~4.4--4.6（Maurer--Cartan変形functorとformality）、
6.4（$L_\infty$関係）、8.2および8.3.3.1（中心のtangent mapとwheel）。

\bibitem{PDolgushev}
V.~Dolgushev,
The Van den Bergh duality and the modular symmetry of a Poisson variety,
\href{https://arxiv.org/abs/math/0612288v3}{arXiv:math/0612288v3} (2007).
Theorem~3 および Proposition~1。

\bibitem{BondalGualtieriPym}
M.~Gualtieri and B.~Pym,
\emph{Poisson modules and degeneracy loci},
Proc. Lond. Math. Soc. \textbf{107} (2013), 627--654.
\href{https://doi.org/10.1112/plms/pds090}{doi:10.1112/plms/pds090};
\href{https://arxiv.org/abs/1203.4293}{arXiv:1203.4293}.
Bondal予想の定式化、四次元Fanoの場合の証明。

\bibitem{BondalDruelPereiraPymTouzet}
S.~Druel, J.~V.~Pereira, B.~Pym and F.~Touzet,
\emph{Bondal's conjecture in dimension five},
\href{https://arxiv.org/abs/2606.03424}{arXiv:2606.03424} (2026).
Fano五次元の全場合、一般の奇数次元の部分結果。

\bibitem{FuHwangCI}
B. Fu and J.-M. Hwang,
\emph{Isotrivial VMRT-structures of complete intersection type},
Asian Journal of Mathematics \textbf{22} (2018), 333--354.
Lemma 4.1 (Br\"uckmannの消滅定理を引用).
\url{https://doi.org/10.4310/AJM.2018.v22.n2.a10}.

\bibitem{BorelRemmert}
A. Borel and R. Remmert,
\emph{\"Uber kompakte homogene K\"ahlersche Mannigfaltigkeiten},
Mathematische Annalen \textbf{145} (1962), 429--439, Theorem I.
\url{https://wglobke.github.io/publications/BorelRemmert1962.pdf}.

\bibitem{EG}
P. Etingof and V. Ginzburg,
\emph{Noncommutative del Pezzo surfaces and Calabi--Yau algebras},
Journal of the European Mathematical Society \textbf{12} (2010), 1371--1416.
Sections 3.5--3.6, Theorem 3.6.2, Proposition 8.1.1(i) およびその証明。
\url{https://doi.org/10.4171/JEMS/235}.
\href{https://arxiv.org/abs/0709.3593v4}{arXiv:0709.3593v4}.

\bibitem{LBSVNE12} L. Le Bruyn, S. P. Smith and M. Van den Bergh,
\emph{Central extensions of three dimensional Artin--Schelter regular algebras},
Math. Z. \textbf{222} (1996), 171--212.
Theorem 2.6, Corollary 2.7, Remark 2.8, Theorem 3.1.3.
\url{https://sites.math.washington.edu/~smith/Research/Central-extns.pdf}.

\bibitem{MSNE12} I. Mori and S. P. Smith,
\emph{The classification of 3-Calabi--Yau algebras with 3 generators and 3 quadratic relations},
Math. Z. \textbf{287} (2017), 215--241.
Lemma 4.1, Theorem 4.2, Corollary 4.3; arXiv v2, pp. 15--17.
\url{https://arxiv.org/abs/1502.07403}; doi:10.1007/s00209-016-1824-5.

\bibitem{BergNE12} G. M. Bergman,
\emph{The diamond lemma for ring theory}, Adv. Math. \textbf{29} (1978), 178--218, Theorem 1.2.
\url{https://doi.org/10.1016/0001-8708(78)90010-5}.

\bibitem{BTNE12} R. Berger and R. Taillefer,
\emph{Poincar\'e--Birkhoff--Witt deformations of Calabi--Yau algebras},
J. Noncommut. Geom. \textbf{1} (2007), 241--270.
Theorems 3.1 and 3.6, especially the potential bimodule complex in the proof of Theorem 3.6.
\url{https://arxiv.org/abs/math/0610112}; doi:10.4171/JNCG/6.

\bibitem{SmithTate}
S. P. Smith and J. Tate,
\emph{The center of the 3-dimensional and 4-dimensional Sklyanin algebras},
K-Theory \textbf{8} (1994), 19--63.
Corollary 2.3 (p.~26); Theorem 3.7(c) (p.~38).
\url{https://doi.org/10.1007/BF00962089};
\url{https://sites.math.washington.edu/~smith/Research/tate.pdf}.

\bibitem{DeLaetLeBruyn}
K. De Laet and L. Le Bruyn,
\emph{The geometry of representations of 3-dimensional Sklyanin algebras},
Algebras and Representation Theory \textbf{18} (2015), 761--776.
arXiv:1405.1158v1, Theorem 1 and Section 2 (canonical central cubic formula);
the result is a reformulation of the Artin--Schelter and Smith--Tate formula.
\url{https://arxiv.org/abs/1405.1158}.

\bibitem{ArtebaniDolgachev}
M. Artebani and I. Dolgachev,
\emph{The Hesse pencil of plane cubic curves},
L'Enseignement Math\'ematique (2) \textbf{55} (2009), no.~3--4, 235--273,
Lemma 1 (every smooth plane cubic admits a Hesse canonical form).
\url{https://doi.org/10.4171/LEM/55-3-3};
\url{https://arxiv.org/abs/math/0611590}.

\bibitem{LPTPoissonCubic}
F.~Loray, J.~V.~Pereira and F.~Touzet,
\emph{Foliations with trivial canonical bundle on Fano 3-folds},
Math. Nachr. \textbf{286} (2013), 921--940.
\href{https://arxiv.org/abs/1112.3871}{arXiv:1112.3871}.

\bibitem{PStacksQuotient}
The Stacks Project Authors,
\emph{The Stacks Project},
Criteria for Representability, Theorem~97.17.2
(Tag~\texttt{06FI}), flat locally finitely presented groupoids and algebraic quotient stacks.
\url{https://stacks.math.columbia.edu/tag/06FI}.

\bibitem{PGLefschetz}
A.~Grothendieck,
\emph{Cohomologie locale des faisceaux coh\'erents et th\'eor\`emes
de Lefschetz locaux et globaux (SGA~2)},
Expos\'e~XII, Corollaire~3.7 (Picard groups of sufficiently
high-dimensional smooth projective complete intersections).
\url{https://grothendiecksga.com/read/sga2/en/XII-3.html}.

\bibitem{PMatsumuraMonsky}
H.~Matsumura and P.~Monsky,
\emph{On the automorphisms of hypersurfaces},
J. Math. Kyoto Univ. \textbf{3} (1963/64), 347--361.
\url{https://doi.org/10.1215/kjm/1250524785}.

\bibitem{PAZAmbient}
M.~Artin and J.~J.~Zhang,
\emph{Noncommutative projective schemes},
Adv. Math. \textbf{109} (1994), 228--287,
\href{https://doi.org/10.1006/aima.1994.1087}{doi:10.1006/aima.1994.1087}.

\bibitem{LP-Keller} B. Keller, \emph{On differential graded categories},
International Congress of Mathematicians, Vol. II (2006), 151--190, Theorem 3.8.
\url{https://arxiv.org/abs/math/0601185}.

\bibitem{LP-Orlov} D. Orlov,
\emph{Smooth and proper noncommutative schemes and gluing of DG categories},
Adv. Math. \textbf{302} (2016), 59--105, Section 3.4, Proposition 3.31.
\url{https://arxiv.org/abs/1402.7364}.

\bibitem{PetersenTransfer} D.~Petersen,
\emph{A closer look at Kadeishvili's theorem},
Higher Structures \textbf{4} (2020), no.~2, 211--221.
\url{https://doi.org/10.21136/HS.2020.14}.

\bibitem{BurkeTransfer} J.~Burke,
\emph{Transfer of $A_\infty$ structures to projective resolutions},
arXiv:1801.08933 (2018).
\url{https://arxiv.org/abs/1801.08933}.

\bibitem{AinftyPerry2019} A.~Perry,
\emph{Noncommutative homological projective duality},
Advances in Mathematics \textbf{350} (2019), 877--972,
Lemma 4.15(3).
\url{https://doi.org/10.1016/j.aim.2019.04.052}.

\bibitem{UeyamaIsolated}
K. Ueyama,
\emph{Noncommutative graded Gorenstein isolated singularities},
in \emph{Proceedings of the 45th Symposium on Ring Theory and Representation Theory}
(Shinshu University, September 7--9, 2012),
K. Koike (ed.), February 2013, 158--163.
Theorem 7, printed p.~161.
\url{https://ring-theory-japan.com/ring/oldmeeting/2012/report2012/proceedings.pdf}.

\bibitem{ZhangTwist}
J.~J. Zhang,
\emph{Twisted graded algebras and equivalences of graded categories},
Proceedings of the London Mathematical Society (3) \textbf{72} (1996), no.~2, 281--311.
Theorems 1.1 and 1.3.
\url{https://doi.org/10.1112/plms/s3-72.2.281}.
\url{https://deepblue.lib.umich.edu/bitstream/handle/2027.42/135651/plms0281.pdf}.

\bibitem{RRZ}
M. Reyes, D. Rogalski and J.~J. Zhang,
\emph{Skew Calabi--Yau algebras and homological identities},
Advances in Mathematics \textbf{264} (2014), 308--354.
Theorems 0.3 and 5.4; Example 5.5.
\url{https://doi.org/10.1016/j.aim.2014.07.010}.
\url{https://arxiv.org/abs/1302.0437}.

\bibitem{Griffiths}
P.~A. Griffiths,
\emph{On the periods of certain rational integrals.~I},
Annals of Mathematics (2) \textbf{90} (1969), no.~3, 460--495.
Section 8, especially Theorems 8.1 and 8.3 and formula (8.6).
\url{https://doi.org/10.2307/1970746}.
\url{https://www.cimat.mx/~luis/seminarios/Teoria-de-Hodge/period_integral_I.pdf}.

\bibitem{KuznetsovFCY}
A. Kuznetsov,
\emph{Calabi--Yau and fractional Calabi--Yau categories},
Journal f\"ur die reine und angewandte Mathematik \textbf{753} (2019), 239--267.
Lemmas 2.6, 3.13 and 3.19; Corollary 4.1.
\url{https://doi.org/10.1515/crelle-2017-0004}.

\bibitem{CHL}
C.-H. Cho, H. Hong and S.-C. Lau,
\emph{Noncommutative homological mirror functor},
Memoirs of the American Mathematical Society \textbf{271} (2021), no.~1326.
Chapter 5, especially Theorems 5.15 and 5.16.
\url{https://doi.org/10.1090/memo/1326}.
\url{https://arxiv.org/abs/1512.07128}.

\bibitem{Lee}
S. Lee,
\emph{Noncommutative homological mirror symmetry of elliptic curves},
Kyoto Journal of Mathematics \textbf{61} (2021), no.~3, 723--743.
Theorem 6.3 and the corollary in the introduction.
\url{https://arxiv.org/abs/2004.09145}.

\bibitem{CisinskiTabuada}
D.-C. Cisinski and G. Tabuada,
\emph{Lefschetz and Hirzebruch--Riemann--Roch formulas via noncommutative motives},
Journal of Noncommutative Geometry \textbf{8} (2014), 1171--1190.
Section 4 and Example 5.6.
\url{https://doi.org/10.4171/JNCG/183}.
\url{https://www.math.univ-toulouse.fr/~dcisinsk/ncLHRR.pdf}.

\bibitem{MathewHH}
A. Mathew,
\emph{Kaledin's degeneration theorem and topological Hochschild homology},
Geometry \& Topology \textbf{24} (2020), no.~6, 2675--2708.
Theorem 1.3.
\url{https://doi.org/10.2140/gt.2020.24.2675}.
\url{https://arxiv.org/abs/1710.09045}.

\bibitem{PerryIHC}
A. Perry,
\emph{The integral Hodge conjecture for two-dimensional Calabi--Yau categories},
Compositio Mathematica \textbf{158} (2022), no.~2, 287--333.
Lemmas 3.4 and 3.6; Theorem 3.5.
\url{https://doi.org/10.1112/S0010437X22007266}.

\bibitem{Beilinson}
A. A. Beilinson,
\emph{Coherent sheaves on \(\PP^n\) and problems of linear algebra},
Functional Analysis and Its Applications \textbf{12} (1978), 214--216.
\url{https://doi.org/10.1007/BF01681436}.
原著Section 4の対角Koszul分解を確認。
\url{https://www.mathnet.ru/eng/faa2009}.

\bibitem{DDL}
O. De Deken and W. Lowen,
\emph{Abelian and derived deformations in the presence of \(\ZZ\)-generating geometric helices},
Journal of Noncommutative Geometry \textbf{5} (2011), 477--505.
Theorem 4.5; Propositions 5.5, 5.8, 5.9; Lemma 5.14.
\url{https://doi.org/10.4171/JNCG/83}.

\bibitem{Keller}
B. Keller, \emph{Deriving DG categories},
Annales scientifiques de l'École Normale Supérieure \textbf{27} (1994), 63--102.
\url{https://doi.org/10.24033/asens.1689}.

\bibitem{Mathew}
A. Mathew, \emph{The Galois group of a stable homotopy theory},
Advances in Mathematics \textbf{291} (2016), 403--541.
Proposition 3.28, Remark 3.29, Proposition 3.35.
\url{https://arxiv.org/abs/1404.2156}.

\bibitem{CCKM}
T. Cassidy, A. Conner, E. Kirkman and W. F. Moore,
\emph{Periodic free resolutions from twisted matrix factorizations},
Journal of Algebra \textbf{455} (2016), 137--163.
Proposition 5.7, Theorem 5.8.
\url{https://doi.org/10.1016/j.jalgebra.2016.01.037}.

\bibitem{BFK}
M. Ballard, D. Favero and L. Katzarkov,
\emph{A category of kernels for equivariant factorizations and its implications for Hodge theory},
Publications Mathématiques de l'IHÉS \textbf{120} (2014), 1--111.
Theorem 5.39のKoszul cohomologyによる公式、Theorem 5.15とCorollary 5.18のequivariant外部積.
\url{https://doi.org/10.1007/s10240-013-0059-9}.
確認した公刊本文 \url{https://www.numdam.org/article/PMIHES_2014__120__1_0.pdf}.

\bibitem{Huy}
D. Huybrechts, \emph{The K3 category of a cubic fourfold},
Compositio Mathematica \textbf{153} (2017), 586--620.
Proposition 3.12およびSection 3.5.
\url{https://doi.org/10.1112/S0010437X16008137}.

\bibitem{Lurie}
J. Lurie, \emph{Derived Algebraic Geometry X: Formal Moduli Problems},
2011. Theorems 5.1.9, 5.3.16およびProposition 5.3.21.
\url{https://www.math.ias.edu/~lurie/papers/DAG-X.pdf}.

\bibitem{BKP}
A. Blanc, L. Katzarkov and P. Pandit,
\emph{Generators in formal deformations of categories},
Compositio Mathematica \textbf{154} (2018).
本文で用いるのはarXiv版のDefinition 3.3およびSection 3.3.
\url{https://arxiv.org/abs/1705.00655}.

\bibitem{KL}
B. Keller and W. Lowen,
\emph{On Hochschild cohomology and Morita deformations},
International Mathematics Research Notices (2009).
\url{https://webusers.imj-prg.fr/~bernhard.keller/publ/KellerLowenHochDefo.pdf}.

\bibitem{Walton}
C. Walton, \emph{Representation theory of three-dimensional Sklyanin algebras},
Nuclear Physics B \textbf{860} (2012), 167--185.
Theorem 2.14および式(2.15)--(2.16).
\url{https://doi.org/10.1016/j.nuclphysb.2012.02.015}.

\bibitem{CGWZ}
K. Chan, J. Gaddis, R. Won and J. J. Zhang,
\emph{Ozone groups of Artin--Schelter regular algebras satisfying a polynomial identity},
arXiv:2312.13014, v2 (2024). Introductionの非退化Sklyanin代数の記述。
\url{https://arxiv.org/abs/2312.13014}.

\bibitem{Elagin}
A. Elagin, \emph{On equivariant triangulated categories},
arXiv:1403.7027, Theorem 4.2.
\url{https://arxiv.org/abs/1403.7027}.

\bibitem{ELS}
A. Elagin, V. A. Lunts and O. M. Schn\"urer,
\emph{Smoothness of Derived Categories of Algebras},
Moscow Mathematical Journal \textbf{20} (2020), no. 2, 277--309.
確認した一次資料はarXiv:1810.07626v2 (22 March 2019),
Theorem 4.4 (Verdier局所化), Proposition 4.10とTheorem 4.15 (稠密開集合と生成子),
Theorem 4.18 (外部積生成), Theorem 5.1 (有限中心代数の有界導来圏のsmoothness).
\url{https://arxiv.org/abs/1810.07626};
\url{https://www.mathjournals.org/mmj/2020-020-002/2020-020-002-003.html}.

\bibitem{Stacks}
The Stacks Project Authors, \emph{The Stacks Project}.
Tag 02HX (Lemma 37.12.1, Noetherian底上の無限小smoothness判定),
Tag 055U (Lemma 37.38.6, smooth morphismのétale局所section).
\url{https://stacks.math.columbia.edu/tag/02HX};
\url{https://stacks.math.columbia.edu/tag/055U}.
参照日2026年11月4日。

\bibitem{Pirozhkov}
D. Pirozhkov, \emph{A categorical Torelli theorem for hypersurfaces},
Bulletin of the London Mathematical Society \textbf{56} (2024), 3075--3089.
\url{https://doi.org/10.1112/blms.13117}.
公刊版のabstractで、cubicに対する仮定 \(n>3\) を確認。
本稿の全次元の一意性には \cite[Theorem 1.7]{LZ} を用いる。

\bibitem{HuyUpdate}
D. Huybrechts, \emph{The K3 category of a cubic fourfold: an update},
Beiträge zur Algebra und Geometrie / Contributions to Algebra and Geometry (2025).
Section 9: Unknown deformations.
\url{https://doi.org/10.1007/s13366-025-00788-9}.

\bibitem{ChanMulti}
D. Chan, \emph{Twisted multi-homogeneous coordinate rings},
Journal of Algebra \textbf{223} (2000), 438--456.
Theorem 3.4; proof of Theorem 5.2, Claim 5.3 and Lemmas 5.4--5.6;
Corollaries 5.8--5.9.
\url{https://doi.org/10.1006/jabr.1999.7938}.
確認した著者版 \url{https://web.maths.unsw.edu.au/~danielch/paper/tbcr.ps}, pp.8--9,14--17.

\bibitem{KeelerMulti}
D. S. Keeler, \emph{Noncommutative ampleness for multiple divisors},
Journal of Algebra \textbf{265} (2003), no. 1, 299--311.
Theorem 2.7, Corollary 2.8 (任意有限個のbimoduleのNC-ampleness);
Hypothesis 3.1, Theorem 3.3, Corollary 3.5 (二因子のnoetherian性).
番号はarXiv一次TeXソースで確認し、刊行版との番号一致は未照合。
\url{https://doi.org/10.1016/S0021-8693(03)00126-1};
\url{https://arxiv.org/abs/math/0210417};
\url{https://arxiv.org/e-print/math/0210417}.

\bibitem{WWY3}
C. Walton, X. Wang and M. Yakimov,
\emph{Poisson geometry of PI three-dimensional Sklyanin algebras},
Proceedings of the London Mathematical Society \textbf{118} (2019), 1471--1500.
Theorem 1.3(1),(2a), verified in arXiv:1704.04975v3, p.~3.
\url{https://doi.org/10.1112/plms.12220};
\url{https://arxiv.org/abs/1704.04975}.

\bibitem{ToenFinitude}
B. Toën, \emph{Finitude homotopique des dg-algèbres propres et lisses},
Proceedings of the London Mathematical Society (3) \textbf{98} (2009), 217--240.
Lemmas 3.13--3.14. Lemma 3.14の結論は部分前層のquasi-compactnessであり、一般の開性ではない。
\url{https://doi.org/10.1112/plms/pdn027};
確認した著者版 \url{https://www.math-beb.fr/wp-content/uploads/2026/10/finitesatdgalg.pdf}.

\bibitem{SnowdenDVR}
A. Snowden, \emph{Lecture 8: Elliptic curves over DVRs},
Math 679 lecture notes, University of Michigan, 2013.
Section ``Reduction of torsion points'', good reductionでのprime-to-residue-characteristic torsionの還元の単射性。
\url{https://websites.umich.edu/~asnowden/teaching/2013/679/L08.html}.
参照日2026年11月6日。

\bibitem{HMY}
J.-W. He, X.-C. Ma and Y. Ye, \emph{Generalized Kn\"orrer's Periodicity Theorem},
arXiv:2107.06438v2 (16 August 2021).
対象は中心二次元の商であり、本稿の三次の場合には直接適用しない。
\url{https://arxiv.org/abs/2107.06438}.

\bibitem{OrlovGlue}
D. Orlov,
\emph{Smooth and proper noncommutative schemes and gluing of DG categories},
Advances in Mathematics \textbf{302} (2016), 59--105.
Theorem 3.25 (smoothness of dg gluing; Proposition 3.24はsmoothnessからregularityのみ)。
\url{https://doi.org/10.1016/j.aim.2016.07.014};
\url{https://arxiv.org/abs/1402.7364}.

\bibitem{BurkeStevenson}
J. Burke and G. Stevenson,
\emph{The derived category of a graded Gorenstein ring},
arXiv:1507.00830v1 (2015), Definition 5.1 and Theorem 6.4.
\url{https://arxiv.org/abs/1507.00830}.

\bibitem{MBNegron}
C. Negron, \emph{The cup product on Hochschild cohomology via twisting cochains and applications to Koszul rings},
\href{https://arxiv.org/abs/1304.0527v4}{arXiv:1304.0527v4} (2016), Theorem~5.3.

\bibitem{MBTu}
J. Tu, \emph{Matrix factorizations via Koszul duality},
Compos. Math. \textbf{150} (2014), 1549--1578;
\href{https://arxiv.org/abs/1009.4151v2}{arXiv:1009.4151v2}, Appendix~A, Lemma~A.1.

\bibitem{MBLowen}
W. Lowen, \emph{Hochschild cohomology, the characteristic morphism and derived deformations},
Compos. Math. \textbf{144} (2008), 1557--1580;
\href{https://arxiv.org/abs/0707.2602}{arXiv:0707.2602}, Propositions~3.11, 3.13.

\bibitem{MBGHLR}
A. Guan, J. Holstein, A. Lazarev and K. Rasmussen,
\emph{Hochschild cohomology of the second kind: Koszul duality and Morita invariance},
\href{https://arxiv.org/abs/2312.16645v5}{arXiv:2312.16645v5} (2026),
Theorems~4.7, 6.14 and Remark~6.16.

\bibitem{MBLinZhang}
X. Lin and S. Zhang, \emph{Serre algebra, matrix factorization and categorical Torelli theorem for hypersurfaces},
\href{https://arxiv.org/abs/2310.09927v2}{arXiv:2310.09927v2} (2024),
Theorem~3.5 and Remark~3.6. 本稿では体上の特殊ファイバーの同定に用いる。

\bibitem{MBStacks}
The Stacks Project Authors, \emph{The Stacks Project},
\href{https://stacks.math.columbia.edu/tag/00T7}{Tag 00T7}.
\'Etale近傍の構成には相対Jacobianの可逆性を用いる。

\bibitem{FARStacks}
The Stacks Project Authors,
\emph{The Stacks Project}, Tags 00MC and 0BNH,
\url{https://stacks.math.columbia.edu/tag/00MC},
\url{https://stacks.math.columbia.edu/tag/0BNH}.
完備化の忠実平坦性、有限加群の完備化。2026年11月10日確認。

\bibitem{FARLinZhang}
X. Lin and S. Zhang,
\emph{Serre algebra, matrix factorization and categorical Torelli theorem for hypersurfaces},
arXiv:2310.09927v2 (2024), Theorem 3.5, Remark 3.6, Section 5.
\url{https://arxiv.org/html/2310.09927v2}.

\bibitem{FARHwang}
J.-M. Hwang,
\emph{Symmetrizer group of a projective hypersurface},
J. Math. Soc. Japan \textbf{78} (2026), 55--62;
arXiv:2501.08633v1, Sections 1--3.
\url{https://arxiv.org/html/2501.08633v1}.

\bibitem{GCDBFK}
M.~Ballard, D.~Favero and L.~Katzarkov,
\emph{Orlov spectra: bounds and gaps},
arXiv:1012.0864v3 (2013), Section~2 and Theorem~5.15.
\url{https://arxiv.org/abs/1012.0864}.
参照したPDFの第41頁の定理と第6頁の基礎体・生成段数の規約を確認。

\bibitem{CGTu}
J.~Tu,
\emph{Matrix factorizations via Koszul duality},
Compositio Mathematica \textbf{150} (2014), 1549--1578.
\url{https://doi.org/10.1112/S0010437X14007295}.
\url{https://arxiv.org/abs/1009.4151}.
参照箇所はarXiv v2のSections~2.2, 6.1およびTheorem~6.1。
curved twisted complexとgradingの枠組みとして参照し、
非可換Artin基底上の本稿の比較定理全体を引用するものではない。

\bibitem{ROGStevenson}
G.~Stevenson, \emph{Subcategories of singularity categories via tensor actions},
\href{https://arxiv.org/abs/1105.4698v3}{arXiv:1105.4698v3},
Theorem~6.13 (PDF p.~26),
\href{https://doi.org/10.1112/S0010437X1300746X}{doi:10.1112/S0010437X1300746X}.
原論文の非次数付き分類を引用し、次数シフト軌道への移行は本文で証明する。

\bibitem{ROGDyckerhoff}
T.~Dyckerhoff, \emph{Compact generators in categories of matrix factorizations},
Duke Math. J. \textbf{159} (2011), 223--274;
\href{https://arxiv.org/abs/0904.4713v5}{arXiv:0904.4713v5},
\S2.4, \S5.6, Theorem~5.9. 本文の三次Euler計算は直接も証明する。

\bibitem{ROGStacks}
The Stacks Project Authors, \emph{The Stacks Project},
\href{https://stacks.math.columbia.edu/tag/09SM}{Tag 09SM},
\S13.37, Lemmas~13.37.3--4 and Proposition~13.37.6;
\href{https://stacks.math.columbia.edu/tag/00FV}{Tag 00FV}, Theorem~10.34.1;
\href{https://stacks.math.columbia.edu/tag/0BQQ}{Tag 0BQQ}, \S36.16 (projective subscheme generators).

\bibitem{ROGTV}
B.~To\"en and M.~Vaqui\'e, \emph{Moduli of objects in dg-categories},
\href{https://arxiv.org/abs/math/0503269v5}{arXiv:math/0503269v5},
Lemmas~2.8 and 2.10 (PDF pp.~21--24).
片側perfect性の比較と、perfect対象・射のfiltered-colimit降下に用いる。

\bibitem{RASBFK}
M. Ballard, D. Favero and L. Katzarkov,
\emph{A category of kernels for equivariant factorizations and its implications for Hodge theory},
Publ. Math. IH\'ES \textbf{120} (2014), 1--111.
\href{https://doi.org/10.1007/s10240-013-0059-9}{doi:10.1007/s10240-013-0059-9};
\href{https://arxiv.org/abs/1105.3177}{arXiv:1105.3177}.

\bibitem{RASDyck}
T. Dyckerhoff,
\emph{Compact generators in categories of matrix factorizations},
Duke Math. J. \textbf{159} (2011), 223--274.
\href{https://arxiv.org/abs/0904.4713}{arXiv:0904.4713}.

\bibitem{RASLZ}
X. Lin and S. Zhang,
\emph{Serre algebra, matrix factorization and categorical Torelli theorem for hypersurfaces},
\href{https://arxiv.org/html/2310.09927v2}{arXiv:2310.09927v2},
Theorem~3.5 and Remark~3.6.

\bibitem{KJDito}
G. Dito,
\emph{Kontsevich star-product on the dual of a Lie algebra},
Letters in Mathematical Physics \textbf{48} (1999), 307--322.
\href{https://doi.org/10.1023/A:1007643618406}{doi:10.1023/A:1007643618406};
\href{https://arxiv.org/abs/math/9905080}{arXiv:math/9905080}.
式~(3)（arXiv版の印刷頁5）は一般Poisson tensorに対する二次星積とeye項の係数を与える。
本稿では同論文のbivectorを$\pi/2$に置き換える。Section~3はWeyl対称化を扱う。

\bibitem{KJShoikhet}
B. Shoikhet,
\emph{Vanishing of the Kontsevich integrals of the wheels},
\href{https://arxiv.org/abs/math/0007080}{arXiv:math/0007080} (2000).
Section~1.3はboundary vertexのないtangent mapの偶数wheelの重み$W_{2k}=0$を述べる。
本稿では$W_2=0$を用い、星積の二つのboundary vertexを持つeye項とは区別する。

\bibitem{MNFDolgushev}
V.~Dolgushev,
\emph{Covariant and equivariant formality theorems},
Advances in Mathematics \textbf{191} (2005), 147--177.
\href{https://doi.org/10.1016/j.aim.2004.02.001}{doi:10.1016/j.aim.2004.02.001};
\href{https://arxiv.org/abs/math/0307212}{arXiv:math/0307212}.
有限群不変な接続を使うformalityの同変化。本文のaffine線形作用ではKontsevichのGL同変性を使う。

\bibitem{AROBFK}
M.~Ballard, D.~Favero and L.~Katzarkov,
\emph{A category of kernels for equivariant factorizations and its implications for Hodge theory},
Publ. Math. Inst. Hautes Études Sci. \textbf{120} (2014), 1--111,
\href{https://doi.org/10.1007/s10240-013-0059-9}{doi:10.1007/s10240-013-0059-9},
Theorem~5.39 (printed pp.~82--84).

\bibitem{AROLZ}
X.~Lin and S.~Zhang,
\emph{Serre algebra, matrix factorization and categorical Torelli theorem for hypersurfaces},
\href{https://arxiv.org/abs/2310.09927v2}{arXiv:2310.09927v2},
Theorem~3.5, Remark~3.6.

\bibitem{CDGHLR}
A.~Guan, J.~Holstein, A.~Lazarev and K.~Rasmussen,
\emph{Hochschild cohomology of the second kind: Koszul duality and Morita invariance},
arXiv:2312.16645v5 (2026-09-17).
\href{https://arxiv.org/abs/2312.16645v5}{arXiv:2312.16645v5}.
Theorem~4.7, Theorem~6.14 and Remark~6.16; pp.~32--33 of the PDF.
本稿では包絡代数上の導来圏と曲率付き双加群の導来圏の区別に用いる。

\bibitem{OSTV}
B.~To\"en and M.~Vaqui\'e,
\emph{Moduli of objects in dg-categories},
Ann. Sci. \`Ecole Norm. Sup. (4) \textbf{40} (2007), 387--444.
\href{https://doi.org/10.1016/j.ansens.2007.05.001}{doi:10.1016/j.ansens.2007.05.001};
\href{https://arxiv.org/abs/math/0503269}{arXiv:math/0503269}.
Lemma~2.8(2),(3), Lemma~2.10 (printed pp.~404--407).

\bibitem{OSPopescu}
The Stacks Project Authors, \emph{The Stacks Project},
\href{https://stacks.math.columbia.edu/tag/07GC}{Tag 07GC},
Theorem~16.12.1: regular ring maps are filtered colimits of smooth maps.

\bibitem{OSEtale}
The Stacks Project Authors, \emph{The Stacks Project},
\href{https://stacks.math.columbia.edu/tag/07M7}{Tag 07M7}, Lemma~15.9.14:
smooth代数の剰余切断の\'etale持上げ。

\bibitem{OSStacksSplit}
The Stacks Project Authors, \emph{The Stacks Project},
\href{https://stacks.math.columbia.edu/tag/068U}{Tag 068U}, Lemma~15.78.4:
特殊ファイバーのcohomology消滅によるperfect複体の局所分裂。

\bibitem{OSApprox}
The Stacks Project Authors, \emph{The Stacks Project},
\href{https://stacks.math.columbia.edu/tag/07QZ}{Tag 07QZ}, Theorem~16.13.2:
G-ring上の形式解の\'etale近似。本文はそのPopescuと切断による証明法を用いる。

\bibitem{CRSBS}
J.~Burke and G.~Stevenson, \emph{The derived category of a graded Gorenstein ring},
\href{https://arxiv.org/abs/1507.00830}{arXiv:1507.00830v1} (2015),
Proposition~4.18, Lemmas~6.1--6.2 and Theorem~6.4.
PDF pp.~12--14, 17--19を確認。

\bibitem{CRSTV}
B.~To\"en and M.~Vaqui\'e, \emph{Moduli of objects in dg-categories},
Ann. Sci. \'Ecole Norm. Sup. (4) \textbf{40} (2007), 387--444;
\href{https://arxiv.org/abs/math/0503269}{arXiv:math/0503269v5}, Lemma~2.8(2).
PDF pp.~21--22を確認。pseudo-perfectからperfectへの原理を用いる。

\bibitem{CRSPerry}
A.~Perry, \emph{Noncommutative homological projective duality},
Adv. Math. \textbf{350} (2019), 877--972, Lemma~4.15(3).
\href{https://www.sciencedirect.com/science/article/am/pii/S000187081930221X}{著者受理稿}。
admissible半直交成分とsmoothnessの同値を用いる。

\bibitem{CRSStevenson}
G.~Stevenson, \emph{Subcategories of singularity categories via tensor actions},
Compos. Math. \textbf{150} (2014), 229--272;
\href{https://arxiv.org/abs/1105.4698}{arXiv:1105.4698v3}, Theorem~6.13 (PDF p.~26).
今回再確認したのはaffine hypersurfaceの全thick部分圏の分類である。

\bibitem{DKITu}
J. Tu, \emph{Matrix factorizations via Koszul duality},
Compos. Math. \textbf{150} (2014), 1549--1578;
\href{https://arxiv.org/abs/1009.4151}{arXiv:1009.4151}, Appendix~A, Lemma~A.1.

\bibitem{DKIBFK}
M. Ballard, D. Favero and L. Katzarkov,
\emph{A category of kernels for equivariant factorizations and its implications for Hodge theory},
Publ. Math. Inst. Hautes \`Etudes Sci. \textbf{120} (2014), 1--111;
\href{https://doi.org/10.1007/s10240-013-0059-9}{doi:10.1007/s10240-013-0059-9},
Theorem~5.22, Remark~5.21.

\bibitem{DKIDyckerhoff}
T. Dyckerhoff, \emph{Compact generators in categories of matrix factorizations},
Duke Math. J. \textbf{159} (2011), 223--274;
\href{https://arxiv.org/abs/0904.4713}{arXiv:0904.4713}, Section~5.

\bibitem{TWCrainic}
M. Crainic,
\emph{On the perturbation lemma, and deformations},
arXiv:math/0403266.
\url{https://arxiv.org/abs/math/0403266}.
著者公開版\url{https://www.staff.science.uu.nl/~crain101/pert2.pdf}の
Theorem 2.3とRemark 2.4（2026年11月12日本文確認）。
同版の$ip=1+bh+hb$という符号は、本稿の$h=-s$に対応する。

\bibitem{RDPP}
A.~Polishchuk and L.~Positselski,
\emph{Hochschild (co)homology of the second kind I},
Trans. Amer. Math. Soc. \textbf{364} (2012), no.~10, 5311--5368.
\href{https://doi.org/10.1090/S0002-9947-2012-05667-4}{doi:10.1090/S0002-9947-2012-05667-4};
\href{https://arxiv.org/abs/1010.0982v2}{arXiv:1010.0982v2}.
\S3.3では有限射影次元を持つ対象の第二種Extを扱う。
\S3.5の対角比較条件を、本稿の非可換族について確認済みとはしない。

\bibitem{GRToen}
B.~To\"en, \emph{The homotopy theory of dg-categories and derived Morita theory},
Invent. Math. \textbf{167} (2007), 615--667.
\href{https://arxiv.org/abs/math/0408337}{arXiv:math/0408337};
\S2 (arbitrary commutative base ring), Theorems~8.9 and 8.15.

\bibitem{GRKuznetsov}
A.~Kuznetsov, \emph{Hochschild homology and semiorthogonal decompositions},
\href{https://arxiv.org/abs/0904.4330}{arXiv:0904.4330}.
射影核を使うadmissible部分圏のHochschild理論。

\bibitem{GRStacks}
The Stacks Project Authors, \emph{The Stacks Project},
\href{https://stacks.math.columbia.edu/tag/0B91}{Tag 0B91}, Lemma~36.30.4:
flat proper有限表示射のperfect直接像と任意の導来底変更。

\bibitem{GRBFK}
M.~Ballard, D.~Favero and L.~Katzarkov,
\emph{A category of kernels for equivariant factorizations and its implications for Hodge theory},
Publ. Math. Inst. Hautes \`Etudes Sci. \textbf{120} (2014), 1--111.
\href{https://doi.org/10.1007/s10240-013-0059-9}{doi:10.1007/s10240-013-0059-9}, Theorem~5.39.

\bibitem{GRLZ}
X.~Lin and S.~Zhang, \emph{Serre algebra, matrix factorization and categorical Torelli theorem for hypersurfaces},
\href{https://arxiv.org/abs/2310.09927v2}{arXiv:2310.09927v2} (2024),
Theorem~3.5, Remark~3.6 and \S5. 本節で引用するのは体上の計算。

\bibitem{GRLZ26}
X.~Lin and S.~Zhang, \emph{Categorical Torelli and Fourier--Mukai partners of hypersurfaces},
\href{https://arxiv.org/abs/2611.05915v1}{arXiv:2611.05915v1} (2026-11-04), \S2.1.
相対projected kernelの背景として参照し、本研究の非可換generic同値を与えるとはしない。

\bibitem{GRToenFinite}
B.~To\"en, \emph{Finitude homotopique des dg-alg\`ebres propres et lisses},
\href{https://doi.org/10.1112/plms/pdn027}{doi:10.1112/plms/pdn027}.
有限$A_n$決定性はsmooth proper性を仮定する。今回この定理を非可換族へ適用してはいない。

\bibitem{RSRBS} Jesse Burke and Greg Stevenson,
\emph{The derived category of a graded Gorenstein ring},
arXiv:1507.00830v1 (2015).
\url{https://arxiv.org/abs/1507.00830}.
参照：Definition 5.1, Proposition 4.18, Lemmas 6.1--6.2, Theorem 6.4.
\bibitem{RSRToen} Bertrand Toën,
\emph{Anneaux de définition des dg-algèbres propres et lisses},
arXiv:math/0611546v3 (2010).
\url{https://arxiv.org/abs/math/0611546}.
参照：Théorème 1, F4, F6.
\bibitem{RSRTV} Bertrand Toën and Michel Vaquié,
\emph{Moduli of objects in dg-categories},
arXiv:math/0503269v5 (2007).
\url{https://arxiv.org/abs/math/0503269}.
参照：Lemmas 2.8, 2.10, Corollary 2.13, Theorem 3.6, Corollaries 3.15, 3.17, 3.19.
\bibitem{RSROrlov} Dmitri Orlov,
\emph{Smooth and proper noncommutative schemes and gluing of DG categories},
arXiv:1402.7364v5 (2015).
\url{https://arxiv.org/abs/1402.7364}.
参照：Proposition 3.8, Theorem 3.25（体上）.
\bibitem{RSRRepo} nc-cubic project,
\emph{一般次元の局所代数的模型、非楕円型量子化、およびrotation付き可換復元},
commit \texttt{4f2eb328a2249a173f0c658cdb744277f041ff1d}, 2026年11月10日JST.
\url{https://github.com/uedakazushi/nc-cubic}.
一次原稿：\texttt{papers/sections/local\_projective\_model.tex},
\texttt{nonelliptic\_quantization.tex}, \texttt{rotation\_models.tex}.
局所模型の入力、親原稿の形式的復元と、本稿の双加群による局所実現証明を区別する。



\bibitem{P2Kontsevich}
M.~Kontsevich,
Deformation quantization of Poisson manifolds,
Lett. Math. Phys. \textbf{66} (2003), 157--216.
\bibitem{P2LPT}
F.~Loray, J.~V.~Pereira and F.~Touzet,
Foliations with trivial canonical bundle on Fano 3-folds,
Math. Nachr. \textbf{286} (2013), 921--940.
\bibitem{P2Mayanskiy}
E.~Mayanskiy,
Poisson cohomology of two Fano threefolds,
\url{https://arxiv.org/abs/1303.5916}.

\end{thebibliography}
\end{document}
